\documentclass[12pt,UTF8]{ctexbook}

% 导入设定
\input{../shared/settings}
% 代数部分

% 像空间
\newcommand\Img[0]{
    \mathrm{Img\,}
}
% 核：零空间
\newcommand\Ker[0]{
    \mathrm{Ker\,}
}
% 容：像空间的维数
\newcommand\rang[0]{
    \mathrm{rang\,}
}
% 佚：零空间的维数
\newcommand\nul[0]{
    \mathrm{null\,}
}
% 迹
\newcommand\ji[0]{\operatorname{tr}}
% 函数的限制：限制在定义域的某个子集中的函数
\newcommand\restr[2]{{% we make the whole thing an ordinary symbol
  \left.\kern-\nulldelimiterspace % automatically resize the bar with \right
  #1 % the function
  \vphantom{\big|} % pretend it's a little taller at normal size
  \right|_{#2} % this is the delimiter
  }
}
% 随附矩阵、附阵、伴随矩阵
\newcommand\co[0]{\operatorname{Co}}

% 交换集：交换元构成的集合
\newcommand\Cmu[1]{
    \mathrm{Comm\,\left(#1\right)}
}

% 商结构：两个代数结构的商结构
\newcommand\quotient[2]{
    \mathop{\raisebox{0.65ex}{$#1$}$\mkern-5mu\diagup\mkern-7mu$\raisebox{-0.7ex}{$#2$}}
}
% \usepackage{url}

% 封面
% --- 新增：定义作者和书名的变量，方便统一管理 ---
\newcommand{\mytitle}{平直代数}
\newcommand{\myauthor}{大青花鱼}
% \title{\zihao{0} \bfseries 平直代数}
% \author{\zihao{2} \texttt{大青花鱼}}
% \date{\bfseries\today}
% \date{}
% 正文
\begin{document}
% \maketitle
% --- 版权页（扉页背面）开始 ---
% ========== 第1页（右页，即扉页）==========
\cleardoublepage                % 确保从右页（奇数页）开始
\thispagestyle{empty}           % 本页无页码无页眉
\begin{center}
    \vspace*{2cm}               % 顶部留白
    
    % 书名
    {\zihao{0} \bfseries \mytitle}
    
    \vspace{1cm}              % 书名与作者间距
    
    % 作者名
    {\zihao{2} \texttt{\myauthor}}
    
    \vfill                      % 将后续所有内容推至页面底部
    
    % --- 分隔线（可选，增加精致感） ---
    \rule{0.5\textwidth}{0.4pt} \\[0.5cm]
    
    % --- 版权声明（你要求的完整文字，保证显示） ---
    {\scriptsize
    \copyright\ 2026 \texttt{\myauthor} \\[0.2cm]
    本作品采用 知识共享署名-相同方式共享 4.0 国际许可协议（CC BY-SA 4.0）进行许可。\\
    协议全文详见：https://creativecommons.org/licenses/by-sa/4.0/ \\
    您可以自由地共享、修改本作品，但须遵守“署名”和“相同方式共享”的条件。
    }
    
    % \vspace{1.5cm}              % 底部留白，不贴边
\end{center}
\newpage   
% ========== 扉页结束 ==========

\tableofcontents
\newpage

\chapter{平直空间}
\section{平直空间的定义}
这本书是关于“平直空间”的。什么是平直空间呢？
让我们从“正比例关系”讲起。

正比例关系是我们在生活中常常接触到的一种关系。
如果两个量成正比，一个量增加若干倍，另一个量也随之增加若干倍。

举例来说，一只兔子四条腿，两只兔子就有八条腿，以此类推。

另外，如果张三有三只兔子，李四有五只兔子，那么兔腿数量可以按照各自的算（$3\times4 +5\times4$），
也可以先把兔子数量加起来再算（$(3+5)\times4$）。

这种性质，反映到函数的图像上，是一条直线。
我们把这种性质称为\textbf{直线性}，简称\textbf{直性}。

有直性的映射，称为\textbf{平直映射}。有直性的代数对象，就是\textbf{平直代数}。
本书主要介绍的就是平直空间上的平直映射的基本性质。这些概念是点和直线、平面、空间、平移、对称和旋转等概念的自然推广。

本章我们首先规定一般意义上的平直空间的概念。

研究平面或立体空间中的对象时，我们通过在其中建立坐标轴，把平面或立体空间中的形体对象转化为代数对象。
比如，平面上的一条直线，可以用一个代数式表示：
$$ y = 2x + 1$$

平面或空间中的点，我们用坐标表示，或者，用一个从原点指向该点的向量表示。
每个向量，都对应着空间中的一点。这样，我们可以描述组成空间的任何元素。

从这个概念出发，我们引入平直空间的概念：

\begin{df}\label{df:1_1}
给定域\footnote{域的定义参见定义\ref{df:0_3}。}$\mathbb{F}$，$\mathbb{F}$上的\textbf{平直空间}$\mathbb{V}$是这样一个集合:

首先，$\mathbb{V}$上定义了两种二元运算：$\diamond$和$\bullet$。其中：
\begin{itemize}
    \item $\diamond$运算将两个$\mathbb{V}$中的元素映射为$\mathbb{V}$中的另一个元素。
    \item $\bullet$运算将一个$\mathbb{F}$中的元素和一个$\mathbb{V}$中的元素映射为$\mathbb{V}$中另一个元素。
\end{itemize}
其次，$\mathbb{F}$、$\mathbb{V}$和以上两种运算具有以下性质：
\begin{enumerate}
    \item $\mathbb{V}$中元素对于$\diamond$运算满足交换律：$\forall \,\, \mathbf{u}, \, \mathbf{v} \, \in \mathbb{V}, \,\, \mathbf{u} \diamond \mathbf{v} = \mathbf{v} \diamond \mathbf{u}.$
    \item $\mathbb{V}$中元素对于$\diamond$运算满足结合律：$\forall \,\, \mathbf{u}, \, \mathbf{v}, \, \mathbf{w} \, \in \mathbb{V}, \,\, (\mathbf{u} \diamond \mathbf{v}) \diamond \mathbf{w} = \mathbf{u} \diamond (\mathbf{v} \diamond \mathbf{w}).$
    \item $\mathbb{V}$中元素对于$\diamond$运算存在零元素：$\exists \,\, \mathbf{0} \, \in \mathbb{V}$ 使得 $\forall \,\, \mathbf{u} \, \in \mathbb{V}, \,\, \mathbf{u} \diamond \mathbf{0} = \mathbf{u}.$
    \item $\mathbb{V}$中元素对于$\diamond$运算存在逆元素：$\forall \,\, \mathbf{u} \, \in \mathbb{V}, \,\,\, \exists \,\, \mathbf{v} \, \in \mathbb{V}$ 使得 $\mathbf{u} \diamond \mathbf{v} = \mathbf{0}.$
    \item $\mathbb{F}$中乘法$\cdot$和$\bullet$运算兼容：$\forall \,\, a, \, b \, \in \mathbb{F}, \,\,\, \forall \,\, \mathbf{u} \, \in \mathbb{V}, \,\, (a \cdot b)\bullet\mathbf{u} = a \bullet (b\bullet\mathbf{u}).$
    \item $\mathbb{F}$中有$\bullet$运算的空元$1_{\mathbb{V}}$：$\forall \,\, \mathbf{u} \, \in \mathbb{V}, \,\,\, 1_{\mathbb{V}} \bullet \mathbf{u} = \mathbf{u}. $
    \item $\mathbb{F}$中乘法$\cdot$对$\diamond$运算满足分配律：$\forall \,\, \mathbf{u}, \, \mathbf{v} \, \in \mathbb{V}, \,\,\, \forall \,\, a \, \in \mathbb{F}, \,\,\, a\cdot(\mathbf{u} \diamond \mathbf{v}) = (a \cdot \mathbf{u}) \diamond (a \cdot \mathbf{v}). $
    \item $\bullet$运算对$\mathbb{F}$中加法$+$也满足分配律：$\forall \,\, a, \, b \, \in \mathbb{F}, \,\,\, \forall \,\, \mathbf{u} \, \in \mathbb{V}, \,\,\, (a + b)\bullet\mathbf{u} = (a \bullet \mathbf{u}) + (b \diamond \mathbf{u}). $
\end{enumerate}
\end{df}
以上的条件看起来非常抽象。我们用一个具体的例子来解释一下。
平面形学中，我们处理的对象如直线、线段、抛物线等等都在平面上。我们用$(x,y)$的方式记录平面上的点，其中$x$和$y$是实数。
我们考虑所有形同
$$ \mathbf{v} = (x, y)$$
的对象构成的集合：$\{(x, y) \,\,|\,\, x,y\,\in\,\mathbb{R}\}.$ 把它们看作向量。定义向量的$+$运算（向量的加法）：
$$ (x_1, y_1) + (x_2, y_2) = (x_1+x_2, y_1+y_2).$$
定义实数$a$和向量$ \mathbf{v} = (x, y)$的$\cdot$运算（称为向量的数乘）：
$$ a \cdot (x, y) = (ax, ay).$$
我们把向量的加法$+$和向量的数乘$\cdot$作为以上提到的两个二元运算，对上面提到的8个性质逐个验证一下：
\begin{enumerate}
    \item 向量对于向量的加法$+$运算满足交换律：$\forall \,\, \mathbf{u}, \, \mathbf{v} \, \in \mathbb{R}^2, \,\, \mathbf{u} + \mathbf{v} = \mathbf{v} + \mathbf{u}.$
    \item 向量对于向量的加法$+$运算满足结合律：$\forall \,\, \mathbf{u}, \, \mathbf{v}, \, \mathbf{w} \, \in \mathbb{R}^2, \,\, (\mathbf{u} + \mathbf{v}) + \mathbf{w} = \mathbf{u} + (\mathbf{v} + \mathbf{w}).$
    \item 向量加法有零元素，就是零向量：$\mathbf{0} = (0, 0)$：$\forall \,\, \mathbf{u} \, \in \mathbb{R}^2, \,\, \mathbf{u} + \mathbf{0} = \mathbf{u}.$
    \item 向量加法有逆元素，就是逆向量：$\forall \,\, \mathbf{u} = (x, y) \, \in \mathbb{R}^2, \,\,\, \exists \,\, \mathbf{v} = (-x, -y)\, \in \mathbb{R}^2$ 使得 $\mathbf{u} + \mathbf{v} = (0, 0) = \mathbf{0}.$
    \item 实数乘法$\cdot$和数乘运算兼容：$\forall \,\, a, \, b \, \in \mathbb{R}, \,\,\, \forall \,\, \mathbf{u} = (x, y) \, \in \mathbb{R}^2, \,\, (ab)\mathbf{u} = (abx, aby)= a(b\mathbf{u}).$
    \item 实数1就是数乘运算的幺元：$\forall \,\, \mathbf{u} \, \in \mathbb{R}^2, \,\,\, 1\cdot \mathbf{u} = \mathbf{u}. $
    \item 实数乘法向量的加法运算满足分配律：$\forall \,\, \mathbf{u}, \, \mathbf{v} \, \in \mathbb{R}^2, \,\,\, \forall \,\, a \, \in \mathbf{R}, \,\,\, a(\mathbf{u} + \mathbf{v}) = (a \mathbf{u}) + (a \mathbf{v}). $
    \item 数乘对实数加法$+$满足分配律：$\forall \,\, a, \, b \, \in \mathbb{R}^2, \,\,\, \forall \,\, \mathbf{u} \, \in \mathbf{R}^2, \,\,\, (a + b)\mathbf{u} = (a \mathbf{u}) + (b \mathbf{u}). $
\end{enumerate}

我们可以看到，平直空间的概念，就是从研究平面、立体空间衍生而来。

以后，在平直空间中，我们也会沿用研究平面和立体空间时的术语和记号。
我们把$\mathbb{V}$中的元素称为\textbf{向量}或\textbf{点}，把$\mathbb{F}$中的元素称为\textbf{标量}。
为了和标量区别，代表向量的变量一般用粗体表示。

我们把$\star$和$\circ$运算叫做向量的\textbf{加法}运算和\textbf{数乘}运算（也称为\textbf{标量乘法}运算），在不至于混淆的情况下（也就是一般情况下），
记为$+$和$\cdot$（和标量的加法、乘法一样，但一般情况下，我们都能区别出来）。
这两种运算称为平直空间配备的运算。

我们也把$\mathbb{V}$中的零元素称为\textbf{零向量}，向量加法的逆元素称为\textbf{逆向量}（$\mathbf{u}$的逆向量记为$-\mathbf{u}$）。
我们一般把数乘运算的幺元$1_\mathbb{V}$设为$\mathbb{F}$乘法的幺元$1_\mathbb{F}$（可以思考一下，如果两者不一样的话，会有什么影响\footnote{提示：$u = 1_\mathbb{V}\mathbf{u} = (1_\mathbb{F}1_\mathbb{V})\mathbf{u} = 1_\mathbb{F}(1_\mathbb{V} \mathbf{u}) = 1_\mathbb{F}\mathbf{u}$，使用性质4和性质8，研究$-\mathbf{u}.$}）。

$\mathbb{F}$一般称作$\mathbb{V}$的系数域。注意到$\mathbb{V}$是从$\mathbb{F}$衍生并定义的，
所以当我们需要仔细说明平直空间的性质的时候，我们会强调$\mathbb{V}$是域$\mathbb{F}$上的平直空间。
比如，我们刚才讨论的$\mathbb{R}^2$，是$\mathbb{R}$上的平直空间。
很多情况下，我们关心的性质和$\mathbb{F}$没有太大关系。这时候我们一般不会特意强调$\mathbb{F}$。

下面我们来看几个很容易想到的例子：

首先，只有零向量的空间$\{\mathbf{0}\}$，可以是任何域$\mathbb{F}$上的空间。
但这个空间没有什么有趣的性质。我们一般叫它平凡的平直空间。

给定域$\mathbb{F}$和正整数$n$，$\mathbb{F}$中$n$个元素构成的有序$n$元组的集合$\mathbb{F}^n$是$\mathbb{F}$上的平直空间。
$\mathbb{F}^n$中的向量就是$n$元组$v = (a_1, a_2, \cdots, a_n)$，向量加法就是元组中的每个元素（称为分量）分别对应相加，而数乘就是标量和向量的每个分量对应相乘。
容易验证，它满足平直空间的性质。

当$n=1$的时候，我们发现，$\mathbb{F}^n$就是$\mathbb{F}$自己。所以，任何域都是自身的平直空间。

$n=2,3$的时候，$\mathbb{F}^n$分别是我们熟悉的平面和立体空间。

另一个典型的平直空间，是$\mathbb{F}$中元素构成的（无穷）数列：
$$ a = (a_0, a_1, a_2, \cdots ) $$
构成的空间$\mathbb{F}^\mathbb{N}$。
$\mathbb{F}^\mathbb{N}$的向量加法，就是数列中下标相同的元素相应相加，数乘运算就是用标量乘以数组的每个元素。
容易验证，$\mathbb{F}^\mathbb{N}$也是一个$\mathbb{F}$上的平直空间，满足上述8个性质。

再举一个例子。考虑以$\mathbb{F}$中元素为系数的一元整式\footnote{具体定义见第一章最后一节“整式”。}：
$$ p(z) = a_0 + a_1 z + \cdots + a_m z^m, \,\, m \in \mathbb{N} $$
一元整式集合$F[z]$对于整式的加法和系数乘以整式的乘法，也构成平直空间。

进一步可以思考，所有$k$元整式、所有$m$次整式、所有$k$元$m$次整式、所有整式数组、所有整式数列，是否也构成$\mathbb{F}$上的平直空间？

\section{平直空间的性质}
平直空间是平面和立体空间的自然推广。
它包括平面和立体空间，但不仅仅如此。
我们现在来看一些平面和立体空间具备的基本性质，是否任何平直空间都具备。

我们知道，平面和立体空间都只有一个原点。对一般的平直空间，是否也是如此呢？

\begin{pp}\label{pp:1_1}
同一个平直空间中，零向量是唯一的。
\end{pp} 
\begin{proof} % [{\bfseries 证明\ref{pp:1_1}}]
    设平直空间里有两个零向量$\mathbf{0}$和$\mathbf{0}'$，则由性质3可知：
    $$ \mathbf{0} = \mathbf{0} + \mathbf{0}' = \mathbf{0}'.$$
    第一个等号成立，是因为$\mathbf{0}'$是零向量。第二个等号成立，是因为$\mathbf{0}$是零向量。
    两个零向量相等，说明是同一个零向量。
    因此，同一个平直空间里，只有一个零向量。零向量是唯一的。
\end{proof}
 

平面中，一个向量的逆向量是和它长度相等，方向相反的向量。这样的向量有且仅有一个。
对于一般的平直空间，是否也是如此呢？
\begin{pp}\label{pp:1_2}
同一个平直空间中，每个向量有唯一的逆向量。
\end{pp} 
\begin{proof} % [{\bfseries 证明\ref{pp:1_2}}]
    设$\mathbf{u} \, \in \mathbb{V}$有两个逆向量$\mathbf{v}$和$\mathbf{w}$，则由性质2可知：
    $$ \mathbf{v} = \mathbf{v} + \mathbf{0} = \mathbf{v} + (\mathbf{u}+\mathbf{w}) = (\mathbf{v}+\mathbf{u}) + \mathbf{w} = \mathbf{w}.$$
    两个逆向量是同一个向量。
    因此，同一个平直空间里，每个向量只有一个逆向量。
\end{proof}
 
逆向量是唯一的，我们一般将$\mathbf{u}$的逆向量记作$-\mathbf{u}$，
减去一个向量$\mathbf{u}$，就是加上它的逆向量$-\mathbf{u}$.

为了好说话，我们在讨论平直空间的一般性质的时候，
都假设我们讨论的是一个给定的系数域$\mathbb{F}$和$\mathbb{F}$上的平直空间$\mathbb{V}$。
空间中的向量一般用$\mathbf{u}$、$\mathbf{v}$、$\mathbf{w}$等表示。
如无特殊需要，记号就不再重复声明了。

平面中，零向量就是任何向量长度缩减为0的结果，或者说0倍的向量就是零向量。
一般的平直空间里，也有类似的结论：
\begin{pp}\label{pp:1_3}
    $\forall \mathbf{u} \, \in \mathbb{V}, \,\,\, 0 \mathbf{u} = \mathbf{0}.$
\end{pp} 
\begin{proof} % [{\bfseries 证明\ref{pp:1_3}}]
    $$ 0\mathbf{u} = (0 + 0)\mathbf{u} = 0\mathbf{u} + 0\mathbf{u}.$$
    两边各加上$0\mathbf{u}$的逆向量，左边得到零向量，右边得到$0\mathbf{u}$. 所以
    $$ 0\mathbf{u} = \mathbf{0}.$$
\end{proof}
 

上一个命题说明了系数域上的标量0对一般向量的作用。
零向量对系数域上一般标量的作用也类似：
\begin{pp}\label{pp:1_4}
    $\forall a \, \in \mathbb{F}, \,\,\, a \mathbf{0} = \mathbf{0}.$
\end{pp} 
\begin{proof} % [{\bfseries 证明\ref{pp:1_4}}]
    $$ a\mathbf{0} = a(\mathbf{0}+\mathbf{0}) = a\mathbf{0}+a\mathbf{0}.$$
    两边各加上$a\mathbf{0}$的逆向量，左边得到零向量，右边得到$a\mathbf{0}$. 所以
    $$ a\mathbf{0} = \mathbf{0}.$$
\end{proof}
  

平面和立体空间中，求向量的逆向量，只需要把它的坐标加个负号，也就是乘上$-1$。
一般的平直空间里，也可以这么做：
\begin{pp}\label{pp:1_5}
    $\forall \mathbf{u} \, \in \mathbb{V}, \,\,\, (-1) \mathbf{u} = -\mathbf{u}.$
\end{pp} 
\begin{proof} % [{\bfseries 证明\ref{pp:1_5}}]
    $$ 0 = 0\mathbf{u} = (1 + (-1))\mathbf{u} = 1\mathbf{u} + (-1)\mathbf{u} = \mathbf{u} + (-1)\mathbf{u}.$$
    逆向量是唯一的，所以
    $$ (-1)\mathbf{u} = -\mathbf{u}.$$
\end{proof}
 

\section{子空间}
考虑平面$\mathbb{R}^2$内的直线：$y = 2x.$ 它也可以看作形同$(t, 2t)$的点的集合：
$$ \mathcal{D}_0 = \{ (t, 2t), \, t \in \mathbb{R}\}$$
这个集合是否是一个平直空间呢？
可以验证，对于平直空间$\mathbb{R}^2$配备的向量加法和数乘运算来说，$\mathcal{D}_0$满足平直空间的性质。零元素是点$(0,0)$.

我们再来看$\mathbb{R}^2$内的另一条直线：$y = 2x + 1$. 它也可以看作形同$(t, 2t + 1)$的点的集合：
$$ \mathcal{D}_1 = \{ (t, 2t + 1), \, t \in \mathbb{R}\}$$
这个集合是否是一个平直空间呢？
我们发现，对于平直空间$\mathbb{R}^2$配备的向量加法来说，直线$\mathcal{D}_1$里没有零向量。
其次，我们检查直线上两点$(0,1)$和$(1,3)$，它们的和是$(2,4)$，不在直线上。
所以，$\mathcal{D}_1$不是平直空间。

我们发现，平直空间的子集，有些是平直空间，有些不是。我们引入子空间的概念：
\begin{df}\label{df:1_2}
    若平直空间$\mathbb{V}$的子集$\mathbb{U}$（关于$\mathbb{V}$配备的运算）也是平直空间，则称$\mathbb{U}$是$\mathbb{V}$的{\bfseries 子空间}。
\end{df}

讨论某个平直空间的子集是否是其子空间的时候，我们默认其配备的是同样的运算。
平直空间的子集对于该平直空间配备的运算不构成子空间，但对于其他的运算构成子空间的情况是否存在呢？我们暂且不讨论。

作为最简单的平凡例子：$\{\mathbf{0}\}$是任何空间的子空间。任何空间都是自身的子空间。不等于整个空间的子空间，称为其{\bfseries 真子空间}（这个定义和真子集类似）

怎样判断平直空间间$\mathbb{V}$的子集$\mathbb{U}$是否是它的子空间呢？显然，直接对此子集验证平直空间的所有性质，必定能得出结论。
不过，我们已经掌握了这个集合的一定信息（这个集合是某个平直空间的子集）。是否有更方便的方法判断呢？

我们来看成为平直空间需满足的8个性质。性质1和性质2是“自闭合”的。
也就是说，是元素之间的关系，不涉及新元素。
所以，$\mathbb{U}$中的元素，作为$\mathbb{V}$中元素，必然满足性质1和性质2.

性质3和性质4涉及到新元素的引入，所以是必须的。
$\mathbb{U}$中需要有零向量，而且$\mathbb{U}$中每个向量的逆向量也应该在$\mathbb{U}$里面。

$\mathbb{U}$中的零向量，是否就是$\mathbb{V}$的零元素？可以是另一个向量吗？

我们知道，$\mathbb{U}$中任何向量和$\mathbb{V}$的零向量相加，都等于它自己。
所以，根据零向量的唯一性，$\mathbb{U}$的零向量，就是$\mathbb{}$的零向量。

同理，根据逆向量的唯一性可知，$\mathbb{U}$中向量的逆向量，就是它在$\mathbb{V}$中的逆向量。

性质5至性质8，是系数域$\mathbb{F}$中的标量和$\mathbb{U}$中的向量的关系。
我们发现，只要性质里涉及的对象在$\mathbb{U}$里面，相关的性质必然是成立的。
反之，如果涉及的对象不在$\mathbb{U}$里面，相关的性质就无法成立。

比如，对某个向量$\mathbf{u} \in \mathbb{U}$和标量$a, b$来说，$(ab)\mathbf{u} = a(b\mathbf{u})$无疑是成立的。
这是平直空间$\mathbb{V}$的性质。
所以，只要$(ab)\mathbf{u}$、$b\mathbf{u}$和$a(b\mathbf{u})$都在$\mathbb{U}$里面。我们就可以说性质5对$\mathbb{U}$成立。
反之则不成立。

综上讨论，我们发现，判定平直空间的子集是否是子空间时，我们可以“缩减手续”：

\begin{pp}\label{pp:1_6}
    若平直空间$\mathbb{V}$的子集$\mathbb{U}$满足以下条件，则$\mathbb{U}$是$\mathbb{V}$的子空间。
    \begin{enumerate}
        \item $\mathbb{V}$的零向量属于$\mathbb{U}$
        \item $\mathbb{U}$对向量加法封闭：$\forall \,\, \mathbf{u}, \, \mathbf{v}, \, \in \mathbb{U}$，其和$\mathbf{u} + \mathbf{v} \, \in \mathbb{U}$
        \item $\mathbb{U}$对数乘运算封闭：$\forall \,\, a \, \in \mathbb{F}, \,\, \mathbf{u}, \, \in \mathbb{U}, \,\,\, a\mathbf{u} \, \in \mathbb{U}$
    \end{enumerate}
\end{pp}

注意其中没有明确要求逆向量存在。
不过，向量$\mathbf{u}$的逆向量就是$(-1)\mathbf{u}$，所以第3个条件中其实已经包含了这一点。

我们可以初步从直观的角度理解子空间的性质。子空间要求向量对数乘运算封闭，也就是说，如果一个向量属于子空间，那么该向量放缩之后，仍然在子空间里。
此外，子空间对向量加法封闭。如果空间中有两个不共线的向量，那么以它们为基础“坐标轴”，就能衍生出一个“平面”。
举一个最简单的例子：假设子空间里有向右和向上的向量，那么经过适当的放缩后加起来，就能衍生出上下左右方向的向量。
但如果子空间里没有向前或向后的向量，那么我们永远无法“拼出”前后方向的向量。

我们来看一些常见的例子。前面我们研究了平面内的直线。
那么，以下的集合是否是子空间呢？
\begin{enumerate}
    \item $\mathbb{V} = \mathbb{R}^3$，$\mathbb{U} = \{(x_1, x_2, 0) \,\, | \,\, x_1, x_2 \, \in \mathbb{R}\}$.
    \item $\mathbb{V} = \mathbb{R}^2$，$\mathcal{D}_{a,b} = \{(x, y) \,\, | \,\, x, y \, \in \mathbb{R}, \,\, y = ax + b\}$.
    \item $\mathbb{V} = \mathbb{R}[z]$，$\mathbb{R}_m[z] = \{a_0 + a_1 z + \cdots + a_m z^m \,\, | \,\, a_0,\, a_1,\, \cdots ,\, a_m \, \in \mathbb{R}\}$.
    \item $\mathbb{V} = \mathbb{R}[z]$，$\mathcal{I}_n (\mathbb{R}) = \{p(z) \,\, | \,\, p(1) = p(2) = \cdots = p(n) = 0 \}$.
\end{enumerate}

经过以上研究，我们发现，平面$\mathbb{R}^2$内的一条直线$\mathcal{D}_{a,b}$是子空间，当且仅当$b = 0$。
也就是说，过原点的直线，才是$\mathbb{R}^2$的子空间。
立体空间$\mathbb{R}^3$内的平面和直线，也仅在过原点的时候，成为它的子空间。

我们可以猜测，$\mathbb{R}^2$的所有子空间，就是平凡空间$\{\mathbf{0}\}$、$\mathbb{R}^2$自身，以及所有过原点的直线。
$\mathbb{R}^3$的所有子空间，就是平凡空间$\{\mathbf{0}\}$、$\mathbb{R}^3$自身，以及所有过原点的直线和平面。
我们已经证明了这些对象都是子空间。能否证明，其他子集都不是子空间呢？这需要更多的工具来分析平直空间的性质。

\section{和与直和}
集合论中，我们通过讨论集合的并集和交集，了解集合之间的关系。我们通过把集合划分为不同的不相交子集的并集，来分类讨论集合元素的性质。同样地，我们希望能将平直空间划分为更“小”的空间的“并”。
这种划分是否可能呢？

平直空间和子空间都是集合。我们可以从集合的基本性质来研究。

平直空间$\mathbb{V}$的两个子空间：$\mathbb{U}_1$和$\mathbb{U}_2$，它们的交集是否是子空间呢？
答案是肯定的。我们可以直接验证。

$\mathbb{U}_1$和$\mathbb{U}_2$的并集是否是子空间呢？不一定。我们在后面可以看到，一般来说，不是。
举例来说，直线 $y = 3x$ 和 $y = 2x$ 都是平面的子空间。但它们的并集不是子空间。
比如，向量$(1,3)$和$(1,2)$分别在两条直线上，因此在并集里。但它们的和$(2,5)$不属于任何一条直线。

我们发现，集合的并集无法满足子空间的第2个判断标准：对向量加法封闭。
从上面的例子里可以发现，并集要成为子空间的话，需要“更多的”向量。
比如，如果我们往并集里添加向量，让并集成为子空间。既然向量$(1,3)$和$(1,2)$在里面，那么$(1,3)+(1,2)$、$(1,3)+2*(1,2)$、$2*(1,3)+(1,2)$等向量也应该在里面。

为此，对于平直空间，我们要定义比集合的并更“大”的操作。
\begin{df}\label{df:1_3}
    平直空间$\mathbb{V}$的子空间$\mathbb{U}_1, \mathbb{U}_2, \cdots , \mathbb{U}_m$的\textbf{和}为以下集合：\\
    $$ \mathbb{U}_1 + \mathbb{U}_2 + \cdots +\mathbb{U}_m = \{\mathbf{u}_1+\mathbf{u}_2+\cdots+\mathbf{u}_m \,\, | \,\, \mathbf{u}_1 \in \mathbb{U}_1, \mathbf{u}_2 \in \mathbb{U}_2, \cdots , \mathbf{u}_m \in \mathbb{U}_m \}$$
\end{df}

子空间的和是将子空间里所有向量分别相加。
显然，子空间的并集是子空间之和的子集。子空间之和是比并集更“大”的集合。

\begin{pp}\label{pp:1_7}
    平直空间$\mathbb{V}$的子空间$\mathbb{U}_1, \mathbb{U}_2, \cdots , \mathbb{U}_m$的和仍然是子空间。
\end{pp}
\begin{proof} % [{\bfseries 证明\ref{pp:1_7}}]
我们使用归纳法。$m=2$时，可以直接验证。\\
假设对某个$m \ge 2$，$\mathbb{U}_1, \mathbb{U}_2, \cdots , \mathbb{U}_m$的和为子空间$\mathbb{W}$。
那么，对于$\mathbb{U}_1, \mathbb{U}_2, \cdots , \mathbb{U}_{m+1}$，其和为：
$$ \mathbb{W}' = \mathbb{U}_1 + \mathbb{U}_2 + \cdots +\mathbb{U}_{m+1} = \{\mathbf{u}_1+\mathbf{u}_2+\cdots+\mathbf{u}_m+\mathbf{u}_{m+1} \,\, | \,\, \mathbf{u}_1 \in \mathbb{U}_1, \mathbf{u}_2 \in \mathbb{U}_2, \cdots , \mathbf{u}_{m+1} \in \mathbb{U}_{m+1} \}$$
我们可以证明，$\mathbb{W}'$就是$\mathbb{W}$和$\mathbb{U}_{m+1}$的和。\\
在$\mathbb{W}'$中任取一个向量$\mathbf{w}'$，按照定义，
$$\mathbf{w}' = \mathbf{u}_1+\mathbf{u}_2+\cdots+\mathbf{u}_m+\mathbf{u}_{m+1}$$
而等式右边前$m$项之和按照定义是$\mathbb{W}$中向量，所以$\mathbf{w}'$是$\mathbb{W}$中某向量和$\mathbb{U}_{m+1}$某向量之和。
$$\mathbb{W}' \subseteq \mathbb{W} + \mathbb{U}_{m+1}$$
另一方面，$\mathbb{W} + \mathbb{U}_{m+1}$中任何元素都能写成$\mathbf{u}_1+\mathbf{u}_2+\cdots+\mathbf{u}_m+\mathbf{u}_{m+1}$的形式，因此属于$\mathbb{W}'$。所以：
$$\mathbb{W} + \mathbb{U}_{m+1} \subseteq \mathbb{W}'.$$
这说明两者相等。我们回到$m=2$的情况：$\mathbb{W}' = \mathbb{W} + \mathbb{U}_{m+1}$是子空间。\\
综上可知命题对任意$m \ge 2$成立。若干个子空间的和必然是子空间。\\
\end{proof}


几个子空间的和包含了相加的每个子空间。
不仅如此，如果某个平直空间包含了若干个子空间，那么必然也包含它们的和。
所以子空间的和不仅是子空间，还是包含它们的子空间中“最小”的空间。

让我们来看几个例子。

直线 $\mathcal{D}_1: \,\, y = 3x$ 和 $\mathcal{D}_2: \,\, y = 2x$ 都是
平面的子空间。
$$\mathcal{D}_1 + \mathcal{D}_2 = \{(t+s, 3t+2s) \,\, | \,\, t \in \mathcal{D}_1, \,\, s \in \mathcal{D}_2 \}$$

对任意$(x,y)$，$t = y - 2x, s = 3x - y$使得$(t+s, 3t+2s) = (x, y)$。所以，平面的任何向量都在$\mathcal{D}_1 + \mathcal{D}_2$里面。所以，$\mathcal{D}_1 + \mathcal{D}_2$就是平面$\mathbb{R}^2$。

立体空间$\mathbb{R}^3$中，直线 $\mathcal{D}: \,\, \{(t,-2t,t) \,\, |\,\, t\in\mathbb{R}\}$ 和平面 $\mathcal{P} \,\, : \,\, x + y - z = 0$ 都是平面的子空间。

对任意$(x,y,z) \in \mathbb{R}^3$，我们有
$$(x,y,z) = \left(\frac{z-x-y}{2},x+y-z,\frac{z-x-y}{2}\right) + \left(\frac{3x+y-z}{2},z-x,\frac{x+y+z}{2}\right).$$
而且$(\frac{z-x-y}{2},x+y-z,\frac{z-x-y}{2})\,\in\,\mathcal{D}$，$(\frac{3x+y-z}{2},z-x,\frac{x+y+z}{2})\,\in\,\mathcal{P}$。
所以$\mathcal{D}+\mathcal{P}$就是$\mathbb{R}^3$。\footnote{思考：是否有其他的写法？}

如果把直线换成$\mathcal{D}$换成平面$\mathcal{P}_2 : x = 0$，对任意$(x,y,z) \in \mathbb{R}^3$，我们有：
$$(x,y,z) = \left(0,y+x-z,0\right) + \left(x,z-x,z\right)$$
或者：
$$(x,y,z) = \left(0,y+x,z\right) + \left(x,-x,0\right).$$
等式右侧第一项为$\mathcal{P}_2$中元素，第二项为$\mathcal{P}$中元素。所以$\mathcal{P}+\mathcal{P}_2$就是$\mathbb{R}^3$。

我们发现，有时候平直空间每个元素都可以写成子空间元素之和（我们把这件事简称为向量的“分解”）。这时候，平直空间可以表示成自身的子空间的和。有时候，写法不止一种。
我们再给平直空间和它的子空间定义一种新的关系：

\begin{df}\label{df:1_4}
    若平直空间$\mathbb{V}$中任意元素$\mathbf{u}$都能以唯一的方式写成其子空间$\mathbb{U}_1, \mathbb{U}_2, \cdots , \mathbb{U}_m$中元素的和：
    $$ \mathbf{u} = \mathbf{u}_1+\mathbf{u}_2+\cdots+\mathbf{u}_m ,\,\,\, \mathbf{u}_1 \in \mathbb{U}_1, \mathbf{u}_2 \in \mathbb{U}_2, \cdots , \mathbf{u}_m \in \mathbb{U}_m $$
    则称平直空间$\mathbb{V}$可表示为其子空间$\mathbb{U}_1, \mathbb{U}_2, \cdots , \mathbb{U}_m$的\textbf{直和}，记作：
    $$ \mathbb{V} = \mathbb{U}_1 \oplus \mathbb{U}_2 \oplus \cdots \oplus \mathbb{U}_m.$$
\end{df}

需要注意的是，直和并不是一种运算，而是一种关系的表示形式。当我们说某个空间可以表示为其子空间的直和时，
并不说明我们在对涉及的子空间做某种运算。

我们来看之前的例子：

对平面中任意一点$(x,y)\in\mathbb{R}^2$，可以验证$t = y - 2x, s = 3x - y$可以将其写成$\mathcal{D}_1$和$\mathcal{D}_2$中向量之和。
所以，直线$\mathcal{D}_1$和$\mathcal{D}_2$之和为$\mathbb{R}^2$。

另一方面，如果$(x,y)$能以两种形式写成$\mathcal{D}_1$和$\mathcal{D}_2$中向量之和：
$$ (x, y) = (t_1, 2t_1) + (s_1, 3s_1) = (t_2, 2t_2) + (s_2, 3s_2).$$
那么，
\begin{align*}
    t_1 + s_1 &= t_2 + s_2  \\
    2t_1 + 3s_1 &= 2t_2 + 3s_2 
\end{align*}
把等式右侧看作常数，解这个二元一次方程组，可得：
$$ t_1 = t_2, \quad s_1 = s_2.$$
两种形式其实是同一种。
所以，$\mathbb{R}^2 = \mathcal{D}_1 \oplus \mathcal{D}_2.$

证明子空间的和是直和，要证明每个元素的表达方式都是唯一的。
是否有更简单方便的方法呢？我们可以把验证工作缩减到一个特殊的元素：零向量上。
\begin{pp}\label{pp:1_8}
    设空间$\mathbb{V} = \mathbb{U}_1 + \mathbb{U}_2 + \cdots + \mathbb{U}_m$，则$\mathbb{V} = \mathbb{U}_1 \oplus \mathbb{U}_2 \oplus \cdots \oplus \mathbb{U}_m$当且仅当零向量只能表示为零向量之和，即：

    若零向量$\mathbf{0}$能表示为
    $$ \mathbf{0} = \mathbf{u}_1 + \mathbf{u}_2 + \cdots + \mathbf{u}_m,\,\,\, \mathbf{u}_1\in \mathbb{U}_1,\,\,\mathbf{u}_2\in \mathbb{U}_2,\,\,\cdots,\,\,\mathbf{u}_m\in \mathbb{U}_m, $$
    
    则$\mathbf{u}_1 = \mathbf{u}_2 = \cdots = \mathbf{u}_m = \mathbf{0}.$
\end{pp}
\begin{proof} % [{\bfseries 证明\ref{pp:1_8}}]
    首先，如果$\mathbb{V} = \mathbb{U}_1 \oplus \mathbb{U}_2 \oplus \cdots \oplus \mathbb{U}_m$，
    那么$\mathbb{V}$中每个元素都以唯一方式写成各个子空间中元素的和，
    于是零向量也以唯一方式写成各个子空间元素的和。
    但我们已经知道一种平凡的写法：零向量是每个子空间中都有的元素，所以：
    $$ \mathbf{0} = \mathbf{0} + \mathbf{0} + \cdots + \mathbf{0},\,\,\, \mathbf{0}\in \mathbb{U}_1,\,\,\mathbf{0}\in \mathbb{U}_2,\,\,\cdots,\,\,\mathbf{0}\in \mathbb{U}_m, $$
    因此这就是零向量唯一的写法。也就是说，零向量只能表示成零向量之和。\\
    再证明另一个方向：如果零向量只能表示成零向量之和，那么，考察任一向量$\mathbf{u}$。\\
    $\mathbb{V} = \mathbb{U}_1 + \mathbb{U}_2 + \cdots + \mathbb{U}_m$，所以必然存在$\mathbf{u}_1\in \mathbb{U}_1,\,\,\mathbf{u}_2\in \mathbb{U}_2,\,\,\cdots,\,\,\mathbf{u}_m\in \mathbb{U}_m$使得
    $$ \mathbf{u} = \mathbf{u}_1 + \mathbf{u}_2 + \cdots + \mathbf{u}_m.$$
    而如果有另一组向量：
    $$\mathbf{v}_1\in \mathbb{U}_1,\,\,\mathbf{v}_2\in \mathbb{U}_2,\,\,\cdots,\,\,\mathbf{v}_m\in \mathbb{U}_m, $$
    使得
    \begin{align*} 
        \mathbf{u} &= \mathbf{u}_1 + \mathbf{u}_2 + \cdots + \mathbf{u}_m  \\
        \mathbf{u} &= \mathbf{v}_1 + \mathbf{v}_2 + \cdots + \mathbf{v}_m, 
    \end{align*}
    那么，上面两式相减，可以得到零向量的一种表达方式：
    $$ \mathbf{0} = (\mathbf{u}_1 - \mathbf{v}_1) + (\mathbf{u}_2 - \mathbf{v}_2) + \cdots + (\mathbf{u}_m - \mathbf{v}_m). $$
    注意到$\mathbf{u}_1 - \mathbf{v}_1\in \mathbb{U}_1,\,\,\mathbf{u}_2 - \mathbf{v}_2\in \mathbb{U}_2,\,\,\cdots,\,\,\mathbf{u}_m - \mathbf{v}_m\in \mathbb{U}_m$
    这说明$\mathbf{u}_1 - \mathbf{v}_1$、$\mathbf{u}_2 - \mathbf{v}_2$等都是零向量，也就是说：$ \mathbf{u}_1 = \mathbf{v}_1, \mathbf{u}_2 = \mathbf{v}_2, \cdots \mathbf{u}_m = \mathbf{v}_m.$
    因此，任一向量$\mathbf{u}$都只能以唯一方式写成子空间$\mathbb{U}_1, \mathbb{U}_2, \cdots, \mathbb{U}_m$中向量之和。
    $\mathbb{V} = \mathbb{U}_1 \oplus \mathbb{U}_2 \oplus \cdots \oplus \mathbb{U}_m$。\\
\end{proof}


以上结论说明，要证明子空间的和是直和，可以把对所有向量的验证工作，简化为对零向量的验证。

子空间的和与直和的关系，有些类似子集的并集和不交并。
我们知道，把集合表示为子集的不交并，等同于把集合做划分，分拆成几个部分。
不同部分交集是空集。

子空间的交集不可能是空集，因为零向量是任何子空间必有的元素。而且零向量是唯一的，因此必然是各个子空间公共的元素。
而子空间之和为直和，说明任意两个子空间都没有零向量以外的公共元素。

对于两个子空间的情况，我们可以将这个想法简单地写成：
\begin{pp}\label{pp:1_9}
    设空间$\mathbb{V} = \mathbb{U} + \mathbb{W} $，则$\mathbb{V} = \mathbb{U} \oplus \mathbb{W}$当且仅当$\mathbb{U} \cap \mathbb{W} = \{\mathbf{0}\}$.
\end{pp}

这样我们就有了一个简单的判定法则：如果要验证两个子空间之和为直和，只需验证它们除了零向量再没有公共元素即可。

比如，$\mathbb{R}^2$中两条直线：$\mathcal{D}_1:\,\,y = a_1x$和$\mathcal{D}_2:\,\,y = a_2x$，它们没有非零公共元素，当且仅当$a_1 \neq a_2$。
所以，$\mathcal{D}_1+\mathcal{D}_2 = \mathcal{D}_1\oplus\mathcal{D}_2$当且仅当$a_1 \neq a_2$.

又如：考虑所有$\mathbb{R}^3$中一条过原点的直线$\mathcal{D}$与一个过原点的平面$\mathcal{P}$，
它们都是$\mathbb{R}^3$的子空间。它们显然相交于原点。
我们知道立体空间中，直线和平面要么平行，要么相交，要么直线在平面上。
所以，要么$\mathcal{D} \subset \mathcal{P}$，即$\mathcal{D} + \mathcal{P} = \mathcal{P}$，
要么$\mathcal{D} + \mathcal{P} = \mathcal{D} \oplus\mathcal{P}$.

而对$\mathbb{R}^3$中两个过原点的平面来说，
由于立体空间里平面要么不相交，要么交集是直线，所以两个平面之和必然不是直和。

\sectionext{商空间}

定义单个向量与空间某个子集的和，以及单个标量与空间子集的乘积：
\begin{align*}\forall \,\,\mathbf{v} \in \mathbb{V}, \,\,\,U \subseteq \mathbb{V}, & \quad \mathbf{v} + U = \{\mathbf{v} + \mathbf{u} | \mathbf{u} \in U\}. \\
\forall \,\,\lambda \in \mathbb{F}, \,\,\,U \subseteq \mathbb{V}, & \quad \lambda \cdot U = \{\lambda \cdot \mathbf{u} | \mathbf{u} \in U\}.
\end{align*}
$\mathbf{v} + U$就是将集合$U$里每个向量加上同一个向量$\mathbf{v}$，$\lambda U$就是把$U$里每个向量放缩同样倍数$\lambda$。

举例来说，设$U$是平面中的某个形体，那么$\mathbf{v} + U$就是把它按照向量$\mathbf{v}$平移。
比如，$S = \{(x,y) | 0 \leqslant x, y \leqslant 1\}$是一个正方形，
而$(0.2, 0.4) + S = \{(x,y) | 0.2 \leqslant x \leqslant 1.2, 0.4 \leqslant y \leqslant 1.4\}$是把这个正方形按照$(0.2, 0.4)$平移，得到位置不同的正方形。$1.5 S = \{(x,y) | 0 \leqslant x, y \leqslant 1.5\}$是把这个正方形按1.5比例放大。

同样地，如果考虑直线$P = \{(x,y) | y = 2 x + 0.1394\}$，那么$(0.5, 1) + P = \{(x,y) | y - 1 = 2 (x - 0.5) + 0.1394\}$就是把直线$P$按照$(0.5, 1)$平移，
而$(1, 1) + P = \{(x,y) | y - 1 = 2 (x - 1) + 0.1394\}$则是把直线$P$按照$(1, 1)$平移。

$(0.5, 1) + P$和原来的直线$P$重合，而$(1, 1) + P$与$P$没有交点。
另一方面，$(0.2, 0.4) + S$与$S$部分重合，只有$(0, 0) + S$和$S$完全重合。

我们考虑一种特殊情况：$U$是子空间（以下记为$\mathbb{U}$）。
这时，对于两个向量$\mathbf{v}, \mathbf{w} \in \mathbb{V}$，$\mathbf{v} + \mathbb{U}$和$\mathbf{w} + \mathbb{U}$是什么关系呢？

我们考虑两者的交集。设两个集合有公共元素：$\mathbf{x} = \mathbf{v} + \mathbf{u}_1 = \mathbf{w} + \mathbf{u}_2, \,\,\, \mathbf{u}_1, \mathbf{u}_2 \in \mathbb{U}$，
那么$\mathbf{v} - \mathbf{w} = \mathbf{u}_1 - \mathbf{u}_2 \in \mathbb{U}$. 
而如果$\mathbf{v} - \mathbf{w} \in \mathbb{U}$，那么显然$\{\mathbf{v} + \mathbf{u} | \mathbf{u} \in \mathbb{U}\} = \{\mathbf{w} +  (\mathbf{v} - \mathbf{w}) + \mathbf{u} | \mathbf{u} \in \mathbb{U}\}.$ 
然而$\{\mathbf{v} - \mathbf{w} + \mathbf{u} | \mathbf{u} \in \mathbb{U}\}$就是子空间$\mathbb{U}$。
所以只要$\mathbf{v} + \mathbb{U}$和$\mathbf{w} + \mathbb{U}$有公共元素，两者就是同一个子空间。
从直观的角度来说，要么两者不相交，要么两者重合。

考虑映射：
\begin{align*}
    \tau_\mathbb{U}: \mathbb{V} &\rightarrow \mathcal{P}(\mathbb{V})  \\
    \mathbf{v} &\mapsto \mathbf{v} + \mathbb{U}
\end{align*}
其中的$\mathcal{P}(\mathbb{V})$是$\mathbb{V}$所有子集的集合，称为$\mathbb{V}$的\textbf{幂集}。出发空间$\mathbb{V}$是平直空间，我们希望到达空间也能是平直空间。为此，我们定义以下的运算规则：
\begin{align*}
    \forall \,\,\lambda \in \mathbb{F}, \,\,\, \mathbf{v}, \mathbf{w} \in \mathbb{V}, \\
    \left(\mathbf{v} + \mathbb{U}\right) + \left(\mathbf{w} + \mathbb{U}\right) &= (\mathbf{v} + \mathbf{w}) + \mathbb{U},  \\
    \lambda \left(\mathbf{v} + \mathbb{U}\right) &= (\lambda\mathbf{v}) + \mathbb{U}. 
\end{align*}
读者可以验证，这些定义在逻辑上自洽，而且对给定的子空间$\mathbb{U}$，
$\tau_\mathbb{U}$的像集关于这两种运算也构成一个平直空间。
这个空间的元素都可以写成$\mathbf{v} + \mathbb{U}$的形式。
我们称这个空间为$\mathbb{V}$关于$\mathbb{U}$的\textbf{商空间}，记为$\mathbb{V}/_\mathbb{U}$。
$\mathbb{V}/_\mathbb{U}$的零元素是$\mathbf{0} + \mathbb{U} = \mathbb{U}$。我们把映射重新写为：
\begin{align*}
    \tau_\mathbb{U}: \mathbb{V} &\rightarrow \mathbb{V}/_\mathbb{U}  \\
    \mathbf{v} &\mapsto \mathbf{v} + \mathbb{U}.
\end{align*}
按照定义，$\tau_\mathbb{U}$是满射。

要注意，$\mathbb{V}/_\mathbb{U}$并非$\mathbb{V}$的子空间。
用平面$\mathbb{R}^2$作为例子的话，设子空间$P = \{(x, 3x) |x \in \mathbb{R}\}$，它是一条过原点的直线。
而$\mathbb{V}/_P$就是所有形如$\{(x + a, 3x + b) |x \in \mathbb{R}\}$的直线，
其中$(a, b)$就是对应的向量$\mathbf{v}$。
我们把这些直线写为$\{(x, y) |y = 3x + b - a , \,\,\,x \in \mathbb{R}\}$的形式，可以看出，
这些直线就是所有与直线$P$平行的直线，以及$P$自身。可以说，我们把平面划分为与$P$平行的直线以及$P$自身的不交并集，
或者说，我们用$P$和它的平行线“铺满”了平面。

\chapter{有限维平直空间}

第一章中，我们介绍了平直空间的一些基本性质。
我们看到，在讨论子空间的关系的时候，零向量能否写成某些不全为零的向量之和是很关键的。
本章我们继续讨论向量之间的关系，从这一点出发，进一步探讨平直空间的性质。

\begin{center}
    \fbox{
        \shortstack[l]{
      下文中，除非另有声明，我们总假设$\mathbb{F}$是常见数域如\\
      $\mathbb{Q}$、$\mathbb{R}$或$\mathbb{C}$，
      $\mathbb{V}$是$\mathbb{F}$上的平直空间。
        }
    }
\end{center}

\section{生成空间}
\begin{df}\label{df2_1}
    给定一组向量$\mathbf{v}_1, \mathbf{v}_2 , \cdots , \mathbf{v}_m \in \mathbb{V}$，给定标量$a_1, a_2, \cdots, a_m\in \mathbb{F}$，
    $$a_1\mathbf{v}_1+a_2\mathbf{v}_2 + \cdots + a_m\mathbf{v}_m$$
    称为$\mathbf{v}_1, \mathbf{v}_2 , \cdots , \mathbf{v}_m$的\textbf{直组合}。
\end{df}

\begin{pp}\label{pp:2_1}
    给定一组向量$\mathcal{S} \subset \mathbb{V}$，从$\mathcal{S}$中抽取有限个向量，构成直组合。
    所有这些直组合的集合构成一个子空间，称为它们生成的空间。记作：
    $$ \left \langle \mathbf{v}_1, \mathbf{v}_2 , \cdots , \mathbf{v}_m \right \rangle$$.
\end{pp}
我们可以按平直空间的定义验证，也可以按子空间的判断标准验证。

另一方面，假设某个子空间包含$\mathbf{v}_1, \mathbf{v}_2 , \cdots , \mathbf{v}_m$，由于子空间对向量加法和数乘运算封闭，它自然也包含了所有$\mathbf{v}_1, \mathbf{v}_2 , \cdots , \mathbf{v}_m$的所有直组合。
所以，任何包含$\mathbf{v}_1, \mathbf{v}_2 , \cdots , \mathbf{v}_m$的子空间都包含其生成空间包含$\left\langle\mathbf{v}_1, \mathbf{v}_2 , \cdots , \mathbf{v}_m\right\rangle$. 
生成空间是包含一组向量的“最小”的子空间。它是任何包含这组向量的空间的子空间。

反过来，如果某个空间是一组向量$\mathcal{S}$生成的空间，我们称$\mathcal{S}$是该空间的\textbf{生成集}。

注意：这里$\mathcal{S}$被称为生成集，但它不一定是一个集合，我们也不要求$\mathcal{S}$是一个集合。

如果一个平直空间可以通过有限个向量生成，我们就称它为{\bfseries 有限维平直空间}（在平直代数中简称{\bfseries 有限维空间}）。否则称其为{\bfseries 无限维平直空间}。

由所有$\mathbb{F}$-有序$n$元组构成的空间$\mathbb{F}^n$就是一个有限维空间。因为形如
$$ (1,0,\cdots, 0), \, (0,1,\cdots,0),\cdots ,(0,0,\cdots,1)$$
的$n$个向量可以生成整个空间$\mathbb{F}^n$。$n=2$的时候，我们得到熟悉的平面，$n=3$的时候，我们得到立体空间。

所有整式构成的空间$\mathbb{F}[z]$，则是无限维空间。因为如果我们考虑其中的有限个向量，也就是整式，设它们之中最高的次数是$m$，那么任何$m+1$次的整式都无法表示为它们的直组合。
也就是说，有限个整式没法生成$\mathbb{F}[z]$。
不过，所有次数不超过$n$的整式构成的空间$\mathbb{F}_n[z]$就是有限维空间。
因为$1, z, \cdots, z^n$就能生成整个空间$\mathbb{F}_n[z]$。

如果我们考虑$1,z,z^2$这三个向量，它们生成的空间是所有次数不大于2的整式$a_0+a_1z+a_2z^2$构成的空间$\mathbb{F}_2[z]$。
那么$1,z,2z-1,z^2$这四个向量呢？它们生成的空间仍然是$\mathbb{F}_2[z]$。

我们可以发现，相比$1,z,z^2$，加入$2z-1$这个向量，有点“多余”。为什么呢？因为$2z-1$可以用$1$和$z$生成出来。
$$ 2z-1 = -1\cdot 1+2\cdot z$$
所以如果某个整式可以表示为：$ p = a_1\cdot 1+a_2\cdot z+a_3\cdot (2z-1)+a_4\cdot z^2$，
它也可以表示为$ p = (a_1-a_3)\cdot 1+(a_2+2a_3)\cdot z+a_4\cdot z^2$.
我们不“需要”$2z-1$。$2z-1$在生成空间的时候，并没有独特的“贡献”。

容易发现，用足够多的向量来生成空间并不是难事。比如，空间$\mathbb{V}$自己肯定可以生成$\mathbb{V}$。
所以，我们比较关心的是最小生成集的问题，即：生成空间$\mathbb{V}$，最少要用多少个向量？

在上一个例子中，$\mathbb{F}_2[z]$有$\{1,z,z^2\}$和$\{1,z,2z-1,z^2\}$
这两个生成集。但前者比后者小。我们发现$2z-1$有些“多余”，因为它可以被其它的向量“造出来”。

我们把这个想法推广到一般的一组向量$\mathbf{v}_1, \mathbf{v}_2 , \cdots , \mathbf{v}_m$上。
\begin{df}\label{df:2_2}
    向量$\mathbf{v}_1, \mathbf{v}_2 , \cdots , \mathbf{v}_m$被称为\textbf{直相关}，当且仅当存在不全为零的标量$a_1, a_2 , \cdots , a_m$，使得：
    $$a_1\mathbf{v}_1+a_2\mathbf{v}_2+\cdots+a_m\mathbf{v}_m = \mathbf{0}.$$
    反之，向量$\mathbf{v}_1, \mathbf{v}_2 , \cdots , \mathbf{v}_m$被称为\textbf{直无关}，当且仅当使得
    $$a_1\mathbf{v}_1+a_2\mathbf{v}_2+\cdots+a_m\mathbf{v}_m = \mathbf{0}$$
    的标量$a_1, a_2 , \cdots , a_m$只有一组：$a_1 = a_2 = \cdots = a_m = 0.$
\end{df}
若干个向量直相关，说明其中有一些（至少一个）向量能够写成其他向量的直组合。
反之，直无关的一组向量里，没有任何向量可以表示为其他向量的直组合。

从零向量分解的角度来看，直相关的向量给出了零向量非平凡的分解。
所以，用这些向量表示其他向量的时候，可以通过加上零向量的方式，
给出同一个向量的多种表示方式。

比如，$0 = 1\cdot 1+(-2)\cdot z+1\cdot(2z-1)$，
所以如果$ p = a_1\cdot 1+a_2\cdot z+a_3\cdot(2z-1)$，那么
$p = p + 0 = (a_1+1)\cdot 1+(a_2-2)\cdot z+(a_3+1)\cdot(2z-1)$. 
不严谨地说，$\{1,z,2z-1,z^2\}$把$\mathbb{F}_2[z]$重复生成了多遍；
而$\{1,z,z^2\}$把$\mathbb{F}_2[z]$恰好生成了一遍，因为每个向量都只有一种表示成$\{1,z,z^2\}$直组合的方式。

直觉上来看，向量个数增多会让他们“更相关”，数量减少则使他们更“不相关”。
用更清晰的语言来说，就是：如果一组向量直相关，那么增加一个向量后，新的这组向量仍然直相关。
如果一组向量直无关，那么删去一个向量后，新得到的这组向量仍然直无关。

如果删去某个向量之后，变成了空集呢？我们尚未定义空集是直相关还是无关。
为了让以上这个“准则”适用于所有集合，我们约定：空集是直无关的。

我们知道单个非零向量是直无关的。对于直相关的一组向量$\mathcal{S}$（假设不全为零向量），我们从其中某个非零向量构成的单元素组合开始，把$\mathcal{S}$里的其他向量逐个添加进来，
那么，从某个时刻开始，这组向量会变得直相关。
\begin{pp}\label{pp:2_2}
    设有不全为零的一组向量$\mathbf{v}_1, \mathbf{v}_2 , \cdots , \mathbf{v}_m$直相关。不妨设$\mathbf{v}_1 \neq 0$，那么存在下标$j \in \{2,3, \cdots, m\}$满足以下性质：
    \begin{enumerate}
    \item $\mathbf{v}_j \in \langle \mathbf{v}_1, \mathbf{v}_2 , \cdots , \mathbf{v}_{j-1}\rangle$
    \item 如果从$\mathbf{v}_1, \mathbf{v}_2 , \cdots , \mathbf{v}_m$中移除$\mathbf{v}_j$，剩余向量生成的空间与原来$m$个向量生成的空间相同。
    \end{enumerate}
\end{pp}
\begin{proof} % [{\bfseries 证明\ref{pp:2_2}}]
    按直相关的定义，存在不全为零的标量系数$a_1, a_2, \cdots, a_m$使得：
    $$ a_1\mathbf{v}_1+a_2\mathbf{v}_2+\cdots+a_m\mathbf{v}_m = \mathbf{0}. $$
    由于$\mathbf{v}_1 \neq 0$，所以$a_2, \cdots, a_m$不全为0. 否则可推出$a_1\mathbf{v}_1 = \mathbf{0}$，从而$a_1$也是0，矛盾。\\
    假设$j$是$a_2, \cdots, a_m$中不为0的数的下标中最大的那个，也就是说$\forall k > j, a_k = 0$. 这表示$ a_1\mathbf{v}_1+a_2\mathbf{v}_2+\cdots+a_j\mathbf{v}_j = \mathbf{0}. $
    所以
    $$ \mathbf{v}_j = -\frac{a_1}{a_j}\mathbf{v}_1-\frac{a_2}{a_j}\mathbf{v}_2-\cdots-\frac{a_{j-1}}{a_j}\mathbf{v}_{j-1} \in \langle \mathbf{v}_1, \mathbf{v}_2 , \cdots , \mathbf{v}_{j-1}\rangle. $$
    由于$ \mathbf{v}_j$可以表示为$\mathbf{v}_1, \mathbf{v}_2 , \cdots , \mathbf{v}_{j-1}$的直组合，设某个向量是$\mathbf{v}_1, \mathbf{v}_2 , \cdots , \mathbf{v}_m$的直组合：
    $$ \mathbf{u} = b_1\mathbf{v}_1+b_2\mathbf{v}_2+\cdots+b_m\mathbf{v}_m$$
    把右侧表达式中的$\mathbf{v}_j$替换为$\mathbf{v}_1, \mathbf{v}_2 , \cdots , \mathbf{v}_{j-1}$的直组合，就得到了不含$\mathbf{v}_j$的直组合表达式。\\
    所以，去除$\mathbf{v}_j$，整组向量的生成空间不变\footnote{其实我们只证明了一个方向的归属关系，不过另一个方向比较显然，留给读者思考。}。\\
\end{proof}


我们通常把以上这种命题称为“存在性命题”，它们声称的是某种存在性。我们给出的证明是一种“构造性证明”：直接构造出符合要求的对象，证明其存在。

一组向量是否直相关，与它是否是空间的生成集有没有关系呢？
直觉上，向量越少，越不容易直相关，也越不容易生成整个空间。
如果一组向量能生成整个空间，那么任意一个向量都能表示为这组向量的直组合。
因此往空间的生成集中添加任何向量，都会让它变得直相关。

所以，我们可以猜想，直无关的向量组和空间的生成集相比，元素更少。
比如，我们从某个非零向量$\mathbf{v}$出发，逐个从平直空间里“拿出”向量，把它和$\mathbf{v}$放到一起。
如果新向量在之前向量生成的空间里，就不放入，否则放入。那么这组向量能生成的空间越来越“大”。
全部向量都放完之后，生成的空间应当就是整个平直空间。

用严谨的语言来说，就是：
\begin{pp}\label{pp:2_3}
    有限维平直空间中，任何一组直无关的向量，其元素个数都不多于任意一组生成集的元素个数。
\end{pp}
 
这是一个重要的结论（证明参见附录三\ref{c-10}）。我们现在可以认为，有限维平直空间中直无关的向量组，元素个数有上限，而平直空间的生成集，其元素个数有下限。
我们可以用这个结论来证明很多性质。比如说，直觉告诉我们，有限维平直空间的子空间也是有限维空间。现在我们可以证明这一点：

\begin{pp}\label{pp:2_4}
    有限维平直空间的子空间也是有限维空间。
\end{pp}
\begin{proof} % [{\bfseries 证明\ref{pp:2_3}}]
    设$\mathbb{U}$是有限维平直空间$\mathbb{V}$的子空间。
    $(\mathbf{v}_1, \mathbf{v}_2, \cdots , \mathbf{v}_m )$是$\mathbb{V}$的生成集。\\
    若$\mathbb{U} = \{\mathbf{0}\}$，则$\mathbb{U}$显然为有限维空间。\\
    若$\mathbb{U} \neq \{\mathbf{0}\}$，则从$\mathbb{U}$中取非零向量$\mathbf{u}_1$。
    若$\mathbb{U} = \langle\mathbf{u}_1\rangle$，则$\mathbb{U}$为有限维空间。
    若$\mathbb{U}$中有元素不属于$\langle\mathbf{u}_1\rangle$，则将其中一个取出，和$\mathbf{u}_1$放在一起。\\
    依此操作，第$j$步时，首先检验是否$\mathbb{U} = \langle\mathbf{u}_1, \cdots , \mathbf{u}_{j-1}\rangle$，
    若$\mathbb{U} = \langle\mathbf{u}_1, \cdots , \mathbf{u}_{j-1}\rangle$，则$\mathbb{U}$为有限维空间。
    若$\mathbb{U}$中有元素不属于$\langle\mathbf{u}_1, \cdots , \mathbf{u}_{j-1}\rangle$，取其中一个：
    $\mathbf{u}_j \notin \langle\mathbf{u}_1, \cdots , \mathbf{u}_{j-1}\rangle$，放入$\mathbf{u}_1, \cdots , \mathbf{u}_{j-1}$中。 \\
    每次操作后，向量组增加一个元素，而$\mathbf{u}_1, \cdots , \mathbf{u}_{j}$都是直无关的向量，所以其元素个数不超过$m$。
    所以，有限次（不超过$m$次）操作后，必然有$\mathbb{U} = \langle\mathbf{u}_1,\cdots , \mathbf{u}_{j}\rangle$（$j \leqslant m$）.\\
    所以，$\mathbb{U}$也是有限维空间，而且其最小生成集的元素个数不多于$\mathbb{V}$的最小生成集的元素个数。\\
\end{proof}


\section{空间的基}
通过上一节的讨论，我们已经能得出这样的结论：平直空间中直无关向量的最大个数，和最小生成集的元素个数，刻画了平直空间的重要性质。
现在我们将这个性质明确到一个新的概念上：平直空间的{\bfseries 基}。

\begin{df}\label{df:2_3}
    平直空间$\mathbb{V}$的一组{\bfseries 基}（也称{\bfseries 基底}），是一组生成$\mathbb{V}$且直无关的向量。
\end{df}

举例来说，给定正整数$n$，$n$元数组构成的空间$\mathbb{F}^n$就包含一组基：
$$ \mathcal{B}_c = \left((1,0,\cdots,0),(0,1,\cdots,0),\cdots,(0,0,\cdots,1)\right).$$
这组基也称作$\mathbb{F}^n$的{\bfseries 标准基}。

当然，$\mathbb{F}^n$不止包含$\mathcal{B}_c$这一组基，还有其他的基底。
比如把$\mathcal{B}_c$中的$(1,0,\cdots,0)$换成$(-1,1,\cdots,0)$，就又得到一组基。

令$n = 2$，我们来看$\mathbb{R}^2$的基。可以验证，$\left((1,0),(3,-2)\right)$是一组基，$\left((-1,2),(0,1)\right)$也是一组基。
$\left((1,0),(3,-2),(-2,2)\right)$则不是基，因为$(1,0) = (3,-2) + (-2,2).$ 
$(1,2)$也不是基，因为它无法生成$\mathbb{R}^2.$

基中的向量称为\textbf{基向量}。

注意，以上定义中，似乎并没有给$\{\mathbf{0}\}$定义基。
我们约定：空集是直无关的，并且$\{\mathbf{0}\}$的生成集是空集.

空间的基有什么用处呢？我们来看下面这个命题：
\begin{pp}\label{pp:2_5}
    $ \mathcal{B} = (\mathbf{e}_1, \mathbf{e}_2, \cdots , \mathbf{e}_n )$是空间$\mathbb{V}$的一组基，
    当且仅当空间中每个元素$\mathbf{u}$都能以唯一的方式写成
    $$ \mathbf{u} = a_1\mathbf{e}_1 +a_2 \mathbf{e}_2 + \cdots + a_n\mathbf{e}_n, \,\,\, a_1,\,a_2,\,\cdots,\, a_n \,\in\,\mathbb{F}.$$
    我们称$a_1, a_2, \cdots a_m$为“$\mathbf{u}$的$\mathcal{B}$分解”。
\end{pp}
这个式子似曾相识。我们还记得在空间的直和中有类似的的表述。空间的直和类似于集合的分划，把空间拆分成几个较小的空间。
但每个较小空间的性质仍然不得而知。而空间的基则以有限个直无关的向量给出了空间中每个元素的唯一分解方式\footnote{要注意的是，这里我们没有明确写出来，但隐含的前提是：$\mathbb{V}$是有限维空间。如果$\mathbb{V}$是无限维空间，基底的问题会变得比较复杂。}。

\begin{proof} % [{\bfseries 证明\ref{pp:2_5}}]
    充分性：\\
    $\mathcal{B}$是$\mathbb{V}$的生成集，所以$\mathbb{V}$中每个元素$\mathbf{u}$都能写成$\mathcal{B}$中元素的直组合：
    $$ \mathbf{u} = a_1\mathbf{v}_1 +a_2 \mathbf{v}_2 + \cdots + a_n\mathbf{v}_n, \,\,\, a_1,\,a_2,\,\cdots,\, a_n \,\in\,\mathbb{F}.$$
    只需证明这个写法是唯一的。\\
    假设还有另一种写法：
    $$ \mathbf{u} = a_1'\mathbf{v}_1 +a_2' \mathbf{v}_2 + \cdots + a_n'\mathbf{v}_n, \,\,\, a_1' ,\,a_2' ,\,\cdots,\, a_n' \,\in\,\mathbb{F}.$$
    两式相减，可得：
    $$ \mathbf{0} = (a_1 - a_1')\mathbf{v}_1 + (a_2 - a_2') \mathbf{v}_2 + \cdots + (a_n - a_n')\mathbf{v}_n.$$
    而$\mathcal{B}$是直无关的，因此根据定义，
    $$ a_1 - a_1' = a_2 - a_2' = \cdots = a_n - a_n' = 0.$$
    这说明写法是唯一的。\\
    必要性：\\
    我们要证明，在每个元素都能唯一分解为$\mathcal{B}$中元素直组合的条件下，$\mathcal{B}$是$\mathbb{V}$的直无关的生成集。\\
    分解的存在性，可以说明$\mathcal{B}$是$\mathbb{V}$的生成集。\\ 
    分解的唯一性，可以说明$\mathcal{B}$是直无关的。这是因为，零向量显然可以写成：
    $$ \mathbf{0} = 0\cdot\mathbf{v}_1 + 0\cdot\mathbf{v}_2 + \cdots + 0\cdot\mathbf{v}_n.$$
    由于分解是唯一的，所以不存在其他表示方式。这就证明$\mathcal{B}$是直无关的。\\
\end{proof}


既然每个向量都能唯一表示成基向量的直组合，我们将相应的系数写成有序$n$元组，就得到了每个向量在给定基底上的“坐标”。
这样看来，平直空间和有序$n$元组构成的空间是“等价”的。是否如此呢？我们在后面会作进一步分析。

空间的基是否就是“最大”的直无关向量组、“最小”的生成集呢？我们首先来看生成集和基底的关系：

\begin{pp}\label{pp:2_6}
    可以通过适当地移除元素，把有限维空间$\mathbb{V}$的每个生成集变成一组基，
\end{pp}

\begin{proof} % [{\bfseries 证明\ref{pp:2_6}}]
    设$\mathcal{F}_{\mathbf{v}} = (\mathbf{v}_1, \mathbf{v}_2, \cdots , \mathbf{v}_m )$是$\mathbb{V}$的生成集。
    如果$\mathcal{F}_{\mathbf{v}}$直无关，则按照定义，$\mathcal{F}_{\mathbf{v}}$是一组基。\\
    如果$\mathcal{F}_{\mathbf{v}}$直相关。我们进行如下操作：\\
    首先移除$\mathcal{F}_{\mathbf{v}}$中所有的零向量。如果$\mathcal{F}_{\mathbf{v}}$完全由零向量构成，说明$\mathbb{V} = \{\mathbf{0}\}$。这时，$\mathcal{F}_{\mathbf{v}}$移除全部元素后剩下空集，正好直无关且生成$\mathbb{V}$。\\
    如果$\mathcal{F}_{\mathbf{v}}$中还有若干非零向量，且仍直相关，根据 \ref{pp:2_2}，可以从其中找出一个向量，把它从$\mathcal{F}_{\mathbf{v}}$移除后，生成空间不变，$\mathcal{F}_{\mathbf{v}}$仍旧是生成集。\\
    如果$\mathcal{F}_{\mathbf{v}}$仍然直相关，则重复以上操作。每次操作后，$\mathcal{F}_{\mathbf{v}}$中元素数量减少1. 所以，有限次操作之后，要么$\mathcal{F}_{\mathbf{v}}$变为空集，无法生成空间；要么$\mathcal{F}_{\mathbf{v}}$变为直无关的生成集。
    这时我们就得到了空间的一组基。\\
\end{proof}


我们用一个具体的例子来演示一遍证明里的操作。

$\mathcal{F} = \left((1,0),(8,2),(3,1),(0,4),(0,0)\right)$是$\mathbb{R}^2$的生成集。

我们首先去除$\mathcal{F}$里的零向量，得到$\left((1,0),(8,2),(3,1),(0,4)\right).$ 这组向量直相关：
$$ (8,2) = 2\cdot(1,0) + 2\cdot(3,1).$$
去除$(8,2)$，得到$\left((1,0),(3,1),(0,4)\right).$ 这组向量仍然直相关：
$$ (3,1) = 3\cdot(1,0) + 0.25\cdot(0,4).$$
去除$(3,1)$，得到$\left((1,0),(0,4)\right).$ 现在我们得到了一组直无关的向量，而且能生成$\mathbb{R}^2$，也就是$\mathbb{R}^2$的一组基。

从这个命题可以推出：
\begin{cor}\label{cor:2_1}
    任意有限维空间都包含一组基。
\end{cor}

空间的任意生成集都可以缩减为它的基底。那么任意直无关的向量能否扩充为一组基呢？
\begin{pp}\label{pp:2_7}
    可以通过适当地添加元素，把有限维空间$\mathbb{V}$的每一组直无关的向量变成一组基。
\end{pp}
\begin{proof} % [{\bfseries 证明\ref{pp:2_7}}]
    设$\mathcal{F}_{\mathbf{v}} = (\mathbf{v}_1, \mathbf{v}_2, \cdots , \mathbf{v}_m )$是$\mathbb{V}$中一组直无关的向量。
    如果$\mathcal{F}_{\mathbf{v}}$生成$\mathbb{V}$，则按照定义，$\mathcal{F}_{\mathbf{v}}$是一组基。\\
    如果$\mathcal{F}_{\mathbf{v}}$生成的空间不是$\mathbb{V}$，则是$\mathbb{V}$的真子空间。\\
    空间$\mathbb{V}$是有限维空间，所以存在生成集：$\mathcal{F}_{\mathbf{u}} = (\mathbf{u}_1, \mathbf{u}_2, \cdots , \mathbf{u}_n )$. 我们进行以下操作：\\
    逐个查看$(\mathbf{u}_1, \mathbf{u}_2, \cdots , \mathbf{u}_n )$，其中必然有一个元素不在$\langle\mathcal{F}_{\mathbf{v}}\rangle$生成的空间中。\\
    这是因为，反设$\mathcal{F}_{\mathbf{u}}$中每个元素都能写成$\mathcal{F}_{\mathbf{v}}$中元素的直组合，那么由于$\mathbb{V}$中任一元素都能写成$\mathcal{F}_{\mathbf{u}}$中元素的直组合，从而也能写成$\mathcal{F}_{\mathbf{v}}$中元素的直组合。
    这说明$\mathcal{F}_{\mathbf{v}}$可以生成整个空间，矛盾！\\
    所以，必然存在某个$\mathbf{u}_i$不在$\langle\mathcal{F}_{\mathbf{v}}\rangle$中。我们将它添加到$\mathcal{F}_{\mathbf{v}}$里去，得到的仍然是一组直无关的向量。\\
    每次操作，$\mathcal{F}_{\mathbf{v}}$中增加一个元素。所以，要么在有限次操作后，$\langle\mathcal{F}_{\mathbf{v}}\rangle = \mathbb{V}$，我们得到$\mathbb{V}$的直无关的生成集（也就是基底）；要么$\mathcal{F}_{\mathbf{v}}$中元素个数超过$n$，而这和命题\ref{pp:2_3} 矛盾。\\
\end{proof}


我们还可以从这个命题推出一个更有用的结论：
\begin{pp}\label{pp:2_8}
    可以给有限维空间$\mathbb{V}$的每个子空间$\mathbb{U}$找一个“补空间”：$\mathbb{W}$，也就是使得
    $$ \mathbb{V} = \mathbb{U} \oplus \mathbb{W}.$$
    子空间$\mathbb{U}$的\textbf{补空间}一般记作$\mathbb{U}^c.$
\end{pp}
\begin{proof} % [{\bfseries 证明\ref{pp:2_8}}]
    设$\mathcal{B}_1$是子空间$\mathbb{U}$的一组基。$\mathcal{B}_1$是直无关的。所以可以通过适当添加元素，将$\mathcal{B}_1$扩充为$\mathbb{V}$的基底$\mathcal{B}$。\\
    设$\mathcal{B}_2$是$\mathcal{B}_1$扩充到$\mathcal{B}$过程中添加的所有元素。$\mathcal{B}_2$也是直无关的。考虑$\mathbb{W} = \langle\mathcal{B}_1\rangle$，可以验证：
    $$ \mathbb{V} = \mathbb{U} \oplus \mathbb{W}.$$
    其中$ \mathbb{V} = \mathbb{U} + \mathbb{W}$显然成立。下面证明$ \mathbb{U} \cap \mathbb{W} = \{\mathbf{0}\}$.\\
    设$\mathcal{B}_1 = (\mathbf{v}_1, \mathbf{v}_2, \cdots , \mathbf{v}_m )$，$\mathcal{B}_1 = (\mathbf{v}_{m+1}, \cdots , \mathbf{v}_n )$. 设$ \mathbf{u} = \mathbb{U} \cap \mathbb{W}$，则：
    $$ \mathbf{u} = a_1\mathbf{v}_1 + a_2\mathbf{v}_2 + \cdot + a_m\mathbf{v}_m = a_{m+1}\mathbf{v}_{m+1} + \cdots + a_n\mathbf{v}_n. $$
    $$a_1\mathbf{v}_1 + a_2\mathbf{v}_2 + \cdots + a_m\mathbf{v}_m + (-a_{m+1})\mathbf{v}_{m+1} + \cdots + (-a_{n})\mathbf{v}_{n} = \mathbf{0}. $$
    然而$(\mathbf{v}_1, \mathbf{v}_2, \cdots , \mathbf{v}_n )$是直无关的。所以，
    $$ a_1 = a_2 = \cdots = a_m = -a_{m+1} = \cdots = -a_n = 0.$$
    所以$\mathbf{u} = \mathbf{0}$。
    因此，$ \mathbb{U} \cap \mathbb{W} = \{\mathbf{0}\}$.\\
\end{proof}


\section{空间的维数}
空间的基底使我们能够用空间中的有限个向量来描述空间中任一向量。我们对平直空间的了解又进了一步。
从直无关向量组、生成集到基底，我们认识到，一组向量的性质，和它的元素个数有直接关系。而这个关系是空间内在的性质。
我们自然要问：给定一个空间里，是否所有的基底都有相同的元素个数？如果是的话，这个数反映了空间的什么性质？

\begin{pp}\label{pp:2_9}
    同一个有限维空间$\mathbb{V}$的所有基底都有相等的元素个数，称为此空间的{\bfseries 维数}，记作$\dim\mathbb{V}$。
\end{pp}
\begin{proof} % [{\bfseries 证明\ref{pp:2_9}}]
    设$\mathcal{B}_1, \mathcal{B}_2$是空间$\mathbb{V}$的两组基。
    按照定义，$\mathcal{B}_1$是直无关的，而$\mathcal{B}_2$是该空间的生成集。
    所以，根据命题 \ref{pp:2_2} 的结论，$\mathcal{B}_1$的元素个数不多于$\mathcal{B}_2$的元素个数。
    由对称性可知， $\mathcal{B}_2$的元素个数不多于$\mathcal{B}_1$的元素个数。
    所以，$\mathcal{B}_1$和$\mathcal{B}_2$的元素个数相等。\\
\end{proof}


维数的定义，也解释了本章的名字：有限维平直空间的含义。
每个能够通过有限个向量生成的空间，都有一个内在的性质。这个性质可以用维数来表示。

我们熟悉的平面，其维数是2，所以也常称为二维空间。立体空间维数是3，所以也称为三维空间。

那么，维数对有限维平直空间的意义是什么呢？

维数可以看作是衡量有限维空间“大小”的一个指标。比如，我们可以确定空间和子空间的维数关系：

\begin{pp}\label{pp:2_10}
    有限维空间的子空间的维数不大于其自身的维数；其真子空间的维数小于其自身维数。
\end{pp}

\begin{proof} % [{\bfseries 证明\ref{pp:2_10}}]
    设$\mathbb{U}$是空间$\mathbb{V}$的子空间，
    $\mathcal{B}$是$\mathbb{U}$的一组基。
    按照定义，$\mathcal{B}$是直无关的，所以根据命题 \ref{pp:2_7}，$\mathcal{B}$可以被扩充为$\mathbb{V}$的一组基。 
    这说明$\mathcal{B}$的元素个数小于等于$\mathbb{V}$的维数。\\
    考虑到命题\ref{pp:2_7}证明中扩充的方式，如果$\mathbb{U}$是空间$\mathbb{V}$的真子空间，$\mathbb{V}$中有不属于$\mathbb{U}$的向量，
    所以，$\mathcal{B}$扩充到$\mathbb{V}$的基底的过程中，至少要添加1个向量。这说明，如果$\mathbb{U}$是空间$\mathbb{V}$的真子空间，$\mathcal{B}$的元素个数（严格）小于$\mathbb{V}$的维数。\\
\end{proof}


0维空间只有一个，就是$\{\mathbf{0}\}$，按约定由空集生成的空间。1维空间就是直线概念的推广，形同：
$$ \{k\mathbf{u}\,\,|\,\,k\in\mathbb{F}\}$$
其基底可以是空间中任一非零向量。

设一个$n$维空间$\mathbb{V}$有一组基：$ \mathcal{B} = (\mathbf{v}_1, \mathbf{v}_2, \cdots , \mathbf{v}_n )$.
基中每个元素$\mathbf{v}_i$可以生成一条“直线”，即一维空间：
$$ \mathbf{D}_i =  \{k\mathbf{v}_i\,\,|\,\,k\in\mathbb{F}\}$$
而$\mathbb{V}$可以表示成这$n$条直线的直和：
$$\mathbb{V} = \mathbf{D}_1 \oplus \mathbf{D}_2 \oplus \cdots \oplus \mathbf{D}_n. $$
空间里每个向量都能写成分别来自这$n$条直线的$n$个向量的和。

我们还可以用维数的“语言”重述一些结论。
比如，一组直无关的向量，其元素个数不超过空间的维数。
一组向量能生成空间，其元素个数必然不少于空间的维数。

\begin{pp}\label{pp:2_11}
    若有限维空间的生成集的元素个数等于空间的维数，则它是空间的基。
\end{pp}

\begin{proof} % [{\bfseries 证明\ref{pp:2_11}}]
    只需证明，元素个数等于空间的维数的生成集直无关。\\
    反设其直相关，则可以从其中去除一个向量，使得剩余部分仍然是空间的生成集。
    但这时其元素个数小于空间维数，矛盾！\\
\end{proof}


\begin{pp}\label{pp:2_12}
    若有限维空间中一组直无关向量的元素个数等于空间的维数，则它是空间的基。
\end{pp}

\begin{proof} % [{\bfseries 证明\ref{pp:2_12}}]
    只需证明，这组直无关向量能生成整个空间。\\
    根据命题 \ref{pp:2_7}，这组直无关的向量可以被扩充为空间的基。
    但它已经有和空间的基一样多的元素。
    所以实际上无需扩充。也就是说，它已经是一组基了。\\    
\end{proof}


下面这个结论可以让我们感受到，空间的维数在某些方面和集合元素个数的相似之处：

\begin{pp}\textbf{维数基本定理}\label{pp:2_13}\\
    设$\mathbb{U}$和$\mathbb{W}$是有限维空间$\mathbb{V}$的两个子空间，则它们的维数满足以下关系：
    $$ \dim (\mathbb{U} + \mathbb{W}) = \dim \mathbb{U} + \dim \mathbb{W} - \dim (\mathbb{U} \cap \mathbb{W}).$$
\end{pp}
如果我们把空间的和比作集合的并，把空间的交比作集合的交，把空间的维数比作集合的元素个数，则这个关系和集合论中有限元素集合的关系是一致的。
\begin{proof} % [{\bfseries 证明\ref{pp:2_13}}]
    记$\mathbb{T} = \mathbb{U} \cap \mathbb{W}$. 根据命题 \ref{pp:2_8}，我们可以给$\mathbb{T}$在$\mathbb{U}$和$\mathbb{W}$里各自找一个补空间：
    $$ \mathbb{U} = \mathbb{T} \oplus \mathbb{U}', \,\,\, \mathbb{W} = \mathbb{T} \oplus \mathbb{W}'.$$
    记$\dim \mathbb{T} = h, \,\,\, \dim \mathbb{U}' = f, \,\,\, \dim \mathbb{W}' = g$. 我们考虑这三个空间的基：
    $$ \mathcal{B}_T = (\mathbf{t}_1, \mathbf{t}_2, \cdots, \mathbf{t}_h), \,\,\, \mathcal{B}_{U'} = (\mathbf{u}_1, \mathbf{u}_2, \cdots, \mathbf{u}_f), \,\,\, \mathcal{B}_{W'} = (\mathbf{w}_1, \mathbf{w}_2, \cdots, \mathbf{w}_g). $$  
    其中$\mathcal{B}_{U'}$和$\mathcal{B}_{W'}$分别是$\mathbb{T}$按照命题 \ref{pp:2_8} 证明中的方式扩充为$\mathbb{U}$和$\mathbb{W}$时得到的向量组。
    这表示，$\mathcal{B}_T\cup\mathcal{B}_{U'}$是$\mathbb{U}$的基，$\mathcal{B}_T\cup\mathcal{B}_{W'}$是$\mathbb{W}$的基。
    所以，
    $$\dim \mathbb{U} = h + f,\,\,\,\dim \mathbb{W} = h + g.$$
    只需证明：$ \dim (\mathbb{U} + \mathbb{W}) = h + f + g.$\\
    为此我们证明$\mathcal{B}_T\cup\mathcal{B}_{U'}\cup\mathcal{B}_{W'}$是$\mathbb{U} + \mathbb{W}$的一组基。\\
    首先，$\mathcal{B}_T\cup\mathcal{B}_{U'}\cup\mathcal{B}_{W'}$可以生成$\mathbb{U}$，也可以生成$\mathbb{W}$，所以可以生成$\mathbb{U} + \mathbb{W}$.\\
    其次，证明$\mathcal{B}_T\cup\mathcal{B}_{U'}\cup\mathcal{B}_{W'}$直无关。假设零向量可以分解为$\mathcal{B}_T\cup\mathcal{B}_{U'}\cup\mathcal{B}_{W'}$的直组合，则说明零向量可以表示为：
    $$ \mathbf{0} = \mathbf{t} + \mathbf{u} + \mathbf{w}. \,\,\,\mathbf{t}\in\mathbb{T},\,\,\mathbf{u}\in\mathbb{U}',\,\,\mathbf{w}\in\mathbb{W}'.$$
    所以$\mathbf{u} = -\mathbf{t} - \mathbf{w} \in \mathbb{W}$.
    也就是说，$\mathbf{u}\in \mathbb{U}'\cap\mathbb{W}$.
    然而$\mathbb{U}'\cap\mathbb{W} \subset \mathbb{U}\cap\mathbb{W} = \mathbb{T}$，所以$\mathbf{u}\in\mathbb{T}\cap\mathbb{U}' = \{\mathbf{0}\}$.\\
    因此，$\mathbf{u} = \mathbf{0}.$\\
    同理可得，$\mathbf{w} = \mathbf{0}.$ 所以$\mathbf{t} = \mathbf{0}.$\\
    而$\mathcal{B}_T$、$\mathcal{B}_{U'}$和$\mathcal{B}_{W'}$分别是直无关的。所以零向量分解为$\mathcal{B}_T\cup\mathcal{B}_{U'}\cup\mathcal{B}_{W'}$的直组合时，每个分量系数都是0. 
    这就证明$\mathcal{B}_T\cup\mathcal{B}_{U'}\cup\mathcal{B}_{W'}$直无关，从而是$\mathbb{U} + \mathbb{W}$的一组基。\\
\end{proof}


以上证明中，我们其实还附带证明了这两个结论：
\begin{enumerate}
    \item $\mathbb{U} + \mathbb{W} = \mathbb{U} \oplus \mathbb{W}$当且仅当$\dim (\mathbb{U} + \mathbb{W}) = \dim \mathbb{U} + \dim \mathbb{W}.$
    \item $\mathbb{U} + \mathbb{W} = \left(\mathbb{U} \cap \mathbb{W}\right) \oplus \mathbb{U}' \oplus \mathbb{W}'.$
\end{enumerate}

其中第一个结论可以推广为：
\begin{pp}\label{pp:2_14}
    设$\mathbb{V} = \mathbb{V}_1 + \mathbb{V}_2 + \cdots + \mathbb{V}_n$，则
    $$ \dim \mathbb{V} \leq \dim \mathbb{V}_1 + \dim \mathbb{V}_2 + \cdots + \dim \mathbb{V}_n.$$
    当且仅当$\mathbb{V} = \mathbb{V}_1 \oplus \mathbb{V}_2 \oplus \cdots \oplus \mathbb{V}_n$时上式取等号。\\
\end{pp}

\begin{proof} % [{\bfseries 证明\ref{pp:2_14}}]
    可以使用归纳法。前面已经证明了$n=2$的情况。假设命题对某个$n = n_0 \geq 2$成立，对$n = n_0+1$的情况，我们要证明：
    $$ \dim \mathbb{V} \leq \dim \mathbb{V}_1 + \dim \mathbb{V}_2 + \cdots + \dim \mathbb{V}_{n_0} + \dim \mathbb{V}_{n_0+1}.$$
    等号当且仅当$\mathbb{V} = \mathbb{V}_1 \oplus \mathbb{V}_2 \oplus \cdots \oplus \mathbb{V}_{n_0+1}$时成立。\\
    设$\mathbb{W} = \mathbb{V}_1 + \mathbb{V}_2 + \cdots + \mathbb{V}_{n_0},$ 则$\mathbb{V} = \mathbb{W} + \mathbb{V}_{n_0+1}.$
    应用$n=2$和$n=n_0$时的结论，可知：\\
    \begin{align*}
         \dim \mathbb{V} &\leq \dim \mathbb{W} + \dim \mathbb{V}_{n_0+1}  \\
               &\leq \dim \mathbb{V}_1 + \dim \mathbb{V}_2 + \cdots + \dim \mathbb{V}_n + \dim \mathbb{V}_{n_0+1} 
    \end{align*}
    等号成立当且仅当$n=2$和$n=n_0$的等号成立条件同时成立。而后者等价于$\mathbb{V} = \mathbb{V}_1 \oplus \mathbb{V}_2 \oplus \cdots \oplus \mathbb{V}_{n_0}$且$\mathbb{V} = \mathbb{W} \oplus \mathbb{V}_{n_0+1}$，
    也就是$\mathbb{V} = \mathbb{V}_1 \oplus \mathbb{V}_2 \oplus \cdots \oplus \mathbb{V}_{n_0+1}.$ 
    这样我们就证明了$n = n_0 + 1$的情况。\\
    由归纳假设，可知命题对所有大于等于2的自然数成立。\\
\end{proof}


\chapter{直映射}

通过前两章的讨论，我们对有限维平直空间的基本性质有了初步的了解。
平直空间可以看作研究的平台，在此之上，我们将深入探索空间中定义的对象和映射关系。
但本书不会深入讨论平直空间中的对象。研究平直空间的复杂形体对象需要更深刻的工具。
本书将重点从代数层面，探讨平直空间中最基本一类映射的性质。

\begin{center}
    \fbox{
        \shortstack[l]{
      下文中，除非另有声明，我们总假设$\mathbb{F}$是常见数域如\\
      $\mathbb{Q}$、$\mathbb{R}$或$\mathbb{C}$，
      $\mathbb{V}$和$\mathbb{W}$是$\mathbb{F}$上的平直空间。
        }
    }
\end{center}

\section{直映射的定义}
函数，也称为映射，是描述变换关系的概念。我们定义平直空间中的直映射如下：
\begin{df}\label{df:3_1}
    设有从空间$\mathbb{V}$到$\mathbb{W}$的映射$T$，满足以下条件：
    \begin{enumerate}
        \item $\forall \,\, \mathbf{u} , \mathbf{v} , \,\, \in \mathbb{V}, \,\,\, T(\mathbf{u} + \mathbf{v}) = T(\mathbf{u}) + T(\mathbf{v}). $
        \item $\forall \,\, \mathbf{u} \,\, \in \mathbb{V}, \,\, a \in \mathbb{F}, \,\,\, T(a\mathbf{u}) = aT(\mathbf{u}). $
    \end{enumerate}
    则称之为{\bfseries （平）直映射}或{\bfseries （平）直变换}。
\end{df}
为了与向量区分，本书中尽量使用大写字母表示直映射。
直映射将$\mathbb{V}$中向量$\mathbf{u}$映射为$\mathbb{W}$中向量$T(\mathbf{u})$。
$T(\mathbf{u})$称为$\mathbf{u}$在变换$T$下的\textbf{值}或\textbf{结果}，也称$\mathbf{u}$的\textbf{像}。
我们沿用一般函数的记法，使用$T(\mathbf{u})$表示。
在不至于混淆的时候，也会简记为$T\mathbf{u}$.

直映射可以看作是正比例函数的推广。
我们知道正比例函数的图像是一条（过原点的）直线。
直映射中的“线”就是指“直线”。直映射是具有“直性”的映射。

具体来说，在定义里，直映射要求映射具有两个性质。
第一个性质称为“{\bfseries 可加性}”，即保持向量加法运算性质不变。
需要注意的是，$T(\mathbf{u} + \mathbf{v})$中的加号表示空间$\mathbb{V}$（也称为出发空间）中的向量加法，而$T(\mathbf{u}) + T(\mathbf{v})$中的加号表示$\mathbb{W}$（也称为到达空间）中的向量加法。
第二个性质称为“{\bfseries 齐次性}”，即保持数乘运算性质不变。
同样，需要注意的是，$T(a\mathbf{u})$中的数乘是$\mathbb{V}$中的数乘，而$aT(\mathbf{u})$中的数乘是$\mathbb{W}$中的数乘。

有了这两个性质，直映射就有效地保存了平直空间$\mathbb{V}$、$\mathbb{W}$的直性。
这种保持代数结构性质的映射，被称为同态（保持相同的形态）。直映射保持平直空间性质不变，因此称为\textbf{直同态}。

所有从$\mathbb{V}$到$\mathbb{W}$的直映射的集合记作$\mathcal{L}(\mathbb{V},\mathbb{W})$.
如果$\mathbb{V} = \mathbb{W}$，则记作$\mathcal{L}(\mathbb{V})$.  
$\mathcal{L}(\mathbb{V})$中的元素则称为\textbf{直自同态}，即映射到自身的同态，也称为\textbf{直算子}。
在不至于混淆的情况下，我们可以省略“直”或“平直”，直接称之为“自同态”或“算子”。

我们先来看几个以后会经常用到的例子（读者可以自己试着验证它们是否满足直映射的性质）：

\paragraph{零映射} 零映射指把$\mathbb{V}$中所有向量映射到$\mathbb{W}$中零向量的映射。
在不至于混淆的情况下，我们一般用$\mathrm{0}$表示零映射。
\begin{align*}
    \mathrm{0}: \,\, \mathbb{V} & \to \mathbb{W}  \\
    \mathbf{v} & \mapsto \mathbf{0}_W 
\end{align*}

\paragraph{恒等映射} 恒等映射指把$\mathbb{V}$中向量映射到$\mathbb{V}$中同一个向量的映射。
在不至于混淆的情况下，我们一般用$\mathrm{I}$表示恒等映射，否则，我们使用$\mathrm{Id}$表示恒等映射。
\begin{align*}
    \mathrm{I}: \,\, \mathbb{V} & \to \mathbb{V}  \\
    \mathbf{v} & \mapsto \mathbf{v} 
\end{align*}

\paragraph{缩放映射}
缩放映射指把$\mathbb{V}$中向量映射到$\mathbb{V}$中同一个向量的若干倍的映射。其效果为将每个向量都乘以同一个标量倍数$a$。$a$也称为缩放系数。
我们一般用$\mathrm{H}_a$表示$a$倍缩放映射。
\begin{align*}
    \mathrm{H}_a: \,\, \mathbb{V} & \to \mathbb{V}  \\
    \mathbf{v} & \mapsto a\mathbf{v} 
\end{align*}

以上的映射在任何空间中都存在。下面再来看一些特定空间里的例子：

\paragraph{导数映射} 考虑一元整式空间$\mathbb{F}[z]$到$\mathbb{F}[z]$的映射：
\begin{align*}
    \mathbf{D}: \,\, \mathbb{F}[z] & \to \mathbb{F}[z]  \\
    p & \mapsto \mathbf{D}p 
\end{align*}
其中$\mathbf{D}p$是整式$p$的导数。$\mathbf{D}$是直映射。
比如，$p(z) = 1 - 3z + 2z^2$，则$\mathbf{D}p(z) = -3 + 4z.$

\paragraph{积分映射} 考虑一元整式空间$\mathbb{F}[z]$到$\mathbb{F}[z]$的映射：
\begin{align*}
    \mathcal{I}: \,\, \mathbb{F}[z] & \to \mathbb{F}[z]  \\
    p & \mapsto \mathcal{I}p 
\end{align*}
其中$\mathcal{I}p$是整式$p$的积分。$\mathcal{I}$是直映射。
比如，$p(z) = 1 - 4z + 3z^2$，则$\mathcal{I}p(z) = z - 2z^2+ z^3.$

\paragraph{定积分映射} 考虑区间$[a,b]$上可积函数的集合$\mathbf{L}^1([a,b])$到$\mathbb{R}$的映射：
\begin{align*}
    \mathcal{I}_{[a,b]}: \,\, \mathbf{L}^1([a,b]) & \to \quad\mathbb{R}  \\
    f\quad & \mapsto \int_a^b f(t)\mathrm{d}t 
\end{align*}
$\mathcal{I}_{[a,b]}$是直映射。这是从另一个角度描述积分的直性。

\paragraph{整式乘法映射} 给定整式$q\in\mathbb{F}[z]$，考虑$\mathbb{F}[z]$到$\mathbb{F}[z]$的映射：
\begin{align*}
    \mathbf{C}_{q}: \,\, \mathbb{F}[z] & \to \mathbb{F}[z]  \\
    p\quad & \mapsto qp 
\end{align*}
$\mathbf{C}_{q}$是直映射。

\paragraph{整式插值映射} 给定$m$元数组：$\mathcal{S} = (a_1, a_2, \cdots, a_m)$，考虑$\mathbb{F}[z]$到$\mathbb{F}^m$的映射：
\begin{align*}
    \mathbf{I}_{\mathcal{S}}: \,\, \mathbb{F}[z] & \to \qquad\qquad\mathbb{F}^m  \\
    p\quad & \mapsto \left(p(a_1), p(a_2), \cdots, p(a_m)\right) 
\end{align*}
$\mathbf{I}_{\mathcal{S}}$是直映射。这个映射将整式和它在有限个点上的取值联系起来，在处理插值问题的时候经常用到。

\paragraph{数列的移位映射} 给定自然数$m$，考虑所有数列构成的空间$\mathbb{F}^{\mathbb{N}}$到$\mathbb{F}^{\mathbb{N}}$的映射：
\begin{align*}
    \mathbf{T}_m: \,\, \qquad\qquad\mathbb{F}^{\mathbb{N}}\qquad\qquad & \to \quad\quad \mathbb{F}^{\mathbb{N}}  \\
    (a_1, \cdots, a_m, a_{m+1}, \cdots ) & \mapsto (a_{m+1}, a_{m+2}, \cdots ) 
\end{align*}
$\mathbf{T}_m$是直映射。这个映射把数列中每个数往前移动$m$（前$m$个数去掉）。同理有向后移位的映射，也是直映射。

\section{直映射的代数性质}
给定空间$\mathbb{V}$和$\mathbb{W}$，$\mathcal{L}(\mathbb{V},\mathbb{W})$是否拥有某种代数结构呢？
我们知道，代数结构由配备的运算法则决定。我们给$\mathcal{L}(\mathbb{V},\mathbb{W})$配备两种运算：加法和数乘运算。
\begin{enumerate}
    \item $\forall A, B \,\, \in \,\, \mathcal{L}(\mathbb{V},\mathbb{W})$，定义映射$A+B$： $\forall u \,\, \in \mathbb{V}, \,\,\,\, (A+B)(u) = A(u) + B(u).$
    \item $\forall A \,\, \in \,\, \mathcal{L}(\mathbb{V},\mathbb{W}),  \,\,\, \forall k \in \mathbb{F}$，定义映射$kA$：$\forall u \,\, \in \mathbb{V}, \,\,\,\, (kA)(u) = k\cdot A(u).$
\end{enumerate}
可以验证，$\mathcal{L}(\mathbb{V},\mathbb{W})$配备了这两种运算后，构成一个平直空间。

如果出发空间$\mathbb{V}$和到达空间$\mathbb{W}$是同一个空间，除了以上两种运算，我们还可以为$\mathcal{L}(\mathbb{V})$配备另一种运算：映射的复合。
$\forall A, B, \,\, \in \,\, \mathcal{L}(\mathbb{V})$，定义其复合映射$A\circ B$为：$\forall u \in \mathbb{V}, \,\,\, (A\circ B)(u) = A(B(u)).$

我们可以验证，$\mathcal{L}(\mathbb{V})$关于映射的加法运算和复合运算构成一个\textbf{环}\footnote{见定义\ref{df:0_2}.}。环的乘法幺元是恒等映射$\mathrm{I}$。
复合运算并不满足交换律，所以这个环并不是交换环。

此外，映射的复合运算关于加法运算是\textbf{双平直}（\textbf{双直}）的：
\begin{align*}
    \forall A, B, C \,\, \in \,\, \mathcal{L}(\mathbb{V}), \,\, &\forall k, l \,\, \in \,\, \mathbb{F},  \\
    A\circ(k\cdot B+l\cdot C) &= k\cdot A\circ B + l\cdot A\circ C,  \\
    (k\cdot A+l\cdot B)\circ C &= k\cdot A\circ C + l\cdot B\circ C. 
\end{align*}

把这三种运算都算上，我们称$\mathcal{L}(\mathbb{V})$关于映射的加法、数乘和复合运算构成一个\textbf{域上场}，或\textbf{域上的结合代数}。

\section{有限维空间的直映射}
我们从有限维空间到有限维空间的映射开始，对直映射进行研究。
设出发空间$\mathbb{V}$和到达空间$\mathbb{W}$的维数分别是$n$和$m$，
设$T$是$\mathcal{L}(\mathbb{V}, \mathbb{W})$中元素，
给定空间$\mathbb{V}$的一组基$\mathcal{B}_{\mathbb{V}} = (\mathbf{e}_1,\mathbf{e}_2,\cdots,\mathbf{e}_n)$，按照定义，$\mathbb{V}$中每个元素$\mathbf{u}$都可以唯一写成：
$$ \mathbf{u} = u_1\mathbf{e}_1+u_2\mathbf{e}_2+\cdots+u_n\mathbf{e}_n.$$
所以，
$$T\mathbf{u} = u_1T\mathbf{e}_1+u_2T\mathbf{e}_2+\cdots+u_nT\mathbf{e}_n.$$
这说明，$T$的性质可以完全由$(\mathbf{e}_1,\mathbf{e}_2,\cdots,\mathbf{e}_n)$和$(T\mathbf{e}_1,T\mathbf{e}_2,\cdots,T\mathbf{e}_n)$确定。
只要知道了这两者，我们就可以求出$\mathbb{V}$中任一向量在映射$T$下的结果。

\begin{pp}\label{pp:3_0}
    有限维空间上的直映射由它在一组基上的像确定。
\end{pp}

为了描述直映射在$\mathcal{B}_{\mathbb{V}}$上的像，我们设有$\mathbb{W}$的一组基：$\mathcal{B}_{\mathbb{W}} = (\mathbf{f}_1,\mathbf{f}_2,\cdots,\mathbf{f}_m)$。
那么，直映射在$\mathcal{B}_{\mathbb{V}}$上的像必然可以表示为$\mathcal{B}_{\mathbb{W}}$中元素的直组合。

考虑一个最基本的例子。当$\mathbb{W}$是有序$n$元组$(x_1, x_2, \cdots, x_n)$构成的空间$\mathbb{F}^n$，
映射：
\begin{align*}
    \mathrm{I}_{\mathbb{V}}^{\mathbb{F}^n} : \mathbb{V} &\rightarrow \mathbb{F}^n  \\
    u_1\mathbf{e}_1+u_2\mathbf{e}_2+\cdots+u_n\mathbf{e}_n &\mapsto (u_1, u_2, \cdots , u_n). 
\end{align*}
就是$\mathbb{V}$和$\mathbb{F}^n$之间的直同态。它把基向量$\mathbf{e}_i$映射到$\mathbb{F}^n$的标准基里的基向量：
$$ (0, \cdots, 0, 1, 0, \cdots, 0)\quad\mbox{（第i个位置为1，其余为0）}.$$
它的逆映射：
\begin{align*}
    \mathrm{I}_{\mathbb{F}^n}^{\mathbb{V}} : \mathbb{F}^n &\rightarrow \mathbb{V}  \\
    (u_1, u_2, \cdots , u_n) &\mapsto u_1\mathbf{e}_1+u_2\mathbf{e}_2+\cdots+u_n\mathbf{e}_n. 
\end{align*}
也存在。这说明$\mathrm{I}_{\mathbb{V}}^{\mathbb{F}^n}$是直同构。

到这里，我们就对有限维平直空间完成了第一阶段的探索：
有限维平直空间的“大小”可以用维数来表示。
$n$维平直空间可以表示成$n$条相互“独立”的“直线”的直和，并且和$\mathbb{F}^n$直同构。

所以，以后我们谈论有限维平直空间的时候，可以直接把$\mathbb{F}^n$作为例子来思考。
适用于$\mathbb{F}^n$的命题自然也适用于任何$n$维平直空间。
反过来，讨论任何$n$维平直空间$\mathbb{V}$及其中向量的时候，也可以通过$\mathrm{I}_{\mathbb{V}}^{\mathbb{F}^n}$将其与$\mathbb{F}^n$对应起来。
对于$\mathbb{V}$的基底$\mathcal{B} = (\mathbf{e}_1, \mathbf{e}_2, \cdots , \mathbf{e}_n)$和向量：
$$\mathbf{u} = u_1\mathbf{e}_1+u_2\mathbf{e}_2+\cdots+u_n\mathbf{e}_n,$$
我们可以把$\mathbf{u}$等同于$\mathbb{F}^n$中的向量$(u_1, u_2, \cdots , u_n)$。
当$n=2,3$的时候，这就是我们熟悉的坐标。我们把$(u_1, u_2, \cdots , u_n)$称为
$\mathbf{u}$在$\mathcal{B}$上的\textbf{坐标}，或者说它是向量$\mathbf{u}$的\textbf{坐标形式}。

接下来，我们讨论直映射构成的空间。考虑这样的直映射$E_{i,j}$：给定$1 \leq i \leq n, \,\, \, 1 \leq j \leq m$，$ E_{i,j} (\mathbf{v}_i)= \mathbf{w}_j$，
而对任意$k \neq i$，$E_{i,j}(\mathbf{v}_k) = \mathbf{0}.$\footnote{读者可以验证这样的直映射$E_{i,j}$唯一存在。}

这样定义的$E_{i,j}$合共有$n\times m$个。而这$n\times m$个直映射在直映射构成的平直空间$\mathcal{L}(\mathbb{V},\mathbb{W})$中是是直无关的，
并且可以生成整个空间。这说明集合$\{E_{i,j}\}_{1 \leq i \leq n,  1 \leq j \leq m}$构成$\mathcal{L}(\mathbb{V},\mathbb{W})$的一组基，
$\mathcal{L}(\mathbb{V},\mathbb{W})$的维数是$n\times m$。

最后，我们也可以通过子空间来研究有限维空间上直映射的性质。我们定义直映射限制在子空间上的限制：
\begin{df}\label{df:3_4}
    给定平直空间$\mathbb{V}$、$\mathbb{W}$，直映射$T \in \mathcal{L}(\mathbb{V},\mathbb{W})$及$\mathbb{V}$的子空间$\mathbb{U}$，
    $T$在子空间$\mathbb{U}$上的\textbf{限制}是满足：
    $$ \forall \mathbf{u} \in \mathbb{U}, \,\,\, S\mathbf{u} = T\mathbf{u}. $$
    的直映射$S \in \mathcal{L}(\mathbb{U},\mathbb{W})$。一般将其记作$\restr{T}{\mathbb{U}}.$
\end{df}

\begin{pp}\label{pp:3_1}
    设有限维空间$\mathbb{V}$为其子空间的直和：
    $$ \mathbb{V} = \mathbb{V}_1 \oplus \mathbb{V}_2 \oplus \cdots \oplus \mathbb{V}_m, $$
    给定这些子空间上的一组直映射$\{T_i:\mathbb{V}_i\mapsto\mathbb{W}\}_{i\in\left[1..m\right]}$，则唯一存在$T\in\mathcal{L}(\mathbb{V},\mathbb{W})$，使得：
    $$ \forall i\in[1..m],\,\,\, \restr{T}{\mathbb{V}_i} = T_i.$$

\end{pp}
\begin{proof} % [{\bfseries 证明\ref{pp:3_1}}]
    根据直和定义，$\mathbb{V}$中任一向量$\mathbf{u}$都可以唯一表示为：
    $$ \mathbf{u} = a_1\mathbf{u}_1 + a_2\mathbf{u}_2 + \cdots + a_m\mathbf{u}_m,\,\,\, \forall i\in[1..m],\,\,a_i\in\mathbb{F},\,\mathbf{u}_i\in\mathbb{V}_i.$$
    若有$T$满足题设条件，则$\forall i \in[1..m],\,\,\,T(\mathbf{u}_i) = T_i(\mathbf{u}_i)$.\\
    所以
    \begin{align*}
        T(\mathbf{u}) &= a_1T(\mathbf{u}_1) + a_2T(\mathbf{u}_2) + \cdots + a_mT(\mathbf{u}_m)  \\
                      &= a_1T_1(\mathbf{u}_1) + a_2T_2(\mathbf{u}_2) + \cdots + a_mT_m(\mathbf{u}_m). 
    \end{align*}
    这说明，如果$T$存在，则必然是唯一的。下面证明存在性：

    构造以下映射：
    \begin{align*}
        T: \,\, \mathbb{V} \,\,& \to \,\,\mathbb{W}  \\
        \mathbf{u} \,\,& \mapsto \,\, a_1T_1(\mathbf{u}_1) + a_2T_2(\mathbf{u}_2) + \cdots + a_mT_m(\mathbf{u}_m) 
    \end{align*} 
    由于子空间的和为直和，可以验证映射$T$是直映射，且$\forall i\in[1..m],\,\,\, \restr{T}{\mathbb{V}_i} = T_i.$ 
    这就证明了存在性。

\end{proof}


\section{零空间与像空间}
下面我们来了解与直映射本质相关的两个概念。
\begin{pp}\label{pp:3_2}
    给定从空间$\mathbb{V}$到$\mathbb{W}$的直映射$T$，所有使得$T(\mathbf{u}) = \mathbf{0}$的向量$\mathbf{u}$构成$\mathbb{V}$的子空间，
    称为它的\textbf{零空间}，也称为它的\textbf{核}，记为$\Ker T$。
    零空间的维数称为直映射的\textbf{佚}，记为$\nul T$。
\end{pp}
\begin{proof} % [{\bfseries 证明\ref{pp:3_2}}]
    按照子空间的验证标准：
    \begin{enumerate}
        \item $T(\mathbf{0}) = \mathbf{0}$，所以$\mathbf{0} \in \Ker T$.
        \item $\forall \mathbf{u}, \mathbf{v} \in \Ker T, \,\,\, T(\mathbf{u} + \mathbf{v}) = T(\mathbf{u}) + T(\mathbf{v}) = \mathbf{0}.$ 所以$\mathbf{u} + \mathbf{v} \in \Ker T.$
        \item $\forall \mathbf{u} \in \Ker T,\,\,\, a\in\mathbb{F},\,\,\, T(a\mathbf{u}) = aT(\mathbf{u}) = \mathbf{0}.$ 所以$a\mathbf{u} \in \Ker T.$
    \end{enumerate}
\end{proof}


\begin{pp}\label{pp:3_3}
    给定从空间$\mathbb{V}$到$\mathbb{W}$的直映射$T$，$\mathbb{W}$中所有能写成$\mathbf{v} = T(\mathbf{u}),\,\,\mathbf{u}\in\mathbf{V}$的向量$\mathbf{v}$构成$\mathbb{W}$的子空间，
    称为它的\textbf{像空间}，记为$\Img T$。
    像空间的维数称为直映射的\textbf{容}或\textbf{秩}，记为$\rang T$。
\end{pp}
\begin{proof} % [{\bfseries 证明\ref{pp:3_3}}]
    按照子空间的验证标准：
    \begin{enumerate}
        \item $T(\mathbf{0}) = \mathbf{0}$，所以$\mathbf{0} \in \Img T$.
        \item $\forall \mathbf{u}_1, \mathbf{u}_2 \in \Img T, \,\,\, \exists T(\mathbf{v}_1) = \mathbf{u}_1,\,\,T(\mathbf{v}_2) = \mathbf{u}_2$ 所以$\mathbf{u}_1 + \mathbf{u}_2 = T(\mathbf{v}_1) + T(\mathbf{v}_2) = T(\mathbf{v}_1 + \mathbf{v}_2) \in \Img T.$
        \item $\forall  a\in\mathbb{F},\,\,\,\mathbf{u} = T(\mathbf{v}) \in \Img T\,\, \Rightarrow \,\, a\mathbf{u} = aT(\mathbf{v}) = T(a\mathbf{v}).$ 所以$a\mathbf{u} \in \Img T.$
    \end{enumerate}
\end{proof}


以上我们定义了两个新概念：零空间和像空间。
零空间是$\mathbb{W}$中零元素关于$T$的原像集合：$\Ker T = T^{-1}(\{\mathbf{0}\})$。
像空间是$\mathbb{V}$关于$T$的像的集合：$\Img T = T(\mathbb{V})$。

这两个概念对于理解直映射$T$的性质有重要的帮助。像空间比较好理解，
打个比方来说，像空间告诉你直映射在到达空间$\mathbb{W}$里闹出了多大的“动静”，有多少$\mathbb{W}$中的元素被“牵扯”进来了。
零空间则从另一个不那么直观的角度告诉你这件事：有多少出发空间$\mathbb{V}$里的元素被$T$“默默掩盖”了，映射成零向量，
参与不到$T$在$\mathbb{W}$闹出的“动静”里头去。

那么，零空间和像空间有什么关联呢？以下这个定理揭露了它们和出发空间维数之间的关系：
\begin{pp}\textbf{容佚定理}\label{pp:3_3_1}\\
    给定从空间$\mathbb{V}$到$\mathbb{W}$的直映射$T$，则：
    $$ \nul T + \rang T = \dim \mathbb{V}. $$
\end{pp}
\begin{proof} % [{\bfseries 证明\ref{pp:3_3_1}}]
    设$\nul T = m$，考虑$\Ker T$的一组基：
    $$ \mathcal{B}_k = \{\mathbf{u}_1, \mathbf{u}_2, \cdots, \mathbf{u}_m\}$$
    我们将它扩充为$\mathbb{V}$的一组基：
    $$ \mathcal{B} = \{\mathbf{u}_1, \mathbf{u}_2, \cdots, \mathbf{u}_m, \mathbf{v}_1, \mathbf{v}_2, \cdots, \mathbf{v}_n\}$$
    $(\mathbf{v}_1, \mathbf{v}_2, \cdots, \mathbf{v}_n)$生成的空间$\mathbb{U}$是$\Ker T$的补空间：
    $$ \mathbb{V} = \Ker T \oplus \mathbb{U}.$$
    这时$\dim \mathbb{V} = m+n.$ 我们只需证明$\rang T = n.$\\
    考虑这组向量：$(T\mathbf{v}_1, T\mathbf{v}_2, \cdots T\mathbf{v}_n).$ 
    我们证明它是$\Img T$的基。\\
    首先，$\forall \mathbf{u} = T\mathbf{v} \in \Img T,$
    \begin{align*}
        \mathbf{u} &= T\mathbf{v}  \\
                   &= T(a_1\mathbf{u}_1 + a_2\mathbf{u}_2 + \cdots + a_m\mathbf{u}_m + b_1\mathbf{v}_1 + b_2\mathbf{v}_2 + \cdots + b_n\mathbf{v}_n)  \\
                   &= b_1T(\mathbf{v}_1) + b_2T(\mathbf{v}_2) + \cdots + b_nT(\mathbf{v}_n) 
    \end{align*}
    其中$\mathbf{v} = a_1\mathbf{u}_1 + a_2\mathbf{u}_2 + \cdots + a_m\mathbf{u}_m + b_1\mathbf{v}_1 + b_2\mathbf{v}_2 + \cdots + b_n\mathbf{v}_n$是$\mathbf{v}$的$\mathcal{B}$分解。
    而由于$\mathbf{u}_1, \mathbf{u}_2, \cdots, \mathbf{u}_m$在零空间中，所以$T(a_1\mathbf{u}_1 + a_2\mathbf{u}_2 + \cdots + a_m\mathbf{u}_m) = \mathbf{0}.$
    于是，我们推出$(T\mathbf{v}_1, T\mathbf{v}_2, \cdots T\mathbf{v}_n).$可以生成$\Img T.$\\
    再证明$(T\mathbf{v}_1, T\mathbf{v}_2, \cdots T\mathbf{v}_n)$直无关。\\
    假设存在$c_1, c_2, \cdots, c_n$使得$c_1T(\mathbf{v}_1) + c_2T(\mathbf{v}_2) + \cdots + c_nT(\mathbf{v}_n) = \mathbf{0}$，这说明
    $\mathbf{v} = c_1\mathbf{v}_1 + c_2\mathbf{v}_2 + \cdots + c_n\mathbf{v}_n\in\Ker T.$ 而$\mathbf{v}$显然属于$\mathbb{U}$，所以$\mathbf{v}\in\Ker T \cap \mathbb{U} = \{\mathbf{0}\}.$
    这说明$c_1=c_3=\cdots=c_n=0$，因为$(\mathbf{v}_1, \mathbf{v}_2, \cdots, \mathbf{v}_n)$是直无关的。
    这就证明了$(T\mathbf{v}_1, T\mathbf{v}_2, \cdots T\mathbf{v}_n)$直无关，因而是$\Img T$的基。\\
    $\rang T = n$，所以
    $$ \nul T + \rang T = m + n = \dim \mathbb{V}. $$
\end{proof}

这个定理告诉我们，直映射的零空间和像空间虽然“分处两地”，但它们是“一体两面”的关系。
直映射把出发空间$\mathbb{V}$的“一部分”投射到$\mathbb{W}$中，成为像空间，另一部分“掩盖”起来，是为零空间。
零空间越小，像空间越大；零空间越大，像空间越小。它们“合并”起来，总是和$\mathbb{V}$“一样大”。

这里我们所说的“大小”可以用维数衡量。比如，我们还可以得出：
\begin{enumerate}
    \item 像空间的维数不可能大于出发空间的维数。
    \item 若直映射是单射，则出发空间的维数必然小于等于到达空间的维数\footnote{零空间的维数是零，像空间的维数等于出发空间的维数。}。
    \item 若直映射是满射，则出发空间的维数必然大于等于到达空间的维数\footnote{到达空间维数等于像空间的维数。}。
    \item 若直映射是双射，则出发空间的维数必然等于到达空间的维数。
    \item 若直映射是双射，则它将出发空间的基映射为到达空间的基（从 \ref{pp:3_3} 的证明可知）；反之亦然。
    \item 若直映射的出发空间的维数等于到达空间的维数，则直映射是单射等价于它是满射，也等价于它是双射。
\end{enumerate}
如果直映射是双射，则称其为\textbf{可逆映射}\footnote{全称应为“可逆直映射”，在这里我们省略“直”字。}。
$T$的逆映射是使得$T\circ S = \mathrm{I}_{\mathbb{W}}$，$S\circ T = \mathrm{I}_{\mathbb{V}}$的映射$S$，
一般记为$T^{-1}$。由上可知，可逆映射的出发空间和到达空间维度相同。

\begin{pp}\label{pp:3_3_2}
    直映射的逆映射如果存在，也是直映射。
\end{pp}
\begin{proof} % [{\bfseries 证明\ref{pp:3_3_2}}]
    设从空间$\mathbb{V}$到$\mathbb{W}$的直映射$T$是可逆映射，存在映射$S$使得
    \begin{align*}
        \forall \mathbf{v} \in \mathbb{V} ,& \quad S(T(\mathbf{v})) = \mathbf{v},   \\
        \forall \mathbf{w} \in \mathbb{W} ,& \quad T(S(\mathbf{w})) = \mathbf{w} 
    \end{align*}
    那么，$\forall \mathbf{w}_1, \mathbf{w}_2 \in \mathbb{W}$，$\forall \lambda \in \mathbb{F}$，
    存在$\mathbf{v}_1, \mathbf{v}_2 \in \mathbb{V}$，使得$T(\mathbf{v}_1) = \mathbf{w}_1$，$T(\mathbf{w}_2) = \mathbf{v}_2$。
    因此：
    \begin{align*}
        S(\mathbf{w}_1 + \lambda\mathbf{w}_2) &= S(T(\mathbf{v}_1) + \lambda T(\mathbf{v}_2))  \\
        &= S(T(\mathbf{v}_1 + \lambda \mathbf{v}_2))  \\
        &= \mathbf{v}_1 + \lambda \mathbf{v}_2  
    \end{align*}
    另一方面，$T(S(\mathbf{w}_1)) = \mathbf{w}_1$，而$T$是双射，因此是单射。所以$S(\mathbf{w}_1) = \mathbf{v}_1$. 
    同理，$S(\mathbf{w}_2) = \mathbf{v}_2$。故而
    $$ S(\mathbf{w}_1 + \lambda\mathbf{w}_2) = \mathbf{v}_1 + \lambda \mathbf{v}_2 = S(\mathbf{w}_1) + \lambda S(\mathbf{w}_2).$$
    这说明$S$是一个直映射。
\end{proof}


我们还可以把直映射的逆映射的条件放宽：
\begin{pp}\label{pp:3_3_3}
    设有从$\mathbb{V}$到$\mathbb{W}$的直同构$T$。如果存在$\mathbb{W}$到$\mathbb{V}$的映射$S$，
    使得$T\circ S = \mathrm{I}_{\mathbb{W}}$，那么$S$也是直同构，$S\circ T = \mathrm{I}_{\mathbb{V}}$。
\end{pp}
\begin{proof} % [{\bfseries 证明\ref{pp:3_3_3}}]
    $T$是直同构，说明$\dim \mathbb{V} = \dim \mathbb{W}$.\\
    根据容佚定理，$\rang S \leqslant \dim \mathbb{W}$，所以$\dim T(\Img S) \leqslant \rang S \leqslant \dim \mathbb{W}$。
    也就是说，$\rang (T\circ S) \leqslant \dim \mathbb{W}$。而$T\circ S$是恒等映射，所以上述的不等号都是等号：$\rang S = \dim \mathbb{W}$，$\nul S = 0$，
    $S$是单射，$S$也是满射，故是双射。\\
    对任意$\mathbf{v} \in \mathbb{V}$，存在$\mathbf{w} \in \mathbb{W}$使得$S(\mathbf{w}) = \mathbf{v}$。所以，
    $$ S(T(\mathbf{v})) = S(T(S(\mathbf{w}))) = S(\mathbf{w}) = \mathbf{v}. $$
    即$S\circ T = \mathrm{I}_{\mathbb{V}}$。\\
    最后，由命题\ref{pp:3_3_2}可知$S$也是直映射。所以，$S$是直同构。

\end{proof}


由对称性可知，以上结论中，$S\circ T = \mathrm{I}_{\mathbb{V}}$和$T\circ S = \mathrm{I}_{\mathbb{W}}$可以互换。
从两者之一即可推出另外一个。

\section{投影映射和倒影映射}

本节我们研究两类和互补子空间密切相关的直映射。
\begin{df}\label{df:3_2}
    设$\mathbb{V} = \mathbb{U} \oplus \mathbb{W}$，我们知道：
    $\forall \mathbf{x}\in \mathbb{V}$，$x$可以用唯一的方式写成$\mathbb{U}$和$\mathbb{W}$中元素的和：
    $$ \mathbf{x} = \mathbf{x}_1 + \mathbf{x}_2, \,\,\mathbf{x}_1\,\in\mathbb{U},\,\,\mathbf{x}_2\,\in\mathbb{W}.$$
    我们称直映射：
    \begin{align*}
        p_{\mathbb{U}}: \,\, \mathbb{V} & \to \mathbb{V}  \\
        \mathbf{x} & \mapsto \mathbf{x}_1 
    \end{align*}    
    为$\mathbb{U}$上平行于$\mathbb{W}$的\textbf{投影映射}。
\end{df}

我们来看一个例子：
设$\mathcal{D}_x = \{(x,0)|x\in\mathbf{R}\}$，$\mathcal{D}_y = \{(0,y)|y\in\mathbf{R}\}$，
它们分别是平面上的$x$轴和$y$轴所在直线。它们都是子空间，而且
我们知道$\mathbb{R}^2 = \mathcal{D}_x \oplus \mathcal{D}_y.$

我们把$\mathcal{D}_x$上平行于$\mathcal{D}_y$的投影映射简记为$p_x$。
设有$\mathbf{u} = (x_0, y_0) \in\mathbb{R}^2$，则$p_x(\mathbf{u}) = (x_0,0)$.
因为$\mathbf{u}$唯一分解为：
$$\mathbf{u} = (x_0,0) + (0,y_0), \,\,(x_0,0)\in\mathcal{D}_x,\,\,(0,y_0)\in\mathcal{D}_y.$$
我们可以看到，$p_x$将每个点映射到它在$x$轴上的投影点上，故名投影映射。

我们把$\mathcal{D}_y$换成$\mathcal{D}_{3} = \{(x, y) | y - 3x = 0\}.$
这也是一个子空间，而且$\mathbb{R}^2 = \mathcal{D}_x \oplus \mathcal{D}_3.$
设有$\mathbf{u} = (x_0, y_0) \in\mathbb{R}^2$，则$p_x(\mathbf{u}) = (x_0-\frac{y_0}{3},0)$.
这是因为$\mathbf{u}$唯一分解为：
$$\mathbf{u} = (x_0-\frac{y_0}{3},0) + (\frac{y_0}{3},y_0), \,\,(x_0-\frac{y_0}{3},0)\in\mathcal{D}_x,\,\,(\frac{y_0}{3},y_0)\in\mathcal{D}_3.$$
注意到这里$p_x$将每个点映射到它在$x$轴上的投影点上，但这个投影不再是垂直的，而是沿着直线$y - 3x = 0$的方向。
上一个例子里面我们轻易地就认出了“投影”的性质，是因为垂直投影就是沿着$y$轴的方向进行投影。
这也是我们说“$p_x$是平行于$\mathcal{D}_{3}$的投影映射”的含义。

用简单的说法，我们沿着向量指向的点$(x_0, y_0)$出发作$y - 3x = 0$的平行线，和$\mathcal{D}_x$交于一点：$(x_0-\frac{y_0}{3},0).$
这点就是$(x_0, y_0)$平行于$\mathcal{D}_3$投影在$\mathcal{D}_x$上的投影点。

投影映射把我们熟悉的投影关系，推广到一般的平直空间里。
这也让我们对子空间直和的性质有了更直观的了解。

投影映射有什么性质呢？设$p_{\mathbb{U}}$为投影映射，我们可以验证：
$$ \forall \mathbf{x} \in \mathbb{V},\,\,\,p_{\mathbb{U}}(p_{\mathbb{U}}(\mathbf{x})) = p_{\mathbb{U}}(\mathbf{x}).$$
这是因为我们知道，$\mathbb{U}$中元素写成$\mathbb{U}$和$\mathbb{W}$中元素之和的方式就是$\mathbf{x}_1 = \mathbf{x}_1 + \mathbf{0}.$
所以$\forall \mathbf{x}_1\in\mathbb{U},\,\,\, p_{\mathbb{U}}(\mathbf{x}_1) = \mathbf{x}_1.$

我们把具有以上性质的直映射称为\textbf{幂等映射}。那么幂等映射是否就是投影映射呢？
\begin{pp}\label{pp:3_4}
    设$p$为$\mathbb{V}$到$\mathbb{V}$的直映射，则
    $p$为投影映射当且仅当$p\circ p = p.$
\end{pp}
\begin{proof} % [{\bfseries 证明\ref{pp:3_4}}]
    我们刚刚已经证明了，若$p$为投影映射，则$p\circ p = p$。\\
    下面证明，若$p\circ p = p$，则$p$为投影映射。\\
    $$\forall \mathbf{u}\in\mathbb{V}, \,\,\, \mathbf{u} = p(\mathbf{u}) + \mathbf{u} - p(\mathbf{u}).$$
    $p(\mathbf{u})\in\Img p$，$p(\mathbf{u} - p(\mathbf{u})) = p(\mathbf{u}) - p(p(\mathbf{u})) = \mathbf{0}.$
    所以$\mathbf{u} - p(\mathbf{u})\in\Ker p.$
    这说明$\mathbb{V}=\Img p + \Ker p.$\\
    另外，若$\mathbf{u} \in \Img p \cap \Ker p$，$\mathbf{u} \in \Img p\mathbf{u} = p(\mathbf{v}), \,\,\mathbf{v}\in\mathbb{V}.$ 而从$\mathbf{u}\in\Ker p$可推出：
    $$ \mathbf{0} = p(\mathbf{u}) = p(p(\mathbf{v})) = p(\mathbf{v}) = \mathbf{u}.$$
    所以$\Img p \cap \Ker p = \{\mathbf{0}\}.$\\
    这说明$\mathbb{V}=\Img p \oplus \Ker p$，$p$是$\Img p$上的投影映射，平行于$\Ker p$.
\end{proof}


\begin{df}\label{df:3_2_1}
    设$\mathbb{V} = \mathbb{U} \oplus \mathbb{W}$，
    $\forall \mathbf{x}\in \mathbb{V}$，$x$可以用唯一的方式写成$\mathbb{U}$和$\mathbb{W}$中元素的和：
    $$ \mathbf{x} = \mathbf{x}_1 + \mathbf{x}_2, \,\,\mathbf{x}_1\,\in\mathbb{U},\,\,\mathbf{x}_2\,\in\mathbb{W}.$$
    我们称直映射：
    \begin{align*}
        s_{\mathbb{U}}: \,\, \mathbb{V} & \to \mathbb{V}  \\
        \mathbf{x} & \mapsto \mathbf{x}_1 - \mathbf{x}_2 
    \end{align*}    
    为关于$\mathbb{U}$的且平行于$\mathbb{W}$的\textbf{倒影映射}。
\end{df}

我们来看两个例子：
首先仍然是平面上的$x$轴和$y$轴所在直线
$\mathcal{D}_x = \{(x,0)|x\in\mathbf{R}\}$，$\mathcal{D}_y = \{(0,y)|y\in\mathbf{R}\}$，
我们把关于$\mathcal{D}_x$、平行于$\mathcal{D}_y$的倒影映射简记为$s_x$。

设有$\mathbf{u} = (x_0, y_0) \in\mathbb{R}^2$，则$s_x(\mathbf{u}) = (x_0,-y_0)$.
因为$\mathbf{u}$唯一分解为：
$$\mathbf{u} = (x_0,0) + (0,y_0), \,\,(x_0,0)\in\mathcal{D}_x,\,\,(0,y_0)\in\mathcal{D}_y.$$
我们可以看到，$s_x$将每个点映射到它关于$x$轴对称的点上，故名倒影映射。

我们把$\mathcal{D}_y$换成$\mathcal{D}_{3} = \{(x, y) | y - 3x = 0\}.$
这时$s_x(\mathbf{u}) = (x_0-\frac{2y_0}{3},-y_0)$.
这是因为$\mathbf{u}$唯一分解为：
$$\mathbf{u} = (x_0-\frac{y_0}{3},0) + (\frac{y_0}{3},y_0), \,\,(x_0-\frac{y_0}{3},0)\in\mathcal{D}_x,\,\,(\frac{y_0}{3},y_0)\in\mathcal{D}_3.$$
注意到这里$s_x(\mathbf{u})$不再是垂直“倒映”点$\mathbf{u}$，而是沿着直线$y - 3x = 0$的方向进行“倒映”。
这时候$s_x(\mathbf{u})$更像是点对称和轴对称的结合。
可以验证，$(x_0,y_0)$和$(x_0-\frac{2y_0}{3},-y_0)$关于$(x_0-\frac{y_0}{3},0)$对称。
它们所在直线是$y - 3x = y_0 - 3x_0$，斜率和$\mathcal{D}_{3}$相同。

用简单的说法，我们沿着向量指向的点$P = (x_0, y_0)$出发作$y - 3x = 0$的平行线，
和$\mathcal{D}_x$交于一点：$M = (x_0-\frac{y_0}{3},0).$
$(x_0, y_0)$关于点$M$的对称点$Q$就是$(x_0, y_0)$在倒影映射下的对称点。
$M$是$PQ$的中点，$\mathcal{D}_x$是三角形$POQ$关于顶点$O$的中线$OM$所在的直线。

倒影映射把平面中的对称关系，推广到一般的平直空间里。它保持了对称关系的一般性质。
比如，平面中，轴对称点的轴对称点就是原来的点。在一般平直空间中，设$s_{\mathbb{U}}$为倒影映射，我们可以验证：
$$ \forall \mathbf{x} \in \mathbb{V},\,\,\,s_{\mathbb{U}}(s_{\mathbb{U}}(\mathbf{x})) = \mathbf{x}.$$
这是因为我们知道，$\mathbb{V}$中元素写成$\mathbb{U}$和$\mathbb{W}$中元素之和的方式是唯一的。
$$s_{\mathbb{U}}(\mathbf{x}) = \mathbf{x}_1 - \mathbf{x}_2,$$
所以$s_{\mathbb{U}}(s_{\mathbb{U}}(\mathbf{x})) = \mathbf{x}_1 - (- \mathbf{x}_2) =  \mathbf{x}_1 + \mathbf{x}_2 = \mathbf{x}.$

$s_{\mathbb{U}}$连续两次作用在一个向量上，就得到原来的向量。我们说，倒影映射是一种\textbf{对合映射}。
对合映射是指作用两次后变为恒等映射的映射。

对称关系是对合的，这一点很好理解：对称点的对称点就是原来的点。
那么，对合映射是否都是倒影映射呢？

\begin{pp}\label{pp:3_5}
    设$s$为$\mathbb{V}$到$\mathbb{V}$的直映射，则
    $s$为倒影映射当且仅当$s\circ s = \mathrm{I}_{\mathbb{V}}$\footnote{这里$\mathrm{I}_{\mathbb{V}}$表示$\mathbb{V}$上的恒等映射：$\mathrm{I}_{\mathbb{V}}(\mathbf{x}) = \mathbf{x}.$}.
\end{pp}
\begin{proof} % [{\bfseries 证明\ref{pp:3_5}}]
    我们刚刚已经证明了，若$s$为倒影映射，则$s\circ s = \mathrm{I}_{\mathbb{V}}$。下面证明，若$s\circ s = \mathrm{I}_{\mathbb{V}}$，则$s$为倒影映射。
    $$\forall \mathbf{u}\in\mathbb{V}, \,\,\, \mathbf{u} = \frac{\mathbf{u} + s(\mathbf{u})}{2} + \frac{\mathbf{u} - s(\mathbf{u})}{2}.$$
    这说明$\mathbb{V}=\Img \frac{\mathrm{I}_{\mathbb{V}} + s}{2} + \Img \frac{\mathrm{I}_{\mathbb{V}} - s}{2}.$

    同时，
    $$\forall \mathbf{u}\in\mathbb{V}, \,\,\, s(\mathbf{u}) = \frac{\mathbf{u} + s(\mathbf{u})}{2} - \frac{\mathbf{u} - s(\mathbf{u})}{2}.$$
    所以，只要证明$\mathbb{V}=\Img \frac{\mathrm{I}_{\mathbb{V}} + s}{2} \oplus \Img \frac{\mathrm{I}_{\mathbb{V}} - s}{2}$，
    我们就证明了$s$是关于子空间$\Img \frac{\mathrm{I}_{\mathbb{V}} + s}{2}$的、平行于$\Img \frac{\mathrm{I}_{\mathbb{V}} - s}{2}$的倒影映射。

    我们证明这两个子空间交集为$\{\mathbf{0}\}$：

    若$\mathbf{u} \in \Img \frac{\mathrm{I}_{\mathbb{V}} + s}{2} \cap \Img \frac{\mathrm{I}_{\mathbb{V}} - s}{2}$，
    则存在$\mathbf{v},\,\mathbf{w} \in \mathbb{V}$，使得$2\mathbf{u} = \mathbf{v} + s(\mathbf{v}) = \mathbf{w} - s(\mathbf{w}).$

    $$ 2s(\mathbf{u}) = s(\mathbf{v}) + s(s(\mathbf{v})) = s(\mathbf{w}) - s(s(\mathbf{w}))$$
    $$ s(\mathbf{v}) + \mathbf{v} = s(\mathbf{w}) - \mathbf{w} = \mathbf{w} - s(\mathbf{w}).$$
    这说明$s(\mathbf{v}) + \mathbf{v} = \mathbf{0}.$ 所以$\mathbf{u} = \frac{s(\mathbf{v}) + \mathbf{v}}{2} = \mathbf{0}.$

    因此，$\Img \frac{\mathrm{I}_{\mathbb{V}} + s}{2} \cap \Img \frac{\mathrm{I}_{\mathbb{V}} - s}{2} = \{\mathbf{0}\}.$

    这说明，$s$是关于$\Img \frac{\mathrm{I}_{\mathbb{V}} + s}{2}$的、平行于$\Img \frac{\mathrm{I}_{\mathbb{V}} - s}{2}$倒影映射。

\end{proof}


这个证明看起来可能有些“绕”。这是因为其叙述方式是“不自然”的。
面对这个命题的时候，我们思考的顺序其实是这样的：
要证明$s$是倒影映射，我们要把$\mathbb{V}$表示成两个子空间$\mathbb{U}$和$\mathbb{W}$的直和，
并且使得：
$$ \forall \mathbf{u} \,\in \mathbb{V},\,\,\, \mathbf{u} = \mathbf{u}_1 + \mathbf{u}_2, \,\, s(\mathbf{u}) = \mathbf{u}_1 - \mathbf{u}_2,\,\,\, \mathbf{u}_1 \in\mathbb{U},\,\,\mathbf{u}_2 \in\mathbb{W}.$$
反推得：
$$ \mathbf{u}_1 = \frac{\mathbf{u} + s(\mathbf{u})}{2},\,\,\, \mathbf{u}_2 = \frac{\mathbf{u} - s(\mathbf{u})}{2}.$$ 
于是，我们猜测这两个子空间就是$\Img \frac{\mathrm{I}_{\mathbb{V}} + s}{2}$和$\Img \frac{\mathrm{I}_{\mathbb{V}} - s}{2}$。
然后我们证明$\mathbb{V}$是这两个空间的直和，就完成了整个证明。

很多时候，解决问题的思路和最后给出的证明方式并不完全一致。
在读懂证明的同时，也应该想得更深一步，想清楚证明背后的解题思路。

投影映射、倒影映射和复数都有密切关联。
考虑复数$z = x+y\imath $，它可以看作是$\mathbb{R}^2$上的一个点。
而$z$的实部就是$z$在$x$轴上的投影，虚部就是$z$在$y$轴上的投影，
$z$的共轭复数$\bar{z}=x-y\imath $则是$z$关于$x$轴的对称点。
所以，相应的变换都可以看作是$\mathbb{R}^2$上的直映射。

\section{直型和对偶空间}
现在我们来了解另一类特殊的直映射。考虑这样一类直映射$T$：
$$  T : \mathbb{V} \rightarrow \mathbb{F}  $$
$T$把$n$维空间$\mathbb{V}$中的向量映射到其系数域$\mathbb{F}$。
这种直映射称为\textbf{直型}。

由于$\mathbb{F}$也可以看作$\mathbb{F}$上的一维空间，所以直型也可以看作是到一维空间，也就是我们理解的“直线”上的直映射。

直型的像空间是一维空间$\mathbb{F}$的子空间，所以要么是$\mathbb{F}$，要么是$\{\mathbf{0}\}$。
后者是平凡的零映射，而前者是我们需要继续了解的。

假设直型不是零映射，那么它的像空间就是$\mathbb{F}$。根据容佚定理，它的零空间维数是$n-1$，比$\mathbb{V}$的维数少1.

用三维空间$\mathbb{R}^3$作为例子。设有映射：
$$T : (x,y,z) \mapsto ax + by + cz$$
其中$a,b,c$是给定的实数。那么，可以验证$T$是一个直型。$T$的像空间是$\mathbb{R}$，而它的零空间是$\{(x,y,z) | ax+by+cz = 0\}$. 
这是一个过原点的二维平面。反之，任意过原点的二维平面都可以写成$\{(x,y,z) | ax+by+cz = 0\}$的形式，从而对应一个$\mathbb{R}^3$上的直型。

对于一般的有限维空间$\mathbb{V}$，我们将$\mathbb{V}$上的直型的零空间称为$\mathbb{V}$的\textbf{超平面}。
超平面是维数比$\mathbb{V}$少1的子空间。

给定一个$n-1$维的子空间$\mathbb{H}$，能否构造一个以它为零空间的直型呢？我们给$\mathbb{H}$找一组基：$\mathbf{e}_1, \mathbf{e}_2, \cdots, \mathbf{e}_{n-1}$。
可以把这组向量补全，成为$\mathbb{V}$的基。补全的方法是增加一个向量$\mathbf{e}_n$。于是，我们构造直型$T$：
$$ \forall 1 \leqslant i < n, \,\,\, T(\mathbf{e}_i)= 0, \quad T(\mathbf{e}_n)= 1. $$
这就完成了构造。我们也可以让$T(\mathbf{e}_n + a\mathbf{e}_1)= 1$或让$T(\mathbf{e}_n + \mathbf{e}_1 - b\mathbf{e}_3)= 1$等等，构造出各种各样的直型。

以上构造说明，超平面的概念和$n-1$维子空间是等价的。所有超平面都是$n-1$维的子空间，所有$n-1$维的子空间都是超平面。

给定$\mathbb{V}$的一组基：$\mathcal{B} = \{\mathbf{e}_1, \mathbf{e}_2, \cdots, \mathbf{e}_{n}\}$，
非零的直型$T$将某个向量$\mathbf{x} = \sum_{i=1}^n x_i \mathbf{e}_i$映射为$\mathbb{F}$中的一个数：
$$ T(\mathbf{x}) = \sum_{i=1}^n x_i T(\mathbf{e}_i)$$
注意，每个$T(\mathbf{e}_i)$都是$\mathbb{F}$中的一个数。所以，记$a_i = T(\mathbf{e}_i)$，就有：
$$ T(\mathbf{x}) = \sum_{i=1}^n a_i x_i$$
$\sum_{i=1}^n a_i x_i$是一个$n$元一次整式。这种整式在古时候被称为齐次一次型或直型，也是直型名称的由来。
$T$对应的超平面就是方程$ \sum_{i=1}^n a_i x_i = 0$的解。

考虑这样的直型$\mathbf{e}_i^*$：它将$\mathbf{e}_i$映射到$1$，将其他的基向量映射到$0$，也就是说：
$$ \mathbf{e}_i^*(\mathbf{e}_j) = \delta(i,j) $$
其中$\delta$称为\textbf{示等函数}\footnote{也称为克罗内克$\delta$函数，是一种示性函数，即用来标示某个性质的函数：满足该性质时函数值为1，否则为0.}。$\delta(i,j)=1$当且仅当$i=j$，$i\neq j$的时候$\delta(i,j)=0$。

给定$\mathbb{V}$中向量$\mathbf{x} = \sum_{i=1}^n x_i \mathbf{e}_i$，$\mathbf{e}_i^*(\mathbf{x}) = x_i$。
也就是说，$\mathbf{e}_i^*$把$\mathbb{V}$中向量映射为它的第$i$个坐标分量。

\begin{pp}\label{pp:3_5_1}
给定$\mathbb{V}$的一组基：$\mathcal{B} = \{\mathbf{e}_1, \mathbf{e}_2, \cdots, \mathbf{e}_{n}\}$，
$\mathbf{e}_1^*, \mathbf{e}_2^*, \cdots , \mathbf{e}_n^*$构成$\mathcal{L}(\mathbb{V}, \mathbb{F})$的基，
称为$\mathcal{B}$的\textbf{对偶基}，记为$\mathcal{B}^*$。
\end{pp}
\begin{proof} % [{\bfseries 证明\ref{pp:3_5_1}}]
$\forall T \in \mathbb{V}$，记$a_i = T(\mathbf{e}_i)$，则由之前讨论知道，$T = \sum_{i=1}^n a_i \mathbf{e}_i^*.$ 
所以，$\mathcal{B}^*$是$\mathcal{L}(\mathbb{V}, \mathbb{F})$的生成集。\\
另一方面，如果有系数$\lambda_1, \lambda_2, \cdots, \lambda_n \in \mathbb{F}$使得直型
$$ \lambda_1 \mathbf{e}_1^* + \lambda_2 \mathbf{e}_2^* + \cdots + \lambda_n \mathbf{e}_n^* = 0,$$
那么，将其作用于$\mathcal{B}$中各个基向量上，就得到：
$$ \forall 1 \leqslant i \leqslant n, \quad 0 = \sum_{j=1}^n \lambda_j \mathbf{e}_j^*(\mathbf{e}_i) = \lambda_i.$$
这说明$\mathcal{B}^*$中元素直无关。\\
综上可得，$\mathcal{B}^*$是$\mathcal{L}(\mathbb{V}, \mathbb{F})$的基。
\end{proof}


这说明，$\mathcal{L}(\mathbb{V}, \mathbb{F})$也是一个$n$为空间。
我们一般将空间$\mathcal{L}(\mathbb{V}, \mathbb{F})$简记为$\mathbb{V}^*$，称为$\mathbb{V}$的\textbf{对偶空间}。
对$\mathbb{V}$的每一组基，我们都可以找到$\mathbb{V}^*$中与它对应的对偶基。
$\mathbb{V}^*$中的每个直型，都可以对应$\mathbb{V}$中某个元素。具体来说，给定基底$\mathcal{B}$，考虑映射：

\begin{align*} 
    \phi : \mathbb{V} &\rightarrow \mathbb{V}^*  \\
    \mathbf{a} = \sum_{i=1}^n a_i \mathbf{e}_i &\mapsto  T_{\mathbf{a}} : \sum_{i=1}^n a_i \mathbf{e}_i^* 
\end{align*}

可以验证，映射$\phi$是直映射，而且是双射。
这说明，$\mathbb{V}$和$\mathbb{V}^*$同构。

$\mathbb{V}^*$是否也有对偶空间呢？我们考察$\mathcal{L}(\mathbb{V}^*, \mathbb{F})$，也就是所有$\mathbb{V}$上的直型到系数域$\mathbb{F}$的直映射，或者说直型的直型构成的空间。
把前面论证中的$\mathbb{V}$替换为$\mathbb{V}^*$，可知$\mathcal{L}(\mathbb{V}^*, \mathbb{F})$和$\mathbb{V}^*$同构，
所以也和$\mathbb{V}$同构。具体可以构造这样的映射：
\begin{align*} 
    \xi : \mathbb{V} &\rightarrow \mathcal{L}(\mathbb{V}^*, \mathbb{F}) \\
    \mathbf{a}  &\mapsto  U_{\mathbf{a}} : T\mapsto T({\mathbf{a}}) 
\end{align*}
可以验证$\xi$也是直映射，而且是双射。我们把空间$\mathcal{L}(\mathbb{V}^*, \mathbb{F})$简记为$\mathbb{V}^{**}$，
则$\mathbb{V}^{**}\cong\mathbb{V}.$

\section{对偶映射}

平直空间$\mathbb{V}$和$\mathbb{V}^* = \mathcal{L}(\mathbb{V}, \mathbb{F})$同构，这说明
$\mathcal{L}(\mathbb{V}, \mathbb{F})$里的直型和$\mathbb{V}$里的向量可以看作是等价的。
给定另一个平直空间$ \mathbb{W}$，可以为直映射$T \in \mathcal{L}(\mathbb{V}, \mathbb{W})$构造一个对偶的映射：
$T^* \in \mathcal{L}(\mathbb{W}^*, \mathbb{V}^*)$。其定义为：
$$ \forall \varphi \in \mathbb{W}^*, \quad T^*(\varphi) = \varphi \circ T $$
具体来说，给定$\varphi$，任取$\mathbf{u} \in \mathbb{V}$，$T(\mathbf{u}) \in \mathbb{W}$，
所以$(\varphi \circ T)(\mathbf{u}) = \varphi(T(\mathbf{u}))\in \mathbb{F}$。
映射的复合与映射的加法满足分配律，所以容易验证$T^*$是直映射。
从更高的角度来说，$\mathcal{D}^{\Theta}:\,\,\, T \mapsto T^*$也是一个直映射，
从而是$\mathcal{L}(\mathbb{V}, \mathbb{W})$到$\mathcal{L}(\mathbb{W}^*, \mathbb{V}^*)$的同态。

下面我们来阐明这一点：

\begin{pp}\label{pp:3_6}
    \begin{itemize}
        \item $\forall S, T \in \mathcal{L}(\mathbb{V}, \mathbb{W}), \quad (S + T)^* = S^* + T^*$
        \item $\forall T \in \mathcal{L}(\mathbb{V}, \mathbb{W}), \,\,\, \lambda \in \mathbb{F}, \quad (\lambda T)^* = \lambda T^*$
        \item $T=0 \Leftrightarrow T^* = 0$
        \item $\forall S \in \mathcal{L}(\mathbb{U}, \mathbb{V}), \,\,\, T \in \mathcal{L}(\mathbb{V}, \mathbb{W}), \quad (T \circ S)^* = S^* \circ T^*$
    \end{itemize}
\end{pp}
\begin{proof} % [{\bfseries 证明\ref{pp:3_6}}]
前两条容易验证：给定$ \varphi \in \mathbb{W}^*$，$(S + T)^*(\varphi) = \varphi \circ (S + T) = \varphi \circ S + \varphi \circ T$，$(\lambda T)^*(\varphi) = \varphi \circ (\lambda T) = \lambda \varphi \circ T $。这两条可以说明，$\mathcal{D}^{\Theta}:\,\,\, T \mapsto T^*$是一个直映射，从而是$\mathcal{L}(\mathbb{V}, \mathbb{W})$到$\mathcal{L}(\mathbb{W}^*, \mathbb{V}^*)$的同态。

对于第三条：若$T = 0$，则$\forall \varphi \in \mathbb{W}^*$，$T^*(\varphi) = \varphi \circ 0 = 0$。所以$T=0 \Rightarrow T^* = 0$。另一方面，若$T^* = 0$，则$\forall \varphi \in \mathbb{W}^*$，$\varphi\circ T = 0$。选定$\mathbb{V}$的基$(\mathbf{e}_1, \cdots , \mathbf{e}_n)$和$\mathbb{W}^*$的基$(\varphi_1, \cdots , \varphi_m)$，
$$ \forall 1 \leqslant i \leqslant n, \,\,\,  1 \leqslant j \leqslant m, \quad \varphi_j (T(\mathbf{e}_i)) = \mathbf{0}.$$
这说明$\forall 1 \leqslant i \leqslant n$，$T(\mathbf{e}_i) = \mathbf{0}$，也就是说，$T$是零映射。

前三条性质说明，$\mathcal{D}^{\Theta}:\,\,\, T \mapsto T^*$是$\mathcal{L}(\mathbb{V}, \mathbb{W})$到$\mathcal{L}(\mathbb{W}^*, \mathbb{V}^*)$的同构映射。

对于第四条：给定$ \varphi \in \mathbb{W}^*$，
\begin{align*}
(T \circ S)^* (\varphi) &= \varphi \circ (T \circ S)  = (\varphi \circ T) \circ S  \\
&= (T^*(\varphi)) \circ S  = S^*(T^*(\varphi))  \\
&= (S^*\circ T^*)(\varphi). 
\end{align*}
从第三条可知，对偶和复合并不满足分配律。

\end{proof}


给定直映射$T \in \mathcal{L}(\mathbb{V}, \mathbb{W})$，它的对偶映射$T^*$的性质与$T$有什么关系呢？
来看两者的像空间与零空间。

设$\varphi \in \Ker T^*$，那么按照定义$\varphi \circ T = 0$。
也就是说，$\forall \mathbf{u} \in \mathbb{V}, \,\,\, \varphi(T(\mathbf{u})) = 0.$ 
反过来，只要后者成立，自然有$\varphi \in \Ker T^*$。于是，我们可以把$T^*$的零空间写成：
$$ \Ker T^* = \{\varphi \in \mathbb{W}^* | \forall \mathbf{v} \in \Img T, \,\, \varphi(\mathbf{v}) = 0 \} = \{\varphi \in \mathbb{W}^* |\Img T \subseteq \Ker \varphi \}. $$
引入这样一个新概念：

\begin{df}\label{df:3_3}
给定$\mathbb{V}$的子空间$\mathbb{U}$，
所有满足$\forall \mathbf{u} \in \mathbb{U} , \,\,\, \varphi(\mathbf{u}) = 0$的直型$\varphi  \in \mathbb{V}^*$的集合，
称为$\mathbb{U}$的\textbf{归零集}，记为$\mathbb{U}^\circ$。
\end{df}
于是，我们可以把结论写为：

\begin{pp}\label{pp:3_7}
$T^*$的零空间是$T$的像空间的归零集：$\Ker T^* = \left(\Img T\right)^\circ$。
\end{pp}
子空间的归零集和子空间本身有什么联系呢？我们有如下性质：

\begin{pp}\label{pp:3_8}
子空间的归零集是平直空间。
\end{pp}
\begin{proof} % [{\bfseries 证明\ref{pp:3_8}}]
如果$\varphi , \gamma  \in \mathbb{V}^*$满足：$\forall \mathbf{u} \in \mathbb{U} , \,\,\, \varphi(\mathbf{u}) = \gamma (\mathbf{u})= 0$，那么
$$ \forall \mathbf{u} \in \mathbb{U} , \,\,\,\forall \lambda \in \mathbb{F}, \,\,\, (\varphi + \lambda \gamma) (\mathbf{u}) = \varphi (\mathbf{u}) + \lambda \gamma(\mathbf{u})  = 0.$$
\end{proof}

子空间的归零集是平直空间。因此，以后在相关之处，我们用“\textbf{归零空间}”的说法代替“归零集”。

\begin{pp}\label{pp:3_9}
如果$\mathbb{U}$是$\mathbb{V}$的子空间，那么$\dim\mathbb{U} + \dim\mathbb{U}^\circ = \dim \mathbb{V}.$
\end{pp}
\begin{proof} % [{\bfseries 证明\ref{pp:3_9}}]
考虑映射$\iota \in \mathcal{L}(\mathbb{U}, \mathbb{V}) : \mathbf{u} \mapsto \mathbf{u}$。它类似于恒等映射，只是出发空间和到达空间不同。考虑$\iota$的对偶映射$\iota^* \in \mathcal{L}(\mathbb{V}^*, \mathbb{U}^*)$。$\iota^*$的零空间就是$U^\circ$：
$$ \forall \varphi \in \mathbb{V}^*, \quad \varphi \in \Ker \iota^* \Leftrightarrow \varphi \circ \iota = 0 \Leftrightarrow \forall \mathbf{u} \in\mathbb{U},\,\,\, \varphi (\mathbf{u}) = \varphi (\iota(\mathbf{u})) = 0.$$
而容佚定理告诉我们，$\rang \iota^* + \dim \mathbb{U}^\circ = \rang \iota^* + \nul \iota^* = \dim \mathbb{V}^* = \dim \mathbb{V}. $ 所以，只需证明$\rang \iota^* = \dim \mathbb{U}$ . 为此，我们证明$\Img \iota^* = \mathbb{U}^*$.

注意到$\Img \iota^*$是$\mathbb{U}^*$的子空间，所以只需证明$\mathbb{U}^* \subseteq \Img \iota^*$. \\
任取$\varphi \in \mathbb{U}^* $，我们构造一个直型$\gamma \in \mathbb{V}^*$，使得$\iota^*(\gamma) = \varphi$。\\
将空间$\mathbb{V}$写成$\mathbb{U}$和它的补空间的直和：$\mathbb{V} = \mathbb{U} \oplus \mathbb{U}^c$，
构造映射$p \in \mathcal{L}(\mathbb{V}, \mathbb{U})$：
$$ \forall  \mathbf{v} \in \mathbb{V}, \quad p(\mathbf{v}) = \mathbf{u} \in \mathbb{U},$$
其中$\mathbf{u}$是向量$\mathbf{v}$关于$\mathbb{U}$和$\mathbb{U}^c$的唯一分解：
$\mathbf{v} = \mathbf{u} + \mathbf{w}$，$\mathbf{u} \in \mathbb{U}$，$\mathbf{w} \in \mathbb{U}^c$.

考虑映射$\gamma = \varphi \circ p$，它是$\mathbb{V}^*$中元素，而且由于映射$p$在$\mathbb{U}$上是恒等映射，
$$ \forall \mathbf{u} \in \mathbb{U}, \quad \gamma(\mathbf{u}) = \varphi(p(\mathbf{u})) = \varphi(u).$$
所以
$$ \forall \mathbf{u} \in \mathbb{U}, \quad (\iota^*(\gamma))(\mathbf{u}) = (\gamma \circ \iota)(\mathbf{u}) = \gamma(\iota(\mathbf{u})) = \gamma(\mathbf{u}) = \varphi(u)。$$
这就说明，$\iota^*(\gamma) = \varphi$. 于是$\mathbb{U}^* \subseteq \Img \iota^*$，也就是说，$\Img \iota^* = \mathbb{U}^*$。
\end{proof}

子空间的归零空间，维数和它的补空间是一样的。所以，直映射$T \in \mathcal{L}(\mathbb{V}, \mathbb{W})$的像空间的归零空间的维数，
是$T$的像空间的补空间的维数：$\dim \left(\Img T\right)^\circ = \dim \mathbb{W} - \rang T$。因此：
$$ \nul T^* = \dim \left(\Img T\right)^\circ = \dim \mathbb{W} - \rang T = \dim \mathbb{W} - \dim \mathbb{V} + \nul T.$$
这就是$T^*$和$T$的佚的关系。另一方面，也可以得到容的关系：
$$ \rang T = \dim \mathbb{W} - \nul T^* =  \rang T^*. $$

对偶映射$T^*$的容与$T$的容相等，因而和$T$的零空间的归零空间的维数相等。
$T^*$的像空间是$\mathbb{V}^*$的子空间，$T$的零空间的归零空间也是$\mathbb{V}^*$的子空间，它们维数相等。
那么，它们是否是同一个子空间呢？
\begin{pp}\label{pp:3_10}
    $T^*$的像空间是$T$的零空间的归零空间：$\Img T^* = \left(\Ker T\right)^\circ$。
\end{pp}
\begin{proof} % [{\bfseries 证明\ref{pp:3_10}}]
两者维数相等，只需证明一个是另一个的子空间。下证$\Img T^* \subseteq \left(\Ker T\right)^\circ$。

对任意$\varphi \in \Img T^*$，按定义存在$\psi \in \mathbb{W}^*$，使得$\varphi = T^*(\psi) = \psi\circ T$ 。因此，
$$\forall \mathbf{v} \in \Ker T , \quad \varphi (\mathbf{v}) = \psi(T(\mathbf{v})) = \psi(\mathbf{0}) = 0.$$
这说明$\varphi \in (\Ker T)^\circ$。这就证明了$\Img T^* \subseteq \left(\Ker T\right)^\circ$。
\end{proof}

如果$T$是单射，那么$\Ker T = \{\mathbf{0}\}$，因此它的归零空间为$\mathbb{V}^*$，
这说明$T^*$的像空间是$\mathbb{V}^*$，$T^*$是满射。反过来，如果$T^*$是满射，也可以推出$T$是单射。

类似地，可以推出$T$是满射当且仅当$T^*$是单射。

\sectionext{直映射与商空间}
前面我们定义了商空间和映射$\tau_\mathbb{U}: \mathbf{v} \mapsto \mathbf{v} + \mathbb{U}$。
可以验证，$\tau_\mathbb{U}$是直映射，因此，根据容佚定理，
$$\nul \tau_\mathbb{U} + \rang \tau_\mathbb{U} = \dim \mathbb{V}.$$
由于$\tau_\mathbb{U}$是满射，$\rang \tau_\mathbb{U} = \dim \mathbb{V}/_\mathbb{U}.$ 
另一方面，$\tau_\mathbb{U}$的零空间是映射为$\mathbb{U}$的元素的集合：$\Ker \tau_\mathbb{U} = \{\mathbf{v} | \mathbf{v} + \mathbb{U} = \mathbb{U}\} $。
显然，如果$\mathbf{v} \in \mathbb{U}$，那么$\mathbf{v} + \mathbb{U} = \mathbb{U}$。
反之，如果$\mathbf{v} + \mathbb{U} = \mathbb{U}$，
那么$\mathbf{v} = \mathbf{v} + \mathbf{0}  \in \mathbf{v} + \mathbb{U} =\mathbb{U}$. 
所以，$\{\mathbf{v} | \mathbf{v} + \mathbb{U} = \mathbb{U}\} $就是$\mathbb{U}$本身。
因此，$\nul \tau_\mathbb{U} = \dim \mathbb{U}$. 这样，上面的等式可以写为：
$$\dim \mathbb{U} + \dim \mathbb{V}/_\mathbb{U} = \dim \mathbb{V}.$$
于是，我们知道了商空间的维数。注意到，商空间的维数和补空间的维数是一样的。

商空间和补空间，都提供了把平直空间按照某个子空间“拆分”的方法。
不同在于，补空间是空间的子空间，而商空间不是空间的子空间。
补空间是把子空间看作空间的部分，提供对剩余部分的描述，两者位于同一层面。
商空间是把子空间看作衡量空间的“单位”，提供对空间的概括描述，把空间分为了不同的层面。

给定直映射$T \in \mathcal{L}(\mathbb{V}, \mathbb{W})$，容佚定理描述了$T$的像空间维度和零空间维度的关系。
下面我们用另一种方法刻画$T$的像空间。考虑映射$\widetilde{T}$：
\begin{align*} 
    \widetilde{T} : \mathbb{V}/_{\Ker T} &\rightarrow \mathbb{W}  \\
    \mathbf{v} + \Ker T &\mapsto T(\mathbf{v}) 
\end{align*}
首先，以上定义把$\mathbf{v} + \Ker T$映射到$T(\mathbf{v})$。这么定义是可行的，因为如果$\mathbf{v} + \Ker T = \mathbf{w} + \Ker T$，
则$\mathbf{v} - \mathbf{w} \in \Ker T$，也就是说，$T(\mathbf{v} - \mathbf{w}) = \mathbf{0}$，即$T(\mathbf{v}) =T(\mathbf{w})$. 
这说明一个$\mathbb{V}/_{\Ker T}$中元素只会映射到一个值。

其次可以验证，这个映射是直映射。此外，它还是单射，因为
$$ T(\mathbf{v}) = \mathbf{0} \Rightarrow \mathbf{v} \in \Ker T \Rightarrow \mathbf{v} + \Ker T = \Ker T. $$
即$\widetilde{T}^{-1}(\{\mathbf{0}\}) = \{\Ker T\}.$

映射$\widetilde{T}$可以看作把空间$\mathbb{V}$以$\Ker T$为“单位”划分的新空间映射到$\mathbb{W})$。
从$\widetilde{T}$的定义可知，它的像空间和$T$的像空间是一样的。
因此，作为单射，$\widetilde{T}$是把$\mathbb{V}/_{\Ker T}$映射到$\Img T$上的双射。
\begin{pp}{\bfseries 直同态基本定理 }\label{pp:3_11}\\
    设$T$是空间$\mathbb{V}$到$\mathbb{W}$的直映射，则
    $$\mathbb{V}/_{\Ker T} \cong \Img T.$$
\end{pp}

直观的理解就是：$T$把$\mathbb{V}$中“一群”向量映射到$\mathbb{W}$中同一个向量。
$\mathbb{W}$中每个向量要么没有原像，要么原像是$\Ker T$“大小”的集合。
集合中任意两个向量相差一个$\Ker T$中的元素，因而不影响它们经过$T$映射为同一个值。

另一个角度来看，$T$定义了$\mathbb{V}$中一个等价关系：
向量$\mathbf{v}$等价于$\mathbf{w}$当且仅当$T\mathbf{v} = T\mathbf{w}$。
商空间$\mathbb{V}/_{\Ker T}$中每个向量都对应一个等价类。

\chapter{矩阵}

前面我们提到，给定有限维空间$\mathbb{V}$的一组基$\mathcal{B} = (\mathbf{u}_1, \mathbf{u}_2, \cdots, \mathbf{u}_n)$，
那么从$\mathbb{V}$到$\mathbb{W}$的直映射$T$就被$\mathcal{B}$在$T$下的像$(T\mathbf{u}_1, T\mathbf{u}_2, \cdots, T\mathbf{u}_n)$完全确定。
本章我们将引入一个新的概念：\textbf{矩阵}。矩阵可以方便有效地表示直映射，让我们更好地了解直映射的性质。

\section{直映射的矩阵} 

什么是矩阵？把若干行数字从上到下排在一起，每行数字个数一样多，就是一个矩阵。
设$m$和$n$是正整数，一个$m\times n$的矩阵就是以下的样子：
$$
A = \left[
\begin{matrix}
 a_{1,1}     & a_{1,2}   & \cdots & a_{1,n} \\
 a_{2,1}     & a_{2,2}   & \cdots & a_{2,n} \\
 \vdots      & \vdots    & \ddots & \vdots  \\
 a_{m,1}     & a_{m,2}   & \cdots & a_{m,n} \\
\end{matrix}
\right]
$$
可以看到$m\times n$个实数排成了“长方形”的“数阵”，所以叫做矩阵。
我们用双下标指示矩阵的元素，第一个下标为元素所在的行数，第二个下标为元素所在的列数。
比如$a_{3,2}$表示位于矩阵第$3$行第$2$列的元素。

本书中，直映射对应的矩阵通常用大写字母表示，而其元素用对应的小写字母表示。
矩阵元素有两种展示方法。一种以数阵的方式展示矩阵元素的位置关系，即如上方公式一样，
另一种简约的方法用于在行文中定义或展示矩阵元素，
比如矩阵$A$也可以定义为$A = [a_{i,j}]_{1\leqslant i\leqslant m, 1\leqslant j\leqslant n}$。

$m$行单列的矩阵有$m \times 1 = m$个元素，单行$n$列的矩阵有$1 \times n = n$个元素。
如果把$m$行单列的矩阵看作有序$m$元数组，所有这样的数组关于数组的加法和数乘运算构成一个$m$维的平直空间。
同理，如果把单行$n$列的矩阵看作有序$m$元数组，所有这样的数组关于数组的加法和数乘运算构成一个$n$维的平直空间。
因此，它们可以用来对应$m$维或$n$维空间中的向量。比如，给定$m$维空间$\mathbb{V}$和一组基底$\mathcal{B} = \left(\mathbf{e}_i\right)_{1\leqslant i\leqslant m}$，
那么$\mathbb{V}$中某个向量$\mathbf{v}$可以唯一写成：
$$ \mathbf{v} = u_1\mathbf{e}_1 + u_2\mathbf{e}_2 + \cdots + u_m\mathbf{e}_m.$$
于是它可以写成$m$行单列的矩阵：
$$
\begin{bmatrix}
    u_{1} \\
    u_{2} \\
    \vdots  \\
    u_{m}  \\
   \end{bmatrix}
$$
我们称它为\textbf{列向量}，记为$[\mathbf{v}]_{\mathcal{B}}$。
列向量的元素就是向量的坐标，所以，列向量其实是向量坐标形式以外的又一种表达方式。

当然，我们也可以把向量写成单行$m$列的矩阵：
$$
\begin{bmatrix}
    u_{1} & u_{2} &\cdots & u_{m}
   \end{bmatrix}
$$
我们称这种写法为\textbf{行向量}。

一个$m$行$n$列的矩阵有$m\times n$个元素，如果将它看作是$m\times n$元数组，
所有这样的数组关于数组的加法和数乘运算构成一个$m\times n$维的平直空间。
而我们在前一章提到，
$n$维空间到$m$空间的所有直映射构成一个$m\times n$维的空间：$\mathcal{L}(\mathbb{V},\mathbb{W})$。
两者维数都是$m\times n$，“大小”一样，这让我们想到，可以用矩阵对应直映射。

矩阵如何表示直映射呢？我们首先定下空间$\mathbb{V}$和$\mathbb{W}$中各一组基：
$\mathcal{B}_{\mathbb{V}} = (\mathbf{e}_1, \mathbf{e}_2, \cdots, \mathbf{e}_n)$、
$\mathcal{B}_{\mathbb{W}} = (\mathbf{f}_1, \mathbf{f}_2, \cdots, \mathbf{f}_m)$.
对$\mathcal{B}_{\mathbb{V}}$中任一元素$\mathbf{e}_i$，$T\mathbf{e}_i$在$\mathcal{B}_{\mathbb{W}}$下有唯一分解，
设其为：
$$ T\mathbf{e}_i = t_{1,i}\mathbf{f}_1 + t_{2,i}\mathbf{f}_2 + \cdots + t_{m,i}\mathbf{f}_m.$$
对每个$1\leq i \leq n$，我们都有一个有序$m$元组：
$$ (t_{1,i}, t_{2,i}, \cdots, t_{m,i}) $$ 
将这$n$个数组竖着从左到右排起来，就成了一个矩阵\footnote{从这里可以看到，矩阵的内容和基底向量的顺序是有关系的，不能把基底当成普通的集合看待。}.
\begin{align*}
  \quad  & \quad \,\, \mathbf{e}_1 \quad \, \cdots \quad \mathbf{e}_i \quad \cdots \quad  \mathbf{e}_n  \\
\begin{matrix} \mathbf{f}_1 \\ \mathbf{f}_2 \\ \vdots \\ \mathbf{f}_m \end{matrix} & \left[
\begin{matrix}
 t_{1,1}     & \cdots & t_{1,i}   & \cdots & t_{1,n} \\
 t_{2,1}     & \cdots & t_{2,i}   & \cdots & t_{2,n} \\
 \vdots      & \cdots & \vdots    & \ddots & \vdots  \\
 t_{m,1}     & \cdots & t_{m,i}   & \cdots & t_{m,n} \\
\end{matrix}
\right]   
\end{align*}
通过矩阵中的$m\times n$个元素，我们就记录了一个直映射$T$.

注意：这个矩阵自身并不能代表一个直映射，它仅仅记录了映射$T$在给定的基底$\mathcal{B}_{\mathbb{V}}$和$\mathcal{B}_{\mathbb{W}}$时的效果。
只有把它和$\mathcal{B}_{\mathbb{V}}$和$\mathcal{B}_{\mathbb{W}}$一起才能完全确定一个直映射$T$。
我们把这个矩阵称为$T$在$\mathcal{B}_{\mathbb{V}}$和$\mathcal{B}_{\mathbb{W}}$上的矩阵，记作$[T]_{\mathcal{B}_{\mathbb{V}}}^{\mathcal{B}_{\mathbb{W}}}$或
$\mathcal{M}_{\mathcal{B}_{\mathbb{V}}}^{\mathcal{B}_{\mathbb{W}}}(T).$

映射$T$作用于向量$\mathbf{u}$，结果为：
\begin{align*}
    T\mathbf{u} &= u_1T\mathbf{e}_1 + u_2T\mathbf{e}_2 + \cdots + u_nT\mathbf{e}_n  \\
                &= u_1\left(\sum_{k=1}^m t_{k,1}\mathbf{f}_k\right) + u_2\left(\sum_{k=1}^m t_{k,2}\mathbf{f}_k\right) + \cdots + u_n \left(\sum_{k=1}^m t_{k,n}\mathbf{f}_k\right)  \\
                &= \left(\sum_{i=1}^n u_i t_{1,i}\right) \mathbf{f}_1 + \left(\sum_{i=1}^n u_i t_{2,i}\right) \mathbf{f}_2 + \cdots + \left(\sum_{i=1}^n u_i t_{m,i}\right) \mathbf{f}_m.  \\
                &= \sum_{j=1}^m \sum_{i=1}^n u_i t_{j,i} \mathbf{f}_j 
\end{align*} 

如何通过矩阵的方式来表示以上的运算呢？我们可以使用$\mathbf{u}$对应的列向量和$T$对应的矩阵$[T]_{\mathcal{B}_{\mathbb{V}}}^{\mathcal{B}_{\mathbb{W}}}$。

首先从矩阵$[T]_{\mathcal{B}_{\mathbb{V}}}^{\mathcal{B}_{\mathbb{W}}}$中取出上起第一横行，
将横行左起第一个元素$t_{1,1}$和列向量上起第一个元素$u_1$相乘，
将第二个元素$t_{1,2}$和第二个元素$u_2$相乘，依此类推，然后将这些乘积加起来，
就得到了$\sum_{i=1}^n u_i t_{1,i}$，也就是$T\mathbf{u}$在$\mathcal{B}_{\mathbb{W}}$分解中$\mathbf{f}_1$的系数。

\begin{align*}
& \begin{bmatrix}u_1\\u_2\\ \vdots \\u_n \end{bmatrix} & \\
   [t_{1,1}, t_{1,2} , \cdots , t_{1,n}] & \qquad = \sum_{i=1}^n u_i t_{1,i} 
\end{align*}

类似地，我们可以对矩阵的第二行、第三行等等进行同样操作。最后我们得到：
\begin{align*}
&  \begin{bmatrix}u_1\\u_2\\ \vdots \\u_n \end{bmatrix} &  \\
    \begin{bmatrix}
        t_{1,1}    & t_{1,2}   & \cdots & t_{1,n} \\
        t_{2,1}    & t_{2,2}   & \cdots & t_{2,n} \\
        \vdots     & \vdots    & \ddots & \vdots  \\
        t_{m,1}    & t_{m,2}   & \cdots & t_{m,n} \\
       \end{bmatrix} \,& \quad\quad = \begin{bmatrix}\sum_{i=1}^n u_i t_{1,i}\\\sum_{i=1}^n u_i t_{2,i}\\ \vdots \\\sum_{i=1}^n u_i t_{m,i} \end{bmatrix} 
\end{align*}
这就是矩阵的乘法运算。我们也可以将其简单记为$[T]_{\mathcal{B}_{\mathbb{V}}}^{\mathcal{B}_{\mathbb{W}}}[\mathbf{u}]_{\mathcal{B}_{\mathbb{V}}}$。

给定出发空间和到达空间的基底$\mathcal{B}_{\mathbb{V}}$和$\mathcal{B}_{\mathbb{W}}$，确定了$T$在这些基底上对应的矩阵，
我们就可以给出$\mathbb{V}$中任何向量经过$T$映射的值。

举例来说，要描述$\mathbf{e}_1$经过映射得到的向量，我们可以用矩阵计算。$\mathbf{e}_1$在$\mathcal{B}_{\mathbb{V}}$上的坐标为
$(1, 0, \cdots, 0)$：

\begin{align*}
    & \begin{bmatrix}1\\ 0 \\ \vdots \\0 \end{bmatrix} &  \\
        \begin{bmatrix}
            t_{1,1}    & t_{1,2}   & \cdots & t_{1,n} \\
            t_{2,1}    & t_{2,2}   & \cdots & t_{2,n} \\
            \vdots     & \vdots    & \ddots & \vdots  \\
            t_{m,1}    & t_{m,2}   & \cdots & t_{m,n} \\
           \end{bmatrix} \,& \quad\quad = \begin{bmatrix} t_{1,1}\\ t_{2,1}\\ \vdots \\ t_{m,1} \end{bmatrix} 
\end{align*}
得到的坐标$(t_{1,1},  t_{2,1},  \cdots , t_{m,1})$表示向量$ t_{1,1}\mathbf{f}_1 + t_{2,1}\mathbf{f}_2 + \cdots + t_{m,1}\mathbf{f}_m$。
这正是$ T\mathbf{e}_1$的$\mathcal{B}_{\mathbb{W}}$分解。
也就是说，矩阵$[T]_{\mathcal{B}_{\mathbb{V}}}^{\mathcal{B}_{\mathbb{W}}}$的第一列系数，
表示$\mathbf{e}_1$经$T$映射的结果$T(\mathbf{e}_1)$在$\mathcal{B}_{\mathbb{W}}$上的坐标。

一般来说，矩阵的第$j$列对应着$T(\mathbf{e}_j)$这个向量在$\mathcal{B}_{\mathbb{W}}$上的坐标。
这就是我们前面定义直映射矩阵时要竖着排列$T(\mathbf{e}_1)$坐标的原因。
这样设计的排布方式和矩阵乘法正好相容，使得我们能够用矩阵乘法表达直映射。

对$m \times n$矩阵$T$来说，我们可以把它和一个$\mathbb{F}^n$到$\mathbb{F}^m$的直映射对应起来。
考虑$\mathbb{F}^n$和$\mathbb{F}^m$的标准基$\mathcal{C}_{n}$和$\mathcal{C}_{m}$，
则它们和矩阵$T$一同可以确定一个直映射。我们把这个映射称为$T$（默认）对应的直映射，也记为$T$。
矩阵$T$的第$i$列可以看作$\mathbb{F}^m$里的向量，我们将它记为$\mathbf{t}_i^|$：
$$\mathbf{t}_i^| = (t_{1,i}, t_{2,i}, \cdots, t_{m,i})$$
$\mathbf{t}_i^|$可以看作$\mathcal{C}_{n}$的第$i$个基向量$\mathbf{c}_i$经过映射后，
在$\mathcal{C}_{m}$上的坐标形式。

同样地，我们还可以把矩阵每行的元素看作$\mathbb{F}^n$里的向量：
$$ \mathbf{t}_i^- = (t_{i,1}, t_{i,2}, \cdots, t_{i,n}).$$
\begin{df}\label{df:4_1_0}
    矩阵的\textbf{列空间}是它的各个纵列对应的$\mathbb{F}^m$里的向量$\mathbf{t}_1^|, \mathbf{t}_2^|, \cdots , \mathbf{t}_n^|$生成的空间：$\langle \mathbf{t}_1^|, \mathbf{t}_2^|, \cdots , \mathbf{t}_n^| \rangle$. 
    矩阵的\textbf{行空间}是它的各个横行对应的$\mathbb{F}^n$里的向量$\mathbf{t}_1^-, \mathbf{t}_2^-, \cdots , \mathbf{t}_m^-$生成的空间：$\langle \mathbf{t}_1^-, \mathbf{t}_2^-, \cdots , \mathbf{t}_n^- \rangle$. 
\end{df}
$\mathbf{t}_i^|$就是$T\mathbf{c}_i$的坐标形式，所以矩阵的列空间就是矩阵对应的直映射的像空间。
矩阵的行空间则比较难理解，后面我们会证明，它可以看作该直映射的对偶映射的像空间。

我们来看一个具体的例子：

考虑二维平面$\mathbb{R}^2$到自身的一个映射$H$：
$$ \forall (x, y) \in \mathbb{R}^2, \,\,\, H(x,y) = (2x, y).$$
这个映射的效果是把点的横坐标放大到2倍，而纵坐标不变。
可以验证这是一个直映射。它的矩阵是什么呢？

我们首先考虑在哪个基上构建矩阵。
为了简单起见，我们使用标准基$\left((1,0),(0,1)\right)$作为出发空间（也是到达空间）的基。
根据$H$的定义，我们需要找到如下的矩阵$H$：
$$
H = \begin{bmatrix}
    h_{1,1} & h_{1,2}  \\
    h_{2,1} & h_{2,2}  \\
  \end{bmatrix}
$$
使得
\begin{align*}
&  \,\begin{bmatrix}x\\y\end{bmatrix}  \\
    \begin{bmatrix}
         h_{1,1} & h_{1,2}  \\
         h_{2,1} & h_{2,2}  \\
       \end{bmatrix} \,&\qquad = \begin{bmatrix}2x\\y \end{bmatrix} 
\end{align*}
令$y = 0$，我们得到：$h_{1,1}x = 2x$和$h_{1,1}x = 0$.

所以$h_{1,1} = 2, \,\,\, h_{2,1} = 0.$

类似地，令$x = 0$，可以推出：$h_{1,2} = 0, \,\,\, h_{2,2} = 1.$

也就是说：
$$
H = \begin{bmatrix}
    2 & 0  \\
    0 & 1  \\
  \end{bmatrix}.
$$

一些基本的直映射的矩阵可以很容易推出。

比如，零映射的矩阵就是所有元素都是$0$的矩阵，也叫\textbf{零矩阵}。
可以证明，它的样子和基底的选择没有关系。我们将其记为$0$.

$n$维空间到自身的恒等映射的矩阵是：
$$
\begin{bmatrix}
 1     & 0  & \,\cdots \, & 0 \\
 0     & 1  & \,\cdots \,& 0 \\
 \vdots      & \,\vdots  \,  & \ddots & \vdots  \\
 0     & 0   & \,\cdots \,& 1 \\
\end{bmatrix}
$$
这里我们设定出发空间和到达空间的基底是同一个基底，否则矩阵会变成别的样子。
除此之外，它的样子和基底的选择没有关系。我们称其为\textbf{单位矩阵}，记为$\mathbb{I}$。

上一章我们提到过$\mathcal{L}(\mathbb{V}, \mathbb{W})$的基底。基底中每个元素对应着一个矩阵，称为\textbf{基矩阵}或\textbf{基本矩阵}。
给定$1\leq i \leq n$和$1\leq j \leq m$，基底中的元素$E_{i,j}$将$\mathbb{V}$的基底中第$i$个元素$\mathbf{v}_i$
映射到$\mathbb{W}$基底的第$j$个元素：$\mathbf{w}_j$，并将其他元素都映射为零向量。
所以，$E_{i,j}$对应的矩阵是第$j$行第$i$列为1，其余元素为0的矩阵。

\section{矩阵的性质}

这一节我们介绍矩阵的一些基本性质。

\subsection{矩阵的运算}
$m\times n$的矩阵本质上是$m\times n$个标量按顺序排列。
矩阵的代数性质，关键在于为其定义的运算。

我们很自然地会想到为矩阵配备两个基本的运算：矩阵之间的加法和矩阵的数乘，分别对应直映射的加法和数乘运算。这两种运算分别如下：

\begin{align*}
&\begin{bmatrix}
 a_{1,1}     & a_{1,2}   & \cdots & a_{1,n} \\
 a_{2,1}     & a_{2,2}   & \cdots & a_{2,n} \\
 \vdots      & \vdots    & \ddots & \vdots  \\
 a_{m,1}     & a_{m,2}   & \cdots & a_{m,n} \\
\end{bmatrix}
+
\begin{bmatrix}
    b_{1,1}     & b_{1,2}   & \cdots & b_{1,n} \\
    b_{2,1}     & b_{2,2}   & \cdots & b_{2,n} \\
    \vdots      & \vdots    & \ddots & \vdots  \\
    b_{m,1}     & b_{m,2}   & \cdots & b_{m,n} \\
\end{bmatrix}  \\
=&
\begin{bmatrix}
    a_{1,1} + b_{1,1}     & a_{1,1} + b_{1,2}   & \cdots & a_{1,n} + b_{1,n} \\
    a_{2,1} + b_{2,1}     & a_{2,2} + b_{2,2}   & \cdots & a_{2,n} + b_{2,n} \\
    \vdots                & \vdots              & \ddots &           \vdots  \\
    a_{m,1} + b_{m,1}     & a_{m,2} + b_{m,2}   & \cdots & a_{m,n} + b_{m,n} \\
\end{bmatrix} ,
\end{align*}

$$
    k\cdot\begin{bmatrix}
     a_{1,1}     & a_{1,2}   & \cdots & a_{1,n} \\
     a_{2,1}     & a_{2,2}   & \cdots & a_{2,n} \\
     \vdots      & \vdots    & \ddots & \vdots  \\
     a_{m,1}     & a_{m,2}   & \cdots & a_{m,n} \\
    \end{bmatrix}
    =
    \begin{bmatrix}
        k\cdot a_{1,1}     & k\cdot a_{1,2}   & \cdots & k\cdot a_{1,n} \\
        k\cdot a_{2,1}     & k\cdot a_{2,2}   & \cdots & k\cdot a_{2,n} \\
        \vdots      & \vdots    & \ddots & \vdots  \\
        k\cdot a_{m,1}     & k\cdot a_{m,2}   & \cdots & k\cdot a_{m,n} \\
    \end{bmatrix} 
$$

可以验证，这两个运算对应了直映射的加法和数乘运算。

给定$\mathcal{L}(\mathbb{V}, \mathbb{W})$中两个映射$A$和$B$，
设$\mathcal{B}_\mathbb{V} = (\mathbf{e}_1, \cdots, \mathbf{e}_n)$和$\mathcal{B}_\mathbb{W} = (\mathbf{f}_1, \cdots \mathbf{f}_m)$分别是$\mathbb{V}$和$\mathbb{W}$的基，
设$A$和$B$在基底$\mathcal{B}_\mathbb{V}$和$\mathcal{B}_\mathbb{W}$下的矩阵是
$\mathcal{M}_{\mathcal{B}_\mathbb{V}}^{\mathcal{B}_\mathbb{W}}(A) = [a_{i,j}]_{1\leqslant i\leqslant m, 1\leqslant i\leqslant n}, \mathcal{M}_{\mathcal{B}_\mathbb{V}}^{\mathcal{B}_\mathbb{W}}(B) = [b_{i,j}]_{1\leqslant i\leqslant m, 1\leqslant i\leqslant n}$。
我们来计算映射$A+B$的矩阵。

对$\mathbb{V}$中任意元素$\mathbf{u}$，设它的$\mathcal{B}_\mathbb{V}$分解为$\mathbf{u} = u_1 \mathbf{e}_1 + u_2 \mathbf{e}_2 + \cdots + u_n\mathbf{e}_n$，
则：
\begin{align*}
    A(\mathbf{u}) &= \left(\sum_{i=1}^n u_i a_{1,i}\right) \mathbf{f}_1 + \left(\sum_{i=1}^n u_i a_{2,i}\right) \mathbf{f}_2 + \cdots + \left(\sum_{i=1}^n u_i a_{m,i}\right) \mathbf{f}_m  \\
    B(\mathbf{u}) &= \left(\sum_{i=1}^n u_i b_{1,i}\right) \mathbf{f}_1 + \left(\sum_{i=1}^n u_i b_{2,i}\right) \mathbf{f}_2 + \cdots + \left(\sum_{i=1}^n u_i b_{m,i}\right) \mathbf{f}_m 
\end{align*}
所以
\begin{align*}
    (A+B)(\mathbf{u}) &= A(\mathbf{u}) + B(\mathbf{u})  \\
    &= \left(\sum_{i=1}^n u_i a_{1,i}\right) \mathbf{f}_1 + \left(\sum_{i=1}^n u_i a_{2,i}\right) \mathbf{f}_2 + \cdots + \left(\sum_{i=1}^n u_i a_{m,i}\right) \mathbf{f}_m  \\
    & \,\,\,+ \left(\sum_{i=1}^n u_i b_{1,i}\right) \mathbf{f}_1 + \left(\sum_{i=1}^n u_i b_{2,i}\right) \mathbf{f}_2 + \cdots + \left(\sum_{i=1}^n u_i b_{m,i}\right) \mathbf{f}_m  \\
    &= \left(\sum_{i=1}^n u_i (a_{1,i} + b_{1,i}) \right) \mathbf{f}_1 + \left(\sum_{i=1}^n u_i (a_{2,i} + b_{2,i})\right) \mathbf{f}_2 + \cdots  \\
    & \,\,\,+ \left(\sum_{i=1}^n u_i (a_{m,i} + b_{m,i})\right) \mathbf{f}_m  
\end{align*}
这说明映射$A+B$在对应基底上的矩阵就是$A$的矩阵与$B$的矩阵相加的结果。

类似地，可以验证映射$kA$对应的矩阵就是标量$k$乘以$A$的矩阵的结果。

所有$m\times n$矩阵的集合，在配备了这两种运算后，成为一个平直空间，一般记作$\mathfrak{M}_{m,n}(\mathbb{F})$。

直映射的复合，则对应着矩阵的乘法：给定两个矩阵$A$和$B$，其乘积定义为：
\begin{align*}
& \begin{bmatrix}
    a_{1,1}    & a_{1,2}   & \cdots & a_{1,n} \\
    a_{2,1}    & a_{2,2}   & \cdots & a_{2,n} \\
    \vdots     & \vdots    & \ddots & \vdots  \\
    a_{m,1}    & a_{m,2}   & \cdots & a_{m,n} \\
   \end{bmatrix} \cdot 
   \begin{bmatrix}
    b_{1,1}     & b_{1,2}   & \cdots & b_{1,p} \\
    b_{2,1}     & b_{2,2}   & \cdots & b_{2,p} \\
    \vdots      & \vdots    & \ddots & \vdots  \\
    b_{n,1}     & b_{n,2}   & \cdots & b_{n,p} \\
\end{bmatrix}  \\
    =&\begin{bmatrix}
        \sum_{i=1}^n a_{1,i} b_{i,1}     & \sum_{i=1}^n a_{1,i} b_{i,2}   & \cdots & \sum_{i=1}^n a_{1,i} b_{i,p} \\
        \sum_{i=1}^n a_{2,i} b_{i,1}     & \sum_{i=1}^n a_{2,i} b_{i,2}   & \cdots & \sum_{i=1}^n a_{2,i} b_{i,p} \\
        \vdots      & \vdots    & \ddots & \vdots  \\
        \sum_{i=1}^n a_{m,i} b_{i,1}     & \sum_{i=1}^n a_{m,i} b_{i,2}   & \cdots & \sum_{i=1}^n a_{m,i} b_{i,p} \\
    \end{bmatrix}. 
\end{align*}

读者可以自行验证，复合映射$A\circ B$作用在任意向量$\mathbf{u}$上的结果，
就是矩阵$A$与$B$的乘积对应的映射，作用在$\mathbf{u}$上的结果。

注意每个矩阵的下标。$A$是一个$m\times n$的矩阵，$B$是一个$n\times p$的矩阵，
而它们的乘积是一个$m\times p$的矩阵。其中的$m,n,p$是正整数。

我们在前面的例子里就可以看出这一点。
当我们查看一个$n$维空间中的向量经过直映射变换的矩阵表示时，
我们计算了表示直映射的$m\times n$的矩阵和表示向量的$n\times 1$矩阵的乘积，结果是一个$m\times 1$矩阵。

所以，从我们上面定义的矩阵平直空间的角度来看，
两个矩阵相乘，得到第三个矩阵，这三个矩阵很可能不在同一个空间里。

显然，只有当我们讨论的对象在同一个空间里，讨论才能继续下去。
要使得它们在同一个空间里，我们就要求$m=n$且$n=p$，即这三个数相等。
这时，我们讨论的空间就变成了$n\times n$矩阵的空间\footnote{一般称$n\times n$矩阵为\textbf{方阵}。}。
我们把这个空间记为$\mathfrak{M}_n(\mathbb{F})$，表示系数为$\mathbb{F}$中元素的所有$n\times n$矩阵的空间。

在$\mathfrak{M}_n(\mathbb{F})$中，可以为矩阵之间配备乘法。这时两个$\mathfrak{M}_n(\mathbb{F})$中的矩阵可以相乘，乘积也在$\mathfrak{M}_n(\mathbb{F})$中。
我们还可以验证，矩阵乘法和矩阵空间的矩阵加法和数乘满足如下的关系：
\begin{enumerate}
    \item 兼容数乘：$\forall A,B \,\in \mathfrak{M}_n(\mathbb{F}), \,\, \forall \lambda\in\mathbb{F}, \,\,\, (\lambda\cdot A)\cdot B =  A \cdot (\lambda\cdot B) = \lambda\cdot (A\cdot B). $
    \item 分配律：$\forall A,B,C \,\in \mathfrak{M}_n(\mathbb{F}), \,\, (A + B)\cdot C = A\cdot C + B\cdot C,\,\,\, C\cdot (A + B) = C\cdot A + C\cdot B.$
    \item 结合律：$\forall A,B,C \,\in \mathfrak{M}_n(\mathbb{F}), \,\, (A + B)\cdot C = A\cdot C + B\cdot C,\,\,\, A\cdot(B\cdot C) = (A\cdot B)\cdot C.$
\end{enumerate}
我们说$\mathfrak{M}_n(\mathbb{F})$关于矩阵加法和乘法构成一个\textbf{环}。
$\mathfrak{M}_n(\mathbb{F})$中的矩阵可以相加、相乘，并乘上任意系数。
和前面一样，我们把这种结构称为域$\mathbb{F}$上的\textbf{域上场}，或域$\mathbb{F}$上的\textbf{结合代数}。

注意到这个结构和直映射构成的空间是一样的，两者同构。

矩阵的乘法满足结合律，它是否也满足交换律呢？

用以下的例子试一试：
\begin{align*}
    A = \begin{bmatrix}
        1     & 1  \\
        0     & 1 \\
    \end{bmatrix}, 
    B = \begin{bmatrix}
        0     & 1  \\
        1     & 0  \\
    \end{bmatrix}, 
\end{align*}
则
\begin{align*}
    A \cdot B = \begin{bmatrix}
        1     & 1  \\
        1     & 0 \\
    \end{bmatrix}, 
    B \cdot A = \begin{bmatrix}
        0     & 1  \\
        1     & 1  \\
    \end{bmatrix} 
\end{align*}
所以$A\cdot B \neq B \cdot A$。

可以验证，任意两个矩阵$A$和$B$，$A\cdot B$和$B \cdot A$不一定相等，
所以，矩阵乘法不满足交换律。这一点和直变换的复合也是一样的，我们用矩阵的语言验证了这一点。

% 最后，我们就矩阵作为直映射的表示形式再说两句。

% 前面的例子里，我们总把表示直映射的矩阵放在左边，
% 把表示向量的矩阵放在右边，让两者相乘，表达直映射作用在向量上。
% 能否把它们的位置调换过来呢？

% 如果要将它们的位置对调，而仍然得出有意义的结果，
% 我们会发现，要把直映射和向量的相应系数横着写，并且从上到下排起来。
% 这样排布的向量和矩阵，通过矩阵乘法，依然可以表达直映射作用在向量上的过程。公式为：
% \begin{align*}
%     & \begin{bmatrix}
%         t_{1,1}    & t_{2,1}   & \cdots & t_{m,1} \\
%         t_{1,2}    & t_{2,2}   & \cdots & t_{m,2} \\
%         \vdots     & \vdots    & \ddots & \vdots  \\
%         t_{1,n}    & t_{2,n}   & \cdots & t_{m,n} \\
%        \end{bmatrix}  &  \\
%     \begin{bmatrix}u_1 & u_2 & \cdots & u_n \end{bmatrix}    &  \\
%     &= \begin{bmatrix}\sum_{i=1}^n u_i t_{1,i} & \sum_{i=1}^n u_i t_{2,i} & \cdots  & \sum_{i=1}^n u_i t_{m,i} \end{bmatrix} 
% \end{align*}
% 这个公式和上面详尽写出的第一种记法和公式相比，恰好是“对称”的：
% 关于从左上到右下的$45^\circ$直线呈轴对称关系。

% 具体来说，列向量变成了行向量，$m \times n$的矩阵变成了$n \times m$的矩阵。
% 原来第$i$行第$j$列的元素$t_{i,j}$出现在新矩阵的第$j$行第$i$列。
% 而我们很确定这个等式是正确的，因为我们除了把原来的公式“翻转”以外，什么都没有做，里面的代数运算丝毫不受影响！

% 我们把这种表达方式称为矩阵右乘行向量，把之前的表达方式称为矩阵左乘列向量。
% 实际上，矩阵右乘行向量的表达方式并无问题，甚至在很多情况下更为自然。
% 本书后面仍将使用矩阵左乘列向量的方式，以便和传统的映射记法（$f(x)$，函数符号在左，自变量在右）兼容，便于读者统一理解直映射和矩阵的关系。
% 但我们也鼓励读者适当使用矩阵右乘行向量的方式理解某些概念。

\subsection{矩阵的转置}
把$m\times n$矩阵$A$从第一行第一列起，自左上到右下画一条线，这条线称为矩阵的\textbf{主对角线}。以主对角线为轴把矩阵“翻转”，得到一个新矩阵，称为$A$的\textbf{转置矩阵}，记为$A^{\tr}$或$A'$。具体来说：
$$ \forall 1\leqslant i \leqslant m, \; 1 \leqslant j \leqslant n, \;\; A^{\tr}_{i,j} = A_{j, i}.$$

矩阵的转置，是将矩阵元素沿着主对角线做轴对称变换。
另一种说法是：矩阵的转置，是将矩阵自左向右的纵列，变为自上而下的横行。

根据轴对称的性质，我们很容易就能确定：
\begin{enumerate}
    \item 矩阵转置的转置是它自身；
    \item 矩阵主对角线上的元素，不受矩阵转置的影响；
\end{enumerate}

从上面“翻转”矩阵乘法公式的动作中，我们还能得到如下的关系式：
$$ (A \cdot B)^{\tr} = B^{\tr} \cdot A^{\tr}.$$

我们可以通过计算$(A\cdot B)^{\tr}$中任一元素来验证这个等式。设$A,B$分别是$m\times n$和$n\times p$的矩阵。记$C = A\cdot B$，则$C$是$m \times p$矩阵，$(A\cdot B)^{\tr} = C^{\tr}$是$p\times m$矩阵。根据矩阵乘法：
$$ \forall \; 1\leq i, \leq m,\; \; 1 \leqslant j \leqslant p,\;\, c_{i, j} = \sum_{k=1}^n a_{i,k} b_{k,j},$$
因此，$(A\cdot B)^{\tr}$的第$i$行第$j$列元素$x_{i,j}$就是$C$的第$j$行第$i$列元素$c_{j,i}$，它等于
\begin{align*}
    x_{i,j} &= c_{j,i} = \sum_{k=1}^n a_{j,k} b_{k,i}  \\
    &= \sum_{k=1}^n b^{\tr}_{i, k} a^{\tr}_{k, j}  
\end{align*}
即$B^{\tr}$第$i$行乘$A^{\tr}$第$j$列的结果。其中$a^{\tr}_{k, j} , b^{\tr}_{i, k}$表示$A^{\tr}$和$B^{\tr}$中相应的元素。

这个关系式和交换律有点像，但并非交换律，只是同一个矩阵计算的不同写法。这一点需要读者注意。

最后还可以定义$\mathfrak{M}_{m,n}(\mathbb{F})$上的转置映射$A \mapsto A^{\tr}$。不难验证它是一个直映射。

\subsection{矩阵的迹}
最后介绍方形矩阵的\textbf{迹}。给定$n\times n$方阵$A$，定义$A$的迹是$A$的对角线元素的和，记为$\ji(A)$：
\begin{align*}
    \ji:\;\; \mathfrak{M}_n(\mathbb{F})\; &\rightarrow \; \mathbb{F} \\
    A \quad &\mapsto \; \ji(A) = \sum_{i=1}^n a_{i,i}
\end{align*}
按定义来算一些简单矩阵的迹：零矩阵的迹是$0$。$n$阶单位矩阵的迹是$n$。下面给出一些具体矩阵的迹：
\begin{align*}
    \ji\left(\begin{bmatrix}
        1 & -2 \\ 2 & 1
    \end{bmatrix}\right)
    &= \phantom{-}1 + 1 = 2, \quad 
    \ji\left(\begin{bmatrix}
        1 & 2 \\ -3 & 4
    \end{bmatrix}\right)
    = 1 + 4 = 5, \\
    \ji\left(\begin{bmatrix}
        -2 & 1 \\ 0 & 5
    \end{bmatrix}\right)
    &= -2 + 5 = 3,\quad
    \ji\left(\begin{bmatrix}
        4 & -3 \\ 1 & -4
    \end{bmatrix}\right)
    = 4 - 4 = 0.
\end{align*}

迹有什么基本性质呢？

显然，矩阵的转置不改变迹，因为对角线上的元素不变。
$$ \ji(A^{\tr}) = \ji(A).$$
另外，$\ji$自身是直映射：
\begin{align*}
    \forall \; A, B \; \in \, \mathfrak{M}_n(\mathbb{F}), \;\; \forall s, t \; \in \, \mathbb{F}, \quad \ji(sA + tB) = s\ji(A) + t\ji(B).
\end{align*}
而且满足“乘法交换律”：
$$ \forall \; A, B \; \in \, \mathfrak{M}_n(\mathbb{F}), \quad \ji(AB) = \ji(BA). $$
证明：$AB$的对角线上第$i$个元素是$\sum_{k=1}^n a_{ik} b_{ki}$，所以其对角线元素之和是
$$ \sum_{i=1}^n \sum_{k=1}^n a_{ik} b_{ki}. $$
注意到这个式子关于下标对称。把两个求和的顺序调转，就得到
\begin{align*}
    \sum_{i=1}^n \sum_{k=1}^n a_{ik} b_{ki} = \sum_{k=1}^n \sum_{i=1}^n a_{ik} b_{ki} = \sum_{k=1}^n \sum_{i=1}^n b_{ki} a_{ik}.
\end{align*}
后者就是$BA$对角线元素之和。

\section{一次方程组}
当我们把向量经过直映射得到另一个向量的过程用矩阵表示时，
会发现其表达式的形式和多元一次方程组相同。
比如，在$\mathbb{R}^3$中，
向量$\mathbf{x} = (x_1, x_2, x_3)$经过直映射$A = [a_{i,j}]_{1\leqslant i \leqslant 3}$，
得到另一个向量$\mathbf{b} = (b_1, b_2, b_3)$，
用矩阵乘法的形式，$A\mathbf{x} = \mathbf{b}$可以表示为：
$$
    \begin{bmatrix}
        a_{1,1}     & a_{1,2}   & a_{1,3} \\
        a_{2,1}     & a_{2,2}   & a_{2,3} \\
        a_{3,1}     & a_{3,2}   & a_{3,3} \\
    \end{bmatrix} \cdot 
    \begin{bmatrix}
         x_{1} \\ x_{2} \\ x_{n} 
    \end{bmatrix} = 
    \begin{bmatrix}
        b_{1} \\ b_{2} \\ b_{n} 
    \end{bmatrix}
$$
计算左侧矩阵乘法，得到
$$
    \begin{bmatrix}
        a_{1,1} x_1 + a_{1,2} x_2 +  a_{1,3} x_3 \\
        a_{2,1} x_1 + a_{2,2} x_2 +  a_{2,3} x_3 \\
        a_{3,1} x_1 + a_{3,2} x_2 +  a_{3,3} x_3 \\
    \end{bmatrix} = 
    \begin{bmatrix}
        b_{1} \\ b_{2} \\ b_{n} 
    \end{bmatrix} 
$$
因此，对各个$b_i$来说，有等式：
$$
\left\{
    \begin{aligned}
        a_{1,1} x_1 + a_{1,2} x_2 +  a_{1,3} x_3 = b_1 \\
        a_{2,1} x_1 + a_{2,2} x_2 +  a_{2,3} x_3 = b_2 \\
        a_{3,1} x_1 + a_{3,2} x_2 +  a_{3,3} x_3 = b_3 \\
    \end{aligned}
\right.
$$
如果把$x_1, x_2, x_3$看作未知数，把$a_{i,j}$和$b_i$看作已知数，这就是一个三元一次方程组。

一般来说，给定基底，对应$m\times n$矩阵$A$的直映射$A \in \mathcal{L}(\mathbb{V}, \mathbb{W})$把向量$\mathbf{x}$映射到向量$\mathbf{b}$，
就对应一个含有$m$个方程的$n$元一次方程组。
$$
\left\{
    \begin{aligned}
        a_{1,1} x_1 + a_{1,2} x_2 + \cdots + a_{1,n} x_n &= b_1 \\
        a_{2,1} x_1 + a_{2,2} x_2 + \cdots + a_{2,n} x_n &= b_2 \\
        &\vdots \\
        a_{m,1} x_1 + a_{m,2} x_2 + \cdots + a_{m,n} x_n &= b_m \\
    \end{aligned}
\right.\qquad (\mathcal{S})
$$
求解一个这样的方程组，等价于找出一个经过直映射$A$能得到向量$\mathbf{b}$的向量$\mathbf{x}$。
我们把两者统称为解\textbf{一次方程组}，把满足一次方程组的未知数向量$\mathbf{x}$称为它的解，
把所有解的集合称为一次方程组的\textbf{解集}。

什么情况下，我们能找到一个这样的向量$\mathbf{x}$呢？我们从一个特殊情况：$\mathbf{b} = \mathbf{0}$出发。
这时，一次方程组变为$A\mathbf{x} = \mathbf{0}$，我们将这个特殊情况称为“解\textbf{齐次一次方程组}”。
什么样的向量$\mathbf{x}$使得$A\mathbf{x} = \mathbf{0}$呢？显然，当且仅当它是直映射$A$零空间的元素。
也就是说，$A\mathbf{x} = \mathbf{0}$解的集合是$A$的零空间：$\Ker A$。
显然，$\mathbf{x} = \mathbf{0}$是任何映射零空间的元素，所以必然是解。
我们把它称为一次方程组的\textbf{平凡解}。如果$A$的零空间维数大于0，那么除了$\mathbf{0}$之外，方程还有别的解。

如果$\mathbf{b} \neq \mathbf{0}$，我们希望找到$A\mathbf{x} = \mathbf{b}$的所有解。
假设我们已经知道某个向量$\mathbf{x_0}$是一次方程组的解：$A\mathbf{x_0} = \mathbf{b}$，
那么，其它的解与$\mathbf{x_0}$是什么关系呢？

假设$\mathbf{x_1}$是一次方程组的另一个解：$A\mathbf{x_1} = \mathbf{b}$，那么，两式相减，得到：
$$A(\mathbf{x_1} - \mathbf{x_0}) = \mathbf{0}$$
我们回到了前面的特殊情况。可以推出：
$ \mathbf{x_1} - \mathbf{x_0} \in \Ker A.$
反过来，如果$\exists \mathbf{v} \in \Ker A$，使得$\mathbf{x_1} = \mathbf{x_0} + \mathbf{v}$，
那么，
$$A\mathbf{x_1} = A\mathbf{x_0} + A\mathbf{v} = \mathbf{b} + \mathbf{0} = \mathbf{b}.$$
因此$\mathbf{x_1}$是一次方程组$A\mathbf{x} = \mathbf{b}$的解。

以上论证说明，一般情况下，一次方程组$A\mathbf{x} = \mathbf{b}$的解集，如果不是空集，
总能写成$\mathbf{x_0} + \Ker A$的形式。其中$\mathbf{x_0}$称为一次方程组的一个\textbf{特解}。

什么时候，特解存在呢？按照直映射的定义，存在$A\mathbf{x_0} = \mathbf{b}$，
当且仅当$\mathbf{b} \in \Img A$。另一方面，我们可以把
$$
    \begin{bmatrix}
        a_{1,1} x_1 + a_{1,2} x_2 + \cdots + a_{1,n} x_n \\
        a_{2,1} x_1 + a_{2,2} x_2 + \cdots + a_{2,n} x_n \\
        \vdots \\
        a_{m,1} x_1 + a_{m,2} x_2 + \cdots + a_{m,n} x_n \\
    \end{bmatrix}
$$
写为：
$$
    x_1 \begin{bmatrix}
        a_{1,1} \\
        a_{2,1} \\
        \vdots \\
        a_{m,1} \\
    \end{bmatrix} + x_2 \begin{bmatrix}
        a_{1,2} \\
        a_{2,2} \\
        \vdots \\
        a_{m,2}  \\
    \end{bmatrix} + \cdots + x_n \begin{bmatrix}
        a_{1,n} \\
        a_{2,n} \\
        \vdots \\
        a_{m,n}  \\
    \end{bmatrix}
$$
其中每个向量由矩阵$A$一列的元素组成。它们就是列向量$A_1^|, A_2^|, \cdots A_n^|$。
而上式就是列向量的直组合。所以，解一次方程组$A\mathbf{x} = \mathbf{b}$又可以看作寻找适当的系数
$x_1, x_2, \cdots , x_n$，使得直组合$x_1 A_1^| + x_2 A_2^| + \cdots + x_n A_n^| = \mathbf{b}$的过程。
所以，一次方程组有解，等于说$\mathbf{b}$在列向量生成的空间$\langle A_1^|, A_2^|, \cdots, A_n^| \rangle$里。

由容佚定理，$\nul A + \rang A = \dim \mathbb{V} = n$，所以，
\begin{enumerate}
    \item 如果$\rang A < n$，那么$\nul A = n - \rang A > 0$，这时一次方程组$A\mathbf{x} = \mathbf{0}$有无穷多组非平凡解。一次方程组$A\mathbf{x} = \mathbf{b}$要么无解，要么有无穷多组解。
    \item 如果$\rang A = n$，那么$\nul A = 0$，即$A$为单射。一次方程组$A\mathbf{x} = \mathbf{0}$没有非平凡解，$A\mathbf{x} = \mathbf{b}$要么无解，要么有唯一解。如果$m=n$，那么$A$也是满射，这时$A\mathbf{x} = \mathbf{b}$总有唯一解。
\end{enumerate}

一次方程组和向量的直相关有密切关系。一组向量$\mathbf{a}_1, \cdots , \mathbf{a}_n \in \mathbb{F}^m$直相关，
说明存在不全为$0$的系数$x_1, \cdots , x_n \in \mathbb{F}$，使得$x_1\mathbf{a}_1 + \cdots + x_n\mathbf{a}_n = \mathbf{0}.$
而这个等式等价于一次方程组：
$$ A\mathbf{x} = \mathbf{0}.$$
其中矩阵$A$的第$i$列是向量$\mathbf{a}_i$（在标准基上的）坐标，而向量$\mathbf{x}$是系数$x_1, \cdots , x_n$构成的列向量。
$\mathbf{a}_1, \cdots , \mathbf{a}_n$直相关，等价于说一次方程组有非零解。反之，它们直无关，则等价于一次方程只有平凡解。

如何求一次方程组的解呢？我们从中学学过的消元法出发，给出一个基于归纳法的算法。

设要解的方程组$(\mathcal{S})$有$k$个方程。第$i$个方程可以写成$\sum_{j=1}^m a_{i,j} x_j = b_{i}$的形式。


\begin{pp}\label{pp:5_0_0}
通过有限次四则运算，可以得到多元一次方程组的一组解，或证明其无解。
\end{pp}
\begin{proof} % [{\bfseries 证明\ref{pp:5_0_0}}]
使用归纳法，证明命题：对于$m$元一次方程组，可以通过有限次四则运算得到一组解，或证明其无解。\\
$m=1$时，每个方程变为$a_{i,1} x_1 = b_{i}$。作分类讨论：\\
如果任一个方程中$a_{i,1} = 0$，$b_{i} \neq 0$，则方程无解。\\
如果所有方程都是$a_{i,1} = b_{i} = 0$的情况，即$0\cdot x_1 = 0$的形式，这时，$x_1$可以取任意值。
我们把这个情形称为“任意填充”。如果某个方程满足$a_{i,1} \neq 0$，则对应解得$x_1 = \frac{b_{i}}{a_{i,1}} $。
如果对所有$a_{i,1} \neq 0$的方程，$\frac{b_{i}}{a_{i,1}}$都等于同一个值$c$，则方程组有唯一解$x_1 = c $。
否则方程组无解。\\
综上，求解或判断无解用到至多$k$次四则运算。\\
假设$m < n$的时候命题成立。对于$m = n$的情形，不失一般性，我们从$x_1$的系数出发，作分类讨论。\\
如果$k$个方程中，$x_1$的系数都是$0$，那么方程组至多有$n-1$个未知数。由归纳假设可知命题成立。
这时，$x_1$的值可以任意填充。\\
如果至少有一个方程中$x_1$的系数不是$0$，不失一般性，假设$a_{1,1} \neq 0$。
我们尝试构造一个未知数个数不超过$n-1$的方程组。\\
如果方程组$(\mathcal{S})$只有一个方程：
$$ \sum_{j=1}^n a_{1,j} x_j = b_{1}. $$
则我们可以考虑方程组$(\mathcal{S}')$：
$$ \sum_{j=2}^n a_{1,j} x_j = 0. $$
其未知数个数不超过$n-1$。\\
如果方程组$(\mathcal{S})$不只有一个方程，假如还有其他方程中$x_1$的系数$a_{i,1} \neq 0$，执行操作：
把第$1$个方程乘以系数$-\frac{a_{i,1}}{a_{1,1}}$，
加到第$i$个方程上。这样，新得到的第$i$个方程中，$x_1$的系数就变成$0$了。经过至多$k-1$次操作，
我们可以确保方程组中恰有一个方程（第$1$个方程）中$x_1$系数不为$0$，其余的都为$0$。
此时，考虑第$2$至第$k$个方程构成的方程组$(\mathcal{S}')$，则其未知数个数不超过$n-1$。\\
由归纳假设可知，我们可以通过有限次四则运算得到$(\mathcal{S}')$的一组解或证明它无解。\\
如果$(\mathcal{S}')$无解，那么$(\mathcal{S})$也无解。如果$(\mathcal{S}')$有一组解：
$x_2 = c_2, \cdots , x_n = c_n$，那么将其代入第$1$个方程，就得到：
$$ a_{1,1} x_1 = b_{1} - \sum_{j=2}^n a_{1,j} c_j. $$
于是解得：
$$ \left\{
    \begin{array}{rl}
    x_1 &= \frac{b_{1} - \sum_{j=2}^n a_{1,j} c_j}{a_{1,1}} \\
    x_2 &= c_2  \\
     \quad & \,\,\,\vdots   \\
    x_n &= c_n 
    \end{array}
\right.$$
所以$(\mathcal{S})$的每一组解都对应$(\mathcal{S}')$的一组解，且可以通过至多$2n-1$次四则运算从$(\mathcal{S}')$得到。\\
综上所述，命题当$m = n$也成立。因而，根据归纳法，命题对任意正整数$m$成立。

\end{proof}

我们来看一个例子。求解以下方程组：
$$ (\mathcal{S}):\quad \left\{
    \begin{array}{rl}
x + 2y + z &= 3 \quad \quad (1)\\
2x - y + 2z &= 1  \quad \quad (2)
\end{array}
\right. $$
它对应的是由有理数系数的矩阵$A \in \mathfrak{M}_{2,3}(\mathbb{Q})$
（对应一个从$\mathbb{Q}^3$到$\mathbb{Q}^2$的直映射）和向量$\mathbf{b}\in \mathbb{Q}^2$构成的一次向量方程：
$$ A\mathbf{x} = \mathbf{b}, \quad A = \begin{bmatrix}1 & 2 & 1 \\ 2 & -1 & 2\end{bmatrix}, \,\,\, \mathbf{b} = \begin{bmatrix} 3 \\ 1 \end{bmatrix}, \,\,\, \mathbf{x} = \begin{bmatrix} x \\ y \\ z \end{bmatrix} \in \mathbb{Q}^3.$$
根据算法，消去方程$(2)$中的$x$项，$(\mathcal{S})$变为：
$$ \left\{
    \begin{array}{rl}
x + 2y + z &= 3 \\
 - 5y + 0z &= -5  
\end{array}
\right. $$
接着求解二元一次方程组：
$$ (\mathcal{S}'):\quad  - 5y + 0z = -5 $$
按照算法，再转为求解一元一次方程组：
$$ (\mathcal{S}''):\quad  0z = 0 $$
$z$可以任意填充。对每个解$z = z_0$，代入$(\mathcal{S}')$，得到解：
$$ \left\{
    \begin{array}{rl}
y &=\frac{-5 - 0\cdot z}{-5} = 1 \\
 z &= z_0 
\end{array}
\right. $$
再代入$(\mathcal{S})$得到：
$$ \left\{
    \begin{array}{rl}
x &= 1 - z_0 \\
y &= 1 \\
 z &= z_0
\end{array}
\right. $$
比如，取$z_0 = 0.4$，就得到一组解：
$$ \left\{
    \begin{array}{rl}
x &= 0.6  \\
y &= 1 \\
z &= 0.4  
\end{array}
\right. $$
我们注意到，$z_0$的值不一定要在系数域中。比如，我们取$z_0 = \sqrt{2}$或$z_0 = 1 - \imath $，都可以得到$(\mathcal{S})$的一组解。但我们不能说方程$A\mathbf{x} = \mathbf{b}$有解$(\imath , 1, 1 - \imath )$，因为向量$(\imath , 1, 1 - \imath )$不属于$\mathbb{Q}^3$。

如果我们把$A$看作$\mathfrak{M}_{2,3}(\mathbb{R})$或$\mathfrak{M}_{2,3}(\mathbb{C})$里的矩阵，
我们得到的解集会产生变化。从消元法的角度来看，变化的原因在“任意填充”这一步。如果我们规定任意填充的时候只能填有理数，
那么得到的就是有理数解。如果我们规定可以填充实数或复数，那么就可以得到实数或复数的解。如果求解的时候没有“任意填充”的步骤，
那么所有解都通过对系数的有限次四则运算得到。而域结构对四则运算是封闭的，所以得到的解必然也在系数域里面。

如何发现算法中是否出现了“任意填充”呢？出现“任意填充”，表示方程组有无穷组解。反之，如果方程组只有一组解或无解，
那么算法中就不可能出现“任意填充”。所以，比照前面我们对解集的讨论，我们可以总结出现“任意填充”的条件。

例如，如果映射为单射，那么其零空间为$\{0\}$，方程组有唯一解。这时，求解算法中不会出现“任意填充”，因此，
唯一解的值必然在系数域中。

从这里可以引出一个重要的结论：
\begin{pp}\label{pp:5_0}
    给定域$\mathbb{K}$和它的子域$\mathbb{F}$，
    如果$\mathbb{F}^n$中的一组向量$\mathbf{v}_1, \mathbf{v}_2, \cdots , \mathbf{v}_m$直无关，
    那么它们作为$\mathbb{K}^n$的向量仍然直无关。
\end{pp}
这个命题触及了直相关性的核心：如果我们无法在$\mathbb{F}$中找到一组系数，使对应向量的直组合等于零向量，
那么，我们把寻找的范围扩大，能否找到这样的系数呢？答案是否定的。本质的原因是：直性涉及的运算只有加法和数乘，
仅仅靠这两种运算，我们无法跳出域结构的“五指山”。

\begin{proof} % [{\bfseries 证明\ref{pp:5_0}}]
设有$x_1, \cdots, x_m$使得
$$ x_1\mathbf{v}_1 + x_2 \mathbf{v}_2 + \cdots + x_m \mathbf{v}_m = \mathbf{0}. $$
那么，$\mathbf{x} = (x_1, \cdots, x_m)$是一次方程组：
$$ (\mathcal{S}): \quad V\mathbf{x} = \mathbf{0} $$
的解。其中矩阵$V$的第$i$列是向量$\mathbf{v}_i$在标准基上的坐标。\\
当我们把矩阵$V$看作$\mathfrak{M}_{m,n}(\mathbb{F})$中元素时，
由于$\mathbf{v}_1, \mathbf{v}_2, \cdots , \mathbf{v}_m$作为$\mathbb{F}^n$的向量直无关，
所以$(\mathcal{S})$有唯一解：$\mathbf{x} = \mathbf{0}$。
这说明，我们用消元法求解时不会出现“任意填充”。\\
现在，我们把$V$作为$\mathfrak{M}_{m,n}(\mathbb{K})$中的矩阵看待，
重新用消元法求解$(\mathcal{S})$，由于过程中没有“任意填充”，所以得到的仍然是唯一解：
$\mathbf{x} = \mathbf{0}$。这说明$\mathbf{v}_1, \mathbf{v}_2, \cdots , \mathbf{v}_m$作为$\mathbb{K}^n$的向量，
仍然直无关。\\
\end{proof}


\section{基变换}
直映射在给定基底的时候，可以用矩阵来表示。
那么，不同的基底上的直映射的矩阵有什么联系呢？

首先来看向量在不同基底上的坐标有什么联系。
设$\mathcal{B}_1 = (\mathbf{e}_1, \mathbf{e}_2, \cdots, \mathbf{e}_n)$和$\mathcal{B}_2= (\mathbf{e}_1', \mathbf{e}_2', \cdots, \mathbf{e}_n')$是空间$\mathbb{V}$的两组基。
任意给定向量$\mathbf{u} \in \mathbb{V}$，设它的$\mathcal{B}_1$和$\mathcal{B}_2$分解分别为$\sum_{i=1}^n u_i\mathbf{e}_i $与$\sum_{i=1}^n u_i'\mathbf{e}_i' $。
如果可以将$\sum_{i=1}^n u_i'\mathbf{e}_i' $表达为$\mathbf{e}_1, \mathbf{e}_2, \cdots, \mathbf{e}_n$的直组合，由基底的直无关性，
这个直组合与$u_1, u_2, \cdots, u_n$吻合，于是我们就可以得出$u_i$和$u_i'$之间的联系。为此，我们需要知道每个$\mathbf{e}_i'$如何用$\mathbf{e}_1, \mathbf{e}_2, \cdots, \mathbf{e}_n$的直组合来表示。

$\mathbf{e}_i'$是$\mathbb{V}$中的向量，所以必然有唯一的方式写成$\mathbf{e}_1, \mathbf{e}_2, \cdots, \mathbf{e}_n$的直组合。换句话说，
$$ \forall \,\, 1 \leqslant i \leqslant n, \quad \exists g_{1, i}^{1\rightarrow 2}, g_{2, i}^{1\rightarrow 2} ,\cdots , g_{n, i}^{1\rightarrow 2} \in \mathbb{F}, \quad \mbox{使得} \,\,\, \mathbf{e}_i' = \sum_{k=1}^n g_{k,i}^{1\rightarrow 2} \mathbf{e}_k$$
其中的系数是由$\mathbf{e}_i'$自身的性质决定的。将这些系数按照上面给出的下标写成矩阵：
$$G^{1\rightarrow 2} = \left[g_{i,j}^{1\rightarrow 2}\right]_{1\leqslant i , j \leqslant n}$$
我们将这个矩阵称为从$\mathcal{B}_1$到$\mathcal{B}_2$的基变换矩阵。注意：矩阵里第$i$列是$\mathbf{e}_i'$的$\mathbf{e}_1, \mathbf{e}_2, \cdots, \mathbf{e}_n$分解，而不是第$i$行。

我们现在来计算$\sum_{i=1}^n u_i'\mathbf{e}_i' $对应的$\mathbf{e}_1, \mathbf{e}_2, \cdots, \mathbf{e}_n$的直组合：
$$\sum_{i=1}^n u_i'\mathbf{e}_i' = \sum_{i=1}^n \sum_{k=1}^n g_{k,i}^{1\rightarrow 2} u_i' \mathbf{e}_k = \sum_{k=1}^n \left(\sum_{i=1}^n g_{k,i}^{1\rightarrow 2} u_i'\right) \mathbf{e}_k$$
所以，
$$\forall \,\, 1 \leqslant k \leqslant n, \quad u_k = \sum_{i=1}^n g_{k,i}^{1\rightarrow 2} u_i'.$$
用$u_1', u_2', \cdots, u_n'$表示$u_1, u_2, \cdots, u_n$的方法，写成矩阵形式，就是：
$$\begin{bmatrix}g_{1,1}^{1\rightarrow 2} & g_{1,2}^{1\rightarrow 2} & \cdots & g_{1,n}^{1\rightarrow 2} \\ g_{2,1}^{1\rightarrow 2} & g_{2,2}^{1\rightarrow 2} & \cdots & g_{2,n}^{1\rightarrow 2}\\ \vdots & \vdots & \ddots & \vdots \\ g_{n,1}^{1\rightarrow 2} & g_{n,2}^{1\rightarrow 2} & \cdots & g_{n,n}^{1\rightarrow 2}\end{bmatrix} \begin{bmatrix}u_1' \\ u_2'\\ \vdots\\ u_n'\end{bmatrix} = \begin{bmatrix}u_1 \\ u_2\\ \vdots\\ u_n\end{bmatrix}$$
我们把向量$\mathbf{u}$在原基底$\mathcal{B}_1$和新基底$\mathcal{B}_2$上的坐标分别记为$[\mathbf{u}]_{\mathcal{B}_1}$和$[\mathbf{u}]_{\mathcal{B}_2}$，因此上式也可以简写为：
$$ G^{1\rightarrow 2} [\mathbf{u}]_{\mathcal{B}_2} = [\mathbf{u}]_{\mathcal{B}_1}. $$
直观上，矩阵$G^{1\rightarrow 2}$把向量在基底$\mathcal{B}_2$上的坐标转化为$\mathcal{B}_1$上的坐标。

同理，如果我们将$\mathbf{e}_i$用$\mathbf{e}_1', \mathbf{e}_2', \cdots, \mathbf{e}_n'$的直组合表示，记其系数为$g_{i,j}^{2\rightarrow 1}$，那么，构成的矩阵$G^{2\rightarrow 1} = \left[g_{i,j}^{2\rightarrow 1}\right]_{1\leqslant i , j \leqslant n}$就满足
$$\begin{bmatrix}g_{1,1}^{2\rightarrow 1} & g_{1,2}^{2\rightarrow 1} & \cdots & g_{1,n}^{2\rightarrow 1} \\ g_{2,1}^{2\rightarrow 1} & g_{2,2}^{2\rightarrow 1} & \cdots & g_{2,n}^{2\rightarrow 1}\\ \vdots & \vdots & \ddots & \vdots \\ g_{n,1}^{2\rightarrow 1} & g_{n,2}^{2\rightarrow 1} & \cdots & g_{n,n}^{2\rightarrow 1}\end{bmatrix} \begin{bmatrix}u_1 \\ u_2\\ \vdots\\ u_n\end{bmatrix} = \begin{bmatrix}u_1' \\ u_2'\\ \vdots\\ u_n'\end{bmatrix}$$
以上就是把向量在一个基底上的坐标转为另一个基底上的坐标的方法。下面我们来看坐标转换如何影响直映射的矩阵。

设$T$是从空间$\mathbb{V}$到$\mathbb{W}$的直映射。$\mathcal{B}_1 = (\mathbf{e}_1, \mathbf{e}_2, \cdots, \mathbf{e}_n)$和$\mathcal{C}_1 = (\mathbf{f}_1, \mathbf{f}_2, \cdots, \mathbf{f}_m)$分别是$\mathbb{V}$和$\mathbb{W}$的基底。$T$在这两组基底上的矩阵是：
$$[T]_{\mathcal{B}_1}^{\mathcal{C}_1} = [t_{i, j}]_{1\leqslant i \leqslant m, 1\leqslant j \leqslant n}$$
如果把$\mathcal{B}_1$换成另一组基：$\mathcal{B}_2 = (\mathbf{e}_1', \mathbf{e}_2', \cdots, \mathbf{e}_n')$，$T$在$\mathcal{B}_2$和$\mathcal{C}$上的矩阵是怎样的呢？

我们将新的基底$\mathcal{B}_2$上的矩阵记为：
$$[T]_{\mathcal{B}_2}^{\mathcal{C}_1} = [t_{i, j}']_{1\leqslant i \leqslant m, 1\leqslant j \leqslant n}$$

取向量$\mathbf{u} \in \mathbb{V}$，设它的$\mathcal{B}_1$和$\mathcal{B}_2$分解分别为$\sum_{i=1}^n u_i\mathbf{e}_i $与$\sum_{i=1}^n u_i'\mathbf{e}_i' $。在原来的基底$\mathcal{B}_1$上，$\mathbf{u}$的像可以写成：
$$T(\mathbf{u}) = \sum_{j=1}^m \sum_{i=1}^n t_{j,i} u_i \mathbf{f}_j.$$
在新的基底$\mathcal{B}_2$上，$\mathbf{u}$的像可以写成：
$$T(\mathbf{u}) = \sum_{j=1}^m \sum_{i=1}^n t_{j,i}' u_i' \mathbf{f}_j.$$
把在前面的讨论中用$u_1', u_2', \cdots, u_n'$表示$u_1, u_2, \cdots, u_n$的方式代入上式，就得到
$$T(\mathbf{u}) = \sum_{j=1}^m \sum_{i=1}^n t_{j,i}\left(\sum_{k=1}^n g_{i,k}^{1\rightarrow 2} u_k'\right) \mathbf{f}_j = \sum_{j=1}^m \sum_{k=1}^n \left(\sum_{i=1}^n t_{j,i} g_{i,k}^{1\rightarrow 2}\right) u_k' \mathbf{f}_j.$$

上式右边的括号里的求和，与矩阵乘法的结果很像。我们可以验证，如果设矩阵$S = [s_{i, j}]_{1\leqslant i \leqslant m, 1\leqslant j \leqslant n}$为$T$在原基底上的矩阵右乘基变换矩阵$G^{1\rightarrow 2}$的积：$S = [T]_{\mathcal{B}_1}^{\mathcal{C}_1} \cdot G^{1\rightarrow 2}$，那么：
$$\forall \,\,1\leqslant j \leqslant m, 1\leqslant k \leqslant n, \quad s_{j, k} = \sum_{i=1}^n t_{j,i} g_{i,k}^{1\rightarrow 2}$$
所以，
$$\sum_{j=1}^m \sum_{k=1}^n \left(\sum_{i=1}^n t_{j,i} g_{i,k}^{1\rightarrow 2}\right) u_k' \mathbf{f}_j = \sum_{j=1}^m \sum_{i=1}^n s_{j,i} u_i' \mathbf{f}_j.$$
这说明，$S$就是$T$在新基底$\mathcal{B}_2$上的矩阵\footnote{这一步略去证明，读者可以自行使用特殊值法证明$S$就是$[T]_{\mathcal{B}_2}^{\mathcal{C}_1}$.}：
$$[T]_{\mathcal{B}_2}^{\mathcal{C}_1} = [T]_{\mathcal{B}_1}^{\mathcal{C}_1} \cdot G^{1\rightarrow 2}.$$

这就是直映射矩阵的“换底公式”。

由对称性，我们可以立刻得出：
$$[T]_{\mathcal{B}_1}^{\mathcal{C}_1} = [T]_{\mathcal{B}_2}^{\mathcal{C}_1} \cdot G^{2\rightarrow 1}.$$
两个式子相乘，可以消掉$[T]_{\mathcal{B}_1}^{\mathcal{C}_1}$和$[T]_{\mathcal{B}_2}^{\mathcal{C}_1}$，得到\footnote{这里并不能直接“消掉”两个矩阵，之所以可以这么做，是因为以上关系对任意的$T$都成立。读者可以使用$E_{i,j}$矩阵做一个严谨的证明。}：
$$G^{1\rightarrow 2} \cdot G^{2\rightarrow 1} = G^{2\rightarrow 1} \cdot G^{1\rightarrow 2} = \mathrm{I}_n.$$
这说明，基变换是可逆的，从旧基底变换到新基底，再变换回来，等同于不做任何变换。

从另一个方面看，如果令空间$\mathbb{W} = \mathbb{V}$，$T = \mathrm{I}$，$\mathcal{C}_1 = \mathcal{B}_1$，那么，按照定义，
$$[T]_{\mathcal{B}_1}^{\mathcal{C}_1} = [\mathrm{I}]_{\mathcal{B}_1}^{\mathcal{B}_1} = \mathrm{I}_n.$$
所以，
$$[\mathrm{I}]_{\mathcal{B}_2}^{\mathcal{B}_1} = [\mathrm{I}]_{\mathcal{B}_1}^{\mathcal{B}_1} \cdot G^{1\rightarrow 2} = G^{1\rightarrow 2}.$$
从$\mathcal{B}_1$到$\mathcal{B}_2$的基变换矩阵，可以理解为一个“转换插头”，把新基底上的坐标转化为原基底上的坐标。然后，原基底上的坐标构成的列向量与映射在旧基底上的矩阵相乘，得出结果。
用矩阵的简记法，可以看得更清楚：
\begin{align*}
[T]_{\mathcal{B}_2}^{\mathcal{C}_1} \cdot [\mathbf{u}]_{\mathcal{B}_2} &= [T]_{\mathcal{B}_1}^{\mathcal{C}_1} \cdot [\mathrm{I}]_{\mathcal{B}_2}^{\mathcal{B}_1} \cdot [\mathbf{u}]_{\mathcal{B}_2} \\
&= [T]_{\mathcal{B}_1}^{\mathcal{C}_1} \cdot \left([\mathrm{I}]_{\mathcal{B}_2}^{\mathcal{B}_1} \cdot [\mathbf{u}]_{\mathcal{B}_2}\right) \\
&= [T]_{\mathcal{B}_1}^{\mathcal{C}_1} \cdot [\mathbf{u}]_{\mathcal{B}_1}.
\end{align*}
我们还可以探讨另一个对称的性质：把到达空间$\mathcal{C}_1$换成另一组基$\mathcal{C}_2$，对直映射矩阵有什么影响？用类似的方式，可以算出：
$$[T]_{\mathcal{B}_1}^{\mathcal{C}_2} = [\mathrm{I}]_{\mathcal{C}_1}^{\mathcal{C}_2} \cdot [T]_{\mathcal{B}_1}^{\mathcal{C}_1} .$$
矩阵$[\mathrm{I}]_{\mathcal{C}_1}^{\mathcal{C}_2}$在另一端起到“转换插头”的作用，把映射结果从原基底坐标转为新基底坐标：
\begin{align*}
[T]_{\mathcal{B}_1}^{\mathcal{C}_2} \cdot [\mathbf{u}]_{\mathcal{B}_1} &= [\mathrm{I}]_{\mathcal{C}_1}^{\mathcal{C}_2} \cdot [T]_{\mathcal{B}_1}^{\mathcal{C}_1} \cdot [\mathbf{u}]_{\mathcal{B}_1} \\
&= [\mathrm{I}]_{\mathcal{C}_1}^{\mathcal{C}_2} \cdot \left([T]_{\mathcal{B}_1}^{\mathcal{C}_1} \cdot [\mathbf{u}]_{\mathcal{B}_1}\right) \\
&= [\mathrm{I}]_{\mathcal{C}_1}^{\mathcal{C}_2} \cdot [T\mathbf{u}]_{\mathcal{C}_1} \\
&= [T\mathbf{u}]_{\mathcal{C}_2}.
\end{align*}

\section{逆矩阵}

计算基变换矩阵的时候，我们得到了一个“副产品”：
$$G^{1\rightarrow 2} \cdot G^{2\rightarrow 1} = G^{2\rightarrow 1} \cdot G^{1\rightarrow 2} = \mathrm{I}_n.$$
我们把满足$AB=BA=\mathrm{I}_n$的两个矩阵$A$、$B$称为\textbf{互逆矩阵}，把其中一个叫做另一个的\textbf{逆矩阵}。上式中，$G^{1\rightarrow 2}$是$G^{2\rightarrow 1}$的逆矩阵，$G^{2\rightarrow 1}$是$G^{1\rightarrow 2}$的逆矩阵。如果$B$是$A$的逆矩阵，一般将$B$记作$A^{-1}$.

不是所有的矩阵都有逆矩阵。首先，要使得$AB=BA=\mathrm{I}_n$成立，$A$和$B$必须都是$n\times n$的方阵。
其次，如果我们将这两个方阵看成$\mathbb{F}^n$上的自同态在标准基上的矩阵，那么，它们必然都是双射，也就是自同构。
我们把有逆矩阵的矩阵称为\textbf{可逆矩阵}。

另一方面，如果矩阵$A$是某个$\mathbb{F}^n$上自同构（记为$f$）在某个基底$\mathcal{B}$上的矩阵：$A = [f]_{\mathcal{B}}^{\mathcal{B}}$，那么，由于$f$是双射，必然存在逆映射，而且逆映射也是直映射，即$f^{-1}$也是自同构。设它在同一个基底上的矩阵为$B$，那么由于
$$\forall \mathbf{u} \in \mathbb{F}^n, \quad \mathbf{u} = f^{-1}(f(\mathbf{u})) = f(f^{-1}(\mathbf{u})),$$
所以相应地，矩阵$A$和$B$也满足：
$$\forall \mathbf{u} \in \mathbb{F}^n, \quad [\mathbf{u}]_{\mathcal{B}} = BA[\mathbf{u}]_{\mathcal{B}} = AB[\mathbf{u}]_{\mathcal{B}},$$
这说明$AB=BA=\mathrm{I}_n.$

可以说，逆映射和逆矩阵是同一个概念在自同态和矩阵上的表现。
给定基底后，可逆映射对应着可逆矩阵，而它的逆映射也对应着它的矩阵的逆矩阵。

因此，可逆矩阵$AB=BA=\mathrm{I}_n$的条件可以和可逆映射一样，进一步放宽，去掉其中一个等式约束。

\begin{pp}\label{pp:5_1}
    给定方阵$A$，如果存在$B$，使得$AB = \mathrm{I}_n$，则$BA = \mathrm{I}_n$；如果存在$B$，使得$BA = \mathrm{I}_n$，则$AB = \mathrm{I}_n$。
\end{pp}
\begin{proof} % [{\bfseries 证明\ref{pp:5_1}}]
给定$\mathbb{F}^n$的某个基底，考虑$A$和$B$对应的自同态：$A = [f]_{\mathcal{B}}^{\mathcal{B}}, B = [g]_{\mathcal{B}}^{\mathcal{B}}$。
$AB = \mathrm{I}_n$等价于说复合映射$f\circ g$是恒等映射。
因此，根据命题\ref{pp:3_3_3}，$f$和$g$都是自同构。任取$\mathbf{u} \in \mathbb{F}^n$，存在$\mathbf{v} \in \mathbb{F}^n$使得$g(\mathbf{v}) = \mathbf{u}$，因此，
$$g(f(\mathbf{u})) = g(f(g(\mathbf{v}))) = g(\mathbf{v}) = \mathbf{u}.$$
因此，$g\circ f$是恒等映射。对应到矩阵，即$BA = \mathrm{I}_n$。类似地，$BA = \mathrm{I}_n$说明$g\circ f$是恒等映射，因此$f\circ g$是恒等映射，$AB = \mathrm{I}_n$。
\end{proof}


和逆映射一样，逆矩阵如果存在，必然是唯一的。给定矩阵$A$，如果有$B$和$C$满足$AB=BA=\mathrm{I}_n$，$AC=CA=\mathrm{I}_n$，
那么$C = CAB = B$。

如何判断矩阵是否可逆，乃至于获得矩阵的逆矩阵呢？

讨论一次方程组的时候，我们证明了，如果方程组有解，使用消元法，可以在有限步内得到一次方程组$ A\mathbf{x} = \mathbf{b}$的一组解。
从矩阵的角度来看，如果方阵$A$是可逆矩阵，那么可以在两边左乘$A$的逆矩阵，得到
$$ \mathbf{x} = A^{-1}\mathbf{b}.$$
所以一次方程组有唯一解。我们通过消元法得到的解，就是$A^{-1}\mathbf{b}$的值。
另一方面，如果$A$不是可逆矩阵，说明对应的自同态不是双射，从而不是满射也不是单射。
对于不同的向量$\mathbf{b}$，一次方程组不一定有解，或者有多于一组解。

于是，我们有了一种获取逆矩阵的算法：让$\mathbf{b}$取遍$\mathbb{F}^n$的标准基$\mathcal{C} = (\mathbf{c}_i)_{1\leqslant i \leqslant n}$，求解$n$个一次方程组：
$$ A\mathbf{x} = \mathbf{c}_i$$
如果它们都有唯一解，我们就得到所有的$A^{-1}\mathbf{c}_i$，也就是矩阵$A^{-1}$每一列的系数，从而得到$A^{-1}$的值。
如果某些方程组没有解或有多于一组解，我们就知道$A$对应的自同态不是满射或不是单射，从而不是双射，
也就是说，$A$不是可逆矩阵。

此外，如果$A$是可逆矩阵，那么每次的一次方程组都有唯一解，所以不会出现“任意填充”。
因此，用消元法求出的逆矩阵的元素，也必然是系数域中的元素。
不会出现矩阵可逆，但逆矩阵不在$\mathfrak{M}_n(\mathbb{F})$中的情况。

\section{平直群}
可逆映射，或者说自同构，及其对应的可逆矩阵，让我们想到一种称为群的代数结构。
群的定义可参见定义\ref{df:0_1}。简单来说，配备了二元运算$\,\,\cdot\,\,$的集合$G$如果
满足结合律、存在单位元和普遍可逆性，就构成一个群。如果我们考虑配备了映射的复合的$n$维空间$\mathbb{V}$上自同构
的集合，或配备了矩阵乘法的$n\times n$可逆矩阵集合，可以验证它们都构成群。这里我们用矩阵集合作为例子：
\begin{enumerate}
    \item 可逆矩阵的乘积仍然是可逆矩阵，集合对于矩阵乘法运算封闭。
    \item 矩阵乘法满足结合律。
    \item 单位矩阵乘以任何可逆矩阵$P$等于它自身：$P \cdot \mathrm{I}_n = \mathrm{I}_n \cdot P = P.$
    \item 可逆矩阵的逆矩阵是它关于矩阵乘法的逆元。$P \cdot P^{-1} = P^{-1} \cdot P = \mathrm{I}_n.$
\end{enumerate}
类似地，也可以验证配备了映射的复合的$n$维空间自同构的集合是群（参见附录一定义\ref{df:0_1}）。

不难看出，这两个群是同构的。给定了空间的基底$\mathcal{B}$后，可以构造群同构：
$$ \pi : T \mapsto [T]_\mathcal{B}.$$
因此我们将其视为同一个群来讨论。我们称这个群为（系数域$\mathbb{F}$上的）\textbf{$n$阶平直群}或\textbf{$n$阶一般平直群}，记为$\mathrm{GL}_n(\mathbb{F})$。

\section{相似矩阵}

我们知道，一般情况下，基变换前后的矩阵有如下关系：
$$[T]_{\mathcal{B}_2}^{\mathcal{C}_2} = [\mathrm{I}]_{\mathcal{C}_1}^{\mathcal{C}_2} \cdot [T]_{\mathcal{B}_1}^{\mathcal{C}_1} \cdot [\mathrm{I}]_{\mathcal{B}_2}^{\mathcal{B}_1} .$$
如果我们只考虑自同态，令$\mathcal{C}_1 = \mathcal{B}_1$，$\mathcal{C}_2 = \mathcal{B}_2$，那么，上式变为：
$$[T]_{\mathcal{B}_2}^{\mathcal{B}_2} = [\mathrm{I}]_{\mathcal{B}_1}^{\mathcal{B}_2} \cdot [T]_{\mathcal{B}_1}^{\mathcal{B}_1} \cdot [\mathrm{I}]_{\mathcal{B}_2}^{\mathcal{B}_1} .$$
而我们还知道，$[\mathrm{I}]_{\mathcal{B}_1}^{\mathcal{B}_2}$和$ [\mathrm{I}]_{\mathcal{B}_2}^{\mathcal{B}_1}$互为逆矩阵。不失一般性，让我们把$[\mathrm{I}]_{\mathcal{B}_1}^{\mathcal{B}_2}$记作$P$，那么上式变为：
$$[T]_{\mathcal{B}_2}^{\mathcal{B}_2} = P \cdot [T]_{\mathcal{B}_1}^{\mathcal{B}_1} \cdot P^{-1} .$$
这就是自同态在不同基底上的矩阵的关系。我们把这个关系称为\textbf{相似关系}。

\begin{df}\label{df:5_1}
给定$n\times n$矩阵$A$、$B$，如果存在可逆矩阵$P$，使得$A = PBP^{-1}$，就说$A$和$B$是\textbf{相似}的矩阵，或者说$A$与$B$之间存在相似关系。
\end{df}
可以验证，相似关系是一种等价关系。同一个自同态在不同基底上的矩阵相似。反过来，相似的矩阵是否一定是某个自同态在不同基底上的矩阵呢？这等价于问：可逆矩阵是否一定是基变换矩阵。我们可以直接证明后者。

\begin{pp}\label{pp:5_2}
    可逆矩阵一定是基变换矩阵。
\end{pp}
\begin{proof} % [{\bfseries 证明\ref{pp:5_2}}]
给定可逆矩阵$P = [p_{i,j}]_{1\leqslant i , j \leqslant n}$，
按定义，存在矩阵$Q = [q_{i,j}]_{1\leqslant i , j \leqslant n}$，
使得$PQ = QP = \mathrm{I}_n$.

给定$n$维平直空间$\mathbb{V}$和它的基底$\mathcal{B} = (\mathbf{e}_1, \mathbf{e}_2, \cdots , \mathbf{e}_n)$，定义一组向量$\mathcal{B}' = (\mathbf{e}_1', \mathbf{e}_2', \cdots , \mathbf{e}_n')$：
$$\forall 1 \leqslant i \leqslant n, \quad \mathbf{e}_i' = \sum_{k=1}^n p_{k,i} \mathbf{e}_k.$$
如果$(\mathbf{e}_1', \mathbf{e}_2', \cdots , \mathbf{e}_n')$是$\mathbb{V}$的另一组基底，那么按照基变换矩阵的构造方法，$P$就是一个基变换矩阵：$P = [\mathrm{I}]_{\mathcal{B}'}^{\mathcal{B}}$.

下面证明$(\mathbf{e}_1', \mathbf{e}_2', \cdots , \mathbf{e}_n')$是$\mathbb{V}$的基底。

$PQ = \mathrm{I}_n$，所以，$\forall \mathbf{u} \in \mathbb{V}$，$\mathbf{u} = (PQ)(\mathbf{u})$ ，即：
$$\mathbf{u} = \sum_{j=1}^m \sum_{i=1}^n \left(\sum_{k=1}^n p_{j,k} q_{k, i} \right) u_i \mathbf{e}_j.$$
整理可得：
\begin{align*}
    \mathbf{u} &= \sum_{k=1}^n \sum_{i=1}^n q_{k, i}  u_i \left( \sum_{j=1}^n p_{j,k} \mathbf{e}_j\right)  \\
    &= \sum_{k=1}^n \left(\sum_{i=1}^n q_{k, i}  u_i\right) \mathbf{e}_k'. 
\end{align*}

也就是说，$\mathbf{u}$总能写成$(\mathbf{e}_1', \mathbf{e}_2', \cdots , \mathbf{e}_n')$的直组合。
$\mathcal{B}'$生成空间$\mathbb{V}$，且有$n$个元素，所以$\mathcal{B}'$是$\mathbb{V}$的基底。
\end{proof}

相似矩阵刻画了直自同态在不同基底上的矩阵之间的关系。每个直自同态，都对应一个$n\times n$矩阵在相似关系下的等价类。
相似的矩阵可以看作是同一个直自同态在不同基底上的“画像”。画像的角度可以不同，但对象是一样的。

关于相似矩阵的另一个性质：相似矩阵的迹相同。设$A$为矩阵，$P$为可逆矩阵。
$$ \ji(P^{-1}AP) = \ji(P^{-1}(AP)) = \ji(APP^{-1}) = \ji(A). $$
这说明直自同态在任何基底上的迹都相同。因此，迹是直自同态的一个内蕴性质，我们可以定义直自同态的迹：
\begin{df}{\textbf{直自同态的迹}}
    有限维平直空间上的直自同态的迹，是它在任意基底上的矩阵的迹。
\end{df}

\sectionext{交换空间}
一般来说，直映射的复合和矩阵乘法都不满足交换律。但是，从个例来说，对某些直映射和矩阵，总存在一些直映射或矩阵，与它们满足交换律。
比如，对可逆映射$T$，$TT^{-1} = T^{-1}T = \mathrm{I}$，所以它的逆映射与它满足交换律。显然，要讨论交换，直映射必须是直自同态，矩阵必须是方阵。因此，以下讨论中可能用到的“矩阵”、“自同态”或“映射”，都指方阵和直自同态。

\begin{df}\label{df:5_2}
    给定直自同态$T\in \mathcal{L}(\mathbb{V})$，所有与$T$可交换的直自同态$S\in\mathcal{L}(\mathbb{V})$的集合称为$A$的\textbf{交换集}，记为$\Cmu{T}$.\\
    给定直自同态的集合$\mathcal{S}$，直自同态若与$\mathcal{S}$中每个元素都交换，就称为与集合$\mathcal{S}$交换。\\
    所有与$\mathcal{S}$交换的直自同态的集合称为$\mathcal{S}$的\textbf{交换集}，记为$\Cmu{\mathcal{S}}$。\\
    对于$\mathfrak{M}_n(\mathbb{R})$中的矩阵，可以类似定义其交换集。
\end{df}
两个直自同态可交换，则在给定的基底上，它们的矩阵也可交换。两个矩阵可交换，那么可以构造两个直映射，
使得它们在特定基底上的矩阵是这两个矩阵，从而这两个自同态也可交换。
矩阵可交换和自同态可交换的概念可以说是等价的。
因此，我们接下来讨论交换关系和交换集的性质可能只用到两者之一，但对另一者也自动成立。

最简单的可交换映射是恒等映射和零映射，对应单位矩阵$\mathrm{I}_n$和零矩阵$0$。它们和任意映射交换，所以它们的交换集都是全体自同态，也属于任何映射的交换集。

对其他的映射来说，交换集有哪些元素呢？

每个映射都和自己交换。给定映射$A$，$A^2, A^3, \cdots ...$也和$A$交换，其中
$$ A^k = \overbrace{A\circ A \circ \cdots \circ A}^{k\mbox{\scriptsize 次}}$$
称为$A$的$k$次幂。约定$A$的$0$次幂是单位矩阵$\mathrm{I}_n$\footnote{如果$A$是可逆映射，
还可以定义$A^{(-k)} = \left(A^{-1}\right)^k$。}。
所以，任何映射的交换集都包括它自身和它的乘方。

如果映射$B$和$C$都和$A$交换，那么
$$ (B+C)\circ A = B\circ A + C\circ A = A\circ B + A\circ C = A\circ (B+C).$$
也就是说，它们的和仍然与$A$交换。另外，如果映射$B$与$A$交换，那么对任意系数$\lambda$，$\lambda B$仍然与$A$交换。这说明，映射或映射之集合的交换集是平直空间。所以，我们也可以称交换集为\textbf{交换空间}。

此外，如果映射$B$和$C$都和$A$交换，那么
$$ B\circ C\circ A = B\circ A\circ C = A\circ B\circ C, \,\,\,\, C\circ B\circ A = C\circ A\circ B = A\circ C\circ B.$$
所以$B\circ C$、$C\circ B$都与$A$交换。
这说明交换空间也是域上场。我们称之为\textbf{交换域上场}或\textbf{交换代数}。

要注意的是，映射的复合（矩阵乘法）的交换性并没有传递性。具体来说，如果映射$B$与$A$交换，$C$与$B$交换，$C$与$A$不一定交换。我们可以构造一个简单的反例：设矩阵$B$为恒等映射或零映射，则对任意映射$A$和$C$，$B$与$A$交换，$C$与$B$交换，但我们知道$C$与$A$不一定交换。

从以上几个例子出发，给定$A \in \mathcal{L}(\mathbb{V})$，
如果考虑以$\mathbb{F}$中元素为系数的一元整式$p(z) = c_0 + c_1 z + \cdots + c_m z^m$，
$p(z)$对应一个由$A$经过加、减、复合得到的矩阵：
$$ p(A) = c_0 \mathrm{I}_n + c_1 A + \cdots + c_m A^m.$$
根据前面的例子，我们知道，$p(A) \in \Cmu{A}$。

可以验证，映射$\Psi_A: p \mapsto p(A)$是$\mathbb{F}$上的一元整式环$\mathbb{F}[z]$到$\mathfrak{M}_n(\mathbb{F})$的直映射，同时也是一个环同态\footnote{见定义\ref{df:0_5_2}.}。也就是说，所有“$A$构成的整式”构成一个域上场，一般记为$\mathbb{F}[A]$。
$\mathbb{F}[A]$是交换域上场，而且$\mathbb{F}[A] \subseteq \Cmu{A}$。

相似矩阵可以看作同一个直映射在不同基底上的“画像”，那么相似矩阵的交换空间有什么关系呢？
\begin{pp}\label{pp:5_3}
    如果$A, B \in \mathfrak{M}_n(\mathbb{F})$是相似矩阵，存在可逆矩阵$P$使得$AP = PB$，那么$\Cmu{A}$与$\Cmu{B}$同构。
\end{pp}
\begin{proof} % [{\bfseries 证明\ref{pp:5_3}}]
    给定$M \in \Cmu{A}$，我们有$ AM = MA$。
    将$A = PBP^{-1}$代入上式，得到：
    $$ PBP^{-1}M = MPBP^{-1}.$$
    所以等式两边的同时左乘$P^{-1}$，右乘$P$，就得到：
    $$  BP^{-1}MP = P^{-1}MPB. $$
    也就是说，$ P^{-1}MP \in \Cmu{B}.$\\
    于是，构造映射：
    $$ f : M \mapsto P^{-1}MP.$$
    容易验证$f$是直映射。另外，可以构造映射：
    $$ g : M \mapsto PMP^{-1}.$$
    不难验证$g$是$f$的逆映射，把$\Cmu{B}$中的元素映射到$\Cmu{A}$中元素。\\
    综上可知，$f$是双射，即同构映射，因此$\Cmu{A}$与$\Cmu{B}$同构。\\
\end{proof}

用具体的例子来理解这个同构：如果$A$和$B$是相似矩阵，$AP = PB$，考虑$A$的整式：
$$p(A) = c_0\mathrm{I}_n + c_1 A + \cdots + c_m A^m,$$
将它用$B$和$P$表示。
$$A^i = \left(PBP^{-1}\right)^i = PBP^{-1} PBP^{-1} \cdots PBP^{-1} = PB^iP{-1}.$$
所以，
\begin{align*}
    p(A) &= c_0\mathrm{I}_n + c_1 PBP^{-1} + \cdots + c_m PB^mP{-1}  \\
    &= P\left(c_0\mathrm{I}_n + c_1 B + \cdots + B^m\right)P^{-1}  \\
    &= Pp(B)P^{-1}.  
\end{align*}
也就是说，$p(A)$和$p(B)$也是相似矩阵。

\section{对偶映射的矩阵}
设$T \in \mathcal{L}(\mathbb{V}, \mathbb{W})$，$T$的对偶映射把$\mathbb{W}^*$中元素$\varphi$映射为$\varphi \circ T$。
给定空间$\mathbb{V}$上的基$\mathcal{B}_\mathbb{V} = (\mathbf{e}_1, \cdots, \mathbf{e}_n)$
和$\mathbb{W}$上的基$\mathcal{B}_\mathbb{W} = (\mathbf{f}_1, \cdots, \mathbf{f}_m)$，又记它们在各自对偶空间的对偶基为$\mathcal{B}_\mathbb{V}^* = (\mathbf{e}_1^*, \cdots, \mathbf{e}_n^*)$和$\mathcal{B}_\mathbb{W}^* = (\mathbf{f}_1^*, \cdots, \mathbf{f}_m^*)$ 。
考虑任意$\mathbf{u} = \sum_{j=1}^n u_j \mathbf{e}_j \in \mathbb{V}$，它经过映射$T$得到：
$$ T(\mathbf{u}) =  \sum_{i=1}^m \sum_{j=1}^n u_j t_{i,j} \mathbf{f}_i.$$
其中$(t_{i, j})_{1\leqslant i\leqslant m, 1\leqslant j\leqslant n}$是映射$T$在对应基底上的矩阵。所以，
$$ \forall 1\leqslant i\leqslant m, \quad f_i^*\circ T(\mathbf{u}) = \sum_{j=1}^n u_j t_{i,j}.$$
而$\mathbf{e}_j^*(\mathbf{u}) = u_j$，所以，
$$ \forall 1\leqslant i\leqslant m, \quad f_i^*\circ T(\mathbf{u}) = \sum_{j=1}^n u_j t_{i,j} = \sum_{j=1}^n t_{i,j} \mathbf{e}_j^*(\mathbf{u}).$$
这说明$T^*(f_i^*) = \sum_{j=1}^n t_{i,j} \mathbf{e}_j^* $. 
考虑$T^*$在基底$\mathcal{B}_\mathbb{W}^*$和$\mathcal{B}_\mathbb{V}^*$上的矩阵表示。
它第$i$行第$j$列的元素就是$t_{j,i}$。也就是说，这个矩阵可以记为$[t_{j,i}]_{1\leqslant i\leqslant m, 1\leqslant j\leqslant n}$。
我们在前面已经定义过，这个矩阵称为$[t_{i,j}]_{1\leqslant i\leqslant m, 1\leqslant j\leqslant n}$的转置矩阵。

以上的推导说明，
\begin{pp}\label{pp:5_4}
对偶映射的矩阵是原映射的转置矩阵。
\end{pp}

我们知道，给定基底时，直映射可以用矩阵来表示。前面提到，矩阵的列空间就是对应直映射的像空间。而对偶映射的矩阵是原映射的转置矩阵。
所以，使用对偶映射的概念，我们可以这么理解矩阵的行空间：矩阵的行空间就是对应直映射的对偶映射的像空间。
这里矩阵和直映射的关系可以是特定设定的基底，也可以直接在$\mathbb{F}^n$、$\mathbb{F}^m$及其对应的标准基框架下讨论。

前面已经证明，对偶映射的容和原映射的容相等。所以，考虑对应的矩阵的话，我们可以得出：

\begin{pp}\label{pp:5_5}
矩阵的行空间与列空间的维数相等。
\end{pp}
我们把这个维数称为\textbf{矩阵的容}。

矩阵的容是从直映射的容抽象而来。给定直映射和空间基底，它的矩阵的容就是它的容。

\chapter{本征值和自同态约简}
从向量到平直空间，再到直映射，我们对平直空间的认识不断加深。
这一章中，我们会研究从平直空间到它自身的直映射，也就是（直）自同态。
我们会通过称为本征分析的手段，约简自同态，也就是说，找出用比较简单的形式表示自同态的办法。
我们会看到，约简自同态可以帮助我们更好地看清楚它的内在性质。

\begin{center}
    \fbox{
        \shortstack[l]{
      下文中，除非另有声明，我们总假设$\mathbb{F}$是常见数域如\\
      $\mathbb{Q}$、$\mathbb{R}$或$\mathbb{C}$，
      $\mathbb{V}$是$\mathbb{F}$上的有限维平直空间，维数$n$大于等于一。
        }
    }
\end{center}

\section{子空间中的自同态}
\subsection{不变子空间}
考虑从平直空间$\mathbb{V}$映射到自身的直映射：$T \in \mathcal{L}(\mathbb{V})$. 
如果我们能够将$\mathbb{V}$分解成一些子空间的直和：
$$ \mathbb{V} = \mathbb{U}_1 \oplus \mathbb{U}_2 \oplus \cdots \oplus \mathbb{U}_k. $$
而映射$T$可以限制在每个子空间内部：
$$ T(\mathbb{U}_i) \subseteq \mathbb{U}_i$$
那么，我们可以把$T$分解，只研究$T$在维数更小的空间中的性质。
然而，不是每一个把$\mathbb{V}$分解成子空间直和的方式，
都能对应让$T$的作用恰好限制在子空间内部。

我们把能够让$T$的作用恰好限制在其内部的子空间$\mathbb{U}$称为$T$的不变子空间。

\begin{df}\label{df:6_1}
    设$T$为从$\mathbb{V}$映射到其自身的直映射。
    如果$\mathbb{V}$的某个子空间$\mathbb{U}$满足： 
    $$ \forall \,\, \mathbf{u} \,\, \in \,\, \mathbb{U}, \,\,\, T\mathbf{u} \,\, \in \,\, \mathbb{U} $$
    就称$\mathbb{U}$为$T$（作用下）的\textbf{不变子空间}，简称$T$-不变子空间。
    或者说$\mathbb{U}$在$T$的作用下不变。
    $T$在$\mathbb{U}$上的限制$\restr{T}{\mathbb{U}}$是一个从$\mathbb{U}$到$\mathbb{U}$的自同态：
    $$ \forall \,\, \mathbf{u} \,\, \in \,\, \mathbb{U}, \,\,\, \restr{T}{\mathbb{U}}\mathbf{u} = T\mathbf{u}. $$
\end{df}

$\{\mathbf{0}\}$和$\mathbb{V}$显然是任意$T$的不变子空间。

此外，直映射$T$的零空间$\Ker T$是它的不变子空间，
因为
$$\forall \,\, \mathbf{u} \,\, \in \,\, \Ker T, \,\,\, T\mathbf{u} = \mathbf{0} \,\, \in \,\, \Ker T. $$
直映射$T$的像空间$\Img T$也是它的不变子空间，因为
$$\forall \,\, \mathbf{u} \,\, \in \,\, \Img T, \,\,\, T\mathbf{u} \in \,\, \Img T. $$

我们知道，直映射把平直空间映射为平直空间。
所以，如果$\mathbb{V}$的某个子空间$\mathbb{U}$在$T$的作用下不变，
那么它在$T$下的像是它自己的子空间。

我们用一个具体的例子来感受不变子空间。考虑$\mathbb{R}^2$上的自同态$A$，它在标准基上的矩阵为：
$$
A = \begin{bmatrix} 4 & 3 \\ -2 & -1 \end{bmatrix}.
$$
考虑向量$\mathbf{u}_1 = (1,-1)$：
$$
A \mathbf{u}_1 = \begin{bmatrix} 4 & 3 \\ -2 & -1 \end{bmatrix}\begin{bmatrix} 1 \\ -1 \end{bmatrix} = \begin{bmatrix} 4-3 \\ -2+1 \end{bmatrix} = \begin{bmatrix} 1 \\ -1 \end{bmatrix} = \mathbf{u}_1.
$$
所以$A\mathbf{u}_1 = 1\cdot \mathbf{u}_1 \in \mathbb{R}\mathbf{u}_1$。
由直性可知$A$把$\mathbb{R}\mathbf{u}_1$中任何向量仍映射到$\mathbb{R}\mathbf{u}_1$中。因此$\mathbb{R}\mathbf{u}_1$是$A$的不变子空间。

再考虑子空间$\mathbf{u}_2 = (0,1)$：
$$
A \mathbf{u}_2 = \begin{bmatrix} 4 & 3 \\ -2 & -1 \end{bmatrix} \begin{bmatrix} 0 \\ 1 \end{bmatrix} = \begin{bmatrix} 3 \\ -1 \end{bmatrix}.
$$
向量$(3,-1)$不在$\mathbb{R}\mathbf{u}_2$中，所以$\mathbb{R}\mathbf{u}_2$不是$A$的不变子空间。

可以看到，不是所有子空间都是给定自同态的不变子空间。寻找不变子空间，是约简自同态的第一步。

\subsection{本征值和本征向量}
直映射的不变子空间有什么性质呢？我们首先从1维的不变子空间开始研究。

1维空间就是我们通常理解的直线。$\mathbb{V}$的所有1维子空间$\mathbb{U}$可以这么描述：
$$\mathbb{U} = \{a\mathbf{u} | a \in \mathbb{F}\}.$$
其中$\mathbf{u}$是$\mathbb{U}$中某个非零向量。

\begin{pp}\label{pp:6_1}
    如果某个1维子空间$\mathbb{U}$在$T$作用下不变，那么必然存在某个$\lambda \in \mathbb{F}$，使得
    $$ \forall \,\, \mathbf{u} \,\, \in \,\,\mathbb{U} ,\,\,\, T\mathbf{u} = \lambda\mathbf{u}.$$
\end{pp}
也就是说$\restr{T}{\mathbb{U}}$可以看作缩放映射。
\begin{proof} % [{\bfseries 证明\ref{pp:6_1}}]
    设有非零向量$\mathbf{u}, \,\, \mathbf{v} \in \mathbb{U}$。记$T(\mathbf{u}) = \lambda_1 \mathbf{u} , \,\,\, T(\mathbf{v}) = \lambda_2 \mathbf{v}$。\\
    $\mathbb{U}$是1维空间，所以$\exists ! a \neq 0$使得$\mathbf{v} = a\mathbf{u}$. 一方面，
    $$ T(\mathbf{v}) = T(a\mathbf{u}) = a T(\mathbf{u}) = a\lambda_1 \mathbf{u}$$
    另一方面
    $$ T(\mathbf{v}) = \lambda_2 \mathbf{v} = a\lambda_2 \mathbf{u}$$
    这说明$\lambda_1 = \lambda_2$。记这个参数为$\lambda$，则
    $$\forall \,\, \mathbf{u} \in \mathbb{U}, \,\,\,T(\mathbf{u}) = \lambda\mathbf{u}$$
    它就是我们要找的系数（至于零向量，显然$T(\mathbf{0}) = \mathbf{0} = \lambda \mathbf{0}$）。\\
\end{proof}

从以上证明，我们还可以顺便推出：如果两个非零向量$\mathbf{u}_1$、$\mathbf{u}_2$使得$T(\mathbf{u}_1) = \lambda_1\mathbf{u}_1, \,\,\, T(\mathbf{u}_2) = \lambda_2\mathbf{u}_2$，
$\lambda_1 \neq \lambda_2$，那么这两个向量必然直无关。它们对应的1维不变子空间交集为$\{\mathbf{0}\}$.

1维不变子空间对我们了解直自同态很有帮助。只要映射有1维不变子空间，映射在其中就是简单的缩放变换。
我们用1维不变子空间中的非零向量和对应的“缩放倍数”来代表它。

\begin{df}\label{df:6_2}
    设$T$为从$\mathbb{V}$映射到其自身的直映射（$\mathbb{V}$上的自同态）。
    如果$\exists \mathbf{u} \in \mathbb{V}, \mathbf{u} \neq \mathbf{0}$，以及$\lambda \in \mathbb{F}$，使得： 
    $$ T(\mathbf{u}) = \lambda \mathbf{u} $$
    就称$\mathbf{u}$为$T$的\textbf{本征向量}，
    对应\textbf{本征值}$\lambda$。
\end{df}
要注意的是，本征向量必须是非零向量。

自同态$T$有本征值$\lambda$，说明$T - \lambda \mathrm{I}$不是单射，反之亦然。
其零空间$\Ker (T - \lambda \mathrm{I})$是$T$作用下的不变子空间，
维数至少是1，但不一定是1. 比如，考虑缩放映射$T = 2\mathrm{I}$，
其零空间$\Ker (T - 2 \mathrm{I})$是整个空间$\mathbb{V}$。

所以，1维不变子空间可以用本征值和本征向量来代表。但这些概念并不是等价的。
一个本征向量对应唯一的本征值，产生唯一的1维不变子空间。
一个1维不变子空间对应着唯一的本征值，但对应着不止一个本征向量（相差一个标量系数）。
一个本征值$\lambda$可以对应多个本征向量，甚至多个1维不变子空间，或者说维数大于1的不变子空间。
考虑$\Ker (T - \lambda \mathrm{I})$，设其维数为$k$，则通过它的每一组基底，都可以产生$k$个1维不变子空间，
其直和为$\Ker (T - \lambda \mathrm{I})$。

我们来看具体的例子。仍然考虑$\mathbb{R}^2$上的自同态$A$，其矩阵为：
$$
A = \begin{bmatrix} 4 & 3 \\ -2 & -1 \end{bmatrix}.
$$
我们验证向量$\mathbf{u}_1 = (1,-1)$是否是本征向量。直接计算：
$$
A\mathbf{u}_1 = \begin{bmatrix} 4 & 3 \\ -2 & -1 \end{bmatrix}\begin{bmatrix} 1 \\ -1 \end{bmatrix} = \begin{bmatrix} 4-3 \\ -2+1 \end{bmatrix} = \begin{bmatrix} 1 \\ -1 \end{bmatrix} = 1 \cdot \mathbf{u}_1.
$$
$A$把$\mathbf{u}_1$映射为它自身的$1$倍。所以$\mathbf{u}_1$是$A$的本征向量，对应本征值$\lambda = 1$。

再验证向量$\mathbf{u}_2 = (3,-2)$：
$$
A \mathbf{u}_2 = \begin{bmatrix} 4 & 3 \\ -2 & -1 \end{bmatrix} \begin{bmatrix} 3 \\ -2 \end{bmatrix} = \begin{bmatrix} 12-6 \\ -6+2 \end{bmatrix} = \begin{bmatrix} 6 \\ -4 \end{bmatrix} = 2 \cdot \mathbf{u}_2.
$$
$A$把$\mathbf{u}_2$映射为它自身的$2$倍。所以$\mathbf{u}_2$也是$A$的本征向量，对应本征值$\lambda = 2$。

$A$有两个不同的本征值$1$和$2$，各自对应一个1维本征空间。
两个本征向量$(1,-1)$和$(3,-2)$对应不同的本征值，按照命题\ref{pp:6_2}的结论，它们直无关，可以构成$\mathbb{R}^2$的基。

再来看另一个自同态$B$，其矩阵为：
$$
B = \begin{bmatrix} 3 & -4 \\ 1 & -1 \end{bmatrix}.
$$
验证向量$\mathbf{v}=(2,1)$：
$$
B\mathbf{v} =  \begin{bmatrix} 3 & -4 \\ 1 & -1 \end{bmatrix}\begin{bmatrix} 2 \\ 1 \end{bmatrix} = \begin{bmatrix} 6-4 \\ 2-1 \end{bmatrix} = \begin{bmatrix} 2 \\ 1 \end{bmatrix} = 1 \cdot \mathbf{v}.
$$
所以$\mathbf{v}$是$B$的本征向量，对应本征值$\lambda = 1$。

现在来看$B$关于本征值$1$的本征空间。按照定义，它是$\Ker(B - \mathrm{I})$。计算：
$$
B - \mathrm{I} = \begin{bmatrix} 2 & -4 \\ 1 & -2 \end{bmatrix}.
$$
解方程$(B - \mathrm{I})\mathbf{u} = \mathbf{0}$：
$$
\begin{bmatrix} 2 & -4 \\ 1 & -2 \end{bmatrix} \begin{bmatrix} u_1 \\ u_2 \end{bmatrix} = \begin{bmatrix} 0 \\ 0 \end{bmatrix}.
$$
两个方程分别是$2u_1 - 4u_2 = 0$和$u_1 - 2u_2 = 0$，它们是等价的。
所以解得$u_1 = 2u_2$，本征空间为$\mathbb{R}\mathbf{v}$，维数为$1$。

$B$还有没有$\mathbb{R}\mathbf{v}$以外的本征向量呢？假设有$\mathbb{R}\mathbf{v}$以外的非零向量$\mathbf{v}'=(v_1', v_2')$和对应的本征值$\lambda'$使得$B \mathbf{v}' = \lambda' \mathbf{v}'$，代入得到：
\begin{align*}
    \mathbf{0} &= B\mathbf{v}' - \lambda'\mathbf{v}' = (B - \lambda' \mathrm{I}_2)\mathbf{v}' \\
    \begin{bmatrix} 0 \\ 0 \end{bmatrix} &=  \begin{bmatrix} 3 - \lambda' & -4 \\ 1 & -1 - \lambda' \end{bmatrix} \begin{bmatrix}
       v_1'\\ v_2'
    \end{bmatrix} =  \begin{bmatrix}
       (3 - \lambda')v_1' - 4v_2' \\ v_1' - (1+\lambda')v_2'
    \end{bmatrix}
\end{align*}
即：
$$
\left\{
    \begin{aligned}
        (3 - \lambda')v_1' = v_2' \\
        v_1'= (1+\lambda')v_2' 
    \end{aligned}
\right.
$$
$\mathbf{v}'$是非零向量，所以$v_1', v_2'$不全为零。不妨设$v_2'\neq 0$，就得到
\begin{align*}
    v_2' &= (3 - \lambda')v_1' &= (3 - \lambda')(1+\lambda')v_2', \\
    1 - (3 - \lambda')(1+\lambda') &= 0, \\
    \lambda'^2 - 2\lambda' + 1 &= 0.
\end{align*}
解得$\lambda' = 1$。然而我们已经知道，以$1$为本征值的本征空间为$\mathbb{R}\mathbf{v}$，维数为$1$。

这里出现了一个关键的区别：$A$有两个不同的本征值，对应两个1维本征空间，
它们的直和恰好是整个$\mathbb{R}^2$。
而$B$虽然也有本征值，但本征空间只有1维，无法铺满2维的$\mathbb{R}^2$。
也就是说，$B$没有足够多的直无关的本征向量来构成空间的基底。
这个差别，将决定两个自同态能否对角化。

后面我们会学到系统的方法来寻找自同态的所有本征值。
不过现在，通过直接计算$T\mathbf{u}$并观察它是否是$\mathbf{u}$的标量倍，
我们已经可以验证本征值和本征向量了。

\subsection{对角矩阵和对角化}
1维不变子空间性质良好，而如果能够把这种良好性质“推广”到整个空间，我们就可以轻松了解整个自同态的结构。

假设存在非零向量$\mathbf{e}_1, \mathbf{e}_2, \cdots, \mathbf{e}_n \in \mathbb{V}$以及标量$\lambda_1, \lambda_2, \cdots, \lambda_n$，使得：
$$ T(\mathbf{e}_1) = \lambda_1 \mathbf{e}_1, T(\mathbf{e}_2) = \lambda_2 \mathbf{e}_2, \cdots , T(\mathbf{e}_n) = \lambda_n \mathbf{e}_n$$
并且向量组$\mathcal{B} = (\mathbf{e}_1, \mathbf{e}_2, \cdots, \mathbf{e}_n)$直无关，从而构成空间$\mathbb{V}$的基。
那么在基底$\mathcal{B}$上，任意向量经过$T$映射的像为：
$$ \forall \mathbf{u} = a_1 \mathbf{e}_1 + \cdots + a_n \mathbf{e}_n, \,\,\, T(\mathbf{u}) = a_1 \lambda_1 \mathbf{e}_1 + \cdots + a_n \lambda_n \mathbf{e}_n$$
如果我们写出$T$在这个基底上的矩阵：
$$ [T]_{\mathcal{B}} = \begin{bmatrix} \lambda_1 & 0 & \cdots & 0 \\ 0 & \lambda_2 & \cdots & 0 \\ \vdots & \vdots & \ddots & \vdots \\ 0 & 0 & \cdots & \lambda_n \end{bmatrix}$$
这个矩阵只在对角线上有非零元素，其他位置上都是零，只比恒等映射、缩放映射的矩阵再复杂一点。这样的矩阵叫做\textbf{对角矩阵}。
可以表达成对角矩阵的的自同态称为\textbf{可对角化}的自同态。

这样的形式简单明了，我们能够一眼理解映射$T$的性质。比如，我们可以轻易计算复合函数$T^2 = T\circ T$：
$$  \forall \mathbf{u} = a_1 \mathbf{e}_1 + \cdots + a_n \mathbf{e}_n, \,\,\, T(T(\mathbf{u})) = a_1 \lambda_1^2 \mathbf{e}_1 + \cdots + a_n \lambda_n^2 \mathbf{e}_n$$
而一般来说，自同态乘幂$T^k$的矩阵是：
$$ [T^k]_{\mathcal{B}} = \begin{bmatrix} \lambda_1^k & 0 & \cdots & 0 \\ 0 & \lambda_2^k & \cdots & 0 \\ \vdots & \vdots & \ddots & \vdots \\ 0 & 0 & \cdots & \lambda_n^k \end{bmatrix}$$

反之，如果自同态在某个基底$\mathcal{B} = (\mathbf{e}_1, \mathbf{e}_2, \cdots, \mathbf{e}_n)$上的矩阵是对角矩阵，
$$ [T]_{\mathcal{B}} = \begin{bmatrix} \lambda_1 & 0 & \cdots & 0 \\ 0 & \lambda_2 & \cdots & 0 \\ \vdots & \vdots & \ddots & \vdots \\ 0 & 0 & \cdots & \lambda_n \end{bmatrix}$$
那么
$$ T(\mathbf{e}_1) = \lambda_1 \mathbf{e}_1, T(\mathbf{e}_2) = \lambda_2 \mathbf{e}_2, \cdots , T(\mathbf{e}_n) = \lambda_n \mathbf{e}_n$$
这说明每个基向量都是自同态的本征向量，对应的本征值就是对角线上的对应元素。
换句话说，自同态可对角化，等价于说自同态有一组本征向量能构成空间的基底。

那么，什么样的一组本征向量可以构成空间的基底呢？前面我们已经证明，两个本征向量对应的本征值不同，
那么它们直无关。现在我们把这个性质推广到多个向量：
\begin{pp}\label{pp:6_2}
    如果自同态$T$的一组本征向量$\mathbf{u}_1, \mathbf{u}_2, \cdots, \mathbf{u}_k$对应的本征值$\lambda_1, \lambda_2, \cdots, \lambda_k$两两不等，
    那么这组向量直无关。
\end{pp}
\begin{proof} % [{\bfseries 证明\ref{pp:6_2}}]
    设有系数$a_1, \cdots, a_k$使得：
    \begin{align*}
         a_1 \mathbf{u}_1 + a_2 \mathbf{u}_2 + \cdots + a_k \mathbf{u}_k = \mathbf{0}\tag{1}
    \end{align*}
    经过$T$映射后，有：
    \begin{align*} 
        a_1 \lambda_1\mathbf{u}_1 + a_2 \lambda_2 \mathbf{u}_2 + \cdots + a_k \lambda_k \mathbf{u}_k = \mathbf{0}\tag{2}
    \end{align*}
    (1) 式乘以$\lambda_k$之后减去 (2) 式，可以消去式子左边最后一项：
    $$ a_1 (\lambda_k - \lambda_1)\mathbf{u}_1 + a_2 (\lambda_k - \lambda_2) \mathbf{u}_2 + \cdots + a_{k-1} (\lambda_k - \lambda_{k-1}) \mathbf{u}_{k-1} = \mathbf{0}.$$
    式子左边只剩下$k-1$项。而由于本征值两两不等，以上系数中原来不为零的$a_i$，乘以$\lambda_k - \lambda_i$之后仍不为零。\\
    重复以上操作，每次操作式子左边项数减少1，而非零系数仍然非零。操作$k-1$次后，式子左边只剩下一项，
    其系数为$a_1$乘以某个非零的标量\footnote{这个量可以用连乘表示出来。}。而式子右边为零向量。所以$a_1 = 0$。\\
    将这个结果代入倒数第二次操作得到的式子。式子左边有两项，而第一项由于$a_1=0$成为零向量，所以类似可得到$a_2 = 0$。\\
    依此类推，可以得出$a_1 = \cdots = a_k = 0.$ 这说明这组向量直无关。 
    
\end{proof}

所以，如果我们找到了$n$个两两不同的本征值，把它们对应的本征向量组成基底，自同态$T$在这组基底上就是对角矩阵。

我们用一个完整的例子来展示对角化的过程。
仍然考虑自同态$A$，其矩阵为：
$$
A = \begin{bmatrix} 4 & 3 \\ -2 & -1 \end{bmatrix}.
$$
前面我们已经验证，$A$有两个本征值$\lambda_1 = 1$和$\lambda_2 = 2$，
对应的本征向量分别是$\mathbf{e}_1 = (1,-1)$和$\mathbf{e}_2 = (3,-2)$。
它们对应不同的本征值，所以直无关，构成$\mathbb{R}^2$的一组基$\mathcal{B} = (\mathbf{e}_1, \mathbf{e}_2)$。

在这组基上，$A$的矩阵为：
$$
[A]_{\mathcal{B}} = \begin{bmatrix} 1 & 0 \\ 0 & 2 \end{bmatrix}.
$$
这是因为$A\mathbf{e}_1 = 1\cdot\mathbf{e}_1 + 0\cdot\mathbf{e}_2$，$A\mathbf{e}_2 = 0\cdot\mathbf{e}_1 + 2\cdot\mathbf{e}_2$。
$A$在基底$\mathcal{B}$上是对角矩阵，所以$A$可对角化。

对角化带来什么好处呢？比如，我们要计算$A$的乘幂$A^k$。
在基底$\mathcal{B}$上，$A^k$的矩阵就是：
$$
[A^k]_{\mathcal{B}} = \begin{bmatrix} 1^k & 0 \\ 0 & 2^k \end{bmatrix} = \begin{bmatrix} 1 & 0 \\ 0 & 2^k \end{bmatrix}.
$$
也就是说，$A^k\mathbf{e}_1 = \mathbf{e}_1$，$A^k\mathbf{e}_2 = 2^k\mathbf{e}_2$。
对任意向量$\mathbf{u} = a\mathbf{e}_1 + b\mathbf{e}_2$，有：
$$
A^k\mathbf{u} = a\mathbf{e}_1 + 2^k b\mathbf{e}_2.
$$
乘幂的计算变得十分简单。

反过来，我们来看$B$为什么不能对角化。
$B$的唯一本征值是$\lambda = 1$，本征空间$\Ker(B - \mathrm{I}) = \mathbb{R} (2,1)$只有$1$维。
因此，不存在由本征向量构成的基底。所以$B$不能对角化。

最后看一个$3$维的例子。考虑$\mathbb{R}^3$上的自同态$C$，其矩阵为：
$$
C = \begin{bmatrix} 5 & 0 & 1 \\ 1 & 3 & 2 \\ 0 & 0 & 3 \end{bmatrix}.
$$
我们可以直接验证：
$$
C \begin{bmatrix} 2 \\ 1 \\ 0 \end{bmatrix} = \begin{bmatrix} 10 \\ 5 \\ 0 \end{bmatrix} = 5 \cdot \begin{bmatrix} 2 \\ 1 \\ 0 \end{bmatrix},
\qquad
C \begin{bmatrix} 0 \\ 1 \\ 0 \end{bmatrix} = \begin{bmatrix} 0 \\ 3 \\ 0 \end{bmatrix} = 3 \cdot \begin{bmatrix} 0 \\ 1 \\ 0 \end{bmatrix}.
$$
所以$(2,1,0)$是$C$关于本征值$5$的本征向量，$(0,1,0)$是$C$关于本征值$3$的本征向量。

反之，考虑$C\mathbf{v} = \lambda \mathbf{v}$。如前展开得到：
\begin{align*}
    \mathbf{0} &= C\mathbf{v} - \lambda\mathbf{v} = (C - \lambda \mathrm{I}_3)\mathbf{v} \\
    \begin{bmatrix} 0 \\ 0 \\ 0 \end{bmatrix} &=  \begin{bmatrix} 5- \lambda & 0 & 1 \\ 1 & 3- \lambda & 2 \\ 0 & 0 & 3- \lambda \end{bmatrix} \begin{bmatrix}
       v_1\\ v_2 \\ v_3
    \end{bmatrix} =  \begin{bmatrix}
       (5 - \lambda)v_1 +v_3 \\ v_1 + (3 - \lambda)v_2 + 2v_3 \\ (3 - \lambda)v_3
    \end{bmatrix}
\end{align*}
即：
$$
\left\{
    \begin{aligned}
        (5 - \lambda)v_1 + v_3 = 0 \\
        v_1 + (3 - \lambda)v_2 + 2v_3 = 0 \\
        (3 - \lambda)v_3 = 0
    \end{aligned}
\right.
$$
首先考虑$\lambda = 3$的情况。这时第三个方程变为$0 = 0$。前两个方程变为
$$
\left\{
    \begin{aligned}
        2v_1 + v_3 = 0 \\
        v_1  + 2v_3 = 0 
    \end{aligned}
\right.
$$
解得$v_1 = v_3 = 0$。于是$\mathbf{v} \in \mathbb{R}(0,1,0)$。这说明本征值$3$对应的本征空间就是一维空间$\mathbb{R}(0,1,0)$。

再考虑$\lambda\neq 3$的情况。这时第三个方程解出$v_3 = 0$。代入前两个方程得到
$$
\left\{
    \begin{aligned}
        (5 - \lambda)v_1 = 0 \\
        v_1 + (3 - \lambda)v_2 = 0 
    \end{aligned}
\right.
$$
第一个方程解得$\lambda = 5$或$v_1 = 0$。如果$\lambda = 5$，那么解得$v_1 -2 v_2=0 $，得到本征值$5$对应的一维本征空间$\mathbb{R}(2,1,0)$。如果$v_1 = 0$，那么解得$v_2 = 0$，$\mathbf{v}$为零向量，排除。

综上，我们知道$C$只有两个本征值$3,5$分别对应$1$维的本征空间$\mathbb{R}(2,1,0)$和$\mathbb{R}(0,1,0)$.和$B$一样，它的本征向量无法构成$\mathbb{R}^3$的基底，于是无法对角化。

\section{自同态与整式}

\subsection{极小归零式}

上一节我们讨论了不变子空间是$1$维空间的情形。$1$维不变子空间总对应一个本征值，其中的非零向量都是本征向量。
现在我们来研究维数大于$1$的不变子空间。为此，我们要从自同态的整式开始着手。

设$T\in\mathcal{L}(\mathbb{V})$，考虑$T$与自身的复合得到的自同态、它们之间的加法与数乘运算，我们会发现：
每个以$\mathbb{F}$中元素为系数的一元整式$p(z) = a_0 + a_1 z + \cdots + a_m z^m$，都可以对应一个由$T$衍生出的映射：
$$ p(T) = a_0 \mathrm{I} + a_1 T + \cdots + a_m T^m.$$
这个映射具体为：
$$ \forall \mathbf{u} \in \mathbb{V}, \,\,\, p(T)(\mathbf{u}) = a_0 \mathbf{u} + a_1 T(\mathbf{u}) + \cdots + a_m \overbrace{T(T(\cdots(T}^{m\mbox{\scriptsize 次}}(\mathbf{u})))). $$
我们用恒等映射$\mathrm{I}$对应整式的常数项，或者说，约定$T$的“0次幂”就是恒等映射。 

可以验证，映射$\Psi_T: p \mapsto p(T)$是$\mathbb{F}$上的整式环$\mathbb{F}[z]$到$\mathcal{L}(\mathbb{V})$的直映射，同时也是一个环同态。
也就是说，所有“$T$构成的整式”构成一个域上场，一般记为$\mathbb{F}[T]$。这是一个交换域上场，也称为交换代数。

可以验证，$\mathbb{F}[T]$是$\mathcal{L}(\mathbb{V})$的子空间，所以维数有限。考虑一组映射$\mathrm{I}, T, T^2, \cdots , T^k$，
它们都是空间$\mathcal{L}(\mathbb{V})$里的向量，而当$k\geqslant n^2$的时候，
$\mathrm{I}, T, T^2, \cdots , T^k$直相关，所以必然有不全为$0$的系数$a_0, a_1, \cdots a_m$使得
$$a_0 \mathrm{I} + a_1 T + \cdots + a_m T^m = 0.$$
也就是说，存在整式$p(z) = c_0 + c_1 z + \cdots + c_{k} z^k$使得$\Psi_T(p) = 0$。这里右侧的0指的是$\mathcal{L}(\mathbb{V})$中的零映射。这表明，$\Psi_T$并非单射。也就是说，作为直映射，$\Psi_T$的零空间维数大于0。

我们将$\Psi_T$的零空间记为$\mathfrak{G}(T)$，称之为$T$的\textbf{归零空间}。$\mathfrak{G}(T)$里的元素称为$T$的\textbf{归零式}。

容易验证，归零空间有以下性质：
\begin{enumerate}
    \item 如果$p, q \in \mathfrak{G}(T)$，那么$p + q \in \mathfrak{G}(T).$
    \item 如果$\lambda \in \mathbb{F}, \,\, p \in \mathfrak{G}(T)$，那么$\lambda p \in \mathfrak{G}(T).$
    \item 如果$p \in \mathfrak{G}(T), q \in \mathbb{F}[z]$，那么$pq \in \mathfrak{G}(T).$
\end{enumerate}
这说明$\mathfrak{G}(T)$是整式环$\mathbb{F}[z]$的藏\footnote{见定义\ref{df:0_3_4}.}。$\mathbb{F}[z]$是单藏环，所以$\mathfrak{G}(T)$是单藏\footnote{关于单藏环的定义与性质，参见命题\ref{pp:0_3_5}。}，即由某个首一整式生成：$\mathfrak{G}(T) = \langle m_T \rangle$。在首一的要求下，这个整式是唯一的。我们称其为$T$的\textbf{极小归零式}，简称\textbf{极小式}，$\mathfrak{G}(T)$中所有元素的都是它的倍式。

作为例子，我们来看对角矩阵对应的直映射的极小归零式。
设对角矩阵$T$的对角线上元素为$\lambda_1, \cdots , \lambda_n$，
它们就是对应直映射的本征值。实际上，将$T$表达为对角矩阵的基底，就是由对应的本征向量组成的。
设这个基为$\mathcal{B} = (\mathbf{e}_1, \cdots, \mathbf{e}_n)$，那么：
$$ \forall 1 \leqslant i \leqslant n, T(\mathbf{e}_i) = \lambda_i \mathbf{e}_i. $$

将对角线元素中重复的值去掉，得到本征值集合：
$\lambda_{(1)}, \cdots , \lambda_{(k)}$。考虑整式：
$$ \pi(z) = \prod_{i=1}^k (z - \lambda_{(i)}).$$
我们可以证明$\pi(T)$是零映射：
$$ \forall \mathbf{u} \in \mathbb{V}, \,\,\, \mathbf{u} = \sum_{i=1}^n u_i \mathbf{e}_i, $$
因此
\begin{align*}
     \pi(T)(\mathbf{u}) &= \sum u_i \pi(T) (\mathbf{e}_i)  \\
     &= \sum_{i=1}^n u_i Q_i(T) (T(\mathbf{e}_i) - \lambda_i \mathbf{e}_i)  \\
     &= \sum_{i=1}^n u_i Q_i(T) (\mathbf{0}) = \mathbf{0}.  
\end{align*}
其中$Q_i$是本征值集合中不为$\lambda_i$的本征值对应的$z - \lambda_{(j)}$的乘积。

如果$T$的本征值两两不等，那么极小归零式是$n$次整式。
如果其中有相等的，那么极小归零式的次数小于$n$。

我们把$m_T$的次数记为$d_T$。对任意整式$p\in\mathbb{F}[z]$，
它对$m_T$的带余除法得到：
$$p = m_T q + r,$$
则$p(T) = q(T)m_T(T) + r(T) = r(T)$。也就是说，所有$T$构成的整式的值都等于次数小于$d_T$的整式，于是可以表示为$\mathcal{B} = (\mathrm{I}, T, \cdots , T^{d_T - 1})$的直组合，这说明$\mathrm{I}, T, \cdots , T^{d_T - 1}$生成$\mathbb{F}[T]$。

我们来看上一节中三个自同态的极小归零式。
首先考虑自同态$A$，其矩阵为：
$$
A = \begin{bmatrix} 4 & 3 \\ -2 & -1 \end{bmatrix}.
$$
第一节中我们已经验证，$A$有本征值$1$和$2$。考虑整式$p(z) = (z-1)(z-2) = z^2 - 3z + 2$。
我们验证$p(A)$是否为零映射。首先计算$A^2$：
$$
A^2 = \begin{bmatrix} 4 & 3 \\ -2 & -1 \end{bmatrix}\begin{bmatrix} 4 & 3 \\ -2 & -1 \end{bmatrix} = \begin{bmatrix} 16-6 & 12-3 \\ -8+2 & -6+1 \end{bmatrix} = \begin{bmatrix} 10 & 9 \\ -6 & -5 \end{bmatrix}.
$$
于是：
$$
p(A) = A^2 - 3A + 2\mathrm{I} = \begin{bmatrix} 10 & 9 \\ -6 & -5 \end{bmatrix} - \begin{bmatrix} 12 & 9 \\ -6 & -3 \end{bmatrix} + \begin{bmatrix} 2 & 0 \\ 0 & 2 \end{bmatrix} = \begin{bmatrix} 0 & 0 \\ 0 & 0 \end{bmatrix}.
$$
所以$p(z) = z^2 - 3z + 2$是$A$的归零式。

它是否是极小归零式呢？$p$的次数为$2$。一次归零式意味着$A = \lambda\mathrm{I}$，即$A$是缩放映射。
但$A$显然不是缩放映射（$A\mathbf{e}_1 = \mathbf{e}_1$而$A\mathbf{e}_2 = 2\mathbf{e}_2$，两个本征值不等），
所以一次归零式不存在。因此$A$的极小归零式为：
$$
m_A(z) = (z-1)(z-2) = z^2 - 3z + 2.
$$

再考虑自同态$B$，其矩阵为：
$$
B = \begin{bmatrix} 3 & -4 \\ 1 & -1 \end{bmatrix}.
$$
第一节中我们验证了$B$的唯一本征值为$1$，本征空间为$\mathbb{R}(2,1)$。
考虑整式$q(z) = (z-1)^2 = z^2 - 2z + 1$。首先计算$B^2$：
$$
B^2 = \begin{bmatrix} 3 & -4 \\ 1 & -1 \end{bmatrix}\begin{bmatrix} 3 & -4 \\ 1 & -1 \end{bmatrix} = \begin{bmatrix} 9-4 & -12+4 \\ 3-1 & -4+1 \end{bmatrix} = \begin{bmatrix} 5 & -8 \\ 2 & -3 \end{bmatrix}.
$$
于是：
$$
q(B) = B^2 - 2B + \mathrm{I} = \begin{bmatrix} 5 & -8 \\ 2 & -3 \end{bmatrix} - \begin{bmatrix} 6 & -8 \\ 2 & -2 \end{bmatrix} + \begin{bmatrix} 1 & 0 \\ 0 & 1 \end{bmatrix} = \begin{bmatrix} 0 & 0 \\ 0 & 0 \end{bmatrix}.
$$
所以$q(z) = (z-1)^2$是$B$的归零式。

显然$B$也没有一次归零式。因此$B$的极小归零式为：
$$
m_B(z) = (z-1)^2 = z^2 - 2z + 1.
$$

最后考虑自同态$C$，其矩阵为：
$$
C = \begin{bmatrix} 5 & 0 & 1 \\ 1 & 3 & 2 \\ 0 & 0 & 3 \end{bmatrix}.
$$
第一节中我们验证了$C$有本征值$5$和$3$，对应的本征空间都只有$1$维。
首先验证$(z-3)(z-5)$是否是归零式。计算：
$$
C - 3\mathrm{I} = \begin{bmatrix} 2 & 0 & 1 \\ 1 & 0 & 2 \\ 0 & 0 & 0 \end{bmatrix}, \quad C - 5\mathrm{I} = \begin{bmatrix} 0 & 0 & 1 \\ 1 & -2 & 2 \\ 0 & 0 & -2 \end{bmatrix}.
$$
它们的乘积为：
$$
(C-3\mathrm{I})(C-5\mathrm{I}) = \begin{bmatrix} 0 & 0 & 0 \\ 0 & 0 & -3 \\ 0 & 0 & 0 \end{bmatrix} \neq \begin{bmatrix} 0 & 0 & 0 \\ 0 & 0 & 0 \\ 0 & 0 & 0 \end{bmatrix}.
$$
所以$(z-3)(z-5)$不是$C$的归零式。

再试$(z-3)^2(z-5)$。先计算$(C-3\mathrm{I})^2$：
$$
(C-3\mathrm{I})^2 = \begin{bmatrix} 2 & 0 & 1 \\ 1 & 0 & 2 \\ 0 & 0 & 0 \end{bmatrix}^2 = \begin{bmatrix} 4 & 0 & 2 \\ 2 & 0 & 1 \\ 0 & 0 & 0 \end{bmatrix}.
$$
于是：
$$
(C-3\mathrm{I})^2(C-5\mathrm{I}) = \begin{bmatrix} 4 & 0 & 2 \\ 2 & 0 & 1 \\ 0 & 0 & 0 \end{bmatrix}\begin{bmatrix} 0 & 0 & 1 \\ 1 & -2 & 2 \\ 0 & 0 & -2 \end{bmatrix} = \begin{bmatrix} 0 & 0 & 0 \\ 0 & 0 & 0 \\ 0 & 0 & 0 \end{bmatrix}.
$$
所以$(z-3)^2(z-5)$是$C$的归零式。

能否找到更低次的归零式呢？
一次归零式要求$C = \lambda\mathrm{I}$，显然不成立。二次归零式呢？我们假设有首一的二次多项式$p$使得$p(C) $为零矩阵，那么，用$p$除$(z-3)^2(z-5)$，余式$r$至多是一次式。按定义有：
$$ 0 = (C-3\mathrm{I})^2(C-5\mathrm{I}) + p(C)q(C) + r(C) = r(C).$$
于是$r(C) = 0$。这说明$r$只能是$0$。于是$p$是$(z-3)^2(z-5)$的因式。也就是说，$p$只能是$(z-3)^2$或$(z-3)(z-5)$。
$(z-3)(z-5)$已经排除。而$(z-3)^2$也不成立，因为只要把本征值$5$对应的任何本征向量$\mathbf{u}$代入就得到：
$$ 
(C - 3\mathrm{I})^2\mathbf{u} = (C - 3\mathrm{I})(C\mathbf{u} - 3\mathbf{u}) = (C - 3\mathrm{I})(2\mathbf{u}) = 2(C\mathbf{u} - 3\mathbf{u}) = 4\mathbf{u} \neq \mathbf{0}.
$$
于是$C$的极小归零式为：
$$
m_C(z) = (z-3)^2(z-5).
$$

将三个极小归零式对比：
\begin{center}
\begin{tabular}{|c|c|c|c|}
\hline
 & $A$ & $B$ & $C$ \\
\hline
极小归零式 & $(z-1)(z-2)$ & \phantom{.2)}$(z-1)^2$\phantom{.2)} & $(z-3)^2(z-5)$ \\
次数 & $2$ & $2$ & $3$ \\
重根 & 无 & 有 & 有 \\
\hline
\end{tabular}
\end{center}

回忆第一节中的结论：$A$可对角化，$B$和$C$不可对角化。
这里我们看到，可对角化的$A$的极小归零式没有重根，
而不可对角化的$B$和$C$的极小归零式都有重根。这并非巧合。

另一方面，如果$\mathrm{I}, T, \cdots , T^{d_T - 1}$的直组合等于零映射，那么其系数就对应了一个次数小于$d_T$的归零式。根据$m_T$的定义，该整式为零，因此所有系数为零。于是$\mathrm{I}, T , \cdots , T^{d_T - 1}$直无关。

综上可得，$\mathbb{F}[T]$是维数为$d_T$的平直空间。

\subsection{整式和零空间}
对一般的自同态来说，本征值和极小归零式有什么联系呢？
\begin{pp}\label{pp:6_3}
    设$T$为空间$\mathbb{V}$的自同态，$m_T$是它的极小归零式，那么，
    $\lambda$为$T$的本征值当且仅当它是$m_T$的根。
\end{pp}
\begin{proof} % [{\bfseries 证明\ref{pp:6_3}}]
    分两个方向证明。首先，如果$\lambda$是$T$的本征值，那么存在非零向量$\mathbf{v}$
    使得
    $$ T(\mathbf{v}) = \lambda \mathbf{v}.$$
    则对任意正整数$k$，都有$T^k(\mathbf{v}) = \lambda^k \mathbf{v}$。
    进而可以推出，对任意$p\in\mathbb{F}[z]$，都有$p(T)(\mathbf{v}) = p(\lambda) \mathbf{v}.$\\
    极小归零式是归零式，所以$m_T(T)(\mathbf{v}) = \mathbf{0}$，这说明$m_T(\lambda) = 0$，$\lambda$是$m_T$的根。\\
    反过来，如果$\lambda$是极小归零式的根，那么存在整式$Q$使得：
    $$ m_T(z) = (z - \lambda)Q(z).$$
    整式$Q$次数比$m_T$低，所以存在向量$\mathbf{u}$使得$Q(T)(\mathbf{u}) \neq \mathbf{0}$。\\
    而$m_T(T)(\mathbf{u}) = \mathbf{0}$，这说明：
    $$ (T - \lambda \mathrm{I})(Q(T)(\mathbf{u})) = \mathbf{0}.$$
    设$\mathbf{v} = Q(T)(\mathbf{u})$，则：
    $$ T(\mathbf{v}) = \lambda \mathbf{v}.$$
    这就说明，$\lambda$是$T$的本征值。

\end{proof}

我们用第一节中的三个自同态来验证。对$A$来说，$m_A(z) = (z-1)(z-2)$，根为$1$和$2$。
第一节中我们已经验证，$1$和$2$确实是$A$的本征值。

对$B$来说，$m_B(z) = (z-1)^2$，唯一的根为$1$。第一节中我们已经验证，$1$确实是$B$的唯一本征值。

对$C$来说，$m_C(z) = (z-3)^2(z-5)$，根为$3$和$5$。第一节中我们已经验证，$3$和$5$确实是$C$的本征值。

反过来，我们也可以利用这个命题来排除不是本征值的数。
比如，$0$不是$m_A$的根（$m_A(0) = 2 \neq 0$），所以$0$不是$A$的本征值。
同理，$4$不是$m_B$的根（$m_B(4) = 9 \neq 0$），所以$4$不是$B$的本征值。

自同态的整式可以作为研究不变子空间的切入点。首先，对任意非零整式$P$，
考虑$P(T)$的零空间$\Ker P(T)$。$\forall \,\, \mathbf{u} \,\, \in \Ker P(T)$，有：
$$ P(T)(T(\mathbf{u})) = T(P(T)(\mathbf{u})) = T(\mathbf{0}) = \mathbf{0}.$$
所以$T(\mathbf{u})\,\,\in\,\,\Ker P(T)$。也就是说，关于$T$的整式的零空间，
总是$T$作用下的不变子空间。

进一步，我们可以验证，如果整式$Q$是$P$的因式，$P = Q\cdot R$，那么
$$ \forall \,\, \mathbf{u} \,\,\in \,\, \Ker Q(T),\,\,\, P(T)(\mathbf{u}) = R(T)(Q(T)(\mathbf{u})) = \mathbf{0}.$$
所以，$\Ker Q(T)$是$\Ker P(T)$的$T$-不变子空间。
如果$Q_1, Q_2, \cdots ,Q_k$是一系列整式，前者总是后者的因式，那么，它们对应的$T$-整式映射\footnote{即整式在$\Psi_T$下的像。}的零空间也都是
$T$-不变子空间，而且前者是后者的$T$-不变子空间。

反过来，设$\mathbb{U}$是$\mathbb{W}$的$T$-不变子空间，
则$T$限制在它们中的部分：$\restr{T}{\mathbb{U}}$、$\restr{T}{\mathbb{W}}$各自有一个极小归零式$m_{T,\mathbb{U}}$、$m_{T,\mathbb{W}}$。
考虑$\mathbf{u} \in \mathbb{U}$，
按照定义，$\restr{T}{\mathbb{U}}$也是$\restr{T}{\mathbb{W}}$限制在子空间$\mathbb{U}$中的部分，所以
$$m_{T,\mathbb{W}}(\restr{T}{\mathbb{U}})(\mathbf{u}) = m_{T,\mathbb{W}}(\restr{T}{\mathbb{W}})(\mathbf{u}) = \mathbf{0}.$$
所以$m_{T,\mathbb{W}}$是$\restr{T}{\mathbb{U}}$的归零式，
从而$m_{T,\mathbb{U}}$是$m_{T,\mathbb{W}}$的因式。

如果$\mathbb{V}$能够写成若干个$T$-不变子空间的直和，
那么$T$限制在每个子空间中的映射对应的极小归零式，都是$m_T$的因式。也就是说，$m_T$是它们的公倍式。
按照$m_T$的定义，$m_T$是它们的最小公倍式。

\begin{pp}\textbf{零空间引理}\label{pp:6_4}
    设$T$为空间$\mathbb{V}$的自同态，整式$P\in\mathbb{F}[z]$可分解为若干个两两互素的整式的乘积：
    $$ P = P_1 P_2\cdots P_k$$
    那么
    $$ \Ker P(T) = \Ker P_1(T) \oplus \Ker P_2(T) \oplus \cdots \oplus \Ker P_k(T). $$
\end{pp}
\begin{proof} % [{\bfseries 证明\ref{pp:6_4}}]
    对整式的个数$k$使用归纳法。$k=1$情况显然。考虑$k=2$的情况。\\
    $P_1$和$P_2$互素，所以根据倍和析因定理，存在整式$Q_1, Q_2$使得：$Q_1 P_1 + Q_2 P_2 = 1.$\\
    对$\Ker P(T)$中元素$\mathbf{u}$，设$\mathbf{u}_1 = Q_1(T)(P_1(T)(\mathbf{u})), \,\,\, \mathbf{u}_2 = Q_2(T)(P_2(T)(\mathbf{u}))$，
    $P = P_1 P_2$，所以$P_1(T)(\mathbf{u}_2) = \mathbf{0}, \,\,\, P_2(T)(\mathbf{u}_1) = \mathbf{0}$。
    $\mathbf{u}_1 \in \Ker P_2(T), \,\,\, \mathbf{u}_2 \in \Ker P_1(T).$
    另一方面，$Q_1(T) \circ P_1(T) + Q_2(T) \circ P_2(T) = \mathrm{I}$，所以：
    $$ \mathbf{u} = Q_1(T)(P_1(T)(\mathbf{u})) + Q_2(T)(P_2(T)(\mathbf{u})) = \mathbf{u}_1 + \mathbf{u}_2.$$
    这说明$\Ker P(T) = \Ker P_1(T) + \Ker P_2(T).$
    而$\forall \,\, \mathbf{u} \,\, \in \Ker P_1(T) \cap \Ker P_2(T)$，
    $$ \mathbf{u} = Q_1(T)(P_1(T)(\mathbf{u})) + Q_2(T)(P_2(T)(\mathbf{u})) = \mathbf{0}. $$
    所以$\Ker P_1(T) \cap \Ker P_2(T) = \{\mathbf{0}\}$，$\Ker P(T) = \Ker P_1(T) \oplus \Ker P_2(T).$\\
    假设命题对$k$成立，对于$k+1$来说，$P = P_1 Q$，其中$Q = P_2 \cdots P_{k+1}$和$P_1$互素。根据$k=2$的情况， 
    $$ \Ker P(T) = \Ker P_1(T) \oplus \Ker Q(T)$$
    而根据归纳假设，
    $$ \Ker Q(T) = \Ker P_2(T) \oplus \cdots \oplus \Ker P_{k+1}.$$
    所以
    $$ \Ker P(T) = \Ker P_1(T) \oplus \Ker P_2(T) \oplus \cdots \oplus \Ker P_{k+1}.$$
    这说明命题对$k+1$也成立。根据归纳原理，命题对任意正整数$k$成立。
    
\end{proof}

我们用自同态$C$来具体感受一下零空间引理。
$C$的极小归零式为$m_C(z) = (z-3)^2(z-5)$。
其中$(z-3)^2$和$(z-5)$互素。按照零空间引理：
$$
\mathbb{R}^3 = \Ker (C-3\mathrm{I})^2 \oplus \Ker (C-5\mathrm{I}).
$$
我们来验证这一点。

首先计算$\Ker(C-5\mathrm{I})$。第一节中我们已经知道，
$C-5\mathrm{I}$的零空间由$(2,1,0)$生成，维数为$1$。

再计算$\Ker(C-3\mathrm{I})^2$。前面已经算出：
$$
(C-3\mathrm{I})^2 = \begin{bmatrix} 4 & 0 & 2 \\ 2 & 0 & 1 \\ 0 & 0 & 0 \end{bmatrix}.
$$
解方程$(C-3\mathrm{I})^2\mathbf{u} = \mathbf{0}$：
$$
\begin{bmatrix} 4 & 0 & 2 \\ 2 & 0 & 1 \\ 0 & 0 & 0 \end{bmatrix}\begin{bmatrix} u_1 \\ u_2 \\ u_3 \end{bmatrix} = \begin{bmatrix} 0 \\ 0 \\ 0 \end{bmatrix}.
$$
两个方程分别是$4u_1 + 2u_3 = 0$和$2u_1 + u_3 = 0$，它们等价。
解得$u_3 = -2u_1$，$u_2$自由。所以$\Ker(C-3\mathrm{I})^2$的维数为$2$，
一组基为$(1,0,-2)$和$(0,1,0)$。

验证直和：$\Ker(C-3\mathrm{I})^2$维数为$2$，$\Ker(C-5\mathrm{I})$维数为$1$，
两者维数之和为$3 = \dim \mathbb{R}^3$。
只需验证交集为$\{\mathbf{0}\}$。设$\mathbf{u} \in \Ker(C-3\mathrm{I})^2 \cap \Ker(C-5\mathrm{I})$，
则$\mathbf{u} = a(2,1,0)$且$4a + 2\cdot 0 = 0$（代入$(C-3\mathrm{I})^2$的第一个方程），解得$a = 0$。
所以交集确实为$\{\mathbf{0}\}$。

于是$\mathbb{R}^3 = \Ker(C-3\mathrm{I})^2 \oplus \Ker(C-5\mathrm{I})$，
两个子空间都是$C$-不变子空间。
在$\Ker(C-5\mathrm{I})$中，$C$的作用是乘以$5$；
在$\Ker(C-3\mathrm{I})^2$中，$C$的作用更复杂，不是简单的缩放。
取$\Ker(C-3\mathrm{I})^2$的基$(0,1,0)$和$(1,0,-2)$，计算：
$$
C\begin{bmatrix} 0 \\ 1 \\ 0 \end{bmatrix} = \begin{bmatrix} 0 \\ 3 \\ 0 \end{bmatrix} = 3\begin{bmatrix} 0 \\ 1 \\ 0 \end{bmatrix},
\quad
C\begin{bmatrix} 1 \\ 0 \\ -2 \end{bmatrix} = \begin{bmatrix} 3 \\ -3 \\ -6 \end{bmatrix}.
$$
验证$(3,-3,-6)$是否在$\Ker(C-3\mathrm{I})^2$中：
$4\cdot 3 + 2\cdot(-6) = 0$，确实在其中。
将它用基$(0,1,0)$和$(1,0,-2)$表示：
$(3,-3,-6) = -3(0,1,0) + 3(1,0,-2)$。
所以$C$限制在$\Ker(C-3\mathrm{I})^2$上，在基$((0,1,0),(1,0,-2))$上的矩阵为：
$$
\begin{bmatrix} 3 & -3 \\ 0 & 3 \end{bmatrix}.
$$
这不是对角矩阵。$C$在$\Ker(C-3\mathrm{I})^2$上的限制不能对角化。
这和$B$的情形类似：本征值$3$的极小归零式含有重根$(z-3)^2$，
对应的不变子空间内部无法进一步分解为$1$维不变子空间的直和。

综合以上，如果我们取$\mathbb{R}^3$的基底为$((0,1,0),\, (1,0,-2),\, (2,1,0))$，
则$C$在这组基上的矩阵为：
$$
\begin{bmatrix} 3 & -3 & 0 \\ 0 & 3 & 0 \\ 0 & 0 & 5 \end{bmatrix}.
$$
我们可以把这样的矩阵叫做\textbf{分块对角矩阵}\footnote{也称为准对角矩阵。}。左上角的$2\times 2$块对应$\Ker(C-3\mathrm{I})^2$，右下角的$1\times 1$块对应$\Ker(C-5\mathrm{I})$。

根据零空间引理，只要$P\in\mathfrak{G}(T)$，那么我们就可以根据$P$的素因式分解，
把空间$\mathbb{V}$分解为若干个$T$-不变子空间的直和。
$$ \mathbb{V} = \Ker P_1(T) \oplus \Ker P_2(T) \oplus \cdots \oplus \Ker P_k(T). $$
在每个$T$-不变子空间$\Ker P_i(T)$中找一组基：$(\mathbf{e}_{i,1}, \mathbf{e}_{i,2}, \cdots , \mathbf{e}_{i, d_i})$.
那么$k$组子空间的基可以拼成空间$\mathbb{V}$的一组基。
$$ \mathcal{B} = (\mathbf{e}_{1,1}, \cdots , \mathbf{e}_{1, d_1},\mathbf{e}_{2,1}, \cdots , \mathbf{e}_{2, d_2}, \cdots , \mathbf{e}_{k,1}, \cdots , \mathbf{e}_{k, d_k})$$
$T$在这组基上的矩阵是分块对角矩阵：
$$ [T]_{\mathcal{B}} = \begin{bmatrix} A_1 & 0 & \cdots & 0 \\ 0 & A_2 & \cdots & 0 \\ \vdots & \vdots & \ddots & \vdots \\ 0 & 0 & \cdots & A_k\end{bmatrix}$$

对$A$来说，$m_A(z) = (z-1)(z-2)$，两个因式$z-1$和$z-2$互素。零空间引理给出：
$$
\mathbb{R}^2 = \Ker(A-\mathrm{I}) \oplus \Ker(A-2\mathrm{I}).
$$
第一节中我们已经知道，$\Ker(A-\mathrm{I}) = \mathbb{R}(1,-1)$，$\Ker(A-2\mathrm{I}) = \mathbb{R}(3,-2)$。
在每个$1$维不变子空间中，$A$的作用就是乘以对应的本征值。
取基底$((1,-1),\, (3,-2))$，$A$在其上的矩阵为：
$$
\begin{bmatrix} 1 & 0 \\ 0 & 2 \end{bmatrix}.
$$
这正是对角矩阵。

对$B$来说，$m_B(z) = (z-1)^2$只有一个素因式$z-1$。
零空间引理在这里不能将空间进一步分解。
$\Ker(B-\mathrm{I})^2 = \mathbb{R}^2$是整个空间，无法拆分为更小的不变子空间的直和。
这和$B$不可对角化是一致的。

可以看到，极小归零式的因式分解决定了空间能分解为多少个不变子空间的直和，
而每个因式是否有重根，则决定了对应的不变子空间内部能否继续对角化。

\section{自同态的约简}
\subsection{三角矩阵和三角化}
我们研究直自同态的一个主要目的是化简。我们希望将直自同态用简单的形式表示，以了解自同态的本质。
具体来说，有两个思路。一个思路是将所有的对象转化，不断用更简单的形式表达。
另一个思路是先将一部分对象用简单形式表达，然后探讨它们和一般对象的关系。

上一节中，我们证明了，在特定的基底上，所有自同态都能表达为分块对角矩阵。这类约简对应一类不变子空间。
这一节中，我们将研究一类特殊的矩阵，并考察哪些自同态能够表达为这类矩阵。它们对应另一类不变子空间。

\begin{df}\label{df:6_3}
    如果$n\times n$矩阵$A = [a_{i,j}]_{1 \leqslant i , j \leqslant n}$的系数满足：
    $$ \forall \,\, 1 \leqslant i < j \leqslant n, \,\,\, a_{i,j} = 0,$$
    则称其为\textbf{下三角矩阵}。若它的系数满足：
    $$ \forall \,\, 1 \leqslant j < i \leqslant n, \,\,\, a_{i,j} = 0,$$
    则称其为\textbf{上三角矩阵}。两者统称为\textbf{三角矩阵}。
\end{df}

比如，以下这个矩阵是上三角矩阵：
$$ \begin{bmatrix} 2 & 0 & 4 & 9 \\ 0 & 7 & -1 & 3 \\ 0 & 0 & 1 & 6 \\ 0 & 0 & 0 & 1\end{bmatrix}$$
它第2行第1列、第3行第2列等位置（所有行数比列数大的位置）上的系数都是0。而下面这个矩阵是下三角矩阵。
$$ \begin{bmatrix} 3 & 0 & 0 & 0 \\ 0 & 7 & 0 & 0 \\ 5 & 4 & 1 & 0 \\ -4 & 2 & 0 & 1\end{bmatrix}$$
它第1行第3列、第2行第4列等位置（所有行数比列数小的位置）上的系数都是0。

显然，上三角矩阵的转置矩阵是下三角矩阵，反之亦然。
所有的上三角矩阵关于矩阵加法和数乘运算构成平直空间，维数为$\frac12 n(n+1)$。
所有的下三角矩阵也构成一个维数相同的平直空间。
两者有高度对称性。

怎样的自同态可以表达为上三角矩阵呢？我们可以用嵌套的不变子空间来解释上三角矩阵对应的自同态。
\begin{pp}\label{pp:6_5}
    设$T$是空间$\mathbb{V}$的自同态，$\mathcal{B} = (\mathbf{e}_1, \cdots , \mathbf{e}_n)$是空间的一组基。
    以下命题等价：
    \begin{enumerate}
        \item $T$在$\mathcal{B}$上的矩阵是上三角矩阵。
        \item $\forall \,\, 1 \leqslant i \leqslant n, \,\,\, T(\mathbf{e}_i) \in \mathbb{E}_i = \langle \mathbf{e}_1, \cdots , \mathbf{e}_i\rangle .$
        \item $\forall \,\, 1 \leqslant i \leqslant n, \,\,\, \mathbb{E}_i $是$T$-不变子空间。
    \end{enumerate}
\end{pp}
\begin{proof} % [{\bfseries 证明\ref{pp:6_5}}]
    很容易得出命题(1)和(2)等价，而且命题(3)可以推出命题(2)。下面证明(2)可以推出(3)。\\
    要证明$\mathbb{E}_i $是$T$-不变子空间，只需证明
    $\forall \,\, 1 \leqslant j \leqslant i$，$T(\mathbf{e}_j) \in \mathbb{E}_i$。
    而由(2)，我们已知$T(\mathbf{e}_j) \in \mathbb{E}_j$，后者是$\mathbb{E}_i$的子空间。
    所以$\forall \,\, 1 \leqslant j \leqslant i$，$T(\mathbf{e}_j) \in \mathbb{E}_i$。
    
\end{proof}

我们称一组这样的基为关于$T$的\textbf{三角基}，称把$T$用上三角矩阵表示为\textbf{三角化}，
将可以在某组基上表达为上三角矩阵的自同态称为\textbf{可三角化自同态}。
如果某个矩阵对应一个可三角化自同态，就称它为\textbf{可三角化矩阵}。
矩阵是可三角化矩阵，表示它和某个上三角矩阵是相似矩阵。

可三角化自同态对应着上三角矩阵，对应着一组嵌套的不变子空间：$\mathbb{E}_1 \subseteq \mathbb{E}_2 \subseteq \cdots \subseteq \mathbb{E}_n$.

我们用第一节的自同态来具体感受三角化的过程。$A$可以对角化，从而可以三角化。
$B$的矩阵为：
$$
B = \begin{bmatrix} 3 & -4 \\ 1 & -1 \end{bmatrix}.
$$
第一节中我们已经验证，$B$的唯一本征值为$1$，对应的本征向量为$\mathbf{e}_1 = (2,1)$。
也就是说，$B\mathbf{e}_1 = 1 \cdot \mathbf{e}_1$。
所以$1$维子空间$\mathbb{R} \mathbf{e}_1$是$B$-不变子空间。

要三角化$B$，按命题\ref{pp:6_5}，我们需要找一组基$(\mathbf{e}_1, \mathbf{e}_2)$，
使得$\langle \mathbf{e}_1 \rangle \subseteq \langle \mathbf{e}_1, \mathbf{e}_2 \rangle = \mathbb{R}^2$都是$B$-不变子空间。
第一个条件已经满足。第二个条件自动满足（整个空间总是不变子空间）。
所以，只需取任何与$\mathbf{e}_1$直无关的向量作为$\mathbf{e}_2$即可。

取$\mathbf{e}_2 = (1, 0)$。计算$B\mathbf{e}_2$：
$$
B\mathbf{e}_2 = \begin{bmatrix} 3 & -4 \\ 1 & -1 \end{bmatrix}\begin{bmatrix} 1 \\ 0 \end{bmatrix} = \begin{bmatrix} 3 \\ 1 \end{bmatrix} = 1 \cdot \mathbf{e}_1 + 1 \cdot \mathbf{e}_2.
$$
综合两个结果：
$$
B\mathbf{e}_1 = 1 \cdot \mathbf{e}_1 + 0 \cdot \mathbf{e}_2, \quad B\mathbf{e}_2 = 1 \cdot \mathbf{e}_1 + 1 \cdot \mathbf{e}_2.
$$
因此，$B$在基底$(\mathbf{e}_1, \mathbf{e}_2) = ((2,1), (1,0))$上的矩阵为：
$$
[B]_{\mathcal{B}} = \begin{bmatrix} 1 & 1 \\ 0 & 1 \end{bmatrix}.
$$
这是一个上三角矩阵。$B$被三角化了。

注意这个上三角矩阵不是对角矩阵。对角线上两个元素都是$1$（本征值），
但第$1$行第$2$列有一个非零元素$1$。这个非零元素反映了$B$不可对角化的事实：
$B$只有一个$1$维的本征空间，无法找到由本征向量构成的基底。

再来看$3$维的例子，对于第一节提到的自同态$C$，矩阵为：
$$
C = \begin{bmatrix} 5 & 0 & 1 \\ 1 & 3 & 2 \\ 0 & 0 & 3 \end{bmatrix}.
$$
上一节中我们算得：取$\mathbb{R}^3$的基底为$((0,1,0),\, (1,0,-2),\, (2,1,0))$，
则$C$在这组基上的矩阵为：
$$
\begin{bmatrix} 3 & -3 & 0 \\ 0 & 3 & 0 \\ 0 & 0 & 5 \end{bmatrix}.
$$
这是一个上三角矩阵，$C$被上三角化了。我们的处理方式是针对不同的本征值分别处理。$(2,1,0)$对应一维本征空间。$(0,1,0),\, (1,0,-2)$将本征值$3$对应的空间三角化。这也让我们看到：如果分块对角矩阵的每一对角块都是上（下）三角矩阵，那么整个矩阵也会是上（下）三角矩阵。

\subsection{上三角矩阵的性质}
如果某个直自同态$T$在某个基底上可以表示为上三角矩阵，
我们可以从中推出它的一些性质：

\begin{pp}（判定可逆映射）\label{pp:6_6}
    设自同态$T$可三角化，$[T]_\mathcal{B}$是对应的上三角矩阵。
    那么$T$是可逆映射，当且仅当$[T]_\mathcal{B}$对角线上所有元素不为0.
\end{pp}
\begin{proof} % [{\bfseries 证明\ref{pp:6_6}}]
    我们证明其等价命题：$T$不是可逆映射，当且仅当$[T]_\mathcal{B}$对角线上有某个元素是0.\\
    注意到直自同态是可逆映射（双射）等价于它是单射，我们转为证明：\\
    $T$不是单射，当且仅当$[T]_\mathcal{B}$对角线上有某个元素是0.\\
    设$\mathcal{B} = (\mathbf{e}_1, \cdots , \mathbf{e}_n)$是对应的三角基，$(\lambda_1, \cdots , \lambda_n)$是上三角矩阵$[T]_\mathcal{B}$对角线上的元素。\\
    首先证明：如果$[T]_\mathcal{B}$对角线上某个元素为0，则$T$不是单射。\\
    如果$\lambda_1 = 0$，那么$T(\mathbf{e}_1) = \mathbf{0}$，$T$不是单射。\\
    如果对某个$k>1$，$\lambda_k = 0$，那么$T(\mathbf{e}_k) \in \mathbb{E}_{k-1}$。
    这说明$T(\mathbb{E}_k) \subseteq \mathbb{E}_{k-1}$，也就是说，
    $T$限制在$\mathbb{E}_k$中的部分将一个$k$维空间映射到$k-1$维空间。
    所以，$\restr{T}{\mathbb{E}_k}$作为$\mathbb{E}_k$上的自同态不是满射，因此不是单射。
    可以找到非零向量$\mathbf{u} \in \mathbb{E}_k$使得$T(\mathbf{u}) = \mathbf{0}$，
    即是说，$T$不是单射。\\
    接下来证明：如果$T$不是可逆映射，那么$[T]_\mathcal{B}$对角线上某个元素为0.\\
    若$T$不是单射。存在非零向量$\mathbf{u} \in \mathbb{E}_k$使得$T(\mathbf{u}) = \mathbf{0}$。
    将$\mathbf{u}$在$\mathcal{B}$中分解：
    $$ \mathbf{u} = u_1 \mathbf{e}_1 + \cdots + u_n \mathbf{e}_n.$$
    设$u_k$是$u_1, \cdots u_n$之中最后一个不为0的系数，则：
    $$ \mathbf{u} = u_1 \mathbf{e}_1 + \cdots + u_k \mathbf{e}_k.$$
    如果$k=1$，则$T(\mathbf{e}_1) = \mathbf{0}$，即$\lambda_1 = 0$。如果$k > 1$，则有：
    $$ \mathbf{0} = T(\mathbf{u}) = T(u_1 \mathbf{e}_1 + \cdots + u_{k-1} \mathbf{e}_{k-1}) +  u_k T(\mathbf{e}_k).$$
    由于$\mathbb{E}_{k-1}$是$T$-不变子空间，等式右边第一项还在$\mathbb{E}_{k-1}$里，这说明
    $$ T(\mathbf{e}_k) \in \mathbb{E}_{k-1}$$
    因此$\lambda_k = 0.$

\end{proof}


\begin{pp}（可三角化映射的本征值）\label{pp:6_7}
    设自同态$T$可三角化，$[T]_\mathcal{B}$是对应的上三角矩阵。
    那么$[T]_\mathcal{B}$对角线上元素的集合就是$T$的本征值的集合。
\end{pp}

\begin{proof} % [{\bfseries 证明\ref{pp:6_7}}]
    对任意给定的$\lambda$，考虑映射$T - \lambda\mathrm{I}$在$\mathcal{B}$上的矩阵。\\
    此矩阵的对角线元素为$\lambda_1 - \lambda, \lambda_2 - \lambda, \cdots , \lambda_n - \lambda$，其中的$\lambda_1, \cdots, \lambda_n$是
    $T$三角化后对角线上的元素。\\
    $\lambda$是本征值，当且仅当$T - \lambda\mathrm{I}$不是可逆映射，当且仅当$\lambda$等于某个$\lambda_i$。
    也就是说，$T$的本征值集合就是$[T]_\mathcal{B}$对角线上元素的集合。

\end{proof}

我们用前面三角化的结果来验证这两个命题。

对$B$来说，三角化后的矩阵为：
$$
[B]_{\mathcal{B}} = \begin{bmatrix} 1 & 1 \\ 0 & 1 \end{bmatrix}.
$$
对角线元素为$1, 1$。按命题\ref{pp:6_7}，$B$的本征值集合就是$\{1\}$。
这与我们第一节中直接计算的结果一致。
对角线元素都不为$0$，按命题\ref{pp:6_6}，$B$是可逆映射。
验证：$\det B = 3 \times (-1) - (-4) \times 1 = 1 \neq 0$，确实可逆。

对$C$来说，三角化后的矩阵为：
$$
[C]_{\mathcal{B}} = \begin{bmatrix} 3 & -3 & 0 \\ 0 & 3 & 0 \\ 0 & 0 & 5 \end{bmatrix}.
$$
对角线元素为$3, 3, 5$。按命题\ref{pp:6_7}，$C$的本征值集合就是$\{3, 5\}$。
这也与第一节中的结果一致。注意$3$在对角线上出现了两次，
对应本征值$3$的全重数为$2$。

对$A$来说，第一节中我们已经将它对角化。
在基底$((1,-1), (3,-2))$上，$A$的矩阵为：
$$
[A]_{\mathcal{B}} = \begin{bmatrix} 1 & 0 \\ 0 & 2 \end{bmatrix}.
$$
对角矩阵是上三角矩阵的特殊情况。对角线元素$1, 2$就是$A$的本征值。

可以看到，可三角化的自同态，和可对角化的自同态一样，其本征值和对角元素挂钩。

\subsection{三角化与对角化}
那么，什么样的矩阵可以三角化呢？

\begin{pp}\label{pp:6_8}
    自同态$T$可三角化，当且仅当其极小归零式$m_T$是一次式的乘积。
\end{pp}

\begin{proof} % [{\bfseries 证明\ref{pp:6_8}}]
    首先证明：如果自同态$T$的极小归零式是一次式的乘积，那么$T$可三角化。\\
    我们对空间的维数$\dim \mathbb{V}$使用归纳法。$\dim \mathbb{V} = 1$时，结论显然成立。\\
    假设$\dim \mathbb{V} \leqslant n$时结论成立。当$\dim \mathbb{V} = n+1$的时候，考虑$T$的本征值。
    $m_T$是一次式的乘积，所以必然有根。由命题\ref{pp:6_3}，$T$必然有本征值（即$m_T$的根）$\lambda$。
    这说明$\Ker (T - \lambda\mathrm{I})$是至少1维的子空间，从而$\mathbb{W} = \Img (T - \lambda\mathrm{I})$至多是$n$维子空间。\\
    设$\dim \mathbb{W} = d \leqslant n$，$T$在$\mathbb{W}$上的限制$\restr{T}{\mathbb{W}}$的极小归零式$m_{T,\mathbb{W}}$是$m_T$的因式，因此也是一次式的乘积。
    所以$\restr{T}{\mathbb{W}}$可三角化。设$(\mathbf{e}_1, \cdots, \mathbf{e}_d)$是其三角基，将其补全为$\mathbb{V}$的基。考察$T$在这组基上的矩阵。\\
    对$1 \leqslant i \leqslant d$，我们已经有$T(\mathbf{e}_i) \in \langle \mathbf{e}_1, \cdots, \mathbf{e}_i\rangle$，而对$d < i \leqslant n+1$，
    $$T(\mathbf{e_i}) = (T - \lambda\mathrm{I})(\mathbf{e}_i) + \lambda \mathbf{e}_i$$
    前者属于$\mathbb{W}$，后者属于$\langle \mathbf{e}_i\rangle$，所以$T(\mathbf{e_i}) \in \mathbb{W}+\langle \mathbf{e}_i\rangle \subset  \langle \mathbf{e}_1, \cdots, \mathbf{e}_i\rangle.$
    根据命题\ref{pp:6_5}，$T$在这组基上是三角矩阵。\\
    因此，根据归纳原理，对任意正整数$n$，结论成立。\\
    接下来证明；如果$T$可三角化，那么$m_T$是一次式的乘积。\\
    设$A = (a_{ij})_{1\leqslant i, j\leqslant n}$是$T$三角化后的上三角矩阵，其对角元素是$a_{11}, a_{22}, \cdots, a_{nn}$，
    $\mathcal{B} = (\mathbf{e}_1, \cdots, \mathbf{e}_n)$是其三角基。
    那么，对任意$1\leqslant i \leqslant n$，$(T - a_{ii}\mathrm{I})(\mathbf{e}_i) = \sum_{j=1}^{i-1}a_{ji}\mathbf{e}_j$。这表示
    $$ \forall \,\,1\leqslant j \leqslant i \leqslant n, \,\,\, (T - a_{11}\mathrm{I})(T - a_{22}\mathrm{I})\cdots(T - a_{ii}\mathrm{I})(\mathbf{e}_j) = \mathbf{0}. $$
    从而当$i=n$时，有
    $$\forall \,\,1\leqslant i \leqslant n, \,\,\, \prod_{j=1}^n(T - a_{jj}\mathrm{I})(\mathbf{e}_i) = \mathbf{0}.$$
    定义整式$c_T$：
    $$
    c_T(x) = \prod_{j=1}^n(x - a_{jj}).
    $$
    则我们可以把前面的结论写作：
    $$
    c_T(T) = 0.
    $$
    我们把$c_T$称作$T$的\textbf{本征式}。上式说明，$c_T$是$T$的归零式，而且它是一次式的乘积。所以，$m_T$作为$c_T$的因式，也是一次式的乘积。 

\end{proof}

我们用第一节中的三个自同态来验证这个命题。

$A$的极小归零式为$m_A(z) = (z-1)(z-2)$，是两个一次式的乘积。
按命题\ref{pp:6_8}，$A$可三角化。事实上，$A$可以对角化，
对角化当然也是三角化。

$B$的极小归零式为$m_B(z) = (z-1)^2$，也是一个一次式的乘积（$(z-1)$出现两次）。
按命题\ref{pp:6_8}，$B$可三角化。前面我们已经具体找到了$B$的三角基。

$C$的极小归零式为$m_C(z) = (z-3)^2(z-5)$，同样是一次式的乘积。
按命题\ref{pp:6_8}，$C$可三角化。前面我们也具体找到了$C$的三角基。

三个自同态的极小归零式都能分解为一次式的乘积，所以它们都可三角化。

自同态能否三角化，取决于它的极小归零式能否写成一次式的乘积。
那么，自同态的极小归零式是否总能写成一次式的乘积呢？这和它的系数有关。
我们知道，实系数的整式不总能写成一次式的乘积，
比如$z^2+1$、$z^4+2$等等。因此，如果$T$是实系数域的平直空间上的自同态，
$T$不一定能三角化。

要使得极小归零式总能写成一次式的乘积，系数域$\mathbb{F}$要能使$\mathbb{F}[z]$中的任意非常整式总有根。
我们把这样的$\mathbb{F}$称为\textbf{代数闭域}或称它\textbf{代数封闭}，意思是，
解$\mathbb{F}$为系数的整式方程，根总还是在$\mathbb{F}$里面。这个词源自“代数数”，
表示$\mathbb{F}$对于代数操作（当时主要是指解整式方程）封闭。

那么，有哪些数域是代数闭域呢？我们接触过的例子里，只有复数域$\mathbb{C}$是代数闭域。
这个结论被称为\textbf{代数基本定理}，由高斯等人在19世纪初严格证明。

\begin{pp}\textbf{代数基本定理}\label{pp:6_9}
    任何非常数的复系数整式至少有一个根。
\end{pp}
从中容易推出：对任何正整数$n$，$n$次复系数整式在$\mathbb{C}$中有$n$个根（重根重复计算），可以写成$n$个一次式的乘积。

因此，对于复系数有限维平直空间\footnote{以下也简称为“复平直空间”。同理，我们把“实系数有限维平直空间”简称为“实平直空间”。}的自同态，我们可以宣称：
\begin{pp}\label{pp:6_10}
    任何复系数有限维平直空间的自同态都可三角化。
\end{pp}

我们来看一个在实空间中不可三角化的例子。
考虑$\mathbb{R}^2$上的自同态$D$，其矩阵为：
$$
D = \begin{bmatrix} 0 & -1 \\ 1 & 0 \end{bmatrix}.
$$
验证$D$的极小归零式。计算：
$$
D^2 = \begin{bmatrix} 0 & -1 \\ 1 & 0 \end{bmatrix}\begin{bmatrix} 0 & -1 \\ 1 & 0 \end{bmatrix} = \begin{bmatrix} -1 & 0 \\ 0 & -1 \end{bmatrix} = -\mathrm{I}.
$$
所以$D^2 + \mathrm{I} = 0$，即整式$z^2 + 1$是$D$的归零式。
而$D$不是缩放映射（$\forall\;\lambda \in \mathbb{R},\;\;D \neq \lambda \mathrm{I}$），
所以一次归零式不存在。因此$D$的极小归零式为：
$$
m_D(z) = z^2 + 1.
$$
整式$z^2 + 1$在$\mathbb{R}$中没有根，不能分解为实系数一次式的乘积。
按命题\ref{pp:6_8}，$D$在$\mathbb{R}^2$上不可三角化。

从直观上看，$D$的效果是把平面$\mathbb{R}^2$里的向量逆时针旋转$90°$。
旋转$90°$后，没有非零向量与自身共轴，所以没有本征向量，也没有本征值，
于是也没有$1$维不变子空间。没有$1$维不变子空间，就不可能有三角基。

然而，如果我们把$D$看作$\mathbb{C}^2$上的自同态，情况就不同了。
$z^2 + 1 = (z - \imath)(z + \imath)$在$\mathbb{C}$中可以分解为一次式的乘积。
按命题\ref{pp:6_10}，$D$在$\mathbb{C}^2$上可三角化。
实际上，$D$在$\mathbb{C}^2$上可以对角化：本征值为$\imath$和$-\imath$，
对应的本征向量分别是$(1, -\imath)$和$(1, \imath)$。

这个例子说明，同一个矩阵，在不同的系数域上，可能有完全不同的约简性质。

自同态能否三角化，和极小归零式的性质有关。
那么自同态能否对角化，是否也和极小归零式的性质有关呢？

显然，可对角化的自同态必然可三角化（因为对角矩阵就是一种上三角矩阵）。
换言之，可对角化的自同态如果可以和极小归零式的性质挂钩，
要求的必然是一种比“极小归零式可分解为一次式乘积”更严格的条件。

\begin{pp}\label{pp:6_11}
    自同态$T$可对角化，当且仅当其极小归零式$m_T$是一次式的乘积，且没有重根。我们称这样的整式为\textbf{散式}，或说整式的根\textbf{散在域中}，或\textbf{在域中散开}。
\end{pp}

\begin{proof} % [{\bfseries 证明\ref{pp:6_11}}]
    首先证明，如果$T$可对角化，则$m_T$是一次式乘积，并且没有重根。\\
    设$T$在某个基上为对角矩阵，其对角元素为$\lambda_1, \cdots, \lambda_n$。
    设其中不重复的元素为$\lambda_{(1)}, \cdots, \lambda_{(k)}$（$k \leqslant n$），在本章中我们已经验证过，
    $ \pi(z) = \prod_{i=1}^k (z - \lambda_{(i)})$是$T$的归零式。所以$m_T$是$\pi$的因式。
    $\pi$是一次式的乘积，且没有重根，所以$m_T$也是一次式的乘积，并且没有重根。\\
    接下来证明，如果$m_T$是一次式乘积，并且没有重根，则$T$可对角化。\\
    设$m_T(z) = \prod_{i=1}^k(z - \lambda_i)$，其中$\lambda_1, \cdots , \lambda_k$两两不等。
    根据零空间引理，
    $$\mathbb{V} = \Ker (T - \lambda_1\mathrm{I}) \oplus \cdots \oplus \Ker (T - \lambda_k\mathrm{I}).$$
    考虑其中任意一个$T$-不变子空间$\Ker (T - \lambda_i\mathrm{I})$，设$\mathbf{e}_{i,1}, \cdots , \mathbf{e}_{i,d_i}$是它的一组基底，
    则$\forall 1 \leqslant j \leqslant d_i, \,\,\, T(\mathbf{e}_{i,j}) = \lambda_i \mathbf{e}_{i,j}$。
    这说明，在这组基上，$\Ker (T - \lambda_i\mathrm{I})$是对角矩阵：
    $$ \begin{bmatrix}
        \lambda_i & 0 & \cdots & 0 \\
        0 & \lambda_i & \cdots & 0 \\
        \vdots & \vdots & \ddots & \vdots \\
        0 & \cdots & 0 & \lambda_i
    \end{bmatrix}$$
    对所有$k$个子空间，都可以找出这样的基底。它们合并起来，得到的基底$\mathcal{B}$就是使$[T]_\mathcal{B}$为对角矩阵的基底。

\end{proof}

最后，我们用前面提到的四个自同态来总结三角化与对角化的关系。

$A$的极小归零式$m_A(z) = (z-1)(z-2)$是一次式的乘积，且没有重根。
按命题\ref{pp:6_11}，$A$可对角化。事实上，第一节中我们已经将$A$对角化了。

$B$的极小归零式$m_B(z) = (z-1)^2$是一次式的乘积，但有重根（$z=1$是二重根）。
按命题\ref{pp:6_8}，$B$可三角化；但按命题\ref{pp:6_11}，$B$不可对角化。
这正好对应了我们前面的计算：$B$可以化为上三角矩阵$\begin{bmatrix} 1 & 1 \\ 0 & 1 \end{bmatrix}$，
但无法化为对角矩阵。

$C$的极小归零式$m_C(z) = (z-3)^2(z-5)$也是一次式的乘积，但有重根（$z=3$是二重根）。
同理，$C$可三角化但不可对角化。

$D$的极小归零式$m_D(z) = z^2 + 1$不是一次式的乘积（在$\mathbb{R}$上）。
$D$既不可三角化，也不可对角化。

将四个自同态对比：
\begin{center}
\begin{tabular}{|c|c|c|c|c|}
\hline
 & $A$ & $B$ & $C$ & $D$ \\
\hline
极小归零式 & $(z-1)(z-2)$ & $(z-1)^2$ & $(z-3)^2(z-5)$ & $z^2+1$ \\
\hline
一次式乘积 & 是 & 是 & 是 & 否（在$\mathbb{R}$上） \\
\hline
重根 & 无 & 有 & 有 & -- \\
\hline
三角化 & 能 & 能 & 能 & 否 \\
\hline
对角化 & 能 & 否 & 否 & 否 \\
\hline
\end{tabular}
\end{center}

可以看到，对角化是比三角化更强的要求。
可三角化只要求极小归零式是一次式的乘积，
可对角化要求极小式不仅是一次式的乘积，而且不能有重根。
而极小归零式不是一次式的乘积时，连三角化都做不到。

\section{本征空间}
\subsection{和本征值有关的子空间}
在前面的讨论中，我们发现，自同态能否对角化，与其本征值作为极小归零式的根的重数有关。
我们继续对此进行研究。

设$\lambda$是自同态$T$的一个本征值，它也是$T$的极小归零式的根。由于
$T - \lambda \mathrm{I}$不是单射，$\Ker (T - \lambda \mathrm{I})$是至少1维的子空间。
我们把它叫做$T$关于本征值$\lambda$的\textbf{本征空间}，记为$E_T(\lambda)$。空间中的（非零）向量都是$T$的本征向量，对应本征值$\lambda$。

零空间引理告诉我们，整个空间$\mathbb{V}$可以写成极小归零式分解成的两两互素的因式映射的零空间的直和：
$$ \mathbb{V} = \Ker P_1(T) \oplus \Ker P_2(T) \oplus \cdots \Ker P_k(T)$$
而一次式$z - \lambda$是其中某个整式$P_{k,\lambda}(z)$的因式。由于这些整式两两互素，
所以$z - \lambda$不再是其他整式的因式了。此外，设$P_{k,\lambda}(z) = (z - \lambda)^{d} Q(z)$。
由于$z - \lambda$是素整式，$(z - \lambda)^{d}$和$Q(z)$互素，所以我们总能把上面整个空间的直和分解直接写成：
$$ \mathbb{V} = \Ker (T - \lambda\mathrm{I})^{d_\lambda} \oplus P(T),$$
也就是$(T - \lambda\mathrm{I})^{d_\lambda}$的零空间和另一个与$z - \lambda$互素的整式对应映射的零空间的直和。
其中$d_\lambda$就是本征值作为极小归零式的根的重数。

这让我们想到研究形如$\Ker (T - \lambda\mathrm{I})^i$的不变子空间。我们把这类空间称为\textbf{广义本征空间}。

\begin{df}\label{df:6_4}
    设$\lambda$为自同态$T$的本征值，$i$为自然数，则$\Ker (T - \lambda\mathrm{I})^i$称为$T$关于$\lambda$的（$i$阶）\textbf{广义本征空间}。
    $i$称为广义本征空间的\textbf{阶}。广义本征空间中的非零向量称为$T$关于$\lambda$的\textbf{广义本征向量}。
\end{df}

广义本征空间之间有什么关联呢？我们首先观察到如下的性质：
\begin{pp}\label{pp:6_12_1}
    如果整式$Q$是$P$的因式，则$\Ker Q(T) \subseteq \Ker P(T)$.
\end{pp}
\begin{proof} % [{\bfseries 证明\ref{pp:6_12_1}}]
    设$P = R \cdot Q$，则对$\Ker Q(T)$中向量$\mathbf{v}$，总有
    $$P(T)(\mathbf{v}) = (R(T)\circ Q(T))(\mathbf{v}) = R(T)(Q(T)(\mathbf{v})) = R(T)(\mathbf{0}) = \mathbf{0}. $$
    所以$\mathbf{v} \in \Ker P(T).$

\end{proof}

根据这个性质，容易推出
\begin{align*}
    \{\mathbf{0}\} = \Ker \mathrm{I} &= \Ker (T - \lambda\mathrm{I})^0 \subseteq \Ker (T - \lambda\mathrm{I})^1  \\
    &\subseteq \Ker (T - \lambda\mathrm{I})^2 \subseteq \cdots \subseteq \Ker (T - \lambda\mathrm{I})^i \subseteq \Ker (T - \lambda\mathrm{I})^{i+1} \subseteq \cdots 
\end{align*}
也就是说，从0维空间开始，广义本征空间总是更高阶的广义本征空间的子空间。随着阶数增大，空间维数也不断增大，或维持不变。
这是一个类似于俄罗斯套娃的结构，我们称为\textbf{空间套}。三角基也对应一个空间套，但我们会看到，两者性质有所差异。

以前面提到的自同态$B$和$C$为例。对$B$来说，本征值为$1$。我们已经知道：
$$
\Ker(B - \mathrm{I}) = \mathbb{R}(2,1),
$$
维数为$1$。接下来计算$\Ker(B - \mathrm{I})^2$。首先：
$$
(B - \mathrm{I})^2 = \begin{bmatrix} 2 & -4 \\ 1 & -2 \end{bmatrix}^2 = \begin{bmatrix} 4-4 & -8+8 \\ 2-2 & -4+4 \end{bmatrix} = \begin{bmatrix} 0 & 0 \\ 0 & 0 \end{bmatrix}.
$$
所以$\Ker(B - \mathrm{I})^2 = \mathbb{R}^2$，维数为$2$。空间套为：
$$
\{\mathbf{0}\} \subset \mathbb{R}(2,1) \subset \Ker(B - \mathrm{I})^2 = \mathbb{R}^2.
$$
维数从$0$增长到$1$，再增长到$2$，之后不再增长（因为已经是整个空间了）。

对$C$来说，我们关注本征值$3$。我们已经知道：
$$
\Ker(C - 3\mathrm{I}) = \mathbb{R}(0,1,0),
$$
维数为$1$。接下来计算$\Ker(C - 3\mathrm{I})^2$。首先：
$$
(C - 3\mathrm{I})^2 = \begin{bmatrix} 2 & 0 & 1 \\ 1 & 0 & 2 \\ 0 & 0 & 0 \end{bmatrix}^2 = \begin{bmatrix} 4 & 0 & 2 \\ 2 & 0 & 1 \\ 0 & 0 & 0 \end{bmatrix}.
$$
解方程$(C - 3\mathrm{I})^2\mathbf{u} = \mathbf{0}$：
$$
\begin{bmatrix} 4 & 0 & 2 \\ 2 & 0 & 1 \\ 0 & 0 & 0 \end{bmatrix}\begin{bmatrix} u_1 \\ u_2 \\ u_3 \end{bmatrix} = \begin{bmatrix} 0 \\ 0 \\ 0 \end{bmatrix}.
$$
两个方程分别是$4u_1 + 2u_3 = 0$和$2u_1 + u_3 = 0$，它们等价。
解得$u_3 = -2u_1$，$u_2$自由。所以$\Ker(C - 3\mathrm{I})^2$的维数为$2$，一组基为$(1,0,-2)$和$(0,1,0)$。

再计算$\Ker(C - 3\mathrm{I})^3$。由于
$$
(C - 3\mathrm{I})^3 = (C - 3\mathrm{I})^2(C - 3\mathrm{I}) = \begin{bmatrix} 4 & 0 & 2 \\ 2 & 0 & 1 \\ 0 & 0 & 0 \end{bmatrix}\begin{bmatrix} 2 & 0 & 1 \\ 1 & 0 & 2 \\ 0 & 0 & 0 \end{bmatrix} = \begin{bmatrix} 8 & 0 & 4 \\ 4 & 0 & 2 \\ 0 & 0 & 0 \end{bmatrix},
$$
解方程$(C - 3\mathrm{I})^3\mathbf{u} = \mathbf{0}$得到$8u_1 + 4u_3 = 0$，即$u_3 = -2u_1$，$u_2$自由。
所以$\Ker(C - 3\mathrm{I})^3$的维数仍为$2$，与$\Ker(C - 3\mathrm{I})^2$相同。
空间套为：
$$
\{\mathbf{0}\} \subset \mathbb{R}(0,1,0) \subset \langle (1,0,-2),\, (0,1,0) \rangle = \Ker(C - 3\mathrm{I})^2 = \Ker(C - 3\mathrm{I})^3 \cdots
$$
维数从$0$增长到$1$，再增长到$2$，之后不再增长。事实上，我们知道对于更大的$i$，$\Ker(C - 3\mathrm{I})^i$的维数也不会增长了。因为如果$\Ker(C - 3\mathrm{I})^i$增长到$3$维，就说明$(C - 3\mathrm{I})^i$是零变换。但我们知道本征值$5$对应的本征向量$\mathbf{u}$使得$(C - 3\mathrm{I})^i\mathbf{u} = 2^i\mathbf{u}$，永远不是零向量。所以$(C - 3\mathrm{I})^i$永远不会是零变换。

实际上，由于整个空间维数有限，所以广义本征空间的维数必然会在某个时刻停止增长。
假设$i_0$是广义本征空间第一次停止维数增长的地方，即$\Ker (T - \lambda\mathrm{I})^{i_0} = \Ker (T - \lambda\mathrm{I})^{i_0+1}$，
我们可以证明，广义本征空间自此就不再“扩大”了；之后所有更高阶的广义本征空间，都是同一个空间。
\begin{pp}\label{pp:6_12_2}
    设$T$为自同态，如果对某个自然数$i$有$\Ker T^{i} = \Ker T^{i+1}$，那么
    $$ \forall j \geqslant i, \,\,\, \Ker T^{i} = \Ker T^{j}.$$
\end{pp}
\begin{proof} % [{\bfseries 证明\ref{pp:6_12_2}}]
    只需证明，对任意正整数$k$，$\Ker T^{i+k} = \Ker T^{i+k+1}$。而且，由于$T^{i+k}$是$T^{i+k+1}$的因式，
    $\Ker T^{i+k} \subseteq \Ker T^{i+k+1}$，所以只需证明$\Ker T^{i+k+1} \subseteq \Ker T^{i+k}.$\\
    设有$\mathbf{v}\in\Ker T^{i+k+1}$，则$T^{i+1}(T^k(\mathbf{v})) = \mathbf{0}$。但由于 $\Ker T^{i} = \Ker T^{i+1}$，
    所以$T^k(\mathbf{v}) \in \Ker T^{i+1} $，这说明$T^k(\mathbf{v}) \in \Ker T^{i}$，也就是说，
    $T^{i}(T^k(\mathbf{v})) = \mathbf{0}$，这就证明了$\mathbf{v}\in\Ker T^{i+k}$。\\
    所以，$\Ker T^{i+k+1} \subseteq \Ker T^{i+k}.$

\end{proof}

把这个一般结论用在$(T - \lambda\mathrm{I})$这个自同态上，就可以得到前面的说法。按照这个说法，
广义本征空间的维数必然是从一开始不断增长，然后从某个$i$开始一直稳定不变，形成一个稳定的$T$-不变子空间。
这个空间包含了$T$关于$\lambda$的所有广义本征向量。
我们把这个$i$记为$d_\lambda$。
\begin{df}\label{df:6_4_2}
    设$\lambda$为自同态$T$的本征值，$T$关于$\lambda$的\textbf{广义本征向量}生成的空间称为$T$关于$\lambda$的\textbf{全本征空间}，记为$H_T(\lambda)$。
    $$ H_T(\lambda) = \bigcup_{i\in\mathbb{N}} \Ker (T - \lambda\mathrm{I})^i.$$
\end{df}
$H_T(\lambda)$就是$T$中所有与本征值$\lambda$“有关”的向量形成的空间了。

由于广义本征空间最初增长的时候，每次增长维数至少增加1，而最终形成的子空间维数不超过$n$，
所以增长次数不会超过$n$次，也就是说，$d_\lambda\leqslant n$。

举例来说，假设$n=10$，$\nul (T - \lambda\mathrm{I}) = 3$，那么从此开始，
要么$\nul (T - \lambda\mathrm{I})^2 = 3$，停止扩张，维数不再增长，$d_\lambda=1$，
要么$\nul (T - \lambda\mathrm{I})^2 \geqslant 4$。
假设$\nul (T - \lambda\mathrm{I})^2 = 4$，接下来，
要么$\nul (T - \lambda\mathrm{I})^3 = 4$，停止扩张，维数不再增长，$d_\lambda=2$，
要么$\nul (T - \lambda\mathrm{I})^3 \geqslant 5$。以此类推，在增长得“最慢”而又“最久”，
即每次维数恰好增加1，且最终达到了
$\nul (T - \lambda\mathrm{I})^8 = 10$的情况下，接下来$\nul (T - \lambda\mathrm{I})^9$的维数不可能超过$10$，
所以$\Ker (T - \lambda\mathrm{I})^8 = \Ker (T - \lambda\mathrm{I})^9$，维数不再增长，这时$d_\lambda=8$。
而如果维数停止增长的情况发生在更早，就有$d_\lambda<8$。

既然$d_\lambda \leqslant n$，我们可以用$\Ker (T - \lambda\mathrm{I})^n$来指代$H_T(\lambda)$。

对于$H_T(\lambda)$中某个向量$\mathbf{u}$来说，考察$(T - \lambda\mathrm{I})(\mathbf{u}), (T - \lambda\mathrm{I})^2(\mathbf{u}), \cdots$。
必然存在$d$使得$(T - \lambda\mathrm{I})^{d}(\mathbf{u}) \neq \mathbf{0}$而$(T - \lambda\mathrm{I})^{d+1}(\mathbf{u}) = \mathbf{0}$，
所有这样的自然数$d$都不会大于前面提到的使广义本征空间稳定不变的$d_\lambda$，所以这样的$d$的集合是有限的。里面存在最大的一个$d$。
对所有$j>d$，都有$(T - \lambda\mathrm{I})^{j}(\mathbf{u}) = \mathbf{0}$。

我们知道，不同的本征值对应的本征向量直无关。对于广义本征向量，也有同样的性质。

\begin{pp}\label{pp:6_12_0}
设$T$为自同态，$T$的本征值$\lambda_1, \lambda_2, \cdots, \lambda_k$两两不相等。
设$\mathbf{u}_1, \cdots , \mathbf{u}_k$分别是它们对应的广义本征向量。
那么，$\mathbf{u}_1, \cdots , \mathbf{u}_k$直无关。
\end{pp}
\begin{proof} % [{\bfseries 证明\ref{pp:6_12_0}}]
给定广义本征向量$\mathbf{u}$，对应本征值$\lambda$。
设$d$是使得$(T - \lambda\mathrm{I})^{d}(\mathbf{u}) \neq \mathbf{0}$的最大自然数。
设$(T - \lambda\mathrm{I})^{d}(\mathbf{u}) = \mathbf{x}$，则
$$T\mathbf{x} - \lambda \mathbf{x} = (T - \lambda\mathrm{I})(\mathbf{x}) = (T - \lambda\mathrm{I})^{d+1}(\mathbf{u}) = \mathbf{0}.$$
$\mathbf{x}$是本征值为$\lambda$的本征向量。而对于其他标量$\mu$，
$(T - \mu\mathrm{I})(\mathbf{x}) = T\mathbf{x} - \mu\mathbf{x} = (\lambda - \mu)\mathbf{x}.$

设有直组合：
$$ a_1\mathbf{u}_1 + a_2\mathbf{u}_2 + \cdots + a_k\mathbf{u}_k = \mathbf{0},$$
设$d$是使得$(T - \lambda_1\mathrm{I})^{d}(\mathbf{u}_1) \neq \mathbf{0}$的最大自然数。
记$\mathbf{x} = (T - \lambda_1\mathrm{I})^{d}(\mathbf{u}_1)$，则$\mathbf{x}$为$\lambda_1$的本征向量。
在上式两边施以映射：
$(T - \lambda_1\mathrm{I})^{d}\prod_{i=2}^k (T - \lambda_i\mathrm{I})^{n}$.  
对左侧直组合的第一项，有：
\begin{align*}
(T - \lambda_1\mathrm{I})^{d}\prod_{i=2}^k (T - \lambda_i\mathrm{I})^{n}(a_1\mathbf{u}_1) 
&= a_1\prod_{i=2}^k (T - \lambda_i\mathrm{I})^{n}(T - \lambda_1\mathrm{I})^{d}(\mathbf{u}_1)  \\
&= a_1\prod_{i=2}^k (T - \lambda_i\mathrm{I})^{n}(\mathbf{x})  \\
&= a_1\prod_{i=2}^k (\lambda_1 - \lambda_i)^{n}\mathbf{x}.
\end{align*}
而对于其它项来说，由于$d_{\lambda_i} \leqslant n$，
$$ \forall i \neq 1, \quad (T - \lambda_i\mathrm{I})^{n}(\mathbf{u}_i) = \mathbf{0},$$
所以$a_i\mathbf{u}_i$经过映射的结果是$\mathbf{0}$.
所以经过映射$(T - \lambda_1\mathrm{I})^{d}\prod_{i=2}^k (T - \lambda_i\mathrm{I})^{n}$，
我们得到：
$$a_1\prod_{i=2}^k (\lambda_1 - \lambda_i)^{n}\mathbf{x} = \mathbf{0}.$$
本征值两两不等，因此可推出系数$a_1 = 0$。类似地，对其他的$1 < i \leqslant k $施以对应的映射，
可以证明其他系数$a_i$也是$0$. 所以$\mathbf{u}_1, \cdots , \mathbf{u}_k$直无关。

\end{proof}

\subsection{本重数与全重数}
结合以上讨论，我们引入以下的术语和记法：

\begin{df}\label{df:6_4_3}
    设$\lambda$为自同态$T$的本征值，称$E_T(\lambda)$的维数$\mu_g(\lambda)$为$T$关于$\lambda$的\textbf{本重数}，
    $H_T(\lambda)$的维数$\mu_a(\lambda)$为$T$关于$\lambda$的\textbf{全重数}。
    广义本征空间第一次停止增长的阶数称为本征值$\lambda$的\textbf{阶}，记为$d_\lambda$。
    \begin{align*}
        E_T(\lambda) &= \Ker (T - \lambda\mathrm{I})  \\
        \mu_g(\lambda) &= \dim E_T(\lambda) = \nul (T - \lambda\mathrm{I})  \\
        H_T(\lambda) &= \Ker (T - \lambda\mathrm{I})^{d_\lambda}  \\
        \mu_a(\lambda) &= \dim H_T(\lambda) = \nul (T - \lambda\mathrm{I})^{d_\lambda} 
    \end{align*}
\end{df}
显然，本征值的本重数总是小于等于全重数。
如果本重数等于全重数，那么广义本征空间从一开始就稳定，所以本征值的阶为1，本征空间就是全本征空间。
如果本征值的本重数小于全重数，那么它的阶至少为2。
更具体来说，设
$$ \forall \; i\,\in\, \mathbb{N},\quad b_i = \nul (T - \lambda\mathrm{I})^{i+1} - \nul (T - \lambda\mathrm{I})^{i},$$
则$b_0 = \mu_g(\lambda)$；$\forall i<d_{\lambda}$，$b_i>0$；$\forall i\geqslant d_{\lambda}$，$b_i=0$，于是
$$ \mu_a(\lambda) = \sum_{i\in\mathbb{N}} b_i = \sum_{i=0}^{d_{\lambda}-1} b_i = \mu_g(\lambda) + \sum_{i=1}^{d_\lambda-1} b_i.$$
另一方面，既然$\forall i<d_{\lambda}$，$b_i \geqslant 1$，就有：
$$ \mu_a(\lambda) \geqslant d_\lambda.$$

可三角化的自同态可以对角化，当且仅当它所有的本征值的本重数都等于全重数，阶为1，
也当且仅当所有本征值的本征空间就是全本征空间。

我们把三个自同态的本重数、全重数和阶整理如下：

对$A$来说：
\begin{center}
\begin{tabular}{c|c|c|c|c}
本征值$\lambda$ & $\mu_g(\lambda)$ & $\mu_a(\lambda)$ & $d_\lambda$ & 各$b_i$ \\
\hline
$1$ & $1$ & $1$ & $1$ & $b_0 = 1$ \\
$2$ & $1$ & $1$ & $1$ & $b_0 = 1$ \\
\end{tabular}
\end{center}
两个本征值的本重数都等于全重数，阶都是$1$。
空间套从第一步就稳定，本征空间就是全本征空间。

对$B$来说：
\begin{center}
\begin{tabular}{c|c|c|c|c}
本征值$\lambda$ & $\mu_g(\lambda)$ & $\mu_a(\lambda)$ & $d_\lambda$ & 各$b_i$ \\
\hline
$1$ & $1$ & $2$ & $2$ & $b_0 = 1,\; b_1 = 1$ \\
\end{tabular}
\end{center}
本征值$1$的本重数为$1$，全重数为$2$，阶为$2$。
空间套为$\{\mathbf{0}\} \subset \Ker(B-\mathrm{I}) \subset \Ker(B-\mathrm{I})^2$，
维数分别为$0, 1, 2$。每步增长$1$维：$b_0 = 1$，$b_1 = 2 - 1 = 1$。
全重数$\mu_a(1) = b_0 + b_1 = 1 + 1 = 2$。

对$C$来说：
\begin{center}
\begin{tabular}{c|c|c|c|c}
本征值$\lambda$ & $\mu_g(\lambda)$ & $\mu_a(\lambda)$ & $d_\lambda$ & 各$b_i$ \\
\hline
$3$ & $1$ & $2$ & $2$ & $b_0 = 1,\; b_1 = 1$ \\
$5$ & $1$ & $1$ & $1$ & $b_0 = 1$ \\
\end{tabular}
\end{center}
本征值$3$的本重数为$1$，全重数为$2$，阶为$2$，与$B$的情形类似。
本征值$5$的本重数等于全重数，阶为$1$，与$A$的情形类似。
$C$可以看作"$B$式结构"和"$A$式结构"的组合。

从表中可以观察到：
$A$的所有本征值都满足$\mu_g = \mu_a$，这对应$A$可对角化。
$B$和$C$都有某个本征值满足$\mu_g < \mu_a$，这对应它们不可对角化。

本征值$\lambda$总是极小归零式的根。下面，我们给出$\lambda$作为根在极小归零式中的重数和它的阶的关系。

\begin{pp}\label{pp:6_12_5}
    自同态$T$的本征值$\lambda$作为极小归零式$m_T$的根，对应的重数就是它的阶$d_\lambda$。
\end{pp}
\begin{proof} % [{\bfseries 证明\ref{pp:6_12_5}}]
设$T$的极小归零式写为：
$$ m_T(z) = (z - \lambda)^d q(z),$$
其中整式$q$中不含因式$z - \lambda$（因此也与$z - \lambda$互素）。我们使用反证法，
证明$d < d_{\lambda}$以及$d > d_{\lambda}$都会导致矛盾。\\
一方面，如果$d < d_{\lambda}$，这说明$\Ker (T - \lambda\mathrm{I})^{d} \subsetneq \Ker (T - \lambda\mathrm{I})^{d+1}$，
也就是说，$\exists \mathbf{u}\in \Ker (T - \lambda\mathrm{I})^{d+1}$，
使得$\mathbf{x} = (T - \lambda\mathrm{I})^{d}(\mathbf{u}) \neq \mathbf{0}$，
而$(T - \lambda\mathrm{I})(\mathbf{x}) = \mathbf{0}$。也就是说，$\mathbf{x}$是$\lambda$的本征向量。
然而，
$$q(T)(\mathbf{x}) = q(T)(T - \lambda\mathrm{I})^{d}(\mathbf{u}) = m_T(T)(\mathbf{u}) = \mathbf{0}.$$ 
由于整式$q$与$(z - \lambda)$互素，存在整式$r_1$和$r_2$使得$r_1 q + r_2 (z - \lambda) = 1$，对应的映射为：
$$ r_1(T) q(T) + r_2(T) (T - \lambda \mathrm{I}) = \mathrm{I}. $$
作用在向量$\mathbf{x}$上，就得到：
$$ \mathbf{x} = r_1(T) (q(T)(\mathbf{x})) + r_2(T)((T - \lambda \mathrm{I})(\mathbf{x})) = \mathbf{0} + \mathbf{0} = \mathbf{0}. $$
这与$\mathbf{x}$的定义矛盾了。 \\
另一方面，如果$d > d_{\lambda}$，则说明$\Ker (T - \lambda\mathrm{I})^{d - 1} = \Ker (T - \lambda\mathrm{I})^{d}$。那么考虑整式
$$p(z) = (z - \lambda)^{d-1} q(z)$$
$\forall \mathbf{v} \in \mathbb{V}$，零空间引理说明，可以把$\mathbf{v}$写成
$$\mathbf{v} = \mathbf{v}_1 + \mathbf{v}_2, \,\,\, \mathbf{v}_1 \in \Ker (T - \lambda\mathrm{I})^{d}, \,\,\, \mathbf{v}_2 \in \Ker q(T) $$
由于$\Ker (T - \lambda\mathrm{I})^{d- 1} = \Ker (T - \lambda\mathrm{I})^{d}$，
$(T - \lambda\mathrm{I})^{d- 1}(\mathbf{v}_1) = \mathbf{0}$。
所以，将$p(T)$作用到$\mathbf{v}$上，就有
\begin{align*} 
    p(T)(\mathbf{v}) &= p(T)(\mathbf{v}_1) + p(T)(\mathbf{v}_2)  \\
    &= q(T)((T - \lambda\mathrm{I})^{d- 1}(\mathbf{v}_1)) + (T - \lambda\mathrm{I})^{d- 1}(q(T)(\mathbf{v}_2))  \\
    &= \mathbf{0} + \mathbf{0}  \\
    &= \mathbf{0}. 
\end{align*}
这又与极小归零式的定义矛盾了。\\
以上两方面论证说明，$d$就是$\lambda$的阶$d_{\lambda}$。

\end{proof}

以前面提到的$B$和$C$为例。$B$的极小归零式为$m_B(z) = (z-1)^2$。
本征值$1$在$m_B$中的重数为$2$。
前面我们算出$B$关于本征值$1$的阶$d_1 = 2$（因为$\Ker(B-\mathrm{I}) \subsetneq \Ker(B-\mathrm{I})^2$，
而$\Ker(B-\mathrm{I})^2 = \Ker(B-\mathrm{I})^3 = \cdots$）。
两者相等，与命题\ref{pp:6_12_5}一致。

$C$的极小归零式为$m_C(z) = (z-3)^2(z-5)$。
本征值$3$在$m_C$中的重数为$2$，本征值$5$在$m_C$中的重数为$1$。
前面我们算出$C$关于本征值$3$的阶$d_3 = 2$，关于本征值$5$的阶$d_5 = 1$。
两者分别相等，与命题\ref{pp:6_12_5}一致。

$A$的极小归零式为$m_A(z) = (z-1)(z-2)$。
本征值$1$和$2$在$m_A$中的重数都是$1$。
$A$关于两个本征值的阶也都是$1$。同样一致。

可以看到，极小归零式中某个本征值的重数，恰好刻画了该本征值对应的空间套增长了“多少步”才稳定。

如果$T$的极小归零式是一次式的乘积，那么可以表示成：
$$ m_T(z) = (z - \lambda_1)^{d_1} (z - \lambda_2)^{d_2} \cdots  (z - \lambda_k)^{d_k}.$$
其中$\lambda_1, \lambda_2,\cdots$就是$T$的不同的本征值。零空间引理告诉我们：
$$ \mathbb{V} = \Ker (T - \lambda_1\mathrm{I})^{d_1} \oplus \Ker (T - \lambda_2\mathrm{I})^{d_2} \oplus \cdots \oplus \Ker (T - \lambda_k\mathrm{I})^{d_k}.$$
对每个本征值来说，它作为根在极小归零式里的重数$d_i$就是它的阶$d_{\lambda_i}$，
$\Ker (T - \lambda_i\mathrm{I})^{d_i}$就是$H_T(\lambda_i)$。
所以从维数上统计，
$$ \sum_{i=1}^k \mu_a(\lambda_i) = n.$$
\begin{pp}\label{pp:6_12_2_1}
    所有本征值的全重数之和等于空间的维数。
\end{pp}

另外有$d_i \leqslant \mu_a(\lambda_i)$，加起来就得到：
$$ \deg m_T = \sum_{i=1}^k d_i \leqslant \sum_{i=1}^k \mu_a(\lambda_i) = n.$$
\begin{pp}\label{pp:6_12_2_2}
    极小归零式的次数不超过空间的维数。
\end{pp}

以前面提到的自同态$A,B,C$为例。

对$A$来说，两个本征值的全重数之和为$\mu_a(1) + \mu_a(2) = 1 + 1 = 2 = \dim\mathbb{R}^2$。
极小归零式$m_A(z) = (z-1)(z-2)$的次数为$2$，不超过$\dim\mathbb{R}^2 = 2$。

对$B$来说，唯一本征值的全重数为$\mu_a(1) = 2 = \dim\mathbb{R}^2$。
极小归零式$m_B(z) = (z-1)^2$的次数为$2$，不超过$\dim\mathbb{R}^2 = 2$。

对$C$来说，两个本征值的全重数之和为$\mu_a(3) + \mu_a(5) = 2 + 1 = 3 = \dim\mathbb{R}^3$。
极小归零式$m_C(z) = (z-3)^2(z-5)$的次数为$3$，不超过$\dim\mathbb{R}^3 = 3$。

三个例子中，极小归零式的次数恰好等于空间维数。
但这并非总是如此。比如，缩放映射$T = 2\mathrm{I}$在$\mathbb{R}^n$上的极小归零式是$z - 2$，
次数为$1$，可以远小于空间维数$n$。
极小归零式的次数等于空间维数，当且仅当自同态的每个本征值的阶都等于其全重数。
对$A$、$B$、$C$来说，恰好都满足这个条件。

\section{幂零自同态与标准形式}
\subsection{幂零自同态}
接下来，我们来看看$T$能否在$H_T(\lambda)$中约简。由前面的定义，
自同态$(T - \lambda\mathrm{I})^{d_\lambda}$限制在$H_T(\lambda)$上的部分就是零映射。这让我们想到，研究这样一类自同态：
\begin{df}\label{df:6_4_4}
    给定自同态$N$，如果存在自然数$d$，使得$N^d$为零映射，就称$N$为\textbf{幂零}的。
    使得$N^d$为零映射的最小自然数称为$N$的\textbf{阶}，称$N$为$d$阶幂零自同态。
\end{df}

显然，$z^d$是$d$阶幂零自同态的归零式。它的极小归零式必然是$z^d$的因式，因此就是$z^d$。它的本征值只有0。

举例来说，$\mathbb{R}^3$上的自同态$T: (x_1, x_2, x_3) \mapsto (x_2, x_3, 0)$就是一个幂零映射。验证如下：
\begin{align*}
    T^2(x_1, x_2, x_3) &= T(x_2, x_3, 0) = (x_3, 0, 0),  \\
    T^3(x_1, x_2, x_3) &= T(x_3, 0, 0) =  (0, 0, 0).
\end{align*}
它的阶是3。

再来看前面提到的几个自同态的幂零部分。对自同态$B$来说，它的唯一本征值为$1$，极小归零式为$(z-1)^2$。
考虑$N_B = B - \mathrm{I}$：
$$
N_B = \begin{bmatrix} 3 & -4 \\ 1 & -1 \end{bmatrix} - \begin{bmatrix} 1 & 0 \\ 0 & 1 \end{bmatrix} = \begin{bmatrix} 2 & -4 \\ 1 & -2 \end{bmatrix}.
$$
验证$N_B$是幂零的：
$$
N_B^2 = \begin{bmatrix} 2 & -4 \\ 1 & -2 \end{bmatrix}\begin{bmatrix} 2 & -4 \\ 1 & -2 \end{bmatrix} = \begin{bmatrix} 4-4 & -8+8 \\ 2-2 & -4+4 \end{bmatrix} = \begin{bmatrix} 0 & 0 \\ 0 & 0 \end{bmatrix}.
$$
而$N_B \neq 0$，所以$N_B$是$2$阶幂零自同态。

对自同态$C$来说，本征值$3$的全本征空间为$H_C(3) = \Ker(C-3\mathrm{I})^2$，维数为$2$。
考虑$N_C = C - 3\mathrm{I}$限制在$H_C(3)$上的部分。前面我们已经算出：
$$
C - 3\mathrm{I} = \begin{bmatrix} 2 & 0 & 1 \\ 1 & 0 & 2 \\ 0 & 0 & 0 \end{bmatrix}, \quad (C-3\mathrm{I})^2 = \begin{bmatrix} 4 & 0 & 2 \\ 2 & 0 & 1 \\ 0 & 0 & 0 \end{bmatrix}.
$$
$H_C(3) = \Ker(C-3\mathrm{I})^2$的一组基为$(0,1,0)$和$(1,0,-2)$。
验证$N_C$在$H_C(3)$上是幂零的：
$$
N_C(0,1,0) = (0, 0, 0), \quad N_C(1,0,-2) = (2-2,\, 1-4,\, 0) = (0,-3,0).
$$
而$(0,-3,0) = -3(0,1,0) \in H_C(3)$，所以$N_C$确实把$H_C(3)$映射到自身。
再算一步：
$$
N_C^2(1,0,-2) = N_C(0,-3,0) = (0, 0, 0).
$$
所以$N_C$限制在$H_C(3)$上是$2$阶幂零自同态。

对自同态$A$来说，它的极小归零式$(z-1)(z-2)$没有重根，
两个本征值的阶都是$1$。这意味着$A - \mathrm{I}$和$A - 2\mathrm{I}$
各自的本征空间就是全本征空间，不需要进一步考虑幂零部分。
$A$可以直接对角化，不涉及非平凡的幂零结构。

从广义本征空间的角度来看，幂零自同态$N$其实就是$N - 0\mathrm{I}$，
空间套$\Ker N, \Ker N^2, \cdots$的维数逐渐增长到整个空间的维数$n$，
以上定义的阶$d$就是对应本征值$0$的阶，所以$d$不超过$n$。

幂零自同态可以看作是一类较简单的矩阵。它有比一般矩阵更“良好”的性质。

给定$\mathbb{V}$中向量$\mathbf{v}$，考虑$\mathbf{v}$自己以及它经过若干次$N$映射得到的向量：
$$ \mathbf{v}, N\mathbf{v}, N^2 \mathbf{v}, \cdots , N^i \mathbf{v} ,\cdots $$
我们只考虑其中的非零向量。注意到$i \geqslant d$时，$N^i \mathbf{v} = \mathbf{0}$。而且如果$N^i \mathbf{v} = \mathbf{0}$，那么$N^{i+1} \mathbf{v} = \mathbf{0}$。所以以上序列剩下：
$$ \mathbf{v}, N\mathbf{v}, N^2 \mathbf{v}, \cdots , N^{d_\mathbf{v} - 1} \mathbf{v}, $$
其中$d_\mathbf{v}$是第一个让$N^i \mathbf{v} = \mathbf{0}$的自然数$i$。

如果$\mathbf{v}$是某个元素经$N$映射的像，那么可以“反向延展”这个序列，直到无法再这样做为止。于是，我们得到一个包含$\mathbf{v}$的如下形式的序列：
$$ \mathbf{u}, N\mathbf{u}, N^2 \mathbf{u}, \cdots , N^{d_\mathbf{u} - 1} \mathbf{u} $$
我们称它为$\mathbf{v}$（所在）的\textbf{轨迹}，或$N$在$\mathbf{v}$上留下的轨迹。轨迹中第一个向量称为它的出发向量，最后一个向量称为它的结束向量。
\begin{pp}\label{pp:6_12_3}
    同一条轨迹中的向量直无关。
\end{pp}
\begin{proof} % [{\bfseries 证明\ref{pp:6_12_3}}]
设有直组合：
$$ a_1 \mathbf{u} + a_2 N\mathbf{u} +\cdots + a_{d_\mathbf{u}} N^{d_\mathbf{u} - 1} \mathbf{u} = \mathbf{0}.$$
那么，两边同时经过映射$N^{d_\mathbf{u} - 1}$，可得$a_1 N^{d_\mathbf{u} - 1}\mathbf{u} + a_2 N^{d_\mathbf{u}}\mathbf{u} +\cdots + a_{d_\mathbf{u}} N^{2d_\mathbf{u} - 2} \mathbf{u} = \mathbf{0}$。注意到$N^{d_\mathbf{u}}\mathbf{u}= \mathbf{0}$，所以式子里第二项开始都是零向量。因此只剩下$a_1 N^{d_\mathbf{u} - 1}\mathbf{u} = 0$，这说明$a_1 = 0$。上式变成：
$$ a_2 N\mathbf{u} + a_3 N^2\mathbf{u} +\cdots + a_{d_\mathbf{u}} N^{d_\mathbf{u} - 1} \mathbf{u} = \mathbf{0}. $$
类似地，两边同时经过映射$N^{d_\mathbf{u} - 2}$，可得到$a_2 = 0$。以此类推，可以证明所有系数都为0。这就说明$\mathbf{u}, N\mathbf{u}, N^2 \mathbf{u}, \cdots , N^{d_\mathbf{u} - 1} \mathbf{u}$直无关。

\end{proof}

以自同态$B$和$C$为例。对$N_B = B - \mathrm{I} = \begin{bmatrix} 2 & -4 \\ 1 & -2 \end{bmatrix}$来说
取向量$\mathbf{u} = (1,0)$，追踪它在$N_B$下的轨迹：
$$
N_B(1,0) = (2, 1), \quad N_B(2,1) = (4-4,\, 2-2) = (0,0).
$$
所以$(1,0)$所在的轨迹为：
$$
(1,0) \;\longrightarrow\; (2,1) \;\longrightarrow\; \mathbf{0}.
$$
轨迹长度为$2$。出发向量是$(1,0)$，结束向量是$(2,1)$。
按命题\ref{pp:6_12_3}，$(1,0)$和$(2,1)$直无关。
它们构成$\mathbb{R}^2$的一组基。

对$N_C = C - 3\mathrm{I}$限制在$H_C(3)$上的部分来说，
取向量$\mathbf{u} = (-1, 0, 2)$，追踪它在$N_C$下的轨迹：
$$
N_C(-1,0,2) = (0, 3, 0), \quad N_C(0,3,0) = (0, 0, 0).
$$
所以$(-1,0,2)$所在的轨迹为：
$$
(-1,0,2) \;\longrightarrow\; (0,3,0) \;\longrightarrow\; \mathbf{0}.
$$
轨迹长度为$2$。出发向量是$(-1,0,2)$，结束向量是$(0,3,0)$。
注意结束向量$(0,3,0) = 3(0,1,0)$是本征值$3$的本征向量的$3$倍，
所以$N_C$把结束向量映射为零向量。
按命题\ref{pp:6_12_3}，$(-1,0,2)$和$(0,3,0)$直无关。
它们构成$H_C(3)$的一组基。

再来看$H_C(3)$中另一个向量$(0,1,0)$。它本身就是本征向量：
$$
N_C(0,1,0) = (0,0,0).
$$
所以$(0,1,0)$所在的轨迹只有它自己：
$$
(0,1,0) \;\longrightarrow\; \mathbf{0}.
$$
轨迹长度为$1$。这条轨迹是"退化"的：出发向量就是结束向量。
它对应的循环子空间就是$1$维的本征空间$\mathbb{R}(0,1,0)$。

可以看到，$H_C(3)$中向量$(-1,0,2)$的轨迹"覆盖"了$(0,1,0)$的轨迹：
前者的结束向量是后者的$3$倍。所以，我们不需要把两条轨迹都取出来作为基底。
只取较长的那条轨迹$(-1,0,2), (0,3,0)$就够了。

我们把$N$在某向量上留下的一条轨迹生成的空间称为它的$N$-\textbf{循环子空间}。由上可知，这个空间的基就是对应的轨迹，
而且是$N$-不变子空间。所以也称其为\textbf{循环不变子空间}，称轨迹为它的\textbf{循环基}。
从$\mathbf{u}$出发的轨迹生成的循环子空间记为$C_N(\mathbf{u})$。

我们来写出自同态$B,C$具体的循环子空间。

对$N_B$来说，向量$(1,0)$的轨迹为$(1,0), (2,1)$，所以：
$$
C_{N_B}((1,0)) = \langle (1,0),\, (2,1) \rangle = \mathbb{R}^2.
$$
整个空间就是一条轨迹生成的循环子空间。$N_B$只有一个循环子空间。

对$N_C$限制在$H_C(3)$上来说，向量$(-1,0,2)$的轨迹为$(-1,0,2), (0,3,0)$，所以：
$$
C_{N_C}((-1,0,2)) = \langle (-1,0,2),\, (0,3,0) \rangle = H_C(3).
$$
同样，$H_C(3)$整体就是一条轨迹生成的循环子空间。

这两个例子中，空间恰好等于一条轨迹生成的循环子空间。
但一般来说，空间可能需要多条轨迹的直和来表示。
比如考虑$\mathbb{R}^4$上的幂零自同态$T: (x_1,x_2,x_3,x_4) \mapsto (x_2, 0, x_4, 0)$，
则有两条轨迹：
$$
(0,1,0,0) \;\longrightarrow\; (1,0,0,0) \;\longrightarrow\; \mathbf{0}, \quad \mbox{和}\; (0,0,0,1) \;\longrightarrow\; (0,0,1,0) \;\longrightarrow\; \mathbf{0}.
$$
两条轨迹各长$2$，分别生成$2$维循环子空间。它们的直和就是$\mathbb{R}^4$。

在每个循环子空间$C_N(\mathbf{u}_i)$中，轨迹为$\mathbf{u}_i, N\mathbf{u}_i, \cdots, N^{e_i - 1}\mathbf{u}_i$。考察$N^{k+1}$的零空间相比$N^k$有什么变化。
$N^k$的零空间由使得$N^k(N^j\mathbf{u}_i)=\mathbf{0}$的向量$N^j\mathbf{u}_i$生成。因此这些向量的下标是使得$k+j\geqslant e_i$的$j$：$j=e_i-k, \cdots, e_i - 1$。
当$k$变为$k+1$时：
\begin{itemize}
\item 如果$e_i > k$，则$e_i-k>0$，$e_i-(k+1)\geqslant 0$。因此$\mathbf{u}_i, N\mathbf{u}_i, \cdots, N^{e_i - 1}\mathbf{u}_i$中使得$N^{k+1}N^j\mathbf{u}_i = \mathbf{0}$的向量会多一个。零空间维数增加$1$。
\item 如果$e_i \leqslant k$，则$e_i - k \leqslant 0$。即$N^k$已经把整个$C_N(\mathbf{u}_i)$映射为零向量，零空间维数已经是$e_i$，不再增加。
\end{itemize}
因此，$b_k(\lambda)$作为$p$个$C_N(\mathbf{u}_i)$的零空间增长的总维数，恰好等于轨迹长度严格大于$k$的循环子空间的个数：

\begin{pp}\label{pp:6_13_1}
$$ b_k(\lambda) = \left|\{1\leqslant i \leqslant p\mid e_i > k\}\right|. $$
\end{pp}

\begin{figure}[h]
    \vspace{4pt}
    \centering
    \def\UnitLength{0.7cm}        % 单位长度（x,y 等长）
    \begin{tikzpicture}[x=\UnitLength, y=\UnitLength]

    % ==================== 可调参数 ====================
    \def\AxisXMax{13}             % x 轴正半轴右端点
    \def\AxisYMax{10}             % y 轴正半轴上端点

    % 从下往上每排单位方格数；0 表示该排不画方格（本图第5排为圆点）
    % 列表中不要加入空格
    \def\SquareCounts{12,11,11,8,0,4,3,3,2}
    \def\SquareLineWidth{0.5pt}
    \def\AxisLineWidth{0.6pt}

    % 小黑圆点参数（原第5排半整点位置）
    \def\DotsRowIndex{5}          % 圆点所在排号（从下往上，1-based）
    \pgfmathsetmacro\DotsY{\DotsRowIndex - 0.5} % 该排中心 y
    \def\DotsCount{6}             % 圆点个数
    \def\DotsStartX{0.5}          % 第一个圆点中心 x（原第1列中心）
    \def\DotsSpacing{1}           % 相邻圆点中心间距
    \def\DotRadius{1.5pt}         % 圆点半径
    % =================================================
    
    \draw[fill=blue, fill opacity=0.3] (8, 0) rectangle (9, 3);
    % 先画方格
    \foreach \cnt [count=\row] in \SquareCounts {
    \ifnum\cnt>0
        \foreach \col in {1,...,\cnt} {
        \draw[line width=\SquareLineWidth]
            (\col-1,\row-1) rectangle (\col,\row);
        }
    \fi
    }

    % 再画小黑圆点
    \ifnum\DotsCount>0
    \pgfmathtruncatemacro\DotsLast{\DotsCount-1}
    \foreach \k in {0,...,\DotsLast} {
        \pgfmathsetmacro\DotX{\DotsStartX + \k*\DotsSpacing}
        \fill[black] (\DotX,\DotsY) circle[radius=\DotRadius];
    }
    \fi

    % 最后画坐标轴（只画正半轴，盖在方格边界上保持清晰）
    \draw[-latex, line width=\AxisLineWidth] (0,0) -- (\AxisXMax,0);
    \draw[-latex, line width=\AxisLineWidth] (0,0) -- (0,\AxisYMax);

    % ---------- 标签相关可调参数 ----------
    \def\LabelXSep{0.15}        % 左侧标签到 y 轴的距离（单位与坐标轴相同）
    \def\LabelXBelowSep{0.12}   % 下方标签到 x 轴的距离
    \def\LabelFont{\footnotesize} % 标签字体大小，可改为 \small、\scriptsize 等

    \def\VdotsY{5.0}            % 竖省略号 \vdots 的 y 坐标，可调
    \def\CdotsX{6.5}            % 横省略号 \cdots 的 x 坐标，可调

    \def\TotalRows{9}           % 总排数，用于计算最后一排中心
    \def\MaxCols{12}            % 最大列数，用于计算最后一列中心

    % ---------- y 轴左侧标签 ----------
    \pgfmathsetmacro\LastRowCenter{\TotalRows - 0.5}
    \node[anchor=east, font=\LabelFont] at (-\LabelXSep, 0.5) {$e_1$};
    \node[anchor=east, font=\LabelFont] at (-\LabelXSep, 1.5) {$e_2$};
    \node[anchor=east, font=\LabelFont] at (-\LabelXSep-0.2, \VdotsY) {$\vdots$};
    \node[anchor=east, font=\LabelFont] at (-\LabelXSep-0.2, 1.5+1) {$\vdots$};
    \node[anchor=east, font=\LabelFont] at (-\LabelXSep-0.2, \LastRowCenter-1) {$\vdots$};

    \node[anchor=east, font=\LabelFont] at (-\LabelXSep, \LastRowCenter) {$e_p$};
    % ---------- x 轴下方标签 ----------
    \node[anchor=north, font=\LabelFont] at (0.5, -\LabelXBelowSep) {$b_0$};
    \node[anchor=north, font=\LabelFont] at (1.5, -\LabelXBelowSep) {$b_1$};
    \node[anchor=north, font=\LabelFont] at (1.5+1, -\LabelXBelowSep-0.2) {$\cdots$};
    \node[anchor=north, font=\LabelFont] at (\CdotsX, -\LabelXBelowSep-0.2) {$\cdots$};

    % 最后一列 b_{d_\lambda-1} 用 \MaxCols 自动计算中心
    \pgfmathsetmacro\LastColCenter{\MaxCols - 0}
    \node[anchor=north, font=\LabelFont] at (\LastColCenter-1.5, -\LabelXBelowSep-0.2) {$\cdots$};
    \node[anchor=north, font=\LabelFont] at (\LastColCenter, -\LabelXBelowSep)
    {$b_{d_{\lambda}-1}$};

    \node at (11,5) {$b_k = |\{e_i > k\}|$};
    \draw[-{Triangle[scale=0.8]}] (9.3,4.6) -- (8.5, 2.5);

    \end{tikzpicture}
    \caption*{%
        \begin{tabular}{@{}l@{}}
            \texttt{用横排的方格表示轨道长度$e_1\geqslant e_2\geqslant \cdots \geqslant e_p$，}\\
            \texttt{则$b_k$（左起第$k+1$列高度）就是长度超过$k$的轨道个数。}
        \end{tabular}%
    }
\end{figure}

特别地，$b_0(\lambda) = \mu_g(\lambda) = p$，因为所有轨迹长度都大于$0$。具体来说，我们有：
\begin{enumerate}
\item 本重数$\mu_g(\lambda) = p = b_0(\lambda)$。
\item 全重数$\mu_a(\lambda) = e_1 + e_2 + \cdots + e_p = \sum_{k=0}^{d_\lambda - 1} b_k(\lambda)$。
\item 阶$d_\lambda = e_1$。
\end{enumerate}
\begin{proof}
(1) 每个单位标准形式恰好对应$1$维本征空间（轨迹中只有末向量是本征向量），
所以本征空间维数等于单位标准形式的个数，也等于$b_0(\lambda)$。

(2) 全本征空间的维数等于所有单位标准形式的阶数之和。
另一方面，由$b_k(\lambda) = \left|\{1\leqslant i \leqslant p\mid e_i > k\}\right|$，
$$ \sum_{k=0}^{d_\lambda - 1} b_k(\lambda) = \sum_{k=0}^{d_\lambda - 1} \left|\{1\leqslant i \leqslant p\mid e_i > k\}\right| = \sum_{i=1}^p e_i. $$
（每个轨迹长度为$e_i$的单位标准形式，在$k = 0, 1, \cdots, e_i - 1$时各被计数一次。可以理解用两种方法统计方格的总数。）

(3) 极小归零式中$(z - \lambda)$的重数等于最大轨迹长度$e_1$（命题\ref{pp:6_12_5}），即$d_\lambda = e_1$。
\end{proof}

\begin{et}\label{et:6_nilpotent}
考虑$\mathbb{R}^5$上的自同态$M$，
$$M(x_1, x_2, x_3, x_4, x_5) = (x_2 + x_4,\; x_3,\; 0,\; x_5,\; 0).$$
\begin{enumerate}
\item 写出$M$在标准基上的矩阵。
\item 计算$M^2$和$M^3$，验证$M$是幂零自同态，并确定其阶。
\item 求解$\Ker M$、$\Ker M^2$、$\Ker M^3$，写出各自的维数和一组基。验证空间套：
$$\{\mathbf{0}\} \subset \Ker M \subset \Ker M^2 \subset \Ker M^3,$$
并指出空间套首次稳定的阶数。写出各$b_k$的值。
\item 验证$\mathbf{e}_3 \to \mathbf{e}_2 \to \mathbf{e}_1 \to \mathbf{0}$是$M$的一条轨迹，并求其长度。
\item 在$\Ker M^2$中寻找一个不属于$\langle \mathbf{e}_1, \mathbf{e}_2, \mathbf{e}_3\rangle$的向量$\mathbf{v}$，使得$M\mathbf{v} \neq \mathbf{0}$而$M^2\mathbf{v} = \mathbf{0}$。写出$\mathbf{v}$所在的轨迹及其长度。
\item 验证两条轨迹中的向量合起来构成$\mathbb{R}^5$的一组基，并写出对应的两个循环子空间。
\end{enumerate}
\end{et}

\begin{so}
\textbf{(1)} .\;
$M$在标准基上的矩阵为：
$$
M = \begin{bmatrix}
0 & 1 & 0 & 1 & 0 \\
0 & 0 & 1 & 0 & 0 \\
0 & 0 & 0 & 0 & 0 \\
0 & 0 & 0 & 0 & 1 \\
0 & 0 & 0 & 0 & 0
\end{bmatrix}.
$$
\textbf{(2)} .\;
计算$M^2$。按照矩阵乘法：
$$
M^2 = \begin{bmatrix}
0 & 0 & 1 & 0 & 1 \\
0 & 0 & 0 & 0 & 0 \\
0 & 0 & 0 & 0 & 0 \\
0 & 0 & 0 & 0 & 0 \\
0 & 0 & 0 & 0 & 0
\end{bmatrix}.
$$
具体验证第一行：$M$的第一行为$(0,1,0,1,0)$，乘以$M$得到$M$的第二行加第四行：
$(0,0,1,0,0) + (0,0,0,0,1) = (0,0,1,0,1)$。其余各行类似可得。

再计算$M^3 = M^2 \cdot M$。$M^2$仅第一行非零，为$(0,0,1,0,1)$，乘以$M$得到$M$的第三行加第五行：
$(0,0,0,0,0) + (0,0,0,0,0) = (0,0,0,0,0)$。因此：
$$M^3 = \begin{bmatrix} 0 & 0 & 0 & 0 & 0 \\ 0 & 0 & 0 & 0 & 0 \\ 0 & 0 & 0 & 0 & 0 \\ 0 & 0 & 0 & 0 & 0 \\ 0 & 0 & 0 & 0 & 0 \end{bmatrix}.$$
而$M^2 \neq 0$（第一行有非零元素），所以$M$是\textbf{3阶幂零自同态}，极小归零式为$z^3$。

\textbf{(3)} 求$\Ker M$：解$M\mathbf{x} = \mathbf{0}$，即：
$$x_2 + x_4 = 0, \quad x_3 = 0, \quad x_5 = 0.$$
自由变量为$x_1$和$x_4$（取$x_2 = -x_4$）。一组基为：
$$(1,0,0,0,0), \quad (0,-1,0,1,0).$$
$\dim \Ker M = 2$。

求$\Ker M^2$：解$M^2\mathbf{x} = \mathbf{0}$，即：
$$x_3 + x_5 = 0.$$
自由变量为$x_1, x_2, x_4, x_5$（取$x_3 = -x_5$）。一组基为：
$$(1,0,0,0,0), \quad (0,1,0,0,0), \quad (0,0,-1,0,1), \quad (0,0,0,1,0).$$
$\dim \Ker M^2 = 4$。

$\Ker M^3 = \mathbb{R}^5$，$\dim \Ker M^3 = 5$。

空间套为：
$$\{\mathbf{0}\} \subsetneq \Ker M \subsetneq \Ker M^2 \subsetneq \Ker M^3 = \mathbb{R}^5,$$
维数分别为$0, 2, 4, 5$。空间套在$i=3$时首次稳定（$\Ker M^3 = \Ker M^4 = \cdots = \mathbb{R}^5$）。

各$b_k$值为：
\begin{align*}
    b_0 &= \nul M = 2, \\
    b_1 &= \nul M^2 - \nul M = 2, \\
    b_2 &= \nul M^3 - \nul M^2 = 1
\end{align*}

\textbf{(4)} 逐步计算：
$$M(0,0,1,0,0) = (0,1,0,0,0) = \mathbf{e}_2,$$
$$M(0,1,0,0,0) = (1,0,0,0,0) = \mathbf{e}_1,$$
$$M(1,0,0,0,0) = (0,0,0,0,0) = \mathbf{0}.$$
所以$\mathbf{e}_3 \to \mathbf{e}_2 \to \mathbf{e}_1 \to \mathbf{0}$是一条轨迹。\textbf{长度为3}。

\textbf{(5)} 取$\mathbf{v} = (0,0,-1,0,1)$。首先验证$\mathbf{v} \in \Ker M^2$：
$M^2\mathbf{v} = (x_3+x_5, 0, 0, 0, 0) = (-1+1, 0, 0, 0, 0) = \mathbf{0}$。成立。

计算$M\mathbf{v}$：
$$M(0,0,-1,0,1) = (0+0,\; -1,\; 0,\; 1,\; 0) = (0,-1,0,1,0) \neq \mathbf{0}.$$
计算$M^2\mathbf{v}$：
$$M(0,-1,0,1,0) = (-1+1,\; 0,\; 0,\; 0,\; 0) = (0,0,0,0,0) = \mathbf{0}.$$
所以$\mathbf{v}$所在的轨迹为：
$$(0,0,-1,0,1) \longrightarrow (0,-1,0,1,0) \longrightarrow \mathbf{0}.$$
\textbf{长度为2}。

\textbf{(6)} 两条轨迹的向量合起来，作为列向量排成矩阵：
$$\begin{bmatrix} 1 & 0 & 0 & 0 & 0 \\ 0 & 1 & 0 & -1 & 0 \\ 0 & 0 & 1 & 0 & -1 \\ 0 & 0 & 0 & 1 & 0 \\ 0 & 0 & 0 & 0 & 1 \end{bmatrix}.$$
这是上三角矩阵，对角线元素全为$1$，效（行列式）为$1 \neq 0$。因此这$5$个向量直无关，构成$\mathbb{R}^5$的一组基。

对应的两个循环子空间为：
$$C_M(\mathbf{e}_3) = \langle \mathbf{e}_3,\, \mathbf{e}_2,\, \mathbf{e}_1 \rangle = \langle (1,0,0,0,0),\, (0,1,0,0,0),\, (0,0,1,0,0)\rangle,$$
$$C_M(\mathbf{v}) = \langle (0,0,-1,0,1),\, (0,-1,0,1,0)\rangle.$$
两者维数之和为$3+2=5=\dim\mathbb{R}^5$，且交集为$\{\mathbf{0}\}$（因为五个向量直无关），所以：
$$\mathbb{R}^5 = C_M(\mathbf{e}_3) \oplus C_M(\mathbf{v}).$$
\end{so}

% 反之，从$b_k(\lambda)$可以唯一恢复出轨迹长度：

% \begin{pp}\label{pp:6_13_3}
% 序列$b_0(\lambda), b_1(\lambda), \cdots, b_{d_\lambda - 1}(\lambda)$唯一确定了轨迹长度$e_1 \geq e_2 \geq \cdots \geq e_p$。
% \end{pp}
% \begin{proof}
% 轨迹长度恰好等于$j$的个数为：
% $$ b_{j-1}(\lambda) - b_j(\lambda), $$
% 其中约定$b_{d_\lambda}(\lambda) = 0$。
% 这是因为$b_{j-1}$计算的是$e_i > j-1$（即$e_i \geq j$）的个数，
% $b_j$计算的是$e_i > j$（即$e_i \geq j+1$）的个数，
% 两者之差恰好是$e_i = j$的个数。
% \end{proof}

\subsection{标准形式}
考察$N$限制在这个子空间上的部分在循环基上的矩阵：
$[n_{i,j}]_{1 \leqslant i, j \leqslant d_\mathbf{u}}$，它的形式十分简洁：$n_{2,1} = n_{3,2} = \cdots = n_{d_\mathbf{u}, d_\mathbf{u}-1} = 1$，而其他位置上都是0：
$$ \begin{bmatrix}
     0 & 0 & 0 & \cdots & 0 \\
     1 & 0 & 0 & \cdots & 0 \\
     0 & 1 & 0 & \cdots & 0  \\
     \vdots & \vdots & \ddots & \ddots & \vdots  \\
     0 & 0 & \cdots & 1 & 0 \\
    \end{bmatrix} $$
这是一个下三角矩阵。大多数文献把基底的顺序倒过来：$N^{d_\mathbf{u} - 1} \mathbf{u}, N^{d_\mathbf{u} - 2}\mathbf{u}, \cdots , N \mathbf{u}, \mathbf{u}$. 这时矩阵也变成上三角矩阵：
$$ \begin{bmatrix}
     0 & 1 & 0 & \cdots & 0 \\
     0 & 0 & 1 & \cdots & 0 \\
     \vdots & \vdots & \ddots & \ddots & \vdots  \\
     0 & 0 & 0 & \cdots & 1  \\
     0 & 0 & 0 & \cdots & 0 \\
    \end{bmatrix} $$
如果能够把整个空间都表示成这种形式，我们就给出了幂零自同态的一种优美的表示。为此，我们希望把整个空间表示成几个循环不变子空间的直和，或者说找到几组轨迹，并且把它们拼接成整个空间的基底。事实证明我们确实能做到这一点（证明参见附录三\ref{c-20}）：
\begin{pp}\label{pp:6_12_4}
    给定空间$\mathbb{V}$上幂零自同态$N$，空间可以表示为：
    $$ \mathbb{V} = C_N(\mathbf{u}_1) \oplus C_N(\mathbf{u}_2) \oplus \cdots \oplus C_N(\mathbf{u}_p),$$
    其中$p$为正整数，$\mathbf{u}_1, \mathbf{u}_2, \cdots , \mathbf{u}_p$为$N$决定的向量。
\end{pp}

把空间表示成循环不变子空间的直和，把对应的轨迹作为基底，我们可以将幂零自同态$N$表示为分块的单位标准形式。
回到广义本征空间$H_T(\lambda)$，我们知道$T-\lambda \mathrm{I}$限制在$H_T(\lambda)$上就是幂零自同态，
而$\lambda \mathrm{I}$的矩阵形式在任何基上都是一样的：对角线元素全为$\lambda$的对角矩阵。
所以，在上述基上，$T$在每个循环子空间$C_{T-\lambda \mathrm{I}}(\mathbf{u}_i)$里的矩阵形式是：
$$
J_i(\lambda) = \begin{bmatrix}
     \lambda & 1 & 0 & \cdots & 0 \\
     0 & \lambda & 1 & \cdots & 0 \\
     \vdots & \vdots & \ddots & \ddots & \vdots  \\
     0 & 0 & 0 & \cdots  & 1  \\
     0 & 0 & 0 & \cdots & \lambda \\
    \end{bmatrix}
$$
而$T$在$H_T(\lambda)$里的矩阵形式是：
$$
J(\lambda) = \begin{bmatrix}
     J_1(\lambda) & 0 & \cdots & 0 \\
     0 & J_2(\lambda) & \cdots & 0 \\
     \vdots & \vdots & \ddots  & \vdots  \\
     0 & 0 & \cdots  & J_p(\lambda)  \\
    \end{bmatrix}
$$
我们再次得到一个分块对角矩阵，而这一次我们对每个“分块”的内容也有确定了解。
我们把$J_i(\lambda)$这样的矩阵形式称为自同态的\textbf{单位标准形式}，把形同$J(\lambda)$的矩阵称为自同态的\textbf{标准形式}。
单位标准形式的对角线元素为同一个数，第$i$行第$i-1$列值为1，其余位置都是0。标准形式是由单位标准形式的矩阵构成的分块对角矩阵，
其对角线元素可以不为0，不同的块里对角线元素可以是不同数值，第$i$行第$i-1$列值为0或1，其余位置都是0。
幂零自同态的本征值是0，所以这里的标准形式和前面定义的是一致的。

我们来把$B$和$C$化为标准形式。

\paragraph{自同态$B$的标准形式}

$B$的唯一本征值为$1$，极小归零式为$(z-1)^2$，全重数$\mu_a(1) = 2$。
$N_B = B - \mathrm{I}$是$2$阶幂零，只有一条轨迹：
$$
(1,0) \;\longrightarrow\; (2,1) \;\longrightarrow\; \mathbf{0}.
$$
取轨迹基为$((2,1),\, (1,0))$（结束向量在前，出发向量在后）。
在这组基上：
$$
B(2,1) = 1\cdot(2,1) + 0\cdot(1,0), \quad B(1,0) = 1\cdot(2,1) + 1\cdot(1,0).
$$
所以$B$在这组基上的矩阵为：
$$
[B] = \begin{bmatrix} 1 & 1 \\ 0 & 1 \end{bmatrix} = J_2(1).
$$
这就是$B$的标准形式：一个$2\times 2$的单位标准形式，对应本征值$1$。

\paragraph{自同态$C$的标准形式}

$C$有两个本征值：$5$（全重数$1$）和$3$（全重数$2$）。
极小归零式为$(z-3)^2(z-5)$。

对本征值$5$：全本征空间$H_C(5) = \mathbb{R}(2,1,0)$是$1$维的。
$C$在其上的限制就是乘以$5$。对应的单位标准形式为$J_1(5) = [5]$。

对本征值$3$：全本征空间$H_C(3) = \Ker(C-3\mathrm{I})^2$是$2$维的。
$N_C = C - 3\mathrm{I}$在其上是$2$阶幂零，轨迹为：
$$
(-1,0,2) \;\longrightarrow\; (0,3,0) \;\longrightarrow\; \mathbf{0}.
$$
取轨迹基为$((0,3,0),\, (-1,0,2))$。在这组基上：
$$
C(0,3,0) = 3\cdot(0,3,0) + 0\cdot(-1,0,2),
$$
$$
C(-1,0,2) = 1\cdot(0,3,0) + 3\cdot(-1,0,2).
$$
验证第二个等式：$C(-1,0,2) = (5\cdot(-1)+0+2,\, 1\cdot(-1)+0+4,\, 0) = (-3, 3, 0)$。
而$(0,3,0) + 3(-1,0,2) = (-3, 3, 0)$。成立。

所以$C$在$H_C(3)$的这组基上的矩阵为：
$$
\begin{bmatrix} 3 & 1 \\ 0 & 3 \end{bmatrix} = J_2(3).
$$

综合两个本征值，取$\mathbb{R}^3$的基底为：
$$
\mathcal{B} = \big((0,3,0),\; (-1,0,2),\; (2,1,0)\big).
$$
前两个向量是$H_C(3)$的轨迹基，第三个是$H_C(5)$的基向量。
$C$在这组基上的矩阵为分块对角矩阵：
$$
[C]_{\mathcal{B}} = \begin{bmatrix} 3 & 1 & 0 \\ 0 & 3 & 0 \\ 0 & 0 & 5 \end{bmatrix} = J_2(3) \oplus J_1(5).
$$
这就是$C$的标准形式。

\paragraph{自同态$A$的标准形式}

$A$有两个不同的本征值$1$和$2$，各自全重数为$1$。
$A$可对角化，标准形式就是对角矩阵：
$$
[A] = \begin{bmatrix} 1 & 0 \\ 0 & 2 \end{bmatrix} = J_1(1) \oplus J_1(2).
$$
每个单位标准形式都是$1\times 1$的，没有非对角的$1$。
这对应极小归零式$(z-1)(z-2)$没有重根的事实。

至此，我们对自同态在某个本征值相关的空间里的性质有了深刻的了解。给定本征值$\lambda$，我们可以找到一组基底，
使得自同态在它对应的全本征空间里有唯一的形式。

如果自同态$T$的极小归零式是一次式的乘积，那么可以表示成：
$$ m_T(z) = (z - \lambda_1)^{d_1} (z - \lambda_2)^{d_2} \cdots  (z - \lambda_k)^{d_k}$$
其中$\lambda_1, \lambda_2,\cdots$就是$T$的不同的本征值。零空间引理告诉我们：
$$ \mathbb{V} = \Ker (T - \lambda_1\mathrm{I})^{d_1} \oplus \Ker (T - \lambda_2\mathrm{I})^{d_2} \oplus \cdots \oplus \Ker (T - \lambda_k\mathrm{I})^{d_k}.$$
对任意$i$来说，$\Ker (T - \lambda_i\mathrm{I})^{d_i}$是$T$-不变子空间，
且$T - \lambda_i\mathrm{I}$在其中是幂零自同态。所以，可以选择合适的基底，
使$T$在$\Ker (T - \lambda_i\mathrm{I})^{d_i}$中的部分可以表示为标准形式的矩阵。
从整个空间来说，也是如此。

综上所述，如果自同态$T$的极小归零式是一次式的乘积，$T$不仅可以三角化，而且可以约简为标准形式。
像我们之前提到的复系数有限维平直空间，比如$\mathbb{C}^n$中，任何自同态都可以约简为标准形式。标准形式下的矩阵是分块对角矩阵，对角线由若干个标准块$J_i(\lambda)$构成，分别对应不同的本征值。每个标准块都是一个方阵，它的维数等于本征值的全重数。

标准形式可以帮助我们更好地理解自同态。比如，我们定义过自同态的迹是它在任意基底下的矩阵的对角元素之和。如果自同态可以约简为标准形式，那么其对角线上的元素就是各个$J_i(\lambda)$的对角线元素，分别是不同的本征值$\lambda$。于是我们得到迹和本征值的关系：
\begin{tm}
    如果自同态可约简为标准形式，那么它的迹就是它的本征值按照全重数相加的和。
\end{tm}
然而自同态的迹是它在任意基底下的矩阵的对角元素之和，因此必然在系数域里。这说明，只要自同态可约简为标准形式，它的本征值按照全重数相加的和必然在系数域里。

另外我们在证明命题\ref{pp:6_8}时定义过本征式：
$$
c_T(x) = \prod_{j=1}^n(x - a_{jj}).
$$
它是关于自同态矩阵的$n$次整式，是三角化后矩阵的对角元素对应的一次式的乘积。将自同态化为标准形式后，$c_T$就是所有本征值按照全重数对应的一次式的乘积：
$$
c_T(x) = \prod_{\lambda}(x - \lambda)^{\mu_{a}(\lambda)}.
$$
本征值的全重数大于等于其阶，因此$c_T$是$m_T$的倍式，从而是归零式。我们把这个结论称为\textbf{全本征归零定理}。
\begin{tm}{\textbf{全本征归零定理}}
    对可三角化的自同态$T$，有：
    $$ c_T(T) = 0. $$
\end{tm}


\begin{et}\label{et:6_canonical}
继续考虑上一题中的幂零自同态$M$：
$$
M = \begin{bmatrix}
0 & 1 & 0 & 1 & 0 \\
0 & 0 & 1 & 0 & 0 \\
0 & 0 & 0 & 0 & 0 \\
0 & 0 & 0 & 0 & 1 \\
0 & 0 & 0 & 0 & 0
\end{bmatrix}.
$$
上一题中已经验证，$M$是$3$阶幂零自同态，有两条轨迹：
\begin{align*}
    \mathbf{e}_3 &\longrightarrow \mathbf{e}_2 \longrightarrow \mathbf{e}_1 \longrightarrow \mathbf{0},\\
    \mathbf{f}_2 = (0,0,-1,0,1) &\longrightarrow \mathbf{f}_1 = (0,-1,0,1,0) \longrightarrow \mathbf{0}.
\end{align*}
\begin{enumerate}
\item 取轨迹基：
$$\mathcal{B} = \big(\mathbf{e}_1,\; \mathbf{e}_2,\; \mathbf{e}_3,\; \mathbf{f}_1,\; \mathbf{f}_2\big).$$
逐一计算$M$作用在每个基向量上的结果，将其用$\mathcal{B}$分解，写出$M$在$\mathcal{B}$上的矩阵$[M]_{\mathcal{B}}$。
\item 验证$[M]_{\mathcal{B}}$是$M$的标准形式，并指出它由哪些单位标准形式构成。
\item 写出从标准基到$\mathcal{B}$的基变换矩阵$P$，并验证$P^{-1}MP = [M]_{\mathcal{B}}$。
\item 写出$M$的极小归零式和本征式$c_M(x)$，验证$c_M(M) = 0$（即全本征归零定理在此例中成立）。
\end{enumerate}
\end{et}

\begin{so}
\textbf{(1)} 逐一计算$M$作用在$\mathcal{B}$中每个基向量上的结果。
\begin{align*}
    M\mathbf{e}_1 &= M(1,0,0,0,0) \phantom{-}= (0,0,0,0,0) \phantom{-}= 0\cdot\mathbf{e}_1 + 0\cdot\mathbf{e}_2 + 0\cdot\mathbf{e}_3 + 0\cdot\mathbf{f}_1 + 0\cdot\mathbf{f}_2.\\
    M\mathbf{e}_2 &= M(0,1,0,0,0) \phantom{-}= (1,0,0,0,0) \phantom{-}= 1\cdot\mathbf{e}_1 + 0\cdot\mathbf{e}_2 + 0\cdot\mathbf{e}_3 + 0\cdot\mathbf{f}_1 + 0\cdot\mathbf{f}_2.\\
    M\mathbf{e}_3 &= M(0,0,1,0,0) \phantom{-}= (0,1,0,0,0) \phantom{-}= 0\cdot\mathbf{e}_1 + 1\cdot\mathbf{e}_2 + 0\cdot\mathbf{e}_3 + 0\cdot\mathbf{f}_1 + 0\cdot\mathbf{f}_2.\\
    M\mathbf{f}_1 &= M(0,-1,0,1,0) = (0,0,0,0,0) \phantom{-}= 0\cdot\mathbf{e}_1 + 0\cdot\mathbf{e}_2 + 0\cdot\mathbf{e}_3 + 0\cdot\mathbf{f}_1 + 0\cdot\mathbf{f}_2.\\
    M\mathbf{f}_2 &= M(0,0,-1,0,1) = (0,-1,0,1,0) = 0\cdot\mathbf{e}_1 + 0\cdot\mathbf{e}_2 + 0\cdot\mathbf{e}_3 + 1\cdot\mathbf{f}_1 + 0\cdot\mathbf{f}_2.
\end{align*}

因此，$M$在$\mathcal{B}$上的矩阵为：
$$
[M]_{\mathcal{B}} = \begin{bmatrix}
0 & 1 & 0 & 0 & 0 \\
0 & 0 & 1 & 0 & 0 \\
0 & 0 & 0 & 0 & 0 \\
0 & 0 & 0 & 0 & 1 \\
0 & 0 & 0 & 0 & 0
\end{bmatrix}.
$$

\textbf{(2)} 上式可以写为分块对角形式：
$$
[M]_{\mathcal{B}} = \begin{bmatrix}
0 & 1 & 0 \\
0 & 0 & 1 \\
0 & 0 & 0
\end{bmatrix} \oplus \begin{bmatrix}
0 & 1 \\
0 & 0
\end{bmatrix} = J_3(0) \oplus J_2(0).
$$
这正是$M$的\textbf{标准形式}，由一个$3$阶单位标准形式$J_3(0)$和一个$2$阶单位标准形式$J_2(0)$构成。
$J_3(0)$对应第一条轨迹（长度$3$），$J_2(0)$对应第二条轨迹（长度$2$）。

\textbf{(3)} 基变换矩阵$P$的列为$\mathcal{B}$中各向量在标准基上的坐标：
$$
P = \begin{bmatrix}
1 & 0 & 0 & 0 & 0 \\
0 & 1 & 0 & -1 & 0 \\
0 & 0 & 1 & 0 & -1 \\
0 & 0 & 0 & 1 & 0 \\
0 & 0 & 0 & 0 & 1
\end{bmatrix}.
$$
$P$是上三角矩阵，对角线元素全为$1$，所以$P$可逆。其逆矩阵为：
$$
P^{-1} = \begin{bmatrix}
1 & 0 & 0 & 0 & 0 \\
0 & 1 & 0 & 1 & 0 \\
0 & 0 & 1 & 0 & 1 \\
0 & 0 & 0 & 1 & 0 \\
0 & 0 & 0 & 0 & 1
\end{bmatrix}.
$$
（验证：$P$中$(2,4)$位置为$-1$，$(3,5)$位置为$-1$；$P^{-1}$中对应位置为$+1$，其余不变。读者可以直接计算$PP^{-1}=\mathrm{I}_5$确认。）

下面验证$P^{-1}MP = [M]_{\mathcal{B}}$。首先计算$MP$：
$$
MP = \begin{bmatrix}
0 & 1 & 0 & 0 & 0 \\
0 & 0 & 1 & 0 & -1 \\
0 & 0 & 0 & 0 & 0 \\
0 & 0 & 0 & 0 & 1 \\
0 & 0 & 0 & 0 & 0
\end{bmatrix}.
$$

再计算$P^{-1}(MP)$：
$$
P^{-1}MP = \begin{bmatrix}
0 & 1 & 0 & 0 & 0 \\
0 & 0 & 1 & 0 & 0 \\
0 & 0 & 0 & 0 & 0 \\
0 & 0 & 0 & 0 & 1 \\
0 & 0 & 0 & 0 & 0
\end{bmatrix} = [M]_{\mathcal{B}}. 
$$

\textbf{(4)} $M$是$3$阶幂零自同态，唯一本征值为$0$，全重数为$5$，阶为$3$。

极小归零式为$m_M(x) = x^3$（次数等于最大轨迹长度$e_1 = 3$）。

本征式为所有本征值按全重数对应的一次式的乘积：
$$c_M(x) = (x - 0)^5 = x^5.$$

验证全本征归零定理：$m_M(x) = x^3$确实是$c_M(x) = x^5$的因式，$c_M(M) = M^5$。由于$M^3 = 0$，
$$M^5 = M^3 \cdot M^2 = 0 \cdot M^2 = 0.$$
所以$c_M(M) = 0$，全本征归零定理成立。

另外可以验证$M$的效为$\det(M) = 0$（等于所有本征值按全重数的乘积$0^5 = 0$），迹为$\mathrm{tr}(M) = 0$（等于所有本征值按全重数的和$0 \times 5 = 0$）。
\end{so}

\chapter{自同态的效}
以上各章中，我们对有限维空间中自同态的性质进行了初步探讨，将一般自同态表示成更简单、更好理解的形式。我们使用矩阵来记录自同态在特定基底上的形态。有了矩阵，我们可以量化地讨论自同态的性质。通过讨论自同态的本征值和本征向量，我们证明了在某些特定基底上，自同态能够用形式比较简单的矩阵表示。本章我们探讨自同态的另一个重要本征量。

考虑平面上两个向量$\mathbf{u}, \mathbf{v}$张成的平行四边形。它的面积如何计算呢？

首先定义单位面积。给定平面基底$\mathbf{e}_1, \mathbf{e}_2$，定义$\mathbf{e}_1, \mathbf{e}_2$张成的平行四边形面积为$1$。按照我们对面积的理解，如果$\mathbf{u}, \mathbf{v}$分别是$\mathbf{e}_1, \mathbf{e}_2$的$u,v$倍，那么$\mathbf{u}, \mathbf{v}$张成的平行四边形的面积应该是单位面积的$uv$倍，即面积为$uv$。其次，如果$\mathbf{u}_1, \mathbf{v}$和$\mathbf{u}_2, \mathbf{v}$张成的平行四边形面积为$S_1, S_2$，那么$\mathbf{u}_1+\mathbf{u}_2, \mathbf{v}$张成的平行四边形面积应该是$S_1+S_2$。最后，如果$\mathbf{u}, \mathbf{v}$共轴，那么张成的是退化的平行四边形，面积为$0$。

按照以上性质，我们可以算出任意向量$\mathbf{u}, \mathbf{v}$张成的平行四边形面积。把$\mathbf{u}, \mathbf{v}$分别表示成基向量的直组合：
$$\mathbf{u} = u_1\mathbf{e}_1 + u_2\mathbf{e}_2, \quad \mathbf{v} = v_1\mathbf{e}_1 + v_2 \mathbf{e}_2. $$
那么，张成的平行四边形面积$S(\mathbf{u}, \mathbf{v})$就是：
\begin{align*}
    S(\mathbf{u}, \mathbf{v}) &= S(u_1\mathbf{e}_1 + u_2\mathbf{e}_2, \; v_1\mathbf{e}_1 + v_2 \mathbf{e}_2) \\
    &= S(u_1\mathbf{e}_1, \; v_1\mathbf{e}_1 + v_2 \mathbf{e}_2) + S(u_2\mathbf{e}_2, \; v_1\mathbf{e}_1 + v_2 \mathbf{e}_2) \\
    &= S(u_1\mathbf{e}_1, \; v_1\mathbf{e}_1) + S(u_1\mathbf{e}_1, \; v_2 \mathbf{e}_2) + S(u_2\mathbf{e}_2, \; v_1\mathbf{e}_1) + S(u_2\mathbf{e}_2, \; v_2 \mathbf{e}_2) \\
    &= u_1v_2 S(\mathbf{e}_1, \;\mathbf{e}_2) + u_2v_1 S(\mathbf{e}_2, \;\mathbf{e}_1).
\end{align*}
其中$S(\mathbf{e}_1, \;\mathbf{e}_2)$和$S(\mathbf{e}_2, \;\mathbf{e}_1)$分别是$\mathbf{e}_1, \mathbf{e}_2$张成的平行四边形与$\mathbf{e}_2, \mathbf{e}_1$张成的平行四边形的面积。直观上来看这是平面上两个一模一样的图形。但从向量的观点来看，为了方便统一计算各种平行四边形的面积，我们需要规定它们面积相反。如果认定$S(\mathbf{e}_1, \;\mathbf{e}_2) = 1$，那么$S(\mathbf{e}_2, \;\mathbf{e}_1)$就等于$-1$。实际上，这个性质可由前面“认定”的面积“公理”推出：
\begin{align*}
    0 &= S(\mathbf{e}_1 + \mathbf{e}_2, \;\mathbf{e}_1 + \mathbf{e}_2) \\
    &= S(\mathbf{e}_1, \;\mathbf{e}_1) + S(\mathbf{e}_1, \;\mathbf{e}_2) + S(\mathbf{e}_2, \;\mathbf{e}_1) + S(\mathbf{e}_2, \;\mathbf{e}_2) \\
    &= S(\mathbf{e}_1, \;\mathbf{e}_2) + S(\mathbf{e}_2, \;\mathbf{e}_1).
\end{align*}
于是有
\begin{align*}
    S(\mathbf{u}, \mathbf{v}) &= u_1v_2 S(\mathbf{e}_1, \;\mathbf{e}_2) + u_2v_1 S(\mathbf{e}_2, \;\mathbf{e}_1)\\
    &= u_1v_2 - u_2v_1.
\end{align*}
这样定义的面积称为\textbf{有向面积}。研究把两个平面向量映射到有向面积的映射$S$，它满足前面提到的基本性质：
\begin{enumerate}
    \item 双直性：$S$是关于任一变量平直的双变量映射。
    \item 反称性：$\forall \,\mathbf{u}, \, \in \, \mathbb{V}, \;\;S(\mathbf{u}, \mathbf{u}) = 0.$
    \item 单一性：$S(\mathbf{e}_1, \;\mathbf{e}_2) = 1.$
\end{enumerate}
而$S(\mathbf{u}, \mathbf{v})$的绝对值就是我们平常说的平行四边形的面积。

立体空间中，要计算三个向量形成的平行六面体的体积，也有以上的向量方法。首先按照同样的方法定义有向体积，然后求绝对值得到一般说的体积。具体来说，给定三维空间中基底$\mathbf{e}_1, \mathbf{e}_2, \mathbf{e}_3$，定义满足有向体积映射$V$满足：
\begin{enumerate}
    \item $\forall \, \mathbf{u}_1, \mathbf{u}_2, \mathbf{v}, \mathbf{w}\,\in\,\mathbb{V}$，$\forall \,\lambda,\mu\,\in\mathbb{F}$，都有$S(\lambda\mathbf{u}_1+\mu\mathbf{u}_2, \mathbf{v},\mathbf{w}) = \lambda S(\mathbf{u}_1, \mathbf{v},\mathbf{w}) + \mu S(\mathbf{u}_2, \mathbf{v},\mathbf{w})$
    \item $\forall \, \mathbf{u},\mathbf{v}_1, \mathbf{v}_2, \mathbf{w}\,\in\,\mathbb{V}$，$\forall \,\lambda,\mu\,\in\mathbb{F}$，都有$S(\mathbf{u}, \lambda\mathbf{v}_1+\mu\mathbf{v}_2, \mathbf{w}) = \lambda S(\mathbf{u}, \mathbf{v}_1, \mathbf{w}) + \mu S(\mathbf{u}, \mathbf{v}_2,\mathbf{w})$.
    \item $\forall \, \mathbf{u},\mathbf{v}, \mathbf{w}_1, \mathbf{w}_2\,\in\,\mathbb{V}$，$\forall \,\lambda,\mu\,\in\mathbb{F}$，都有$S(\mathbf{u}, \mathbf{v}, \lambda\mathbf{w}_1+\mu\mathbf{w}_2) = \lambda S(\mathbf{u}, \mathbf{v}, \mathbf{w}_1) + \mu S(\mathbf{u}, \mathbf{v},\mathbf{w}_2)$.
    \item $\forall \, \mathbf{u},\mathbf{v}, \mathbf{w}\,\in\,\mathbb{V}$，如果其中两个相等，就有$S(\mathbf{u}, \mathbf{v},\mathbf{w}) = 0.$
    \item $S(\mathbf{e}_1, \;\mathbf{e}_2,\;\mathbf{e}_3) = 1.$
\end{enumerate}
并规定$V(\mathbf{e}_1, \mathbf{e}_2, \mathbf{e}_3) = 1$。

\begin{wrapfigure}[6]{r}{0.42\textwidth} %this figure will be at the right
    \vspace{-40pt}
    \flushright
    \includegraphics[width=0.4\textwidth]{tu/平行六面体1.png}
\end{wrapfigure}

于是，给定三个向量$\mathbf{u}, \mathbf{v},\mathbf{w}$，同样可以通过在基底上分解而得到：
\begin{align*}
    S(\mathbf{u}, \mathbf{v},\mathbf{w}) &= S(u_1\mathbf{e}_1 + u_2\mathbf{e}_2+u_3\mathbf{e}_3, \; v_1\mathbf{e}_1 + v_2 \mathbf{e}_2+v_3\mathbf{e}_3, \;w_1\mathbf{e}_1 + w_2\mathbf{e}_2 + w_3\mathbf{e}_3) \\
    &= u_1v_2w_3 S(\mathbf{e}_1, \;\mathbf{e}_2, \;\mathbf{e}_3) + u_2v_1w_3 S(\mathbf{e}_2, \;\mathbf{e}_1,\;\mathbf{e}_3)+ u_1v_3w_2 S(\mathbf{e}_1, \;\mathbf{e}_3, \;\mathbf{e}_2) \\
    &\;\; + u_3v_1w_2 S(\mathbf{e}_3, \;\mathbf{e}_1,\;\mathbf{e}_2) + u_2v_3w_1 S(\mathbf{e}_2, \;\mathbf{e}_3, \;\mathbf{e}_1) + u_3v_2w_1 S(\mathbf{e}_3, \;\mathbf{e}_2,\;\mathbf{e}_1) .
\end{align*}
和平面情形的区别是，我们要处理$6$个看起来都是单位平行六面体的有向体积。不过我们可以利用和之前一样的方法讨论：
\begin{align*}
    0 &= S(\mathbf{e}_1 + \mathbf{e}_2, \;\mathbf{e}_1 + \mathbf{e}_2,\; \mathbf{e}_3) \\
    &= S(\mathbf{e}_1, \;\mathbf{e}_1,\; \mathbf{e}_3) + S(\mathbf{e}_1, \;\mathbf{e}_2,\; \mathbf{e}_3) + S(\mathbf{e}_2, \;\mathbf{e}_1,\; \mathbf{e}_3) + S(\mathbf{e}_2, \;\mathbf{e}_2,\; \mathbf{e}_3) \\
    &= S(\mathbf{e}_1, \;\mathbf{e}_2,\; \mathbf{e}_3) + S(\mathbf{e}_2, \;\mathbf{e}_1,\; \mathbf{e}_3).
\end{align*}
因此$S(\mathbf{e}_2, \;\mathbf{e}_1,\; \mathbf{e}_3) = -S(\mathbf{e}_1, \;\mathbf{e}_2,\; \mathbf{e}_3)$。也就是说，交换$\mathbf{e}_1,\mathbf{e}_2$的位置，有向体积就变为相反数。这个结论显然也适用于其他位置。因此我们得到：
\begin{align*}
    S(\mathbf{e}_2, \;\mathbf{e}_1,\; \mathbf{e}_3) &= -S(\mathbf{e}_1, \;\mathbf{e}_2,\; \mathbf{e}_3) = -1,\\
    S(\mathbf{e}_1, \;\mathbf{e}_3,\; \mathbf{e}_2) &= -S(\mathbf{e}_1, \;\mathbf{e}_2,\; \mathbf{e}_3) = -1,\\
    S(\mathbf{e}_3, \;\mathbf{e}_2,\; \mathbf{e}_1) &= -S(\mathbf{e}_1, \;\mathbf{e}_2,\; \mathbf{e}_3) = -1, \\
    S(\mathbf{e}_2, \;\mathbf{e}_3,\; \mathbf{e}_1) &= -S(\mathbf{e}_2, \;\mathbf{e}_1,\; \mathbf{e}_3) = 1,\\
    S(\mathbf{e}_3, \;\mathbf{e}_1,\; \mathbf{e}_2) &= -S(\mathbf{e}_3, \;\mathbf{e}_2,\; \mathbf{e}_1) = 1.
\end{align*}
代入前式得到：
\begin{align*}
    S(\mathbf{u}, \mathbf{v},\mathbf{w}) &= S(u_1\mathbf{e}_1 + u_2\mathbf{e}_2+u_3\mathbf{e}_3, \; v_1\mathbf{e}_1 + v_2 \mathbf{e}_2+v_3\mathbf{e}_3, \;w_1\mathbf{e}_1 + w_2\mathbf{e}_2 + w_3\mathbf{e}_3) \\
    &= u_1v_2w_3 + u_3v_1w_2 + u_2v_3w_1 - u_2v_1w_3 - u_1v_3w_2 - u_3v_2w_1  .
\end{align*}
这就是立体空间中三个向量张成的平行六面体的有向体积。它的绝对值就是我们一般说的体积。

在一般的$n$维平直空间$\mathbb{V}$中，给定基底$\mathcal{B} = \mathbf{e}_1,\mathbf{e}_2,\cdots,\mathbf{e}_n$，我们也可以给$n$个向量定义“有向体积”$\varPsi$：
\begin{align*}
    \varPsi : \;\mathbb{V}^n &\rightarrow \mathbb{F}  
\end{align*}
它把$n$个向量映射到系数域$\mathbb{F}$，并满足：
\begin{enumerate}
    \item $n$重直性：固定其余位置变量，向量组的效对每个位置上的变量来说都是直映射。
    $$\varPsi(\mathbf{v}_1, \cdots, a\mathbf{u} + b\mathbf{w}, \cdots \mathbf{v}_n) = a\varPsi(\mathbf{v}_1, \cdots, \mathbf{u}, \cdots \mathbf{v}_n) + b\varPsi(\mathbf{v}_1, \cdots, \mathbf{w}, \cdots \mathbf{v}_n).$$
    \item 反称性：如果$n$个向量中有两个向量相同，那么向量组的效为$0$。
    $$ \varPsi(\mathbf{v}_1, \cdots, \mathbf{u}, \cdots, \mathbf{u}, \cdots \mathbf{v}_n) = 0. $$
    \item 单一性：$\varPsi(\mathbf{e}_1, \mathbf{e}_2,\cdots,\mathbf{e}_n) = 1.$
\end{enumerate}
我们说这样定义的有向体积，就是$\mathbb{V}$中$n$个向量张成的有向平行体关于由基底张成的单位平行体的倍数。

而对于直自同态$T$来说，给定基底$\mathbf{e}_1,\mathbf{e}_2, \cdots, \mathbf{e}_n$后，自同态$T$由它在基底上的像决定。我们说$\varPsi(T\mathbf{e}_1, \,T\mathbf{e}_2,\cdots,\,T\mathbf{e}_n)$就反映了自同态把单位平行体转为了特定的倍数。这个倍数可以看作自同态对空间尺度的扩张或收缩效果。因此我们称它为自同态$T$（在基底$\mathcal{B}$上）关于$\varPsi$的\textbf{效}。而$\varPsi(\mathbf{u}_1, \mathbf{u}_2,\cdots,\mathbf{u}_n)$也称为向量组$\mathbf{u}_1, \mathbf{u}_2,\cdots,\mathbf{u}_n$在基底$\mathcal{B}$上的\textbf{效}。

下面我们来计算效的表达式。仿照前面在平面和立体空间中的方法，按照定义，我们可以把$n$个向量$\mathbf{u}_1, \mathbf{u}_2,\cdots,\mathbf{u}_n$按照基底分解，然后使用$n$重直性展开，得到$n^n$项，每一项的形式为：
$$ u_{i_1,1}u_{i_2,2}\cdots u_{i_n,n}\varPsi(\mathbf{e}_{i_1}, \mathbf{e}_{i_2},\cdots,\mathbf{e}_{i_n}).$$
其中$i_1,i_2,\cdots,i_n$是$1$到$n$之间的整数，$u_{i,j}$表示向量$\mathbf{u}_j$在基底上分解的第$i$个分量。

由于反称性，如果某项的$i_1,i_2,\cdots,i_n$有两个相同，该项就等于$0$。因此，展开化简后可得：
$$
\varPsi(\mathbf{u}_1, \mathbf{u}_2,\cdots,\mathbf{u}_n) = \sum_{i_1,i_2,\cdots,i_n\;\mbox{两两不等}} u_{i_1,1}u_{i_2,2}\cdots u_{i_n,n}\varPsi(\mathbf{e}_{i_1}, \mathbf{e}_{i_2},\cdots,\mathbf{e}_{i_n}).
$$
而其中$\varPsi(\mathbf{e}_{i_1}, \mathbf{e}_{i_2},\cdots,\mathbf{e}_{i_n})$的值应该总是$1$或$-1$。具体如何确定呢？根据前面的经验可以证明，如果交换两个基向量的位置，$\varPsi(\mathbf{e}_{i_1}, \mathbf{e}_{i_2},\cdots,\mathbf{e}_{i_n})$的值就会变为相反数。所以，我们要深入研究$i_1,i_2,\cdots,i_n$的顺序性质。

考虑$1,2,\cdots, n$到$i_1,i_2,\cdots,i_n$的映射$\sigma$。由于$i_1,i_2,\cdots,i_n$两两不等，所以它是单射。但$i_1,i_2,\cdots,i_n$都属于$1,2,\cdots,n$，所以$\sigma$必然也是满射。也就是说，$i_1,i_2,\cdots,i_n$是$1,2,\cdots,n$到自身的双射，$1,2,\cdots,n$中每个数在$i_1,i_2,\cdots,i_n$里恰好出现一次。从另一个角度来看，$i_1,i_2,\cdots,i_n$就是把$1,2,\cdots,n$里的数重新排列得到的，只有位置换了。我们把这样的映射$\sigma$称为\textbf{置换}。下面我们来研究置换的性质。

\section{置换群}

一般来说，置换指一种特殊的映射。
\begin{df}\label{df:10_1}
    给定集合$X$，从$X$到$X$的双射称为（$X$上的）\textbf{置换}。
\end{df}
顾名思义，置换就是把元素的位置换了。它不会把不同的元素映射到同一个元素，也不会遗漏任一个元素。
集合$X$上所有置换的集合关于映射的复合构成\textbf{群}（参见附录一定义\ref{df:0_1}），称为$X$上的\textbf{置换群}或\textbf{对称群}，一般记为$\mathfrak{S}_X$。
举例来说，设集合$X = \{1, 2, 3\}$，我们可以列举$X$上所有置换：
\begin{align*}
    \begin{pmatrix}
    1 & 2 & 3 \\ 1 & 2 & 3
    \end{pmatrix}, \,\,\,\begin{pmatrix}
        1 & 2 & 3 \\ 1 & 3 & 2
    \end{pmatrix}, \,\,\,\begin{pmatrix}
        1 & 2 & 3 \\ 2 & 1 & 3
    \end{pmatrix},  \\
    \begin{pmatrix}
        1 & 2 & 3 \\ 2 & 3 & 1
    \end{pmatrix}, \,\,\,\begin{pmatrix}
        1 & 2 & 3 \\ 3 & 1 & 2
    \end{pmatrix}, \,\,\,\begin{pmatrix}
        1 & 2 & 3 \\ 3 & 2 & 1
    \end{pmatrix}.  
\end{align*}
每个括号的内容表示一个置换。比如，第一行左起第二个括号中的内容表示把$1$映射到$1$，
把$2$映射到$3$，把$3$映射到$2$的置换。可以发现，这些置换，就对应了我们分解立体空间的效时得到的各项里的下标。

一般把集合$\{1,2,\ldots, n\}$上的置换群记为$\mathfrak{S}_n$。
容易证明，如果$X$有$n$个元素，那么$X$上置换的个数等于
$n$个元素的全排列个数，即$n$的阶乘：$n! = n\times (n-1)\times \cdots\times 2 \times  1$，
且$\mathfrak{S}_X$同构于$\mathfrak{S}_n$。因此，集合元素个数有限的时候，依它的元素个数$n$，我们只需研究$\mathfrak{S}_n$，就掌握了$\mathfrak{S}_X$的性质。

% 上一章中我们讨论了自同态极小归零式的根置换。它把极小归零式的根映射到另一个根。一般来说，我们可以对任何域扩张定义类似的置换：
% \begin{df}\label{df:10_2}
%     设有域扩张$\mathbb{F}\subseteq \mathbb{L}$，$\mathbb{L}$中所有在$\mathbb{F}$上为恒等映射的域自同构（称为$\mathbb{F}$-\textbf{不变自同构}）关于映射的复合构成群，
%     称为域扩张$\mathbb{F}\subseteq \mathbb{L}$的\textbf{自同构群}或$\mathbb{L}$的$\mathbb{F}$-\textbf{不变自同构群}，记作$\mathfrak{g}_\mathbb{F}(\mathbb{L})$。
% \end{df}
% 显然，$\mathbb{F} = \mathbb{L}$的时候，$\mathfrak{g}_\mathbb{F}(\mathbb{L})$是只有空元的群$\{\mathrm{I}_\mathbb{L}\}$。域扩张$\mathbb{R}\subseteq\mathbb{C}$是
% 二维扩张，它的自同构群中元素只有$\mathbb{C}$上的恒等映射$\mathrm{I}_\mathbb{C}$和共轭映射$\mathfrak{c}$：
% $$\mathfrak{g}_\mathbb{F}(\mathbb{L}) = \{\mathrm{I}_\mathbb{C}, \mathfrak{c}\}$$
% 一般来说，如果$\mathbb{L}$是$\mathbb{F}$上$m$次整式$p$的分裂域，
% $p$的$m$个根都在$\mathbb{L}$中。考虑任意$\mathbb{F}$-不变自同构$\varphi$，对$\mathbb{F}$中任意整式$q(z) = a_0 + a_1z + \cdots + a_{m-1}z^{m-1} + a_m z^m$的某个根$\lambda$，有以下等式：
% \begin{align*}
%     0 &= \varphi(0) = \varphi(q(\lambda))  \\
%       &= \varphi(a_0) + \varphi(a_1)\varphi(\lambda) + \cdots + \varphi(a_{m-1})\varphi(\lambda)^{m-1} + \varphi(a_m)\varphi(\lambda)^m  \\
%       & = q(\varphi(\lambda)) 
% \end{align*}
% 这个等式说明，$\varphi(\lambda)$也是$q$的根。取$q = p$，则$\varphi$是关于$p$的$m$个根的置换。也就是说，
% 整式的分裂域的自同构必然是根置换自同构。反之，由于$\mathbb{L}$就是往$\mathbb{F}$中添加$p$的根得到的扩域，所以，只要给定了关于$p$的根的置换，
% 就能将它延拓成$K$-不变自同构。因此结论为：
% \begin{pp}\label{pp:10_1}
%     给定域$\mathbb{F}$上整式$p$，$p$的分裂域的$\mathbb{F}$-不变自同构群由关于$p$的根置换自同构构成。
% \end{pp}
% 显然，$\mathbb{F}$上$n$次整式的分裂域，对应的$\mathbb{F}$-不变自同构群同构于$\mathfrak{S}_n$，
% 于是，问题又回到了对$\mathfrak{S}_n$的研究上。
\subsection{环道与轮换}
给定置换$\varphi$，对集合$\{1,2,\ldots, n\}$中每个元素$i$，我们追踪$j = \varphi(i)$，
如果$j\neq i$，那么继续追踪$\varphi(j)$，如果某次置换结果为$i$，则追踪结束。
如此得到一串不重复的数（最后的$i$不计入）\footnote{读者可以思考为什么没有重复。提示：置换是双射。}，
我们把这一串数称为$\varphi$-\textbf{环道}：顾名思义，
指它们在$\varphi$的作用下轮流出现，最终回到开头，构成一个闭环。
要展示环道，最好的方法是把它们画成一个环，把每个元素用箭头指向经$\varphi$映射后得到的下一个元素。
% 插入图片：环道示意图
行文中，为了方便，我们把这串数从$i$开始按顺序列出，用括号括起来，作为记录。比如考虑把$1$映射到$3$，把$3$映射到$2$，把$2$映射到$1$的环道，就记为$(1,3,2)$，最后把$2$映射到$1$不需要记录。

环道告诉我们，它的元素经过$\varphi$映射能变成哪些元素，以及需要经过几次映射。比如，环道$(1,3,2)$可以把$1$映射到$2$和$3$。其中映射到$3$只需要$1$次映射，$2$则需要两次。

如何找到$\varphi$的所有环道？我们从元素$1$出发，找出$1$所属的环道。如果它包括了$1$到$n$所有元素，
那么显然$\varphi$只有这一个环道。如果有元素不属于这个环道，我们就从里面选最小的，
继续找出它所属的环道。每个环道至少包含一个元素。所以，至多$n$次操作后，我们就得到了$\varphi$的所有环道。
两个元素同属一个环道，表示其中任一个都能经过若干次$\varphi$映射变为另一个，反之，
则说明任一个都不能经过有限次$\varphi$映射变成另一个。因此，不同的环道没有公共元素。
从环道中任何元素出发，经过$\varphi$的不断映射，直到变回自己，最少所需的映射次数是一样的。
我们称之为环道的长度。

我们可以用环道来记录置换。比如，前面提到的把$1$映射到$1$，把$2$映射到$3$，把$3$映射到$2$的置换，就可以记为：$(1)(2, 3)$。
为了简化记号，我们可以忽略长度为$1$的环道。于是这个置换也可以简记为$(2, 3)$。

只有一个环道长度大于$1$的置换称为\textbf{轮换}，它的效果是把这个环道中的元素映射到下一个，而保持其它元素不变。

环道长度为$2$的轮换称为\textbf{对换}，它的效果是交换一对元素的位置。$(2, 3)$是轮换，而且是对换。
可以看出它的效果是交换$2$和$3$的位置。

简记法唯一的问题是恒等变换：它没有长度大于$1$的环道。我们约定恒等变换简记为$(1)$。

设$n=4$，考虑$\mathfrak{S}_4$中的元素：
$$ \begin{pmatrix}
    1 & 2 & 3 & 4 \\ 2 & 1 & 4 & 3\\
\end{pmatrix} $$
按照上面的方法，它可以记为$(1, 2)(3, 4)$。它可以分解为两个环道：$(1, 2)$和$(3, 4)$。而$(1, 2)$和$(3, 4)$也是置换，
而且它们的复合就是置换$(1, 2)(3, 4)$。也就是说，$(1, 2)(3, 4)$可以按环道分解为对应轮换的复合。

值得注意的是，这时复合的顺序不影响结果，也就是说，轮换$(1, 2)$和$(3, 4)$可交换。
用归纳法不难证明，一般情况下结论也成立：如果置换$\varphi$可以分解为若干个环道，则对应轮换复合的结果就是$\varphi$，而且这些轮换两两交换。反之，任何置换都是若干个两两交换轮换的复合。
从轮换自己的角度来看，只要两个轮换涉及的环道交集为空集，这两个轮换就可以交换。

置换可以有不止一种分解成轮换的方式。比如：
$$(1,2,3) = (1,3)(1,2) = (1,2)(2,3)$$
注意到分解出来的轮换有公共元素。我们称它们相交。不考虑排列顺序，置换分解成互不相交的轮换的方式只有一种，
就是按上述找出所有环道的方式。

\section{置换的符号}
置换可以拆分为不相交的轮换。轮换能否进一步分解呢？除了恒等映射，最简单的轮换是只涉及两个元素的轮换，也就是对换。

从我们研究的目的来看，对换能够改变单位平行体的有向体积的正负，所以，我们希望把任何置换都看作对换的复合，然后研究它们作为对换复合的性质。我们尝试把轮换分解为对换的复合。
\begin{pp}\label{pp:10_5}
    给定轮换$\sigma = (i_1, i_2, \cdots, i_k)$，则$\sigma$可以分解为$k-1$个对换：
    $$ \sigma = (i_1, i_2) \circ (i_2, i_3) \circ \cdots \circ (i_{k-1}, i_k). $$
\end{pp}
\begin{proof} % [{\bfseries 证明\ref{pp:10_5}}]
    对轮换的长度$k$使用归纳法证明。$k=2$的情况命题显然成立。假设$k=m$的时候命题成立，
    考虑长度为$m+1$的轮换：$\sigma = (i_1, i_2, \cdots, i_{m+1})$，我们证明
    $$ \sigma = (i_1, i_2, \cdots, i_m) \circ (i_m, i_{m+1}).$$
    这样通过归纳假设就完成了证明。\\
    为了方便，记置换$(i_1, i_2, \cdots, i_m)$为$\alpha$，$(i_m, i_{m+1})$为$\beta$。
    考虑置换$\alpha \circ \beta$，由于$\beta$只涉及$i_m$和$i_{m+1}$，$i_1$到$i_{m-1}$经过$\beta$作用后不变，
    所以$\alpha$可以接着把$i_1$映射到$i_2$，把$i_2$映射到$i_3$，……，把$i_{m-1}$映射到$i_{m}$。\\
    同理，$\alpha$不涉及$i_{m+1}$，所以$i_m$被$\beta$映射到$i_{m+1}$后，也不再受$\alpha$影响。\\
    至于$i_{m+1}$，它被$\beta$映射到$i_m$，再被$\alpha$映射到$i_1$。\\
    综上所述，$(i_1, i_2, \cdots, i_m) \circ (i_m, i_{m+1})$就是$\sigma$。\\
\end{proof}

既然轮换能用对换复合而成，那么一切置换都能用对换复合而成。也就是说，\textbf{所有对换生成了置换群}。

任何置换都是对换复合而成，于是，我们要研究置换作为对换的复合，和恒等变换的关系。前面的计算中我们发现，对换一次下标，单位平行体的有向体积会变号，所以，我们希望研究任何置换变成恒等变换所需的对换次数。如果这个次数有某种规律，我们就可以据此确定所有单位平行体的有向体积了。

为此，研究对换对置换的影响。设有置换$\sigma$，对换$\tau = (a, b)$。
如果$a$、$b$属于$\sigma$的同一个环道\footnote{这里假设$l,r>0$，如果有等于$0$的，说明$a,b$在环道中相邻，被$\tau$分开。}：$(a=x_0, x_1, \cdots, x_l, b=y_0, y_1, \cdots, y_r)$，
那么$\tau\circ\sigma$对$a,x_1, \cdots, x_{l-1}$的影响和$\sigma$一样，但把$x_l$映射到$a$；
同时，对$b,  y_1, \cdots, y_{r-1}$的影响和$\sigma$一样，但把$y_r$映射到$b$。于是，这个环道断成了两个：
$(a, x_1, \cdots, x_l)$和$(b, y_1, \cdots, y_r)$。
如果$a$、$b$分别属于$\sigma$的两个环道\footnote{这里假设$l,r>0$，如果有等于$0$的，说明为单环道，被$\tau$并入另一个环道。}：$(a=x_0, x_1, \cdots, x_l)$和$(b=y_0, y_1, \cdots, y_r)$，
那么$\tau\circ\sigma$把$x_l$映射到$b$，把$y_r$映射到$a$，而对其他元素，影响和$\sigma$一样。于是，两个环道合为一个。
总之，对换$\tau$和置换$\sigma$复合，要么把$\sigma$的某个环道一分为二，要么把它的两个环道合并为一个。

可以验证，给定置换时，集合元素“同属一个环道”是一个等价关系，每个环道就是一个等价类。
对恒等变换来说，每个元素就是一个环道，元素个数就是环道个数。对一般置换来说，
要了解它和恒等变换的“差距”，可以用环道个数的差距来理解。对换$(a, b)$把两个元素归为一个环道，
所以环道个数和恒等变换差$1$。我们称置换$\sigma$的环道个数和恒等变换的差为\textbf{环差}。前面已经说明，不考虑排列顺序，置换分解成互不相交的轮换的方式只有一种，它的环道数和恒等变换的环道数的差是固定的。也就是说，置换的环差是置换固有的，可以记为$\delta(\sigma)$。

对换的环差是$1$。如果置换$\sigma$的环差是$k$，那么根据前面的结论对换$\tau$和$\sigma$复合后，要么某个环道一分为二，要么两个环道合并为一，所以环差要么变成$k+1$，要么变成$k-1$。考虑映射：$\mathfrak{s}:\;\sigma \mapsto\, (-1)^{\delta(\sigma)}$，
则$\mathfrak{s}$的值总是$\pm 1$，且对任何对换$\tau$和置换$\sigma$，总有
$$\mathfrak{s}(\tau\circ\sigma) = (-1)^{\delta(\sigma)\pm 1} = -\mathfrak{s}(\sigma).$$
我们把$\mathfrak{s}(\sigma)$称为置换$\sigma$的\textbf{符号}。

对换和置换复合，会改变置换的符号。前面我们提到，对换生成了置换群。
如果将置换$\sigma$分解为$k$个对换：$\sigma = \tau_1 \circ \tau_2 \circ \cdots \circ \tau_k$，那么置换的符号就是$(-1)^k$：
\begin{align*}
    \mathfrak{s}(\sigma) &= \mathfrak{s}(\tau_1 \circ \tau_2 \circ \cdots \circ \tau_k \circ (1))  \\
    &= (-1) \cdot (-1) \cdots (-1) \cdot \mathfrak{s}((1))  \\
    &= (-1)^k \cdot 1 = (-1)^k.
\end{align*}
置换的符号是由轨道数确定的，而轨道数是置换固有的性质，因此符号$\mathfrak{s}(\sigma)$也是置换固有的性质。也就是说，$(-1)^k=\mathfrak{s}(\sigma)$为固定值，与分解方法无关。这说明置换要么总是分解为奇数个对换，要么总是分解为偶数个对换。同一个置换，不可能既表示为奇数个对换的复合，
又表示成偶数个对换的复合。我们
把$\mathfrak{s}(\sigma) = -1$的置换$\sigma$称为\textbf{奇置换}，
把$\mathfrak{s}(\sigma) = 1$的置换$\sigma$称为\textbf{偶置换}，分别对应奇数和偶数个对换复合。

这对我们研究单位平行体的有向体积的意义是什么呢？以上结论告诉我们，我们可以把前面列出的$n!$个单位平行体分为两类，一类的下标对应奇置换，一类对应偶置换。对应偶置换的单位平行体，有向体积和$\varPsi(\mathbf{e}_{1}, \mathbf{e}_{2},\cdots,\mathbf{e}_{n})$相同；对应奇置换的单位平行体，有向体积和$\varPsi(\mathbf{e}_{1}, \mathbf{e}_{2},\cdots,\mathbf{e}_{n})$相反。

进一步来说，给定两个置换$\sigma_1$和$\sigma_2$，设$\sigma_1$可以分解为$k_1$个对换，
$\sigma_2$可以分解为$k_2$个对换，那么$\sigma_1\circ\sigma_2$可以分解为$k_1+k_2$个对换，于是：
$$\mathfrak{s}(\sigma_1\circ\sigma_2) = (-1)^{k_1+k_2} = \mathfrak{s}(\sigma_1)\mathfrak{s}(\sigma_2).$$
我们说$\mathfrak{s}$是从$\mathfrak{S}_n$到$\left(\{-1,1\}, \times \right)$的群同态\footnote{群同态的定义参见附录一定义\ref{df:0_5}。}。

给定任一个对换，比如$\tau = (1,2)$，考虑映射：
\begin{align*}
    \Phi_\tau : \,\, \mathfrak{S}_n &\rightarrow \mathfrak{S}_n  \\
    \sigma &\mapsto \tau\sigma 
\end{align*} 
容易验证它是双射。它把任一奇置换映射为偶置换，把任一偶置换映射到奇置换。
这说明，$\mathfrak{S}_n$中的奇置换和偶置换一样多。换句话说，前面算得的展开式中的$n!$个单位平行体的有向体积里，$1$和$-1$一样多。这和我们在平面和立体空间的例子里观察到的情形吻合。

\section{效的表示}

按照前面对置换的研究结果，我们可以把$\varPsi(\mathbf{u}_1, \mathbf{u}_2,\cdots,\mathbf{u}_n)$表示为：
$$
\varPsi(\mathbf{u}_1, \mathbf{u}_2,\cdots,\mathbf{u}_n) = \sum_{\sigma\in\mathfrak{S}_n} u_{\sigma(1),1}u_{\sigma(2),2}\cdots u_{\sigma(n),n}\varPsi(\mathbf{e}_{\sigma(1)}, \mathbf{e}_{\sigma(2)},\cdots,\mathbf{e}_{\sigma(n)}).
$$
如果把$n$个向量$\mathbf{u}_1, \mathbf{u}_2,\cdots,\mathbf{u}_n$在基底上的坐标作为列向量，排成矩阵，那么矩阵的第$i$行第$j$列元素就是$u_{i,j}$，也就是$\mathbf{u}_j$的第$i$个分量。于是我们也可以定义矩阵$M = u_{i,j}$的效为$\varPsi(\mathbf{u}_1, \mathbf{u}_2,\cdots,\mathbf{u}_n)$。要注意的是单纯讨论矩阵的效的话，不依赖于基底。但现在如果要讨论向量组的效对应的矩阵的效，就要先确定基底，然后讨论在该基底上的矩阵的效。另外，注意到向量组$T\mathbf{e}_1, \,T\mathbf{e}_2,\cdots,\,T\mathbf{e}_n$的坐标作为列向量排成的矩阵就是$T$在基底上的矩阵，所以前面定义$T$在基底上的效，就是它在该基底上的矩阵的效。我们的记法是前后一致的，可以放心使用。

前面的研究中我们知道，对换复合一个置换，会改变置换的符号，从而把奇置换改为偶置换，反之亦然。也就是说，记对换为$\tau$、置换为$\sigma$，则
$$\varPsi_{\tau\circ\sigma} = -\varPsi_{\sigma}.$$
如果我们把恒等变换对应的有向体积$\varPsi_{(1)}$看作“范本”，
那么所有置换$\sigma$对应的项$\varPsi_{\sigma}$和它都只差一个正负号。而且不难推出
这个正负号就是置换的符号。偶置换对应正号，奇置换对应负号。换句话说，我们可以把$n!$项归并为一项：
\begin{align*}
    \varPsi(\mathbf{u}_1, \mathbf{u}_2,\cdots,\mathbf{u}_n) &= \sum_{\sigma\in\mathfrak{S}_n} u_{\sigma(1),1}u_{\sigma(2),2}\cdots u_{\sigma(n),n} \varPsi\left(\mathbf{e}_{\sigma(1)}, \mathbf{e}_{\sigma(2)}, \cdots, \mathbf{e}_{\sigma(n)}\right).  \\
    &= \left(\sum_{\sigma\in\mathfrak{S}_n} \mathfrak{s}(\sigma) \prod_{i=1}^n u_{\sigma(i),i} \right) \varPsi\left(\mathbf{e}_{1}, \mathbf{e}_{2}, \cdots, \mathbf{e}_{n}\right).  
\end{align*}
而我们规定了$\varPsi\left(\mathbf{e}_{1}, \mathbf{e}_{2}, \cdots, \mathbf{e}_{n}\right) = 1$，所以
$$ 
\varPsi(\mathbf{u}_1, \mathbf{u}_2,\cdots,\mathbf{u}_n) = \sum_{\sigma\in\mathfrak{S}_n} \mathfrak{s}(\sigma) \prod_{i=1}^n u_{\sigma(i),i} .
$$
这就是向量组张成的平行体的有向体积。要注意的是，这个公式和$\varPsi$没有任何关系。也就是说，满足以上定义的任何$\varPsi$，都会定义唯一的映射。
\begin{df}\label{df:10_3}
给定$n$维平直空间$\mathbb{V}$及基底$\mathcal{B}$，关于基底$\mathcal{B}$的效映射$\det_\mathcal{B}$是唯一把$\mathbb{V}$中向量组$(\mathbf{u}_1, \mathbf{u}_2,\cdots,\mathbf{u}_n)$映射为系数域元素：
$$
\det(\mathbf{u}_1, \mathbf{u}_2,\cdots,\mathbf{u}_n) = \sum_{\sigma\in\mathfrak{S}_n} \mathfrak{s}(\sigma) \prod_{i=1}^n u_{\sigma(i),i}
$$
并将$\mathcal{B}$映射到域的乘法单位元$1_{\mathbb{F}}$的共轴退化$n$重直型。其中$u_{i,j}$是$\mathbf{u}_j$在$\mathbb{B}$中的第$i$个坐标分量。

矩阵$M = [u_{i,j}]_{1\leqslant i,j\leqslant n}$的效是其列向量组在$\mathbb{F}^n$标准基底下的效：
$$
\det(M) = \sum_{\sigma\in\mathfrak{S}_n} \mathfrak{s}(\sigma) \prod_{i=1}^n u_{\sigma(i),i}.
$$
其中$u_{i,j}$是$M$第$i$行第$j$列的元素。
\end{df}

\begin{df}\label{df:10_4}
自同态$T$关于基底$\mathcal{B} = \mathbf{e}_{1}, \mathbf{e}_{2}, \cdots, \mathbf{e}_{n}$的效就是它作用在$\mathcal{B}$上得到的向量组$T\mathbf{e}_1, T\mathbf{e}_2, \cdots , T\mathbf{e}_n$的效。
记为$\det_\mathcal{B}(T)$，它也是$T$在$\mathbb{B}$上的矩阵$[T]_{\mathbb{B}}$的效。
\end{df}

乍一看，$\det_\mathcal{B}$和基底$\mathcal{B}$的选择有关。但我们定义的效是“有向体积”的推广。而面积、体积的放缩比例，并不依赖基底的选择。从直觉来说，我们希望证明效也不依赖基底的选择。

如果我们能够证明自同态$T$在不同基底上的矩阵都有相同的效，那么它的效就不依赖于基底了。为此，我们考虑基变换对效的影响。
\begin{pp}\textbf{基变换下向量组的效 }\label{pp:10_7}
    给定$n$维平直空间$\mathbb{V}$中两个基底$\mathcal{B}$、$\mathcal{B}'$，
    $\mathbb{V}$中向量组$\mathbf{A} = (\mathbf{a}_1, \mathbf{a}_2, \cdots , \mathbf{a}_n)$在基底$\mathcal{B}$上的效和$\mathcal{B}'$上的效满足以下关系：
    $$ \det_{\mathcal{B}'}(\mathbf{A}) = \det_{\mathcal{B}'}(\mathcal{B}) \times \det_\mathcal{B} (\mathbf{A}). $$
\end{pp}
\begin{proof} % [{\bfseries 证明\ref{pp:10_7}}]
    考虑$\det_{\mathcal{B}'}$，它和$\det_{\mathcal{B}}$的区别只在于规定不同的基底张成的平行体的有向体积为1，其他完全一致，因此两者作为映射只相差一个系数。即
    $$\forall \mathbf{A},\;\; \det_{\mathcal{B}'}(\mathbf{A}) = c\cdot \det_\mathcal{B} (\mathbf{A}). $$
    其中$c$是系数。

    用$A=\mathcal{B}$代入得到：
    $$ \det_{\mathcal{B}'}(\mathcal{B}) = c \cdot \det_\mathcal{B} (\mathcal{B}) = c\cdot 1 = c. $$
    这就确定了系数$c$。也就得到：
    $$\forall \mathbf{A},\;\;  \det_{\mathcal{B}'}(\mathbf{A}) = \det_{\mathcal{B}'}(\mathcal{B}) \times \det_\mathcal{B} (\mathbf{A}). $$
\end{proof}

据此可以推断出：
\begin{pp}\label{pp:10_10}
    给定$n$维平直空间$\mathbb{V}$，
    $\mathbb{V}$中自同态$T$的效不依赖于基底。
\end{pp}
\begin{proof} % [{\bfseries 证明\ref{pp:10_10}}]
    给定空间$\mathbb{V}$中两个基底$\mathcal{B}$、$\mathcal{B}'$，
    下面证明自同态$T$在两个基底上的矩阵的效相同。\\
    考虑映射：
    $$ d_{T, \mathcal{B}} : (\mathbf{v}_1, \mathbf{v}_2, \cdots , \mathbf{v}_n) \mapsto \det_\mathcal{B} (T\mathbf{v}_1, T\mathbf{v}_2, \cdots , T\mathbf{v}_n) = \det([T]_\mathcal{B}).$$
    它也是反称的$n$重直型，和$\det_\mathcal{B}$的唯一不同在于$d_{T, \mathcal{B}}(\mathcal{B})$的值不是$1$。
    所以$d_{T, \mathcal{B}}$和$\det_{\mathcal{B}}$只相差一个系数：$d_{T, \mathcal{B}} = \lambda \cdot \det_{\mathcal{B}}$。\\
    把基底$\mathcal{B}$作为向量组代入，得到$\lambda = \det([T]_\mathcal{B})$。
    也就是说，
    $$ \det_\mathcal{B} (T\mathbf{v}_1, T\mathbf{v}_2, \cdots , T\mathbf{v}_n) = \det([T]_\mathcal{B}) \cdot \det_\mathcal{B} (\mathbf{v}_1, \mathbf{v}_2, \cdots , \mathbf{v}_n). $$
    同理，对另一组基底：
    $$ \det_{\mathcal{B}'} (T\mathbf{v}_1, T\mathbf{v}_2, \cdots , T\mathbf{v}_n) = \det([T]_{\mathcal{B}'}) \cdot \det_{\mathcal{B}'} (\mathbf{v}_1, \mathbf{v}_2, \cdots , \mathbf{v}_n). $$
    而通过基变换对向量组的效的作用，我们知道：
    \begin{align*}
        \det_{\mathcal{B}'} (T\mathbf{v}_1, T\mathbf{v}_2, \cdots , T\mathbf{v}_n) &= \det_{\mathcal{B}'}(\mathcal{B}) \cdot \det_\mathcal{B} (T\mathbf{v}_1, T\mathbf{v}_2, \cdots , T\mathbf{v}_n),  \\
        \det_{\mathcal{B}'} (\mathbf{v}_1, \mathbf{v}_2, \cdots , \mathbf{v}_n) &= \det_{\mathcal{B}'}(\mathcal{B}) \cdot \det_\mathcal{B} (\mathbf{v}_1, \mathbf{v}_2, \cdots , \mathbf{v}_n) 
    \end{align*}
    所以，等式两边代入并约去$\det_{\mathcal{B}'}(\mathcal{B})$，就得到$\det([T]_\mathcal{B}) = \det([T]_{\mathcal{B}'})$。\\
\end{proof}
至此我们得到结论：
\begin{df}
    自同态$T\in\mathcal{L}(\mathbb{V})$的效是它在任意基底上的矩阵的效，也是任何基底关于它的像在该基底上的效，记为$\det(T)$。
\end{df}

上面的证明里，我们用了任意的向量组$(\mathbf{v}_1, \mathbf{v}_2, \cdots , \mathbf{v}_n)$。显然，用特定的向量组，比如两个基底中的一个，
能够更快给出证明，但使用任意的向量组也让我们对自同态的效有了更深刻的了解。比如，如果我们把$(\mathbf{v}_1, \mathbf{v}_2, \cdots , \mathbf{v}_n)$看作
另一个自同态$S$作用在基底上的结果，那么我们就得到了关于自同态复合的效的性质：
\begin{pp}
    $\forall \;T,S\,\in\,\mathcal{L}(\mathbb{V})$，有：
    $$ \det (T\circ S) = \det(T) \cdot \det(S). $$
\end{pp}

总结一下，以上我们定义了向量组的效、自同态的效和矩阵的效。向量组的效依赖于基底的选择，自同态的效、矩阵的效不依赖于基底。
平直空间$\mathbb{V}$中自同态的效等于它在任意基底上的矩阵的效，等于它的本征式的常数项乘以$(-1)^{\dim \mathbb{V}}$，也等于它的所有（隐）本征值按照全重数乘方后的乘积。

更一般地，给定$n\times n$个数，只要按一定顺序排列成矩阵，我们也能定义它的效。如果已知矩阵的元素，
那么可以用以下的形式表示它的效：
$$ \begin{vmatrix}
    a_{1,1} & a_{1,2}  & \cdots & a_{1,n} \\
    a_{2,1} & a_{2,2} &  \cdots & a_{2,n} \\
    \vdots & \vdots & \ddots & \vdots \\
    a_{n,1} & a_{n,2} & \cdots & a_{n,n} 
\end{vmatrix}$$
也就是把矩阵的方括号改为竖线。这样的表现形式称为矩阵的\textbf{行列式}。一般的中文文献使用“行列式”一词指代效、效函数和行列式。
本书仅使用“行列式”指代矩阵的效的某种表现形式：“行列式”之于“效”，正如“竖式”之于“除法”。

\section{效的基本性质}
上一节，我们定义了自同态的效和矩阵的效。现在我们来研究效的性质。上一节末尾我们得到：
\begin{pp}
    $\forall \;T,S\,\in\,\mathcal{L}(\mathbb{V})$，有：
    $$ \det (T\circ S) = \det(T) \cdot \det(S). $$
\end{pp}
可以这么理解：效函数把自同态复合运算（矩阵乘法）映射为系数域中的标量乘法。只要系数域中的标量乘法符合交换律，就有：
$$
\forall \;T,S\,\in\,\mathcal{L}(\mathbb{V}), \quad \det (T\circ S) = \det(T) \cdot \det(S) = \det(S) \cdot \det(T) = \det(S\circ T).
$$
即自同态交换复合顺序等效。效映射关于自同态的复合运算交换。这对一般的数域都成立。
一直以来，我们不知道$T\circ S$和$S\circ T$的关系。现在我们知道，两者有相同的效。

据此又可以推得：
\begin{align*}
\forall \;A,B,C\,&\in\,\mathcal{L}(\mathbb{V}), \\
 \det (A\circ B\circ C) &= \det ((A\circ B)\circ C) = \det(C\circ(A\circ B)), \\
 \det (A\circ B\circ C) &= \det (A\circ (B\circ C)) = \det((B\circ C)\circ A),
\end{align*}

效函数是共线退化的$n$重直型，据此我们可以推出：
\begin{pp}\label{pp:10_30}
    设$T$是$n$维空间中自同态，则对标量$k$，有$\det(kT) = k^n \det(T)$.
\end{pp}
\begin{proof} % [{\bfseries 证明\ref{pp:10_30}}]
    这是$n$重直性的直接推论。\\
\end{proof}

这个结论反映了平直空间维数的本质性质。

\begin{pp}\label{pp:10_32}
    如果自同态不可逆，则它的效为$0$。如果自同态可逆，则它的效不为$0$。
\end{pp}
\begin{proof} % [{\bfseries 证明\ref{pp:10_32}}]
    从向量组的效考虑。如果向量组$\mathbf{v}_1, \mathbf{v}_2, \cdots, \mathbf{v}_n$中某个向量为零向量，
    不妨设$\mathbf{v}_1$为零向量，则对任何基底$\mathcal{B}$，
    \begin{align*}
        \det_\mathcal{B} (\mathbf{0}, \mathbf{v}_2, \cdots, \mathbf{v}_n) &= \det_\mathcal{B} (0\cdot \mathbf{u}, \mathbf{v}_2, \cdots, \mathbf{v}_n)  \\
        &= 0\cdot \det_\mathcal{B} (\mathbf{u}, \mathbf{v}_2, \cdots, \mathbf{v}_n) = 0.  
    \end{align*}
    其中$\mathbf{u}$是任意非零向量。\\
    再考虑自同态的效。自同态不可逆，等价于说它的零空间有非零向量，即它将某个非零向量映射为零向量。
    把这个向量扩充为空间的基底，则自同态的效是它作用在这组基底得到的向量组的效。这组向量组包括零向量，于是效为$0$。
    反之，如果自同态可逆，则它不把非零向量映射到零向量，于是任何基底经它映射得到的向量组都不含零向量，因此它的效不为$0$。\\
\end{proof}

这个结论给我们证明自同态可逆或不可逆提供了新的方法。更进一步，我们可以证明，自同态的效和它的逆的效有如下关系：
\begin{pp}\label{pp:10_33}
    如果自同态$T$可逆，则$\det(T^{-1}) = \frac{1}{\det(T)}$。
\end{pp}
\begin{proof} % [{\bfseries 证明\ref{pp:10_33}}]
    根据命题\ref{pp:10_31}，
    $$1 = \det(\mathrm{I}) = \det(T\circ T^{-1}) = \det(T) \det(T^{-1}).$$
\end{proof}

命题\ref{pp:10_31}和\ref{pp:10_33}说明，效函数$\det$其实是从$n$阶平直群$(\mathrm{GL}_n(\mathbb{F}), \circ)$到系数域乘法群$(\mathbb{F}, \cdot)$上的群同态。

\subsectionext{效与空间定向}
我们再来看一些与自同态的直观性质相关的讨论。效的出发点就与自同态的直观性质相关。
它表示自同态将单位平行体映射后的放缩比例。在讨论效的放缩特性时，我们引入了有向面积和有向体积的概念，来解释效为负值的情形。
在推导向量组的效的表达式的过程中，我们已经感受到基向量的排列顺序对效的影响。我们已经证明，对换（矩阵两列位置互换）会改变效的正负号。
这到底代表什么呢？

效的正负号，可以作为空间的“定向”。

经典平面中，给定两个向量 $\mathbf{u}, \mathbf{v}$，它们的效：$\det(\mathbf{u}, \mathbf{v})$给出了它们张成的平行四边形的面积。这里 $\det(\mathbf{u}, \mathbf{v})$ 可以取正值，也可以取负值。交换 $\mathbf{u}$ 和 $\mathbf{v}$ 的顺序，面积就变为原来的相反数。这说明，除了描述传统上面积的大小，效还携带了另一个信息：这两个向量的先后顺序。

从更基础的情况来分析。取定基底 $\mathfrak{B} = (\mathbf{e}_1, \mathbf{e}_2)$，考虑基向量对换后的基底 $\mathfrak{B}' = (\mathbf{e}_2, \mathbf{e}_1)$。它们张成的“平行四边形”完全重合，但按照定义，它们作为向量组的效必然一个是 $1$，一个是 $-1$。至于谁是$1$，谁是$-1$，并没有先天的条件来决定，而是由我们自由决定。但一旦决定了，就等于给平面选定了一个“方向”。确定以后，平面的图形面积都以这个“定向”来计算。

以上两种定向的关系是一种对称性。如果我们把$\mathbf{e}_1+\mathbf{e}_2$为方向的直线作为对称轴，那么两个基底关于它轴对称。于是，对应的直变换就反转这两种定向。但光凭平移和旋转是无法把一个基底变成另一个基底的。

三维立体空间中，这种关系可以看得更清楚。我们会说某个基底$\mathfrak{B} = (\mathbf{e}_1, \mathbf{e}_2, \mathbf{e}_3)$是右手系，而把其中两个基向量（比如$\mathbf{e}_1$和$\mathbf{e}_2$）调换顺序，得到的就是左手系。反之亦然。这样的一对左手系和右手系基底是镜面对称的，关于$\mathbf{e}_1+\mathbf{e}_2$和$\mathbf{e}_3$张成的平面对称。而通过平移和旋转无法把其中一个基底变成另一个。

\paragraph{有序基的定向等价类}
我们可以把上面的思路直接推广到 $n$ 维平直空间。设 $\mathbb{V}$ 是 $n$ 维实平直空间。考虑 $\mathbb{V}$ 中所有基底构成的集合。任取两组基 $\mathcal{B} = (\mathbf{e}_1, \dots, \mathbf{e}_n)$ 和 $\mathcal{B}' = (\mathbf{e}'_1, \dots, \mathbf{e}'_n)$。它们之间的基变换矩阵为 $P$，其效 $\det_{\mathcal{B}}(\mathcal{B}') = \det P$ 是一个非零实数。

我们说 $\mathcal{B}$ 与 $\mathcal{B}'$ \textbf{同定向}，当且仅当 $\det_{\mathcal{B}}(\mathcal{B}') > 0$。容易验证这是一个等价关系：自反性源于 $\det_{\mathcal{B}}(\mathcal{B}) = 1 > 0$；对称性源于 $\det_{\mathcal{B}'}(\mathcal{B}) = (\det_{\mathcal{B}}(\mathcal{B}'))^{-1}$，两者同号；传递性源于 $\det_{\mathcal{B}}(\mathcal{B}'') = \det_{\mathcal{B}}(\mathcal{B}') \cdot \det_{\mathcal{B}'}(\mathcal{B}'')$，正数乘以正数仍然是正数。

这个等价关系将 $\mathbb{V}$ 中所有的有序基恰好分为两个等价类。取定其中任意一个等价类，就称为 $\mathbb{V}$ 的一个\textbf{定向}。选定了一个定向的平直空间，称为\textbf{定向平直空间}。通常，我们约定由某个熟知的基（例如 $\mathbb{R}^n$ 的标准基）所代表的等价类为正定向，另一类为负定向。

\paragraph{效函数与定向}
在选定了定向（即选定了一个参照基 $\mathcal{B}$ 作为正定向代表）之后，效函数 $\det_{\mathcal{B}}$ 的意义变得更加鲜明。对于任意一个向量组 $\mathbf{v} = (\mathbf{v}_1, \dots, \mathbf{v}_n)$：
\begin{itemize}
    \item 若 $\det_{\mathcal{B}}(\mathbf{v}) > 0$，则这组向量构成一个与参照基同定向的有序基。它所张成的平行体具有正的“有向体积”。
    \item 若 $\det_{\mathcal{B}}(\mathbf{v}) < 0$，则它构成一个反定向的基，有向体积为负。
    \item 若 $\det_{\mathcal{B}}(\mathbf{v}) = 0$，则这组向量线性相关，不构成基，所张成的图形退化，体积为零。
\end{itemize}

\paragraph{自同态对定向的作用}
设 $T$ 是一个可逆自同态。取定空间的一组有序基 $\mathcal{B}$，$T$ 将它映射为另一组基 $T(\mathcal{B}) = (T\mathbf{e}_1, \dots, T\mathbf{e}_n)$。我们问：$T(\mathcal{B})$ 与 $\mathcal{B}$ 同定向吗？

将 $T(\mathcal{B})$ 中的每个向量用 $\mathcal{B}$ 展开，其系数矩阵正是 $T$ 在基 $\mathcal{B}$ 下的矩阵 $[T]_{\mathcal{B}}$。因此，从基 $\mathcal{B}$ 到基 $T(\mathcal{B})$ 的基变换矩阵就是 $[T]_{\mathcal{B}}$ 本身。根据同定向的定义，两者同定向当且仅当 $\det [T]_{\mathcal{B}} > 0$。而 $\det [T]_{\mathcal{B}}$ 正是自同态 $T$ 的效 $\det T$，它不依赖于基的选择。

于是，我们得到一条简洁的判则：
\begin{itemize}
    \item 若 $\det T > 0$，则 $T$ 将任意一组基变为与其同定向的基。我们称 $T$ 为\textbf{保持定向}的自同态。
    \item 若 $\det T < 0$，则 $T$ 将任意一组基变为与其反定向的基。我们称 $T$ 为\textbf{反转定向}的自同态。
\end{itemize}

恒等映射的效为 $1 > 0$，它当然保持定向。效为负的可逆算子，例如关于超平面的反射，则将正定向的基映射到负定向的基，在代数的意义上“翻转”了空间。

\section{效与矩阵}

自同态或矩阵的效，最初来自对行列式的计算。17世纪，日本算学家关孝和在解一次方程组时，最早引入了行列式的概念。他将方程组的系数按未知数分行分列排放。比如，求解以下方程组：
$$
\left\{
    \begin{aligned}
        3 x + 4 y + \phantom{2}z = 1 \\
        \phantom{2} x - 5 y - 2z = 3 \\
        -7 x - \phantom{2}y + 6z = 2 \\
    \end{aligned}
\right.
$$
就将其系数排列成：
$$
\begin{vmatrix}
    \phantom{-}3 & \phantom{-}4 & \phantom{-}1 \\
    \phantom{-}1 & -5 & -2 \\
    -7 & -1 & \phantom{-}6
\end{vmatrix}
$$
而这个式子实际上等同于计算$3\times(-5)\times6+4\times(-2)\times 7 + 1\times 1\times(-1) - (-7)\times(-5)\times 1 - (-1)\times(-2)\times 3 - 6\times 1\times 4$。这种排列方法和现在的矩阵的排列方法一致，但主要是为了计算效。

实际上，我们可以发现，通过矩阵，可以更好地理解效的计算和意义。

由于矩阵的效可以用向量组的效表示，因此，通过效函数的定义，我们可以得到以下的性质：
\begin{pp}\label{pp:10_36}
    矩阵的两列元素相同，则效等于$0$。
\end{pp}
\begin{proof} % [{\bfseries 证明\ref{pp:10_36}}]
    矩阵的两列元素其实就是向量组的两个向量。\\
    具体来说，设矩阵是向量组$\mathbf{v} = (\mathbf{v}_1, \mathbf{v}_2, \cdots , \mathbf{v}_n)$在某个基底$\mathcal{B}$上的列向量拼合而成，
    则它的效就是向量组$\mathbf{v}$的效。所以，共线退化性质“翻译”成矩阵的“语言”，就是两列元素相同的矩阵效等于$0$。
\end{proof}

\begin{pp}\label{pp:10_37}
    交换矩阵的任意两列元素，则效变号。
\end{pp}
\begin{proof} % [{\bfseries 证明\ref{pp:10_37}}]
    同样地，设矩阵是向量组$\mathbf{v} = (\mathbf{v}_1, \mathbf{v}_2, \cdots , \mathbf{v}_n)$在某个基底$\mathcal{B}$上的列向量拼合而成，
    则它的效就是向量组$\mathbf{v}$的效。交换矩阵的两列，比如说第$i$列和第$j$列，等于说第$i$个向量和第$j$个向量互换位置。
    前面我们已经做过类似的计算：
    \begin{align*}
        & \,\,\,\det_\mathcal{B} (\mathbf{v}_1,  \cdots , \mathbf{v}_i, \cdots , \mathbf{v}_j, \cdots , \mathbf{v}_n) + \det_\mathcal{B} (\mathbf{v}_1,  \cdots , \mathbf{v}_j, \cdots , \mathbf{v}_i, \cdots , \mathbf{v}_n)  \\
        =& \,\,\, \det_\mathcal{B} (\mathbf{v}_1,  \cdots , \mathbf{v}_i, \cdots , \mathbf{v}_j, \cdots , \mathbf{v}_n) + \det_\mathcal{B} (\mathbf{v}_1,  \cdots , \mathbf{v}_j, \cdots , \mathbf{v}_i, \cdots , \mathbf{v}_n)  \\
        & \,\, + \det_\mathcal{B} (\mathbf{v}_1,  \cdots , \mathbf{v}_i, \cdots , \mathbf{v}_i, \cdots , \mathbf{v}_n) + \det_\mathcal{B} (\mathbf{v}_1,  \cdots , \mathbf{v}_j, \cdots , \mathbf{v}_j, \cdots , \mathbf{v}_n)  \\
        =& \,\,\,\det_\mathcal{B} (\mathbf{v}_1,  \cdots , \mathbf{v}_i + \mathbf{v}_j, \cdots , \mathbf{v}_i + \mathbf{v}_j, \cdots , \mathbf{v}_n)  \\
        =& \,\,\, 0. 
    \end{align*}
    所以$ \det_\mathcal{B} (\mathbf{v}_1,  \cdots , \mathbf{v}_i, \cdots , \mathbf{v}_j, \cdots , \mathbf{v}_n) = - \det_\mathcal{B} (\mathbf{v}_1,  \cdots , \mathbf{v}_j, \cdots , \mathbf{v}_i, \cdots , \mathbf{v}_n). $
\end{proof}

\begin{pp}\label{pp:10_38}
    把矩阵的一列乘以一个常数，加到另一列上，矩阵的效不变。
\end{pp}
\begin{proof} % [{\bfseries 证明\ref{pp:10_38}}]
    同上，我们可以把矩阵的列“翻译”成向量组的向量。
    设矩阵是向量组$\mathbf{v} = (\mathbf{v}_1, \mathbf{v}_2, \cdots , \mathbf{v}_n)$在某个基底$\mathcal{B}$上的列向量拼合而成，
    则它的效就是向量组$\mathbf{v}$的效。把矩阵的一列乘以一个常数，加到另一列上，就是把向量组的某个向量乘以一个常数，
    加到另一个位置的向量上。因此，新向量组的效：
    \begin{align*}
        & \,\,\,\det_\mathcal{B} (\mathbf{v}_1,  \cdots , \mathbf{v}_i, \cdots , k \mathbf{v}_i + \mathbf{v}_j, \cdots , \mathbf{v}_n)  \\
        =& \,\,\, k\det_\mathcal{B} (\mathbf{v}_1,  \cdots , \mathbf{v}_i, \cdots , \mathbf{v}_i, \cdots , \mathbf{v}_n) + \det_\mathcal{B} (\mathbf{v}_1,  \cdots , \mathbf{v}_i, \cdots , \mathbf{v}_j, \cdots , \mathbf{v}_n)  \\
        =& \,\,\, k\cdot 0 + \det_\mathcal{B} (\mathbf{v}_1,  \cdots , \mathbf{v}_i, \cdots , \mathbf{v}_j, \cdots , \mathbf{v}_n)  \\
        =& \,\,\, \det_\mathcal{B} (\mathbf{v}_1,  \cdots , \mathbf{v}_i, \cdots , \mathbf{v}_j, \cdots , \mathbf{v}_n)  
    \end{align*}
    等于原向量组的效。\\
\end{proof}

此外，从矩阵的角度来看，效其实是关于矩阵元素的整式。从此出发，我们也可以得到效的一些基本性质：
\begin{pp}\label{pp:10_34}
    矩阵的共轭的效等于它的效的共轭：
    $$ \det(\overline{M}) = \overline{\det(M)}.$$
\end{pp}
\begin{proof} % [{\bfseries 证明\ref{pp:10_34}}]
    效是关于矩阵元素的整式。设矩阵$M$第$i$行第$j$列的元素记为$m_{i,j}$，则矩阵$M$的效可以明确写出：
    $$ \det (M) = \sum_{\sigma\in\mathfrak{S}_n} \mathfrak{s}(\sigma) \prod_{i=1}^n m_{\sigma(i),i}. $$
    矩阵的共轭指将矩阵所有元素取共轭得到的矩阵。所以，它的效为：
    $$ \det(\overline{M}) = \sum_{\sigma\in\mathfrak{S}_n} \mathfrak{s}(\sigma) \prod_{i=1}^n \overline{m_{\sigma(i),i}} = \overline{\sum_{\sigma\in\mathfrak{S}_n} \mathfrak{s}(\sigma) \prod_{i=1}^n m_{\sigma(i),i}} = \overline{\det(M)}. $$
\end{proof}

\begin{pp}\label{pp:10_35}
    矩阵$M$的效等于其转置的效：
    $$ \det(M) = \det(M^{\tr}).$$
\end{pp}
\begin{proof} % [{\bfseries 证明\ref{pp:10_35}}]
    设矩阵$M$第$i$行第$j$列的元素记为$m_{i,j}$，则矩阵$M$的效可以明确写出：
    $$ \det (M) = \sum_{\sigma\in\mathfrak{S}_n} \mathfrak{s}(\sigma) \prod_{i=1}^n m_{\sigma(i),i}. $$
    它的转置$M^{\tr}$的效也可以类似写出：
    $$ \det (M^{\tr}) = \sum_{\sigma\in\mathfrak{S}_n} \mathfrak{s}(\sigma) \prod_{i=1}^n m_{\sigma(i), i}. $$
    考虑其中某个置换$\sigma$，令$j = \sigma(i)$，可以得到$\prod_{i=1}^n m_{\sigma(i), i} = \prod_{j=1}^n m_{\sigma^{-1}(j), j}$。
    所以$M^{\tr}$的效可以写为
    $$ \det (M^{\tr}) = \sum_{\sigma\in\mathfrak{S}_n} \mathfrak{s}(\sigma) \prod_{i=1}^n m_{\sigma^{-1}(i), i}. $$
    $\sigma^{-1}$的符号和$\sigma$的符号是相同的，因为
    $$ \mathfrak{s}(\sigma) \cdot \mathfrak{s}(\sigma^{-1}) = \mathfrak{s}(\sigma \circ \sigma^{-1}) = \mathfrak{s}((1)) = 1. $$
    于是$M^{\tr}$的效可以写为
    $$ \det (M^{\tr}) = \sum_{\sigma\in\mathfrak{S}_n} \mathfrak{s}(\sigma^{-1}) \prod_{i=1}^n m_{\sigma^{-1}(i), i}. $$
    最后，$\sigma$遍历置换群$\mathfrak{S}_n$时，$\sigma^{-1}$恰好也将$\mathfrak{S}_n$遍历一遍。所以，
    $$ \sum_{\sigma\in\mathfrak{S}_n} \mathfrak{s}(\sigma^{-1}) \prod_{i=1}^n m_{\sigma^{-1}(i), i} = \sum_{\sigma\in\mathfrak{S}_n} \mathfrak{s}(\sigma) \prod_{i=1}^n m_{\sigma(i), i}. $$
    也就是说，
    $$ \det(M) = \det(M^{\tr}).$$
\end{proof}

这个结论告诉我们，矩阵的效也可以看作行向量构成的向量组的效。所以，前面关于矩阵列元素变换操作的结论，
可以推广到行元素变换操作上。比如，矩阵某一行的元素乘以一个常数，加到另一行上，矩阵的效不变。

\section{效的计算和展开}
根据效的表达式，计算$n$维空间自同态或$n$阶方阵的效，需要计算$n!$项乘积之和。$n=2$时需要计算两项，$n=3$时需要计算$6$项，$n=4$时要算$24$项。$n!$随着$n$增大而指数级增长，计算量巨大。如何简便地计算自同态的效呢？

我们知道自同态在任意基底上的矩阵都有相同的效。那么，如果自同态在某些基底上形式很简单，那么就能方便地计算矩阵的效。

首先来看自同态可对角化，即在某个基底上是对角矩阵的情况。写出效的表达式：
$$
\det(T) = \sum_{\sigma\in\mathfrak{S}_n} \mathfrak{s}(\sigma) \prod_{i=1}^n t_{\sigma(i),i}.
$$
其中$t_{i,j}$是$T$在基底上的矩阵$[T]$。由于$[T]$是对角矩阵，上面的项中只要有$\sigma(i)\neq i$，就有$t_{\sigma(i),i} = 0$。于是$n!$项中只有一项不为零，就是$\sigma$为恒等变换时的情况。即：
$$
\det(T) = \prod_{i=1}^n t_{i,i}.
$$
这也很好理解。因为对角线上的元素就是作为基向量的本征向量的本征值，也是$T$在该方向上的放缩比例。于是所有方向上放缩比例的乘积就是总的放缩效果。比如$n=2$时，矩形长宽的放缩系数的乘积，就是矩形面积的总变化比例。

如果自同态可三角化，即在某个基底上是上（下）三角矩阵，我们也可以简化效的计算。不妨设自同态在某基底下是上三角矩阵，也就是说，$i>j$时，$t_{i,j} = 0$。于是，在效的表达式里，只要$\sigma(i) > i$，$\sigma$对应的乘积项就是$0$了。而$\sigma$是双射，只要$\sigma$不是恒等变换，必然要把某个数映射到比它大的数。因此，和对角化的情形一样，$n!$项中只有一项不为零。，就是$\sigma$为恒等变换时的情况。即：
$$
\det(T) = \prod_{i=1}^n t_{i,i}.
$$
直观理解：自同态的矩阵是三角矩阵，说明在基底上展现的效果是某种错切变换，也就是说，除了对角线上的系数说明了它在各个方向上的放缩比例，其他系数只负责把平行体“平推”，做“切变”，于是不改变“体积”。具体来说，根据我们在自同态约简一章中的结论，这时自同态的效是所有本征值按照全重数的乘积：
$$
\det(T) = \prod\lambda_i^{\mu_a(\lambda_i)}.
$$

我们在自同态约简一章中证明了，$\mathbb{C}$平直空间中的自同态总可以三角化。于是，只要平直空间的系数域是$\mathbb{C}$的子域，比如$\mathbf{\mathbb{R}}$，总可以将其中自同态三角化，然后求其对角线元素的乘积。于是，对于实平直空间中的矩阵或自同态，我们可以将它看作复平直空间中的矩阵或自同态，将其三角化，然后求其对角线元素的乘积，就得到自同态的效。

\subsection{用效表示本征式}
给定自同态$T$，考察以下的整式：
$$
\det(x \mathrm{I} - T).
$$
它是$x \mathrm{I} - T$在任意基底上的矩阵的效。注意这个整式不依赖基底的选择，我们可以任选一个基底，计算其矩阵的效：
\begin{align*}
    \det(x \mathrm{I} - T) = \det([x \mathrm{I} - T]_{\mathcal{B}}) = \det(x\mathrm{I}_n - [T]_{\mathcal{B}}).
\end{align*}
不定元$x$恰好在矩阵$x\mathrm{I}_n - [T]_{\mathcal{B}}$的对角线上出现。具体来说，它的对角线上是$n$个一次式：$x - t_{1,1}, x - t_{2,2}, \cdot, t_{n, n}$。因此，考虑效的表达式中各项里$x$的次数。当$\sigma$是恒等变换时，$\prod_{i=1}^n (x - t_{\sigma(i),i})$是关于$x$的首一$n$次整式。对于其他的$\sigma$，对应项里的$x - t_{i,i}$出现次数都小于$n$，因此作为整式的次数小于$n$。因此，将所有项加起来，仍然是一个关于$x$的首一$n$次整式。也就是说，对于一般矩阵，$\det(x \mathrm{I} - T)$是一个关于其元素的$n$次整式。

如果$T$在某个基底上是三角矩阵或对角矩阵，那么$\det(x \mathrm{I} - T)$就是矩阵$x\mathrm{I}_n - [T]$的对角线元素的乘积。也就是：
$$
\det(x \mathrm{I} - T) = \prod_{i=1}^n (x - t_{i,i}).
$$
也就是说，$\det(x \mathrm{I} - T)$就是我们已经定义过的\textbf{本征式}$c_T$。我们已经证明了，对于可三角化矩阵，我们可以将其化为标准形式，于是它的对角线元素就是它的本征值按照全重数排列。于是，$c_T$可以写成$n$个关于$T$的本征值的一次式的乘积：
$$
\det(x \mathrm{I} - T) = c_T = \prod_{\lambda} (x - \lambda)^{\mu_a(\lambda)}
$$
其中$\lambda_i$是不同的本征值，而$\mu_a(i)$是其全重数。
另一方面，$\det(x \mathrm{I} - T)$不依赖于基底。因此，对于实系数空间的自同态，即便它作为实自同态无法三角化，我们可以考虑它作为复自同态。复自同态总可以三角化。于是可以写出
$$c_T = \prod_{\lambda} (x - \lambda)^{\mu_a(\lambda)}.$$
其中的本征值$\lambda$是复数。但$T$又能写成实系数矩阵，于是$\det(x \mathrm{I} - T)$可以是实系数整式。也就是说，我们可以得到复系数整式和实系数整式的等式：
$$ \prod_{\lambda} (x - \lambda)^{\mu_a(\lambda)} = \det(x \mathrm{I} - T) $$
这能让我们获取很多信息。比如，作为复自同态，我们知道$c_T(T) = 0$。但既然$c_T$是实系数整式，那么对于任何实向量，显然也有$c_T(T) = 0$。这说明只要按照以上方法用效定义$c_T$，全本征归零定理对实自同态也成立。
又比如，令$x = 0$，就得到$\det(T) = \prod_{\lambda}^{\mu_a(\lambda)}$，也就是说，实自同态的效是它在复空间中的本征值按全重数的乘积。即便这时的本征值是复数，它们按全重数的乘积是实数。这个结论我们会在后面再详细讨论。

\subsection{按行列展开效}
另一种计算效的方法是所谓的效的按行（列）展开。

$n\times n$矩阵$A$的效的每一项，可以看作从每一列（行）中取一个元素，使得恰好遍历所有的行（列）。因此，我们可以这样把这$n!$个项分为$n$组：所有从第一列取第一个元素构成的项作为一组，所有从第一列取第二个元素构成的项作为一组，等等。这样每组都有$(n-1)!$个元素。我们记为：
\begin{align*}
    \det(A) &= \sum_{\sigma\in\mathfrak{S}_n} \mathfrak{s}(\sigma) \prod_{i=1}^n a_{\sigma(i),i} = \sum_{k=1}^n a_{k,1}\left( \sum_{\sigma(1)=k} \mathfrak{s}(\sigma) \prod_{i=2}^n a_{\sigma(i),i}\right) \\
    &= \sum_{k=1}^n a_{k,1} S_k(A).
\end{align*}
其中$S_k(A)$是$(n-1)!$项之和。

考虑$S_k(A)$里的$(n-1)!$个元素。它恰好是把矩阵的第一列和该元素所在行删去之后，剩余的$(n-1)\times (n-1)$矩阵的行列式（可能差一个正负号）。比如，$S_3(A)$里的$(n-1)!$个元素，恰好是把矩阵的第一列和第三行删去后，$(n-1)\times (n-1)$矩阵的行列式。

这是因为除了第一列的元素已经选定，这$(n-1)!$个元素遍历了剩余$(n-1)$列元素的所有合格的取法，这就和直接计算删除了第一列和第$k$行以外的$(n-1)\times (n-1)$矩阵的效一样。唯一可能的差别是正负号：该项在$(n-1)\times (n-1)$矩阵的效里的正负号，是否和它在原矩阵的效里的正负号一样？

我们用置换的语言来重构“删除第一列和第$k$行”这件事。从置换的角度来看，$(n-1)\times (n-1)$矩阵里的某一项对应的置换，就是把原矩阵里的置换的“号码”都换了：原本的第二列变成了第一列，第三列变成了第二列；原本的第$k+1$行变成了第$k$行，第$k+2$行变成了第$k+1$行，等等。换句话说，我们可以视为把第一列和第$k$行挪到了最后（第$n$），而把原本排在它们后面的行列都往前挪了一位。因此，如果原本某一项对应的置换是$\sigma$，那么它就对应：
$$
\sigma' = (n,n-1,\cdots,k+1,k)\circ\sigma \circ(n,n-1,\cdots, 2,1).
$$
这个$\sigma'$的轨道表示和$\sigma$在$(n-1)\times (n-1)$矩阵的效里对应的置换的轨道表示一样，因为$\sigma'(n)=n$，$n$不会出现在轨道里。因此，两者的符号相同。我们只需要研究$\sigma'$的符号。根据符号的性质：
$$
\mathfrak{s} (\sigma') = \mathfrak{s}(n,n-1,\cdots,k+1,k)\cdot \mathfrak{s}(\sigma)\cdot \mathfrak{s}(n,n-1,\cdots, 2,1)
$$
轮换的符号和长度相关，$(n,n-1,\cdots,k+1,k)$是长为$n-k+1$的环道可以拆为$n-k$个对换，于是符号是$(-1)^{(n-k)}$。$(n,n-1,\cdots,2,1)$是长为$n$的环道，可以拆为$n-1$个对换，于是符号是$(-1)^{(n-1)}$。也就是说：
\begin{align*}
    \mathfrak{s} (\sigma') &= (-1)^{(n-k)}\cdot \mathfrak{s}(\sigma) \cdot (-1)^{(n-1)} \\
    &= (-1)^{(2n-k-1)}\mathfrak{s}(\sigma) = (-1)^{k+1}\mathfrak{s}(\sigma).
\end{align*}
也就是说，
\begin{align*}
    S_k(A) &= \sum_{\sigma(1)=k} \mathfrak{s}(\sigma) \prod_{i=2}^n a_{\sigma(i),i} = \sum_{\sigma(1)=k} \mathfrak{s}(\sigma) \prod_{i=1}^{n-1} a_{\sigma'(i),i} \\
    &= \sum_{\sigma'(n)=n} (-1)^{k+1}\mathfrak{s}(\sigma') \prod_{i=1}^{n-1} a_{\sigma'(i),i} = (-1)^{k+1}\sum_{\tau\in\mathfrak{S}_{n-1}} \mathfrak{s}(\tau) \prod_{i=1}^{n-1} a_{\tau(i),i} \\
    &= (-1)^{k+1} \det(A_{k,1}).
\end{align*}
其中$A_{k,1}$就是把$A$的第一列和第$k$行去掉得到的$(n-1)\times (n-1)$矩阵，称为$A$的$(k,1)-$\textbf{余子阵}。而$A$的效可以表示为：
$$
\det(A) = \sum_{k=1}^n (-1)^{k+1} a_{k,1} \det(A_{k,1}).
$$
注意到这是对第一列进行讨论得到的结果。从上面的分析来看，把第一列换成第$l$列，只会影响计算$\sigma'$，即为了计算$\sigma'$我们要将$(n,n-1,\cdots, 2,1)$改为$(n,n-1,\cdots,l+1,l)$，从而
\begin{align*}
    \mathfrak{s} (\sigma') &= (-1)^{(n-k+1-1)}\cdot \mathfrak{s}(\sigma) \cdot (-1)^{(n-l+1-1)} \\
    &= (-1)^{(2n-k-l)}\mathfrak{s}(\sigma) = (-1)^{k+l}\mathfrak{s}(\sigma).
\end{align*}
也就是说，最终$A$的效可以表示为：
$$
\det(A) = \sum_{k=1}^n (-1)^{k+l} a_{k,l} \det(A_{k,l}^{\times}).
$$
其中$A_{k,l}^{\times}$就是把$A$的第$l$列和第$k$行去掉得到的$(n-1)\times (n-1)$矩阵，称为$A$的$(k,l)-$\textbf{余子阵}。我们把这个结果称为“按第$l$列展开$A$的效”。它把$n\times n$矩阵的效拆分为$n$个$(n-1)\times (n-1)$矩阵的效。

由于矩阵的转置不改变效，所以同样有按行展开的方法，和以上形式对称。

\begin{tm}{\textbf{按行（列）展开矩阵的效}}
    我们可以把$n\times n$矩阵$A$的效按第$i$列展开。具体为：
    $$
    \det(A) = \sum_{k=1}^n (-1)^{k+i} a_{k,i} \det(A_{k,i}^{\times}).
    $$
    其中$A_{k,i}^{\times}$就是把$A$的第$i$列和第$k$行去掉得到的$(n-1)\times (n-1)$矩阵，称为$A$的$(k,i)-$\textbf{余子阵}。
    也可以把矩阵$A$的效按第$i$行展开。具体为：
    $$
    \det(A) = \sum_{k=1}^n (-1)^{i+k} a_{i,k} \det(A_{i,k}^{\times}).
    $$
    其中$A_{i,k}^{\times}$就是把$A$的第$i$行和第$k$列去掉得到的$(n-1)\times (n-1)$矩阵，称为$A$的$(i,k)-$\textbf{余子阵}。
\end{tm}

把矩阵的效按行（列）展开，有两个用处。首先，这是一个关于效函数的递推关系。我们可以通过$n-1$阶方阵的效推出$n$阶方阵的效。所以，我们可以用这种递推关系定义所有矩阵的效。从$1$阶矩阵的效等于自己唯一的元素开始，递推给出所有矩阵的效。其次，按行（列）展开给了我们一种简化计算矩阵的效的方法。表面看来，把矩阵的效按行（列）展开并没有减少计算量，我们还是要算$n!$项的加法。不过，如果某行（列）几乎全是$0$，那么按该行（列）展开后，$0$元素对应的余子阵所有项都不用计算了。即便没有这样的行（列），首先通过把一行（列）的倍数加到另一行（列）的方法，让矩阵的某行（列）变为几乎全是$0$，再按该行（列）展开，就能一口气减少大量的计算。

举例来说，我们想计算以下$4$阶矩阵的效：
\begin{align*}
    A = \begin{bmatrix}
        5    & 3  & -3 & 0 \\
        -2   & 0   & 1 & 4 \\
        1     & 0    & -2 & -1  \\
        4    & 0   & 0 & 3 \\
       \end{bmatrix}
\end{align*}
原本要计算$4!= 24$项之和。但我们观察到第$2$列几乎都是$0$。于是按第二列展开：
\begin{align*}
    \det(A) &= (-1)^{1+2} \cdot 3\cdot \det\begin{pmatrix}
        -2      & 1 & 4 \\
        1        & -2 & -1  \\
        4       & 0 & 3 \\
       \end{pmatrix}
       + (-1)^{2+2} \cdot 0\cdot \det\begin{pmatrix}
        5     & -3 & 0 \\
        1    & -2 & -1  \\
        4    & 0 & 3 \\
       \end{pmatrix}
       \\
       &+ (-1)^{3+2} \cdot 0\cdot \det\begin{pmatrix}
        5     & -3 & 0 \\
        -2    & 1 & 4 \\
        4    & 0 & 3 \\
       \end{pmatrix}+ (-1)^{4+2} \cdot 0\cdot \det\begin{pmatrix}
        5    & -3 & 0 \\
        -2    & 1 & 4 \\
        1    & -2 & -1  \\
       \end{pmatrix}
\end{align*}
注意到后三项其实都是$0$，所以根本不用计算余子阵的效。直接得到：
\begin{align*}
    \det(A) = (-1)^{1+2} \cdot 3\cdot \det\begin{pmatrix}
        -2      & 1 & 4 \\
        1        & -2 & -1  \\
        4       & 0 & 3 \\
       \end{pmatrix}.
\end{align*}
这样只需算$3! = 6$项之和，减少大量计算。又比如在上式中我们看到$3$阶方阵里没有那么“优秀”的行列，但可以通过把第一行乘以$2$加到第二行上（这不改变矩阵的效），得到：
\begin{align*}
    \det(A) = (-1)^{1+2} \cdot 3\cdot \det\begin{pmatrix}
        -2      & 1 & 4 \\
        -3       & 0 & 7  \\
        4       & 0 & 3 \\
       \end{pmatrix} = (-1)^{1+2} \cdot 3\cdot (-1)^{1+2} \cdot 1\cdot \det\begin{pmatrix}
        -3  & 7  \\
        4   & 3 \\
       \end{pmatrix}.
\end{align*}
后面一步是把$3$阶矩阵的效再次按第二列展开。最后只需计算$2$阶矩阵的效，可以轻松获得结果。

当然，这种简化计算的效果并不是最好的。比如我们可以直接通过消元法，把矩阵变为上三角矩阵，那么矩阵的效就是对角线元素的乘积。对于无法变为三角矩阵的情况，我们有三角分解算法，以和消元法类似的思路，把任意方阵化为上三角矩阵和下三角矩阵的乘积，于是其效就等于两者乘积，可以通过分别算对角线元素的乘积得到。具体内容本书不做展开，读者可以参考矩阵计算方面的论述。

以上讨论中我们建立了矩阵的余子阵$A_{i,j}^{\times}$的概念。它可以看作关于位置信息$(i,j)$的映射。我们可以把这个概念推广：
\begin{df}
    给定$m\times n$的矩阵$A = [a_{i,j}]_{1\leqslant i\leqslant m,1\leqslant j\leqslant n}$和正整数$p< m, \;q< n$。设有严格递增下标数组$I=(i_1,i_2,\cdots,i_p)$、$J=(j_1,j_2,\cdots,j_q)$\footnote{即是说$1\leqslant i_1<i_2<\cdots<i_p\leqslant m$、$1\leqslant j_1<j_2<\cdots<j_q\leqslant n$。}，则我们可以定义矩阵关于行号$i_1,i_2,\cdots,i_p$和$j_1,j_2,\cdots,j_q$的\textbf{子阵}$A_{I,J}$和它的\textbf{余子阵}$A_{I,J}^{\times}$。$A_{I,J}$是$I,J$中的行号、列号对应的行列交叉的部分取出来组成的$p\times q$矩阵。$A_{I,J}^{\times}$是把$A$的相关行号、列号的行列移除后剩下的$(m-p)\times(n-q)$矩阵。
\end{df}
这个定义在符号和记录方法上很繁琐，但意思很好理解，就是把删一行一列改为删多行多列。这样的子阵和余子阵可以用来定义分块矩阵。给定$m\times n$矩阵，首先指定正整数$p< m, \;q< n$，我们可以按照$p,q$来给$A$分块。指定$I=\hin{1}{p}$、$J=\hin{1}{q}$，那么子阵$A_{I,J}$就是$A$的前$p$行$q$列的元素构成的矩阵（左上角部分），而$A_{I,J}^{\times}$就是后$(m-p)$行$(n-q)$列的元素构成的矩阵（右下角部分）。同理，指定$I'=\hin{p+1}{n}$、$J=\hin{1}{q}$，那么我们可以得到$A$的后$(m-p)$行前$q$列的元素构成的矩阵（左下角部分）和前$p$行后$(n-q)$列的元素构成的矩阵（右上角部分）。可以说我们在第$p$行、第$q$列后把矩阵$A$切成四块，那么就得到四个子阵。它们两两成对，互为余子阵。我们可以把$A$写成：
$$
A = \begin{bmatrix}
    A_{I,J} & A_{I',J}^{\times} \\ A_{I',J} &  A_{I,J}^{\times}
\end{bmatrix}.
$$
分块矩阵在矩阵论里有极大的应用，能得出关于直映射的很多结果。想要了解相关知识可以学习矩阵论的内容。

\subsection{附阵与方程组的解}

回顾按矩阵第$i$列展开效的公式：
$$
\det(A) = \sum_{k=1}^n (-1)^{k+i} a_{k,i} \det(A_{k,i}^{\times}).
$$
其中$A_{k,i}^{\times}$的元素都不在第$i$列。如果把公式中的$a_{k,i}$替换成别的$n$个数，就等于把矩阵的第$i$列换掉以后，再按第$i$列展开，计算其效。那么，如果把公式中的$a_{k,i}$替换成$a_{k,j}$，即另一列，会怎么样呢？由于第$i$列被替换成第$j$列，所以只要$j\neq i$，矩阵就有两列一模一样。按照效的基本性质，这时候矩阵的效等于零。

于是我们可以设计这样一个矩阵$\co(A) = [c_{i,j}]_{1\leqslant i,j\leqslant n}$，它和$A$一样也是$n\times n$矩阵，它的第$i$行第$j$列的元素是：
$$
c_{i,j} = (-1)^{j+i} \det(A_{j,i}^{\times})
$$
这样当我们计算乘积$M = \co(A)\cdot A$，就有：
$$
m_{i,j} = \sum_{k=1}^n c_{i,k} a_{k,i} = (-1)^{k+i} \det(A_{k,i}^{\times}) a_{k,i}.
$$
由上面讨论可知，$i=j$时，$m_{i,j} = \det(A)$，$i\neq j$时，$m_{i,j} = 0$。因此，$M$是一个对角矩阵，对角线上所有的元素都是$\det(A)$，而其余位置的元素都是$0$。即：
$$
\co(A)\cdot A = \det(A)I_n.
$$
同理，我们可以计算$M' = A\co(A)$的第$i$行第$j$列：
\begin{align*}
    m'_{i,j} = \sum_{k=1}^n a_{i,k} c_{k,j} = (-1)^{j+k}a_{i,k} \det(A_{j,k}^{\times}) .
\end{align*}
这正是在把$A$按第$j$行展开的公式里，把$a_{j,k}$换成$a_{i,k}$的结果。按同样的论据，$i\neq j$时，$m'_{i,j}=0$；$i=j$时，$m'_{i,j} = \det(A)$。于是得到：
$$
A\cdot \co(A) = \det(A)I_n = \co(A)A.
$$
也就是说，如果$\det(A)\neq 0$（即矩阵$A$可逆），那么我们能用余子阵具体地写出$A$的逆矩阵：
$$
A^{-1} = \frac{1}{\det(A)}\co(A).
$$
其第$i$行第$j$列的元素是
$$\left(A^{-1}\right)_{i,j} = \frac{(-1)^{i+j} \det(A_{j,i}^{\times})}{\det(A)}.$$
我们把$\co(A)$称为$A$的\textbf{随附矩阵}或\textbf{附阵}\footnote{我们把随附矩阵$\co(A)$的转置称为$A$的\textbf{余阵}，记为$\operatorname{Com}(A)$：$\co(A) = \operatorname{Com}(A)^{\tr}$。余阵的第$i$行第$j$列恰好就是原矩阵移除第$i$行第$j$列得到的余子阵，而随附矩阵则是其转置。}。注意：
\begin{enumerate}
    \item 随附矩阵是关于原矩阵$A$的元素的多项式。$A$可逆时，$A$的随附矩阵说明$A^{-1}$是关于$A$的元素的有理函数。
    \item $\co(A)\cdot A = \det(A)I_n$这个关系式对任意矩阵$A$都成立，不需要$A$可逆；而后面$A^{-1}$的表达式则显然需要$A$可逆的条件。两者的成立条件不同，不可混淆。
\end{enumerate}


关于矩阵的附阵，有以下基本性质，读者可以自行验证或参见附录三证明\ref{c-25}：
\begin{pp}\label{pp:10_50}
    设$A,B$为$n\times n$矩阵，则其附阵满足以下性质：
    \begin{enumerate}
    \item $\co(I_n) = I_n.$
    \item $\co(A^{\tr}) = \co(A)^{\tr}.$
    \item $\co(kA) = k^{n-1}\co(A).$
    \item $\co(AB) = \co(B)\co(A). $
    \item 如果$A$可逆，那么$\co(A^{-1}) = \co(A)^{-1} = \frac{A}{\det(A)}.$
    \item $\det(\co(A)) = \det(A)^{n-1}.$
\end{enumerate}
\end{pp}

附阵的一个基本应用是解一次方程组：
$$ A\mathbf{x} = \mathbf{b}.$$
如果系数矩阵$A$可逆，那么我们知道：
$$ \mathbf{x} = A^{-1}\mathbf{b}.$$
但之前我们没有明确给出求矩阵的逆的方法。而使用附阵我们可以给出方程的解的直接表达式：
$$
    \mathbf{x} =  \frac{1}{\det(A)}\co(A)\mathbf{b}.
$$
这个等式成立，因为$A$可逆当且仅当$\det(A)\neq 0$。具体来说，$\mathbf{x}$的第$i$个分量
$$
x_i = \frac{\co(A)\mathbf{b}}{\det(A)} = \frac{\det(A_i(\mathbf{b}))}{\det(A)}. 
$$
其中矩阵$A_i(\mathbf{b})$是把$A$的第$i$列换成向量$\mathbf{b}$得到的矩阵。这个结论也叫\textbf{克莱姆法则}。

如果$A$不可逆而解$\mathbf{x}$存在，那么有：
$$
\co(A)\mathbf{b} = \co(A)A\mathbf{x} = \det(A)I_n\mathbf{x} = \mathbf{0}.
$$
这说明方程组在$A$不可逆时仍有解的必要条件是$\co(A)\mathbf{b} =\mathbf{0}$（但并非充分条件）。


\chapter{内积空间}
我们定义平直空间时，把它看作是解析几何中平面和立体空间的推广。
平直空间“继承”了它们的“直性”，但还缺少了一些重要的元素。
比如说，为了研究平面和空间中的几何对象，我们常常需要计算距离、长度和角度。
然而在平直空间里，我们还没有类似的概念。
从本章开始，我们会为平直空间“配备”距离、长度、角度等概念，
不过，我们会用一种乍看起来不太一样的方法来定义它们。
这就是我们接下来要介绍的内积的概念。

\begin{center}
    \fbox{
        \shortstack[l]{
      下文中，除非另有声明，我们总假设$\mathbb{F}$是常见数域如\\
      $\mathbb{Q}$、$\mathbb{R}$或$\mathbb{C}$，
      $\mathbb{V}$是$\mathbb{F}$上的有限维平直空间，维数$n$大于等于一。
        }
    }
\end{center}

\section{内积}
\subsection{共同的不变量}
平面中，两个向量$(a_1, b_1)$和$(a_2, b_2)$如果满足$a_1a_2 + b_1b_2 = 0$，这两个向量就相互垂直。
如果我们考察$a_1a_2 + b_1b_2$这个量，会发现它和距离、长度、角度等概念密切相关。

比如，向量$(a_2, b_2)$对应的点到法向量为$(a_1, b_1)$的直线$a_1x + b_1y = 0$的距离，是：
$$ d = \frac{a_1a_2 + b_1b_2}{\sqrt{a_1^2 + b_1^2}}. $$

如果$a_1 = a_2$，$b_1 = b_2$，那么$ a_1a_2 + b_1b_2$就变成$a_1^2 + b_1^2$，也就是向量$(a_1, b_1)$的长度的平方。

如果固定$(a_1, b_1)$不变，那么$(a_2, b_2)$和它形成的夹角$\theta$也和它相关：
$$ \cos \theta = \frac{a_1a_2 + b_1b_2}{\sqrt{a_1^2 + b_1^2}\sqrt{a_2^2 + b_2^2}}. $$

立体空间中，我们也用类似的方法定义垂直、长度和夹角。这个定义有什么本质特点呢？

考察函数$\digamma$：
\begin{align*}
    \sigma: \,\,\,\mathbb{R}^2 \times \mathbb{R}^2 &\rightarrow \mathbb{R}  \\
    \left((a_1, b_1) , (a_2, b_2)\right) &\mapsto a_1a_2 + b_1b_2 
\end{align*}
$\digamma$将两个$\mathbb{R}^2$中的向量映射为一个实数。如果固定其中一个变量，函数关于另一个变量是直映射。
我们将这种映射称为双直映射，即关于两个变量中的任一个都是直映射。
由于映射结果是一个实数，仿照“直型”的说法，我们称这种映射为\textbf{双直型}。

此外，如果调换两个向量的位置，得到的结果不变。我们说$\digamma$是（关于变量）对称的映射。

如果两个向量相等，映射结果为：
$$\digamma\left((a_1, b_1) , (a_1, b_1)\right) = a_1^2 + b_1^2.$$
所以映射结果总是非负实数，当且仅当$a_1 = b_1 = 0$（零向量）的时候映射结果为0，其余时候都严格大于0.
我们称这类映射为\textbf{正定映射}。

我们将这些特点集合起来，为一般的有限维实空间$\mathbb{R}^n$定义一种类似的映射：

\begin{df}\label{df:7_1}
    给定空间$\mathbb{R}^n$，从$\mathbb{R}^n \times \mathbb{R}^n$映射到$\mathbb{R}$的对称正定双直型称为$\mathbb{R}^n$上的\textbf{内积}。
    配置了内积的空间$\mathbb{R}^n$称为\textbf{内积空间}$\mathbb{R}^n$。
\end{df}
由于$n$维实平直空间和$\mathbb{R}^n$同构，所以我们实际上也定义了一般有限维实空间上的内积。

以下我们暂且约定，不另作说明时，提到的平直空间都是有限维实空间$\mathbb{R}^n$。

我们一般将内积用$(\,\cdot \,|\, \cdot\,)$符号表示。比如，给定实空间中两个向量：$\mathbf{u}, \mathbf{v}$，它们的内积记作$(\mathbf{u} | \mathbf{v})$.

\subsection{内积的性质}
首先来看一个最典型的例子。我们定义以下映射：
\begin{align*}
    \mathbb{R}^n \times \mathbb{R}^n &\rightarrow \mathbb{R}  \\
   \left((u_1, u_2, \cdots, u_n) , (v_1, v_2, \cdots v_n)\right) &\mapsto u_1v_1 + u_2v_2 + \cdots + u_n v_n 
\end{align*}

可以按内积的定义验证，这个映射是一个内积。
我们称它为\textbf{点积}。不至于混淆的情况下，一般记作$\mathbf{u} \cdot \mathbf{v}$。配备了点积的空间$\mathbb{R}^n$称为\textbf{欧几里得空间}。
很多我们在平面几何或空间几何中学习过的结论，都可以推广到欧几里得空间里面来。

$\mathbb{R}^n$上还可以有什么其他的内积呢？我们可以定义
\begin{align*}
    \mathbb{R}^n \times \mathbb{R}^n &\rightarrow \mathbb{R}  \\
   \left((u_1, u_2, \cdots, u_n) , (v_1, v_2, \cdots v_n)\right) &\mapsto 2u_1v_1 + 2u_2v_2 + \cdots + 2u_n v_n 
\end{align*}
可以验证，这也是一个内积。

更一般地，给定$n$个正实数：$c_1, c_2, \cdots, c_n$，以下的映射
\begin{align*}
    \mathbb{R}^n \times \mathbb{R}^n &\rightarrow \mathbb{R}  \\
   \left((u_1, u_2, \cdots, u_n) , (v_1, v_2, \cdots v_n)\right) &\mapsto c_1u_1v_1 + c_2u_2v_2 + \cdots + c_nu_n v_n 
\end{align*}
也是一个内积。可以看作是“加权”的点积。
为了满足正定性质，“权重”不可以是0或复数。

内积给我们带来了什么呢？作为映射，内积将两个向量映射到一个实数。实数有一个重要的基本性质：完全有序。
用通俗的语言来说，任意两个实数，都可以比较大小。这对于一般的域来说不成立，因为不是所有的域都配备了\textbf{全序关系}。

\begin{df}\label{df:7_2}
    给定集合$X$，$X$上的\textbf{全序}指满足以下条件的二元关系$\leqslant :  X \times X \rightarrow X$：
    \begin{itemize}
        \item 反对称性：$ \forall a, b \,\,\, \in X, \,\,\, a \leqslant b, b\leqslant a \Rightarrow a = b.$
        \item 传递性：$ \forall a, b, c \,\,\, \in X, \,\,\, a \leqslant b, b\leqslant c \Rightarrow a \leqslant c.$
        \item 完全性：$ \forall a, b \,\,\, \in X$，或者$a \leqslant b$，或者$b\leqslant a$，没有其他可能性\footnote{但两者不互斥，可以同时成立。}。
    \end{itemize}
    
\end{df}

实数集上的大于等于关系$\leqslant$就是典型的全序，于是，任意内积的结果可以比较大小，我们可以像在平面和立体空间中那样来使用它。

\section{模}
\subsection{什么是长度？}
有了内积，我们可以用它来定义“长度”概念，将平面和三维空间中的长度概念推广到实空间中。

首先来看“长度”这个概念，当我们提到、使用“长度”的时候，我们希望用它有什么性质？

我们把“长度”看作一个映射$N$，它将平直空间中每个向量都有定义，映射的结果称为这个向量的长度。

首先，长度应该是一个非负的数，用来衡量向量的“长短”，也就是说，每个向量都对应一个非负的数，可以互相比较。
我们希望只有零向量的长度是0，其他向量的长度都大于0.
这个性质也称为\textbf{正定性}。

此外，我们希望长度兼容向量的数乘运算。也就是说，向量乘以一个系数$a$后，它的长度也变为原来的$|a|$倍\footnote{注意这里要加绝对值。}。

最后，我们希望长度满足三角不等式：
$$ \forall \,\, \mathbf{u}, \mathbf{v}  \,\, \in \,\, \mathbb{V}, \,\,\, N(\mathbf{u}) + N(\mathbf{v}) \geqslant N(\mathbf{u}+\mathbf{v}). $$
当且仅当$\mathbf{u}, \mathbf{v}$直相关时，等号成立\footnote{实际上，最后这个要求并不是必须的。关于不满足三角不等式的“长度”，也有很多研究。本书不讨论这方面的知识。}。

满足以上几个性质的映射$N : \,\, \mathbb{V} \rightarrow \mathbb{R}$，我们称其为\textbf{模}或\textbf{范定义数}。
模是长度概念在实空间上的推广。配备了模的平直空间，称为\textbf{有长平直空间}或\textbf{有长空间}。

\begin{pp}\label{pp:7_1}
    每个内积都可以导出一个模，称为其（自然）对应的模。
\end{pp}

\begin{proof} % [{\bfseries 证明\ref{pp:7_1}}]\label{pf:7_1}
    给定$\mathbb{R}^n$上的内积$(\cdot | \cdot)$，定义映射：
    \begin{align*}
        N : \mathbb{R}^n &\rightarrow \mathbb{R}  \\
        \mathbf{u} & \mapsto  \sqrt{(\mathbf{u} | \mathbf{u})} 
    \end{align*}
    下面证明$N$是一个模。\\
    首先，内积是正定映射，由定义可以推出$N$是正定的。\\
    其次，内积是双直映射，所以：
    \begin{align*}
        \forall \,\, a \,\, \in \mathbb{R}, \,\,\,  \mathbf{u} \,\, \in \mathbb{R}^n, \,\,\,  \\
        N(a\mathbf{u}) &= \sqrt{(a\mathbf{u} | a\mathbf{u})}  \\
        &= \sqrt{a^2(\mathbf{u} | \mathbf{u})}  \\
        &= |a| \sqrt{(\mathbf{u} | \mathbf{u})} = |a|N(\mathbf{u}))
    \end{align*}
    最后，证明三角不等式：
    \begin{align*}
        \forall \,\, \mathbf{u} ,  \mathbf{v} \,\, \in \mathbb{R}^n, \,\,\,  \\
        (N(\mathbf{u}) + N(\mathbf{v}))^2 &= (\mathbf{u} | \mathbf{u}) + (\mathbf{v} | \mathbf{v}) + 2\sqrt{(\mathbf{u} | \mathbf{u})(\mathbf{v} | \mathbf{v})}  \\
        N(\mathbf{u}+\mathbf{v})^2 &= (\mathbf{u} | \mathbf{u}) + (\mathbf{v} | \mathbf{v}) + 2(\mathbf{u} | \mathbf{v}) 
    \end{align*}
    所以，只需证明
    $$ (\mathbf{u} | \mathbf{u})(\mathbf{v} | \mathbf{v}) \geqslant (\mathbf{u} | \mathbf{v})^2. $$
    注意到$\mathbf{u}$或$\mathbf{v}$为零向量时，上式显然成立。对于两者都不为零向量的情况，
    考虑以下关于实数$\lambda$的表达式：
    $$ (\mathbf{u}+ \lambda\mathbf{v} |\mathbf{u}+ \lambda\mathbf{v} ). $$
    它是一个关于$\lambda$的二次整式：
    $$ (\mathbf{u}+ \lambda\mathbf{v} |\mathbf{u}+ \lambda\mathbf{v} ) = (\mathbf{v} | \mathbf{v}) \lambda^2 + 2 (\mathbf{u} | \mathbf{v}) \lambda + (\mathbf{u} | \mathbf{u}). \quad\quad (*)$$
    另一方面，根据内积的定义：
    $$ (\mathbf{u}+ \lambda\mathbf{v} |\mathbf{u}+ \lambda\mathbf{v} ) \geqslant 0.$$
    所以，$(*)$式对应的二次方程
    $$  (\mathbf{v} | \mathbf{v}) \lambda^2 + 2 (\mathbf{u} | \mathbf{v}) \lambda + (\mathbf{u} | \mathbf{u}) = 0  $$
    的判别式小于等于0. 而它的判别式就是：
    $$ \Delta = 4 (\mathbf{u} | \mathbf{v})^2 - 4 (\mathbf{u} | \mathbf{u}) (\mathbf{v} | \mathbf{v}). $$
    这就说明，
    $$ (\mathbf{u} | \mathbf{u})(\mathbf{v} | \mathbf{v}) \geqslant (\mathbf{u} | \mathbf{v})^2. $$
    取等号时，$(1)$式对应的二次方程恰有重根，也就是说存在唯一的$\lambda$使得$\mathbf{u}+ \lambda\mathbf{v}$为零向量。
    也就是说，$\mathbf{u}, \mathbf{v}$直相关。\\     
\end{proof}

每个内积都可以导出一个模，所以内积空间自然也是赋范空间。

$\mathbb{R}^n$上的点积对应导出的是以下模：
$$ N(u_1, u_2, \cdots , u_n) = \sqrt{u_1^2 + u_2^2 + \cdots + u_n^2}. $$
可以看出，它是平面几何和空间几何中长度概念的自然推广，一般称为\textbf{欧几里得模}。

\subsection{长度与距离}
有了长度概念，我们可以继续定义“距离”。
给定集合$X$，我们定义距离为满足以下性质的映射：
\begin{align*}
    X \times X &\rightarrow \mathbb{R}  \\
   (a, b) &\mapsto \mathrm{d}(a, b)
\end{align*}
\begin{itemize}
    \item 正定性：$ \forall a, b \,\,\, \in X, \,\,\, \mathrm{d}(a, b) \geqslant 0$，而且$\mathrm{d}(a, b) = 0 \Leftrightarrow a = b.$
    \item 对称性：$ \forall a, b \,\,\, \in X, \,\,\, \mathrm{d}(a, b) = \mathrm{d}(b, a).$
    \item 三角不等式：$ \forall a, b, c \,\,\, \in X, \,\,\, \mathrm{d}(a, b) + \mathrm{d}(b, c) \geqslant \mathrm{d}(a, c).$
\end{itemize}
不难验证，$\mathrm{d}(\mathbf{u}, \mathbf{v}) = N(\mathbf{u} - \mathbf{v})$就定义了一个距离函数。
所以，和平面、三维空间的几何一样，（原点出发的）向量之差的长度就是向量对应点之间的距离。
而从内积可以导出模，所以距离函数也可以直接从$\mathrm{d}(\mathbf{u}, \mathbf{v}) = (\mathbf{u} - \mathbf{v}|\mathbf{u} - \mathbf{v})$定义。

欧几里得空间$\mathbb{R}^n$的距离是：
\begin{align*}
    &\mathrm{d}\left((u_1, u_2, \cdots , u_n), (v_1, v_2, \cdots , v_n) \right)  \\
  = &\sqrt{(u_1-v_1)^2 + (u_2-v_2)^2 + \cdots + (u_n-v_n)^2}. 
\end{align*}

要注意的是，模和距离也可以脱离内积的定义存在。赋范空间不一定是内积空间。
比如，我们可以定义以下模：
$$ N(u_1, u_2, \cdots , u_n) = \max\left(|u_1|, |u_2|, \cdots, |u_n|\right). $$
这个模不能被任何内积导出。

而距离也可以脱离模存在。定义了距离的空间称为\textbf{度量空间}。度量空间不一定是赋范空间。
比如，距离函数不一定满足平移不变性：
$$ \forall \,\, a, b, c, \,\, \in X \,\,\, \mathrm{d}(a, b) = \mathrm{d}(a + c, b + c). $$
但模导出的距离函数总满足平移不变性\footnote{读者可以试着构造一个不满足平移不变性的距离函数。}。

以下，我们将内积导出的模用$\|\,\cdot\,\|$表记。
比如，向量$\mathbf{u}$的模记作：$\|\mathbf{u}\|$。

\section{夹角与正交}
\subsection{向量的夹角}
定义了内积导出的长度、距离等概念后，可以继续定义夹角的概念。
为此，我们先回顾一下平面几何中向量夹角的定义。
给定平面向量$(a_1, b_1)$和$(a_2, b_2)$，它们的夹角$\theta$满足：
$$ \cos \theta = \frac{a_1a_2 + b_1b_2}{\sqrt{a_1^2 + b_1^2}\sqrt{a_2^2 + b_2^2}}. $$

这个等式总成立，是因为$|a_1a_2 + b_1b_2| \leqslant \sqrt{a_1^2 + b_1^2}\sqrt{a_2^2 + b_2^2}.$ 
这样，上式右边的值总在区间$[-1,1]$上，能够定义为夹角的余弦值。
对于一般的实内积空间，是否可以这个关系呢？
我们注意到，证明内积导出的模满足三角不等式的过程中，我们附带证明了以下不等式：
$$ |(\mathbf{u} | \mathbf{v})| \leqslant \|\mathbf{u}\| \cdot \|\mathbf{v}\|$$
这就是我们所需的关系。它叫做\textbf{内积不等式}。有了这个不等式，我们可以定义两个向量的夹角为满足：
$$ \cos \theta = \frac{(\mathbf{u} | \mathbf{v})}{\|\mathbf{u}\|  \|\mathbf{v}\|} $$
的$\theta$。

给定向量$\mathbf{u}$，考察函数$ f: \,\, \mathbf{v} \mapsto \frac{(\mathbf{u} | \mathbf{v})}{\|\mathbf{u}\|  \|\mathbf{v}\|}$，可以发现：
\begin{itemize}
    \item $f$是一个“奇函数”：$f( -\mathbf{v}) = -f(\mathbf{v})$
    \item 两个非零向量直相关（$\mathbf{u} = a\mathbf{v}$）时，如果$a>0$，$f(\mathbf{v}) = 1$，如果$a<0$，$f(\mathbf{v}) = -1.$
    \item 两个非零向量直无关时，如果$f(\mathbf{v}) \neq 0$，那么$f\left(\mathbf{u}-\frac{(\mathbf{u} | \mathbf{u})}{(\mathbf{u} | \mathbf{v})}\mathbf{v}\right) = 0.$
\end{itemize}
从最后一点可以得出，给定向量$\mathbf{u}$，总能找出一个非零向量$\mathbf{v}$，使得它们夹角的余弦值为0. 具体来说，就是
$$ (\mathbf{u} | \mathbf{v}) = 0.$$
我们把这样的向量称为和$\mathbf{u}$\textbf{正交}的向量。
两个向量$\mathbf{u}, \mathbf{v}$满足$(\mathbf{u} | \mathbf{v}) = 0$，就说它们\textbf{正交}。
如果两个向量集合$U, V$满足：
$$ \forall \,\, \mathbf{u} \,\, \in \,\, U, \, \,\, \mathbf{v} \,\, \in \,\, V, \,\,\, (\mathbf{u} | \mathbf{v}) = 0,$$
就说这两个集合\textbf{正交}。

正交就是平面几何中垂直概念的推广。

按定义不难验证，对于正交的向量，勾股定理成立：
\begin{pp}\label{pp:7_2}
    两个向量$\mathbf{u}, \mathbf{v}$正交，则：
    $$ \|\mathbf{u}\|^2 + \|\mathbf{v}\|^2 = \|\mathbf{u} + \mathbf{v}\|^2.$$
\end{pp}

\subsection{正交分解}
给定非零向量$\mathbf{u}, \mathbf{v}$，我们可以把$\mathbf{u}$写成一个与$\mathbf{v}$共线的非零向量
和一个与$\mathbf{v}$正交的向量$\mathbf{w}$之和。具体来说，我们可以这么找向量$\mathbf{w}$：

对非零实数$a$，我们把$\mathbf{u}$写成
$$\mathbf{u} = a\mathbf{v} + (\mathbf{u} - a\mathbf{v}) $$
要找一个$a$，使得
$$ (a\mathbf{v} | \mathbf{u} - a\mathbf{v}) = 0.$$
由于我们设$a \neq 0$，上式可以解得：
$$ a = \frac{(\mathbf{u} | \mathbf{v})}{(\mathbf{v} | \mathbf{v})}.$$
也就是说，$\mathbf{u}$可以写成
$$\mathbf{u} = \frac{(\mathbf{u} | \mathbf{v})}{(\mathbf{v} | \mathbf{v})}\mathbf{v} + \left(\mathbf{u} - \frac{(\mathbf{u} | \mathbf{v})}{(\mathbf{v} | \mathbf{v})}\mathbf{v}\right). $$
我们把它称为向量$\mathbf{u}$依$\mathbf{v}$的\textbf{正交分解}。

直观来说，$\frac{(\mathbf{u} | \mathbf{v})}{(\mathbf{v} | \mathbf{v})}\mathbf{v}$是向量$\mathbf{u}$在$\mathbf{v}$所在直线上的投影，
而$\left(\mathbf{u} - \frac{(\mathbf{u} | \mathbf{v})}{(\mathbf{v} | \mathbf{v})}\mathbf{v}\right)$是投影余量（从投影点指向$\mathbf{u}$的端点的向量）。

如果我们计算$\left(\mathbf{u} - \frac{(\mathbf{u} | \mathbf{v})}{(\mathbf{v} | \mathbf{v})}\mathbf{v}\right)$ 的模，会得到：

\begin{align*}
    \|\left(\mathbf{u} - \frac{(\mathbf{u} | \mathbf{v})}{(\mathbf{v} | \mathbf{v})}\mathbf{v}\right)\|^2 &= 
    \|\mathbf{u}\|^2 + \frac{(\mathbf{u} | \mathbf{v})^2}{\|\mathbf{v} \|^4} \|\mathbf{v} \|^2 - 2 \frac{(\mathbf{u} | \mathbf{v})}{\|\mathbf{v} \|^2} \mathbf{u} | \mathbf{v})  \\
    &= \|\mathbf{u}\|^2 - \frac{(\mathbf{u} | \mathbf{v})^2}{\|\mathbf{v} \|^2}  
\end{align*}

这个模大于等于0，恰好对应内积不等式。所以，内积不等式的直观解释为：
向量$\mathbf{u}$依$\mathbf{v}$正交分解时，投影余量的长度大于等于0. 
或者说，向量的长度总大于等于它在任何方向上投影的长度。

\subsection{内积空间中的形体}
作为例子，我们可以看一些基本几何对象在内积空间中的推广。

平面解析几何中，圆心位于$(a,b)$、半径为$R$的圆的方程是：
$$ (x - a)^2 + (y - b)^2 = R^2.$$
它表示圆上的点到圆心的距离是$R$。立体解析几何中，球心位于$(a,b,c)$、半径为$R$的球面的方程是：
$$ (x - a)^2 + (y - b)^2 + (z - c)^2= R^2.$$
它表示球面上的点到球心的距离是$R$。欧几里得空间$\mathbb{R}^n$中，一点到点$\mathbf{c}$的距离为$R$这件事可以用以下方程描述：
$$ \| \mathbf{x} - \mathbf{c}\| = R. $$
记$\mathbf{x} = (x_1, \cdots , x_n)$，$\mathbf{c} = (c_1, \cdots , c_n)$，则上式还可以写成：
$$ (x_1 - c_1)^2 + (x_2 - c_2)^2 + \cdots + (x_n - c_n)^2 = R^2.$$
可以看到，$n=2,3$时，上式就分别是圆和球面的方程。
$\mathbb{R}^n$中到某点距离为定值的点的集合，一般称为超球面或$n-1$维球面，
记为$\mathcal{S}^{n-1}(\mathbb{R})$.

平面几何中，三角形是最基础的多边形。各种多边形都可以看作是三角形拼合起来的。
立体几何中，最基础的多面体则是四面体。各种多面体可以看作是四面体拼接的结果。

三角形可以由两个（不共线的）向量定义。同样地，四面体可以用三个（不共面的）向量定义。
在$n$维空间中，我们也可以类似定义这样的几何对象。

设$\mathcal{H} = \mathbf{u}_1, \mathbf{u}_2, \cdots , \mathbf{u}_n$是一组$n$个直无关的向量，考虑以下集合：
$$ \mathfrak{C}_n (\mathcal{H}) = \left\{ \theta_1\mathbf{u}_1 + \theta_2\mathbf{u}_2 + \cdots + \theta_n\mathbf{u}_n \, \middle\vert \, \theta_1, \theta_2, \cdots, \theta_n \geqslant 0, \,\, \sum_{i=1}^n \theta_i = 1 \right\}$$
我们把$\mathfrak{C}_n$称为$n$阶\textbf{单纯形}。2阶单纯形就是三角形。比如，设$\mathbf{u}_1 =(0,1), \mathbf{u}_2 = (1,0)$，
就得到以$(0,0), (0,1), (1,0)$为顶点的三角形。注意到原点必然是顶点，另外两个顶点就是$\mathbf{u}_1$和$\mathbf{u}_2$。
同理可以用三个向量构造三维空间中的四面体。同样，原点是它的一个顶点，其他顶点分别是设定的向量。

如果在定义中增加一个平移向量$\mathbf{u}_0$，可以得到空间中任意位置的单纯形。
$$ \mathfrak{C}_n (\mathcal{H}) + \mathbf{u}_0 = \left\{ \mathbf{u}_0 + \theta_1\mathbf{u}_1 +  \cdots + \theta_n\mathbf{u}_n \, \middle\vert \, \theta_1,  \cdots, \theta_n \geqslant 0, \,\, \sum_{i=1}^n \theta_i = 1 \right\}$$
比如在前面三角形的例子中增加平移向量$(-0.4, -0.5)$，就得到以$(-0.4, -0.5)$，$(-0.4, 0.5)$，$(0.6, -0.5)$为顶点的三角形。

\section{复空间的内积}
\subsection{从复数到实数}
前面我们介绍了实空间的内积。对实空间来说，定义内积和模是很自然的，因为它的系数域是实数。
对复空间来说，我们能否用同样的方式定义内积和模呢？

问题在于，我们希望内积导出的模是一个非负实数。
对于$\mathbb{R}^n$来说，我们用平方和处理了这个问题。
对于复数来说，复数的平方和未必是非负的，甚至未必是实数。
所以，如果按同样的方式定义模，我们的“长度”就没法比较大小了，失去了意义。

怎样的复数是实数呢？等于自己的共轭的复数是实数。我们利用共轭来解决这个问题。我们给复空间的内积定义如下：
\begin{df}\label{df:7_3}
    给定空间$\mathbb{C}^n$，从$\mathbb{C}^n \times \mathbb{C}^n$映射到$\mathbb{C}$的映射如果满足以下条件，就称为$\mathbb{C}^n$上的\textbf{内积}。
    \begin{itemize}
        \item 共轭对称性：$\forall \,\, \mathbf{u}, \mathbf{v} \,\, \in \,\, \mathbb{C}^n, \,\,\, (\mathbf{v}|\mathbf{u}) = \overline{(\mathbf{u}|\mathbf{v})}.$
        \item 左平直（关于左变量平直）：$\forall \,\, \mathbf{u}, \mathbf{v}, \mathbf{w} \,\, \in \,\, \mathbb{C}^n$，$\forall \,\,a, b \,\, \in \,\, \mathbb{C}$，$ (a\mathbf{u}+b\mathbf{w}|\mathbf{v}) = a(\mathbf{u}|\mathbf{v}) + b(\mathbf{w}|\mathbf{v}).$
        \item 正定性：$\forall \,\, \mathbf{u} \,\, \in \,\, \mathbb{C}^n, \,\,\, (\mathbf{u}|\mathbf{u}) \geqslant 0$。等号成立当且仅当$\mathbf{u}$是零向量。
    \end{itemize}
    
    配置了内积的空间$\mathbb{C}^n$称为\textbf{内积空间}$\mathbb{C}^n$。
\end{df}
同样的，通过同构，我们也定义了一般的有限维复系数平直空间上的内积。

定义中，由共轭对称性可以推出，$(\mathbf{u}|\mathbf{u}) = \overline{(\mathbf{u}|\mathbf{u})}$。
所以$(\mathbf{u}|\mathbf{u})$是实数。这样才能够定义正定性。
复空间的内积左平直，对于右变量，有：
$$\forall \,\, \mathbf{u}, \mathbf{v}, \mathbf{w} \,\, \in \,\, \mathbb{C}^n, \,\,\, \forall \,\,a, b \,\, \in \,\, \mathbb{C}, \,\,\, (\mathbf{u}|a\mathbf{v}+b\mathbf{w}) = \overline{a}(\mathbf{u}|\mathbf{v}) + \overline{b}(\mathbf{u}|\mathbf{w}).$$
读者可以自行验证这一点。这个性质也叫\text{共轭平直}。
由对称性，我们也可以把复空间内积定义为右平直（关于右变量平直），关于左变量共轭平直。
不过，为了和大多数文献接轨，我们还是采用关左平直的定义。这种一边平直另一边共轭平直的性质也称为\textbf{雔直性}，我们说复空间的内积是\textbf{正定雔直型}。

通过如此定义的内积，导出的模为：
$$ N(\mathbf{u}) = \sqrt{(\mathbf{u}|\mathbf{u})} \geqslant 0.$$

具体来说，在$\mathbb{C}^n$里我们也可以定义点积：
$$ \left((u_1, \cdots, u_n)|(v_1, \cdots, v_n)\right) = u_1\overline{v_1} + \cdots + u_n\overline{v_n}.$$
导出的模为：
$$ \| (u_1, \cdots, u_n)\| = \sqrt{u_1\overline{u_1} + \cdots + u_n\overline{u_n}} = \sqrt{\sum_{i=1}^n |u_i|^2}. $$
对应的距离则是：
$$ \mathrm{d} \left((u_1, \cdots, u_n), (v_1, \cdots, v_n)\right)= \sqrt{\sum_{i=1}^n |u_i - v_i|^2}. $$

\subsection{两种内积的区别}
复空间上的内积和实空间的内积一样吗？
实际上，如果按以上标准去定义实空间的内积，结果和之前不会有区别。
因为对于实数而言，它的共轭就是自己，所以共轭平直就是平直，共轭对称就是对称。
我们可以统一用复空间上的内积定义来描述复空间和实空间上的内积。

前面实空间中模和距离的结论对复空间内积导出的模和距离仍然适用，
但我们需要重新审视一下内积不等式。在其证明中，我们考察了关于实数$\lambda$的表达式：
$$ (\mathbf{u}+ \lambda\mathbf{v} |\mathbf{u}+ \lambda\mathbf{v} ) \geqslant 0 $$
现在我们当然把$\lambda$扩展到复数，上式展开后变成：
$$ \|\mathbf{v} \|^2 |\lambda|^2 + \lambda(\mathbf{v} | \mathbf{u}) + \overline{\lambda}(\mathbf{u} | \mathbf{v}) + \|\mathbf{u} \|^2.$$
其中$\|\mathbf{v} \|^2$和$\|\mathbf{u} \|^2$都是实数。而
$$ \overline{\lambda(\mathbf{v} | \mathbf{u})} = \overline{\lambda}(\mathbf{u} | \mathbf{v}).$$
所以
$$\lambda(\mathbf{v} | \mathbf{u}) + \overline{\lambda}(\mathbf{u} | \mathbf{v}) = 2\Re(\left(\lambda(\mathbf{v} | \mathbf{u})\right)) $$
因此也是实数。$\lambda$模长不变的条件下，这个式子什么时候达到最小值呢？是$\lambda$和$(\mathbf{v} | \mathbf{u})$的共轭同方向的时候。设
$$ \lambda = t (\mathbf{u} | \mathbf{v})$$
其中$t$为实数，则
$$ \Re(\left(\lambda(\mathbf{v} | \mathbf{u})\right)) = t| (\mathbf{u} | \mathbf{v})|^2$$
原式变为：
$$ \|\mathbf{v} \|^2 \| (\mathbf{u} | \mathbf{v})\|^2 t^2 + 2| (\mathbf{u} | \mathbf{v})|^2 t + \|\mathbf{u} \|^2 \quad \quad (*)$$
上式为实系数二次整式。它总大于等于0，说明判别式非正。计算判别式即可得出：
$$ | (\mathbf{u} | \mathbf{v})| \leqslant \|\mathbf{u} \| \|\mathbf{v} \|.$$
再来看不等式取等号的情况。不等式取等号，等价于上面的判别式为0，即$(*)$式恰有一个实数根。
这等价于存在（唯一的）$\lambda$使得
$$ (\mathbf{u}+ \lambda\mathbf{v} |\mathbf{u}+ \lambda\mathbf{v} ) = 0 $$
也就等价于$\mathbf{u}$和$\mathbf{v}$直相关。综上，内积不等式对复内积仍然成立。

\section{正交基}
\subsection{直角坐标系}
解析几何中，我们总会建立直角坐标系。
直角坐标系由两个相互垂直的向量构成，一般记为$\mathbf{e}_1 = (1,0)$和$\mathbf{e}_2 = (0,1)$。

一般的内积空间里，我们也可以定义“直角坐标系”，其中需要用到我们刚刚引入的正交概念。

如果一组非零向量两两正交，就称之为一组\textbf{正交}的向量。
如果一组正交的向量中每个向量模都是1，就称之为一组\textbf{归一正交}的向量。
如果一组正交的向量是内积空间的基底，就称之为\textbf{正交基}。
如果一组归一正交的向量是内积空间的基底，就称之为\textbf{幺正基}，简称\textbf{幺正基}。

从解析几何的经验，我们可以想见，如果用归一正交的向量做基底，
我们可以方便地处理各种关于长度和距离的计算。

比如，如果$\mathcal{B} = \mathbf{e}_1, \cdots , \mathbf{e}_n$是普通的基底，给定向量$\mathbf{v}$，
将它在$\mathcal{B}$上分解可能很麻烦。如果$\mathcal{B}$是归一正交基，那么可以容易地通过内积来做到这一点。

\begin{pp}\label{pp:7_3}
    给定幺正基$\mathcal{B} = \mathbf{e}_1, \cdots , \mathbf{e}_n$，向量$\mathbf{v}$在$\mathcal{B}$上分解为：
    $$\mathbf{v} = (\mathbf{v}|\mathbf{e}_1)\mathbf{e}_1 + (\mathbf{v}|\mathbf{e}_2)\mathbf{e}_2 + \cdots + (\mathbf{v}|\mathbf{e}_n)\mathbf{e}_n.$$
    并且有
    $$\|\mathbf{v}\|^2 = |(\mathbf{v}|\mathbf{e}_1)|^2 + |(\mathbf{v}|\mathbf{e}_2)|^2 + \cdots + |(\mathbf{v}|\mathbf{e}_n)|^2.$$
\end{pp}

\begin{proof} % [{\bfseries 证明\ref{pp:7_3}}]\label{pf:7_3}
    设向量$\mathbf{v}$分解为
    $$ \mathbf{v} = v_1\mathbf{e}_1 + v_2\mathbf{e}_2 + \cdots + v_n \mathbf{e}_n, $$
    用$\mathbf{e}_i$与之做内积，计算分量$v_i$：
    $$ (\mathbf{v} | \mathbf{e}_i) = (v_1\mathbf{e}_1 + v_2\mathbf{e}_2 + \cdots + v_n \mathbf{e}_n | \mathbf{e}_i) = v_i \|\mathbf{e}_i\|^2 = v_i. $$
    而$\mathbf{v}$的模平方为：
    \begin{align*}
        \|\mathbf{v}\|^2 &=  (v_1\mathbf{e}_1 + v_2\mathbf{e}_2 + \cdots + v_n \mathbf{e}_n |v_1\mathbf{e}_1 + v_2\mathbf{e}_2 + \cdots + v_n \mathbf{e}_n)  \\
        &= \sum_{1 \leqslant i, j \leqslant n} v_i \overline{v_j} (\mathbf{e}_i | \mathbf{e}_j)  \\
        &= \sum_{i=1}^n | v_i|^2 \|\mathbf{e}_i \|^2 + \sum_{i \neq j} v_i \overline{v_j} (\mathbf{e}_i | \mathbf{e}_j)  \\
        &= \sum_{i=1}^n | v_i|^2 = \sum_{i=1}^n |(\mathbf{v} | \mathbf{e}_i)|^2 
    \end{align*}
\end{proof}


也就是说，如果$\mathbf{e}_1, \cdots , \mathbf{e}_n$是一组归一正交向量的话，就有：
$$ \|a_1\mathbf{e}_1 + \cdots + a_n\mathbf{e}_n\|^2 = a_1^2 + \cdots + a_n^2.$$

归一正交的向量，是否适合做基底呢？比如，是否会直相关呢？
直觉上来说，两两“垂直”的向量，不大可能直相关。我们可以严格证明这个结论。

\begin{pp}\label{pp:7_4}
    一组正交的向量必然直无关。
\end{pp}

\begin{proof} % [{\bfseries 证明\ref{pp:7_4}}]\label{pf:7_4}
    设$\mathbf{e}_1, \mathbf{e}_2, \cdots , \mathbf{e}_n$是一组正交向量。
    如果有$a_1, a_2, \cdots, a_n$使得
    $$a_1\mathbf{e}_1 + a_2\mathbf{e}_2 + \cdots + a_n \mathbf{e}_n = \mathbf{0}$$
    那么，用$\mathbf{e}_i$与之做内积，就得到：
    $$\forall \, 1 \leqslant i \leqslant n, \,\,\, 0 = (\mathbf{0} |\mathbf{e}_i) = ( a_1\mathbf{e}_1 + a_2\mathbf{e}_2 + \cdots + a_n \mathbf{e}_n | \mathbf{e}_i) = a_i \|\mathbf{e}_i\|^2.$$
    所以$a_1 = a_2 = \cdots = a_n = 0$（注意，一组两两正交的非零向量才能称为正交）。
    这就说明，$\mathbf{e}_1, \mathbf{e}_2, \cdots , \mathbf{e}_n$直无关。
\end{proof}


正交向量组总是直无关的，所以只要向量组元素个数等于空间的维数，我们就找到了一组正交基。

\subsection{构建直角坐标系}
具体来说，如何得到一组归一正交的向量呢？

从一组正交的向量不难得到一组归一正交向量，只需将每个向量除以它的模（一个标量）即可。
因为把一个向量乘以一个非零系数$k$，不会改变它和别的向量是否正交这件事。它和另一个向量的内积也变为原来的$k$倍。
原来与它正交的向量仍然正交，原来与它不正交的向量仍然不正交。

所以，怎样从一组（直无关的）向量得到一组正交向量呢？

我们先从两个向量开始。举例来说，设三维实空间$\mathbb{R}^3$里有向量$\mathbf{e}_1 = (1.5,0,0)$和$\mathbf{e}_2 = (0.2,1.42,0)$。
它们的点积是$1.5\cdot 0.2 + 0\times 1.42 + 0\times 0 = 0.3$，不是0。如何得到正交向量呢？
$\mathbf{e}_2$和$\mathbf{e}_1$，表现为它依$\mathbf{e}_1$的投影不是0. 
我们可以把这个投影分量从$\mathbf{e}_2$里减掉，剩下的部分就和$\mathbf{e}_1$正交了。

我们新造一个向量：
$$ \mathbf{e}_2' = \mathbf{e}_2 - \frac{(\mathbf{e}_2|\mathbf{e}_1)}{\|\mathbf{e}_1\|^2}\mathbf{e}_1 = (0.2,1.42,0) - \frac{0.3}{1.5^2} * (1.5,0,0) = (0, 1.42, 0).$$
它和$\mathbf{e}_1$的内积就是：
$$ (\mathbf{e}_2' | \mathbf{e}_1)= (\mathbf{e}_2 - \frac{(\mathbf{e}_2|\mathbf{e}_1)}{\|\mathbf{e}_1\|^2}\mathbf{e}_1|\mathbf{e}_1) = (\mathbf{e}_2|\mathbf{e}_1) - \frac{(\mathbf{e}_2|\mathbf{e}_1)}{\|\mathbf{e}_1\|^2}\|\mathbf{e}_1\|^2 = 0.$$
这说明$\mathbf{e}_2'$和$\mathbf{e}_1$正交，我们得到一组正交向量。 

如果向量不止两个，能否用同样的手法得到正交向量呢？答案是肯定的。
我们把这个方法叫做\textbf{消影法}或\textbf{格拉姆-施密特方法}。

\begin{pp}\label{pp:7_5}
    从一组直无关的向量，可以得到一组（归一）正交的向量。
\end{pp}

\begin{proof} % [{\bfseries 证明\ref{pp:7_5}}]\label{pf:7_5}
    设$\mathbf{e}_1, \mathbf{e}_2, \cdots , \mathbf{e}_n$是一组直无关的向量。
    我们用以下方法构造一组（归一）正交的向量$\mathbf{f}_1, \mathbf{f}_2, \cdots , \mathbf{f}_n$。\\
    第一步，$\mathbf{f}_1 = \mathbf{e}_1$（如果要得到归一正交向量，则除以其模）。\\
    如果第$k$步后我们得到了一组$k$个（归一）正交向量$\mathbf{f}_1, \mathbf{f}_2, \cdots , \mathbf{f}_k$，
    其中每个元素都是$\mathbf{e}_1, \mathbf{e}_2, \cdots , \mathbf{e}_k$的直组合，那么，在第$k+1$步，我们构造以下向量：
    $$ \mathbf{v} = \mathbf{e}_{k+1} - \sum_{j=1}^k \frac{(\mathbf{e}_{k+1}|\mathbf{f}_j)}{\|\mathbf{f}_j\|^2} \mathbf{f}_j.$$
    $ \mathbf{v}$是$\mathbf{e}_{k+1}$减去$\mathbf{e}_1, \mathbf{e}_2, \cdots , \mathbf{e}_{k}$的直组合，
    而$\mathbf{e}_{k+1}\,\notin \, \langle \mathbf{e}_1, \mathbf{e}_2, \cdots , \mathbf{e}_{k} \rangle$，
    因此$\mathbf{v}$不是零向量。\\
    我们证明$ \mathbf{v}$和前$k$个向量都正交。\\
    将$ \mathbf{v}$和其中任一个做内积：
    \begin{align*} 
        (\mathbf{v} | \mathbf{f}_i) &= (\mathbf{e}_{k+1} | \mathbf{f}_i) - \sum_{j=1}^k \frac{(\mathbf{e}_{k+1}|\mathbf{f}_j)}{\|\mathbf{f}_j\|^2} (\mathbf{f}_j | \mathbf{f}_i)  \\
        &= (\mathbf{e}_{k+1} | \mathbf{f}_i) - \frac{(\mathbf{e}_{k+1}|\mathbf{f}_i)}{\|\mathbf{f}_i\|^2} (\mathbf{f}_i | \mathbf{f}_i)  \\
        &= (\mathbf{e}_{k+1} | \mathbf{f}_i) - (\mathbf{e}_{k+1} | \mathbf{f}_i) = 0 
    \end{align*}
    所以，$\mathbf{v}$和前$k$个向量都正交。\\
    如果要得到正交向量，那么直接取$\mathbf{f}_{k+1} = \mathbf{v}$即可。
    如果要得到归一正交向量，则将$\mathbf{v}$除以其模得到$\mathbf{f}_{k+1}$。
    这样我们就得到了一组$k+1$个（归一）正交向量，每个元素都是原本前$k+1$个向量的直组合。\\
    以上步骤从$k=1$逐步执行到$k=n$，就得到一组$n$个（归一）正交向量。

\end{proof}


对每个有限维内积空间，我们选取一组基。由于基底必然是直无关的，
通过以上方法，可以构造一组元素个数相等的（归一）正交向量。
这组（归一）正交向量是直无关的，从而构成了空间的一组（归一）正交基。
这说明，从内积空间的任一组基出发，都能得到一组（归一）正交基。

不仅如此，一般平直空间里直无关向量可以扩充为基底，内积空间中，
一组（归一）正交向量也可以扩充为（归一）正交基。这是因为（归一）正交向量总是直无关的。
因此可以先将其扩充为基底，然后通过消影法产生（归一）正交基。
注意观察以上证明可以发现，如果前$k$个向量已经（归一）正交，那么前$k$步并不会改变它们的值。
所以（归一）正交向量的扩充从狭义上讲也是成立的。

\subsection{正交基的兼容性}
上一章中，我们认识到，某些条件下，自同态拥有三角基（在此基上的矩阵是上三角矩阵）。
在内积空间里，能否将这个结论强化为（归一）正交基呢？从直觉上来看，
我们可以用消影法处理三角基。由于消影法前$k$步只涉及前$k$个向量，
可以认为，三角基的性质不会被“弄乱”。

\begin{pp}\label{pp:7_6}
    如果自同态$T$在基底$\mathcal{B} = (\mathbf{e}_1, \mathbf{e}_2, \cdots , \mathbf{e}_n)$上的矩阵$[T]_\mathcal{B}$是上三角矩阵，
    那么可以构造一组归一正交基，使得$T$在其上的矩阵仍是上三角矩阵。
\end{pp}

\begin{proof} % [{\bfseries 证明\ref{pp:7_6}}]\label{pf:7_6}
    $[T]_\mathcal{B}$是上三角矩阵，等价于
    $$\forall \, i, \,\,\, T(\mathbf{e}_i) \in \langle \mathbf{e}_1, \cdots , \mathbf{e}_i \rangle.$$ 
    而我们知道，消影法中，$\mathbf{f}_i \in \langle \mathbf{e}_1, \cdots , \mathbf{e}_i \rangle$，
    且$\langle \mathbf{f}_1, \cdots , \mathbf{f}_i \rangle = \langle \mathbf{e}_1, \cdots , \mathbf{e}_i \rangle.$ 
    所以，用消影法处理$\mathcal{B}$得到的新基底$(\mathbf{f}_1, \mathbf{f}_2, \cdots , \mathbf{f}_n)$，仍然满足
    $$\forall \, i, \,\,\, T(\mathbf{f}_i) \in \langle \mathbf{f}_1, \cdots , \mathbf{f}_i \rangle.$$
    这说明$T$在新基底上的矩阵仍是上三角矩阵。

\end{proof}


从此可以推出，有限维复内积空间的自同态，总能在某个幺正基上表示成上三角矩阵。

\section{正交空间}
\subsection{正交子空间和正交补}
不久前我们定义了集合之间的正交关系：集合$U$与$V$正交（简记为$U\perp V$），当且仅当
$$ \forall \,\, \mathbf{u}\,\, \in U, \,\, \mathbf{v} \,\, \in V, \,\,\, (\mathbf{u} | \mathbf{v}) = 0.$$
本节我们进一步研究这个概念。

正交的本质是内积为0. 而内积运算是双平直（共轭平直）的。
所以，容易推出有限维内积空间的集合之间的正交关系有以下的基本性质：
\begin{itemize}
    \item 如果集合$U$与集合$V$正交，那么$U$也和$V$生成的子空间正交。
    \item 两个子空间的基底正交，则这两个子空间正交。
    \item 两个集合正交，则其交集为空集或$\{\mathbf{0}\}$.
    \item 两个子空间正交，则其交集为$\{\mathbf{0}\}$.
\end{itemize}

这说明，如果我们把正交关系应用到子空间上，就可以将空间的直和分解进一步加强：
\begin{pp}\label{pp:7_7}
    给定有限维内积空间$\mathbb{V}$及其子空间$\mathbb{U}$，存在唯一子空间$\mathbb{W}$，
    使$\mathbb{U}$与$\mathbb{W}$正交，且
    $$ \mathbb{V} = \mathbb{U} \oplus \mathbb{W}.$$
    $\mathbb{W}$称为子空间$\mathbb{U}$的\textbf{正交补}，记为$\mathbb{U}^{\perp}$。
\end{pp}

\begin{proof} % [{\bfseries 证明\ref{pp:7_7}}]\label{pf:7_7}
    首先证明存在性：设$\mathbb{U}$的维数为$k$。选$\mathbb{U}$的一组幺正基$(\mathbf{e_1}, \cdots, \mathbf{e}_k)$，
    将其扩充为$\mathbb{V}$的一组幺正基：$\mathcal{B} = (\mathbf{e_1}, \cdots, \mathbf{e}_k, \mathbf{e}_{k+1}, \cdots , \mathbf{e}_n)$。
    则后$n-k$个向量$(\mathbf{e}_{k+1}, \cdots , \mathbf{e}_n)$可以生成一个子空间$\mathbb{W}$。\\
    后$n-k$个向量直无关，构成$\mathbb{W}$的基底，所以$\mathbb{W}$的维数是$n-k$。
    两组向量$(\mathbf{e}_{k+1}, \cdots , \mathbf{e}_n)$和$(\mathbf{e_1}, \cdots, \mathbf{e}_k)$正交，
    所以和$\mathbb{U}$正交，从而$\mathbb{U}$和$\mathbb{W}$正交。易得
    $$ \mathbb{V} = \mathbb{U} \oplus \mathbb{W}.$$

    然后证明唯一性：假设有两个子空间$\mathbb{W}_1$和$\mathbb{W}_2$满足题设条件。
    设有$\mathbf{x} \in \mathbb{W}_1$。由于
    $$ \mathbb{V} = \mathbb{U} \oplus \mathbb{W}_2,$$
    可将$\mathbf{x}$写为$\mathbf{x} = \mathbf{u} + \mathbf{w}_2$，其中$\mathbf{u} \in \mathbb{U}$，$\mathbf{w}_2 \in \mathbb{W}_2$。

    设$(\mathbf{e_1}, \cdots, \mathbf{e}_k)$为$\mathbb{U}$的一组幺正基。
    $\mathbf{u}$可以分解为：
    $$ \mathbf{u} = (\mathbf{u}|\mathbf{e}_1)\mathbf{e}_1 + \cdots + (\mathbf{u}|\mathbf{e}_k)\mathbf{e}_k.$$
    而由于$\mathbf{x} \in\mathbb{W}_1$，$\mathbf{w}_2 \in\mathbb{W}_2$，两个子空间都与$\mathbb{U}$正交。
    所以，
    $$ \forall \, 1 \leqslant i \leqslant k, \,\, (\mathbf{u}|\mathbf{e}_i) = (\mathbf{x} - \mathbf{w}_2|\mathbf{e}_i) = (\mathbf{x}|\mathbf{e}_i) - (\mathbf{w}_2|\mathbf{e}_i) = 0 - 0 = 0.$$
    这说明$\mathbf{u}$是零向量，$\mathbf{x} = \mathbf{w}_2 \in \mathbb{W}_2$，从而$\mathbb{W}_1 \subseteq \mathbb{W}_2$.\\
    由对称性，同理可得$\mathbb{W}_2 \subseteq \mathbb{W}_1$. 所以两个子空间相等。这就证明了唯一性。

\end{proof}


我们曾经定义过子空间的补空间。补空间并不是唯一的。
以带直角坐标系$Oxyz$的三维空间$\mathbb{R}^3$为例子，$xOy$平面是其子空间，
它的补空间可以是竖直的直线$\{(0,0,t) | t \in \mathbb{R}\}$（$z$轴所在直线），
也可以是“斜线”$\{(0,1.12t,t) | t \in \mathbb{R}\}$或$\{(0.15t,0.2t,t) | t \in \mathbb{R}\}$等等。
但如果规定了补空间必须和$xOy$平面正交，那么就只有$\{(0,0,t) | t \in \mathbb{R}\}$满足条件。

从证明里我们可以看到，这是因为正交的条件要求直线在$xOy$上的投影分量为零向量。
“没有影子”的直线只有一条，就是垂直于$xOy$平面的$z$轴所在直线。

对于一般的集合，我们也可以定义其正交补：$U^{\perp} = \{\mathbf{v} \in \mathbb{V} | \mathbf{v} \perp \mathbb{U}\}$.

集合的正交补总是子空间，且为集合生成的子空间的正交补。
\begin{pp}\label{pp:7_7_1}
    设$\mathbb{U}$是空间$\mathbb{V}$的子空间，则
    $$\left(\mathbb{U}^\perp\right)^\perp = \mathbb{U}.$$
\end{pp}
\begin{proof} % [{\bfseries 证明\ref{pp:7_7_1}}]\label{pf:7_7_1}
    首先证明$\mathbb{U}\subseteq \left(\mathbb{U}^\perp\right)^\perp$。\\
    给定向量$\mathbf{u} \in \mathbb{U}$，
    $$\forall \mathbf{w} \, \in \, \mathbb{U}^\perp, \,\,\, (\mathbf{u} | \mathbf{w}) = 0.$$
    所以$\mathbf{u} \in \left(\mathbb{U}^\perp\right)^\perp$.\\
    再证明$\left(\mathbb{U}^\perp\right)^\perp \subseteq \mathbb{U}.$\\
    给定向量$\mathbf{u} \in \left(\mathbb{U}^\perp\right)^\perp$，
    由于$\mathbb{V} = \mathbb{U} \oplus \mathbb{U}^\perp$，$\mathbf{u}$可以写成：
    $$ \mathbf{u} = \mathbf{u}_1 + \mathbf{u}_2 \,\,\, \mathbf{u}_1 \,\in \, \mathbb{U}, \,\,\, \mathbf{u}_2 \,\in \, \mathbb{U}^\perp$$
    据前面证明的部分，$\mathbf{u}_1 \,\in \, \mathbb{U} \subseteq \left(\mathbb{U}^\perp\right)^\perp$，所以
    $$ \mathbf{u}_2 = \mathbf{u} - \mathbf{u}_1 \in \left(\mathbb{U}^\perp\right)^\perp.$$
    因此，$\mathbf{u}_2 \,\in \, \mathbb{U}^\perp \cap \left(\mathbb{U}^\perp\right)^\perp$。
    另一方面，$\mathbb{U}^\perp$也是$\mathbb{V}$的子空间，所以
    $$ \mathbb{V} = \mathbb{U}^\perp \oplus \left(\mathbb{U}^\perp\right)^\perp.$$
    因此$\mathbb{U}^\perp \cap \left(\mathbb{U}^\perp\right)^\perp = \{\mathbf{0}\}$。
    这说明$\mathbf{u}_2 = \mathbf{0}$，于是$\mathbf{u} = \mathbf{u}_1 \in  \mathbb{U}$。

\end{proof}


\subsection{向量到子空间的距离}
前文中我们定义了向量沿另一向量的正交投影。向量$\mathbf{u}$可以分解为沿非零向量$\mathbf{v}$的投影分量，和正交于$\mathbf{v}$的投影余量。
投影余量的长度，其实就是$\mathbf{u}$作为空间中一点到$\mathbf{v}$所在直线的距离。

对于子空间，我们也可以定义向量在上面的正交投影，以及向量到子空间的距离。

\begin{df}\label{df:7_4}
    给定有限维内积空间，向量$\mathbf{x}$在其子空间$\mathbb{U}$的正交投影，
    指满足$\mathbf{x} - \mathbf{u} \perp \mathbf{u}$的$\mathbb{U}$中非零向量$\mathbf{u}$。\\
    $\mathbf{x}$到子空间$\mathbb{U}$的距离，指
    $$\min_{\mathbf{u}\in\mathbb{U}} \|\mathbf{x} - \mathbf{u}\|$$
\end{df}

向量在子空间上的正交投影，和我们之前研究过的投影映射有关。
\begin{pp}\label{pp:7_8}
    设$\mathbb{U}$是空间$\mathbb{V}$的子空间，则向量$\mathbf{x}$在$\mathbb{U}$上的正交投影为$p_\mathbb{U}(\mathbf{x})$，
    其中$p_\mathbb{U}$是$\mathbb{U}$上平行于$\mathbb{U}^\perp$的投影映射，称为关于$\mathbb{U}$的正交投影映射。
\end{pp}

\begin{proof} % [{\bfseries 证明\ref{pp:7_8}}]\label{pf:7_8}
    考虑$\mathbb{U}$上平行于$\mathbb{U}^\perp$的投影映射$p_\mathbb{U}$，对任意向量$\mathbf{x}$，$\mathbf{x}$可以分解为
    $$ \mathbf{x} = \mathbf{u} + \mathbf{v}, \quad \mathbf{u} \in \mathbb{U}, \,\,\, \mathbf{v} \in \mathbb{U}^\perp . $$
    按定义，$p_\mathbb{U}(\mathbf{x}) = \mathbf{u}$，因此，
    $$ \mathbf{x} - p_\mathbb{U}(\mathbf{x}) = \mathbf{x} - \mathbf{u} = \mathbf{v},$$
    而按照定义，$\mathbf{v} \perp \mathbf{u}$。因此，$\mathbf{x}$在$\mathbb{U}$上的正交投影为$p_\mathbb{U}(\mathbf{x})$。

\end{proof}


由于$\mathbf{x} - p_\mathbb{U}(\mathbf{x}) \perp p_\mathbb{U}(\mathbf{x})$，所以由勾股定理有
$$ \| \mathbf{x} - p_\mathbb{U}(\mathbf{x})  \|^2 + \| p_\mathbb{U}(\mathbf{x}) \|^2 = \| \mathbf{x} \|^2.  $$
所以$\| p_\mathbb{U}(\mathbf{x}) \| \leqslant \| \mathbf{x} \|$，等号当且仅当投影余量为零向量时取得。

显然，向量到子空间的距离可以推广为向量到内积空间中一般集合或形体的距离，乃至集合和形体之间的距离。

从上面的等式出发，我们可以得到一个关于向量到子空间距离的重要结论。
对任意$\mathbf{u} \in \mathbb{U}$，考虑向量$\mathbf{x} - \mathbf{u}$。我们将它改写为：
$$ \mathbf{x} - \mathbf{u} = \left(\mathbf{x} - p_\mathbb{U}(\mathbf{x})\right) + \left(p_\mathbb{U}(\mathbf{x}) - \mathbf{u}\right). $$
注意到$\mathbf{x} - p_\mathbb{U}(\mathbf{x}) \in \mathbb{U}^\perp$，而$p_\mathbb{U}(\mathbf{x}) - \mathbf{u} \in \mathbb{U}$（因为$p_\mathbb{U}(\mathbf{x})$和$\mathbf{u}$都是$\mathbb{U}$中元素），
所以两者正交。由勾股定理（命题\ref{pp:7_2}）：
$$ \|\mathbf{x} - \mathbf{u}\|^2 = \|\mathbf{x} - p_\mathbb{U}(\mathbf{x})\|^2 + \|p_\mathbb{U}(\mathbf{x}) - \mathbf{u}\|^2. $$
由于$\|p_\mathbb{U}(\mathbf{x}) - \mathbf{u}\|^2 \geqslant 0$，我们立即得到：
$$ \|\mathbf{x} - \mathbf{u}\|^2 \geqslant \|\mathbf{x} - p_\mathbb{U}(\mathbf{x})\|^2. $$
等号成立当且仅当$\|p_\mathbb{U}(\mathbf{x}) - \mathbf{u}\|^2 = 0$，也就是说$p_\mathbb{U}(\mathbf{x}) = \mathbf{u}$。

\begin{pp}\textbf{最佳逼近定理}\label{pp:7_9}
设$\mathbb{U}$是有限维内积空间$\mathbb{V}$的子空间，$\mathbf{x} \in \mathbb{V}$。则
$$ \forall \,\, \mathbf{u} \,\, \in \,\, \mathbb{U}, \quad \|\mathbf{x} - p_\mathbb{U}(\mathbf{x})\| \leqslant \|\mathbf{x} - \mathbf{u}\|. $$
等号成立当且仅当$\mathbf{u} = p_\mathbb{U}(\mathbf{x})$。
换句话说，正交投影$p_\mathbb{U}(\mathbf{x})$是$\mathbb{U}$中距离$\mathbf{x}$最近的向量，而且是唯一的最近向量。
\end{pp}

这个结论直观上很好理解：从空间中一点到子空间的最短路径，是垂直于子空间的那条线段。想象空中的太阳到大地的光线，其中最短的是垂直于地面的线段。
在平面$\mathbb{R}^2$中，从一点到一条过原点的直线的最短距离是垂线段的长度。
在立体空间$\mathbb{R}^3$中，从一点到一个过原点的平面的最短距离也是垂线段的长度。
最佳逼近定理把这一事实推广到了任意有限维内积空间的子空间上。

有了最佳逼近定理，我们可以给出向量到子空间距离的明确表达式：
$$ \mathrm{d}(\mathbf{x}, \mathbb{U}) = \|\mathbf{x} - p_\mathbb{U}(\mathbf{x})\|. $$
结合勾股定理，还有：
$$ \mathrm{d}(\mathbf{x}, \mathbb{U})^2 = \|\mathbf{x}\|^2 - \|p_\mathbb{U}(\mathbf{x})\|^2. $$

如何计算正交投影呢？如果已知$\mathbb{U}$的一组幺正基$(\mathbf{e}_1, \cdots, \mathbf{e}_k)$，
那么按照命题\ref{pp:7_3}的分解公式，$\mathbf{x}$在$\mathbb{U}$上的正交投影为：
$$ p_\mathbb{U}(\mathbf{x}) = (\mathbf{x}|\mathbf{e}_1)\mathbf{e}_1 + (\mathbf{x}|\mathbf{e}_2)\mathbf{e}_2 + \cdots + (\mathbf{x}|\mathbf{e}_k)\mathbf{e}_k. $$
如果$(\mathbf{e}_1, \cdots, \mathbf{e}_k, \mathbf{e}_{k+1}, \cdots, \mathbf{e}_n)$是$\mathbb{V}$的幺正基，
其中前$k$个向量构成$\mathbb{U}$的幺正基，那么投影余量为：
$$ \mathbf{x} - p_\mathbb{U}(\mathbf{x}) = (\mathbf{x}|\mathbf{e}_{k+1})\mathbf{e}_{k+1} + \cdots + (\mathbf{x}|\mathbf{e}_n)\mathbf{e}_n. $$
于是，向量到子空间的距离为：
$$ \mathrm{d}(\mathbf{x}, \mathbb{U}) = \sqrt{|(\mathbf{x}|\mathbf{e}_{k+1})|^2 + \cdots + |(\mathbf{x}|\mathbf{e}_n)|^2}. $$

\subsection{最小二乘法}

最佳逼近定理有一个重要的应用：处理无解的一次方程组。
第四章中，我们讨论了一次方程组$A\mathbf{x} = \mathbf{b}$的解的存在性。
当$\mathbf{b} \notin \Img A$时，方程组无解。
在实际问题中，我们常常遇到这种情况：观测数据不完全满足某个模型，
或者说，方程的个数多于未知数的个数（$m > n$），方程组"过度约束"。
这时，我们退而求其次，不要求$A\mathbf{x} = \mathbf{b}$严格成立，
而是寻找一个向量$\mathbf{x}_0$，使得$A\mathbf{x}_0$尽可能接近$\mathbf{b}$。
"尽可能接近"的标准，就是让误差向量$A\mathbf{x}_0 - \mathbf{b}$的模最小。
这就是\textbf{最小二乘问题}。

更精确地说：
\begin{df}\label{df:7_5}
给定$m \times n$实矩阵$A$和向量$\mathbf{b} \in \mathbb{R}^m$，
一次方程组$A\mathbf{x} = \mathbf{b}$的\textbf{最小二乘解}是使得
$$ \|A\mathbf{x}_0 - \mathbf{b}\| = \min_{\mathbf{x} \in \mathbb{R}^n} \|A\mathbf{x} - \mathbf{b}\| $$
的向量$\mathbf{x}_0 \in \mathbb{R}^n$。
\end{df}

从几何的角度来看，$A\mathbf{x}$遍历的是$A$的像空间$\Img A$，也就是$A$的列空间。
当$\mathbf{b} \notin \Img A$时，我们在$\Img A$中寻找距离$\mathbf{b}$最近的点。
根据最佳逼近定理，这个最近的点就是$\mathbf{b}$在$\Img A$上的正交投影$p_{\Img A}(\mathbf{b})$。
于是，最小二乘问题等价于：找$\mathbf{x}_0$使得
$$ A\mathbf{x}_0 = p_{\Img A}(\mathbf{b}). $$

正交投影的特征是投影余量与子空间正交。也就是说：
$$ \mathbf{b} - A\mathbf{x}_0 \perp \Img A. $$
而$\Img A$由$A$的列向量$\mathbf{a}_1, \cdots, \mathbf{a}_n$生成。
所以，上式等价于：
$$ \forall \,\, 1 \leqslant i \leqslant n, \quad (\mathbf{b} - A\mathbf{x}_0 \,|\, \mathbf{a}_i) = 0. $$
在实空间中，内积就是点积：$(\mathbf{u}|\mathbf{v}) = \mathbf{u}^{\tr}\mathbf{v}$。
所以上述条件等价于：
$$ \forall \,\, 1 \leqslant i \leqslant n, \quad \mathbf{a}_i^{\tr}(\mathbf{b} - A\mathbf{x}_0) = 0. $$
把$n$个等式合写为矩阵形式。注意到$\mathbf{a}_i$是$A$的第$i$列，所以$\mathbf{a}_i^{\tr}$是$A^{\tr}$的第$i$行。
$n$个等式合起来就是：
$$ A^{\tr}(\mathbf{b} - A\mathbf{x}_0) = \mathbf{0}. $$
整理得到：
$$ A^{\tr} A \mathbf{x}_0 = A^{\tr} \mathbf{b}. $$
我们把上式称为一次方程组$A\mathbf{x} = \mathbf{b}$的\textbf{法方程}
%\footnote{也称为正规方程。对于复空间，将$A^{\tr}$替换为共轭转置$A^H$即可。}。

\begin{pp}\label{pp:7_10}
一次方程组$A\mathbf{x} = \mathbf{b}$的最小二乘解，是法方程$A^{\tr} A \mathbf{x} = A^{\tr} \mathbf{b}$的解。
\end{pp}

\begin{proof} % [{\bfseries 证明\ref{pp:7_10}}]
设$\mathbf{x}_0$是最小二乘解。按照定义，$A\mathbf{x}_0$是$\mathbf{b}$在$\Img A$上的正交投影。
因此，$\mathbf{b} - A\mathbf{x}_0 \perp \Img A$。
对$\Img A$中任意向量$A\mathbf{v}$（$\mathbf{v} \in \mathbb{R}^n$），有：
$$ (A\mathbf{v} \,|\, \mathbf{b} - A\mathbf{x}_0) = 0. $$
用点积的语言写出来，就是：
$$ (A\mathbf{v})^{\tr}(\mathbf{b} - A\mathbf{x}_0) = \mathbf{v}^{\tr} A^{\tr}(\mathbf{b} - A\mathbf{x}_0) = 0. $$
由于上式对任意$\mathbf{v} \in \mathbb{R}^n$成立，所以$A^{\tr}(\mathbf{b} - A\mathbf{x}_0) = \mathbf{0}$，
即$A^{\tr} A \mathbf{x}_0 = A^{\tr} \mathbf{b}$。\\
反过来，如果$\mathbf{x}_0$满足$A^{\tr} A \mathbf{x}_0 = A^{\tr} \mathbf{b}$，
则$A^{\tr}(\mathbf{b} - A\mathbf{x}_0) = \mathbf{0}$。
对$\Img A$中任意向量$A\mathbf{v}$：
$$ (A\mathbf{v} \,|\, \mathbf{b} - A\mathbf{x}_0) = \mathbf{v}^{\tr} A^{\tr}(\mathbf{b} - A\mathbf{x}_0) = \mathbf{v}^{\tr} \mathbf{0} = 0. $$
这说明$\mathbf{b} - A\mathbf{x}_0 \perp \Img A$，因此$A\mathbf{x}_0$是$\mathbf{b}$在$\Img A$上的正交投影。
由最佳逼近定理，$\mathbf{x}_0$是最小二乘解。
\end{proof}

法方程$A^{\tr} A \mathbf{x} = A^{\tr} \mathbf{b}$是一个$n \times n$的一次方程组。
它是否总有解呢？最小二乘解是否唯一呢？这取决于$A^{\tr} A$是否可逆。

\begin{pp}\label{pp:7_11}
设$A$为$m \times n$实矩阵。则$A^{\tr} A$可逆，当且仅当$A$的列向量直无关（即$A$的容为$n$）。
\end{pp}

\begin{proof} % [{\bfseries 证明\ref{pp:7_11}}]
首先证明：如果$A$的列向量直无关，则$A^{\tr} A$可逆。
设$\mathbf{v} \in \Ker(A^{\tr} A)$，则$A^{\tr} A \mathbf{v} = \mathbf{0}$。
用$\mathbf{v}^{\tr}$左乘等式两边，得到：
$$ \mathbf{v}^{\tr} A^{\tr} A \mathbf{v} = (A\mathbf{v})^{\tr}(A\mathbf{v}) = \|A\mathbf{v}\|^2 = 0. $$
所以$A\mathbf{v} = \mathbf{0}$，即$\mathbf{v} \in \Ker A$。
而$A$的列向量直无关等价于$A$是单射（见第四章的讨论），即$\Ker A = \{\mathbf{0}\}$。
所以$\mathbf{v} = \mathbf{0}$。这说明$\Ker(A^{\tr} A) = \{\mathbf{0}\}$，$A^{\tr} A$是单射，从而是可逆的。\\
反过来，如果$A$的列向量直相关，则存在非零向量$\mathbf{v}$使得$A\mathbf{v} = \mathbf{0}$。
于是$A^{\tr} A \mathbf{v} = A^{\tr} \mathbf{0} = \mathbf{0}$，即$\mathbf{v} \in \Ker(A^{\tr} A)$且$\mathbf{v}
eq \mathbf{0}$。
所以$A^{\tr} A$不是单射，不可逆。
\end{proof}

于是，当$A$的列向量直无关时（即$A$的容等于$n$，或者说$n \leqslant m$且$A$的列向量直无关），
法方程有唯一解：
$$ \mathbf{x}_0 = (A^{\tr} A)^{-1} A^{\tr} \mathbf{b}. $$
这时，最小二乘解也是唯一的。

当$A$的列向量直相关时（即$A$的容小于$n$），法方程有无穷多组解。
这时，最小二乘解不唯一。但正交投影$p_{\Img A}(\mathbf{b})$仍然是唯一的，
只是方程$A\mathbf{x} = p_{\Img A}(\mathbf{b})$的解不唯一。
换句话说，"最佳逼近点"是唯一的，但到达这个点的"路径"不止一条。

我们来看一个具体的例子。假设我们要用一条直线$y = ax + b$来拟合三个数据点：
$(1, 2)$、$(2, 3)$、$(3, 5)$。
将三个点代入直线方程，得到一次方程组：
$$ \left\{
\begin{aligned}
a + b &= 2 \\
2a + b &= 3 \\
3a + b &= 5
\end{aligned}
\right. $$
写成矩阵形式$A\mathbf{x} = \mathbf{b}$：
$$ \begin{bmatrix} 1 & 1 \\ 2 & 1 \\ 3 & 1 \end{bmatrix} \begin{bmatrix} a \\ b \end{bmatrix} = \begin{bmatrix} 2 \\ 3 \\ 5 \end{bmatrix}. $$
容易验证，这个方程组无解（三个点不共线）。
我们求它的最小二乘解。首先计算法方程中的矩阵和向量：
$$ A^{\tr} A = \begin{bmatrix} 1 & 2 & 3 \\ 1 & 1 & 1 \end{bmatrix} \begin{bmatrix} 1 & 1 \\ 2 & 1 \\ 3 & 1 \end{bmatrix} = \begin{bmatrix} 14 & 6 \\ 6 & 3 \end{bmatrix}, $$
$$ A^{\tr} \mathbf{b} = \begin{bmatrix} 1 & 2 & 3 \\ 1 & 1 & 1 \end{bmatrix} \begin{bmatrix} 2 \\ 3 \\ 5 \end{bmatrix} = \begin{bmatrix} 23 \\ 10 \end{bmatrix}. $$
法方程为：
$$ \begin{bmatrix} 14 & 6 \\ 6 & 3 \end{bmatrix} \begin{bmatrix} a \\ b \end{bmatrix} = \begin{bmatrix} 23 \\ 10 \end{bmatrix}. $$
$A^{\tr} A$的效（行列式）为$14 \times 3 - 6 \times 6 = 6 
eq 0$，所以可逆。解得：
$$ a = \frac{23 \times 3 - 6 \times 10}{6} = \frac{69 - 60}{6} = \frac{3}{2}, \quad b = \frac{14 \times 10 - 6 \times 23}{6} = \frac{140 - 138}{6} = \frac{1}{3}. $$
于是，最小二乘拟合直线为$y = \frac{3}{2}x + \frac{1}{3}$。
误差向量为：
$$ \mathbf{b} - A\mathbf{x}_0 = \begin{bmatrix} 2 \\ 3 \\ 5 \end{bmatrix} - \begin{bmatrix} 1 & 1 \\ 2 & 1 \\ 3 & 1 \end{bmatrix} \begin{bmatrix} \frac{3}{2} \\[2pt] \frac{1}{3} \end{bmatrix} = \begin{bmatrix} 2 - \frac{11}{6} \\[2pt] 3 - \frac{10}{3} \\[2pt] 5 - \frac{29}{6} \end{bmatrix} = \begin{bmatrix} \frac{1}{6} \\[2pt] -\frac{1}{3} \\[2pt] \frac{1}{6} \end{bmatrix}. $$
读者可以验证，这个误差向量与$A$的两个列向量$(1,2,3)^{\tr}$和$(1,1,1)^{\tr}$的点积都是$0$，
也就是说，误差向量确实正交于$\Img A$。

从最佳逼近定理的角度理解这个例子：
向量$\mathbf{b} = (2, 3, 5)$不在$A$的像空间（$\mathbb{R}^3$中由$(1,2,3)^{\tr}$和$(1,1,1)^{\tr}$张成的二维平面）中。
最小二乘解给出的$A\mathbf{x}_0$是$\mathbf{b}$在这个平面上的正交投影。
误差向量$\mathbf{b} - A\mathbf{x}_0$垂直于这个平面，其模就是$\mathbf{b}$到平面的距离。

最后，我们注意到，如果一次方程组$A\mathbf{x} = \mathbf{b}$本身有解，
那么$\mathbf{b} \in \Img A$，正交投影就是$\mathbf{b}$自身，最小误差为$0$。
这时，法方程$A^{\tr} A \mathbf{x} = A^{\tr} \mathbf{b}$的解就是原方程组的解。
所以，最小二乘法是有解情形的自然推广：有解时它给出精确解，无解时它给出"最接近的解"。


\chapter{内积空间的自同态}
上一章我们介绍了配备内积的平直空间。我们自然要问：配备了内积以后，我们研究平直空间的自同态时候，是否有更多的手段和工具了呢？

\section{内积的矩阵}
\subsection{用矩阵表示内积}
研究空间的自同态乃至直映射的时候，我们从直映射对基底的作用出发。
对于内积来说，我们也想从基底的角度展开研究。

给定实空间$\mathbb{V}$的一组基底$ \mathcal{B} = (\mathbf{e}_1, \mathbf{e}_2, \cdots , \mathbf{e}_n )$，
对空间的某个内积$(\,\cdot \,|\, \cdot\,)$来说，如果它在基底上的作用效果确定的话，我们就可以给出它对任意向量的作用效果。
\begin{align*}
    \forall \,\, \mathbf{u} = \sum_{i=1}^n u_i \mathbf{e}_i , & \,\,\mathbf{v} = \sum_{j=1}^n v_j \mathbf{e}_j  \,\,\in \mathbb{V},  \\
    ( \mathbf{u} |  \mathbf{v} ) &= \nji{\sum_{i=1}^n u_i \mathbf{e}_i}{\sum_{j=1}^n v_j \mathbf{e}_j}  \\
    &= \sum_{i=1}^n \sum_{j=1}^n u_i v_j (\mathbf{e}_i|\mathbf{e}_j) 
\end{align*}
我们把上式中的$(\mathbf{e}_i|\mathbf{e}_j)$排列起来，组成一个矩阵，令矩阵的第$i$行第$j$列的元素为$(\mathbf{e}_j|\mathbf{e}_i)$。
记这个矩阵为
$$
S = [s_{i,j}]_{1 \leqslant i, j \leqslant n} = \begin{bmatrix}
    (\mathbf{e}_1|\mathbf{e}_1) & (\mathbf{e}_2|\mathbf{e}_1) & \cdots & (\mathbf{e}_n|\mathbf{e}_1) \\
    (\mathbf{e}_1|\mathbf{e}_2) & (\mathbf{e}_2|\mathbf{e}_2) & \cdots & (\mathbf{e}_n|\mathbf{e}_2) \\
    \vdots & \vdots & \ddots & \vdots \\
    (\mathbf{e}_1|\mathbf{e}_n) & (\mathbf{e}_2|\mathbf{e}_n) & \cdots & (\mathbf{e}_n|\mathbf{e}_n) \\
\end{bmatrix}
$$
计算向量$\mathbf{u}$和$\mathbf{v}$的内积，就是把$\mathbf{u}$和$\mathbf{v}$对应的列向量进行如下计算：
\begin{align*}
    [\mathbf{v}]^{\tr} S [\mathbf{u}] &= \sum_{i=1}^n \sum_{j=1}^n v_i s_{i,j} u_j   \\
    &= \sum_{i=1}^n \sum_{j=1}^n u_j v_i (\mathbf{e}_j |\mathbf{e}_i )  \\
    &= \nji{\sum_{j=1}^n u_j \mathbf{e}_j}{\sum_{i=1}^n v_i \mathbf{e}_i}  \\
    &= ( \mathbf{u}  | \mathbf{v}  ). 
\end{align*}
我们称$S$为内积$(\,\cdot \,|\, \cdot\,)$在$\mathcal{B}$上的矩阵表示。$[\mathbf{v}]^{\tr} S [\mathbf{u}]$就是对应的计算方式\footnote{可以看出，矩阵乘法中向量的顺序和内积中相反。这是因为我们采用了矩阵左乘列向量的表达方式。如果用矩阵右乘行向量的方式，两者的顺序就相同了。要注意矩阵也随之转置。}。

由于实空间的内积有对称性，所以
$$\forall i, j, \quad s_{i,j} = (\mathbf{e}_j |\mathbf{e}_i ) = (\mathbf{e}_i |\mathbf{e}_j ) = s_{j,i}$$
我们把这样的矩阵叫做\textbf{对称矩阵}。

实空间的内积还有正定性质，即对非零向量$\mathbf{u}$，$(\mathbf{u} | \mathbf{u}) > 0$，因此：
$$\forall (u_1, u_2, \cdots, u_n), \quad \sum_{i=1}^n \sum_{j=1}^n u_i s_{i,j} u_j \geqslant 0,$$
且$\sum_{i=1}^n \sum_{j=1}^n u_i s_{i,j} u_j = 0$当且仅当$u_1 = u_2 = \cdots = u_n = 0$。
我们把这样的矩阵叫做\textbf{正定矩阵}。

可以说，对于给定的基底$\mathcal{B}$，实空间的每个内积，都对应一个对称正定矩阵。

对于复空间，我们也可以作类似定义。复空间的内积关于左变量平直，而关于右变量是共轭平直的。
所以
$$ ( \mathbf{u} |  \mathbf{v} ) = \sum_{i=1}^n \sum_{j=1}^n u_i \overline{v_j} (\mathbf{e}_i|\mathbf{e}_j) $$
计算向量$\mathbf{u}$和$\mathbf{v}$的内积，矩阵形式为$[\overline{\mathbf{v}}]^{\tr} S [\mathbf{u}]$。
对应的矩阵满足：$s_{i,j} = (\mathbf{e}_j|\mathbf{e}_i) = \overline{(\mathbf{e}_i|\mathbf{e}_j)} = \overline{s_{j,i}}$。
我们称这种矩阵为\textbf{共轭对称矩阵}\footnote{也称为厄尔米特矩阵或埃尔米特矩阵。}。
我们定义矩阵的共轭转置为以下映射：
$$ [a_{i,j}]_{1 \leqslant i \leqslant n, 1 \leqslant j \leqslant m} \mapsto [\overline{a_{j,i}}]_{1 \leqslant j \leqslant m, 1 \leqslant i \leqslant n}$$
记矩阵$A$的共轭转置为$A^H$。于是$\mathbf{u}$和$\mathbf{v}$的复内积的矩阵形式变为$[\mathbf{v}]^H S [\mathbf{u}]$。

共轭对称矩阵就是等于自身的共轭转置矩阵的矩阵。

按照定义，对称矩阵也是共轭对称矩阵，所以我们可以将实内积和复内积两种情形的写法合并。
我们说，给定基底后，每个实空间和复空间的内积，都对应一个共轭对称的正定矩阵。

以上是内积在一般基底上的表示形式。如果基底本身和这个内积有关，则矩阵表现为特殊的形式。比如，当基底是关于这个内积的正交基时，
对于$i \neq j$，有$(\mathbf{e}_i |\mathbf{e}_j ) = 0$，所以这时候内积的矩阵是对角矩阵。
它的对角线第$i$个元素是$(\mathbf{e}_i |\mathbf{e}_i) = \| \mathbf{e}_i \|^2$是第$i$个基向量模的平方。
如果基底是关于这个内积的幺正基，那么$\| \mathbf{e}_i \|^2 = 1$，内积的矩阵是单位矩阵$\mathrm{I}_n$。
反之，如果某个内积在某个幺正基上的矩阵是单位矩阵，则定义幺正基的内积和这个内积必然等价。

以二维实空间$\mathbb{R}^2$为例。给定内积：$S: \left((x_1, y_1), (x_2, y_2) \right) \mapsto x_1 x_2 + 2 y_1 y_2$。
可以验证，$\mathcal{B}_1 = \left(\mathbf{e}_1 = (1, 0), \mathbf{e}_2 = \left(0, \frac{\sqrt{2}}{2}\right)\right)$是关于内积$S$的幺正基。
于是，可以算得内积$S$在基底$\mathcal{B}$上的矩阵是：
\begin{align*}
    \begin{bmatrix}
        (\mathbf{e}_1 |\mathbf{e}_1) & (\mathbf{e}_2 |\mathbf{e}_1) \\
        (\mathbf{e}_1 |\mathbf{e}_2) & (\mathbf{e}_2 |\mathbf{e}_2)
    \end{bmatrix} &= \begin{bmatrix}
        1\cdot 1 + 2\cdot 0\cdot 0 &  0\cdot 1 + 2 \cdot \frac{\sqrt{2}}{2} \cdot 0 \\
        1\cdot 0 + 2\cdot 0\cdot \frac{\sqrt{2}}{2} & 0\cdot 0 + 2\cdot \frac{\sqrt{2}}{2} \cdot \frac{\sqrt{2}}{2}
    \end{bmatrix}  \\
    &= \begin{bmatrix}
       1 & 0 \\
       0 & 1
    \end{bmatrix} = \mathrm{I}_2. 
\end{align*}
考虑另一组基底$\mathcal{B}_2 = \left(\mathbf{e}_1' = (\frac{\sqrt{3}}{3}, \frac{\sqrt{3}}{3}), \mathbf{e}_2' = \left(\frac{\sqrt{6}}{3}, -\frac{\sqrt{6}}{6}\right)\right)$。
$\mathcal{B}_2$也是关于内积$S$的幺正基。同样地，可以算得内积$S$在基底$\mathcal{B}_2$上的矩阵是：
\begin{align*}
    \begin{bmatrix}
        (\mathbf{e}_1 |\mathbf{e}_1) & (\mathbf{e}_2 |\mathbf{e}_1) \\
        (\mathbf{e}_1 |\mathbf{e}_2) & (\mathbf{e}_2 |\mathbf{e}_2)
    \end{bmatrix} &= \begin{bmatrix}
        \frac{\sqrt{3}}{3}\cdot \frac{\sqrt{3}}{3} + 2\cdot \frac{\sqrt{3}}{3} \cdot \frac{\sqrt{3}}{3} &  \frac{\sqrt{6}}{3}\cdot \frac{\sqrt{3}}{3} + 2\cdot \left(-\frac{\sqrt{6}}{6}\right) \cdot \frac{\sqrt{3}}{3} \\
        \frac{\sqrt{3}}{3}\cdot \frac{\sqrt{6}}{3} + 2\cdot \frac{\sqrt{3}}{3} \cdot \left(-\frac{\sqrt{6}}{6}\right) & \frac{\sqrt{6}}{3}\cdot \frac{\sqrt{6}}{3} + 2\cdot \left(-\frac{\sqrt{6}}{6}\right) \cdot \left(-\frac{\sqrt{6}}{6}\right)
    \end{bmatrix}  \\
    &= \begin{bmatrix}
       1 & 0 \\
       0 & 1
    \end{bmatrix} = \mathrm{I}_2. 
\end{align*}
而内积$S': \left((x_1, y_1), (x_2, y_2) \right) \mapsto 2 x_1 x_2 + y_1 y_2$在$\mathcal{B}_1$上的矩阵是：
\begin{align*}
    \begin{bmatrix}
        (\mathbf{e}_1 |\mathbf{e}_1) & (\mathbf{e}_2 |\mathbf{e}_1) \\
        (\mathbf{e}_1 |\mathbf{e}_2) & (\mathbf{e}_2 |\mathbf{e}_2)
    \end{bmatrix} &= \begin{bmatrix}
        2\cdot 1\cdot 1 + 0\cdot 0 &  2\cdot 0\cdot 1 + \frac{\sqrt{2}}{2} \cdot 0 \\
        2\cdot 1\cdot 0 + 0\cdot \frac{\sqrt{2}}{2} & 2\cdot 0\cdot 0 + \frac{\sqrt{2}}{2} \cdot \frac{\sqrt{2}}{2}
    \end{bmatrix}  \\
    &= \begin{bmatrix}
       2 & 0 \\
       0 & \frac12
    \end{bmatrix}. 
\end{align*}

\subsection{内积和基变换}
如果基底变了，内积的矩阵会有什么变化呢？之前我们讨论过直映射的矩阵在基变换下的形式。首先，我们考察内积空间中的基变换。
设有基底$ \mathcal{B}_1 = (\mathbf{e}_1, \mathbf{e}_2, \cdots , \mathbf{e}_n )$和$ \mathcal{B}_2 = (\mathbf{f}_1, \mathbf{f}_2, \cdots , \mathbf{f}_n )$，
从$\mathcal{B}_1$到$\mathcal{B}_2$的基变换的矩阵定义为
$$ G^{1\rightarrow 2} = \left[g_{i,j}^{1\rightarrow 2}\right]_{1\leqslant i , j \leqslant n} $$
其中
$$ \mathbf{f}_i = \sum_{k=1}^n g_{k,i}^{1\rightarrow 2} \mathbf{e}_k $$
因此，如果两组基底都是幺正基，就有：
$$ \forall i, j, \quad (\mathbf{f}_i | \mathbf{e}_j) = g_{j,i}^{1\rightarrow 2}. $$
所以$ G^{1\rightarrow 2}$又可以写成：
$$ G^{1\rightarrow 2} = [(\mathbf{f}_j | \mathbf{e}_i)]_{1\leqslant i , j \leqslant n} $$
同理，
$$ G^{2\rightarrow 1} = [(\mathbf{e}_j | \mathbf{f}_i)]_{1\leqslant i , j \leqslant n} $$
$ G^{1\rightarrow 2}$的共轭转置为：
$$ \left(G^{1\rightarrow 2} \right)^H = [\overline{(\mathbf{f}_i | \mathbf{e}_j)}]_{1\leqslant i , j \leqslant n} = [(\mathbf{e}_j | \mathbf{f}_i)]_{1\leqslant i , j \leqslant n} = G^{2\rightarrow 1}. $$
另一方面，我们知道$ G^{1\rightarrow 2} G^{2\rightarrow 1} = \mathrm{I}_n $，所以：
$$ G^{1\rightarrow 2} \left(G^{1\rightarrow 2} \right)^H = \mathrm{I}_n $$
也就是说，$G^{1\rightarrow 2}$的转置共轭矩阵是它的逆矩阵。我们把这样的矩阵称为\textbf{归一正交矩阵}，简称\textbf{幺正矩阵}，把幺正基到幺正基的基变换称为\textbf{归一正交基变换}，简称\textbf{幺正基变换}。

幺正基变换的基变换矩阵是幺正矩阵。反之，幺正矩阵总是幺正基之间的基变换矩阵。
\begin{pp}\label{pp:8_1}
    如果矩阵$A = [a_{i,j}]_{1\leqslant i , j \leqslant n}$满足
    $$ A \cdot A^H = \mathrm{I}_n$$
    那么$A$可以表示成两个幺正基之间的基变换矩阵。
\end{pp}
\begin{proof} % [{\bfseries 证明\ref{pp:8_1}}]\label{pf8_1}
    $A$的行向量为
    $$ \forall 1\leqslant i \leqslant n, \quad A_i^- = [a_{i,1}, \cdots a_{i,n}]$$
    在$A^H$中，第$j$列的列向量就是$A_j^-$的元素共轭转置得到的向量。所以相乘之后，$ A \cdot A^H$第$i$行第$j$列的元素就是：
    $$ \sum_{k=1}^n a_{i,k} \overline{a_{j,k}}. $$
    设$\mathbb{F}$为矩阵$A$系数所在的域，考虑$\mathbb{F}^n$的幺正基$ \mathcal{C} = (\mathbf{e}_1, \mathbf{e}_2, \cdots , \mathbf{e}_n )$，构造一组$n$个向量：
    $$ \forall 1\leqslant i \leqslant n, \quad \mathbf{f}_i = \sum_{k=1}^n a_{i,k} \mathbf{e}_k $$
    计算这组向量两两之间的内积（$ \mathcal{C} $对应的内积）：
    $$ (\mathbf{f}_i | \mathbf{f}_j ) = \sum_{k=1}^n a_{i,k} \overline{a_{j,k}} $$
    $ (\mathbf{f}_i | \mathbf{f}_j )$就是$ A \cdot A^H$第$i$行第$j$列的元素，而$ A \cdot A^H = \mathrm{I}_n$。
    所以当且仅当$i = j$时，$ (\mathbf{f}_i | \mathbf{f}_j ) = 1$，其余时候$ (\mathbf{f}_i | \mathbf{f}_j ) = 0$。
    这说明$\mathcal{B} = (\mathbf{f}_1, \mathbf{f}_2, \cdots , \mathbf{f}_n )$是此内积下的幺正基。
    而由其定义知，矩阵$A$是从$ \mathcal{C} $变换到$ \mathcal{B} $的基变换矩阵。

\end{proof}


回到之前的讨论，我们假设$\mathcal{B}_1$和$\mathcal{B}_2$都是幺正基。
记内积在$ \mathcal{B}_1$上的矩阵为$[S]_{\mathcal{B}_1}$，在$ \mathcal{B}_2$上的矩阵为$[S]_{\mathcal{B}_2} $。
给定向量$\mathbf{u}$和$\mathbf{v}$。设它们在两组基上的分解分别是：
\begin{align*}
    \mathbf{u} &= \sum_{i=1}^n u_i \mathbf{e}_i = \sum_{i=1}^n u_i' \mathbf{f}_i  \\
    \mathbf{v} &= \sum_{i=1}^n v_i \mathbf{e}_i = \sum_{i=1}^n v_i' \mathbf{f}_i  
\end{align*}
设$G^{1\rightarrow 2}$为$\mathcal{B}_1$到$\mathcal{B}_2$的基变换矩阵，则$\mathbf{u}$和$\mathbf{v}$在新旧基底上的列向量满足：
$$ [\mathbf{u}]_{\mathcal{B}_2} = G^{2\rightarrow 1} [\mathbf{u}]_{\mathcal{B}_1}, \quad [\mathbf{v}]_{\mathcal{B}_2} = G^{2\rightarrow 1} [\mathbf{v}]_{\mathcal{B}_1} $$
考虑$ ( \mathbf{u} |  \mathbf{v} )$在两个基底上的矩阵形式。
\begin{align*}
    ( \mathbf{u} |  \mathbf{v} ) &= [\mathbf{v}]_{\mathcal{B}_1}^H [S]_{\mathcal{B}_1} [\mathbf{u}]_{\mathcal{B}_1}  \\
    &= [\mathbf{v}]_{\mathcal{B}_2}^H [S]_{\mathcal{B}_2} [\mathbf{u}]_{\mathcal{B}_2}  \\
    &= [\mathbf{v}]_{\mathcal{B}_1}^H \left(G^{2 \rightarrow 1}\right)^H [S]_{\mathcal{B}_2} G^{2 \rightarrow 1} [\mathbf{u}]_{\mathcal{B}_1}  
\end{align*}
我们把$G^{2\rightarrow 1}$记为$[\mathrm{I}]_{\mathcal{B}_1}^{\mathcal{B}_2}$。于是，从上式可以得出：
$$[S]_{\mathcal{B}_1} = \left([\mathrm{I}]_{\mathcal{B}_1}^{\mathcal{B}_2}\right)^H [S]_{\mathcal{B}_2} [\mathrm{I}]_{\mathcal{B}_1}^{\mathcal{B}_2}.$$
注意到$ \left([\mathrm{I}]_{\mathcal{B}_1}^{\mathcal{B}_2}\right)^H = \left([\mathrm{I}]_{\mathcal{B}_1}^{\mathcal{B}_2}\right)^{-1} = [\mathrm{I}]_{\mathcal{B}_2}^{\mathcal{B}_1}$，因此：
$$[S]_{\mathcal{B}_1} = [\mathrm{I}]_{\mathcal{B}_2}^{\mathcal{B}_1} [S]_{\mathcal{B}_2} [\mathrm{I}]_{\mathcal{B}_1}^{\mathcal{B}_2}.$$
同理，有：
$$[S]_{\mathcal{B}_2} = [\mathrm{I}]_{\mathcal{B}_1}^{\mathcal{B}_2} [S]_{\mathcal{B}_1} [\mathrm{I}]_{\mathcal{B}_2}^{\mathcal{B}_1}.$$
这个矩阵形式和直自同态的基变换公式是一样的，但它们的意义并不相同。方阵$[S]_{\mathcal{B}_1}$在这里表示一个内积，
而非直自同态。另外，这里的$\mathcal{B}_1$和$\mathcal{B}_2$都是幺正基。
也只有当它们都是幺正基的时候，矩阵的共轭转置和矩阵的逆相同，使得两种基变换公式一致。

\section{伴随映射}
\subsection{对偶表示定理}
我们介绍过平直空间的对偶空间。给定空间$\mathbb{V}$，它的对偶空间$\mathbb{V}^*$是所有$\mathbb{V}$上直型的集合。
给定$\mathbb{V}$的基底$\mathcal{B} = (\mathbf{e}_1, \mathbf{e}_2, \cdots , \mathbf{e}_n )$，
它的对偶基$\mathcal{B}^* = (\mathbf{e}_1^*, \mathbf{e}_2^*, \cdots , \mathbf{e}_n^*)$是$\mathbb{V}^*$的基底。给定一个直型$\varphi \in \mathbb{V}^*$，它可以写为：
$$ \varphi = \sum_{i=1}^n a_i \mathbf{e}_i^* $$
如果基底$\mathcal{B}$是幺正基，并且对应一个内积$(\,\cdot \,|\, \cdot\,)$，那么，设$ \mathbf{a} = \sum_{i=1}^n \overline{a_i} \mathbf{e}_i $，则：
$$ \forall \mathbf{x} = \sum_{i=1}^n x_i \mathbf{e}_i, \quad \varphi(\mathbf{x}) = \sum_{i=1}^n a_i x_i = (\mathbf{x} | \mathbf{a}). $$
考虑映射$\mathbf{x} \mapsto (\mathbf{x} | \mathbf{a})$，可以验证这是一个直型。对于确定的基底$\mathcal{B}$，我们能够唯一确定向量$\mathbf{a}$。
$$ \mathbf{a} = \sum_{i=1}^n \overline{a_i} \mathbf{e}_i = \sum_{i=1}^n \overline{\varphi(\mathbf{e}_i)} \mathbf{e}_i.$$
以上讨论说明：
\begin{pp}\textbf{对偶表示定理}\label{pp:8_2}
    $\mathbb{V}$的对偶空间$\mathbb{V}^*$中每个元素$\varphi$，都唯一对应向量$\mathbf{a} \in \mathbb{V}$，使得
    $$ \varphi = \mathbf{x} \mapsto (\mathbf{x} | \mathbf{a}). $$
    给定$\mathbb{V}$的一组幺正基$\mathcal{B} = (\mathbf{e}_1, \mathbf{e}_2, \cdots , \mathbf{e}_n )$，
    向量$\mathbf{a}$可以具体写为：
    $$ \mathbf{a} = \sum_{i=1}^n \overline{\varphi(\mathbf{e}_i)} \mathbf{e}_i.$$
\end{pp}

比起对偶表示定理，更容易验证的是，给定$\mathbb{V}$中任意向量$\mathbf{a}$，$\mathbb{V}^*$中都有唯一的元素$\varphi$，使得
$$ \varphi = \mathbf{x} \mapsto (\mathbf{x} | \mathbf{a}). $$
这表示我们可以定义$\mathbb{V}$到$\mathbb{V}^*$的映射:$\chi : \mathbf{a} \mapsto \varphi$。
可以验证，$\chi$是同构映射。而对偶表示定理说明，同构$\chi$的逆映射存在，也是一个同构映射。

\subsection{自同态的“伴”}
给定自同态$A$和向量$\mathbf{v}$，考虑映射$\mathbf{u} \mapsto (A\mathbf{u} | \mathbf{v})$。可以验证，这是一个直型。
根据对偶表示定理，存在唯一向量$\mathbf{a}$，使得这个直型可以表达为$\mathbf{u} \mapsto (\mathbf{u} | \mathbf{a})$。
当自同态$A$不变而$\mathbf{v}$变化的时候，对应的唯一$\mathbf{a}$也产生变化。这个映射关系可以记为：
$$ A^\star : \mathbf{v} \mapsto \mathbf{a}. $$
我们把这个映射称为自同态$A$的\textbf{伴随映射}。按照伴随映射的定义，
$$ \forall \mathbf{u}, \mathbf{v}, \quad (A\mathbf{u} | \mathbf{v}) = (\mathbf{u} | A^\star(\mathbf{v})). $$
\begin{pp}\label{pp:8_3}
    给定直自同态$A$，唯一存在直自同态$A^\star \in \mathcal{L}(\mathbb{V})$，使得
    $$ \forall \mathbf{u}, \mathbf{v}, \quad (A\mathbf{u} | \mathbf{v}) = (\mathbf{u} | A^\star(\mathbf{v})). $$
    $A^\star$称为$A$的伴随映射。
\end{pp}
\begin{proof} % [{\bfseries 证明\ref{pp:8_3}}]\label{pf8_3}
    据之前的讨论，$A^\star$作为映射唯一存在。只需证明$A^\star$是直映射。
    考虑$\mathbf{u}, \mathbf{v}_1, \mathbf{v}_2 \in \mathbb{V}$，$\lambda \in \mathbb{F}$，
    \begin{align*}
        (\mathbf{u} | A^\star(\mathbf{v}_1 + \lambda \mathbf{v}_2)) &= (A\mathbf{u} | \mathbf{v}_1 + \lambda \mathbf{v}_2)  \\
        &= (A\mathbf{u} | \mathbf{v}_1) + \overline{\lambda} (A\mathbf{u} | \mathbf{v}_2)  \\
        &= (\mathbf{u} | A^\star(\mathbf{v}_1)) + \overline{\lambda} (\mathbf{u} | A^\star(\mathbf{v}_2))  \\
        &= (\mathbf{u} | A^\star(\mathbf{v}_1) + \lambda A^\star(\mathbf{v}_2)) 
    \end{align*}
    即
    $$ \forall \mathbf{u}, \mathbf{v}_1, \mathbf{v}_2 \in \mathbb{V}, \quad (\mathbf{u} | A^\star(\mathbf{v}_1 + \lambda \mathbf{v}_2) - (A^\star(\mathbf{v}_1) + \lambda A^\star(\mathbf{v}_2))) = 0.$$
    以上关系对任意向量$\mathbf{u}$总成立，所以
    $$ A^\star(\mathbf{v}_1 + \lambda \mathbf{v}_2) - (A^\star(\mathbf{v}_1) + \lambda A^\star(\mathbf{v}_2)) = \mathbf{0} $$
    这就证明$A^\star$是直映射。

\end{proof}


我们进一步考虑将自同态映射到其伴随映射的映射：
$$ \bigstar : A \mapsto A^\star $$
不难验证，这个映射满足（对任意$S, T \in \mathcal{L}(\mathbb{V})$）：
\begin{enumerate}
    \item $S^\star + T^\star = (S + T)^\star$
    \item $\forall \lambda \in \mathbb{F}, \quad (\lambda T)^\star = \overline{\lambda} T^\star$
    \item $ (T^\star)^\star = T $
    \item $ (T\circ S)^\star = S^\star \circ T^\star$
    \item $ \mathrm{I}^\star = \mathrm{I}$
\end{enumerate}

作为直自同态，我们来看伴随映射的像空间和零空间和原映射的关系。
\begin{pp}\label{pp:8_4}
    设$T$为内积空间$\mathbb{V}$中自同态，则$T$的伴随映射$T^\star$的像空间和零空间满足：
    \begin{enumerate}[label=\alph*.]   
        \item $\Ker T^\star = (\Img T)^\perp$
        \item $\Img T^\star = (\Ker T)^\perp$
        \item $\Ker T = (\Img T^\star)^\perp$
        \item $\Img T = (\Ker T^\star)^\perp$
    \end{enumerate}
\end{pp}
\begin{proof} % [{\bfseries 证明\ref{pp:8_4}}]\label{pf:8_4}
    首先证明\textit{a}。
    $\forall \mathbf{v} \in \Ker T^\star$，按照定义，
    $$ \forall \mathbf{u}, \,\,\,(T\mathbf{u} | \mathbf{v}) = (\mathbf{u} | T^\star\mathbf{v}) = (\mathbf{u} | \mathbf{0}) = 0.  $$
    所以$\mathbf{v} \in (\Img T)^\perp$。反之，$\forall \mathbf{v} \in (\Img T)^\perp$，
    $$ \forall \mathbf{u}, \,\,\, (\mathbf{u} | T^\star\mathbf{v}) = (T\mathbf{u} | \mathbf{v}) = 0.$$
    所以$T^\star\mathbf{v} = \mathbf{0}$，这说明$\mathbf{v} \in \Ker T^\star$。
    于是我们可以推出$\Ker T^\star = (\Img T)^\perp$，\textit{a}命题成立。\\
    由于$ (T^\star)^\star = T $，在\textit{a}命题中用$T^\star$代替$T$，得到$\Ker T = (\Img T^\star)^\perp$，
    这就证明了\textit{c}命题。\\
    另一方面，根据正交补的性质，$\left((\Img T)^\perp\right)^\perp = \Img T$，所以对\textit{a}命题等号两边取正交补，
    可以得到$\Img T = \left((\Img T)^\perp\right)^\perp = (\Ker T^\star)^\perp$，所以\textit{d}命题成立。\\
    类似地，对\textit{c}命题等号两边取正交补，可以得到$(\Ker T)^\perp = \left((\Img T^\star)^\perp\right)^\perp = \Img T^\star$。
    所以\textit{b}命题也成立。\\
\end{proof}


\subsection{伴随映射的矩阵}
说了这么多，为什么要研究自同态的伴随映射呢？我们先来看伴随映射的矩阵和原映射矩阵的关系。
\begin{pp}\label{pp:8_5}
    设$T$为内积空间$\mathbb{V}$中自同态，给定幺正基$\mathcal{B} = (\mathbf{e}_i)_{1 \leqslant i \leqslant n}$，
    则$T$的伴随映射$T^\star$在$\mathcal{B}$上的矩阵是$T$在$\mathcal{B}$上的矩阵的共轭转置。
    $$ [T^\star]_\mathcal{B} = [T]_\mathcal{B}^H $$
\end{pp}
\begin{proof} % [{\bfseries 证明\ref{pp:8_5}}]\label{pf:8_5}
    按照直映射的矩阵的定义，$[T^\star]_\mathcal{B}$的第$i$列的元素是$T^\star \mathbf{e}_i$的$\mathcal{B}$分解。
    也就是说，如果记$[T^\star]_\mathcal{B} = [s_{i,j}]_{1 \leqslant i, j \leqslant n}$，那么
    $$ T^\star \mathbf{e}_i = \sum_{k=1}^n s_{k,i} \mathbf{e}_k. $$
    把两边和$\mathbf{e}_j$作内积，可以得到：
    $$ (T\mathbf{e}_j | \mathbf{e}_i) = (\mathbf{e}_j | T^\star \mathbf{e}_i) =\nji{ \mathbf{e}_j }{ \sum_{k=1}^n s_{k,i} \mathbf{e}_k } = \overline{s_{j,i}}.$$
    而另一方面，如果记$[T]_\mathcal{B} = [t_{i,j}]_{1 \leqslant i, j \leqslant n}$，那么
    $$ (T\mathbf{e}_j | \mathbf{e}_i) = \nji{ \sum_{k=1}^n t_{k,j} \mathbf{e}_k }{ \mathbf{e}_i } = t_{i,j}. $$
    因此，我们得到：$t_{i,j} = \overline{s_{j,i}}$。这就说明，
    $$ [T^\star]_\mathcal{B} = [T]_\mathcal{B}^H $$
\end{proof}

伴随映射的矩阵就是原映射矩阵的共轭转置。

前面的章节里，我们介绍了对偶映射。给定自同态$T\in\mathcal{L}(\mathbf{V})$，
它的对偶映射是$\mathbb{V}^*$上的自同态：$T^* \in \mathcal{L}(\mathbf{V}^*)$。
伴随映射和对偶映射有一个共同之处。考虑自同态$T$在内积对应的幺正基$\mathcal{B}$上的矩阵：$[T]_\mathcal{B}$，
则其伴随映射在同一个基底上的矩阵是它的共轭转置；另一方面，
$T$的对偶映射在幺正基$\mathcal{B}$的对偶基上是它的转置矩阵。如果$\mathbb{V}$是实内积空间，那么共轭转置就是转置矩阵，
于是对偶映射的矩阵和伴随映射的矩阵相同。如果$\mathbb{V}$是复内积空间，那么两者之间差一个共轭变换。

实际上，如果我们考虑映射：$T^* \circ \chi$和$\chi \circ T^\star$，
其中$\chi$是关联相应内积的同构映射。这两个映射都从$\mathbb{V}$映射到$\mathbb{V}^*$。我们可以证明两者相同：
\begin{pp}\label{pp:8_6}
    设$T$为内积空间$\mathbb{V}$中自同态，
    则$T$的伴随映射$T^\star$和它的对偶映射$T^*$满足以下关系：
    $$ T^* \circ \chi = \chi \circ T^\star $$
\end{pp}
\begin{proof} % [{\bfseries 证明\ref{pp:8_6}}]\label{pf:8_6}
    对任意向量$\mathbf{a} \in \mathbb{V}$，分别计算$\left(T^\star \circ \chi\right)(\mathbf{a})$和$\left(\chi \circ T^*\right)(\mathbf{a})$。\\
    一方面，
    $$ \left(T^* \circ \chi\right)(\mathbf{a}) = T^* \left(\chi(\mathbf{a})\right) = \chi(\mathbf{a}) \circ T, $$
    所以$\forall \mathbf{v} \in \mathbb{V}$，
    \begin{align*}
        \left(T^* \circ \chi\right)(\mathbf{a})(\mathbf{v}) &= \left(\chi(\mathbf{a}) \circ T\right) (\mathbf{v})  \\
        &= \chi(\mathbf{a}) \left(T\mathbf{v} \right)  \\
        &= \left(T\mathbf{v} | \mathbf{a}\right). 
    \end{align*}
    另一方面，
    $$ \left(\chi \circ T^\star\right)(\mathbf{a}) = \chi \left(T^\star\mathbf{a}\right), $$
    所以$\forall \mathbf{v} \in \mathbb{V}$，
    \begin{align*}
        \left(\chi \circ T^\star\right)(\mathbf{a})(\mathbf{v}) &= \left(\chi \left(T^\star\mathbf{a}\right)\right) (\mathbf{v})  \\
        &= \left(\mathbf{v} | T^\star\mathbf{a}\right)  \\
        &= \left(T\mathbf{v} | \mathbf{a}\right). 
    \end{align*}
    因此我们证明了
    $$ T^* \circ \chi = \chi \circ T^\star. $$

\end{proof}

这个关系说明，在内积空间中，自同态的伴随映射和对偶映射之间有对应的关系。
由于对应关系中涉及的都是同构，所以伴随映射和对偶映射可以看作是“等同”的。
也因为如此，我们会看到不少著作中将两者用相同的记号表示。
本书中使用了十分相似的记号：$*$和$\star$稍做区分。

\subsection{自伴自同态与谱定理}
有了伴随映射的定义，我们可以定义这么一种自同态：
\begin{df}\label{df:8_2}
    给定内积空间$\mathbb{V}$中自同态$T$，如果$T$的伴随映射$T^*$等于$T$，
    就称它为\textbf{自伴自同态}或\textbf{自伴算子}。换句话说，如果
    $$ \forall \mathbf{u}, \,\, \mathbf{v} \,\, \in \mathbb{V},\,\,\, (T\mathbf{u} | \mathbf{v}) = (\mathbf{u} | T\mathbf{v}), $$
    就称$T$是自伴的，它就是它自己的伴随映射。
\end{df}
显然，自伴自同态的矩阵是共轭对称矩阵，这就和之前内积的矩阵联系上了。考虑内积空间上的另一个内积$(\, \cdot \, | \, \cdot \,)_{S}$，
设它在原内积的某个幺正基$\mathcal{B}$上的矩阵为$S$，$S$是共轭对称的正定矩阵。
我们可以构造一个自同态，使得它在$\mathcal{B}$上的矩阵也是$S$。为了方便，我们将这个自同态记为$S$。于是，
$$ \forall \mathbf{u}, \,\, \mathbf{v} \,\, \in \mathbb{V},\,\,\, (S\mathbf{u} | \mathbf{v}) = \sum_{i=1}^n \sum_{j=1}^n u_i s_{i,j} \overline{v_j},$$
其中的$s_{i,j}$、$u_i$、$v_j$等分别是$S$的对应元素以及$\mathbf{u}$和$\mathbf{v}$在$\mathcal{B}$上的坐标。
可以看到，这个结果和内积$(\, \cdot \, | \, \cdot \,)_{S}$作用于$\mathbf{u}$和$\mathbf{v}$的效果是一样的。
所以，可以用$(\mathbf{u}, \mathbf{v} ) \mapsto (S\mathbf{u} | \mathbf{v}) $代替内积$(\, \cdot \, | \, \cdot \,)_{S}$。
我们把自伴自同态$S$称为内积$(\, \cdot \, | \, \cdot \,)_{S}$对应的自同态。
这样，从某个特定的内积出发，通过研究内积对应的自伴自同态的性质，就可以了解其他内积的性质。

在“本征值与自同态约简”一章中，我们研究了使用本征值把自同态化为更简单的形式的方法。
自伴自同态有什么性质呢？首先，我们来研究它的本征值。

\begin{pp}\label{pp:8_7}
    设$T$为复内积空间$\mathbb{V}$中的自伴自同态，而$\lambda$是它的本征值，
    那么$\lambda$是实数。
\end{pp}
\begin{proof} % [{\bfseries 证明\ref{pp:8_7}}]\label{pf:8_7}
    设$\mathbf{v}$是$\lambda$对应的本征向量，计算$\lambda \|\mathbf{v}\|^2$。
    它是一个标量。另一方面：
    \begin{align*} 
        \lambda \|\mathbf{v}\|^2 &= \langle \lambda \mathbf{v} | \mathbf{v} \rangle = \langle T\mathbf{v} | \mathbf{v} \rangle  \\
        &= \langle \mathbf{v} | T \mathbf{v} \rangle = \langle \mathbf{v} | \lambda \mathbf{v} \rangle  \\
        &= \overline{\lambda} \langle \mathbf{v} | \mathbf{v} \rangle = \overline{\lambda} \|\mathbf{v}\|^2. 
    \end{align*}
\end{proof}

从特例来看，如果自伴自同态的矩阵是对角矩阵\footnote{容易验证这样的矩阵对应自伴自同态。}，那么，
一方面由于它等于它的共轭转置，这说明对角线上的元素等于它自己的共轭，也就是说，它们都是实数。
另一方面，对角线上的元素也正是矩阵乃至自同态的本征值。

如果自伴自同态的矩阵是上三角矩阵，那么同样可以得出对角线上的元素是实数。此外，由于对角线下方都是0，矩阵又等于其共轭转置，所以对角线上方的元素也是0。这说明，自伴自同态的矩阵，只要是上三角矩阵，
就是对角矩阵，从而对角线上的元素就是它的本征值。

命题\ref{pp:6_10}说明，复空间的自同态总能三角化。而命题\ref{pp:7_6}进一步说明，我们可以找到一组幺正基$\mathcal{B}$，
使得复内积空间的自伴自同态在$\mathcal{B}$上是上三角矩阵。
因此，在$\mathcal{B}$上，自伴自同态是对角矩阵，$\mathcal{B}$就是对应的本征向量。我们把这个结论叫做谱定理。
\begin{pp}\textbf{复内积空间的谱定理 }\label{pp:8_8}
    设$T$为复内积空间$\mathbb{V}$中的自伴自同态，则存在由$T$的本征向量构成的幺正基$\mathcal{B}$，
    $T$在$\mathcal{B}$上是对角矩阵。也称$T$可以\textbf{归一正交对角化}，简称\textbf{幺正对角化}。
\end{pp}

对于实内积空间中的自伴自同态$T$，我们可以考虑它对应的对称矩阵在复内积空间$\mathbb{C}^n$中对应的自同态$T^{\mathbb{C}}$。
对称矩阵当然是共轭对称的，所以$T^{\mathbb{C}}$在复内积空间中也是自伴自同态，因此可以对角化。
这说明它的极小归零式$m$是一次式的乘积，且没有重根。

另一方面，命题\ref{pp:8_7}说明$T^{\mathbb{C}}$的本征值都是实数，于是极小归零式$m$的根都是实根。这说明$m$是实系数整式。
而$m(T) = 0$（两者矩阵相同），这说明$T$的极小归零式$m_T$是$m$的因式，因此也是一次式的乘积，且没有重根。
因此，$T$在实内积空间中也可以对角化。于是根据命题\ref{pp:7_6}，它可以幺正对角化。
\begin{pp}\textbf{实内积空间的谱定理 }\label{pp:8_8_0}
    设$T$为实内积空间$\mathbb{V}$中的自伴自同态，则存在由$T$的本征向量构成的幺正基$\mathcal{B}$，
    $T$在$\mathcal{B}$上是对角矩阵。或者说，$T$可以幺正对角化。
\end{pp}

\section{幺正自同构}
\subsection{幺正矩阵的自同构}
上一节中，我们用矩阵的语言讨论了基变换对内积对应矩阵的影响。其中的基变换矩阵称为幺正矩阵，对应幺正基变换。
现在，我们来研究幺正矩阵对应的自同构。

\begin{df}\label{df:8_1}
配备了内积$(\,\cdot\, | \,\cdot\,)$的空间$\mathbb{V}$中的自同态$A$如果满足$A\circ A^\star = \mathrm{I}$，就称它为\textbf{归一正交自同构}或\textbf{归一正交算子}，简称\textbf{幺正自同构}。
\end{df}
幺正自同构在任意幺正基下的矩阵，显然是幺正矩阵。同理，幺正矩阵对应于标准基的自同构，也是幺正自同构。

幺正自同构是$\mathbb{V}$上的同构，它的逆映射就是它的伴随自同态。由\ref{pp:3_3_3}可知，幺正自同构也可以定义为满足$A^\star\circ A = \mathrm{I}$的自同态。两个定义是等价的。

如果$A$有本征值$\lambda$，对应本征向量$\mathbf{u}$，那么：
$$ \| A\mathbf{u}\|^2 = (A\mathbf{u}|A\mathbf{u}) = (\lambda\mathbf{u}| \lambda\mathbf{u}) = |\lambda|^2 \|\mathbf{u}\|^2. $$
而另一方面，
$$ \| A\mathbf{u}\|^2 = (A\mathbf{u}|A\mathbf{u}) = (A^\star A\mathbf{u}|\mathbf{u}) = (\mathbf{u}|\mathbf{u}) = \|\mathbf{u}\|^2. $$
这说明$|\lambda| = 1$。也就是说$A$的本征值都在复平面的单位圆上。如果$A$是实内积空间的幺正自同构，那么它的本征值只能是$\pm 1$。

从上面的等式，我们还能得到$ \| A\mathbf{u}\| = \|\mathbf{u}\|$。这说明向量经过$A$映射后，长度不变。我们把这个性质称为\textbf{保距}性质：$ \| A\mathbf{u} - A\mathbf{v}\| = \|A(\mathbf{u} - \mathbf{v})\| = \|\mathbf{u} - \mathbf{v}\|$。可以证明，保距性质的自同态，也是幺正自同构。
\begin{pp}\label{pp:8_8_1}
设$A$是配备了内积$(\,\cdot\, | \,\cdot\,)$的空间$\mathbb{V}$中的自同态，则以下三个性质等价：
\begin{enumerate}[label=\alph*.]
    \item $A$是幺正自同构。
    \item $\forall \,\, \mathbf{u}\,\, \in \mathbb{V}$，$ \| A\mathbf{u}\| = \|\mathbf{u}\|.$
    \item $\forall \,\, \mathbf{u}, \mathbf{v} \,\, \in \mathbb{V}$，$ (A\mathbf{u} | A\mathbf{v}) = (\mathbf{u}|\mathbf{v}).$
\end{enumerate}
\end{pp}
\begin{proof} % [{\bfseries 证明\ref{pp:8_8_1}}]
前面已经说明了$(a) \Rightarrow (b)$，现在证明$(b) \Leftrightarrow (c)$：

$(b) \Leftarrow (c)$：如果对任意向量$\mathbf{u}, \mathbf{v}$，$(A\mathbf{u}|A\mathbf{v}) = (\mathbf{u}|\mathbf{v})$，那么令$\mathbf{u} = \mathbf{v}$，就得到：对任意向量$\mathbf{u}$，$\|A\mathbf{u}\| = \|\mathbf{u}\|$。

$(b) \Rightarrow (c)$：对于复空间內积，我们有如下恒等式：
\begin{align*}
    &(\mathbf{x}+\mathbf{y}|\mathbf{x}+\mathbf{y}) - (\mathbf{x}-\mathbf{y}|\mathbf{x}-\mathbf{y}) + \imath \left[(\mathbf{x}+\imath \mathbf{y}|\mathbf{x}+\imath \mathbf{y}) - (\mathbf{x}-\imath \mathbf{y}|\mathbf{x}-\imath \mathbf{y})\right]  \\
    =& 2(\mathbf{x}|\mathbf{y}) + 2(\mathbf{y}|\mathbf{x}) + \imath \left[-\imath (\mathbf{x}|\mathbf{y}) + \imath (\mathbf{y}|\mathbf{x}) + \imath (\mathbf{y}|\mathbf{x}) -\imath (\mathbf{x}|\mathbf{y}) \right]  \\
    =& 2(\mathbf{x}|\mathbf{y}) + 2(\mathbf{y}|\mathbf{x}) + 2(\mathbf{x}|\mathbf{y}) - 2(\mathbf{y}|\mathbf{x})  \\
    =& 4(\mathbf{x}|\mathbf{y}) 
\end{align*}
同理，
\begin{align*}
    &(A(\mathbf{x}+\mathbf{y})|A(\mathbf{x}+\mathbf{y})) - (A(\mathbf{x}-\mathbf{y})|A(\mathbf{x}-\mathbf{y}))  \\
    &\quad + \imath \left[(A(\mathbf{x}+\imath \mathbf{y})|A(\mathbf{x}+\imath \mathbf{y})) - (A(\mathbf{x}-\imath \mathbf{y})|A(\mathbf{x}-\imath \mathbf{y}))\right]  \\
    =& 2(A\mathbf{x}|A\mathbf{y}) + 2(A\mathbf{y}|A\mathbf{x}) + \imath \left[-\imath (A\mathbf{x}|A\mathbf{y}) + \imath (A\mathbf{y}|A\mathbf{x}) + \imath (A\mathbf{y}|A\mathbf{x}) -\imath (A\mathbf{x}|A\mathbf{y}) \right]  \\
    =& 2(A\mathbf{x}|A\mathbf{y}) + 2(A\mathbf{y}|A\mathbf{x}) + 2(A\mathbf{x}|A\mathbf{y}) - (A\mathbf{y}|A\mathbf{x})  \\
    =& 4(A\mathbf{x}|A\mathbf{y}) 
\end{align*}
所以，如果对任意向量$\mathbf{u}$，$\|A\mathbf{u}\| = \|\mathbf{u}\|$，那么对任意向量$\mathbf{u}, \mathbf{v}$，$(A\mathbf{u}|A\mathbf{v}) = (\mathbf{u}|\mathbf{v})$。

对于实空间內积，我们也有恒等式：
\begin{align*}
    &(\mathbf{x}+\mathbf{y}|\mathbf{x}+\mathbf{y}) - (\mathbf{x}-\mathbf{y}|\mathbf{x}-\mathbf{y})  \\
    =& 2(\mathbf{x}|\mathbf{y}) + 2(\mathbf{y}|\mathbf{x}) = 4(\mathbf{x}|\mathbf{y}) 
\end{align*}
同理，
\begin{align*}
    &(A(\mathbf{x}+\mathbf{y})|A(\mathbf{x}+\mathbf{y})) - (A(\mathbf{x}-\mathbf{y})|A(\mathbf{x}-\mathbf{y}))  \\
    =& 2(A\mathbf{x}|A\mathbf{y}) + 2(A\mathbf{y}|A\mathbf{x}) = 4(A\mathbf{x}|A\mathbf{y}) 
\end{align*}

所以，如果对任意向量$\mathbf{u}$，$\|A\mathbf{u}\| = \|\mathbf{u}\|$，那么对任意向量$\mathbf{u}, \mathbf{v}$，$(A\mathbf{u}|A\mathbf{v}) = (\mathbf{u}|\mathbf{v})$。

最后证明$(c) \Rightarrow (a)$：
对幺正自同构来说，
$$(\mathbf{u} | \mathbf{v}) = (A\mathbf{u} | A\mathbf{v}) = (\mathbf{u} | A^\star A\mathbf{v})$$
因此
$$  \forall \mathbf{u}, \,\, \mathbf{v} \,\, \in \mathbb{V},\,\,\, (\mathbf{u} | (A^\star A - \mathrm{I})\mathbf{v}) = 0.$$
这说明$A^\star A = \mathrm{I}$，同理可得$ A A^\star = \mathrm{I}$。\\
\end{proof}


幺正自同构也称为\textbf{保距映射}或\textbf{保距自同构}。

保距映射不仅保持距离和长度，也保持角度。上一章中，我们用
$$\frac{(\mathbf{u}|\mathbf{v})}{\|\mathbf{u}\|\|\mathbf{v}\|}$$
来定义向量张成的角度。而保距映射下，非零向量$\mathbf{u}, \mathbf{v}$变成$A\mathbf{u}, A\mathbf{v}$后，它们张成的角度对应
$$\frac{(A\mathbf{u}|A\mathbf{v})}{\|A\mathbf{u}\|\|A\mathbf{v}\|} = \frac{(\mathbf{u}|\mathbf{v})}{\|\mathbf{u}\|\|\mathbf{v}\|}.$$
也就是说，经过保距映射，向量张成的角度不变\footnote{严谨起见，我们需要说明以上分式有意义，即分母不为零。读者可以补上此处的说明。}。因此，保距映射也称为保角映射。

正交也是角度的一种。对幺正自同构来说，$(\mathbf{u}|\mathbf{v}) = 0$等价于$(A\mathbf{u}|A\mathbf{v}) = 0$，
也就是说，两个向量如果正交，经过映射后仍然正交。一组向量如果两两正交，经过映射后仍然两两正交。
（归一）正交基经过变换后仍然是（归一）正交基。最后这个结论无非是说幺正自同构就是幺正基变换，上一节中我们已经提到了这一点。

\subsection{归一正交群}
我们再来看一下幺正自同构的代数性质。记$\mathcal{O}_n(\mathbb{F})$为$n$维幺正自同构的集合。
显然，恒等映射是幺正自同构，属于$\mathcal{O}_n(\mathbb{F})$。接下来，两个幺正自同构的复合也是幺正自同构。
最后，幺正自同构的逆映射还是幺正自同构。所以，$\mathcal{O}_n(\mathbb{F})$关于映射的复合构成群，我们把它称为归一正交群。
同理可得，所有$n$维幺正矩阵构成的集合关于矩阵乘法也构成同样的群，且与$\mathcal{O}_n(\mathbb{F})$同构，
我们把它也称为\textbf{归一正交群}，简称\textbf{幺正群}。归一正交群是平直群$\mathrm{GL}_n(\mathbb{F})$的子群。

用$2$维的归一正交群举例。设配备了标准基和点积的实內积空间$\mathbb{R}^2$中自同态$A$的矩阵为：
$$ \begin{bmatrix} a & b \\ c & d \end{bmatrix}. $$
根据定义，$A A^H = A^H A = \mathrm{I}_2$，于是可以得到以下方程组：
$$ \left\{
\begin{array}{l}
    a^2 + b^2 = c^2 + d^2 = 1 \\
    ac + bd = 0 \\
    a^2 + c^2 = b^2 + d^2 = 1 \\
    ac + bd = 0
\end{array} 
\right.
$$
记$a=\cos{\theta}, b = \sin{\theta}$，$c = \cos{\varphi}, d = \sin{\varphi}$，那么从上式可得：$\cos{(\theta - \varphi)} = 0$，且$|\cos{\theta}| = |\sin{\varphi}|, |\sin{\theta}| = |\cos{\varphi}|$ 。于是有$\varphi = \theta \pm \frac{\pi}{2} + 2k\pi, \,\,\, k\in\mathbb{Z}.$ 分别对应两种情况：
$$\begin{bmatrix}\cos{\theta} & \sin{\theta} \\ -\sin{\theta} & \cos{\theta}\end{bmatrix} \mbox{ 和 } \begin{bmatrix}\cos{\theta} & \sin{\theta} \\ \sin{\theta} & -\cos{\theta}\end{bmatrix}$$
注意到后者可以表示为前者和固定矩阵的乘积：
$$\begin{bmatrix}\cos{\theta} & \sin{\theta} \\ \sin{\theta} & -\cos{\theta}\end{bmatrix} =  \begin{bmatrix} 1 &0 \\ 0 & -1\end{bmatrix} \begin{bmatrix}\cos{\theta} & \sin{\theta} \\ -\sin{\theta} & \cos{\theta}\end{bmatrix}$$
$$\begin{bmatrix}\cos{\theta} & \sin{\theta} \\ \sin{\theta} & -\cos{\theta}\end{bmatrix} = \begin{bmatrix}\cos{\theta} & -\sin{\theta} \\ \sin{\theta} & \cos{\theta}\end{bmatrix}\begin{bmatrix} 1 &0 \\ 0 & -1\end{bmatrix}
$$
所以，$2$维实空间的幺正自同构总可以用两种矩阵$\mathrm{R}_2(\theta)$和$\Gamma$表达：
$$\mathrm{R}_2(\theta) = \begin{bmatrix}\cos{\theta} & \sin{\theta} \\ -\sin{\theta} & \cos{\theta}\end{bmatrix}, \,\,\, \Gamma = \begin{bmatrix} 1 &0 \\ 0 & -1\end{bmatrix}$$
任意幺正自同构$A\in\mathcal{O}_2(\mathbb{R})$，要么等于某个$\mathrm{R}_2(\theta)$，要么等于$\Gamma$和某个$\mathrm{R}_2(\theta)$的乘积。

分别来看这两种矩阵。$\mathrm{R}_2(\theta)$可以称为\textbf{旋转矩阵}，对应的自同态可称为\textbf{旋转映射}。
直角坐标系（对应点积）中，任意向量$(x, y)$，经过$\mathrm{R}_2(\theta)$的作用，等于以原点为中心，逆时针旋转了$\theta$角度：
$$\begin{bmatrix}\cos{\theta} & \sin{\theta} \\ -\sin{\theta} & \cos{\theta}\end{bmatrix} \begin{bmatrix} x \\ y\end{bmatrix} = \begin{bmatrix}x\cos{\theta} - y\sin{\theta} \\ x\sin{\theta} + y\cos{\theta}\end{bmatrix} $$
而$\Gamma$是沿$y$轴关于$x$轴的倒影映射，或者说关于$x$轴的轴对称映射：把$(x,y)$映射到$(x, -y)$。我们也称这种映射为\textbf{镜面反射}或\textbf{镜像映射}。
镜面反射是沿某个向量生成的一维子空间（直线），平行于它的正交补的倒影映射，是一种特殊的倒影映射。
用几何的语言来说，欧几里得平面上保距的直映射，要么是沿原点旋转，要么是关于$x$轴作轴对称之后，沿原点旋转。
后者也可以看作关于某条过原点的直线作轴对称。

复內积空间$\mathbb{C}^2$上的情况更复杂一点。按定义我们有如下方程组：
$$\left\{
\begin{array}{l}
    |a|^2 + |b|^2 = |c|^2 + |d|^2 = 1 \\
    \bar{a}c + \bar{b}d = 0 \\
    |a|^2 + |c|^2 = |b|^2 + |d|^2 = 1 \\
    \bar{a}b + \bar{c}d = 0
\end{array}
\right.
$$
记$a=\e^{\imath \alpha}\cos{\theta}, b = \e^{\imath \beta}\sin{\theta}$，则$c=\e^{\imath \alpha'}\sin{\theta}, d = \e^{\imath \beta'}\cos{\theta}$，且满足：
$$
\left\{
\begin{array}{l}
    \left(\e^{\imath (\alpha' - \alpha)} + \e^{\imath (\beta' -\beta)}\right)\cos{\theta}\sin{\theta}= 0 \\
    \left(\e^{\imath (\alpha - \beta)} + \e^{\imath (\alpha' -\beta')}\right)\cos{\theta}\sin{\theta}= 0 
\end{array}
\right.
$$
解得$\theta = \frac{k\pi}{2}, \,\,\, k\in\mathbb{Z}$，或$\alpha' = \alpha + \varphi, \;\;\beta' = \beta + \varphi+\pi, \,\,\, \varphi \in \mathbb{R}$。据此可以得到复内积空间上幺正矩阵的形态：
$$
\forall A \in \mathcal{O}_2(\mathbb{C}), \,\,\, A = \begin{bmatrix}\e^{\imath \alpha} \e^{\imath \varphi} \cos{\theta} & \e^{\imath \alpha} \sin{\theta} \\ -\e^{\imath \beta} \e^{\imath \varphi}\sin{\theta} & \e^{\imath \beta}\cos{\theta}\end{bmatrix} 
 \mbox{ 或 } \begin{bmatrix}\e^{\imath \alpha} & 0 \\ 0 & \e^{\imath (\beta+\varphi)} \end{bmatrix} \mbox{ 或 } \begin{bmatrix}0 & \e^{\imath \beta} \\ \e^{\imath (\alpha+\varphi)} & 0\end{bmatrix}
 $$
也可以写成统一的分解形式：
$$\forall A \in \mathcal{O}_2(\mathbb{C}), \,\,\, 
A = \begin{bmatrix}\e^{\imath \alpha} & 0 \\ 0 & -\e^{\imath \beta} \end{bmatrix}
\begin{bmatrix} \cos{\theta} & -\sin{\theta} \\ \sin{\theta} & \cos{\theta}\end{bmatrix}
\begin{bmatrix}\e^{\imath \varphi} & 0 \\ 0 & -1 \end{bmatrix}. $$

可以看到，$\mathbb{R}$上的归一正交群和$\mathbb{C}$上的归一正交群性质不同。为了方便，我们把前者称为\textbf{实正交群}，把其中元素称为\textbf{实正交变换}；把后者称为\textbf{复正交群}，把其中元素称为\textbf{复正交变换}。一般书籍中也把后者称为\textbf{酉群}，把其中元素称为\textbf{酉变换}。

\section{规范自同态}
\subsection{可幺正对角化的矩阵}
复内积空间中的自伴自同态可以幺正对角化，在由它的本征向量构成的幺正基上，它的矩阵是对角矩阵。
不仅如此，由于它的本征值都是实数，对角矩阵的对角元素都是实数。
我们希望把这个性质扩展到更广泛的一类自同态上。
我们希望把内积空间的自伴自同态的性质扩展到更广泛的一类自同态上。
\begin{df}\label{df:8_3}
    给定内积空间$\mathbb{V}$中自同态$T$，如果$T$在一组幺正基上的矩阵是对角矩阵，
    就称它为\textbf{规范自同态}或\textbf{规范算子}。换句话说，如果$T$在某个幺正基上的矩阵:
    $$ [T]_{\mathcal{B}} = D $$
    其中$D$是对角矩阵，就称$T$是规范的。
\end{df}
显然，自伴自同态必然是规范的。那么，怎样判断一个自同态是否是规范自同态呢？

首先来理解一下规范自同态和自伴自同态的关系。假设在某个幺正基上，
某个自伴自同态的矩阵是对角矩阵。我们知道，这时对角元素都是实数。
实际上，在复内积空间中，对角元素显然可以是复数。所以我们猜测，“自伴”的特性给了自伴自同态“对角线为实数”的限制。
也许，自伴自同态和规范自同态的关系，就和它们对角元素一样，是实数和复数的关系。

我们知道，共轭转置是对合变换：两次共轭转置的复合是恒等变换。对复数来说，共轭变换也是对合变换。
给定复数$z$，将其看作$1\times 1$的矩阵，则它的共轭转置就是共轭。
所以，我们可以猜测，是否可以把共轭转置（$\bigstar$）看作复系数矩阵（自伴自同态）的“共轭变换”？
对于实数来说，共轭变换是恒等变换；也可以说，共轭变换的不动点集就是实数集。
而复系数矩阵（复内积空间）中，共轭对称矩阵（自伴自同态）恰好也是共轭转置（$\bigstar$）的不动点集。

具体来说，对于自伴自同态$T$，如果$\lambda$是$T$的本征值，那么$\lambda$也是$T^*$的本征值。
将它幺正对角化后，这个性质十分明显：对角线上的元素都是实数，所以共轭转置后仍然是自身。
如果$T$可以幺正对角化，但对角线上的元素是复数，那么，共轭转置后得到的是其共轭复数。所以，我们应该有如下性质：

如果$\lambda$是$T$的本征值，那么$\overline{\lambda}$是$T^\star$的本征值。

这个命题可以变成：
$$ \| T(\mathbf{v}) - \lambda \mathbf{v}\|^2 = 0 \Leftrightarrow \| T^\star(\mathbf{v}) - \overline{\lambda} \mathbf{v}\|^2 = 0.$$
如果把两个等式的左侧展开，会发现：
\begin{align*}
    \| T(\mathbf{v}) - \lambda \mathbf{v}\|^2 &= \| T(\mathbf{v})\|^2 - \lambda (\mathbf{v} | T(\mathbf{v})) - \overline{\lambda} (T(\mathbf{v}) | \mathbf{v}) + |\lambda|^2,  \\
    \| T^\star(\mathbf{v}) - \overline{\lambda} \mathbf{v}\|^2 &= \| T^*(\mathbf{v})\|^2 -  \overline{\lambda} (\mathbf{v} | T^\star(\mathbf{v})) - \lambda (T^\star(\mathbf{v}) | \mathbf{v}) + |\lambda|^2,  \\
    &= \| T^\star(\mathbf{v})\|^2 -  \overline{\lambda} (T(\mathbf{v}) | \mathbf{v}) - \lambda (\mathbf{v} | T(\mathbf{v})) + |\lambda|^2, 
\end{align*}
观察两者的展开式（右侧），可以注意到其中只有第一项不一样。也就是说，
如果$T$满足$\| T(\mathbf{v})\|^2 = \| T^\star(\mathbf{v})\|^2$，那么
$$ \| T(\mathbf{v}) - \lambda \mathbf{v}\|^2 = 0 \Leftrightarrow \| T^\star(\mathbf{v}) - \overline{\lambda} \mathbf{v}\|^2 = 0.$$
如果我们进一步转化，还会得到：
\begin{align*}
    \| T(\mathbf{v})\|^2 &= (T(\mathbf{v}) | T(\mathbf{v})) = ((T^\star\circ T)(\mathbf{v}) | \mathbf{v})  \\
    \| T^\star(\mathbf{v})\|^2 &= (T^\star(\mathbf{v}) | T^\star(\mathbf{v})) = ((T\circ T^\star)(\mathbf{v}) | \mathbf{v})  
\end{align*}
也就是说，如果自同态$T$和它的伴随自同态交换：$T^\star\circ T = T\circ T^\star$，那么前面提到的性质也成立。

现在，我们来证明以上这些性质和规范矩阵的定义是等价的：
\begin{pp}\textbf{复内积空间的谱定理 }\label{pp:8_9}
    设$T$为复内积空间$\mathbb{V}$中的自同态，则以下命题等价：
    \begin{enumerate}[label=\alph*.]   
        \item 存在由$T$的本征向量构成的幺正基$\mathcal{B}$，使$T$在$\mathcal{B}$上是对角矩阵。
        \item $T^\star\circ T = T\circ T^\star.$
        \item $\forall \,\, \mathbf{v} \,\, \in \mathbb{V}, \,\,\, \| T(\mathbf{v})\| = \| T^\star(\mathbf{v})\|.$
    \end{enumerate}
\end{pp}
\begin{proof} % [{\bfseries 证明\ref{pp:8_9}}]\label{pf:8_9}
    $(a)\Rightarrow (b)$：如果存在幺正基$\mathcal{B} = (\mathbf{e}_1, \mathbf{e}_2, \cdots , \mathbf{e}_n)$使得$T$在其上的矩阵是对角矩阵，
    那么
    $$[T]_{\mathcal{B}} = \begin{bmatrix}
        a_1 & 0 & \cdots & 0 \\
        0 & a_2 & \cdots & 0 \\
        \vdots & \vdots & \ddots & \vdots \\
        0 & 0 & \cdots & a_n \\
    \end{bmatrix}.
    $$
    因而
    $$[T^\star]_{\mathcal{B}} = \begin{bmatrix}
        \overline{a_1} & 0 & \cdots & 0 \\
        0 & \overline{a_2} & \cdots & 0 \\
        \vdots & \vdots & \ddots & \vdots \\
        0 & 0 & \cdots & \overline{a_n} \\
    \end{bmatrix}.
    $$
    于是：
    \begin{align*}
        [T\circ T^\star]_{\mathcal{B}} &= \begin{bmatrix}
            a_1\overline{a_1} & 0 & \cdots & 0 \\
            0 & a_2\overline{a_2} & \cdots & 0 \\
            \vdots & \vdots & \ddots & \vdots \\
            0 & 0 & \cdots & a_n\overline{a_n} \\
        \end{bmatrix} = \begin{bmatrix}
            \overline{a_1}a_1 & 0 & \cdots & 0 \\
            0 & \overline{a_2}a_2 & \cdots & 0 \\
            \vdots & \vdots & \ddots & \vdots \\
            0 & 0 & \cdots & \overline{a_n}a_n \\
        \end{bmatrix}  \\
        &= [T^\star\circ T]_{\mathcal{B}} 
    \end{align*}
    这说明$T^\star\circ T = T\circ T^\star.$\\
    $(b)\Rightarrow (c)$：
    \begin{align*}
        T^*\circ T = T\circ T^\star &\Rightarrow T^\star\circ T - T\circ T^\star = 0  \\
        &\Rightarrow \forall \,\, \mathbf{v} \,\, \in\mathbb{V}, \quad ((T^\star\circ T - T\circ T^\star)(\mathbf{v})|\mathbf{v}) = 0  \\
        &\Rightarrow \forall \,\, \mathbf{v} \,\, \in\mathbb{V}, \quad ((T^\star\circ T)(\mathbf{v})|\mathbf{v}) = ((T\circ T^\star)(\mathbf{v})|\mathbf{v})  \\
        &\Rightarrow \forall \,\, \mathbf{v} \,\, \in\mathbb{V}, \quad (T(\mathbf{v})|T(\mathbf{v})) = (T^\star(\mathbf{v})|T^\star(\mathbf{v}))  \\
        &\Rightarrow \forall \,\, \mathbf{v} \,\, \in\mathbb{V}, \quad \|T(\mathbf{v})\| = \|T^\star(\mathbf{v})\|  
    \end{align*}
    $(c)\Rightarrow (a)$：将$T$归一正交三角化，得到
    $$ [T]_{\mathcal{B}} = U.$$
    其中${\mathcal{B}} = (\mathbf{e}_1, \mathbf{e}_2, \cdots , \mathbf{e}_n)$是归一正交三角基，$U = [u_{i,j}]_{1\leqslant i \leqslant n}$是上三角矩阵：
    $$U = \begin{bmatrix}
        a_{1,1} & a_{1,2} & \cdots & a_{1,n} \\
        0 & a_{2,2} & \cdots & a_{2,n} \\
        \vdots & \vdots & \ddots & \vdots \\
        0 & 0 & \cdots & a_{n,n} \\
    \end{bmatrix}.
    $$ 
    现在证明$U$是对角矩阵，也就是说$i<j$的时候$a_{i,j} = 0$.\\
    根据矩阵$U$可以写出：
    $$ T(\mathbf{e}_1) = a_{1,1} \mathbf{e}_1,$$
    所以$\| T(\mathbf{e}_1)\| = |a_{1,1}|$，另一方面，$T^\star$在基底${\mathcal{B}}$上的矩阵是$U$的共轭转置，所以
    $$ T^\star(\mathbf{e}_1) = \sum_{i=1}^n \overline{a_{1,i}} \mathbf{e}_i,$$
    也就是说$\| T^\star(\mathbf{e}_1)\|^2 = \sum_{i=1}^n |a_{1,i}|^2.$
    而
    $$ \forall \,\, \mathbf{v} \,\, \in\mathbb{V}, \quad \|T(\mathbf{v})\| = \|T^\star(\mathbf{v})\|, $$
    所以
    $$ \sum_{i=1}^n |a_{1,i}|^2 = \| T^\star(\mathbf{e}_1)\|^2 = \| T(\mathbf{e}_1)\|^2 = |a_{1,1}|^2. $$
    这说明$a_{1,2} = a_{1,3} = \cdots = a_{1,n} = 0$.\\
    从$a_{1,2} = 0$可以推出
    $$ T(\mathbf{e}_2) = a_{1,2} \mathbf{e}_1 + a_{2,2} \mathbf{e}_2 = a_{2,2} \mathbf{e}_2,$$
    另一方面，
    $$ T^\star(\mathbf{e}_2) = \sum_{i=2}^n \overline{a_{2,i}} \mathbf{e}_i.$$
    于是，根据类似的推理可知$a_{2,2} = \cdots = a_{2,n} = 0$. 我们又得到$a_{2,3} = 0$. \\
    我们可以这样一步步推下去。一般地，在$k$步，我们已知对任意$i < k$，$i < j$，都有$a_{i,j} = 0$。
    于是，作为其中一部分，$a_{1,k} = a_{2,k} = \cdots = a_{k-1, k} = 0$。所以
    $$ T(\mathbf{e}_{k}) = \sum_{i=1}^k a_{i,k} \mathbf{e}_i = a_{k,k} \mathbf{e}_k,$$
    因此$\| T(\mathbf{e}_{k}) \|^2 = |a_{k,k}|^2$. 而根据伴随自同态的性质，我们又有：
    $$ T^\star(\mathbf{e}_k) = \sum_{i=k}^n \overline{a_{k,i}} \mathbf{e}_i.$$
    于是得到$\| T^\star(\mathbf{e}_{k}) \|^2 = \sum_{i=k}^n |a_{k,i}|^2$，因此我们得到：
    $$ a_{k,k+1} = a_{k,k+2} = \cdots = a_{k,n} = 0,$$
    与$i<k$，$i<j$，$a_{i,j} = 0$合并，就得到：只要$i<k+1$，$i<j$，就有$a_{i,j} = 0$。
    其中就有$a_{k,k+1} = 0$，于是我们又可以进行下一步……\\
    重复以上步骤，可以推出，只要$i<j$，就有$a_{i,j} = 0$。这就证明了$U$是对角矩阵。\\
\end{proof}

我们把这个版本的谱定理称为完整版的复内积空间谱定理。实内积空间中，自伴自同态可以幺正对角化。
复内积空间中，我们找到了所有可以幺正对角化的自同态，也就是与自己的伴随矩阵交换的自同态。

\subsection{规范自同态与幺正自同构}
实内积空间中，是否也能将规范自同态扩展到与自己的伴随矩阵交换的自同态呢？
我们研究前面二维幺正群中的例子：旋转映射。它的矩阵是：
$$\mathrm{R}_2(\theta) = \begin{bmatrix}\cos{\theta} & \sin{\theta} \\ -\sin{\theta} & \cos{\theta}\end{bmatrix}. $$
它的几何意义是以原点为中心逆时针旋转$\theta$。而作为幺正变换，它的伴随自同态就是它的逆映射，矩阵形式为：
$$\mathrm{R}(-\theta) = \begin{bmatrix}\cos{\theta} & -\sin{\theta} \\ \sin{\theta} & -\cos{\theta}\end{bmatrix}. $$
也就是以原点为中心顺时针旋转$\theta$。这两个映射显然交换：
$$\mathrm{R}_2(\theta)\mathrm{R}(-\theta) = \mathrm{R}(-\theta)\mathrm{R}_2(\theta) = \mathrm{I}_2.$$
但从几何意义来说，$\mathrm{R}_2(\theta)$的效果是把向量旋转一个角度，所以对于大多数的$\theta$来说，不可能和原来的向量共线，
也就是说，不会有本征值和本征向量，于是$\mathrm{R}_2(\theta)$作为二维实内积空间的自同态，无法对角化。

当然，幺正自同构作为复内积空间中的规范自同态，可以在复内积空间中幺正对角化。对角化后，其矩阵的对角元素为幺正自同构的本征值，
因此都是模长为1的复数。另一方面，如果自同态$A$在某个幺正基上是对角矩阵，且对角元素$a_{ii}$都是模长为1的复数，
那么它的伴随自同态在对应基底上的矩阵（$A$的共轭转置）也是对角矩阵，且对角线相应位置的值是其共轭$\overline{a_{ii}}$。
于是$A A^\star = A^\star A = \mathrm{I}_n$，即$A$为幺正自同构。于是我们得到幺正矩阵乃至幺正自同构的又一种刻画：
\begin{pp}\label{pp:8_10}
    内积空间中的自同态$A$是幺正的，当且仅当存在一组由$A$的本征向量构成的幺正基，对应的本征值都是模长为1的复数。
\end{pp}
如果把自伴自同态比作实数，把规范自同态比作一般复数，那么$\mathcal{O}_n(\mathbb{F})$相当于复平面的单位圆。

\section{正定自同态}
\subsection{正定性}
研究自伴自同态的时候，我们提到过，通过研究内积对应的自伴自同态，来了解内积的性质。
而我们知道，内积是正定的（共轭）对称双直型。因此，我们也应该研究对应的自同态性质。
\begin{df}\label{df:8_4}
    如果内积空间$\mathbb{V}$中的自同态$S$满足以下条件，就称它是\textbf{正定自同态}：
    $$ \forall \mathbf{v} \in \mathbb{V}, \,\,\, (S\mathbf{v} | \mathbf{v}) \geqslant 0.$$
\end{df}

来看几个正定自同态的例子。

若自同态$D$在内积对应的幺正基上的矩阵是对角矩阵，且对角元素全部大于等于$0$，则$D$是正定自同态。
如果某个对角元素小于$0$，则它不是正定自同态。

设$\mathbb{U}$是$\mathbb{V}$的子空间，关于$\mathbb{U}$的正交投影映射是正定自同态。

设$T$是$\mathbb{V}$中的自伴自同态，$a$为正实数，$b, c$为实数，且满足$b^2 < 4ac$，
则$aT^2 + bT + c\mathrm{I}$是正定自同态。

读者可以分别自行验证。

从第一个例子可以看出，如果自同态可幺正对角化，那么它是正定自同态的充要条件就是对角元素全部大于等于$0$，也就是本征值全部大于等于$0$。
因此，根据谱定理可以得出：
\begin{pp}\label{pp:8_10_1}
内积空间的自伴自同态是正定自同态，当且仅当它的所有本征值都是非负实数。
\end{pp}

那么，内积空间的正定自同态是否可以扩展到自伴自同态以外呢？下面这个命题说明了，复内积空间中，正定自同态的概念只能停留在自伴自同态中。
\begin{pp}\label{pp:8_11}
    如果内积空间中的自同态$S$满足：
    $$ \forall \mathbf{v} \in \mathbb{V}, \,\,\, (S\mathbf{v} | \mathbf{v}) \geqslant 0,$$
    那么肯定是自伴自同态。
\end{pp}
\begin{proof} % [{\bfseries 证明\ref{pp:8_11}}]
    $$ \forall \mathbf{v} \in \mathbb{V}, \,\,\, (S\mathbf{v} | \mathbf{v}) \geqslant 0,$$
    这说明$(S\mathbf{v} | \mathbf{v})$是正实数。因此
    $$ (S\mathbf{v} | \mathbf{v}) = \overline{(S\mathbf{v} | \mathbf{v})} = (\mathbf{v} | S\mathbf{v})$$
    借助类似\ref{pp:8_8_1}中的恒等式，对任意向量$\mathbf{u}, \mathbf{v}$，
    我们可以把$(S\mathbf{u} | \mathbf{v})$表示为：
    \begin{align*}
        &(S\mathbf{u} | \mathbf{v})  \\
        =&\frac14(S(\mathbf{u}+\mathbf{v})|(\mathbf{u}+\mathbf{v})) - \frac14(S(\mathbf{u}-\mathbf{v})|(\mathbf{u}-\mathbf{v}))  \\
        &\quad + \frac{\imath }{4}\left[(S(\mathbf{u}+\imath \mathbf{v})|(\mathbf{u}+\imath \mathbf{v})) - (S(\mathbf{u}-\imath \mathbf{v})|(\mathbf{u}-\imath \mathbf{v}))\right]  \\
        =&\frac14((\mathbf{u}+\mathbf{v})|S(\mathbf{u}+\mathbf{v})) - \frac14((\mathbf{u}-\mathbf{v})|S(\mathbf{u}-\mathbf{v}))  \\
        &\quad + \frac{\imath }{4}\left[((\mathbf{u}+\imath \mathbf{v})|S(\mathbf{u}+\imath \mathbf{v})) - ((\mathbf{u}-\imath \mathbf{v})|S(\mathbf{u}-\imath \mathbf{v}))\right]  \\
        =&(\mathbf{u} | S\mathbf{v}) 
    \end{align*}
    这说明$S$是自伴自同态。\\
\end{proof}

实内积空间的正定自同态，有没有可能不是自伴自同态呢？实内积空间中，上面的恒等式不复存在。
我们还是考察二维的旋转映射$\mathrm{R}_2(\theta)$。$\theta = \frac{\pi}{2}$时，映射把向量逆时针旋转90°，所以和原向量垂直，
$(S\mathbf{v} | \mathbf{v}) = 0$。$\theta = \frac{\pi}{3}$时，映射把向量逆时针旋转60°，
$(S\mathbf{v} | \mathbf{v}) = \frac{1}{2}\|\mathbf{v}\|^2 \geqslant 0$。但我们知道这些旋转映射不是自伴自同态。

\subsection{自同态的平方根}
对内积空间的自同态，如何刻画其正定性呢？我们回到前面自同态和复数的类比。
如果说规范自同态可比作复数，自伴自同态可比作实数，那么正定自同态似乎可以比作非负实数。
从复内积空间来看，正定自同态对角化后，其本征值确实都是非负实数。
于是，我们尝试用从非负实数的性质出发，来刻画正定自同态。

我们知道，非负实数总是某个实数的平方，或某个复数的模长的平方。从对角矩阵这个特例来看，
给定对角矩阵$D$，如果对角元素是实数，那么$D^2$这个对角矩阵的对角元素都是非负实数，因此是正定矩阵；
如果对角元素是复数，那么$D^\star D$这个对角矩阵的对角元素都是非负实数，因此是正定矩阵。

这让我们猜想，如果某个自同态能写成另一个自伴自同态的平方，或某个自同态与其伴随映射的复合，那么它就是正定矩阵。

这个性质不难验证。比如，设自同态$T = Q^\star Q$，其中$Q$是自同态，那么：
\begin{align*}
    \forall \mathbf{v} \in \mathbb{V},  \\
    (T\mathbf{v} | \mathbf{v}) &= (Q^\star Q\mathbf{v} | \mathbf{v})  \\
    &= (Q\mathbf{v} | Q\mathbf{v})  \\
    &= \| Q\mathbf{v} \|^2 \geqslant 0 
\end{align*}
因此$T$是正定自同态。如果$T = R^2$，其中$R$是自伴自同态，那么由于$R = R^\star$，$T = R^\star R$，
同理可得$T$是正定自同态。反之，对任意自同态$Q$，$T = Q^\star Q$总是自伴自同态，根据谱定理，
它必然能幺正对角化。因此，在某个幺正基$\mathcal{B}$上，矩阵$[T]_\mathcal{B}$是对角矩阵。
$$ [T]_\mathcal{B} = \begin{bmatrix}
    d_{1,1} & 0 & \cdots & 0 \\
    0 & d_{2,2} & \cdots & 0 \\
    \vdots & \vdots & \ddots & \vdots \\
    0 & \cdots & 0 & d_{n,n}
\end{bmatrix}.$$
而由于它是正定自同态，所有对角元素都是非负实数。于是，可以构造矩阵：
$$ R = \begin{bmatrix}
    \sqrt{d_{1,1}} & 0 & \cdots & 0 \\
    0 & \sqrt{d_{2,2}} & \cdots & 0 \\
    \vdots & \vdots & \ddots & \vdots \\
    0 & \cdots & 0 & \sqrt{d_{n,n}}
\end{bmatrix}.$$
它满足$R^2 = [T]_\mathcal{B}$。我们把这样的矩阵称为\textbf{矩阵的平方根}，把它和$\mathcal{B}$确定的自同态称为自同态$T$的\textbf{平方根}，记为$\sqrt{T}$。
于是，$\sqrt{T}^2 = T = Q^\star Q$。我们也可以写：$\sqrt{T} = \sqrt{Q^\star Q}$。

显然，上式中的$R$也是正定的。正定的自伴自同态，拥有正定的平方根。这与我们前面的推想吻合：
可以用平方和模长来刻画正定性。
\begin{pp}\label{pp:8_12}
    给定内积空间中的自同态$T$，以下条件等价：
    \begin{enumerate}
        \item $T$是正定的自伴自同态。
        \item $T$可以幺正对角化，且对角元素（本征值）大于等于$0$。
        \item 存在自同态$Q$使得$T = Q^\star Q.$
        \item 存在自伴自同态$R$，使得$T =R^2.$
    \end{enumerate}
\end{pp}

如果在对角矩阵$R$的某个对角元素，比如说$\sqrt{d_{1,1}}$前面加个负号，得到的矩阵仍满足$R^2 = [T]_\mathcal{B}$. 
所以，如果只把自同态$T$的平方根定义为“平方为$T$的自伴自同态”，那么它的平方根不是唯一的。
这就好比数系里，$2$和$-2$的平方都是$4$。要唯一定义自同态（矩阵）的平方根，就必须加上正定性的限制。

以下定理说明，只要加上正定性的限制，自同态（矩阵）的平方根就是唯一的。
\begin{pp}\label{pp:8_13}
    设$T$是正定的自伴自同态。使得$R^2 = T$的自伴自同态$R$中，有且仅有一个是正定自同态，称之为\textbf{自同态的平方根}。
\end{pp}
\begin{proof} % [{\bfseries 证明\ref{pp:8_13}}]
    $T$是自伴自同态，所以根据谱定理，存在由$T$的本征向量构成的基底。
    $T$的正定性质说明它的本征值都是非负实数。
    设$R$是使得$R^2 = T$的正定自伴自同态。对$T$的任意本征值$\lambda$和对应的本征向量$\mathbf{v}$，
    下面我们证明
    $$R\mathbf{v} = \sqrt{\lambda}\mathbf{v}.$$
    这样，通过$T$的本征向量构成的基底，我们就唯一确定了$R$。\\
    $R$也是正定自伴自同态，所以存在由$R$的本征向量构成的基底：$\mathcal{B} = (\mathbf{e}_i)_{1\leqslant i \leqslant n}$。
    设对应的本征值为$R\mathbf{e}_i = \lambda_i \mathbf{e}_i$，其中$\lambda_i$都是非负实数。\\
    把向量$\mathbf{v}$在$\mathcal{B}$上分解：
    $$ \mathbf{v} = v_1 \mathbf{e}_1 + v_2 \mathbf{e}_2 + \cdots + v_n \mathbf{e}_n $$
    则
    $$ R^2 \mathbf{v} = v_1 \lambda_1^2 \mathbf{e}_1 + v_2 \lambda_2^2 \mathbf{e}_2 + \cdots + v_n \lambda_n^2 \mathbf{e}_n. $$
    另一方面，$R^2 \mathbf{v} = T \mathbf{v}$，所以：
    $$ R^2 \mathbf{v} = \lambda \mathbf{v} = v_1 \lambda \mathbf{e}_1 + v_2 \lambda \mathbf{e}_2 + \cdots + v_n \lambda \mathbf{e}_n.$$
    于是：
    $$ v_1 (\lambda - \lambda_1^2) \mathbf{e}_1 + v_2 (\lambda - \lambda_2^2) \mathbf{e}_2 + \cdots + v_n (\lambda - \lambda_n^2) \mathbf{e}_n = \mathbf{0}.$$
    基底是直无关的，所以$\forall \,\, i ,\,\, v_i (\lambda - \lambda_i^2) = 0$. 这说明，对每个$i$，
    或者$v_i = 0$，或者$\lambda_i = \sqrt{\lambda}$，或者两者皆有之。因此，总有$v_i\lambda_i = v_i\sqrt{\lambda}$。
    这说明
    $$ v_1 \lambda_1 \mathbf{e}_1 + v_2 \lambda_2 \mathbf{e}_2 + \cdots + v_n \lambda_n \mathbf{e}_n = v_1 \sqrt{\lambda} \mathbf{e}_1 + v_2 \sqrt{\lambda} \mathbf{e}_2 + \cdots + v_n \sqrt{\lambda} \mathbf{e}_n, $$
    也就是说
    $$ R \mathbf{v} = \sqrt{\lambda} \mathbf{v}.$$
\end{proof}

\chapter{二次型}

\section{双直型与二次型}
我们把实空间上的内积定义为对称正定双直型。正定性对于内积很重要，因为对应定义的模长应该总大于等于零。那么，如果没有正定性，一般的双直型有什么性质呢？

\begin{df}
平直空间$\mathbb{R}^n$上的双直型：
\begin{align*}
    \phi:\;\;\mathbb{R}^n \times \mathbb{R}^n \rightarrow \mathbb{R}      
\end{align*}
是满足：
\begin{enumerate}
    \item $\forall \mathbf{u}, \mathbf{v}, \mathbf{w} \in \mathbb{R}^n$，$\forall \, k \in \mathbb{R}$，$\phi(\mathbf{u} + k\mathbf{v}, \mathbf{w}) = \phi(\mathbf{u}, \mathbf{w}) + k \phi(\mathbf{v}, \mathbf{w})$。
    \item $\forall \mathbf{u}, \mathbf{v}, \mathbf{w} \in \mathbb{R}^n$，$\forall \, k \in \mathbb{R}$，$\phi(\mathbf{w}, \mathbf{u} + k\mathbf{v}) = \phi(\mathbf{w},\mathbf{u}) + k \phi(\mathbf{w},\mathbf{v})$。
\end{enumerate}
的映射。如果$\forall \mathbf{u}, \mathbf{v}$，都有$\phi(\mathbf{u}, \mathbf{v}) = \phi(\mathbf{v}, \mathbf{u})$，就说$\phi$是对称双直型。
\end{df}

给定$\mathbb{R}^n$的基底$\mathcal{B} = (\mathbf{e}_1, \mathbf{e}_2, \cdots, \mathbf{e}_n)$，则按照定义，
\begin{align*}
    \forall \;\mathbf{x} = \sum_{i=1}^n x_i\mathbf{e}_i,& \;\;\mathbf{y} = \sum_{i=1}^n y_i\mathbf{e}_i, \\
    \phi(\mathbf{x},\mathbf{y}) &= \sum_{i=1}^n\sum_{j=1}^n x_i y_j \phi(\mathbf{e}_i,\mathbf{e}_j).
\end{align*}
因此，只要知道了所有$\phi(\mathbf{e}_i,\mathbf{e}_j)$，就能确定双直型$\phi$对任意向量的结果。

让我们定义：
\begin{align*}
    \mathfrak{q}:\;\;\mathbb{R}^n  &\rightarrow \mathbb{R}  \\
            \mathbf{v}\;\; &\mapsto \;\; \phi(\mathbf{v}, \mathbf{v})
\end{align*}
这个映射把$\mathbb{R}^n$的向量映射到实数。但要注意，它并非平直映射，因此不是直型。按上面的分解方法，可以把$\mathfrak{q}$写为：
\begin{align*}
    \forall \;\mathbf{x} = \sum_{i=1}^n x_i\mathbf{e}_i,\quad \mathfrak{q}(\mathbf{x}) = \sum_{i=1}^n\sum_{j=1}^n x_i x_j \phi(\mathbf{e}_i,\mathbf{e}_j).
\end{align*}
这是一个关于$n$个实数自变量$x_1,x_2,\cdots,x_n$的二次齐次整式。所以我们常称其为\textbf{二次型}。

为什么要研究二次型呢？

让我们回顾高中讨论的二次曲线。平面中，形如
$$ ax^2 + bxy + cy^2 + dx + ey + f = 0 $$
的方程描述的曲线称为二次曲线。我们在高中已经知道，通过平移和旋转坐标轴，可以把二次曲线化为标准形式，从而判断它是椭圆、双曲线还是抛物线。实际上，经过平移消去一次项后，二次曲线方程就对应一个二次型。或者说，
$$ ax^2 + bxy + cy^2 $$
恰好可以写成$\mathfrak{q}(\mathbf{v}) = \phi(\mathbf{v}, \mathbf{v})$的形式，其中$\phi$是$\mathbb{R}^2$上的某个双直型。具体来看一个例子。
$$\phi: \,( \mathbf{x}, \mathbf{y} ) \; \mapsto \; x_1y_1 - 3x_1y_2 + x_2y_1 + 2x_2y_2.
$$
可以验证$\phi$分别关于每个自变量（向量）是直映射。于是$\phi$是双直型。而对应的二次型是：
$$ \mathfrak{q}:\; (x_1, x_2) \; \mapsto \; x_1^2 - 3x_1x_2 + x_2x_1 + 2x_2^2 = x_1^2 - 2x_1x_2 + 2x_2^2.
$$
可以看到，$x_1y_2$和$x_2y_1$在二次型中成为了同类项，可以合并。合并后，二次型$\mathfrak{q}$的判别式：
$$
\Delta = (-2)^2 - 4\cdot 1\cdot 2 = -4 < 0.
$$
因此，方程$\mathfrak{q} = 1$是一个椭圆。可以通过幺正基变换，将其变为标准形式。

对于一般的$n$，二次型就对应了高维空间中的某类“曲面”。而研究二次型就是研究高维空间中“曲面”的性质。

如何研究二次型？可以从研究二次曲线的方法出发。研究二次曲线时，消去一次项后，我们面对二次齐次整式，通过旋转坐标轴，得到了二次曲线的标准形式。用现在的语言来说，我们做了一次基变更，在新的基底下，二次曲线方程的“交叉项”：$bxy$被消去了，于是只剩下：
$$ ax^2 + cy^2,$$
而通过讨论$a,c$的正负，我们得到了不同类型的二次曲线。

对于$\mathbb{R}^n$上的二次型，我们也会采用类似的思路将其约简并分类。不过，为了能使用我们之前得到的结论，我们首先要用自同构的语言来表达二次型。

给定二次型：
$$ \mathfrak{q}:\;\; (x_1, x_2, \cdots, x_n) \; \mapsto \; \sum_{i=1}^n \sum_{j=1}^n a_{i,j} x_i x_j, $$
我们可以将其表示为：
$$ \mathbf{x} \; \mapsto \; [\mathbf{x}]^{\tr} A [\mathbf{x}]. $$
其中$A = [a_{i,j}]_{1\leqslant i,j\leqslant n}$是$n\times n$方阵。注意到我们可以把二次型写成
$$ (x_1, x_2, \cdots, x_n) \; \mapsto \; \sum_{i=1}^n \sum_{j=1}^n \frac{a_{i,j} + a_{j,i}}{2} x_i x_j, $$
这不改变二次型的结果，而$x_i x_j$的系数等于$x_j x_i$的系数（或者说合并同类项后，我们重新分配了系数）。这说明我们可以假设矩阵$A$总是关于对角线对称的。

注意到由于$A$是对称矩阵，我们还可以把二次型写成：
$$ \mathfrak{q}(\mathbf{x}) =  \nji{A[\mathbf{x}]}{[\mathbf{x}]}. $$
这里的$\nji{\cdot}{\cdot}$是点积。可以验证，它和前面的定义等价。这说明，给定每个二次型，我们可以根据它在某个特定基底上的矩阵，找到一个自伴自同态（因为矩阵是对称的）$T$，使得
$$ \mathfrak{q}(\mathbf{x}) = \nji{T(\mathbf{x})}{\mathbf{x}}. $$
注意：这里的$T$随着内积的选取而确定，不同的内积对应不同的$T$。一般我们会默认内积是基底对应的点积，这时$T$的矩阵就是$A$。

反过来，给定一个二次型$\mathfrak{q}$，能否定义一个对称双直型$\phi$，使得$\forall \mathbf{v}$，$\mathfrak{q}(\mathbf{v}) = \phi(\mathbf{v}, \mathbf{v})$呢？根据对称双直型总满足的等式：
\begin{align*}
    \phi(\mathbf{u} + \mathbf{v}, \mathbf{u} + \mathbf{v}) = \phi(\mathbf{u}, \mathbf{u}) + \phi(\mathbf{v}, \mathbf{v}) + 2\phi(\mathbf{u}, \mathbf{v})
\end{align*}
我们可以定义对称双直型：
$$
\phi(\mathbf{u}, \mathbf{v}) = \frac{1}{2}\left(\mathfrak{q}(\mathbf{u} + \mathbf{v}) - \mathfrak{q}(\mathbf{u}) - \mathfrak{q}(\mathbf{v})\right)
$$
这个双直型满足我们的要求。注意这样定义的双直型必然是对称的，而且唯一存在。它和$\mathfrak{q}$对应同一个对称矩阵。如果不要求对称，那么可以自由调整相应矩阵非对角线上的元素，得到不同的双直型。

以上的矩阵是二次型在特定基底上的矩阵。如果基底不同，矩阵会如何变化呢？

设有$\mathbb{R}^n$上的基底$\mathcal{B}=(\mathbf{e}_1,\mathbf{e}_2,\cdots,\mathbf{e}_n)$和$\mathcal{B}'=(\mathbf{e}_1',\mathbf{e}_2',\cdots,\mathbf{e}_n')$。设向量$\mathbf{x}\in\mathbb{R}^n$在$\mathcal{B}$和$\mathcal{B}'$中的坐标分别是$(x_i)_{1\leqslant i\leqslant n}$和$(x_i')_{1\leqslant i\leqslant n}$。我们知道，对于直自同态$T$，它在$\mathcal{B}$和$\mathcal{B}'$中的矩阵满足：
$$ 
[T]_{\mathcal{B}'} = P^{-1}[T]_{\mathcal{B}} P.
$$
其中$P$称为$\mathcal{B}$到$\mathcal{B}'$的基变换矩阵，满足：
$$
    (x_i)_{1\leqslant i\leqslant n} = P(x_i')_{1\leqslant i\leqslant n}.
$$
因此，如果记二次型$\mathfrak{q}$在$\mathcal{B}$和$\mathcal{B}'$中对应的矩阵分别是$A, A'$，那么
\begin{align*}
    [\mathbf{x}]_{\mathcal{B}'}^{\tr} A' [\mathbf{x}]_{\mathcal{B}'} = \mathfrak{q}(\mathbf{x}) &= [\mathbf{x}]_{\mathcal{B}}^{\tr} A [\mathbf{x}]_{\mathcal{B}} \\
    &= (P[\mathbf{x}]_{\mathcal{B}'})^{\tr} A P[\mathbf{x}]_{\mathcal{B}'} \\
    &= [\mathbf{x}]_{\mathcal{B}'}^{\tr} P^{\tr} A P [\mathbf{x}]_{\mathcal{B}'}
\end{align*}
因此得到：
$$ 
    A' = P^{\tr} A P.
$$
这个关系和相似矩阵不同。我们称之为\textbf{合同关系}。具体来说，只要有可逆矩阵$P$使得$A' = P^{\tr} A P$，就说矩阵$A$和$A'$合同，是$A'$是$A$的\textbf{合同矩阵}。

与相似关系一样，矩阵的合同关系是等价关系。首先，取$P=I$，则可知任意矩阵与自身合同。其次，如果$A$和$B$合同，那么按定义有可逆矩阵$P$使得
$$ B = P^{\tr} A P , \quad \mbox{即}\;A = \left(P^{\tr}\right)^{-1} B P^{-1}. $$ 
但矩阵转置的逆矩阵是其逆矩阵的转置矩阵。因此，
$$ B = P^{\tr} A P \quad \Longleftrightarrow \quad A = \left(P^{-1}\right)^{\tr} B P^{-1}. $$ 
$P^{-1}$显然也是可逆矩阵，因此$B$和$A$也合同。最后，如果矩阵$A$和$B$合同，$B$和$C$合同，那么按定义有可逆矩阵$P,Q$使得
$$ B = P^{\tr} A P , \quad C = Q^{\tr} B Q, $$ 
于是
$$ C = Q^{\tr} P^{\tr} A P Q = (PQ)^{\tr} A P Q. $$
$PQ$显然也是可逆矩阵，因此$C$和$A$也合同。至此，我们也就可以说两个矩阵互为合同矩阵了。从变换的视角来看，也是理所当然，毕竟合同矩阵就是同一个二次型在不同基底的表现而已。

注意：如果矩阵$P$的转置矩阵等于它的逆矩阵，那么合同关系就和相似关系一致了。或者说，如果$P^{\tr} = P^{-1}$，那么$A'$既可以看作$A$对应的直自同态在新基底的矩阵，也可以看作$A$对应的二次型在新基底对应的矩阵。

\section{二次型约简}
% 二次型的标准型和分类
给定$n$维空间的二次型，我们可以将它写成
$$ \mathfrak{q}:\;\; (x_1, x_2, \cdots, x_n) \; \mapsto \; \sum_{i=1}^n a_{i,i} x_i^2 + \sum_{i=1}^n \sum_{j<i} 2a_{i,j} x_i x_j. $$
它的表达式一共有$\frac{n(n+1)}{2}$项。其中前$n$项是完全平方项，由形如$x_i^2$的项构成；后面$\frac{n(n-1)}{2}$项称为“交叉项”，由$i\neq j$的项$x_i x_j$构成。不难验证，$n$维空间的所有二次型可以看作一个关于映射的加法和数乘运算构成平直空间，共$\frac{n(n+1)}{2}$维。

从研究二次曲线的过程来看，我们认为，没有交叉项，只有完全平方项的二次型是比较简单的。对于二次曲线来说，就是化简后的标准形式。我们把这样的二次型写出来：
$$ \mathfrak{q}:\;\; (x_1, x_2, \cdots, x_n) \; \mapsto \; \sum_{i=1}^n a_{i,i} x_i^2 . $$
我们称这种二次型为\textbf{对角二次型}。
它对应的矩阵就是对角矩阵。对角线元素是$a_{1,1}, \cdots, a_{n,n}$。我们知道，对角线矩阵在矩阵中也是较为简单的矩阵。

来看对应的二次型$\mathfrak{q}$。如果所有的$a_{i,i}$都是正数，那么$\mathfrak{q}$对应着某个正定的对称双直型，也就是内积。它就是该内积对应的模映射的平方。我们称它为\textbf{正定二次型}。反之，如果所有的$a_{i,i}$都是负数，那么$-\mathfrak{q}$对应这样的内积和模映射。这时我们称$\mathfrak{q}$为\textbf{负定二次型}。

如果所有的$a_{i,i}$都非负，就说$\mathfrak{q}$是\textbf{半正定}的二次型。半正定的条件比正定宽松。它表示即便$\mathbf{x}$不是零向量，$\mathfrak{q}(\mathbf{x})$也可能等于$0$。半正定二次型的性质在于$\mathfrak{q}(\mathbf{x})$总大于等于$0$。这一点让它成为不等式和极值判定时的常用条件。比如，如果向量函数$f$能写成$f(\mathbf{x}) = h(\mathbf{x}) + \mathfrak{q}(\mathbf{x})$，其中$\mathfrak{q}(\mathbf{x})$是半正定二次型，就说明$f\geqslant h$总成立。反之，如果所有的$a_{i,i}$都非正，就说$\mathfrak{q}$是\textbf{半负定}的二次型。半正定和半负定二次型统称半定二次型。

而如果以上条件都不满足，那么就说明对角线元素里既有正数也有负数。我们说这样的二次型是不定的。

为了更细致地研究，我们把对角二次型的$n$个系数：$a_{i,i}$中大于零的个数记为$q^+$，小于零的个数记为$q^-$，零的个数记为$l$。按定义显然总有
$$ q^+ + q^- + l = n.$$
正定二次型就是$q^+ = n$，$q^- = l = 0$；半正定二次型就是$q^- = 0$，$q^+ + l = n$；不定二次型就是$q^+ > 0$且$q^- > 0$。

如果允许变换基向量的位置，那么我们可以把对角二次型的$a_{i,i}$整理为所有正数在最前，然后是负数，最后是$0$，这样方便后续讨论。

以上是对角二次型的讨论。对于一般的二次型，我们能否以及如何将其化归为对角二次型呢？我们知道，二次型表达式里的交叉项，对应着矩阵形式下的非对角线元素。也就是说，我们要把一个实系数对称矩阵转为对角矩阵。而这件事情和我们之前对于自伴自同态做的事情是类似的。通过基变换，把自伴自同态的矩阵转为对角矩阵。区别在于，自同态的基变换对应的是矩阵间的相似关系，而二次型的基变换对应矩阵间的合同关系。不过，谱定理说明，实内积空间里的自伴自同态可以幺正对角化。这说明，对称矩阵不仅可以转为对角矩阵，而且其过渡矩阵$P$是幺正变换，$P$的伴随自同态就是它的逆映射。具体来说，给定实对称矩阵$S$，我们把它看作实内积空间中的自伴自同态，那么，谱定理说明该自同态有幺正基，在新基底下它的矩阵是对角矩阵$D$：
$$ D = P^{-1} S P.$$
而$P^{\tr} = P^{-1}$，所以上式又可以写为
$$ D = P^{\tr} S P.$$
换句话说，如果我们把二次型对应的实对称矩阵$S$看作实内积空间的自伴自同态的矩阵，把它幺正对角化，那么，我们可以拿对角化所用的过渡矩阵$P$，作为基变换矩阵，对二次型进行基变换。这样得到的二次型在新基底下的矩阵，就是对角矩阵了。

综上，结论是：我们总可以通过基变换，把实系数二次型转为对角二次型。

来看一个具体的例子。考虑$\mathbb{R}^2$上的二次型：
$$ \mathfrak{q}(x_1, x_2) = x_1^2 + 4x_1x_2 + x_2^2. $$
它对应的对称矩阵$S$是：
$$
S = 
\begin{bmatrix}
    1 & 2 \\ 2 & 1
\end{bmatrix}
$$
其中非对角系数$s_{1,2} = s_{2,1} = 2$，加起来就是$4$。我们的目标是找到一组新的基底，使得
$\mathfrak{q}$在新基底上的矩阵为对角矩阵。按照前面的讨论，我们将$S$看作某个实内积空间（取标准点积）中的自伴自同态，对其进行幺正对角化。首先找到本征值$3,-1$，对应的本征向量分别是$(1,1)$和$(-1,1)$。归一后，我们可以定义过渡矩阵$P$：
$$
P = \frac{\sqrt{2}}{2}\cdot\begin{bmatrix}
    1 & -1 \\ 1 & 1
\end{bmatrix} = 
\begin{bmatrix}
    \displaystyle\frac{\sqrt{2}}{2} &\displaystyle -\frac{\sqrt{2}}{2} \\ \displaystyle\frac{\sqrt{2}}{2} &\displaystyle \frac{\sqrt{2}}{2}
\end{bmatrix},
$$
使得
$$
D = P^{-1} S
P = \begin{bmatrix}
    \frac{\sqrt{2}}{2} & \frac{\sqrt{2}}{2} \\  -\frac{\sqrt{2}}{2} &\frac{\sqrt{2}}{2}
\end{bmatrix} \begin{bmatrix}
    1 & 2 \\ 2 & 1
\end{bmatrix}\begin{bmatrix}
    \frac{\sqrt{2}}{2} & -\frac{\sqrt{2}}{2} \\ \frac{\sqrt{2}}{2} &\frac{\sqrt{2}}{2}
\end{bmatrix}
= \begin{bmatrix}
    3 & 0 \\ 0 & -1
\end{bmatrix}
$$
而我们可以验证$P^{\tr} P = I$，$P$是幺正矩阵。因此，$S$作为二次型矩阵也可以写为：
$$
S = PDP^{-1} = PDP^{\tr} = P \begin{bmatrix}
    3 & 0 \\ 0 & -1
\end{bmatrix}
P^{\tr}, \quad \mbox{或}\; P^{\tr}SP = \begin{bmatrix}
    3 & 0 \\ 0 & -1
\end{bmatrix}. 
$$
这说明通过基变换，二次型$\mathfrak{q}$在新基底上写为：
$$ \mathfrak{q}(x_1, x_2) = 3x_1^2 - x_2^2. $$

以上的讨论可以用于二次曲线的讨论：给定方程$x^2 + 4xy + y^2 = 1$，我们可以通过旋转将其变为标准形式。这个旋转变换就记录在过渡矩阵里。$P$可以写为：
$$
P = 
\begin{bmatrix}
    \displaystyle\frac{\sqrt{2}}{2} & \displaystyle -\frac{\sqrt{2}}{2} \\\displaystyle \frac{\sqrt{2}}{2} & \displaystyle\frac{\sqrt{2}}{2}
\end{bmatrix} = \begin{bmatrix}
    \displaystyle\cos{\frac{\pi}{4}} & \displaystyle -\sin{\frac{\pi}{4}} \\ \displaystyle\sin{\frac{\pi}{4}} & \displaystyle\cos{\frac{\pi}{4}}
\end{bmatrix}.
$$
这是一个逆时针$45^\circ$的旋转。说明原曲线是双曲线$3x^2 - y^2 = 1$逆时针旋转$45^\circ$后得到的，顺时针旋转$45^\circ$后就变成双曲线$3x^2 - y^2 = 1$。

另外，我们在讨论中默认使用了点积作为内积，这样对应的自伴自同态$T$的矩阵就是二次型的矩阵。如果我们不使用点积，而使用别的内积，对应的自伴自同态的矩阵就不再是二次型的矩阵。具体来说，内积本身就是一种对称双直型（正定），它在基底上表现为对称正定矩阵。设$H$是某个内积对应的矩阵，我们把该内积记为$\nji{\cdot}{\cdot}_H$。它的具体形式为：
$$ \forall \mathbf{x}, \; \mathbf{y} \in \mathbb{R}^n , \;\; (\mathbf{x}, \; \mathbf{y}) \; \mapsto [\mathbf{x}]^{\tr} H[\mathbf{y}] = \nji{\mathbf{x}}{\mathbf{y}}_H. $$
这里使用方括号$[\cdot]$表示向量在$\mathbb{R}^n$的默认基底上的分量表示，作为强调，以和其他基底区别。

我们默认设定的自伴自同态$T$是关于经典内积（点积）的，它的矩阵就是二次型的矩阵。关于新内积$H$，我们也能找到自伴自同态$T_H$，按定义它满足：
$$ \forall \mathbf{x}, \; \mathbf{y} \in \mathbb{R}^n , \;\; \nji{T_H\mathbf{x}}{\mathbf{y}}_H = \nji{T\mathbf{x}}{\mathbf{y}}, $$
即
$$ [T_H\mathbf{x}]^{\tr} H[\mathbf{y}] = [T\mathbf{x}]^{\tr} [\mathbf{y}]. $$
这说明
$$ [T_H]^{\tr} H = T^{\tr}.$$
即
$$ T_H = \left(H^{\tr}\right)^{-1} T = H^{-1} T.$$
这就是不同内积对应的自伴自同态之间的关系。

\section{分容定理}
% Sylvester指标和惯性定理
前一节中，我们讨论了如何把二次型转为对角二次型。为了使用谱定理，我们使用幺正基变换。这是为了说明这样的基变换（根据谱定理）必然存在。不过，把二次型转为对角二次型的方法并非只有这一种。比如，在前面的例子里，我们还可以这样做：
$$ \mathfrak{q}(x_1, x_2) = x_1^2 + 4x_1x_2 + x_2^2 = (x_1 + 2x_2)^2 - 3x_2^2. $$
因此，让$u_1 = x_1+2x_2$，$u_2 = x_2$，则$x_1 = u_1 - 2u_2$，$x_2 = u_2$。也就是说，
$$
\begin{bmatrix}
    u_1 \\ u_2
\end{bmatrix} = 
\begin{bmatrix}
    1 & 2 \\ 0 & 1
\end{bmatrix}
\begin{bmatrix}
    x_1 \\ x_2
\end{bmatrix}, \quad 
\begin{bmatrix}
    x_1 \\ x_2
\end{bmatrix} = 
\begin{bmatrix}
    1 & -2 \\ 0 & 1
\end{bmatrix}
\begin{bmatrix}
    u_1 \\ u_2
\end{bmatrix}
$$
后者是我们要的基变换矩阵$P$：
$$
S' = P^{\tr} S P =
\begin{bmatrix}
    1 & 0 \\ -2 & 1
\end{bmatrix}
\begin{bmatrix}
    1 & 2 \\ 2 & 1
\end{bmatrix}
\begin{bmatrix}
    1 & -2 \\ 0 & 1
\end{bmatrix} = 
\begin{bmatrix}
    1 & 0 \\ 0 & -3
\end{bmatrix}.
$$
我们看到，这里的基变换矩阵并非幺正变换，而是一个错切变换。变换后的形式也不同了：
$$ \mathfrak{q}(u_1, u_2) = u_1^2 -  3u_2^2. $$
直观上看，这是因为我们没有旋转，而是把$y$轴往右“拉歪”了。不过，考察$3x^2 - y^2 = 1$和$x^2 - 3y^2 = 1$会发现，这两个曲线都是双曲线。所以，用不同的合同变换得到对角二次型，并没有改变相应曲线的本质。这让我们思考：合同变化之下，二次型的什么东西是不变的？如何构造一个不变量？

从上面的例子来看，不同的基变换不改变二次型对应双曲线的性质，不会把双曲线变为椭圆。换句话说，得到的对角矩阵，对角线元素总是一正一负。用我们前面定义的话来说，就是二次型以不同方式合同的对角矩阵里，$q^+, q^-, l$的数量不变。我们猜想，这就是同一个二次型在合同变化下的不变量。

\begin{tm}{\textbf{分容定理}}
    设$\mathfrak{q}$是$n$维空间的二次型，$q^+, q^-, l$分别是$\mathfrak{q}$在某个基底下对应的对角矩阵中正数、负数、零的个数。则对于任意基底，$\mathfrak{q}$在该基底下对应的对角矩阵中正数、负数、零的个数仍然是$q^+, q^-, l$。我们把$q^+, q^-$分别称为二次型$\mathfrak{q}$的\textbf{正容}和\textbf{负容}，其和称为\textbf{二次型的容}。    
\end{tm}
% 分容定理起名理由：合同变换不改变矩阵的容，但这不是关键。这个定理的关键意义是说明其中的q^+, q^-是定值。也就是说，两者把二次型的容分成两个固定的部分。于是称为分容定理。

\begin{proof}
    首先证明合同变换不改变矩阵的容。实际上，如果$A' = P^{\tr} A P$，则
    $$\rang(A') = \rang(P^{\tr} A P) \leqslant \minn{}\{\rang(P^{\tr}), \rang(A), \rang(P)\}.$$
    而$P$是可逆矩阵，$\rang(P) = n$，所以$\rang(A') \leqslant \rang(A)$。而由于合同关系是等价关系，以上不等式反过来也成立，所以$\rang(A') = \rang(A)$。\\
    这说明二次型在任意基底下的矩阵都有同样的容$r$，我们称之为二次型的容。据此可知，$q^ + q^- = r$是定值，$l = n - r$也是定值。于是，我们只需要证明任意基底下的$q^+$不变即可。\\
    接下来定义“正定子空间”的概念：如果某个子空间$U$使得$\forall \;\mathbf{0}\,\neq\, \mathbf{u}\in U$，都有$\mathfrak{q}(\mathbf{u}) > 0$，就说$U$是正定子空间。考虑所有正定子空间的维数，取其中最大的维数，记为$q^*$。我们要证明，任意基底下的$q^+$都等于$q^*$。\\
    首先证明$q^+ \leqslant q^*$。给定基底$\mathcal{B}$，设$\mathfrak{q}$在$\mathcal{B}$对应的对角矩阵中正数的个数为$q^+$。设$V^+$是由这些对应正数的基向量张成的子空间，则$\dim(U) = q^+$。$\forall \;\mathbf{0}\,\neq\, \mathbf{u}\in V^+$，都有$\mathfrak{q}(\mathbf{u}) > 0$，因此$V^+$是正定子空间。于是$q^+ \leqslant q^*$。\\
    另一方面，考虑任一个正定子空间$U$，考虑$U$和$\mathcal{B}$中对应非正数的基向量张成的子空间$V^-$的交集。$\forall \; 0 \,\neq \,\mathbf{v} \in U \cap V^-$，都有$\mathfrak{q}(\mathbf{v}) > 0$，又有$\mathfrak{q}(\mathbf{v}) \leqslant 0$，导致矛盾！因此$U \cap V^- = \{\mathbf{0}\}$。于是按照维数基本定理，
    \begin{align*}
        n \geqslant \dim(U + V^-) &= \dim(U) + \dim(V^-) - \dim(U \cap V^-) \\
        &= \dim(U) + (n - q^+)\, -\, 0 \; = \dim(U) + n - q^+,\\
        \mbox{即}\;\;\dim(U) &\leqslant q^+.
    \end{align*}
    这个不等关系对任意正定子空间$U$都成立，因此$q^* \leqslant q^+$。\\
    综上，$q^+ = q^*$是定值。于是，任意基底下的$q^+$都等于$q^*$，从而$q^-$和$l$也都是定值。也就是说，二次型在合同变化下的正数、负数、零的个数不变。
\end{proof}

进一步来说，如何确定特定二次型的正负容呢？首先二次型对应的矩阵是对称矩阵。将其看作自伴自同态的矩阵，则我们已经证明过，按照谱定理，可以幺正对角化。这时它的对角线元素就是它的本征值。而分容定理告诉我们，这些本征值中正数、负数和零的个数不随合同变换改变。于是，二次型的正负容就等于其对应自伴自同态的正负本征值的个数。如果我们继续通过缩放变换把这些本征值变为$1,-1,0$，那么就得到二次型的标准形式：
\begin{tm}
    二次型总可以通过合同变换表示为对角元素为$1,-1,0$的对角矩阵。这个对角矩阵称为二次型的\textbf{标准形式}。其中$1$和$-1$的个数分别是二次型的正容和负容。
\end{tm}

\sectionext{同时对角化}
% 可同时对角化的条件
我们讨论过把一个自同态或一个二次型转为对角形式的问题。比这样化简单个自同态或单个二次型，另一个同样有必要而更复杂的问题是：如何把两个或更多自同态或二次型同时转为对角形式？也就是说，是否能找到一组基底，使得多个自同态或二次型能表示为对角矩阵？我们称之为\textbf{同时对角化}。

同时对角化在很多问题中都有应用。

具体来说，对角化可分为自同态的对角化和二次型的对角化。自同态的对角化指为自同态对应的矩阵找到对角的相似矩阵，对应矩阵的相似关系，可称为\textbf{相似对角化}；二次型的对角化指为二次型对应的矩阵找到对角的合同矩阵，对应矩阵的合同关系，可称为\textbf{合同对角化}。一个是$D = P^{-1}SP$，一个是$D = P^{\tr}SP$。这是两种不同的对角化方式。仅在$P$为幺正矩阵时，这两种对角化方式才一致。我们通过对称矩阵把二次型和实内积空间的自伴自同态联系了起来，让二次型幺正对角化，但这并不是唯一的对角化方式。我们可以通过合同变换把二次型转为对角形式，但这不一定是幺正变换。

至于同时将两个乃至多个二次型或自同态对角化，就更需要分别讨论了，因为要找到适用于多个矩阵的幺正对角化难度更大。我们首先来讨论自同态的同时对角化问题。即我们想找到一组基底，使得多个自同态在该基底下的矩阵都是对角矩阵。我们知道，单个自同态可以相似对角化，但两个自同态能否同时相似对角化呢？答案是：不一定。我们需要给出条件。

首先来看一个必要条件。假设$A$和$B$在同一组基底下都是对角矩阵：
$$ D_A = P^{-1}AP, \quad D_B = P^{-1}BP. $$
由于对角矩阵可交换，因此$D_AD_B = D_BD_A$。从而
$$ AB = P D_A P^{-1} P D_B P^{-1} = P D_A D_B P^{-1} = P D_B D_A P^{-1} = B A. $$
也就是说，可以同时对角化的两个矩阵必然可交换。反过来是否成立呢？考虑两个矩阵$A,B$，假设它们可交换，即$AB = BA$。也就是说，对应的自同态$T_A, T_B$满足$T_A\circ T_B = T_B\circ T_A$。也就是说，如果$\lambda$和$\mathbf{v}$分别是$T_A$的本征值和对应的本征向量，那么
$$ T_A(T_B(\mathbf{v})) = T_B(T_A(\mathbf{v})) = T_B(\lambda \mathbf{v}) = \lambda (T_B\mathbf{v}). $$
这说明$T_B\mathbf{v}$也是$T_A$的本征向量。也就是说，$T_B$把$T_A$的本征空间映射到自身。

于是考虑$T_A$的任意本征空间$E_A(\lambda)$（关于本征值$\lambda$的本征空间）。首先里面任意非零向量都是关于$\lambda$的本征向量。其次，$T_B(E_A(\lambda))\subseteq E_A(\lambda)$。考虑$T_B$限制在$E_A(\lambda)$的映射$B_{\lambda}$，它是$E_A(\lambda)$上的自同态。它的极小式是$T_B$的极小式的因式，因此是散式。于是$B_{\lambda}$可对角化。也就是说我们可以在$E_A(\lambda)$中找到一组$T_B$的本征向量作为基底，使得$B_{\lambda}$在该基底下的矩阵是对角矩阵。而由于$T_A$限制在本征空间上就是$\lambda I$，所以无论怎么取基底都是对角矩阵。于是我们就找到了$T_A$的本征空间$E_A(\lambda)$的一组基底，使得$T_A$和$T_B$在该基底下都是对角矩阵。

这件事在$T_A$的每个本征空间中都成立。我们把$T_B$在$T_A$的各个本征空间中对角化的基底合并起来。由于这些本征空间的直和就是整个空间，于是合并后的基底就是整个空间的基底。而$T_A$和$T_B$在该基底下都是对角矩阵。

综上，我们得到结论：
\begin{tm}
    两个可对角化的自同态可同时对角化，当且仅当它们可交换。
\end{tm}

再来看二次型的同时对角化问题。我们知道，二次型对应的矩阵是对称矩阵。一个简单的情况是“是否能同时幺正对角化”。这等于是说基变换矩阵是幺正的，因此和前面的情况相同了。通过视之为实内积空间的自伴自同态。我们可以得出其充要条件：两个二次型可同时幺正对角化，当且仅当它们的矩阵（或者说对应的自伴自同态）可交换。

对于一般的情形，我们想知道对于两个二次型（对应矩阵$A,B$），能否有可逆矩阵$P$，使得$P^{\tr}AP$和$P^{\tr}BP$都是对角矩阵。

按二次型的定义，$A,B$都是对称矩阵，分别可以合同对角化：$A = P^{\tr}D_AP,\; B = Q^{\tr}D_BQ$。如果对角矩阵$D_A$或$D_B$的全部对角元素都是正数，不妨设是$D_A$，那么$D_A$可以写成：$D_A = D^2$。其中对角矩阵$D$的对角元素分别是$D_A$对角元素的平方根。于是$A = P^{\tr}DDP = P^{\tr}D^{\tr}DP = (DP)^{\tr} DP$。记$R=DP$，则$A = R^{\tr} R$。$R$是可逆矩阵的乘积，所以是可逆矩阵。考虑$(R^{-1})^{\tr}BR^{-1}$，它也是对称矩阵，按照谱定理，可以幺正对角化。即有满足$U^{\tr}U=I$的幺正矩阵$U$使得$U^{\tr}(R^{-1})^{\tr}BR^{-1}U$为对角矩阵。再考虑$U^{\tr}(R^{-1})^{\tr}AR^{-1}U$，
\begin{align*}
    U^{\tr}(R^{-1})^{\tr}AR^{-1}U = U^{\tr}(R^{-1})^{\tr}R^{\tr}RR^{-1}U = U^{\tr}U = I.
\end{align*}
这说明可逆矩阵$R^{-1}U$能同时把$A,B$合同对角化。不仅如此，我们还直接把$A$对角化为单位矩阵。总结：
\begin{tm}
    设$A,B$为实对称矩阵，如果其中一个是正定矩阵，那么两者可以同时合同对角化，而且可以将正定矩阵化为单位矩阵。
\end{tm}
注意这是一个可同时对角化的充分条件。如果要求等价条件，则有如下结论（证明参见附录三\ref{c-30}）：\begin{tm}\label{tm:8-2-2}
    设$A,B$为实对称矩阵，其中至少一个可逆。那么以下两命题等价：
    \begin{enumerate}
        \item $A,B$可同时合同对角化。
        \item 有正定对称矩阵$H$使得$AHB=BHA$。
    \end{enumerate}
\end{tm}

% \section{本征值的范围}
% % Rayleigh商和本征值的极值
% 在第一节我们讨论过用内积表示二次型。给定基底的实内积空间里，二次型的矩阵可以看作自伴自同态，即可以写成：
% $$ \mathfrak{q}(\mathbf{x}) = \nji{T(\mathbf{x})}{\mathbf{x}}. $$
% 根据对偶表示定理，这种表示方法必然存在。于是，我们可以借助内积空间的工具来研究二次型。

\section{极分解与表征分解}

\subsection{极分解}
每个正定的自伴自同态对应着一个内积，那么，一般的自同态和正定的自伴自同态有什么差别呢？
回到自同态和复数的类比，如果把一般自同态看作复数，把正定的自伴自同态看作非负实数，那么正定的自伴自同态的平方根可以看作
矩阵的模长。这让我们想起某个复数和模长的关系：复数的极坐标形式：
$$ z = r \e^{\imath \theta} $$
其中$z$是复数，$r$是它的模长，$\e^{\imath \theta}$是复平面单位圆上的点。

我们知道，模长可以对应自同态的平方根，而单位圆上的点可以对应幺正自同构。
因此，可以猜测：自同态可以分解为幺正自同构和自同态平方根的复合：
$$ T = U \circ R,$$
其中$T$是自同态，$U$是幺正自同构，$R$是正定自伴自同态，或者说，是某个正定自伴自同态的平方根。

从上一节可以知道，对任意自同态$T$，$R = \sqrt{T^\star T}$是正定自伴自同态$T^\star T$的平方根。
任意给定向量$\mathbf{v}$，计算$\|T\mathbf{v}\|$：
\begin{align*}
    \|T\mathbf{v}\|^2 &= (T\mathbf{v} | T\mathbf{v}) = (T^\star T\mathbf{v} | \mathbf{v})  \\
    &= (\sqrt{T^\star T}\mathbf{v} | \sqrt{T^\star T}\mathbf{v})  \\
    &= \| \sqrt{T^\star T} \mathbf{v}\|^2 
\end{align*}
这说明，如果存在自同构$S$使得$S(\sqrt{T^\star T} \mathbf{v}) = T\mathbf{v}$，
那么$S$在$\sqrt{T^\star T}$的像空间上是保距映射，因此必须是幺正自同构。
\begin{pp}\label{pp:8_14}
    设$T$是内积空间$\mathbb{V}$中的自同态，则它总可以分解为：
    $$ T = U \circ R$$
    其中$U$是幺正自同构，$R = \sqrt{T^\star T}$是正定自伴自同态。这个分解方式称为自同态的\textbf{极分解}。
\end{pp}
\begin{proof} % [{\bfseries 证明\ref{pp:8_14}}]
    从前面讨论可知，只需证明存在幺正自同构$U$，
    使得$U(\sqrt{T^\star T} \mathbf{v}) = T\mathbf{v}$。\\
    要证明有这样的自同构，有两个关键。首先，要证明我们可以定义$\sqrt{T^\star T} \mathbf{v}$映射到$T\mathbf{v}$的函数$S$满足直性（因而是同态），
    其次，$\mathbf{v}$遍历整个空间$\mathbb{V}$的时候，$\sqrt{T^\star T} \mathbf{v}$只遍历了$\sqrt{T^\star T}$的像空间，
    $T\mathbf{v}$只遍历了$T$的像空间。
    如果$\sqrt{T^\star T}$和$T$都是自同构，它们的像空间就是$\mathbb{V}$，那自然最好。
    如果它们不是同构，就需要把$S$的定义域合理地扩展到整个空间上，变成$\mathcal{L}(\mathbb{V})$中的元素，并且证明它是幺正自同构。\\
    我们先来证明第一点。先证明$S$可以这么定义。
    也就是说，只要存在$\mathbf{v}$，使得$\mathbf{u} = \sqrt{T^\star T} \mathbf{v}$，那么$S\mathbf{u} = T\mathbf{v}$。\\
    为此，我们需要排除$\mathbf{u}$能以多种方式表示成$\sqrt{T^\star T} \mathbf{v}$，而它们在$T$的映射下结果不同的可能性。
    设$\mathbf{u} = \sqrt{T^\star T} \mathbf{v}_1 = \sqrt{T^\star T} \mathbf{v}_2$，我们要证明$T\mathbf{v}_1 = T\mathbf{v}_2.$\\
    这一点可以用前面论证的保距性质得到：
    \begin{align*}
        \|T\mathbf{v}_1 - T\mathbf{v}_2\| &= \|T(\mathbf{v}_1 - \mathbf{v}_2)\| = \| \sqrt{T^\star T}(\mathbf{v}_1 - \mathbf{v}_2)\|  \\
        &= \| \sqrt{T^\star T} \mathbf{v}_1 - \sqrt{T^\star T} \mathbf{v}_2\|  
    \end{align*}    
    这说明，只要$\mathbf{u} = \sqrt{T^\star T} \mathbf{v}_1 = \sqrt{T^\star T} \mathbf{v}_2$，
    我们就可以定义$S\mathbf{u} = T\mathbf{v}_1 = T\mathbf{v}_2.$ \\
    接着证明$S$的直性。设有$\mathbf{u}_1 = \sqrt{T^\star T} \mathbf{v}_1$，$\mathbf{u}_2 = \sqrt{T^\star T} \mathbf{v}_2$，
    那么对任意$\lambda$，
    $$ \mathbf{u}_1 + \lambda \mathbf{u}_2 = \sqrt{T^\star T} \mathbf{v}_1 + \lambda \sqrt{T^\star T} \mathbf{v}_2 = \sqrt{T^\star T} (\mathbf{v}_1 + \lambda \mathbf{v}_2).$$
    于是
    $$ S(\mathbf{u}_1 + \lambda \mathbf{u}_2) = T (\mathbf{v}_1 + \lambda \mathbf{v}_2) = T\mathbf{v}_1 + \lambda T\mathbf{v}_2 = S\mathbf{u}_1 + \lambda S\mathbf{u}_2.$$
    现在我们把$S$的定义域扩展到整个空间上。目前$S$在$\sqrt{T^\star T}$的像空间上有定义。
    任取向量$\mathbf{u} = T\mathbf{v}$，可以构造$\sqrt{T^\star T} \mathbf{v}$，所以$S$必然是满射。
    而且，由于$S$的保距性质，$S$是单射。这说明$S$是双射，$\rang T = \rang \sqrt{T^\star T}$。
    考虑$\Img T$和$\Img \sqrt{T^\star T}$的正交补，两个补空间也有相同的维数。
    于是，我们可以建立一个$\left(\Img \sqrt{T^\star T}\right)^\perp$到$\left(\Img T\right)^\perp$的同构$S^\perp$。\\
    为了使$S^\perp$和$S$连立起来之后变成一个保距同构，我们还需要保证$S^\perp$是保距同构。
    设$\mathcal{B} = (\mathbf{e}_i)_{1 \leqslant i \leqslant m}$是$\left(\Img \sqrt{T^\star T}\right)^\perp$的一组幺正基（$m = \rang \sqrt{T^\star T}$），
    $\mathcal{B}^\perp = (\mathbf{f}_i)_{1 \leqslant i \leqslant m}$是$\left(\Img T\right)^\perp$的一组幺正基，
    定义：
    $$ S^\perp(\sum_{i=1}^m a_i \mathbf{e}_i) = \sum_{i=1}^m a_i \mathbf{f}_i. $$
    直观来说，$S^\perp$就是把$\mathcal{B}$中向量和$\mathcal{B}^\perp$中向量按次序建立映射关系。
    容易证明这个映射是$\left(\Img \sqrt{T^\star T}\right)^\perp$到$\left(\Img T\right)^\perp$的同构，
    而且关于基底$\mathcal{B}$和$\mathcal{B}^\perp$的矩阵是单位矩阵，因此是保距映射。\\
    根据命题\ref{pp:3_1}，$\mathbb{V}$中有唯一的自同构$U$，在$\Img \sqrt{T^\star T}$和$\left(\Img \sqrt{T^\star T}\right)^\perp$上的限制分别是$S$和$S^\perp$。
    任意向量$\mathbf{x}$都可以唯一写成$\Img \sqrt{T^\star T}$和$\left(\Img \sqrt{T^\star T}\right)^\perp$中元素之和：
    $$ \mathbf{x} = \mathbf{u} + \mathbf{u}^\perp, \,\,\, \mathbf{u} \in \Img \sqrt{T^\star T} , \,\, \mathbf{u}^\perp \in \left(\Img \sqrt{T^\star T}\right)^\perp. $$
    于是，它经过$U$映射的结果为：
    \begin{align*}
        U\mathbf{x} = S\mathbf{u} + S^\perp\mathbf{u}^\perp = T\mathbf{v} + S^\perp\mathbf{u}^\perp 
    \end{align*}
    其中$\mathbf{u} = \sqrt{T^\star T} \mathbf{v}$。
    按照定义，$\mathbf{u} \perp \mathbf{u}^\perp$，$T\mathbf{v} \perp S^\perp\mathbf{u}^\perp$，
    所以$\| \mathbf{u} + \mathbf{u}^\perp \|^2 = \| \mathbf{u} \|^2 + \| \mathbf{u}^\perp \|^2$，
    $\| T\mathbf{v} + S^\perp\mathbf{u}^\perp \|^2 = \| T\mathbf{v} \|^2 + \| S^\perp\mathbf{u}^\perp \|^2$，
    因此，
    \begin{align*}
        \| \mathbf{x} \|^2 &= \| \mathbf{u} \|^2 + \| \mathbf{u}^\perp \|^2  \\
        &= \| T\mathbf{v} \|^2 + \| S^\perp\mathbf{u}^\perp \|^2 = \| U\mathbf{x} \|^2  
    \end{align*}
    这说明自同构$U$是幺正自同构。\\
\end{proof}

我们把一般的自同态分解成幺正自同构和正定自伴自同态的复合，而后两者都是具有特定性质的自同态，可以进一步研究。
这也是一种简化自同态的方法。从证明来看，复合的顺序并不重要。
也就是说，可以写成$T = UR$，也可以写成$T = RU$。当然，两个分解中的$U$未必相同。
此外，幺正自同构和正定自伴自同态都可以幺正对角化，但这不代表一般的自同态都可以幺正对角化了。
原因是两者对应的幺正基一般来说不是同一个基底。$U$和$R$在各自的本征向量构成的基底上是对角矩阵，
但两者不一定有共同的本征向量。

\subsection{表征值与表征分解}
研究自同态的平方根时，我们发现，任何自同态$Q$都对应一个正定自伴自同态$\sqrt{Q^\star Q}$。而根据谱定理，
$\sqrt{Q^\star Q}$总可以幺正对角化。这个性质值得我们重视。因为一般来说，自同态不一定有本征值；
即便复内积空间的自同态，也不一定能对角化，
而$\sqrt{Q^\star Q}$总有本征值，而且总可以幺正对角化。所以，我们希望$\sqrt{Q^\star Q}$能帮助我们了解$Q$的性质。

记$R = \sqrt{Q^\star Q}$，$R$可以幺正对角化。设$R$在本征向量构成的基底$\mathcal{B}$上是对角矩阵，
其对角元素是$R$的本征值。$R$作用在一般向量上，结果的模满足以下等式：
\begin{align*}
    \| R\mathbf{v} \|^2 &= (R\mathbf{v} | R\mathbf{v}) = (R^\star R\mathbf{v} | \mathbf{v})  \\
    &= (Q^\star Q\mathbf{v} | \mathbf{v}) = (Q\mathbf{v} | Q\mathbf{v})  \\
    &= \| Q\mathbf{v} \|^2 
\end{align*}
$R\mathbf{v}$和$Q\mathbf{v}$的模相等。把向量在基底$\mathcal{B}$上分解：
$$ \mathbf{u} = u_1\mathbf{e}_1 + \cdots + u_n\mathbf{e}_i,$$
其中$\mathbf{e}_i$是对应$R$的对角元素$d_{i.i}$的基向量，满足$R\mathbf{e}_i = d_{i.i}\mathbf{e}_i$。这样，
\begin{align*}
    \| Q\mathbf{v} \|^2 &= \| R\mathbf{v} \|^2 = \| \sum_{i=1}^n u_i R\mathbf{e}_i \|^2   \\
    &= \| \sum_{i=1}^n u_i d_{i,i}\mathbf{e}_i \|^2 = \sum_{i=1}^n |u_i|^2 |d_{i,i}|^2  
\end{align*}
由于$R$是正定自同态，对角元素$d_{i,i}$都是非负实数。不妨设$R$的对角元素从大到小排列：
$d_{1,1} \geqslant d_{2,2} \geqslant \cdots \geqslant d_{n,n}$，则
$$ d_{n,n} \| \mathbf{v} \| \leqslant \| Q\mathbf{v} \| \leqslant d_{1,1} \| \mathbf{v} \|. $$
可以看到，$R$的对角元素限定了$Q$作用在向量上的“放缩”效果\footnote{并不是真正的放缩，$Q\mathbf{v}$与$\mathbf{v}$不一定直相关。}。
如果向量$\mathbf{v}$取自单位超球面$\mathcal{S}_{n-1} = \{\mathbf{x} | \| \mathbf{x} \| = 1 \}$，即限定模为1，
那么$Q\mathbf{v}$的模总介于$d_{n,n}$与$d_{1,1}$之间。可以说，$R$的对角线元素决定了$Q$作为映射的“大小”。
因此，我们把$R$的本征值（也就是对角元素的值）称为自同态$Q$的\textbf{表征值}\footnote{也称为\textbf{奇异值}。}。
它们并不能直接指示$Q$的本质特征，但可以告诉我们$Q$的“表现”如何。
自同态的表征值一般记为$\sigma_1 \geqslant \cdots \geqslant \sigma_n$。

如何把$Q$表达成与表征值有关的式子呢？我们可以从极分解出发。我们已经知道，
$Q$可以写成$U\circ R$的形式，其中$U$是幺正自同构，$R$就是$\sqrt{Q^\star Q}$。而$R$的幺正对角化可以写成：
$[R]_\mathcal{B} = V^H \Sigma V$，其中$\mathcal{B}$是任意幺正基，$V$是幺正矩阵，$\Sigma$是对角元素为$\sigma_i$的对角矩阵。
因此$[Q]_\mathcal{B} = [U]_\mathcal{B} V^H \Sigma V$。
注意到$[U]_\mathcal{B}$和$V$都是幺正矩阵，所以它们的乘积仍然是幺正矩阵。记$W = [U]_\mathcal{B}V^H$，则$Q$可以写成：
$$ [Q]_\mathcal{B} = W \Sigma V. $$
其中$W$和$V$是幺正自同构，$\Sigma$是对角矩阵，其对角元素称为$Q$的表征值。这个分解称为$Q$的\textbf{表征分解}。

从向量映射的层面，$V$代表将$R$幺正对角化的幺正基变换。也就是说，如果使得$R$为对角矩阵的幺正基为$\mathcal{C}$，
则$V = [\mathrm{I}]_\mathcal{B}^\mathcal{C}$，而$W$表示$\mathcal{C}$到另一个幺正基$\mathcal{B}'$的幺正基变换。
于是，给定向量$\mathbf{u}$，记它在$\mathcal{C}$中的坐标是$(u_1, \cdots, u_n)$，则经过$\Sigma$对应的映射后，
坐标变为$(\sigma_1 u_1, \cdots, \sigma_n u_n)$。而$W$对应的基变换则保持坐标不变，将基底换为$\mathcal{B}'$。于是我们得到结论：
\begin{pp}\label{pp:8_15}
    设$T$是内积空间$\mathbb{V}$中的自同态，其表征值为$\sigma_1 , \cdots , \sigma_n$，则存在幺正基$(\mathbf{e}_i)_{1\leqslant i \leqslant n}$和
    $(\mathbf{f}_i)_{1\leqslant i \leqslant n}$，使得$T$作用在向量$\mathbf{u}$上的结果为：
    $$ T\mathbf{u} = \sum_{i=1}^n \sigma_i (\mathbf{u} | \mathbf{e}_i) \mathbf{f}_i.$$
\end{pp}
\begin{proof} % [{\bfseries 证明\ref{pp:8_15}}]
    只需将前面讨论中的对象和命题中的称呼对应起来就可以了。
    幺正基$(\mathbf{e}_i)_{1\leqslant i \leqslant n}$就是前述的$\mathcal{C}$。
    向量$\mathbf{u}$在$\mathcal{C}$上的坐标就是它与基向量的内积：$u_i = (\mathbf{u} | \mathbf{e}_i)$。
    因此，如前所述，$T$对向量的作用就是将坐标变为$(\sigma_1 u_1, \cdots, \sigma_n u_n)$，
    然后对基底做$W$对应的幺正基变换，得到新的幺正基，也就是命题中的$(\mathbf{f}_i)_{1\leqslant i \leqslant n}$。\\
\end{proof}

虽说这是表征分解的向量表示，但我们注意到，本质上，极分解和表征分解其实是一回事。如果我们写出极分解对应的向量变换公式，
结果与上式是一样的。

\subsectionext{一般同态的情况}
我们希望将以上的结论扩展到一般的直映射。记出发空间$\mathbf{V}$和到达空间$\mathbf{W}$的内积分别为$(\,\cdot\,|\,\cdot\,)_\mathbb{V}$和$(\,\cdot\,|\,\cdot\,)_\mathbb{W}$，
一般的直映射是否能写成类似极分解或表征分解的形式呢？

我们把自同态的伴随映射推广到一般同态。
\begin{pp}\label{pp:8_16}
    给定直映射$A\in\mathcal{L}(\mathbb{V}, \mathbb{W})$，唯一存在直同态$A^\star \in \mathcal{L}(\mathbb{W}, \mathbb{V})$，使得
    $$ \forall \mathbf{u}\in\mathbb{V}, \mathbf{v}\in\mathbb{W}, \quad (A\mathbf{u} | \mathbf{v})_\mathbb{W} = (\mathbf{u} | A^\star(\mathbf{v}))_\mathbb{V}. $$
    $A^\star$称为$A$的伴随映射。
\end{pp}
证明$A^\star$的存在性，方法和自同态的情况是一样的。对向量$\mathbf{v}\in\mathbb{W}$，
考虑$\mathbf{u} \mapsto (A\mathbf{u}|\mathbf{v})_\mathbb{W}$，根据对偶表示定理，
这个直型唯一对应$\mathbb{V}$中向量$\mathbf{b}$。$\mathbf{v}$到$\mathbf{b}$的映射关系就是$A$的伴随映射。
读者可以自行验证它是直同态。

前面几节提到的伴随映射的基本性质，一般同态的伴随映射与自同态的伴随映射也相同。
一般同态伴随映射的矩阵，也是原同态矩阵的共轭转置。当然，这里不再有“自伴”的概念，因为两者属于不同的空间，不可能相等了。
注意到伴随映射的矩阵和对偶映射的矩阵仍然相同，因此\ref{pp:8_6}的结论仍然适用。
只是等式两边的$\chi$要分别改为关联对应空间内积的同构：
$$ T^* \circ \chi_\mathbb{W} = \chi_\mathbb{V} \circ T^\star. $$

伴随矩阵扩展到一般同态后，对于一般的直同态$T$，我们也可以考察$T^\star T$。
它是$\mathbb{V}$中的自同态。不难验证它是自伴自同态，而且
$$ (T^\star T\mathbf{u}|\mathbf{u})_\mathbb{V} = (T\mathbf{u}|T\mathbf{u})_\mathbb{W} = \| T\mathbf{u} \|_\mathbb{W}^2 \geqslant 0. $$
也就是说$T^\star T$也是正定的自伴自同态。于是，我们仍然可以定义映射：
$$ S : \sqrt{T^\star T} \mathbf{u} \mapsto T\mathbf{u}. $$
注意这是一个$\Img \sqrt{T^\star T}\subseteq\mathbb{V}$到$\Img T\subseteq\mathbb{W}$的直映射。按照\ref{pp:8_14}的思路，我们验证它是否有保距性质：
\begin{align*}
    \|\sqrt{T^\star T} \mathbf{u}\|_\mathbb{V}^2 &= (T^\star T\mathbf{u}|\mathbf{u})_\mathbb{V}  \\
    &= (T\mathbf{u}|T\mathbf{u})_\mathbb{W} = \| T\mathbf{u} \|_\mathbb{W}^2  
\end{align*}
有了保距性质，我们就知道$S$是$\Img \sqrt{T^\star T}$到$\Img T$的双射。
不过，由于出发空间和到达空间维数不一定相同，所以我们不能保证把$S$扩展为一个同构了。
比如，如果出发空间的维数大于到达空间的维数，那么$S$扩展的结果显然不可能是单射。
反之，如果出发空间的维数小于到达空间的维数，那么$S$扩展的结果显然不可能是满射。

问题出在哪里呢？我们注意到，$T$本身是映射到$\mathbb{W}$的，$\sqrt{T^\star T}$的像空间却在$\mathbb{V}$中，
这就让$S$变成了一个从$\mathbb{V}$到$\mathbb{W}$的映射。如果我们给$\sqrt{T^\star T}$加一个“转换插头”$\iota$，
让映射$\iota \circ \sqrt{T^\star T}$把向量映射到$\mathbb{W}$里，那么$S$就再次成为一个自同态。
按照\ref{pp:8_14}中的论证，只要我们验证$\iota \circ \sqrt{T^\star T}$的容和$T$的容一样，
且$S$在$\iota \circ \sqrt{T^\star T}$的像空间上拥有保距性质，
就可以把$S$扩展成一个幺正自同构。

下面我们来确定合适的“转换插头”$\iota$。具体来说，我们只需要把对应的基向量“连接”起来就可以了。
我们在子空间$\Img T$中选一组幺正基$(\mathbf{f}_i)_{1 \leqslant i \leqslant r}$（$r$是$T$的容），将它扩展为整个空间的幺正基$(\mathbf{f}_i)_{1 \leqslant i \leqslant m}$。
如果到达空间的维数$m$比出发空间的维数$n$大，
那么选其中（包含$\Img T$的）$n$个基向量$(\mathbf{f}_i)_{1 \leqslant i \leqslant n}$，
把$\sqrt{T^\star T}$幺正对角化的基向量$(\mathbf{e}_i)_{1 \leqslant i \leqslant n}$按次序映射到这$n$个基向量。
如果到达空间的维数$m$比出发空间的维数$n$小，那么把$\sqrt{T^\star T}$幺正对角化的基向量中（包含$\Img \sqrt{T^\star T}$的）
$m$个基向量按次序映射到全部$m$个基向量$(\mathbf{f}_i)_{1 \leqslant i \leqslant m}$，
把剩余的$n-m$个基向量映射到$\mathbb{W}$中的零向量上。

这样定义了$\iota$后，我们来验证$\iota \circ \sqrt{T^\star T}$是否满足前面提到的两个性质。
容易看出，按照定义，$\iota$在$\Img \sqrt{T^\star T}$上的限制：$\restr{\iota}{\Img \sqrt{T^\star T}}$是“恒等”映射。
$\iota$只会在$n>m$的时候，把与$\Img \sqrt{T^\star T}$正交的基向量映射到零向量，
而$\Img \sqrt{T^\star T}$中的向量经过$\iota$映射，坐标不变。设$\Img \sqrt{T^\star T}$中元素：
$\mathbf{u} = \sqrt{T^\star T} \mathbf{v}$在出发空间的分解为：
$$ \mathbf{u} = \sum_{i=1}^r a_i \mathbf{e}_i,$$
按照定义，
$$ \iota\left(\mathbf{u}\right) = \sum_{i=1}^r a_i \mathbf{f}_i,$$
于是它们的模平方都是$\sum_{i=1}^r |a_i|^2$，因而模相等。这说明$\iota$本身在$\Img \sqrt{T^\star T}$上的限制有保距性质，故而是同构：
$$\rang \left(\iota\circ\sqrt{T^\star T}\right) = \rang \sqrt{T^\star T} = \rang T.$$
这证明了第一个性质。其次，$\iota$这种保距性质说明，对出发空间的任意向量$\mathbf{v}$，
$$\|\sqrt{T^\star T} \mathbf{v} \|_\mathbb{V} = \left(\sum_{i=1}^r |a_i|^2\right)^{\frac12} = \|\iota\left(\sqrt{T^\star T} \mathbf{v}\right) \|_\mathbb{W}.$$
而我们已经验证过，$\|\sqrt{T^\star T} \mathbf{v} \|_\mathbb{V} = \|T \mathbf{v} \|_\mathbb{W}$，
因此，$\| \iota\left(\sqrt{T^\star T} \mathbf{v}\right) \|_\mathbb{W} = \|T \mathbf{v} \|_\mathbb{W}$，
由于新定义的$S$为：
$$ S : \iota\left(\sqrt{T^\star T} \mathbf{v}\right) \mapsto T \mathbf{v}$$
$S$在$\iota \circ \sqrt{T^\star T}$的像空间上拥有保距性质。这证明了第二个性质。

综上所述，对于一般的直同态$T$，我们仍然可以将它做表征分解。从矩阵的角度，$T$的矩阵可以分解为：
$$ T = U \Sigma V.$$
其中，$V$是$n \times n$的幺正矩阵，$U$是$m \times m$的幺正矩阵，$\Sigma$是$m \times n$的（广义）对角矩阵。
具体来说，$\Sigma$的对角线为行数与列数相等的位置：第$1$行第$1$列，第$2$行第$2$列，等等。
$\Sigma$的对角元素为$\sqrt{T^\star T}$的本征值，也就是$T$的表征值。
如果$n=m$，则$\Sigma$就是$\sqrt{T^\star T}$幺正对角化得到的方阵$R$；
如果$n>m$，则$\Sigma$是$R$的一部分：从上数起的前$m$行；
如果$n<m$，则$\Sigma$是$R$和$(m-n) \times n$的零矩阵上下拼接得到的$m\times n$矩阵（$R$在上方）。

而从向量变换的角度，我们可以把\ref{pp:8_15}中的公式改一下：
\begin{pp}\label{pp:8_17}
    设$T$是内积空间$\mathbb{V}$到$\mathbb{W}$的直映射，其表征值为$\sigma_1 \geqslant \cdots \geqslant \sigma_n$，
    则存在$\mathbb{V}$的幺正基$(\mathbf{e}_i)_{1\leqslant i \leqslant n}$和$\mathbb{W}$的幺正基
    $(\mathbf{f}_i)_{1\leqslant i \leqslant m}$，使得$T$作用在向量$\mathbf{u}$上的结果为：
    $$ T\mathbf{u} = \sum_{i=1}^r \sigma_i (\mathbf{u} | \mathbf{e}_i) \mathbf{f}_i.$$
    其中$r$是$T$的容。
\end{pp}
从$\sqrt{T^\star T}$的容与$T$的容相等推出，$\sigma_{r+1} = \sigma_{r+2} = \cdots = \sigma_{n} = 0$。
因此我们只需关注前$r$个基向量。
换句话说，$(\mathbf{f}_i)_{1\leqslant i \leqslant r}$是$\Img T$的一组幺正基。按照\ref{pp:8_14}中的做法，可以将它扩展，
得到$\mathbb{W}$的幺正基，但这里我们只关心$\Img T$中的部分。对应到矩阵的角度，
$\Sigma$和$U$都只有左上角的$r\times r$的部分是与$T$相关的。$\Sigma$剩余的部分都是$0$，
而$U$对角线剩余的部分用$1$填充，使得它成为幺正自同构。

可以看出，\ref{pp:8_15}中的公式其实也是以上公式的特例。我们可以把\ref{pp:8_17}中的公式看成一般性的结论，
称它为直映射的表征分解公式。

\chapter{常见数域空间的自同态}
前几章中我们讨论了空间乃至内积空间中自同态的性质。我们可以看到，复空间中的自同态有良好的性质。
比如，复空间中的自同态总能三角化，复内积空间的自伴自同态乃至规范自同态可以规范正交对角化。

究其根本，复空间自同态的良好性质来自复数域$\mathbb{C}$的代数封闭性质，也就是代数基本定理告诉我们的：
（非常数的）复系数整式总有根。自同态的极小归零式总有根，说明存在本征值，
于是我们可以通过本征值和本征空间分析自同态的性质。

至于其它数域（如$\mathbb{Q}$和$\mathbb{R}$），由于它们并非代数闭域，有理系数（实系数）整式未必有有理数根（实根）。所以，如果自同态的极小归零式没有根，我们就无法通过本征值和本征空间来分析自同态的性质。

我们希望把问题化归为前面的情况，也就是说，我们希望自同态的根在系数域里。一个自然的想法是：我们把根“放到”系数域里，然后在这个扩大的系数域里讨论，就能够研究本征值和本征空间了。这方面的理论称为域扩张理论。我们在附录二里给出了域扩张理论的基本知识，以及利用域扩张讨论本章内容的前提结论。如果读者对域扩张的理论不熟悉，可以阅读附录二的相关内容。

\section{自同态的延拓}

给定系数域为$\mathbb{F}$的$n$维平直空间$\mathbb{V}$和系数属于$\mathbb{F}$的非常整式（下称$\mathbb{F}-$整式）$p$。根据域扩张理论，可以构建$\mathbb{F}\subseteq \mathbb{K}$，使得$p$在扩域里能分解为一次式的乘积（参见附录二的定理\ref{tm:9_5}）。换句话说，在$\mathbb{K}$中解方程$p(x)=0$，可以得到所有的根，所有的根都能在$\mathbb{K}$中找到。我们把这样的扩域称为$p$的\textbf{分裂域}。

$\mathbb{V}$中的直自同态$T$不一定有本征值，因为我们无法在$\mathbb{F}$中找到其极小归零式$m_T$的根。但如果我们把系数域扩张到$m_T$的分裂域，那么，由于$m_T$所有的根在分裂域中，我们可以把这些根和$T$联系起来。借助分裂域，我们可以为一般域上平直空间里的自同态定义本征值。

设$\mathbb{V}$中某个自同态$T$的极小归零式为$m_T$，它的分裂域为$\mathbb{K}$。我们希望在系数域为$\mathbb{K}$的平直空间中讨论问题。所以，我们希望把$T$看成$\mathbb{K}^n$中的自同态。

为此，我们首先选定$\mathbb{V}$中某个基底$\mathcal{B}$，则$\mathbb{V}$中的向量和自同态分别可以表示为元素在域$\mathbb{F}$中的$n\times 1$和$n\times n$矩阵。这些元素自然也在$\mathbb{K}$中。
如果我们选定$\mathbb{K}^n$的一组基，比如标准基，以上的矩阵就唯一确定了$\mathbb{K}^n$中的向量$\mathbf{v}^\mathbb{K}$和自同态$T^{\mathbb{K}}$。用这个方法，我们可以把平直空间$\mathbb{V}$中的向量和自同态延拓到$\mathbb{K}^n$中。

\begin{df}\label{df:9_3}
    给定域扩张$\mathbb{F}\subseteq\mathbb{K}$、$\mathbb{F}$上的$n$维平直空间$\mathbb{V}$，
    及$\mathbb{V}$的基底$\mathcal{B} = (\mathbf{e_i})_{1\leqslant i \leqslant n}$，
    可以构造平直空间$\mathbb{V}$到$\mathbb{K}^n$的单同态：
    \begin{align*}
        \Psi : \mathbb{V} &\rightarrow \mathbb{K}^n  \\
        \sum_{i=1}^n u_i\mathbf{e_i} &\mapsto  \sum_{i=1}^n u_i\mathbf{c_i} 
    \end{align*}
    其中$(\mathbf{c_i})_{1\leqslant i \leqslant n}$是$\mathbb{K}^n$的标准基。\\
    于是对任意向量$\mathbf{v} \in \mathbb{V}$，存在唯一向量$\mathbf{v}^\mathbb{K} = \pi(\mathbf{v})\in \mathbb{V}$，
    使得$\mathbf{v}^\mathbb{K}$在$\mathbb{K}^n$标准基上的坐标和$\mathbf{v}$在基底$\mathcal{B}$上坐标一样。\\
    对任意自同态$T \in \mathcal{L}(\mathbb{V})$，存在唯一自同态$T^{\mathbb{K}} \in \mathcal{L}(\mathbb{K}^n)$：
    $$ T^{\mathbb{K}}(\mathbf{c_i}) = \pi\left(T(\mathbf{e_i})\right) $$
    使得$T^{\mathbb{K}}$在$\mathbb{K}^n$标准基上的坐标和$T$在$\mathbb{V}$中某个基底$\mathcal{B}$的矩阵一样。\\
    我们把$\mathbf{v}^\mathbb{K}$、$T^{\mathbb{K}}$分别称为$\mathbf{v}$和$T$基于$\mathcal{B}$的、从$\mathbb{F}$到$\mathbb{K}$的\textbf{延拓}。
\end{df}
这个定义可以自然地推广到子空间上。只要把同态$\Psi$应用到子空间的基底上即可。
自同态的延拓$T^{\mathbb{K}}$作用在向量的延拓$\mathbf{v}^{\mathbb{K}}$上，和$T$作用在$\mathbf{v}$上，效果是一样的。
只需从矩阵的角度来看：同样的$n\times n$矩阵，乘以同样的$n\times 1$列向量，得到同样的$n\times 1$列向量。

$T^{\mathbb{K}}$，作为$T$的“扩展”，也有自己的极小归零式$m_T^{\mathbb{K}}$。
由于$m_T$是$T^{\mathbb{K}}$的归零式，$m_T^{\mathbb{K}}$必然是$m_T$的因式。下面我们证明两者其实相等。

\begin{pp}\label{pp:9_6}
    给定域$\mathbb{F}$上的$n$维平直空间$\mathbb{V}$及其中的自同态$T$，
    给定$\mathbb{F}$的扩域$\mathbb{K}$，给定$\mathbb{V}$的基底$\mathcal{B}$，
    考虑$T$（基于$\mathcal{B}$）的延拓$T^{\mathbb{K}}$，
    则作为$\mathbb{V}$中自同态，$T$的极小归零式$m_T$与$T^{\mathbb{K}}$作为$\mathbb{K}^n$中自同态的极小归零式$m_T^{\mathbb{K}}$相等。
\end{pp}
\begin{proof} % [{\bfseries 证明\ref{pp:9_6}}]
    我们已经说明$m_T^{\mathbb{K}}$是$m_T$的因式，下面证明$m_T$次数不大于$m_T^{\mathbb{K}}$。\\
    设$m_T$次数为$d$。将$\mathrm{I}, T, T^2, \cdots, T^k$看作平直空间$\mathcal{L}(\mathbb{V})$中的向量组，
    则$k \geqslant d$时，$\mathrm{I}, T, T^2, \cdots, T^k$直相关，
    而$k < d$时，$\mathrm{I}, T, T^2, \cdots, T^{k}$直无关。\\
    考虑它们的延拓。
    作为$\mathcal{L}(\mathbb{K}^n)$中的向量组，由命题\ref{pp:5_0}可知，$k < d$时，
    $\mathrm{I}, T^{\mathbb{K}}, \left(T^{\mathbb{K}}\right)^2, \cdots, \left(T^{\mathbb{K}}\right)^{k}$仍然直无关。
    这说明$m_T^{\mathbb{K}}$的次数不小于$d$。\\
    按照定义，$m_T^{\mathbb{K}}$和$m_T$都是首一整式。如果它们的系数不完全相同，则将两者相减，可以得到一个次数小于$d$的非零整式$p$，
    且$p$也是$T^{\mathbb{K}}$的归零式，这与$m_T^{\mathbb{K}}$的定义矛盾。
    所以$m_T^{\mathbb{K}}$和$m_T$系数完全相同，$m_T^{\mathbb{K}} = m_T$.\\
\end{proof}


这个结论让自同态的延拓成为分析自同态极小归零式和本征值的有力工具。
由于自同态的延拓不会改变其极小归零式，我们可以把自同态放到到“足够大”的域里面进行研究。
为了方便，以后我们把自同态延拓后的极小归零式也记为$m_T$。

给定域$\mathbb{F}$上的$n$维平直空间$\mathbb{V}$及其中的自同态$T$，设$\mathbb{K}$是$m_T$作为$\mathbb{F}$上整式的分裂域。
在$\mathbb{K}$中，$m_T$可以分解为一次式的乘积：
$$ m_T = a(z - \alpha_1)(z - \alpha_2)\cdots(z - \alpha_n), \,\, a \in \mathbb{F}, \,\, \alpha_1, \alpha_2, \cdots ,  \alpha_n \in \mathbb{K}, $$
其中的$\alpha_1, \alpha_2$等都是$m_T$的根，因此也是自同态$T$的延拓的本征值，
对应的本征向量在$\mathbb{K}^n$中。这样得到的本征值，如果不在$\mathbb{F}$中，就称为$T$的\textbf{隐本征值}，对应的本征向量为\textbf{隐本征向量}。
原本自同态$T$的本征值，在有隐本征值的上下文中，也可称为\textbf{显本征值}，对应的本征向量也可称为\textbf{显本征向量}。

更进一步，我们还可以定义对应的\textbf{隐广义本征值}、\textbf{隐广义本征向量}、\textbf{隐本征空间}、\textbf{隐广义本征空间}和\textbf{隐全本征空间}。

举例来说，设自同态$A\in\mathcal{L}\left(\mathbb{Q}^2\right)$在标准基上的矩阵为：
$$ [A]_\mathcal{C} = \begin{bmatrix} 0 & 2 \\ 1 & 0 \end{bmatrix}$$
可以验证$[A]_\mathcal{C}^2 - 2\mathrm{I}_2 = 0$，而整式$x^2 - 2$在$\mathbb{Q}$上是既约整式。
所以$A$的极小归零式就是$x^2 - 2$。

另一方面，$x^2 - 2$在$\mathbb{Q}$上的分裂域是$\mathbb{Q}(\sqrt{2})$，
它在$\mathbb{Q}(\sqrt{2})$中可以分解为：
$$ x^2 - 2 = (x - \sqrt{2})(x + \sqrt{2})$$
所以$A$有两个隐本征值：$\sqrt{2}$和$-\sqrt{2}$，相应的隐本征向量分别有$(\sqrt{2}, 1)$和$(\sqrt{2}, -1)$，
各自对应一个1维隐本征空间。

在$\mathbb{K}^n$中，我们可以更好地理解本征值之间的联系。

\section{根置换自同构和本征式}
\subsection{本征值之间的“对等关系”}
自同态$T$延拓到分裂域后，极小归零式$m_T$可以分解为一次式的乘积。
于是，通过研究$m_T$的根，也就是$T$的隐本征值和显本征值，就能够完全了解$T$。

给定两个不同的本征值$\lambda_1$和$\lambda_2$，对应的本征向量$\mathbf{v}_1$和$\mathbf{v}_2$满足：
$$T^{\mathbb{K}}\mathbf{v}_1 = \lambda_1\mathbf{v}_1, \quad T^{\mathbb{K}}\mathbf{v}_2 = \lambda_2 \mathbf{v}_2.$$
考虑这个等式在$\mathbb{K}^n$的标准基上的矩阵形式。本征值和本征向量$\mathbf{v}_1$和$\mathbf{v}_2$的坐标在分裂域$\mathbb{K}$中，
而$T^{\mathbb{K}}$的系数在$\mathbb{F}$中。如果有一个$\mathbb{K}$上的同构映射$\sigma$，把$\mathbb{F}$中元素映射到自身不变，
而把$\lambda_1$映射到$\lambda_2$，那么，$\sigma(T)$仍然是$T$，而$\sigma(\lambda_1) = \lambda_2$，$\sigma(\mathbf{v}_1) = \mathbf{v}_2$。
两个本征值的本征向量之间就有一个同构映射，两者的本征空间也同构。

什么时候这样的同构映射$\sigma$存在呢？我们有这样的结论：
\begin{pp}\label{pp:9_7}
    给定域$\mathbb{F}$上的整式$f$，设$\mathbb{L}$是它的分裂域。如果$\mathbb{F}$上的既约整式$p$在$\mathbb{L}$中有根：
    $$ p(\alpha) = p(\beta) = 0, \quad \alpha, \beta \in \mathbb{L}, $$
    那么存在域自同构$\sigma$，使得它在$\mathbb{F}$上的限制是恒等映射：
    $$ \restr{\sigma}{\mathbb{F}} = \mathrm{I}, $$
    且$\sigma(\alpha) = \sigma(\beta).$ 我们将这个自同构称为$f$的\textbf{根置换自同构}，
    具体来说，是把根$\alpha$映射到$\beta$的$\mathbb{F}$-不变同构。
\end{pp}
这个命题的证明涉及较多的技术性处理。我们将它的证明思路大致介绍一下：
\begin{proof} % [{\bfseries 证明\ref{pp:9_7}}]
    考虑单扩张$\mathbb{F} \subseteq \mathbb{F}(\alpha) \subseteq \mathbb{L}$，根据命题\ref{pp:9_3}，
    $$ \mathbb{F}(\alpha) \simeq \mathbb{F}[z] /_{\langle p \rangle}.$$
    具体来说，可以构造环同构$\varphi_\alpha$。$\varphi_\alpha$将$\mathbb{F}$的零元映射到$\bar{0}$，幺元映射到$\bar{1}$，
    将$\alpha$映射到$\overline{z}$，等等。\\
    同样地，考虑单扩张$\mathbb{F} \subseteq \mathbb{F}(\alpha) \subseteq \mathbb{L}$，根据命题\ref{pp:9_3}，
    $$ \mathbb{F}(\beta) \simeq \mathbb{F}[z] /_{\langle p \rangle}.$$
    可以构造环同构$\varphi_\beta$。$\varphi_\beta$将$\mathbb{F}$的零元映射到$\bar{0}$，幺元映射到$\bar{1}$，
    将$\beta$映射到$\overline{z}$，等等。\\
    于是，$\varphi_{\alpha\rightarrow\beta} = \varphi_\beta^{-1}  \circ \varphi_\alpha$就是一个把$\alpha$映射到$\beta$而在$\mathbb{F}$上是恒等映射的环同构。
    $\varphi_{\alpha\rightarrow\beta}$是从域$\mathbb{F}(\alpha)$到域$\mathbb{F}(\beta)$的环同构，所以是域同构。
    $\mathbb{F}(\alpha)$和$\mathbb{F}(\beta)$都是$\mathbb{L}$的子域，所以可以将它延拓成为$\mathbb{L}$到$\mathbb{L}$的域同构$\sigma$。
    $\sigma$就是我们要找的域自同构。\\
\end{proof}

回到上一节的例子。$\mathbb{L} = \mathbb{Q}(\sqrt{2})$是整式$x^2 - 2$在$\mathbb{Q}$上的分裂域，
而整式$x^2 - 2$是$\mathbb{Q}$上的既约整式，两个根$\sqrt{2}$和$-\sqrt{2}$都在$\mathbb{L}$里。所以，存在域自同构$\sigma$。
我们可以直接把$\sigma$写出来：
$$\sigma (a + b\sqrt{2}) = a - b\sqrt{2}.\quad a, b, \in \mathbb{Q}.$$
它限制在$\mathbb{Q}$上的部分是恒等映射，而且把$\sqrt{2}$映射到$-\sqrt{2}$。

考虑$A$的隐本征向量$(\sqrt{2}, 1)$和$(\sqrt{2}, -1)$。我们有：
$$ \begin{bmatrix} 0 & 2 \\ 1 & 0 \end{bmatrix} \begin{bmatrix} \sqrt{2} \\ 1 \end{bmatrix} = \sqrt{2} \begin{bmatrix} \sqrt{2} \\ 1 \end{bmatrix}.$$
$\sigma$作用在这个等式两端，得到：
$$ \begin{bmatrix} 0 & 2 \\ 1 & 0 \end{bmatrix} \begin{bmatrix} -\sqrt{2} \\ 1 \end{bmatrix} = -\sqrt{2} \begin{bmatrix} -\sqrt{2} \\ 1 \end{bmatrix}.$$
可以看到$\sigma$把$\sqrt{2}$的本征向量映射到$-\sqrt{2}$的本征向量。

一般来说，如果自同态$T$的本征值$\lambda_1$和$\lambda_2$是$m_T$的同一个既约因式的根，
而$\sigma_{1\rightarrow 2}$是命题\ref{pp:9_7}中将$\lambda_1$映射到$\lambda_2$的$\mathbb{F}$-不变同构。
我们可以构想，对全本征空间$H_{T^{\mathbb{K}}}(\lambda_1)$中向量$\mathbf{v}_1 \in \mathbb{K}^n$，设它满足：
$$ (T^{\mathbb{K}} - \lambda_1\mathrm{I})^i \mathbf{v}_1 = \mathbf{0}. $$
考虑这个等式在$\mathbb{K}^n$的标准基$\mathcal{C}$上的矩阵形式。考虑等式中每个坐标和元素经过$\sigma_{1\rightarrow 2}$映射的结果，可以得到：
$$ ([T^{\mathbb{K}}]_\mathcal{C} - \lambda_2\mathrm{I}_n)^i \sigma_{1\rightarrow 2}([\mathbf{v}_1]_\mathcal{C}) = \mathrm{0}_n. $$
这是一个不严谨的写法。其中的“$\sigma_{1\rightarrow 2}([\mathbf{v}_1]_\mathcal{C})$”表示把$\mathbf{v}_1$在标准基上的坐标经过$\sigma_{1\rightarrow 2}$变换得到的新向量。
$\sigma_{1\rightarrow 2}([\mathbf{v}_1]_\mathcal{C})$对应着一个$H_{T^{\mathbb{K}}}(\lambda_2)$中的向量。

我们现在将上面这个想法用严谨的语言表达。具体来说，我们把域自同构“补全”为直同构：$\overline{\sigma}_{1\rightarrow 2}$。

定义$\mathbb{K}^n$到$\mathbb{K}^n$的自同态$\overline{\sigma}_{1\rightarrow 2}$，
它把$\mathbb{K}^n$中的向量在标准基$(\mathbf{c}_1, \mathbf{c}_2, \cdots , \mathbf{c}_n)$上的坐标做$\sigma_{1\rightarrow 2}$变换，
得到新的向量：
$$ \overline{\sigma}_{1\rightarrow 2} \left(\sum_{i=1}^n v_i \mathbf{c}_i\right) = \sum_{i=1}^n \sigma_{1\rightarrow 2} \left( v_i\right) \mathbf{c}_i.$$
可以验证，$\overline{\sigma}_{1\rightarrow 2}$是一个直同构。

前面我们论证的结果就可以写成：
\begin{pp}\label{pp:9_7_0}
设$T$是系数域$\mathbb{F}$上的$n$维空间$\mathbb{V}$中的自同态，$m_T$是它的极小归零式；设$\mathbb{K}$为$m_T$的分裂域，
而$\lambda_1$和$\lambda_2$是$m_T$的同一个$\mathbb{F}$-既约因式的根，那么，两者所有对应的广义本征空间同构：
$$ \forall \,\, i \in \mathbb{N}, \quad \Ker (T^{\mathbb{K}} - \lambda_1\mathrm{I})^i \cong (T^{\mathbb{K}} - \lambda_2\mathrm{I})^i. $$
它们的本征空间同构：
$$E_{T^{\mathbb{K}}}(\lambda_1) \cong E_{T^{\mathbb{K}}}(\lambda_2),$$
它们的全本征空间也同构：
$$H_{T^{\mathbb{K}}}(\lambda_1) \cong H_{T^{\mathbb{K}}}(\lambda_2).$$
它们的本重数、全重数和阶都一样：
$$ \mu_g(\lambda_1) = \mu_g(\lambda_2), \,\,\, \mu_a(\lambda_1) = \mu_a(\lambda_2), \,\,\, d_{\lambda_1} = d_{\lambda_2}.$$
\end{pp}

\subsection{一般域上自同态的极小归零式}
考虑$m_T$在$\mathbb{F}$上的因式分解：
$$ m_T = p_1^{d_1}\cdots p_k^{d_k}.$$
其中$p_1, \cdots , p_k$是$\mathbb{F}$上的既约整式。将$T$延拓到$m_T$的分裂域之后，$m_T$分解为：
$$ m_T = (x - \lambda_1)^{d_{\lambda_1}}\cdots (x - \lambda_m)^{d_{\lambda_m}}.$$
其中$m \geqslant k$，每个$x - \lambda_i$都是某个$p_j$的因式。

从上面的讨论，我们知道，每个$p_j$都可以写成若干个$x - \lambda_i$的乘积。注意到$p_1, \cdots , p_k$两两互素，
所以不同的$p_j$没有共同的一次因式$x - \lambda_i$。于是，每个$p_j$都对应等式：
$$ p_j^{d_j} = \prod_{i\in \mathfrak{s}(j)} (x - \lambda_i)^{d_{\lambda_i}},$$
其中$\mathfrak{s}(j)$表示属于$p_j$的一次因式$x - \lambda_i$的下标集合。

$d_{\lambda_i}$是本征值$\lambda_i$的阶，所以对同属一个既约整式的本征值来说，是同一个数。
记$d_{\lambda_i} = d$，则：
$$ p_j^{d_j} = \left(\prod_{i\in \mathfrak{s}(j)} (x - \lambda_i)\right)^d. $$
由于$\prod_{i\in \mathfrak{s}(j)} (x - \lambda_i)$整除$p_j$，$d_j$整除$d$。
记$d = d_j \cdot \nu_j$，则$\nu_j$整除$\deg p_j$，而且$\mathfrak{s}(j)$有$\deg p_j / \nu_j$个元素:
$$\quad p_j = \left(\prod_{i\in \mathfrak{s}(j)} (x - \lambda_i)\right)^{\nu_j}. $$
零空间引理告诉我们：
$$ \forall \,\,l \in\mathbb{N},\quad  \Ker p_j^{l}(T^{\mathbb{K}}) = \bigoplus_{i\in \mathfrak{s}(j)} \Ker (T^{\mathbb{K}} - \lambda_i\mathrm{I})^{\nu_j l}.$$
设：
$$ H_{T^{\mathbb{K}}}(j,l) = \bigoplus_{i\in \mathfrak{s}(j)} \Ker (T^{\mathbb{K}} - \lambda_i\mathrm{I})^{l},$$
我们知道，随着自然数$l$不断增加，子空间$\Ker (T^{\mathbb{K}} - \lambda_i\mathrm{I})^{l}$不断扩大，最终在$l = d_{\lambda_i}$处稳定。
根据根置换自同构，以上过程对所有$\lambda_i$一致。所以，如果考虑子空间套：
$$ \{\mathbf{0}\} = H_{T^{\mathbb{K}}}(j,0) \subseteq H_{T^{\mathbb{K}}}(j,1) \subseteq H_{T^{\mathbb{K}}}(j,2) \subseteq \cdots $$
它的扩大方式与前者类似。每次扩大，各个$\Ker (T^{\mathbb{K}} - \lambda_i\mathrm{I})^{l}$的维数同时增加至少1维，
所以$H_{T^{\mathbb{K}}}(j,l)$增加$\frac{\deg p_j}{\nu_j}$维，或其倍数。
于是，子空间$\Ker p_j^{l}(T^{\mathbb{K}}) = H_{T^{\mathbb{K}}}(j,\nu_j l)$的维数总是$\deg p_j$或其倍数，且这个维数增加的“步调”与各个子空间$\Ker (T^{\mathbb{K}} - \lambda_i\mathrm{I})^{l}$
的“步调”一致。也就是说，两者都在$l = d_j$处稳定。

最终，子空间$\Ker p_j^{d_j}(T^{\mathbb{K}}) = H_{T^{\mathbb{K}}}(j, \nu_j d_j) = H_{T^{\mathbb{K}}}(j, d)$，而它的维数
$\nul p_j^{d_j}(T^{\mathbb{K}}) = \sum_{i\in \mathfrak{s}(j)} \mu_a(\lambda_i) = \frac{\mu_a(\lambda_i) \deg p_j}{\nu_j} $是$\deg p_j$的倍数。
这说明$\mu_a(\lambda_i)$是$\nu_j$的倍数。记$\mu_a(\lambda_i) = \nu_j \mu_a(j)$，则
$$ \nul p_j^{d_j}(T^{\mathbb{K}}) = \mu_a(j) \deg p_j. $$
与全本征空间的情形一样，考虑$d_j$和$\mu_a(j)$的关系，可以发现：
\begin{align*}
    d_j &= \frac{d_{\lambda_i}}{\nu_j} \leqslant \frac{\mu_a(\lambda_i)}{\nu_j} = \mu_a(j).  
\end{align*}

因此，我们可以这样总结一般域上的自同态$T$的极小归零式$m_T$的性质：
\begin{pp}\label{pp:9_7_1}
不考虑排列顺序，$m_T$可以唯一分解为若干个$\mathbb{F}$上既约整式的乘积。
$$ m_T = p_1^{d_1}\cdots p_k^{d_k}.$$
空间可以分解为：
$$ \mathbb{V} = \Ker p_1^{d_1}(T) \oplus \cdots \oplus \Ker p_k^{d_k}(T). $$
其中，既约整式$p_j$的幂次$d_j$是子空间套：
$$ \{\mathbf{0}\} = \Ker p_j^{0}(T) \subseteq \Ker p_j^{1}(T) \subseteq \Ker p_j^{2}(T) \subseteq \cdots $$
首次稳定时的幂次。\\
子空间$\Ker p_j^{d_j}(T)$的维数是$\deg p_j$的倍数。如果记
$$ \nul p_j^{d_j}(T) = \mu_a(j) \deg p_j, $$
$\mu_a(j)$也称作$p_j$的全重数。所有全重数的和，乘以$\deg p_j$作为权重，等于空间的维数：
$$ \sum_j \mu_a(j) \deg p_j = n . $$
\end{pp}

此外，$d_j  \leqslant \mu_a(j) $，所以极小归零式的次数小于等于空间的维数。

\subsection{本征式}
当自同态$T$可三角化的时候，我们定义过它的本征式$c_T$。对于可三角化的矩阵，其本征式为$n$次首一整式，且等于其所有本征值对应的一次式按全重数的乘积。我们还得到结论：$c_T(T) = 0$。
自同态延拓到分裂域后，它的极小归零式可以分解为一次式的乘积，所以根据命题\ref{pp:6_8}，总可以三角化。同理，
我们可以讨论它的本征式。下面从使用本征值定义本征式出发。

\begin{df}\label{df:9_5}
    给定域$\mathbb{F}$上的$n$维平直空间$\mathbb{V}$以及其中自同态$T$，它的极小归零式对应分裂域$\mathbb{K}$。
    设$T$到$\mathbb{K}$的延拓为$T^{\mathbb{K}}$。则从$T^{\mathbb{K}}$三角化后矩阵的对角线元素$\lambda_1, \cdots , \lambda_n$可构造整式：
    $$ c_T = \prod_{i=1}^n (x - \lambda_i) = \prod_{\lambda} (x - \lambda)^{\mu_a(\lambda)}.$$
    $c_T$称为$T$的\textbf{本征式}。
\end{df}
以上定义中，$c_T$似乎依赖于$\mathbb{K}$，它的系数在$\mathbb{K}$中。但我们可以证明：$c_T$的系数在$\mathbb{F}$中。
这就说明，自同态的本征式是唯一的。

比照我们之前定义本征式的方法：
$$
c_T = \det(x\mathrm{I} - T).
$$
自同态可三角化的时候，我们知道两种定义是等价的。现在，我们通过把自同态延拓到分裂域，就能让$T$可三角化。于是，在分裂域中，我们有
$$
c_T = \prod_{i=1}^n (x - \lambda_i) = \det(x\mathrm{I} - T).
$$
但注意到$\det(x\mathrm{I} - T)$是关于$x$的多项式，而其系数，按照效的定义，就是$x\mathrm{I} - T$在任意基底上的矩阵系数。于是我们选取其在原本系数域$\mathbb{F}$的矩阵，则其系数就是$\mathbb{F}$的元素。因此，$c_T$的系数就在$\mathbb{F}$中。

我们也可以直接根据本征值的性质推出这个结论。以下的证明较为繁琐抽象，仅供读者对比参考。
\begin{pp}\label{pp:9_8}
    自同态$T$的本征式是系数在$\mathbb{F}$中的首一$n$次整式。
\end{pp}
\begin{proof} % [{\bfseries 证明\ref{pp:9_8}}]
    按照定义，$c_T$是首一$n$次整式。只需证明它的系数在$\mathbb{F}$中。\\
    首先，根据本征式的定义，它的根的集合是三角化后对角线元素的集合。根据命题\ref{pp:6_7}，它的根的集合就是$T^{\mathbb{K}}$所有本征值的集合。\\
    其次，根据\ref{pp:6_12_4}及其推论，空间的分解
    $$ \mathbb{K}^n = \Ker (T^{\mathbb{K}} - \lambda_1\mathrm{I})^{d_1} \oplus \Ker (T^{\mathbb{K}} - \lambda_2\mathrm{I})^{d_2} \oplus \cdots \oplus \Ker (T^{\mathbb{K}} - \lambda_k\mathrm{I})^{d_k}$$
    中，每个全本征空间$\Ker (T^{\mathbb{K}} - \lambda_i\mathrm{I})^{d_i}$都可以表示为单位标准形式，因此对应的是对角线上
    $\mu_a(\lambda_i)$个$\lambda_i$。于是本征式可以写为：
    $$c_T = \prod_{i} (x - \lambda_i)^{\mu_a(\lambda_i)}.$$
    注意到这些$x - \lambda_i$也是之前极小归零式的因式分解中的各个既约因式$p_j$在$\mathbb{K}[z]$中的因式，
    $c_T$和$m_T$的不同仅仅在于$x - \lambda_i$的次数不同。同一个既约整式$p_j$分解出的一次因式中的（隐）本征值$\lambda_i$，
    其全重数相同。因此，由于同属一个$p_j$的一次因式可以合写为：
    $$ p_j = \left(\prod_{i\in \mathfrak{s}(j)} (x - \lambda_i)\right)^{\nu_j}, $$
    所以，
    \begin{align*}
        \left(\prod_{i\in \mathfrak{s}(j)} (x - \lambda_i)\right)^{\mu_a(\lambda_i)} &= p_j^{\frac{\mu_a(\lambda_i)}{\nu_j}} = p_j^{\mu_a(j)}. 
    \end{align*}
    于是，我们得到：    
    $$c_T = \prod_{j} p_j^{\mu_a(j)}.$$
    其中的$p_j$是$T$的极小归零式$m_T$在$\mathbb{F}$上因式分解得到的既约因式，$\mu_a(j)$是前文讨论中定义的全重数。\\
    $p_j$是系数在$\mathbb{F}$中的整式，所以$c_T$也是系数在$\mathbb{F}$中的整式。\\
\end{proof}

既然$c_T$总可以按照可三角化的情形定义，那么一般情况的\textbf{全本征归零定理}对任何系数域空间上的自同态都成立。
\begin{tm}{\textbf{全本征归零定理}}
    对一般域上的自同态$T$，定义$c_T = \det(x\mathrm{I} - T)$，则有：
    $$ c_T(T) = 0.$$
\end{tm}

让我们把$c_T$具体写出来：
$$c_T(x) = \prod_{i=1}^n (x - \lambda_i) = (-1)^n a_0 + (-1)^{n-1} a_1 x + \cdots - a_{n-1}x^{n-1} + x^n.$$
其中$\lambda_i$是$T$的（隐）本征值。根据多项式的性质，所有$n$个本征值的和等于$a_{n-1}$，所有$n$个本征值的乘积等于$a_0$。所以$a_{n-1}$就是自同态的迹。而$a_0$是所有$n$个本征值的乘积。另一方面，它等于$(-1)^nc_T(0)$，也就是$\det(T)$。于是$a_0$等于$T$的效。换句话说，自同态的效就是按全重数排列的$n$个（隐）本征值的乘积。

\section{实空间自同态的约简}
本节我们把以上这些结论应用到实系数空间$\mathbb{R}^n$上。

\subsection{实空间自同态的标准形式}
第八章中，我们讨论了内积空间中的自同态。谱定理告诉我们，内积空间中的自伴自同态总可以在某个幺正基上表示为对角矩阵。
这里，我们用自同态延拓的方法，给出实内积空间中谱定理的另一个证明。

我们把实内积空间中的自伴自同态$T$进行延拓，从$\mathbb{R}$出发，将$T$延拓到$\mathbb{C}$，
得到一个复内积空间中的自伴自同态$T^{\mathbb{C}}$。后者可以幺正对角化，说明其极小归零式
是没有重根的一次式乘积：
$$m_{T^{\mathbb{C}}}(z) = (z -\lambda_1)\cdots(z - \lambda_k)$$
而我们已知这个极小归零式就是$T$的极小归零式：
$$m_T = m_{T^{\mathbb{C}}}.$$
这说明$m_T$也可以表示为没有重根的一次式乘积，所以$T$可对角化，从而可幺正对角化。

那么，实空间$\mathbb{R}^n$中的其他自同态，是否能对角化呢？我们考虑$\mathbb{R}^2$中的旋转映射$\mathrm{R}_2(\theta)$。
显然，对于一般的$\theta$，比如$\frac{\pi}{3}$，$\frac{\pi}{5}$，（非零）向量经过映射后不会与原向量直相关。
这说明不存在本征向量，更遑论本征向量构成的基底。这表示，实空间中一般的自同态无法对角化。

那么，我们可以将实空间中一般的自同态约简到什么程度呢？零空间引理告诉我们，自同态可以表示为准对角矩阵。
准对角矩阵由分块的矩阵构成，对角线上每块小矩阵都对应自同态极小归零式的既约因式。
因此，我们希望将实空间中自同态的极小归零式尽可能分解。
我们知道，实系数的二次整式并不总能分解为一次整式的乘积。比如$z^2+2$就是不可约的。
但我们可以证明，更高次的实系数整式总是可约的。

要证明这一点，首先需要看到实数域和复数域之间的关系。
\begin{pp}\label{pp:9_9}
    复数域$\mathbb{C}$是实数域$\mathbb{R}$的二维扩张：$[\mathbb{C} : \mathbb{R}] = 2$.
\end{pp}
这一点容易从$\mathbb{C} = \mathbb{R}(\imath )$推得。
$(1, \imath )$就是$\mathbb{C}$作为$\mathbb{R}$上平直空间的一组基。

设$p$是既约实系数整式，$\deg p = n > 1$。根据代数基本定理，$p$的$n$个根都在$\mathbb{C}$里。
于是根据命题\ref{pp:9_4}，$p$对应实数域的某个单扩张：$\mathbb{R}\subseteq\mathbb{R}(\alpha)$的极小归零式，
且$\mathbb{R}(\alpha)\subseteq\mathbb{C}$。因此，根据塔定理，
$$ 2 = [\mathbb{C} : \mathbb{R}] = [\mathbb{C} : \mathbb{R}(\alpha)] [\mathbb{R}(\alpha) : \mathbb{R}]$$
这说明扩张$\mathbb{R}\subseteq\mathbb{R}(\alpha)$的维数$[\mathbb{R}(\alpha) : \mathbb{R}]$小于等于2。
而这个维数就是$\deg p$。
\begin{pp}\label{pp:9_10}
    实系数既约整式必然是一次或二次整式。
\end{pp}
我们在中学已经了解了一次和二次的实系数整式。一次实系数整式必然是既约整式。
二次实系数整式是既约整式，当且仅当它没有实根。这时它的判别式小于0，有一对共轭的复数根。

根据以上结论，对实空间中的自同态$T$的极小归零式$m_T$做因式分解，结果是一次式和二次整式的乘积。
根据零空间引理，我们可以把空间表示成：
$$ \mathbb{V} = \bigoplus_{i}\Ker (T - \lambda_i I)^{d_i} \bigoplus_{j} \Ker \left(T^2 + b_jT + c_j\mathrm{I}\right)^{d_j} . $$
对于形如$\Ker (T - \lambda I)^{d}$的不变子空间，我们已经有充分了解。
现在我们来研究形如$\Ker \left(T^2 + b T + c\mathrm{I}\right)^{d}$的不变子空间。

首先，如果$d = 1$，那么子空间表示为：
$$ \Ker \left(T^2 + bT + c\mathrm{I}\right) $$
其中$b$和$c$是实数系数。对应的整式$p(z) = z^2 + bz + c$没有实根，也就是说，$b^2 < 4c$。
我们记正实数$c = r^2$，则$|b| < 2r$。记$\cos\theta = -\frac{b}{2r}$，则$p$相应写成$z^2 - 2r\cos\theta z + r^2 = (z - r\cos\theta)^2 + (r\sin\theta)^2$。
对这样一个子空间中的向量$\mathbf{u}$来说，按照定义，$p(T)(\mathbf{u}) = \mathbf{0}$。
设$T\mathbf{u} - r\cos\theta \mathbf{u} = r\sin\theta \mathbf{v}$，
则通过计算可以发现，$T\mathbf{v} = -r\sin\theta \mathbf{u} + r\cos\theta \mathbf{v}$。
于是，向量组$(\mathbf{u}, \mathbf{v})$生成的二维子空间是$T$-不变子空间。$T$在这个空间的限制可以用矩阵表示为：
$$ B = \begin{bmatrix}
    r\cos\theta & -r\sin\theta \\ r\sin\theta & r\cos\theta
\end{bmatrix}.$$
它又可以写成：
$$ \begin{bmatrix}
    r\cos\theta & -r\sin\theta \\ r\sin\theta & r\cos\theta
\end{bmatrix} = \begin{bmatrix}
    \cos\theta & -\sin\theta \\ \sin\theta & \cos\theta
\end{bmatrix} \cdot \begin{bmatrix}
    r & 0 \\ 0 & r
\end{bmatrix}.$$
这说明$T$在这个二维不变子空间上的限制可以看作一个旋转映射和一个放缩映射的复合。
这个子空间同构于$\mathbb{R}^2$，也就是我们熟知的平面。
映射$T$把平面中的向量放缩$r$倍，然后旋转$\theta$角度\footnote{这里我们并没有引入内积，所以这里的“角度”并不能理解为直角坐标系中的角度。}。

如果我们把$T$延拓为复空间$\mathbb{C}^n$中的映射$T^\mathbb{C}$，
就可以看得更清楚。$T^\mathbb{C}$在复空间中也有对应的二维不变子空间。
$z^2 + r\cos\theta z + r^2 = (z - r\e^{\imath \theta})(z - r\e^{-\imath \theta})$是没有复根的一次式乘积，
所以它在这个子空间里可以对角化，表示为对角矩阵：
$$ \begin{bmatrix}
    r\e^{\imath \theta} & 0 \\ 0 & r\e^{-\imath \theta}
\end{bmatrix}.$$
可以验证，两个隐本征值是共轭复数，它们对应的隐本征向量分别就是$\mathbf{u}+\imath \mathbf{v}$
和$\mathbf{u}-\imath \mathbf{v}$，也是两个共轭向量。
这个现象可以用我们前面提到的根置换自同构解释。这里涉及的同构就是共轭变换$z \mapsto \bar{z}$。
它作用在实数系数上是恒等变换，且把$\imath $的极小归零式$z^2+1$的根$\imath $映射到另一个根$-\imath $。
这个同构表明，如果某个复数是自同态的本征值，那么它的共轭也是本征值。这说明，实系数自同态的（隐）本征值要么是实数，
要么是一对共轭复数。前者对应形如$\Ker (T - \lambda I)^{d}$的不变子空间，后者对应形如$\Ker \left(T^2 + b T + c\mathrm{I}\right)^{g}$的不变子空间。

这让我们想到，如果$d > 1$，我们面对的其实是一对同构的隐广义本征空间：$H_{T^\mathbb{C}}(\lambda)$和$H_{T^\mathbb{C}}(\bar{\lambda})$。
它们对应着$T$延拓后对应的不变子空间\footnote{读者可以用零空间引理验证。}：
$$\Ker \left(\left(T^\mathbb{C}\right)^2 + b T^\mathbb{C} + c\mathrm{I}\right)^{d} = H_{T^\mathbb{C}}(\lambda) \oplus H_{T^\mathbb{C}}(\bar{\lambda}). $$
设$\mathcal{B}_\lambda = (\mathbf{e}_1, \cdots , \mathbf{e}_k)$是让$T$在$H_T(\lambda)$表示为单位标准形式的基底。
在$\mathcal{B}_\lambda$上，$T$在不变子空间$H_{T^\mathbb{C}}(\lambda)$上的限制可以用单位标准形式表示：
$$ J(\lambda) = \begin{bmatrix}
    \lambda & 1 & 0 & \cdots & 0 \\
    0 & \lambda & 1 & \cdots & 0 \\
    \vdots & \vdots & \ddots & \ddots & \vdots  \\
    0 & 0 & 0 & \cdots & 1  \\
    0 & 0 & 0 & \cdots & \lambda \\
   \end{bmatrix} $$
对应的，我们有基底$\mathcal{B}_{\bar{\lambda}} = (\overline{\mathbf{e}_1}, \cdots , \overline{\mathbf{e}_k})$，
对应$H_{T^\mathbb{C}}(\bar{\lambda})$上的单位标准形式：
$$ J(\bar{\lambda}) = \begin{bmatrix}
    \bar{\lambda} & 1 & 0 & \cdots & 0 \\
    0 & \bar{\lambda} & 1 & \cdots & 0 \\
    \vdots & \vdots & \ddots & \ddots & \vdots  \\
    0 & 0 & 0 & \cdots & 1  \\
    0 & 0 & 0 & \cdots & \bar{\lambda} \\
   \end{bmatrix} $$
从前面二维子空间的例子可知，$\mathbf{e}_1$和$\overline{\mathbf{e}_1}$构成一个二维子空间。
它对应两个实向量$\mathbf{u}_1 = \frac{\mathbf{e}_1 + \overline{\mathbf{e}_1}}{2} = \Re\mathbf{e}_1$
和$\mathbf{u}_1 = \frac{\mathbf{e}_1 - \overline{\mathbf{e}_1}}{2\imath } = \Im\mathbf{e}_1$。
而对应的矩阵就是前面看到的旋转-放缩复合矩阵$B$。其中的$r$和$\theta$分别对应复数$\lambda$的模长和幅角。
我们把它记作$B_\lambda$

接下来看$\mathbf{e}_2$和$\overline{\mathbf{e}_2}$。从单位标准形式可得：
\begin{align*}
    T^\mathbb{C}\mathbf{e}_2 = \mathbf{e}_1 + \lambda \mathbf{e}_2  \\
    T^\mathbb{C}\overline{\mathbf{e}_2} = \overline{\mathbf{e}_1} + \bar{\lambda} \overline{\mathbf{e}_2} 
\end{align*}
设$\mathbf{e}_2 = \mathbf{u}_2 + \imath \mathbf{v}_2$，则我们可以通过计算验证，
\begin{align*}
    T \mathbf{u}_2 = \mathbf{u}_1 + r\cos\theta \mathbf{u}_2 + r\sin\theta \mathbf{v}_2  \\
    T \mathbf{v}_2 = \mathbf{v}_1 - r\sin\theta \mathbf{u}_2 + r\cos\theta \mathbf{v}_2 
\end{align*}
其中的$r$和$\theta$分别对应复数$\lambda$的模长和幅角。
这样，我们如果查看$\mathbf{u}_1, \mathbf{v}_1, \mathbf{u}_2, \mathbf{v}_2$构成的子空间，会发现它也是
$T$-不变子空间，且$T$在它上面的限制可以表示为矩阵：
$$ \begin{bmatrix}
    r\cos\theta & -r\sin\theta & 1 & 0 \\ r\sin\theta & r\cos\theta & 0 & 1 \\
    0 & 0 & r\cos\theta & -r\sin\theta \\ 0 & 0 & r\sin\theta & r\cos\theta
\end{bmatrix} = \begin{bmatrix}
    B_\lambda & \mathrm{I}_2 \\ 0_2 & B_\lambda
\end{bmatrix}.$$
以此类推，我们可以得到一个$2k \times 2k$的矩阵：
$$ J_\mathbb{R}(\lambda) = \begin{bmatrix}
    B_\lambda & \mathrm{I}_2 & 0 & \cdots & 0 \\
    0 & B_\lambda & \mathrm{I}_2 & \cdots & 0 \\
    \vdots & \vdots & \ddots & \ddots & \vdots  \\
    0 & 0 & 0 & \cdots & \mathrm{I}_2  \\
    0 & 0 & 0 & \cdots & B_\lambda \\
   \end{bmatrix} $$
这个矩阵形式上和单位标准形式非常相似，只是“大”了一倍，恰好对应原来的两个单位标准形式$J(\lambda)$和$J(\bar{\lambda})$。
它是一个分块矩阵。单位标准形式对角线的位置从本征值变成了$2\times 2$的矩阵$B_\lambda$，
而主对角线上一格的元素从$1$变成了二维单位矩阵$\mathrm{I}_2$。它对应的基底$\mathcal{B}_\lambda^\mathbb{R}$就是
$J(\lambda)$对应基底$\mathcal{B}_\lambda$中基向量的实部和虚部向量按顺序成对排列起来：
$$\mathcal{B}_\lambda^\mathbb{R} = (\Re\mathbf{e}_1, \Im\mathbf{e}_1, \Re\mathbf{e}_2, \Im\mathbf{e}_2, \cdots \Re\mathbf{e}_k, \Im\mathbf{e}_k)$$
我们把矩阵$J_\mathbb{R}(\lambda)$称为\textbf{二次单位标准形式}。于是，我们得到了实空间中一般自同态的标准形式：
\begin{pp}\textbf{实空间自同态的标准形式 }\label{pp:9_11}
    若$T$是实空间$\mathbb{R}^n$中的自同态，则它可以在恰当的基底上表示成准对角矩阵。
    其对角元素要么是单位标准形式，要么是二次单位标准形式。
\end{pp}

\subsection{实内积空间的自同态}
复内积空间中，我们证明了自同态能幺正对角化的充要条件是它与它的伴随自同态交换。我们把这样的自同态称为规范自同态。
实内积空间中，这样的规范自同态对应怎样的约简形式呢？

给定实内积空间$\mathbb{R}^n$，记其中内积为$(\,\cdot\,|\,\cdot\,)$，定义其中规范自同态为与其伴随自同态交换的自同态。
我们把这样定义的规范自同态$T$延拓到复空间$\mathbb{C}^n$，希望利用复空间的规范自同态的性质来研究它。为此，我们需要把实空间的内积也延拓到$\mathbb{C}^n$上。
我们定义$(\,\cdot\,|\,\cdot\,)_\mathbb{C}$：
\begin{align*}
    \forall \,\, \mathbf{a}, \mathbf{b}, \mathbf{c}, \,\, \in \mathbb{R}^n , \,\,\, (\mathbf{a} + \imath  \mathbf{b}|\mathbf{c})_\mathbb{C} = (\mathbf{a}|\mathbf{c}) + \imath (\mathbf{b}|\mathbf{c})  \\
    \forall \,\, \mathbf{a}, \mathbf{b}, \mathbf{c}, \,\, \in \mathbb{R}^n , \,\,\, (\mathbf{a}|\mathbf{b} + \imath  \mathbf{c})_\mathbb{C} = (\mathbf{a}|\mathbf{b}) - \imath (\mathbf{a}|\mathbf{c}) 
\end{align*}
不难验证这样定义的$(\,\cdot\,|\,\cdot\,)_\mathbb{C}$满足复空间内积的定义。它对应的模$\| \, \cdot \, \|_\mathbb{C}$
满足$\| \mathbf{u} + \imath  \mathbf{v}\|_\mathbb{C}^2 = \|\mathbf{u}\|^2 + \|\mathbf{v}\|^2$。

注意到$T$延拓后得到的$T^\mathbb{C}$保持了与$T$相同的矩阵，所以它们在各自空间中的伴随矩阵也对应着这个矩阵的共轭转置。
这说明$T^\star$和$T^{\mathbb{C}\star}$有相同的矩阵表示。或者说，$T^{\mathbb{C}\star} = T^{\star\mathbb{C}}$。
$T$作为规范自同态，满足$T\circ T^\star = T^\star \circ T$，用矩阵的方式描述，就是两者的矩阵乘积可以交换顺序。
而它们延拓后的复自同态各自对应相同的矩阵，所以可以推出，延拓后的复自同态$T^\mathbb{C}$也是复内积空间中的规范自同态。
根据谱定理\ref{pp:8_9}，$T^\mathbb{C}$可以幺正对角化，其对角元素是$T^\mathbb{C}$的本征值。
而根据命题\ref{pp:9_10}，它的本征值要么是实数，要么是成对的共轭复数。
我们前面已经论证过，一对共轭本征值对应一个$\mathbb{R}^n$中的二维不变子空间，且$T$在其上的限制为旋转-放缩矩阵。
于是，我们可以把实内积空间中的规范自同态约简：
\begin{pp}\textbf{实内积空间的谱定理 }\label{pp:9_12}
    若$T$是实空间$\mathbb{R}^n$中的规范自同态，则它可以在恰当的基底上表示成准对角矩阵。
    其对角元素要么是$1\times 1$的矩阵，要么是$2\times 2$的矩阵：
    $$ \begin{bmatrix}
        r\cos\theta & -r\sin\theta \\ r\sin\theta & r\cos\theta
    \end{bmatrix}.$$
\end{pp}
可以说，实内积空间中的规范自同态也可以近似“对角化”。

作为规范自同态的特例，实内积空间中的幺正自同态也可以表示成这样的准对角矩阵。
我们已经讨论过二维实内积空间中的幺正自同态。
它要么是旋转，要么是镜面反射。旋转的矩阵就是$\mathrm{R}_2(\theta)$，
而镜面反射的矩阵总能在某个正交基上表示为：
$$ \begin{bmatrix}
    1 & 0 \\ 0 & -1
\end{bmatrix}.$$
我们现在可以讨论一般$n$维实内积空间上的幺正自同态$U$。$U$可以在适当的基底下表示为准对角矩阵，
其对角元素要么是$1\times 1$的矩阵，要么是$2\times 2$的矩阵。由于幺正自同态的本征值模长总是$1$，
对应的$1\times 1$的矩阵要么是$1$，要么是$-1$，而$2\times 2$的矩阵总是旋转矩阵。所以它总可以写成：
$$ \begin{bmatrix}
    \mathrm{I}_{m^+} & 0 & 0 \\
    0 & -\mathrm{I}_{m^-} & 0 \\
    0 & 0 & D_2 \\
   \end{bmatrix} $$
其中$m^+$和$m^-$是自然数，$D_2$是一个以$2\times 2$旋转矩阵为对角元素的准对角矩阵。

从几何角度来看，$U$对向量的作用效果可以分为三种。在某些方向上，$U$是恒等变换；
在某些方向上，它使向量方向反转，“大小”不变；在另外某些“平面”上，$U$把向量旋转一定角度。
这样，我们对实内积空间的幺正自同态就有了直观的理解。

上一章中，我们研究过二维实内积空间中的幺正自同态。现在我们来看三维实内积空间中的幺正自同态。
根据以上讨论，三维实内积空间中的幺正自同态要么是对角元素绝对值为1的对角矩阵，
要么是一个$1\times 1$的矩阵和一个$2\times 2$的矩阵为对角元素的准对角矩阵。

首先来看对角矩阵的情况。如果忽略基向量排列的顺序，我们可以按$-1$的个数将其分为$4$个矩阵：
$$ \begin{bmatrix}
    1 & 0 & 0 \\
    0 & 1 & 0 \\
    0 & 0 & 1 \\
\end{bmatrix}, \,\,\, \begin{bmatrix}
    -1 & 0 & 0 \\
    0 & 1 & 0 \\
    0 & 0 & 1 \\
\end{bmatrix}, \,\,\, \begin{bmatrix}
    -1 & 0 & 0 \\
    0 & -1 & 0 \\
    0 & 0 & 1 \\
\end{bmatrix}, \,\,\, \begin{bmatrix}
    -1 & 0 & 0 \\
    0 & -1 & 0 \\
    0 & 0 & -1 \\
\end{bmatrix}$$
而右边两个矩阵无非分别是左边两个矩阵乘以$-1$。记左起第二个矩阵为$\mathrm{Q}_3$，
则对角矩阵的情形可以归结为$\mathrm{I}_3$，$\mathrm{Q}_3$，$-\mathrm{Q}_3$和$-\mathrm{I}_3$。
它们分别对应在配备了直角坐标系的三维空间中的恒等变换和沿各个坐标轴（一个或多个）的镜面反射。

再来看准对角矩阵的情况。我们不妨设$1\times 1$的矩阵总在左上方。它的值只能是$1$或$-1$。
右下方的$2\times 2$矩阵只能是旋转矩阵$\mathrm{R}_2(\theta)$。于是我们得到两种矩阵：
$$ \begin{bmatrix}
    1 & 0 & 0 \\
    0 & \cos\theta & -\sin\theta \\
    0 & \sin\theta & \cos\theta \\
\end{bmatrix}, \,\,\, \begin{bmatrix}
    -1 & 0 & 0 \\
    0 & \cos\theta & -\sin\theta \\
    0 & \sin\theta & \cos\theta \\
\end{bmatrix}$$
而后者是前者乘以$\mathrm{Q}_3$。所以，记前者为$\mathrm{R}_3(\theta)$，则准对角矩阵分为
$\mathrm{R}_3(\theta)$和$\mathrm{Q}_3\mathrm{R}_3(\theta)$两类。前者是以$x$轴为轴心的旋转（也称\textbf{涡旋}），
后者是在此基础上沿$x$轴反射。这些映射，加上最简单的平移，就构成了三维实内积空间中保持距离和夹角不变的映射。

\subsection{有理空间自同态的标准形式}
以上对标准形式的推广可以再进一步。我们以有理空间$\mathbb{Q}^n$为例。$\mathbb{Q}^n$中自同态的极小归零式可以
有任意次数的既约因式。我们考察形如$\Ker p(T)^{d}$的不变子空间，其中$p = a_0 + a_1z + \cdots + a_{m-1}z^{m-1} + z^m$是既约整式。

如果$d = 1$，在空间中任取非零向量$\mathbf{v}$，
考察使得$q(T)(\mathbf{v}) = \mathbf{0}$的非零整式$q$中次数最小的，记为$p_\mathbf{v}$，称为向量$\mathbf{v}$（关于$T$）的\textbf{极小归零式}。
不难证明$p_\mathbf{v}$必然是$p$的因式，而$p$是既约整式，所以$p_\mathbf{v} = p$。
考虑$T$-不变子空间$\langle T^{m-1}\mathbf{v}, T^{m-2}\mathbf{v}, \cdots , \mathbf{v} \rangle$，$T$在其上表示为矩阵：
$$ W_p = \begin{bmatrix}
    -a_{m-1} & 1 & 0 & \cdots & 0 \\
    -a_{m-2} & 0 & 1 & \cdots & 0 \\
    \vdots & \vdots & \ddots & \ddots & \vdots  \\
    -a_1 & 0 & 0 & \cdots & 1  \\
    -a_0 & 0 & 0 & \cdots & 0 \\
   \end{bmatrix} $$
这是一个$m\times m$的矩阵（$m$为$p$的次数）。它的第$i$行第$i+1$列是1，第一列从上到下是$p$除最高次项之外系数的相反数。
这个矩阵形式和我们第六章中分析幂零自同态时得到的形式是一样的。给定$d$阶幂零自同态$N$，如果把空间看作$\Ker N^d$，
那么上面的矩阵其实就是从某个向量轨迹出发得到的循环不变子空间。所以，这里的矩阵无非是一般自同态关于向量$\mathbf{v}$的循环不变子空间$C_T(\mathbf{v})$。
我们称这个形式为向量$\mathbf{v}$的\textbf{显形式}或\textbf{显形矩阵}，也记为$C_T(\mathbf{v})$。向量的极小归零式系数显露在它的矩阵元素中。

$C_T(\mathbf{v})$是$T$-不变子空间，且$C_T(\mathbf{v})\subseteq\Ker p(T)$。实空间的情形中，我们得到二维不变子空间$B_\lambda$，
而在有理空间中，我们得到的是$m$维不变子空间。对一般的情形$\Ker p(T)^{d}$，如果把$T$延拓到复空间，那么$\Ker p(T)^{d}$对应若干个
单位标准形式，每个单位标准形式中的对角元素是$p$的复根。考虑$p$的复根之间的关系，命题\ref{pp:9_7}说明任意两个根之间都存在根置换自同构，
于是，我们可以比照实空间的情形处理。

和实空间的情形一样，我们有
$$\Ker p(T)^{d} = \bigoplus_{i=1}^m (T - \lambda_i \mathrm{I})^d. $$
其中$\lambda_i$是$p$的复根，而$(T - \lambda_i \mathrm{I})^d$在复空间中对应单位标准形式和相应的基向量。
对每个单位标准形式左上角对角元素对应的基向量，都可以通过根置换找到一组共$m$个向量。
它们都是隐本征向量，构成的$m$维不变子空间就是某个循环不变子空间，可以写成显形矩阵$W_p$的形式，对应某个轨迹$C_1$。
于是，单位标准矩阵左起第二列的向量，也可以写成相关的$m$个向量$C_2$与$C_1$的直组合。
$C_1$和$C_2$合起来构成一个$2m$维的$T$-不变子空间，$T$在其中的$2m\times 2m$矩阵为
$$ \begin{bmatrix}
    W_p & \mathrm{I}_m \\
    0 & W_p \\
   \end{bmatrix} $$
以此类推，我们最终得到如下形式的矩阵：
$$ J_\mathbb{Q}(p) = \begin{bmatrix}
    W_p & \mathrm{I}_m & 0 & \cdots & 0 \\
    0 & W_p & \mathrm{I}_m & \cdots & 0 \\
    \vdots & \vdots & \ddots & \ddots & \vdots  \\
    0 & 0 & 0 & \cdots & \mathrm{I}_m  \\
    0 & 0 & 0 & \cdots & W_p \\
   \end{bmatrix} $$
这就是有理空间中自同态的单位标准形式。和零空间引理直接得到的分解相比，我们对每块对角元素的内容都有了确定的了解。
与复空间、实空间不同的是，有理空间中的既约整式没有特定的形式，
我们没有把复空间标准形式的基向量转为有理空间中对应基向量的统一方法。


%%%%%%%%%%%%%%%%%%%%%%%%%%%%%%%%%%%
%%%         end of text         %%%
%%%%%%%%%%%%%%%%%%%%%%%%%%%%%%%%%%%

\begin{appendix}

\chapter{代数基础知识}
开始本书的主题之前，我们首先介绍一些重要的基本概念。它们在接下来的章节中会不时出现，作为定义概念的工具、帮助理解的例子，或在习题中用到。
这些概念，读者在高中数学学习中或有触及。如果你还没有接触过这些概念，我们可以一起来了解一下。

\section{数系的扩张}
什么是数？数是人类在劳动和生活中对一类逻辑概念的抽象。
在各种场景下，人们发现，需要一套抽象系统来精确描述度量关系。
比如，买卖中，我们需要点清货物有多少，好决定交易是否让自己满意。
战争中，将军需要衡量敌我的力量和物资的配备，才能赢得战争。
数，是人们在各种生产活动、交流中提炼出的一套抽象逻辑系统，
能准确描述各种度量关系。

最基本的数，是从数数过程中抽象出来的：一、二、三…… 我们把这些数称为\textbf{自然数}。
各个早期文化都发展出了自然数的概念，并用不同的方式记载在自己的语言中，是为数字和数词。
伴随自然数出现的是加减乘除运算。我们一般称为四则运算。

自然数的除法，衍生出了分数，后来称为\textbf{有理数}。自然数也可以看作是分母为1的分数。
从自然数到分数，是数的概念的第一次扩张。

分数对乘法和除法是封闭的（分数乘以或除以分数，仍然是分数），除了一点：0不能作分母。

对于许多古代文明，自然数和分数已经够用了。比如古埃及文明就长期满足于分数的概念。
古希腊的毕达哥拉斯学派更是认为“万物皆数”，所有的度量关系，都可以表示为“两个自然数的比”。
数构成的系统，称为“数系”。分数构成的系统现在称为有理数系，其中包含了自然数。
从自然数系到有理数系，是上古时期人类对度量关系逐步理解的成果。

然而，随着几何学的发展，人们在处理形体对象的度量关系时发现，一些有理数乃至自然数开方的结果不再是有理数了。
这个发现让公元前4世纪的古希腊人十分困惑，开启了数学史上的“第一次数学危机”。
古希腊人不知道如何解释：面积为2的正方形，
其边长无法表示为两个自然数的比，甚至不是任何已知的数进行四则运算的结果。
这对当时的毕达哥拉斯学派造成了很大冲击。

毕达哥拉斯学派的解决方法，是把这种不是自然数之比的数命名为“量”。
“数”是离散的，不连续的，而“量”是连续的。
“数”具备“不可分性”，从“一”衍生而来，而“量”则是无限可分的。
“数”存在于计数和比例中，而“量”存在于直观形体关系如边长、角度中。
可以表达为自然数之比的量，称为“可公比量”，不能表达成自然数之比的，称为“不可公比量”。

在公元前3世纪左右成书的《几何原本》中，欧几里得也沿用了这个做法，
采用直观形体的“量”作为构建数学的基本概念。
《几何原本》标志着“万物皆数”的神秘主义哲学转变为以“量”和演绎逻辑为中心的公理体系。

既然“量”和“数”是区别开的，并且“量”似乎更为本质。而“量”总是从基于“数”和四则运算的等式中得到。
那么，这类被称为“方程”的等式，就成了数学家研究的对象。“量”被看作是“满足方程的对象”，被称为方程的\textbf{根}。

中国古代的数学家更注重计算方法，即如何找出满足方程的数。
当遇到需要开方的时候，他们把结果分为“开尽”和“开不尽”两种情况，
类比为“除尽”和“除不尽”的区别，注重于找出近似的数。

欧亚大陆的另一边，数学家则注重研究“数”和“量”的关系。
丢番图对方程的研究被认为是代数学的雏形。所谓代数学，就是研究“数”和“代数”关系的学问。
什么是“代数”关系呢？我们来看以下的几个式子：

$$ \frac{\sqrt[2]{2} - 1}{2}, \quad \frac{\sqrt[2]{5} - 1}{2}, \quad \frac{\sqrt[3]{5} - 1}{2}$$

其中的$-1$和$\frac{\cdot}{2}$是“数”和四则运算，而$\sqrt[2]{2}, \sqrt[2]{5}, \sqrt[3]{5}$等是“量”。
如果我们要验证以上几个“量”，哪个是方程：$ X\cdot(1 + X) - 1 = 0 $的根，
这时候，我们关心的是将这几个“量”代入方程中是否能使等式成立。
比如：
\begin{align*}
    (\frac{\sqrt[2]{2} - 1}{2})\times(1 + \frac{\sqrt[2]{2} - 1}{2}) - 1 &= \frac{\sqrt[2]{2} - 1}{2} \times \frac{\sqrt[2]{2} + 1}{2} - 1  \\
    &= \frac{\sqrt[2]{2}\times \sqrt[2]{2} - 1}{4} - 1  \\
    &= \frac{\sqrt[2]{2}\times \sqrt[2]{2} - 5}{4}  
\end{align*}
上式要等于$0$，则分子须为$0$，即$\sqrt[2]{2}\times \sqrt[2]{2} = 5$. 而$\sqrt[2]{2}$的定义不满足这个条件。

我们发现，上面的运算和$\sqrt[2]{2}$没有关系。把$\sqrt[2]{2}$换成$\sqrt[2]{5}$或$\sqrt[3]{5}$，我们仍然会进行同样的计算。
所以，把以上运算过程抽象出来，我们做的事情其实是：
\begin{align*}
    (\frac{a - 1}{2})\times(1 + \frac{a - 1}{2}) - 1 &= \frac{a - 1}{2} \times \frac{a + 1}{2} - 1  \\
    &= \frac{a\times a - 1}{4} - 1  \\
    &= \frac{a\times a - 5}{4}  
\end{align*}
上式中的$a$可以是$\sqrt[2]{2}$，也可以是$\sqrt[2]{5}$乃至任何的量。
以上处理的关系与$a$的值无关，
$a$可以代表或代替一个“数”或“量”，参与四则运算；运算完毕后，再将$a$的值代入。
这就是代数关系，以上的运算称为\textbf{代数运算}。

代数学的发展，为数系添加了新的成员。中世纪时期，阿拉伯数学家逐渐认识到：他们研究的“数”和“量”是统一的。
它们都是方程的根。现在我们把这类数称为\textbf{代数数}。

代数数是否就是所有的数呢？对这个问题的探索分为两个方向。
在代数数的基础上，数系又经过了两次扩展：\textbf{实数}和\textbf{复数}。

实数的概念源于对圆周率的研究。由于人们无法把圆周率对应到任何方程中，
十六世纪以后，对圆周率的研究通过构造有理数的无限逼近来进行。这就是无穷级数。

通过对无穷级数、极限和无穷小量的研究，数学家对
“数”的概念有了新的认识。随着微积分的出现，
人们认识到，“通过无穷步计算得到的数”和无穷小量等概念，都需要严格定义。
十九世纪以来，现代的分析学就在这个基础上产生了。随着对极限、无穷计算的严格定义，数学家能够严格地定义“任意大小的数”：
任何数都是某个递增有理数列的极限。这个定义下的数，包括了人们直觉认知中所有“真实存在的数”，所以被称为\textbf{实数}。
自然数、有理数、一部分代数数，都包含在这个概念之中。

然而，还有另一类数，不在实数的概念里。这就是\textbf{虚数}。

虚数最基本的例子，是$-1$的平方根。这个概念在很长一段时间内不被数学家接受。
当时，数学界一般认为，这是一个不存在的概念，因为直观上它要表示面积为$-1$的正方形的边长。
而“面积为$-1$的正方形”是不存在的。现实中的面积不可能是负数。

然而，卡尔丹在求三次方程的根的时候，遇到了虚数。他需要使用方程
$$X(10 - X) = 40$$
的根，来表示某个三次方程的根。这意味着需要使用：$\frac{5 + \sqrt{-15}}{2}$这样的数。
而从这个数出发，经过一番代数运算后，得到的三次方程的根却是可以理解的代数数。

卡尔丹认为这种数是“诡辩量”，这种技巧“既不可捉摸又没有什么用处”。
莱布尼兹称其为“介乎于存在与不存在之间的两栖数”。不过，随着对数系结构的进一步了解，
数学家们对虚数的认识也发生了变化：数是否对应着有现实意义的量并不重要，重要的是
在此之上构造的抽象结构是否逻辑自洽。

这个转变，标志着近代数学研究进入了一个新的阶段，研究的对象从“对现实的抽象”，
转变为单纯“逻辑上的抽象结构”。比如，负数不再仅仅被看作是“减号以及减号后面作为减数的正数”，和“作为被减数的正数”区分，
而被看作独立的概念，和正数一起被纳入数系中。整数、有理数、代数数的概念得以补全。
同样地，\textbf{复数}也作为抽象概念被学者正视。十九世纪起，高斯、阿贝尔、雅可比、柯西等人
对复数进行了深入研究，构建了复数上的分析学。

从接纳负数，到接纳复数，数学家对“数”和“代数结构”的认识突飞猛进，
以抽象代数结构为工具，解决了三等分角问题、五次方程代数解问题等很多长久以来困惑数学界的难题。从复数开始，数学家不断发展出如四元数、$p$进数等新“数”，数系的扩展不再沿着单一方向，百花齐放。

本书中将用以下记号表示各个数系：
\begin{center}
    \fbox{
        \shortstack[l]{
            $\mathbb{N}$：自然数系 \\
            \vspace{2pt}\\
            $\mathbb{Z}$：整数系\\
            \vspace{2pt}\\
            $\mathbb{Q}$：有理数系\\
            \vspace{2pt}\\
            $\mathbb{A}$：代数数系\\
            \vspace{2pt}\\
            $\mathbb{R}$：实数系\\
            \vspace{2pt}\\
            $\mathbb{C}$：复数系
        }
    }
\end{center}

\section{群、环、域}

上一节讲到，人类对数系的认识从自然数到复数，
经历了从实际经验抽象出数学概念，到研究有实际意义的抽象量，再到研究逻辑上的抽象结构的过程。
本节我们对几个基本的抽象代数结构作初步的介绍。

以下，我们把\textbf{代数结构}定义为“配备了运算法则的集合”，由一个集合，及若干个定义在集合元素上的运算法则构成。
举例来说，自然数系$\mathbb{N}$可以看作是集合$\{0,1,2,\ldots\}$和加法构成的代数结构，
或者是集合$\{0,1,2,\ldots\}$和加法、乘法构成的代数结构。
同一个集合，配备不同的运算法则，可以看作不同的代数结构。

\begin{df}\label{df:0_1}
    配备了运算$\circ$，并满足以下条件的集合$G$，称为\textbf{群}：

    \begin{enumerate}
        \item 运算$\circ$为$G \times G \rightarrow G$的二元封闭运算；
        \item 运算$\circ$满足结合律：$\forall a, b, c\,\, \in\,\, G, \,\,\,\, (a\circ b) \circ c = a\circ (b \circ c).$
        \item 运算$\circ$有\textbf{空元}（也称为\textbf{单位元}）：$\exists \,\, e \,\, \in \,\, G$，使得$\forall a \,\, \in \,\, G, \,\,\,\, a \circ e = e \circ a = a. $
        \item 每个元素关于$\circ$都有\textbf{逆元}：$\forall a \,\, \in \,\, G, \,\,\, \exists \,\, b \,\, \in \,\, G$，使得$a \circ b = b \circ a = e.$ $b$称为$a$的\textbf{逆元}。
    \end{enumerate}
    也称$G$关于$\circ$构成群。
\end{df}

群是一种宽泛的代数结构，只要求配备一种运算法则。
我们用整数系$\mathbb{Z}$作为例子：

$\mathbb{Z}$关于整数的加法构成群。验证如下：
\begin{enumerate}
    \item 加法为$\mathbb{Z} \times \mathbb{Z} \rightarrow \mathbb{Z}$的二元封闭运算；
    \item 加法满足结合律：$\forall a, b, c\,\, \in\,\, \mathbb{Z}, \,\,\,\, (a + b) + c = a + (b + c).$
    \item 加法有空元$0$：$\forall a \,\, \in \,\, \mathbb{Z}, \,\,\,\, a + 0 = 0 + a = a. $
    \item 每个整数都有加法逆元：$\forall a \,\, \in \,\, \mathbb{Z}, \,\,\, \exists \,\, b = -a \,\, \in \,\, \mathbb{Z}$，使得$a + b = b + a = 0.$ 
\end{enumerate}

给定一个集合$S$，所有从$S$到$S$的双射（一一映射）组成的集合$\mathcal{F}(S)$，关于映射的复合运算$\circ$构成群。验证如下：

\begin{enumerate}
    \item 双射复合之后仍然是双射，所以运算$\circ$为$\mathcal{F}(S) \times \mathcal{F}(S) \rightarrow \mathcal{F}(S)$的二元封闭运算；
    \item 运算$\circ$满足结合律：$\forall f, g, h\,\, \in\,\, \mathcal{F}(S), \,\,\,\, \forall x \,\, \in \,\, S, \,\,\,\,[(f\circ g) \circ h](x) = (f\circ g)(h(x)) = f(g(h(x))) = f((g\circ h)(x)) = [f\circ (g \circ h)](x).$\\
          所以$(f\circ g) \circ h = f\circ (g \circ h).$
    \item 运算$\circ$有空元：$e = \mathrm{I}_{S} \,\, \in \,\, \mathcal{F}(S)$，其中$\mathrm{I}_{S}$是$S$上的恒等映射，定义为：$\forall x\in S,\,\,\, \mathrm{I}_{S}(x) = x$。\\
          这是因为$\forall f \,\, \in \,\, \mathcal{F}(S), \,\,\,\,\forall \,\, x \,\, \in \,\, S,\,\,\,\, (f \circ e)(x) = f(e(x)) = f(x) = e(f(x)) = (e\circ f)(x). $ \\
          这说明$\forall f \,\, \in \,\, \mathcal{F}(S), \,\,\,\, f\circ e = e\circ f = f.$
    \item $\forall f \,\, \in \,\, \mathcal{F}(S)$，$f$是双射，所以必然有逆映射。$f$的逆映射就是$f$关于函数复合运算的逆元。
\end{enumerate}

如果群运算满足交换律，就称之为\textbf{交换群}。以上例子中，$\mathrm{Z}$关于加法构成交换群，
而双射集合关于映射复合运算一般不构成交换群。

如果集合$G$中不是每个元素关于二元运算$\circ$都有逆元，但满足其他条件，
则称其构成\textbf{中半群}\footnote{也称为\textbf{幺半群}}。
如果连空元都没有，则称其构成\textbf{半群}。

在群的基础上，再配备一种二元运算，可以得到什么结构呢？

\begin{df}\label{df:0_2}
    配备了运算$\circ$和$\star$，并满足以下条件的集合$R$，称为\textbf{环}：

    \begin{enumerate}
        \item $\circ$和$\star$都是$R \times R \rightarrow R$的二元封闭运算；
        \item $R$关于运算$\circ$构成交换群；
        \item $R$关于运算$\star$构成中半群；
        \item 运算$\star$对运算$\circ$满足分配律：$\forall \,\, a, b, c \,\, \in \,\, R, \,\,\,\, (a \circ b) \star c = (a \star c) \circ (b \star c), \,\,\, c\star (a \circ b) = (c \star a) \circ (c \star b).$
    \end{enumerate}
    也称$R$关于$\circ, \star$构成环。
\end{df}

在环的定义中，我们不要求第二种运算满足交换律。
如果环的第二种运算也满足交换律，就称之为\textbf{交换环}。

为了区别，我们将第一种运算的空元称为\textbf{零元}，将第二种运算的空元称为（乘法）\textbf{幺元}。

环有哪些例子呢？$\mathbb{Z}$不仅关于加法构成群，还关于加法、乘法构成环。验证如下：

\begin{enumerate}
    \item 加法和乘法都是$\mathbb{Z} \times \mathbb{Z} \rightarrow \mathbb{Z}$的二元封闭运算；
    \item $\mathbb{Z}$关于加法构成交换群；
    \item $\mathbb{Z}$关于乘法构成中半群：乘法满足结合律，且有幺元$1$.
    \item 乘法对加法满足分配律：$\forall \,\, a, b, c \,\, \in \,\, \mathbb{Z}, \,\,\,\, (a + b) \times c = (a \times c) + (b \times c), \,\,\, c \times (a + b) = (c \times a) + (c \times b).$
\end{enumerate}

由于乘法满足交换律，$\mathbb{Z}$也是交换环。

为了方便起见，讨论环的时候，可以将定义中的第一种运算称为“加法”，
将第二种运算称为“乘法”。
但需要注意的是，这个语境下的“加法”和“乘法”，
和我们通常所说的概念不同。

比如，在环的定义中，我们并没有对乘法和加法的关系提出分配律以外的要求。
对于加法的零元$0$和乘法的幺元$1$，两者是否一定不等呢？

考虑环中任意元素$a$，

$$ 0 = 0\times a - 0\times a = (0 - 0)\times a = 0 \times a.$$

同样地，$0 = a\times 0$。如果$0 = 1$，那么对环中任意元素$a$，$a = 1\times a = 0\times a = 0$.
也就是说，环中所有元素都是$0$。我们称这种情况为“平凡的”。
有的书中会约定$0$和$1$不相等，直接避免这种情况出现。

另一个有趣的情况是\textbf{零因子}的存在性，即是否存在不为零的元素$a, b$使得$a\times b = 0$.
在$\mathbf{Z}$中，这个情况不存在，但在某些环中（比如合数$n$的同余环$\mathrm{Z}/n\mathbb{Z}$中），零因子是存在的。

如果环的乘法没有零因子，则称其为\textbf{无零因子环}。
无零因子的交换环称为\textbf{整环}。$\mathbf{Z}$就是整环。

环中的元素关于乘法不一定有逆元，特别是加法零元$0$。
这是因为，假若存在$a\in R$使得$0\times a = 1$，
那么从之前的结论可以得到：$0 = 0\times a = 1$，我们又回到了平凡的情况。
在这一点上，我们关于整数除法除数不能为零的约定在环的定义里得到了“解释”。
或者说，我们定义的“环”这种结构很好地“再现”或“包含”了整数的代数性质。

我们一般把元素$a$的加法逆元记为$-a$，乘法逆元记为$a^{-1}$.

如果除了加法零元以外，环中所有元素都有乘法逆元，这样的环就称作\textbf{除环}。
顾名思义，我们可以对环中的元素（加法零元除外）作除法了。
这样，我们就来到了最后一个主要定义对象：

\begin{df}\label{df:0_3}
    配备了运算$\circ$和$\star$，并满足以下条件的集合$K$，称为\textbf{域}：

    \begin{enumerate}
        \item $\circ$和$\star$都是$K \times K \rightarrow K$的二元封闭运算；
        \item $K$关于运算$\circ$构成交换群，其空元称为零元，记为$0_{K}$；
        \item $K$中除$0_{K}$以外的所有元素（记为$K^*$），关于运算$\star$构成交换群；
        \item 运算$\star$对运算$\circ$满足分配律：$\forall \,\, a, b, c \,\, \in \,\, K, \,\,\,\, (a \circ b) \star c = (a \star c) \circ (b \star c), \,\,\, c\star (a \circ b) = (c \star a) \circ (c \star b).$
    \end{enumerate}
    也称$K$关于$\circ, \star$构成域。
\end{df}

换句话说，一个环既是交换环又是除环，就称为域。
也可以说，域是所有非零元素都有乘法逆元的无零因子的交换环。

域有哪些例子呢？首先来看$\mathbb{Z}$。$\mathbb{Z}$关于加法和乘法是否是一个域？
答案是否定的，因为$\mathbb{Z}$中元素大多没有乘法逆元。比如不存在使得$2\times a = 1$的$a$。

但我们由此得到启发：可以查看$\mathbb{Q}$是否是域。这次答案是肯定的。验证如下：
\begin{enumerate}
    \item 加法和乘法都是$\mathbb{Q} \times \mathbb{Q} \rightarrow \mathbb{Q}$的二元封闭运算；
    \item $\mathbb{Q}$关于加法构成交换群，其空元为$0$；
    \item $\mathbb{Q}$中除$0$以外的所有元素，关于乘法构成交换群；
    \item 乘法对加法满足分配律：$\forall \,\, a, b, c \,\, \in \,\, \mathbb{Q}, \,\,\,\, (a + b) \times c = (a \times c) + (b \times c), \,\,\, c \times (a + b) = (c \times a) + (c \times b).$
\end{enumerate}

除了$\mathbb{Q}$，我们还可以验证，$\mathbb{R}$和$\mathbb{C}$也关于加法和乘法构成域。

群、环、域，是我们主要介绍的三种代数结构。除此之外，我们还可以定义各种各样的抽象代数结构。
这些代数结构都是从我们熟悉的某个代数概念出发，抽象出其结构特征，以进行研究。我们一般把对象集合和运算法则用括号合起来，表示某个抽象代数结构。
比如：$\left(\mathrm{Z}, +\right)$是一个群，而$\left(\mathrm{Q}, +, \times\right)$是一个域。

代数结构的子集，也可以拥有同样的代数结构。

\begin{df}\label{df:0_3_1}
    给定群$(G, \circ)$及$G$的子集$H$，如果$H$关于$\circ$构成群，就称$(H, \circ)$是$(G, \circ)$的\textbf{子群}，简记为$H\leqslant G$。
\end{df}

显然，每个群自己都是自己的子群，另外，群的空元构成的子集$(\{e\}, \circ)$也必然是它的子群。
这两个子群被称为任何群的平凡子群。

类似地，可以定义环和域的子结构：
\begin{df}\label{df:0_3_2}
    给定环$(R, \circ, \star)$及$R$的子集$S$，如果$S$关于$\circ, \star$构成环，就称$(S, \circ, \star)$是$(R, \circ, \star)$的\textbf{子环}，简记为$S \leqslant R$。
\end{df}

\begin{df}\label{df:0_3_3}
    给定域$(K, \circ, \star)$及$K$的子集$F$，如果$F$关于$\circ, \star$构成域，就称$(F, \circ, \star)$是$(K, \circ, \star)$的\textbf{子域}，简记为$F \leqslant K$。
\end{df}

对于其他的代数结构，我们也可以类似定义其子结构。子结构的运算法则默认为原结构的运算法则，关键在于：此运算法则在子结构中必须是封闭的。
如果某个子集的元素经过运算后变成不属于该子集的元素，那么这个子集便不是子结构。

对于环来说，有一类子环十分重要。
\begin{df}\label{df:0_3_4}
    设$L$是环$R$的子环。如果$L$满足以下条件：
    \begin{itemize}
        \item 左吸收性：$\forall \,\, a \in R,\,\, l \in L, \,\,\, al \in L.$
        \item 右吸收性：$\forall\,\, a \in R,\,\, l \in L, \,\,\, la \in L.$
    \end{itemize}
    就说$L$是$R$的\textbf{藏}，记为$L \blacktriangleleft R$。
\end{df}

给定环中元素$a$，考虑集合$\{rar' | r, r' \in R\}$，这个集合是$R$的藏，称为$a$生成的藏，也称为\textbf{单藏}，记为$\langle a \rangle$。
如果$R$是交换环，那么$\langle a \rangle = \{ar | r \in R\}= \{ra | r \in R\}$。

如果环的所有藏都是单藏，就说环是\textbf{单藏环}。

\begin{pp}\label{pp:0_3_5}
    给定域$\mathbb{F}$，整式环$\mathbb{F}[z]$是单藏环。
\end{pp}
\begin{proof} % [{\bfseries 证明\ref{pp:5_3_0}}]
    我们要证明$\mathbb{F}[z]$的藏都是单藏。设$I$为$\mathbb{F}[z]$的藏。
    如果$I = \{0\}$，这时$I = \langle 0 \rangle$是单藏。\\
    如果$I$有$0$以外的元素。这时，根据最小数原理，$I$去除$0$的集合中有次数最小的非零整式$p_0$。
    考虑$I$中另一个整式$p$，按照藏定义，$p$除以$p_0$得到的余式$r$也在$I$中，但$\deg r < \deg p_0$，所以$r = 0$。
    也就是说$p_0 | p$，这说明$I = \{p_0q | q \in \mathbb{F}[z]\}$是$p_0$生成的藏，即单藏。\\
\end{proof}

如果$L$是环$R$的藏，考虑所有形如$r + L = \{r + l|l\in L\}$构成的集合$R/_L = \{(r+L) | r \in R \}$，则：
$$ \forall a, b \,\, \in R, \,\,\, \{r + s | r \in (a + L), s \in (b + L)\} = (a + b + L).$$
我们把它用来定义$R/_L$元素的加法：
$$ \forall a, b \,\, \in R, \,\,\, (a + L) + (b + L)= (a + b + L).$$
此外：
$$ \forall a, b \,\, \in R, \,\,\, \{r \cdot s | r \in (a + L), s \in (b + L)\} = (a \cdot b + L).$$
我们把它用来定义$R/_L$元素的乘法：
$$ \forall a, b \,\, \in R, \,\,\, (a + L) \cdot (b + L)= (a \cdot b + L).$$
可以验证，这样配备了加法和乘法的$R/_L$是一个环，它的加法零元是$(0_R + L) = L$。我们把这个环称为$R$关于$L$的\textbf{商环}。

\section{代数结构与映射}
我们知道，映射是指把某些对象对应到另一些对象的对应关系。
一般来说，我们不关心这个对应关系和对象本身的性质有什么联系。
比如说，讨论映射$f: x \mapsto x^2$的时候，$x$的定义域可以是全体实数，正整数或其他集合，
但我们不关心映射$f$本身和作为代数结构的定义域的代数性质有什么关系。

了解了数系的概念，以及群、环、域等抽象代数结构后，我们感兴趣的是：
两个抽象代数结构是否“类似”或“其实是一样的”。
比如，我们定义这么一个群$(G,\circ)$：
\begin{enumerate}
    \item $G$中有一个不是空元的元素$a$；
    \item $a$和$a$的逆元$a^{-1}$不相同；
    \item $G$中所有元素都可以表示成若干个$a$的乘积，或若干个$a^{-1}$的乘积。
    \item $a, a\circ a, a\circ a \circ a, \ldots $各不相同。
\end{enumerate}

群$(G,\circ)$看起来和$\left(\mathrm{Z}, +\right)$很像（比如令$a = 1$）。两者是否实际上是同一个群呢？
我们首先来严格定义“实际上是同一个”这类概念。

\begin{df}\label{df:0_4}
    给定两个代数结构$A$和$B$。$f$是从$A$到$B$的映射。
    如果$f$在$B$中的映射结果关于$B$的运算法则构成的代数结构保持了$A$作为代数结构的性质，
    就称$f$是$A$到$B$的\textbf{同态映射}，简称\textbf{同态}，记为$A \simeq B$。
\end{df}

简单来说，存在同态映射，表示$A$的结构可以在$B$中“再现”：我们可以在$B$里面找到和$A$“类似”的部分：$f(A) = \{f(x) | x \in A\}$。
代数学中，我们把$A$中元素、$A$的子集乃至$A$自身经过同态映射得到的值或值的集合称为\textbf{像}。
比如$f(A) = \{f(a) | a \in A\}$是$A$的像。
设$Y$是$B$的子集，则集合$\{a |a \in A, \,\, f(a) \in Y\}$称为$Y$的\textbf{原像}。
加法零元的原像一般称为同态的\textbf{核}。

用常见的数系作为例子，考虑以下映射：
\begin{align*}
    f : \left(\mathbb{Z},+\right) &\rightarrow \left(\mathbb{Q},+\right)  \\
        x          &\mapsto 2 x 
\end{align*}

$f$是否能保存$\left(\mathbb{Z},+\right)$的结构呢？我们按照群的定义验证：

\begin{enumerate}
    \item 加法仍是二元封闭运算：$\forall x \in \mathbb{Z}, \,\,\, f(x) + f(y) = 2x + 2y = 2(x+y) = f(x+y)$；
    \item 加法仍满足结合律：$\forall a, b, c\,\, \in\,\, \mathbb{Z}, \,\,\,\, (f(a) +  f(b)) + f(c) = f(a) + (f(b)+ f(c)).$
    \item 加法仍有空元$f(0) = 0$：$\forall a \,\, \in \,\, \mathbb{Z}, \,\,\,\, f(a) + f(0) = 2a + 0 = f(a)  = 0 + 2a = f(0) + f(a). $
    \item 加法逆元仍然是加法逆元：$\forall a \,\, \in \,\, \mathbb{Z}, \,\,\, f(a) + f(-a) = 2a - 2a = f(b) + f(a) = 0 = f(0).$
\end{enumerate}

以上我们验证了什么呢？我们验证了：$\mathbb{Z}$的像$f(\mathbb{Z}) = \{f(x) | x \in \mathbb{Z}\}$关于$\mathbb{Q}$里的加法，也是一个群。
对群$\left(\mathbb{Z},+\right)$中元素进行群相关的运算，我们在$f(\mathbb{Z})$里能够“再现”。
我们把这种映射称为\textbf{群同态}。

下面给出群同态的严格定义：
\begin{df}\label{df:0_5}
    给定两个群$(G,\circ)$和$(H,\star)$。$f$是从$G$到$H$的映射。
    如果$\forall u, v \,\, \in \,\, G, \,\,\,\, f(u) \star f(v) = f(u\circ v)$，
    就称$f$是群$(G,\circ)$到$(H,\star)$的\textbf{群同态}。
\end{df}

注意到这个定义比我们之前验证的内容要少不少。但我们可以证明，从$ f(u) \star f(v) = f(u\circ v)$这一条性质就能推导出其余。
这说明，保持运算法则是同态的核心。

如果同态是单射，称之为\textbf{单同态}；
如果同态是满射，称之为\textbf{满同态}。
如果同态将代数结构映射到自身集合，则称为\textbf{自同态}。

同态是否就完整保存了代数结构呢？未必。举例来说，将所有元素映射到0的映射：$f(x) \equiv 0$也是同态，
但显然这样的同态保存不了什么代数结构。我们给同态映射再添加一个条件：一一对应。
这样，同态就能完整保存代数结构了。
我们把一一对应的同态映射称为\textbf{同构}，用$\cong$表记。
比如\textbf{群同构}就是一一对应的群同态。
顾名思义，同构的对象具有相同的代数结构，可以看作是“本质一样的对象”。

对于环和域，我们也有类似的同态定义：
\begin{df}\label{df:0_5_2}
    给定两个环$(R, +, \cdot )$和$(S, \star , \times)$。$f$是从$R$到$S$的映射。
    如果
    \begin{enumerate}
        \item $\forall u, v \,\, \in \,\, R, \,\,\,\, f(u) \star f(v) = f(u + v).$
        \item $\forall u, v \,\, \in \,\, R, \,\,\,\, f(u) \times f(v) = f(u \cdot v).$
        \item $f(1_R) = 1_S.$
    \end{enumerate}
    就称$f$是环$(R, +, \cdot )$到$(S, \star , \times)$的\textbf{环同态}。
\end{df}

\begin{df}\label{df:0_6}
    给定两个域$(K, +, \cdot )$和$(F, \star , \times)$。$f$是从$K$到$F$的映射。
    如果
    \begin{enumerate}
        \item $\forall u, v \,\, \in \,\, K, \,\,\,\, f(u) \star f(v) = f(u + v).$
        \item $\forall u, v \,\, \in \,\, K, \,\,\,\, f(u) \times f(v) = f(u \cdot v).$
    \end{enumerate}
    就称$f$是域$(K, +, \cdot )$到$(F, \star , \times)$的\textbf{域同态}。
\end{df}

注意，由于环中的非零元素关于乘法不构成群，所以需要附加乘法幺元映射到幺元的条件。
而域中非零元素关于乘法构成群，所以从前两个条件能推出乘法幺元映射到幺元的性质，不需要作为条件。
这也说明，如果两个域之间存在环同态，则这个环同态就是域同态。

环和域上的同态，如果是单射、满射或双射（一一对应的映射），称作环和域上的单同态、满同态和同构。
如果同构将代数结构映射到自身集合，则称为\textbf{自同构}。

\section{等价关系与等价类}
除了代数结构之间的关系，我们也关心代数结构内部元素之间的关系。等价关系是“相等”概念的推广，具体定义如下：

\begin{df}\label{df:0_6b}
    集合$X$中元素的\textbf{等价关系}为满足以下条件的二元关系$\sim$：
    \begin{enumerate}
        \item 满足排中律：$\forall u, v \,\, \in \,\, X$， $u \sim v$要么成立，要么不成立，两者必居其一。
        \item 自反性：$\forall u \,\, \in \,\, X$，$u \sim u$成立；
        \item 对称性：$\forall u, v \,\, \in \,\, X, \,\,\,\, u \sim v \Rightarrow v \sim u.$
        \item 传递性：$\forall u, v, w \,\, \in \,\, X, \,\,\,\, u \sim v, \,\,\, v \sim w \Rightarrow u \sim w.$
    \end{enumerate}
\end{df}

所有相互等价的元素的集合称为\textbf{等价类}，比如，所有和某个元素$x$等价的元素的集合，称作$x$（所在）的等价类，一般记作$\bar{x}$。

最简单的等价关系，是数量的相等关系。两个数要么相等，要么不相等。每个数都等于自己。容易验证，相等关系也满足对称性和传递性。

整数的同余关系，也是一种等价关系。比如：如果$a$除以$3$的余数与$b$除以$3$的余数相同，就称$a$和$b$模$3$同余，记作$a \equiv b \mod{3}$。
所有模$3$余$1$的整数：$-2, 1, 16, \cdots$的集合称为$1$的等价类，也称为同余类，记作$\bar{1}$。

显然，以上的等价类也是$4$或$22$的等价类，可以记作$\bar{4}$或$\bar{22}$，我们把其中一个元素拿出来，作为代表，称为等价类的代表元素。
一般来说，我们喜欢把所有模$3$的同余类记为：$\bar{0}, \bar{1}, \bar{2}$。我们可以验证，集合$\{\bar{0}, \bar{1}, \bar{2}\}$关于整数的加法构成一个群，
称为（整数）模$3$的加法群。

所有整数的二元有序对$(a, b)$（$b$不为$0$）中存在以下的等价关系：
$$(a, b) \sim (c, d) \mbox{ 当且仅当 } ad = bc.$$
这个关系可以用来定义有理数域。
加法和乘法可定义为：
$$ (a, b) + (c, d) = (ad + bc, bd), \quad (a, b) \cdot (c, d) = (ac, bd)$$
通过验证映射：
$$ f: \,\, (a, b) \mapsto \frac{a}{b}$$
我们可以证明，$\left(\{\bar{(a, b)} | \sim\}, +, \cdot\right)$和$\left(\mathbb{Q}, +, \cdot\right)$是域同构。

% 给定群$(G, \cdot)$和它的子群$H$，定义：
% $$ \forall a \in G, \quad a\cdot H = \{a \cdot h | h \in H\}. $$
% $a\cdot H$可以简记为$aH$，称为。

% 如果考虑另一种关系：
% $$ \forall a, b \in G, \quad a \sim b \Leftrightarrow \exists g \in G, \quad b = g\cdot a\cdot g^{-1}.$$
% 可以验证，关系$\sim$也是等价关系。

\section{整式}
整式似乎是一个简单的概念。我们从初中开始就认识了整式。

整式可以看作是一个或若干个\textbf{不定元}和某个域$K$中元素经过有限次加法和乘法得到的式子。$K$一般称为整式的系数域。不定元是整式的中心概念，是可以和系数域$K$中元素相加、相乘并满足交换律、结合律、分配律的元素。它可以看作一个$K$以外的元素，也可以看作一个等待被赋值的变量。不定元使得整式往往无法归并为一项，而表现为多项之和，因此整式也称为\textbf{多项式}。

比如，给定某个域$(K, +, \times )$，以$K$中元素为系数的一元整式，一般形式为\footnote{这里假定已经合并同类项。}：
$$ p(z) = a_0 + a_1 z + a_2 z^2 + \cdots + a_n z^n, \,\,\, a_0, a_1, \cdots ,a_n \in K, \,\,\, a_n \neq 0.$$
它由有限个单项式$a_0, a_1 z, \cdots, a_n z^n$相加而成。每个单项式是由不定元与$K$中元素相乘得到，乘法中不定元参加的次数，称为单项式的\textbf{次数}。
上面$p(z)$的表达式中，按单项式的次数从低到高排列。$a_0, a_1, \cdots , a_n$称为整式的系数，
次数最高且系数非零的单项式的次数$n$称为整式的\textbf{次数}，记作$\deg{p}$，$a_n$称为其首项系数。
所有以$K$中元素为系数的一元整式的集合记作$K[z]$。

如果首项系数$a_n = 1$，就称$p(z)$为\textbf{首一整式}。如果整式不含不定元，称之为\textbf{常整式}，简称\textbf{常式}。
所有常整式的集合就是$K$。非零常整式的次数是0，0作为常整式，约定其次数为$-\infty$。

$K[z]$中两个整式相加或相乘，结果仍然是整式。整式的和，次数可以比参与加法的整式次数更低，
比如$p(z) + (-p(z)) = 0$。但如果不同次数的整式相加，那么和的次数总等于参与加法的整式里次数最高的那个。
整式乘积的次数，等于参与乘法的整式的次数之和。

可以验证，$K[z]$关于整式的加法和乘法构成一个交换环，我们称其为域$K$上的（一元）\textbf{整式环}。
整式环的加法零元是域$K$的加法零元，乘法幺元也是$K$的乘法幺元。

两个整式相除，结果类似整数的带余除法。给定整式$p(z)$和$q(z)$，存在唯一的整式$s(z)$和$r(z)$，
使得
$$ p(z) = q(z) \cdot s(z) + r(z),  \quad \deg(r) < \deg(q). $$
整式$s$称为$p$除以$q$的\textbf{商}，$r$称为$p$除以$q$的\textbf{余式}。如果$r = 0$，就说$q$整除$p$，是$p$的\textbf{因式}。
$p$是$q$的\textbf{倍式}。

显然，说$q$整除$p$时，$\deg(q) \leqslant \deg(p)$。如果$\deg(q) < \deg(p)$，就称$q$是$p$的\textbf{真因式}。
如果整式$p$的真因式只有常整式，就称$p$为\textbf{既约整式}、\textbf{不可约整式}或\textbf{素整式}（与素数的定义类似），
否则称$p$为\textbf{可约整式}。

一次整式$az + b$显然是既约整式。一个整式是否可约，与系数域$\mathbb{K}$有关。
比如，$z^2 - 2$在$\mathbb{Q}$中不可约，在$\mathbb{R}$中就可以写成$(z - \sqrt{2})(z + \sqrt{2})$。

如果首项系数为$a$的整式$p$是既约整式（可约整式），那么$p$除以$a$之后仍然是既约整式（可约整式）。
所以，也可以假设既约整式总是首一的。

给定一个非零整式构成的集合$\mathcal{S}$，将里面每个整式的次数记下来，形成另一个集合：
$$\deg \mathcal{S} = \{n | \exists p \in \mathcal{S}, \deg p = n \} $$
根据自然数集的最小数原理\footnote{即“自然数的非空子集总有最小元素”。它和归纳公理等价，是自然数的根本属性之一。}，这个集合有最小元素。
记此最小元素为$m$，则按定义，存在$p_0 \in \mathcal{S}$，使得$\deg p_0 = m$。
我们也称$p_0$为集合$\mathcal{S}$的\textbf{极小元素}。

如果整式$q$同时是若干个整式$p_1, p_2, \cdots , p_k$的因式，就称它为$p_1, p_2, \cdots , p_k$的\textbf{公因式}。
设$p_1, p_2, \cdots , p_k$的所有公因式的集合为$S$，集合中每个整式，次数不超过$p_1, p_2, \cdots , p_k$的次数。
所以，其中必然有最大的次数$m$\footnote{这里我们再次用到了最小数原理。}。$m$对应的公因式中有唯一一个首一整式，
称为$p_1, p_2, \cdots , p_k$的\textbf{最大公因式}。
如果最大公因式为常整式1，就称$p_1, p_2, \cdots , p_k$\textbf{互素}。

关于整式，也有和整数一样的结论：
\begin{pp}{\bfseries 倍和析因定理}\label{pp:0_1}
    记$g(z)$为整式$p_1(z)$和$p_2(z)$的最大公因式。则必然存在整式$q_1(z)$和$q_2(z)$，使得
    $$ p_1(z) q_1(z) + p_2(z) q_2(z) = g(z).$$
    其中$q_1$的次数可以小于$\deg(p_2) - \deg(g)$，$q_2$的次数可以小于$\deg(p_1) - \deg(g)$.
\end{pp} 
它实际上是辗转相除法的反推。辗转相除法是求公因式的算法，首次出现于欧几里得的《几何原本》，
中国历史上则首见于秦九韶的《数书九章》，称为“大衍求一术”。因此本书将这个结论称为倍和析因定理\footnote{一般称为裴蜀引理或贝祖引理。}。
如果两个整式互素，那么最大公因式$g(z)=1$，等式变为：
$$ p_1(z) q_1(z) + p_2(z) q_2(z) = 1.$$

从倍和析因定理可以推出整式版的算术基本引理：
\begin{pp}{\bfseries 既约整式整除定理}\label{pp:0_2}
    如果既约整式$q(z)$整除若干个整式的乘积：$p_1(z)p_2(z)\cdots p_k(z)$，那么它整除其中某个整式$p_i(z)$.
\end{pp} 
与整数理论中一样，这个结论也可以推出整式版的算术基本定理：整式有唯一的因式分解。
\begin{pp}{\bfseries 整式唯一分解定理}\label{pp:0_3}
    每个非常整式$p(z)$都可以分解为若干个首一的既约整式与一个常数的乘积：
    $$ p(z) = a\cdot s_1^{k_1}(z)s_2^{k_2}(z)\cdots s_m^{k_m}(z). $$
    如果不考虑因式分解中既约整式的排列顺序，则分解的方式是唯一的。
\end{pp} 

整式$p$的唯一分解中，如果有一次因式$z - z_0$，就称$z_0$是整式$p$的\textbf{根}。
如果一次因式在分解中出现了$k$次，就说$z_0$是整式$p$的$k$\textbf{重根}。

\chapter{域扩张}

\section{域和域扩张}
\subsection{什么是域}
用通俗的说法，域就是“包含加减乘除的集合”。严格定义如下：
\begin{df}\label{df:9_1}
    配备了运算$\circ$和$\bullet$，并满足以下条件的集合$K$，称为\textbf{域}：
    \begin{enumerate}
        \item $\circ$和$\bullet$都是$K \times K \rightarrow K$的二元封闭运算；
        \item $K$关于运算$\circ$构成交换群，其空元称为零元，记为$0_{K}$；
        \item $K$中除$0_{K}$以外的所有元素（记为$K^\bullet$），关于运算$\bullet$构成交换群，其空元称为幺元，记为$1_{K}$；
        \item 运算$\bullet$对运算$\circ$满足分配律：$\forall \,\, a, b, c \,\, \in \,\, K, \,\,\,\, (a \circ b) \bullet c = (a \bullet c) \circ (b \bullet c), \,\,\, c\bullet (a \circ b) = (c \bullet a) \circ (c \bullet b).$
    \end{enumerate}
    严格来说，域包括集合$K$与其配备的运算，记为$(K, \circ, \star)$。
    配备运算已有说明的情况下，为了叙述方便，也可以只用集合$K$指代域。
\end{df}
不至于混淆的情况下，我们一般把$\circ$和$\bullet$分别用$+$和$\times$（或$\cdot$）代替，
称为“加法”和“乘法”\footnote{对应地，$K^\bullet$也改为$K^\times$。}。

容易验证，$(\mathbb{Q}, +, \times)$是一个域。历史上，域的概念也源自有理数。

另一个例子是$(\mathbb{Z}_p, +, \times)$。其中$\mathbb{Z}_p = \{\bar{0}, \bar{1}, \cdots , \overline{p-1}\}$是整数模素数$p$的同余类，
$+$和$\times$是整数加法和乘法在同余类中的自然延拓\footnote{“延拓”指把某个函数的定义域扩大，并且使得新的定义域上的性质与原定义域上的性质有某种一致性或接驳性。}。具体规则为：
\begin{align*}
    \overline{m} + \overline{n} &= \overline{m+n} = m+n \mod{p}  \\
    \overline{m} \times \overline{n} &= \overline{m\times n} = m\times n \mod{p} 
\end{align*}
可以验证，按照以上定义，$(\mathbb{Z}_p, +, \times)$是一个环；而当$p$是素数时，构成一个域。
比如，当$p = 3$的时候，$\mathbb{Z}_p = \{\bar{0}, \bar{1}, \bar{2}\}$。
运算规则为：
\begin{align*}
    \bar{0} + \bar{0} &= \bar{1} + \bar{2} = \bar{2} + \bar{1}= \bar{0}  \\
    \bar{0} + \bar{1} &= \bar{1} + \bar{0} = \bar{2} + \bar{2}= \bar{1}  \\
    \bar{0} + \bar{2} &= \bar{2} + \bar{0} = \bar{1} + \bar{1}= \bar{2}  \\
    \bar{1} \times \bar{1} &= \bar{2} \times \bar{2} = \bar{1}  \\
    \bar{1} \times \bar{2} &= \bar{2} \times \bar{1} = \bar{2}.  
\end{align*}

域的定义来自于对整数及有理数的代数特征的抽象。考虑映射：
\begin{align*}
    \varPsi : \mathbb{Z} &\rightarrow \mathbb{K}   \\
    n &\mapsto \left\{\begin{array}{rr}\underbrace{1_{K} + 1_{K} + \cdots + 1_{K}}_{n\mbox{次}} & n > 0 \\ 0_{K} & n = 0\\ -(\overbrace{1_{K} + 1_{K} + \cdots + 1_{K}}^{n\mbox{次}}) & n < 0  \end{array} \right.
\end{align*}
$\varPsi$俗称“累加映射”，想法来自于把$K$里的幺元不断累加。大致来说，就是把$n$映射到$n$个幺元的和。
% 我们也简记$\varPsi(n)$为$\varPsi(n) = n 1_{K}$。
可以验证，$\varPsi$是环同态。
比如说，$\varPsi(m+n)$就是$m+n$个$1_{K}$的和，所以自然等于$m$个$1_{K}$的和与$n$个$1_{K}$的和相加。
$\varPsi(m\cdot n)$就是$mn$个$1_{K}$的和，所以等于$m$个“$n$个$1_{K}$的和”再相加得到的和。

本质上，这个同态说明，任何域的加法和乘法等同于整数的加法和乘法。接下来我们会看到，从$\varPsi$的基本性质出发，域可以分为两大类。

考察$\varPsi$的核$\Ker \varPsi$（也就是$0_K$的原像），有两种可能的情况：要么$\Ker \varPsi = \{0\}$，也就是说$\varPsi$是单射；
要么$\Ker \varPsi$包含$0$以外的数。

首先讨论第一种情况。如果$\varPsi$是单射，那么$\varPsi$可以延拓为：
$$ \frac{m}{n} \mapsto \varPsi(m) \cdot \varPsi(n)^{-1} $$
可以验证，它是从域$\mathbb{Q}$到$\mathbb{K}$的环同态，因此也是域同态。
新的$\varPsi$仍然是单射。也就是说，$\mathbb{K}$里的某个部分（$\Img \varPsi$）和$\mathbb{Q}$同构。
换句话说，$\mathbb{K}$里“包含”了$\mathbb{Q}$。从同构意义上说，$\mathbb{Q}$是$\mathbb{K}$的子域。

对于第二种情况，设有非零整数$c \in \Ker \varPsi$，即$\varPsi(c) = 0$。显然，$\varPsi(-c) = -\varPsi(c) = 0$。
所以$\Ker \varPsi$包含正数和负数。我们只考虑其中的正数：$\Ker \varPsi \cap \mathbb{Z}^+$，
按照最小数原理，其中有一个最小的数，我们把它记作$p$。考虑$\Ker \varPsi$中的其它正数$c$，
$$\varPsi(c - p) = \varPsi(c) - \varPsi(p) = 0_{K} - 0_{K} = 0_{K},$$
所以，设$c$除以$p$的商和余数分别是$q$和$0\leqslant r<p$：$c = pq + r$，则$\varPsi(r) = 0_{K}$。
这说明$r = 0$。也就是说，$\Ker \varPsi$中的其它正数都是$p$的倍数。
而对于$\Ker \varPsi$中的负数，其相反数是$\Ker \varPsi$中正数，因此也是$p$的倍数。

于是，我们得到$\Ker \varPsi$的结构：$\Ker \varPsi = \{np | n\in \mathbb{Z}\} = \langle p \rangle$。
我们把$p$称为域$K$的\textbf{界点}\footnote{意为幺元累加的“临界点”，“过界”了就再次“从$0$开始”。也称为域的\textbf{特征}。}。

如果$p=1$，那么$\langle p \rangle = \mathbb{Z}$，$\varPsi$是零映射。
这是$K$只有零元素的平凡情况。如果$p>1$，那么$p$是素数。这是因为如果存在$1 < a \leqslant b < p$使得$p = ab$，那么
$$ \varPsi(a) \varPsi(b) = \varPsi(a) = 0_{K}$$
于是$\varPsi(a) = 0_{K}\cdot \varPsi(b)^{-1} = 0_{K}$，与$p$的定义矛盾。这说明$p>1$时只能是素数。

于是，从$\varPsi$可以自然导出另一个同态：
$$ \overline{\varPsi}: \overline{m} \mapsto \varPsi(m)  $$
$\overline{\varPsi}$是从域$\mathbb{Z}_p$到$\mathbb{K}$的环同态，因此也是域同态。
此外，$\overline{\varPsi}$是单射，也就是说，$\mathbb{K}$里的某个部分（$\Img \overline{\varPsi}$）和$\mathbb{Z}_p$同构。
换句话说，$\mathbb{K}$里“包含”了$\mathbb{Z}_p$。从同构意义上说，$\mathbb{Z}_p$是$\mathbb{K}$的子域。

\subsection{域的扩张}
以上讨论得到的两种域，一种“包含”了$\mathbb{Q}$，另一种“包含”了$\mathbb{Z}_p$。
我们通常也称“包含”$\mathbb{Q}$的域的界点是$0$。于是，所有非平凡的域可以分为两类，一类界点为$0$，一类界点为素数$p$。
对于“包含”$\mathbb{Q}$的域$\mathbb{K}$，我们称$\mathbb{K}$是$\mathbb{Q}$的\textbf{扩域}，
把构建$\mathbb{Q}$到$\mathbb{K}$的环单同态\footnote{指它既是环同态，也是单射。}的过程称为从$\mathbb{Q}$到$\mathbb{K}$的（域）\textbf{扩张}，记为$\mathbb{Q} \subseteq \mathbb{K}$。

更一般地，给定两个域$\mathbb{F}$和$\mathbb{K}$，如果能建立$\mathbb{F}$到$\mathbb{K}$的环单同态，那么我们也称$\mathbb{K}$是$\mathbb{F}$的\textbf{扩域}，
称构建这个环单同态的过程为从$\mathbb{F}$到$\mathbb{K}$的（域）\textbf{扩张}，
还可以称$\mathbb{F}$是$\mathbb{K}$的子域\footnote{实际上$\mathbb{F}$同构于$\mathbb{K}$的某个子域，为了叙述方便，我们姑且使用同一个称呼。}。

作为例子，我们考虑集合：$\mathbb{Q}(\sqrt{2}) = \{a + b\sqrt{2} | a, b \in \mathbb{Q}\}$。
不难验证，$(\mathbb{Q}(\sqrt{2}), +, \times)$是一个域，而且恒等映射
\begin{align*}
    \iota: \mathbb{Q} &\rightarrow \mathbb{Q}(\sqrt{2})  \\
    a &\mapsto a  
\end{align*}
就是$\mathbb{Q}$到$\mathbb{Q}(\sqrt{2})$的环单同态。
所以，$\mathbb{Q} \subseteq \mathbb{Q}(\sqrt{2})$是域扩张，$\mathbb{Q}(\sqrt{2})$是$\mathbb{Q}$的扩域。

另一个重要的例子是$\mathbb{R}$和$\mathbb{C}$。照葫芦画瓢，可以验证$\mathbb{R}$到$\mathbb{C}$的恒等映射也是环单同态，
所以$\mathbb{R} \subseteq \mathbb{C}$是域扩张，$\mathbb{C}$是$\mathbb{R}$的扩域。

\section{域扩张与平直代数}
\subsection{扩域是平直空间}
设有域扩张$\mathbb{F} \subseteq \mathbb{K}$，可以验证：
\begin{enumerate}
    \item $\forall \,\, a, b \,\, \in \mathbb{K}, \,\,\, a + b = b + a.$
    \item $\forall \,\, a, b, c \,\, \in \mathbb{K}, \,\,\, a + (b + c) = (a + b) + c.$
    \item $\forall \,\, a \,\, \in \mathbb{K}, \,\,\, 0_K + a = a + 0_K = a.$
    \item $\forall \,\, a \,\, \in \mathbb{K}, \,\,\, a + (-a) = (-a) + a = 0.$
    \item $\forall \,\, \lambda, \mu \,\, \in \mathbb{F}, \,\,\, a \in \mathbb{K}, \,\,\, (\lambda \mu) a = \lambda(\mu a). $
    \item $\forall \,\, a \,\, \in \mathbb{K}, \,\,\, 1_F \cdot a = a.$
    \item $\forall \,\, \lambda \,\, \in \mathbb{F}, \,\,\, a, b \in \mathbb{K}, \,\,\, \lambda(a + b) = \lambda a + \lambda b.$
    \item $\forall \,\, \lambda, \mu \,\, \in \mathbb{F}, \,\,\, a \in \mathbb{K}, \,\,\, (\lambda + \mu) a =  = \lambda a + \mu a.$
\end{enumerate}
这说明$\mathbb{K}$可以看作一个系数域为$\mathbb{F}$的平直空间。其中的向量加法和数乘分别是域中的加法和乘法。

从之前的例子来看，在域扩张$\mathbb{Q} \subseteq \mathbb{Q}(\sqrt{2})$里，把$(1, \sqrt{2})$作为基底，
可以将$\mathbb{Q}(\sqrt{2})$里的元素$a + b\sqrt{2}$看成坐标为$(a, b)$的向量。
$\mathbb{Q}(\sqrt{2})$可以看作系数域为$\mathbb{Q}$的$2$维平直空间，同构于$\mathbb{Q}^2$。
在域扩张$\mathbb{R} \subseteq \mathbb{C}$里，把$(1, \imath )$作为基底，
可以将复数$a + b\imath $看成坐标为$(a, b)$的向量。
$\mathbb{C}$可以看作系数域为$\mathbb{R}$的$2$维平直空间，同构于$\mathbb{R}^2$。

我们把$\mathbb{K}$作为$\mathbb{F}$上平直空间的维数称作扩张的\textbf{维数}，记为$[\mathbb{K} : \mathbb{F}]$。

如果$\mathbb{K} = \mathbb{F}$，那么扩张的维数是1，反之亦然。我们称其为平凡的扩张。

\subsection{塔定理}
我们知道，平直空间可以是有限维或无限维的。如果$\mathbb{F} \subseteq \mathbb{K}$中，
$\mathbb{K}$是$\mathbb{F}$上的有限维平直空间，就说这个扩张是\textbf{有限扩张}，否则就说它是\textbf{无限扩张}。

扩张的维数有以下基本性质：
\begin{pp}\textbf{塔定理 }\label{pp:9_0}
    给定域扩张$\mathbb{F} \subseteq \mathbb{K} \subseteq \mathbb{L}$\footnote{表示$\mathbb{K}$是$\mathbb{F}$的扩域，而$\mathbb{L}$又是$\mathbb{K}$的扩域。}，
    则$[\mathbb{K} : \mathbb{F}] = [\mathbb{K} : \mathbb{F}]\cdot [\mathbb{L} : \mathbb{K}]$。
\end{pp}
\begin{proof} % [{\bfseries 证明\ref{pp:9_0}}]
    如果$\mathbb{F} \subseteq \mathbb{K}$是无限扩张，那么由于$\mathbb{L}$包含$\mathbb{K}$，
    $\mathbb{F} \subseteq \mathbb{L}$必然是无限扩张。同理，如果$\mathbb{K} \subseteq \mathbb{L}$是无限扩张，
    $\mathbb{F} \subseteq \mathbb{L}$必然也是无限扩张。这两种情况下，
    $$[\mathbb{L} : \mathbb{F}] = \infty = [\mathbb{K} : \mathbb{F}]\cdot [\mathbb{L} : \mathbb{K}].$$
    如果$\mathbb{F} \subseteq \mathbb{K}$和$\mathbb{K} \subseteq \mathbb{L}$都是有限扩张，设$[\mathbb{K} : \mathbb{F}] = n$，
    $\mathcal{B}_1 = \{\mathbf{u}_1, \cdots, \mathbf{u}_n\}$是$\mathbb{K}$作为$\mathbb{F}$上平直空间的一组基；$[\mathbb{L} : \mathbb{K}] = m$，
    $\mathcal{B}_2 = \{\mathbf{v}_1, \cdots, \mathbf{v}_m\}$是$\mathbb{L}$作为$\mathbb{K}$上平直空间的一组基。下面证明：
    $\mathcal{B} = \{\mathbf{u}_i\mathbf{v}_j | 1 \leqslant i \leqslant n, \,\, 1 \leqslant j \leqslant m \}$是$\mathbb{L}$作为$\mathbb{F}$上平直空间的一组基。\\
    首先，$\mathcal{B}$中元素$\mathbf{u}_i\mathbf{v}_j$两两不同。反设有$\mathbf{u}_i\mathbf{v}_j = \mathbf{u}_k\mathbf{v}_l$，则
    $ \mathbf{u}_i\mathbf{v}_j - \mathbf{u}_k\mathbf{v}_l = 0.$ 
    而$\mathbf{u}_i$、$\mathbf{u}_k$是$\mathbb{K}$中元素，所以上式可看作$\mathcal{B}_2$中元素的$\mathbb{K}$-直组合\footnote{表示以$\mathbb{K}$中元素为系数的直组合。下同。}。
    因此从基的直无关性推出$\mathbf{u}_i = \mathbf{u}_k = 0$，矛盾\footnote{这里利用了零向量和系数域的零元重合的事实。下同。}。\\
    接着证明$\mathcal{B}$直无关。设有$\mathbb{F}$中系数$a_{i,j}$使得
    $$ \sum_{i,j} a_{i,j}\mathbf{u}_i\mathbf{v}_j = 0,$$
    那么，将其改写为
    $$ \sum_{j=1}^m \left(\sum_{i=1}^n a_{i,j}\mathbf{u}_i \right) \mathbf{v}_j = 0,$$
    由$\mathcal{B}_2$的直无关性可推出
    $$ \forall \,\, 1 \leqslant j \leqslant m , \,\,\, \sum_{i=1}^n a_{i,j}\mathbf{u}_i = 0.$$
    于是由$\mathcal{B}_1$的直无关性可推出
    $$ \forall \,\, 1 \leqslant i \leqslant n, \,\,1 \leqslant j \leqslant m , \,\,\,a_{i,j} = 0.$$
    再证明$\mathcal{B}$生成整个空间。$\mathbb{L}$中任意元素$\mathbf{x}$都可以写成：
    $$\mathbf{x} = \sum_{j=1}^m b_j \mathbf{v}_j,$$ 
    其中每个系数$b_j$都是$\mathbb{K}$中元素，所以可以写成：
    $$ b_j = \sum_{i=1}^n c_{i,j}\mathbf{u}_i , $$
    其中$c_{i,j}$是$\mathbb{F}$中元素。所以，$\mathbf{x}$可以写成
    $$ \mathbf{x} = \sum_{j=1}^m \left(\sum_{i=1}^n c_{i,j}\mathbf{u}_i \right) \mathbf{v}_j = \sum_{i,j} c_{i,j}\mathbf{u}_i\mathbf{v}_j.$$
    这说明$\mathcal{B}$生成$\mathbb{L}$。于是我们证明了$\mathcal{B}$是$\mathbb{L}$作为$\mathbb{F}$上平直空间的一组基。
    于是
    $$[\mathbb{L} : \mathbb{F}] = mn = [\mathbb{K} : \mathbb{F}]\cdot [\mathbb{L} : \mathbb{K}].$$
\end{proof}
 
比如，在有理数域$\mathbb{Q}$中添加$\sqrt{2}$，得到域$\mathbb{Q}(\sqrt{2})$。
在域$\mathbb{Q}(\sqrt{2})$中再添加添加$\sqrt{3}$，得到域$\mathbb{Q}(\sqrt{2})(\sqrt{3})$。
$\mathbb{Q}(\sqrt{2})$是$\mathbb{Q}$上的$2$维平直空间，
$(1, \sqrt{2})$是它的一组基。而$\mathbb{Q}(\sqrt{2})(\sqrt{3})$则是$\mathbb{Q}(\sqrt{2})$上的$2$维平直空间，
$(1, \sqrt{3})$是它的一组基。最后，$\mathbb{Q}(\sqrt{2})(\sqrt{3})$是$\mathbb{Q}$上的$4$维平直空间，
$(1, \sqrt{2}, \sqrt{3},\sqrt{6})$是它的一组基。   

\section{单扩张与代数扩张}
\subsection{最简单的扩张}
对于有限扩张的情形，如何确定相应平直空间的维数呢？我们从最简单的情形开始研究。

给定域$(\mathbb{F}, +, \times)$和一个不属于$\mathbb{F}$的元素$\alpha$，
考虑同时包含域$\mathbb{F}$和$\alpha$的域$\mathbb{K}$。首先，$\mathbb{K}$保留了集合$\mathbb{F}$及它配备的运算。
接下来，我们考虑以$\mathbb{F}$为系数域的整式集合$\mathbb{F}[z]$。
如果$\mathbb{F}$和$\alpha$同在域$\mathbb{K}$中，那么$\forall p \in \mathbb{F}[z]$，$p(\alpha) \in \mathbb{K}$。
所以集合$\mathbb{F}[\alpha] = \{p(\alpha) | p \in \mathbb{F}[z]\}$也在$\mathbb{K}$里面。考虑映射：
\begin{align*}
    \varPhi : \mathbb{F}[z] &\rightarrow \mathbb{K}  \\
    p &\mapsto p(\alpha)  
\end{align*}
$\varPhi$是一个环同态，因此是域同态。考虑$\varPhi$的核$\Ker\varPhi$，如果其中只有零，就称$\alpha$是$\mathbb{F}$上的\textbf{超越元}。
如果$\Ker\varPhi$中有非零整式$p$，说明$p(\alpha) = 0_K$。这时我们称$\alpha$是$\mathbb{F}$上的\textbf{代数元}\footnote{如果$\mathbb{F}$是有理数域，则称$\alpha$为\textbf{代数数}。“代数”表示它是整式方程的解。}。
\begin{df}\label{df:9_0}
    给定域扩张$\mathbb{F} \subseteq \mathbb{K}$，如果$\mathbb{K}$中每个元素都是$\mathbb{F}$上的代数元，就称这个扩张为\textbf{代数扩张}。
\end{df}
代数扩张对我们研究有限扩张很有用处。
\begin{pp}\label{pp:9_0_1}
    有限扩张必然是代数扩张。
\end{pp}
\begin{proof} % [{\bfseries 证明\ref{pp:9_0_1}}]
    设$\mathbb{F} \subseteq \mathbb{K}$的维数是$n$，对$\mathbb{K}$中任意元素$u$，
    考虑$1, u, \cdots, u^n$这$n+1$个元素，它们作为$\mathbb{F}$上平直空间里的向量，必然直相关，
    即存在$\mathbb{F}$中元素$c_0, c_1, \cdots , c_n$使得
    $$ c_0 + c_1 u + \cdots + c_n u^n = 0_K.$$
    这说明存在$\mathbb{F}[z]$中整式$p$，使得$p(u) = 0_K$，$u$是$\mathbb{F}$上代数元。\\
\end{proof}


考虑$\Ker\varPhi$的非零整式里面次数最小的首一整式$p_\alpha$，
我们称它为$\alpha$在$\mathbb{F}$上的\textbf{极小归零式}，记它的次数为$\deg p_\alpha = d_\alpha$，可以证明，
$\Ker\varPhi$中所有元素都是$p_\alpha$的倍式：$\Ker\varPhi = \{p_\alpha q | q \in \mathbb{F}[z]\} = \langle p_\alpha \rangle $。

另一方面，整式$p_\alpha$必然是域$\mathbb{F}$上的既约整式\footnote{即无法分解为$\mathbb{F}$上的非常整式的乘积，简称$\mathbb{F}$-既约整式。}。这是因为如果$p_\alpha$可以分解为非常整式的乘积$p_\alpha = p_1p_2$，
那么$0_K = p_\alpha(\alpha) = p_1(\alpha)p_2(\alpha)$。因此$p_1(\alpha) = 0_K p_2(\alpha)^{-1} = 0_K$，
$p_1\in \Ker\varPhi$。于是，按$p_\alpha$的定义，$\deg p_1 \geqslant d_\alpha$。但对称地，$\deg p_2 \geqslant d_\alpha$，
于是$d_\alpha = \deg p_\alpha = \deg p_1 + \deg p_2 \geqslant 2d_\alpha$，也就是说$d_\alpha \leqslant 0$，这与$p_\alpha$的定义矛盾。

既然$p_\alpha$是既约整式，对任意整式$p \in\mathbb{F}[z]$，$p$和$p_\alpha$的公因式是常整式。
根据倍和析因定理，存在整式$q, q_1$使得$pq + p_\alpha q_1 = 1_K$。因此，
$$1_K = p(\alpha)q(\alpha) + p_\alpha(\alpha)q_1(\alpha) = p(\alpha)q(\alpha).$$
也就是说$p(\alpha)$总有乘法逆元。这说明$\mathbb{F}[\alpha]$构成一个域。我们把这个域记作$\mathbb{F}(\alpha)$。

$\mathbb{F}(\alpha)$包含了$\mathbb{F}$和$\alpha$，而且是$\mathbb{K}$的子域。
换句话说，任何同时包含$\mathbb{F}$和$\alpha$的域，都包含$\mathbb{F}(\alpha)$。
所以$\mathbb{F}(\alpha)$是同时包含$\mathbb{F}$和$\alpha$的“最小”的域，
可以把它看作是向$\mathbb{F}$中“添加”$\alpha$得到的有限扩张。我们把这种扩张称为\text{单扩张}\footnote{本书只讨论代数扩张中的单扩张，以下提及“单扩张”都默认为代数扩张。如果$\alpha$是超越元，也可以定义超越扩张中的单扩张，本书对此不做讨论。}。

\subsection{用单扩张分解代数扩张}
单扩张可以看作是最简单的代数扩张。从单扩张$\mathbb{F} \subseteq \mathbb{F}(\alpha)$出发，
可以继续“添加”$\mathbb{F}$上的代数元$\beta$。得到的扩张是$\mathbb{F}(\alpha) \subseteq \mathbb{F}(\alpha)(\beta)$。
$\mathbb{F}(\alpha)(\beta)$中的元素是所有形同$p(\alpha, \beta)$的数，
其中$p\in \mathbb{F}[z_1, z_2]$是以$\mathbb{F}$中元素为系数的二元整式。

另一方面，如果我们先“添加”$\beta$，得到扩张$\mathbb{F} \subseteq \mathbb{F}(\beta)$，再“添加”$\alpha$得到
扩张$\mathbb{F}(\beta) \subseteq \mathbb{F}(\beta)(\alpha)$，
则域$\mathbb{F}(\beta)(\alpha)$中的元素是所有形同$p(\beta, \alpha)$的数。
而由于二元整式是交换环，
$$\mathbb{F}(\alpha)(\beta) = \mathbb{F}[\alpha, \beta] = \mathbb{F}[\beta, \alpha] = \mathbb{F}(\beta)(\alpha).$$
通过“添加”两个代数元得到的扩域，与“添加”的顺序无关。

类似地，由于多元整式是交换环，
通过若干次“添加”代数元得到的扩域，与“添加”的顺序无关。我们可以把从$\mathbb{F}$开始，通过$n$次单扩张得到的扩域写成
$$\mathbb{F}(\alpha_1, \alpha_2, \cdots, \alpha_n),$$
其中$\alpha_1, \alpha_2, \cdots, \alpha_n$是每次单扩张对应的代数元，排列顺序与先后顺序无关。

我们可以通过单扩张来理解有限扩张。给定非平凡的有限扩张$\mathbb{F} \subseteq \mathbb{K}$，由\ref{pp:9_0_1}知它是代数扩张。
在$\mathbb{K}$中找一个不属于$\mathbb{F}$的元素$x_1$，则可以构造单扩张$\mathbb{F} \subseteq \mathbb{F}(x_1) \subseteq \mathbb{K}$。
在$\mathbb{K}$中再找一个不属于$\mathbb{F}(x_1)$的元素$x_2$，
则可以继续构造单扩张$\mathbb{F} \subseteq \mathbb{F}(x_1) \subseteq \mathbb{F}(x_1, x_2) \subseteq \mathbb{K}$。
依此下去，我们可以构造由一层层域扩张构成的“塔”。

考虑这些扩张的维数。由\ref{pp:9_0}可知，
\begin{align*}
    [\mathbb{K} : \mathbb{F}] &= [\mathbb{F}(x_1) : \mathbb{F}]\cdot [\mathbb{K} : \mathbb{F}(x_1)]  \\
    &= [\mathbb{F}(x_1) : \mathbb{F}]\cdot [\mathbb{F}(x_1,x_2) : \mathbb{F}(x_1)] \cdot [\mathbb{K} : \mathbb{F}(x_1,x_2)]  \\
    &= \cdots  \\
    &= [\mathbb{F}(x_1) : \mathbb{F}]\cdot \prod_{i=1}^{k-1} [\mathbb{F}(x_1,\cdots,x_{i+1}) : \mathbb{F}(x_1,\cdots,x_i)] \cdot [\mathbb{K} : \mathbb{F}(x_1,\cdots,x_k)] 
\end{align*}
每次单扩张都是不平凡的，所以维数都大于1。如果我们把上式右侧的连乘积记为：
$$ g(k) = [\mathbb{F}(x_1) : \mathbb{F}] \prod_{i=1}^{k-1} [\mathbb{F}(x_1,\cdots,x_{i+1}) : \mathbb{F}(x_1,\cdots,x_i)]$$
那么，第$k$次扩张后，
$$g(k+1) = g(k) \cdot [\mathbb{F}(x_1,\cdots,x_{k+1}) : \mathbb{F}(x_1,\cdots,x_k)] > g(k).$$
这个连乘积的值随着每次单扩张严格增大。于是，经过有限次扩张，它的值会达到$\mathbb{F} \subseteq \mathbb{K}$的维数：$[\mathbb{K} : \mathbb{F}]$。
设$m$次扩张后我们达到了这个状态。这时候，$[\mathbb{K} : \mathbb{F}(x_1,\cdots , x_m)] = 1$。
这说明$\mathbb{K}$就是$\mathbb{F}(x_1,\cdots,x_m)$。

综上，我们得到这样的结论：
\begin{pp}\label{pp:9_2}
    所有有限扩张都可以分解为有限次单扩张。
\end{pp}

\section{单扩张的结构}
\subsection{用环同态理解单扩张}
对单扩张$\mathbb{F} \subseteq \mathbb{F}(\alpha)$来说，扩域$\mathbb{F}(\alpha)$有什么结构特征呢？

我们注意到，$\mathbb{F}(\alpha)$的元素和$p_\alpha$有关。
对任意$x \in \mathbb{F}(\alpha)$，我们总可以把它写成$x = p(\alpha)$，
其中整式$p \in\mathbb{F}[z]$。因此，它又可以写成：
$$ x = p(\alpha) = p_\alpha(\alpha)q(\alpha) + r(\alpha) = r(\alpha).$$
$r$就是$p$除以$p_\alpha$得到的余式。换句话说，如果两个整式相差$p_\alpha$的倍式，
那么在$\mathbb{F}(\alpha)$中对应的元素相同。这让我们想到数论中的同余类。
整式也有类似的现象。它们的本质可以用以下关于环的性质来概括。

首先，考虑环$R$的某个藏$L \blacktriangleleft R$，则$R$到商环$R /_L$的映射：
\begin{align*}
    \pi_L : R &\rightarrow R /_L  \\
    r &\mapsto (r + L) 
\end{align*}
是一个环同态。它的核$\Ker \pi_L$就是$(0 + L)$的原像，也就是藏$L$。

这个结论反过来，称为\textbf{环同态基本定理}：
\begin{pp}\textbf{环同态基本定理 }\label{pp:9_1}
    给定两个环$R_1$和$R_2$。如果$\pi: R_1 \rightarrow R_2$是环同态，那么$\pi$的核$\Ker \pi$是$R_1$的藏，
    $\pi$的像$\Img \pi = \pi (R_1)$是$R_2$的子环，且商环$R_1/_{\Ker \pi}$同构于$\pi$的像$\Img \pi$。
\end{pp}
\begin{proof} % [{\bfseries 证明\ref{pp:9_1}}]
    按照定义，容易验证$\pi$的核$\Ker \pi$是$R_1$的藏，$\pi$的像$\Img \pi = \pi (R_1)$是$R_2$的子环。
    下面证明商环$R_1/_{\Ker \pi}$同构于$\pi$的像$\Img \pi$。\\
    考虑映射：
    \begin{align*}
        \varPhi : R_1/_{\Ker \pi} &\rightarrow \Img \pi  \\
        (r + \Ker \pi) &\mapsto \pi(r)  
    \end{align*}
    由于$\Ker \pi$是$R_1$的藏，$\forall r, s \in R_1$，
    \begin{align*}
        (r + \Ker \pi) + (s + \Ker \pi) &= (r + s + \Ker \pi)  \\
        (r + \Ker \pi) \cdot (s + \Ker \pi) &= (r \cdot s + \Ker \pi)  
    \end{align*}
    所以
    \begin{align*}
        \varPhi((r + \Ker \pi) + (s + \Ker \pi)) &= \varPhi(r + s + \Ker \pi) = \pi(r + s)   \\
        &= \pi(r) + \pi (s) = \varPhi(r + \Ker \pi) + \varPhi(s + \Ker \pi)  \\
        \varPhi((r + \Ker \pi) \cdot (s + \Ker \pi)) &= \varPhi(r \cdot s + \Ker \pi) = \pi(r \cdot s)   \\
        &= \pi(r) \cdot \pi (s) = \varPhi(r + \Ker \pi) \cdot \varPhi(s + \Ker \pi).  
    \end{align*}
    这说明$\varPhi$是环同态。此外，如果$r \in R_1$使得$\varPhi(r + \Ker \pi) = 0_{R_2}$，那么$\pi(r) = 0_{R_2}$，
    即$r \in \Ker \pi$，也就是说$(r + \Ker \pi) = (0 + \Ker \pi) = \Ker \pi$。所以$\varPhi$是单射。\\
    另一方面，$\forall s \in \Img \pi$，可以找到$r \in R_1$使得$\pi(r) = s$，所以
    $$ \varPhi(r + \Ker \pi) = \pi(r) = s.$$
    这说明$\varPhi$是满射。于是$\varPhi$是环同构，$R_1/_{\Ker \pi}$同构于$\Img \pi$。
    
\end{proof}


无论是直同态基本定理，还是环同态基本定理，都有一个特点：如果某个映射保持了某种代数性质，
那么它总是把出发集合里的代数对象分成了一些等价类，把每个类映射到到达集合中。

\subsection{单扩张的元素是关于极小归零式的同余类}
回到$R = \mathbb{F}[z]$，$L = \langle p_\alpha \rangle$的设定。可以证明$L \blacktriangleleft R$。
这时，$\mathbb{F}(\alpha)$同构于商环$R/_L$。考虑映射：
\begin{align*}
    \pi : \mathbb{F}[z] &\rightarrow \mathbb{K}  \\
    p &\mapsto p(\alpha)  
\end{align*}
$\pi$是环同态，所以它的核$\Ker\pi = \langle p_\alpha \rangle$是$\mathbb{F}[z]$的藏，
且$\pi$的像：$\mathbb{F}(\alpha)$同构于商环$\mathbb{F}[z] /_{\langle p_\alpha \rangle}$。

商环$\mathbb{F}[z] /_{\langle p_\alpha \rangle}$的结构和数论中的同余类构成的环$\mathbb{Z}_p$是相似的。
因此，我们把其中的元素记为$\overline{r}$，其中$r$是次数小于$p_\alpha$的整式。比如，常整式$1$记为$\bar{1}$，
一次式$z$记为$\overline{z}$\footnote{注意这里不是指复数的共轭。}。$\mathbb{F}(\alpha)$中任意元素$x$，都可以表示为$\varPhi(\overline{r})$，其中$\overline{r}$是商环$\mathbb{F}[z] /_{\langle p_\alpha \rangle}$中的元素，
表示一个次数小于$d_\alpha$的整式$\overline{r} = c_0\bar{1} + c_1\overline{z} + \cdots + c_{d_{\alpha}-1} \overline{z^{d_{\alpha}-1}}$。
因此，$x$可以表达为：
\begin{align*}
    x &= \varPhi(\overline{r}) = c_0\varPhi(\bar{1}) + c_1\varPhi(\overline{z}) + \cdots + c_{d_{\alpha}-1} \varPhi(\overline{z^{d_{\alpha}-1}})  \\
    &= c_0 1_\mathbb{K} + c_1\alpha + \cdots + c_{d_{\alpha}-1} \alpha^{d_{\alpha}-1}.
\end{align*}
于是我们知道，$1_\mathbb{K}, \alpha, \alpha^2, \cdots , \alpha^{d_{\alpha}-1}$可以生成$\mathbb{F}(\alpha)$。
反过来，考虑$\varPhi$的逆映射：
\begin{align*}
    \varPsi = \varPhi^{-1} : \Img \pi = \mathbb{F}(\alpha) &\rightarrow \mathbb{F}[z] /_{\langle p_\alpha \rangle}  \\
     \pi(r)  &\mapsto \overline{r}  
\end{align*}
当$r = 0$时，我们得到$\varPsi(0_\mathbb{K}) = \overline{0}$；
当$r = 1$时，我们得到$\varPsi(1_\mathbb{K}) = \overline{1}$；
当$r = z$时，就得到$\varPsi(\alpha) = \overline{z}$。类似地，有
$$ \varPsi(\alpha^2) = \overline{z^2} , \,\,\, \varPsi(\alpha^3) = \overline{z^3} , \,\,\, \cdots $$
那么，如果$c_0 1_\mathbb{K} + c_1\alpha + \cdots + c_{d_{\alpha}-1} \alpha^{d_{\alpha}-1} = 0_\mathbb{K}$，就有
\begin{align*}
    \bar{0} &= \varPsi(0_\mathbb{K}) = \varPsi\left(c_0 1_\mathbb{K} + c_1\alpha + \cdots + c_{d_{\alpha}-1} \alpha^{d_{\alpha}-1}\right)  \\
    &= c_0\bar{1} + c_1\overline{z} + \cdots + c_{d_{\alpha}-1} \overline{z^{d_{\alpha}-1}}  \\
    &= \overline{c_0 + c_1z + \cdots + c_{d_{\alpha}-1} z^{d_{\alpha}-1}}. 
\end{align*}
因此，$c_0 = c_1 = \cdots = c_{d_{\alpha}-1} = 0_\mathbb{F}$。这说明$1_\mathbb{K}, \alpha, \alpha^2, \cdots , \alpha^{d_{\alpha}-1}$是$\mathbb{F}(\alpha)$作为$\mathbb{F}$上平直空间的一组基，
空间的维数$[\mathbb{F}(\alpha) : \mathbb{F}] = d_{\alpha}$。

把以上的结论总结一下：
\begin{pp}\label{pp:9_3}
    如果$\mathbb{F} \subseteq \mathbb{F}(\alpha)$是单扩张，$p_\alpha$是$\alpha$在$\mathbb{F}$上的极小归零式，
    次数为$d_{\alpha}$，那么
    \begin{enumerate}
        \item $p_\alpha$是$\mathbb{F}$-既约整式。
        \item $ \mathbb{F}(\alpha) \simeq \mathbb{F}[z] /_{\langle p_\alpha \rangle}.$
        \item $1, \alpha, \alpha^2, \cdots , \alpha^{d_{\alpha}-1}$是$\mathbb{F}(\alpha)$作为$\mathbb{F}$上平直空间的一组基。
        \item 扩张的维数$[\mathbb{F}(\alpha) : \mathbb{F}] = d_{\alpha}.$
    \end{enumerate}
\end{pp}

\section{分裂域}
有限扩张可以分解为有限个单扩张，每个单扩张对应一个既约整式。这让我们思考，如果给定一个域$\mathbb{F}$上的整式$p$，
是否能找到一个对应的有限扩张$\mathbb{F} \subseteq \mathbb{K}$，使$p$的根都在其中？

\begin{pp}\label{pp:9_4}
    给定域$\mathbb{F}$，每个$\mathbb{F}$-既约整式$p$都是某个单扩张$\mathbb{F} \subseteq \mathbb{F}(\alpha)$对应的极小归零式$p_\alpha$。
\end{pp}
\begin{proof} % [{\bfseries 证明\ref{pp:9_4}}]
    只需把对单扩张的讨论反过来论证一遍就可以了。\\
    $p$是既约整式，所以环$\mathbb{F}[z] /_{\langle p_\alpha \rangle}$是域：根据倍和析因定理，任何整式$q$与$p$的最大公因式都是$1$，
    所以可以找到$g$、$h$使得$pg + qh = 1$。，在$\mathbb{F}[z] /_{\langle p_\alpha \rangle}$中，上式变为：$\overline{q}\cdot\overline{h} = \overline{1}$。
    于是$\overline{h}$就是$\overline{q}$的乘法逆元。\\
    设$\alpha = \overline{x}$，则$p(\alpha) = \overline{p} = \overline{0}$，因此$\mathbb{F}[\alpha]$就是$\mathbb{F}[z] /_{\langle p_\alpha \rangle}$，
    于是$\mathbb{F}[\alpha]$可以写成$\mathbb{F}(\alpha)$。
    任何次数小于$p$的非常整式$q$都与$p$互素，所以都是$\mathbb{F}(\alpha)$中的可逆元，$q(\alpha) \neq \overline{0}$。
    所以单扩张$\mathbb{F} \subseteq \mathbb{F}(\alpha)$中，$p$是$\alpha$的极小归零式。\\
\end{proof}

对于可约的整式，我们可以定义分裂域的概念。
\begin{df}\label{df:9_2}
    给定域$\mathbb{F}$及$\mathbb{F}$上的非常整式$f$，如果$\mathbb{F}$的扩域$\mathbb{K}$使得$f$可以写为一次式的乘积：
    $$ f = a(z - \alpha_1)(z - \alpha_2)\cdots(z - \alpha_n), \,\, a \in \mathbb{F}, \,\, \alpha_1, \alpha_2, \cdots ,  \alpha_n \in \mathbb{K}, $$
    且$\mathbb{K} = \mathbb{F}(\alpha_1, \alpha_2, \cdots , \alpha_n)$，
    就称$\mathbb{K}$是$f$的\textbf{分裂域}，$f$在$\mathbb{K}$中分裂。
\end{df}

显然，域$\mathbb{F}$本身就是$\mathbb{F}$上的一次式的分裂域。

下面我们证明，每个非常整式都有分裂域。
\begin{tm}\label{tm:9_5}
    给定域$\mathbb{F}$及$\mathbb{F}$上任意$n$次整式$f$（$n > 0$），都有某个有限扩张$\mathbb{F} \subseteq \mathbb{K}$，
    使得$f$在$\mathbb{K}$中分裂，且扩张的维数$[\mathbb{K} : \mathbb{F}] \leqslant n!$。
\end{tm}
\begin{proof} % [{\bfseries 证明\ref{pp:9_5}}]
    对整式的次数$m$使用归纳法。\\
    $m = 1$时，整式是一次式，所以命题成立。\\
    假设$m < n$时命题成立。$m = n$时，设$p$是整式$f$因式分解中的某个$\mathbb{F}$-既约整式，
    由命题\ref{pp:9_4}可知，存在单扩张$\mathbb{F} \subseteq \mathbb{F}(\alpha)$，使得$p$是$\alpha$的极小归零式，$p(\alpha) = 0$。
    所以，$f$在$\mathbb{F}(\alpha)$中可以写为：
    $$ f(z) = (z - \alpha) q(z).$$
    其中整式$q$的系数在$\mathbb{F}(\alpha)$中，且次数为$n-1$。因此，存在有限扩张$\mathbb{F}(\alpha) \subseteq \mathbb{K}$，
    使得$q$在$\mathbb{K} = \mathbb{F}(\alpha)(\alpha_2, \cdots , \alpha_n)$中分裂。其中$\alpha_2, \cdots , \alpha_n$是$q$的根，
    且扩张的维数$[\mathbb{K} : \mathbb{F}(\alpha)] \leqslant (n - 1)!$。\\
    一方面，
    $$\mathbb{K} = \mathbb{F}(\alpha)(\alpha_2, \cdots , \alpha_n) = \mathbb{F}(\alpha, \alpha_2, \cdots , \alpha_n),$$
    另一方面，根据塔定理（\ref{pp:9_0}），
    $$[\mathbb{K} : \mathbb{F}] = [\mathbb{K} : \mathbb{F}(\alpha)]\cdot [\mathbb{F}(\alpha) : \mathbb{F}]. $$
    因此，扩张$\mathbb{F} \subseteq \mathbb{K}$是有限扩张，$p$在$\mathbb{K}$中分裂，
    且扩张的维数$[\mathbb{K} : \mathbb{F}] \leqslant n \cdot (n - 1)! = n!$。$m = n$时命题也成立。\\
    因此根据归纳原理，对任意正整数$m$，命题成立。\\
\end{proof}

代数基本定理说明，复数域$\mathbb{C}$包含了有理数域$\mathbb{Q}$、实数域$\mathbb{R}$和$\mathbb{C}$上任何整式的分裂域。
比如，考虑$\mathbb{Q}$上任一个整式$f$，$f$可以看作$\mathbb{C}$上的整式，因此可以将它写成：
$$ f = a(z - \alpha_1)(z - \alpha_2)\cdots(z - \alpha_n), \,\, a \in \mathbb{Q}, \,\, \alpha_1, \alpha_2, \cdots ,  \alpha_n \in \mathbb{C}, $$
而$\mathbb{K} = \mathbb{Q}(\alpha_1, \alpha_2, \cdots , \alpha_n)$就是$f$的分裂域。

\chapter{详细证明}

收录正文中证明过长、细节和技术过多的命题，给出详尽的证明。

“有限维平直空间”一章，“生成空间”一节的命题\ref{pp:2_3}：
\begin{pp}\label{c-10}
    有限维平直空间中，任何一组直无关的向量，其元素个数都不多于任意一组生成集的元素个数。
\end{pp}
\begin{proof} % [{\bfseries 证明\ref{pp:2_3}}]

    设$\mathfrak{F}_{\mathbf{u}} = (\mathbf{u}_1, \mathbf{u}_2, \cdots , \mathbf{u}_n )$是平直空间$\mathbb{V}$中一组直无关的向量，
    $\mathfrak{F}_{\mathbf{v}} = (\mathbf{v}_1, \mathbf{v}_2, \cdots , \mathbf{v}_m )$是$\mathbb{V}$的生成集。
    我们要证明$n \leq m$. \\    
    我们把$\mathfrak{F}_{\mathbf{v}}$中的元素逐一替换成$\mathfrak{F}_{\mathbf{u}}$中元素，一共替换$n$次。
    如果以上步骤能完成，说明$\mathfrak{F}_{\mathbf{v}}$的长度不少于$n$，就完成了证明。\\    
    第一步替换具体操作为：\\    
    将$\mathbf{u}_1$添加到$(\mathbf{v}_1, \mathbf{v}_2, \cdots , \mathbf{v}_m )$左侧，形成
    $$\mathfrak{F}_{\mathbf{v}}^{(1)} = (\mathbf{u}_1, \mathbf{v}_1, \mathbf{v}_2, \cdots , \mathbf{v}_m ).$$
    由于$\mathfrak{F}_{\mathbf{v}}$是$\mathbb{V}$的生成集，$\mathbf{u}_1$可以表示为$\mathbf{v}_1, \mathbf{v}_2, \cdots , \mathbf{v}_m$的直组合。
    即$\mathfrak{F}_{\mathbf{v}}^{(1)}$是直相关的。所以根据 \ref{pp:2_2}，我们可以去掉某个$\mathbf{v}$，使得其生成空间不变（仍然为$\mathbb{V}$）.
    这样$\mathfrak{F}_{\mathbf{v}}^{(1)}$变为一组元素个数为$m$的向量，并且仍然是$\mathbb{V}$的生成集。\\
    此后第$j$步，具体操作为：\\
    将$\mathbf{u}_j$添加到$\mathfrak{F}_{\mathbf{v}}^{(j-1)}$中$\mathbf{u}_{j-1}$的右侧。
    注意到$\mathfrak{F}_{\mathbf{v}}^{(j-1)}$是由左侧的$\mathbf{u}_1, \mathbf{u}_2, \cdots , \mathbf{u}_{j-1}$和右侧的$\mathfrak{F}_{\mathbf{v}}$中元素拼接而成，
    所以$\mathbf{u}_j$被放置在$\mathfrak{F}_{\mathbf{v}}^{(j-1)}$中原属于$\mathfrak{F}_{\mathbf{u}}$的元素和原属于$\mathfrak{F}_{\mathbf{v}}$的元素的分界上。
    放置之后形成的$\mathfrak{F}_{\mathbf{v}}^{(j)}$，“布局”和$\mathfrak{F}_{\mathbf{v}}^{(j-1)}$一样，仍然是由左侧原属于$\mathfrak{F}_{\mathbf{}}$的元素和右侧原属于$\mathfrak{F}_{\mathbf{v}}$的元素构成。
    这时，通过和第一步同样的逻辑，我们知道$\mathfrak{F}_{\mathbf{v}}^{(j)}$是直相关的，所以根据 \ref{pp:2_2}，可以从中找到某个元素：它属于它之前的元素的生成空间，
    并且去除这个元素后，$\mathfrak{F}_{\mathbf{v}}^{(j)}$生成的空间不变，仍然是$\mathbb{V}$。\\
    注意由于$\mathfrak{F}_{\mathbf{u}}$是直无关的，我们找到的这个元素肯定不会是左侧原属于$\mathfrak{F}_{\mathbf{u}}$的元素，因而只能是原属于$\mathfrak{F}_{\mathbf{v}}$的元素。
    这样，每一步操作，我们都加入一个原属于$\mathfrak{F}_{\mathbf{u}}$的元素，去除一个原属于$\mathfrak{F}_{\mathbf{v}}$的元素。
    元素个数不变，且最后形成的向量组仍然是$\mathbb{V}$的生成集。\\
    那么，是否会在某一步操作时，发现已经没有原属于$\mathfrak{F}_{\mathbf{v}}$的元素可以去除呢？\\
    如果这种情况发生，说明加入一个原属于$\mathfrak{F}_{\mathbf{u}}$的元素后，只有原属于$\mathfrak{F}_{\mathbf{u}}$的元素了。
    这时候向量组一方面仍然是直相关的，但同时元素全是$\mathfrak{F}_{\mathbf{u}}$中元素。说明$\mathfrak{F}_{\mathbf{u}}$也是直相关的。矛盾！\\
    所以每一步操作都可以完整进行，直到$\mathfrak{F}_{\mathbf{u}}$中所有元素放入完毕。
    这样，我们就证明了$\mathfrak{F}_{\mathbf{v}}$的长度不少于$n$，即$n \leq m$.\\
\end{proof}

“本征值和自同态约简”一章，“幂零自同态与标准形式”之“标准形式”的命题\ref{pp:6_12_4}:
\begin{pp}\label{c-20}
    给定空间$\mathbb{V}$上幂零自同态$N$，空间可以表示为：
    $$ \mathbb{V} = C_N(\mathbf{u}_1) \oplus C_N(\mathbf{u}_2) \oplus \cdots \oplus C_N(\mathbf{u}_p),$$
    其中$p$为正整数，$\mathbf{u}_1, \mathbf{u}_2, \cdots , \mathbf{u}_p$为$N$决定的向量。
\end{pp}
\begin{proof} % [{\bfseries 证明\ref{pp:6_12_4}}]
我们对空间的维数$n$作归纳法。$n=1$时，$N$只可能为零映射，因此结论显然成立：
任选非零向量$\mathbf{u}$，$C_N(\mathbf{u}) = \langle \mathbf{u} \rangle = \mathbb{V}$。\\
假设结论对维数小于$n$的空间上的幂零自同态成立。要证明结论对$n$维空间成立，对任意幂零自同态$N$，只需找到一个轨迹，
让整个空间可以表达成其循环子空间$C_N(\mathbf{u})$和另一个$N$-不变子空间$\mathbb{W}$的直和：
$$ \mathbb{V} = C_N(\mathbf{u}) \oplus \mathbb{W}.$$
这时候，$N$限制在$\mathbb{W}$上的部分也是幂零自同态。$\mathbb{W}$的维数小于$n$，所以根据归纳假设，
$\mathbb{W}$可以写成$N$-循环子空间的直和，因而$\mathbb{V}$可以表示成$N$-循环子空间的直和。\\
于是，问题转为证明空间总可以表达成$N$-循环子空间$C_N(\mathbf{u})$和另一个$N$-不变子空间的直和。\\
我们对幂零自同态的阶数$d$使用归纳法。证明命题：$d$阶幂零自同态可以表示为某个循环子空间和某个$N$-不变子空间的直和。\\
$d=1$时，$N$为零映射，每个非零元素的轨迹就是自己，因此结论显然成立：
$$\mathbb{V} = C(\mathbf{e}_1) \oplus \cdots \oplus C(\mathbf{e}_n)$$
其中$\mathbf{e}_1, \cdots , \mathbf{e}_n$是空间的一组基。\\
假设结论对小于$d$阶的幂零自同态成立。给定$d$阶幂零自同态$N$，则总有向量轨迹长度为$d$。
这是因为如果所有向量轨迹长度不超过$d-1$，那么$\forall \mathbf{v}, N^{d-1}\mathbf{v} = \mathbf{0}$，也就是说$N$的阶至多为$d-1$。\\
设这个轨迹为：
$$ \mathbf{u}, N\mathbf{u}, N^2 \mathbf{u}, \cdots , N^{d - 1} \mathbf{u}$$
记$\mathbb{A} = C_N(\mathbf{u})$，$A$是$N$-不变子空间。考虑它在$N$下的像空间$\mathbb{A}_1 = N(\mathbb{A}) \subseteq \Img N$。
$N$限制在$\Img N$上的部分是$d-1$阶的幂零自同态，所以根据归纳假设$\Img N$可以分解为
$$ \Img N = \mathbb{A}_1 \oplus \mathbb{B}_1 $$
其中的$\mathbb{A}_1$在$\restr{N}{\Img N}$上是$\Img N$的循环子空间，因为按定义$\mathbf{u}\notin \Img N$。\\
设$\mathbb{B} = N^{-1}(\mathbb{B}_1)$为$\mathbb{A}_1$补空间的原像。我们证明以下三个性质：\\
1. $\mathbb{A} + \mathbb{B} = \mathbb{V}.$\\
任取$\mathbf{v} \in \mathbb{V}$，由前可知唯一存在$\mathbf{a}_1\in\mathbb{A}_1$，$\mathbf{b}_1\in\mathbb{B}_1$，
使得$N\mathbf{v} = \mathbf{a}_1 + \mathbf{b}_1$，
且存在$\mathbf{a}\in\mathbb{A}$使得$N\mathbf{a} = \mathbf{a}_1$。而
$N(\mathbf{v} - \mathbf{a}) = \mathbf{b}_1\in\mathbb{B}_1$，所以$\mathbf{v} - \mathbf{a} \in \mathbb{B}$。\\
2. $\mathbb{A} \cap \mathbb{B}_1 = \{\mathbf{0}\}.$\\
设有$\mathbf{x} \in \mathbb{A} \cap \mathbb{B}_1$。按定义，$\mathbb{B}_1 \subseteq \Img N$，所以$\mathbf{x} \in \Img N$。
设$\mathbf{x} = \sum_{i=0}^{d-1} \lambda_i N^i\mathbf{u}$，则$\mathbf{x} = \sum_{i=1}^{d-1} \lambda_{i-1} N^i\mathbf{u}$。
所以$\mathbf{x} \in N(\mathbb{A}) = \mathbb{A}_1$。因此$\mathbf{x} \in \mathbb{A}_1 \cap \mathbb{B}_1 = \{\mathbf{0}\}.$\\
3. $\mathbb{B}_1 \subseteq \mathbb{B}.$\\
按定义，$\mathbb{B}_1$是$N$-不变子空间，因此$\forall \mathbf{x}\in\mathbb{B}_1$，$N\mathbf{x} \in \mathbb{B}_1$，所以$\mathbf{x}\in\mathbb{B}.$\\
以上三个性质说明，存在子空间$\mathbb{B}_1 \subseteq \mathbb{W} \subseteq \mathbb{B}$，使得$\mathbb{V} = \mathbb{A} \oplus \mathbb{W}$。
事实上，如果$\mathbb{V} = \mathbb{A}\oplus\mathbb{B}_1$，那么直接选择$\mathbb{W} = \mathbb{B}_1$即可。
否则将$\mathbb{B}_1$的一组基$\mathcal{B}$扩展为$\mathbb{B}$的基。记扩展过程中添加的基向量为$\mathbf{w}_1, \cdots, \mathbf{w}_s$，
则我们将它们一个个添加到$\mathbb{A}\oplus\mathbb{B}_1$中生成新空间，得到空间套$\mathbb{A}\oplus\mathbb{B}_1 \subseteq \cdots \subseteq \mathbb{V}$。
在此过程中，如果添加某个基向量使得空间维数增加，就将它和$\mathcal{B}$放在一起，直到空间维数等于$n$为止。
这样把$\mathcal{B}$扩展为$\mathcal{B}'$。由构造方法，可知$\mathcal{B}'$直无关，且生成的子空间$\mathbb{W}$满足
$\mathbb{B}_1 \subseteq \mathbb{W} \subseteq \mathbb{B}$。
再之，设添加的基向量生成的空间为$\mathbb{U}$，则$\mathbb{W} = \mathbb{B}_1 \oplus \mathbb{U}$，
所以
$$\mathbb{V} = (\mathbb{A}\oplus\mathbb{B}_1)\oplus \mathbb{U} = \mathbb{A}\oplus(\mathbb{B}_1\oplus \mathbb{U}) = \mathbb{A} \oplus \mathbb{W}.$$
最后，$N(\mathbb{W}) \subseteq N(\mathbb{B}) = \mathbb{B}_1 \subseteq \mathbb{W}$。
于是，我们找到了$N$-不变子空间$\mathbb{W}$，使得$\mathbb{V} = \mathbb{A} \oplus \mathbb{W}$。这就完成了归纳证明。

\end{proof}

“自同态的效”一章，“效的计算”之“附阵和一次方程组的解”的命题\ref{pp:10_50}：
\begin{pp}\label{c-25}
    设$A,B$为$n\times n$矩阵，则其附阵满足以下性质：
    \begin{enumerate}
    \item $\co(I_n) = I_n.$
    \item $\co(A^{\tr}) = \co(A)^{\tr}.$
    \item $\co(kA) = k^{n-1}\co(A).$
    \item $\co(AB) = \co(B)\co(A). $
    \item 如果$A$可逆，那么$\co(A^{-1}) = \co(A)^{-1} = \frac{A}{\det(A)}.$
    \item $\det(\co(A)) = \det(A)^{n-1}.$
\end{enumerate}
\end{pp}
\begin{proof}
    1). $I_n$的逆矩阵就是自己，效为$1$，所以$\co(I_n) = I_n$。

    2). 根据$\co(A)A = A\co(A)$，把$A$转置后按行展开公式里的余子阵，就是把$A$按列展开后的公式里的余子阵，所以$A^{\tr}$的附阵自然就是$A$的附阵的转置。

    3). 余子阵是$n-1$阶矩阵，所以$\co(kA) = k^{n-1}\co(A)$。

    4). $A,B$为可逆矩阵时，
    \begin{align*}
        \co(B)\co(A)AB &= \co(B)\det(A)I_nB = \det(B)\det(A)I_n \\
        &= \det(AB)I_n = \co(AB)AB.
    \end{align*}
    于是两边右乘$B^{-1}A^{-1}$就得到$\co(AB) = \co(B)\co(A)$。\\
    对于一般矩阵$A,B$，考虑矩阵$A+tI_n, B+tI_n$，它们的效是关于$t$的$n$次整式，因此合共至多有$2n$个根。当$t$不是这$2n$个根时，$A+tI_n, B+tI_n$都可逆，于是得到
    $$\co((A+tI_n)(B+tI_n)) = \co(B+tI_n)\co(A+tI_n).$$
    把这个等式看作关于$t$的多项式等式，展开整理移项可得$P(t) = 0$，其中$P$是次数不超过$2n$的整式。然而它对不属于之前提到的$2n$个根的任意$t$都成立。这说明只要系数域是实数或复数这样特征为零的数域，就能找到无限个$t$使得$P(t)=0$，从而得到$P$是零多项式，于是$P(0) = 0$，也就是$\co(AB) = \co(B)\co(A)$\footnote{实际上可以在扩域$\mathbb{F}(t)$里讨论，直接得出$\co((A+tI_n)(B+tI_n)) = \co(B+tI_n)\co(A+tI_n)$是矩阵多项式的恒等式，于是$t=0$时成立。}。

    4). 首先$\co(A)\cdot\frac{A}{\det(A)} = I_n$，所以$\co(A)^{-1} = \frac{A}{\det(A)}$。其次$\co(A^{-1})A^{-1} = \det(A^{-1})I_n = \frac{1}{\det(A)}I_n$。所以$\co(A^{-1}) = \frac{A}{\det(A)}$。

    5). $\det(\co(A))\det(A) = \det(\co(A)A) = \det(\det(A)I_n) = \det(A)^{n}$。如果$A$可逆，那么$\det(\co(A)) = \det(A)^{n-1}$。如果$A$不可逆，那么$\co(A)$必然也不可逆。反设$\co(A)$可逆，则在$\co(A)A = \det(A)I_n$两边同时左乘$\co(A)^{-1}$，得到$A = \det(A)\co(A) = 0$（零矩阵，因为$A$不可逆）。要么$A$并非零矩阵从而矛盾，要么$A$是零矩阵从而$\co(A)$也是零矩阵，矛盾。于是证得$\co(A)$不可逆，于是$\det(\co(A)) = 0 = \det(A)^{n-1}$。因此在所有情况下都有$\det(\co(A)) = \det(A)^{n-1}.$
\end{proof}

“二次型”一章，“同时合同对角化”一节的定理\ref{tm:8-2-2}：
\begin{tm}\label{c-30}
    设$A,B$为实对称矩阵，其中至少一个可逆。那么以下两命题等价：
    \begin{enumerate}
        \item $A,B$可同时合同对角化。
        \item 有正定对称矩阵$H$使得$AHB=BHA$。
    \end{enumerate}
\end{tm}

\begin{proof}
    首先证明：如果$A,B$可同时合同对角化，那么有正定对称矩阵$H$使得$AHB=BHA$。\\
    按定义有可逆矩阵$P$使得$D_A = P^{\tr}AP,\; D_B = P^{\tr}BP$都是对角矩阵。则
    $$A = \left(P^{\tr}\right)^{-1}D_AP^{-1}, \quad B = \left(P^{\tr}\right)^{-1}D_BP^{-1}.$$
    设$H = PP^{\tr}$，则有：
    \begin{align*}
        AHB &= \left(P^{\tr}\right)^{-1}D_AP^{-1}PP^{\tr}\left(P^{\tr}\right)^{-1}D_BP^{-1} \\
        &= \left(P^{\tr}\right)^{-1}D_AD_BP^{-1} = \left(P^{\tr}\right)^{-1}D_BD_AP^{-1} \\
        &= \left(P^{\tr}\right)^{-1}D_BP^{-1}PP^{\tr}\left(P^{\tr}\right)^{-1}D_AP^{-1} \\
        &= BHA.
    \end{align*}
    再证明$H$是正定对称矩阵。首先，$H^{\tr} = \left(PP^{\tr}\right)^{\tr} = \left(P^{\tr}\right)^{\tr}P^{\tr} = PP^{\tr} = H$，所以$H$是对称矩阵。其次，$\forall \;\mathbf{x}\in\mathbb{R}^n$，
    \begin{align*}
        [\mathbf{x}]^{\tr}H[\mathbf{x}] &= [\mathbf{x}]^{\tr}PP^{\tr}[\mathbf{x}] = \left(P^{\tr}[\mathbf{x}]\right)^{\tr}\left(P^{\tr}[\mathbf{x}]\right) = \|P^{\tr}[\mathbf{x}]\|^2 \geqslant 0.
    \end{align*} 
    其中$\|\cdot \|$是点积。最后，$P^{\tr}$可逆，所以上式取等号当且仅当$\mathbf{x} = \mathbf{0}$。这就说明$H$是正定对称矩阵。
    
    再证明另一个方向：如果有正定对称矩阵$H$使得$AHB=BHA$，则$A,B$可同时合同对角化。\\
    由于$H$是正定对称矩阵，它可以确定一个内积：
    $$\nji{\mathbf{x}}{\mathbf{y}}_H = [\mathbf{x}]^{\tr}H[\mathbf{y}].$$
    按条件$A,B$中至少有一个可逆，不妨设$B$可逆。考察$AB^{-1}$，下面证明它对应一个关于$\nji{\cdot}{\cdot}_H$的内积空间中的自伴自同态$T$：
    $$ \forall \mathbf{x}, \mathbf{y} \in \mathbb{R}^n,\quad \nji{\mathbf{x}}{T\mathbf{y}}_H = \nji{T\mathbf{x}}{\mathbf{y}}_H. $$
    具体来说，就是要证明：
    $$ [\mathbf{x}]^{\tr}H\left(AB^{-1}[\mathbf{y}]\right) = \left(AB^{-1}[\mathbf{x}]\right)^{\tr}H[\mathbf{y}]. $$
    注意到$A,B$都是对称矩阵，从而$B^{-1}$也是对称矩阵。于是：
    \begin{align*}
        \left(AB^{-1}[\mathbf{x}]\right)^{\tr}H[\mathbf{y}] = [\mathbf{x}]^{\tr}\left(B^{-1}\right)^{\tr}A^{\tr}H[\mathbf{y}] = [\mathbf{x}]^{\tr}B^{-1}AH[\mathbf{y}].
    \end{align*}
    然而，只要在$AHB=BHA$的左右两旁同时乘上$B^{-1}$，就得到
    $$ B^{-1}AH = B^{-1}AHBB^{-1} = B^{-1}BHAB^{-1} =  HAB^{-1} .$$
    于是：
    $$ [\mathbf{x}]^{\tr}HAB^{-1}[\mathbf{y}] = [\mathbf{x}]^{\tr}B^{-1}AH[\mathbf{y}] = \left(AB^{-1}[\mathbf{x}]\right)^{\tr}H[\mathbf{y}]. $$
    这就证明了，$T$关于$\nji{\cdot}{\cdot}_H$是自伴自同态。于是，根据谱定理，$AB^{-1}$可以相似对角化。同理，$B^{-1}A = B^{-1}(AB^{-1})B$相似于$AB^{-1}$，所以也可以相似对角化。
    
    选择$R = B^{-1}A$进行讨论。按定义$B R = A$，而$A = A^{\tr}$，所以
    $$
    BR = A = A^{\tr} = R^{\tr} B^{\tr} = R^{\tr} B.
    $$
    于是，$\forall \; \mathbf{x}, \mathbf{y} \in \mathbf{R}^n$，
    $$
    \left(R[\mathbf{x}]\right)^{\tr}B[\mathbf{y}] =  [\mathbf{x}]^{\tr}R^{\tr}B[\mathbf{y}] = [\mathbf{x}]^{\tr}BR[\mathbf{y}] 
    $$
    记$R$对应的自伴自同态为$T_R$。对于$T_R$的本征值$\lambda, \mu$及其对应的本征向量$\mathbf{u}, \mathbf{v}$，我们有：
    \begin{align*}
        \lambda [\mathbf{x}]^{\tr}B[\mathbf{y}] & = [R\mathbf{x}]^{\tr}B[\mathbf{y}] = [\mathbf{x}]^{\tr}BR[\mathbf{y}] = \mu[\mathbf{x}]^{\tr}B[\mathbf{y}], \\
        (\lambda - \mu) [\mathbf{x}]^{\tr}B[\mathbf{y}] &= 0.
    \end{align*}
    这说明只要$\lambda \neq \mu$，就有$[\mathbf{x}]^{\tr}B[\mathbf{y}] = 0$，我们称其为“关于$B$正交”。

    记$B$对应的双直型为$\phi_B$。按照定义，$\mathbf{x},\, \mathbf{y}$关于$B$正交，就是说$\phi_B(\mathbf{u}, \mathbf{v}) = 0$。考察$\phi_B$限制在$T_R$的本征空间$E_R(\lambda)$中的部分：
    \begin{align*}
        \phi_B^{\lambda}: \; E_R(\lambda)\times E_R(\lambda)\; &\longrightarrow \;\, \mathbb{R} \\
        (\mathbf{x},\,\; \mathbf{y}) \quad\quad &\mapsto \quad [\mathbf{x}]^{\tr}B[\mathbf{y}]
    \end{align*}
    它在$E_R(\lambda)$中仍然是对称双直型。此外，$\phi_B^{\lambda}$还是\textbf{非退化}的：$\forall \mathbf{u}\in E_R(\lambda)$，如果$\forall \mathbf{v} \in E_R(\lambda)$，都有$\phi_B^{\lambda}(\mathbf{u}, \mathbf{v}) = 0$，那么$\mathbf{u} = 0$。这是因为$R$可对角化，拥有由本征向量构成的基底。前面已经证明了，$\mathbf{u}$作为$E_R(\lambda)$中向量和其他本征空间里的向量关于$B$正交。如果$\mathbf{u}$还跟$E_R(\lambda)$中所有向量正交，就是和整个空间关于$B$正交：
    $$ \forall \; \mathbf{v}, \quad [\mathbf{u}]^{\tr}B[\mathbf{v}] = 0. $$
    我们可以取$\mathbf{v}$为$B^{-1}$的各列，拼合起来就有$0 = [\mathbf{u}]^{\tr}BB^{-1} = [\mathbf{u}]^{\tr}$。

    这又进一步说明，在$E_R(\lambda)$中总能找到非零向量$\mathbf{u}$，使得$[\mathbf{u}]^{\tr}B[\mathbf{u}] \neq 0$。反设找不到这样的向量，即$E_R(\lambda)$中任何非零向量$\mathbf{u}$都使得$[\mathbf{u}]^{\tr}B[\mathbf{u}] = 0$。那么根据恒等式：
    $$
    \phi_B^{\lambda}(\mathbf{u}, \mathbf{v}) = \frac{1}{2}\left(\phi_B^{\lambda}(\mathbf{u}+\mathbf{v},\,\mathbf{u}+\mathbf{v}) - \phi_B^{\lambda}(\mathbf{u},\mathbf{u}) - \phi_B^{\lambda}(\mathbf{v},\mathbf{v})\right), 
    $$
    $\forall\;\mathbf{u}, \mathbf{v}\; \in E_R(\lambda)$，$[\mathbf{u}]^{\tr}B[\mathbf{v}] = 0$。于是必然有非零向量$\mathbf{u}$和$E_R(\lambda)$中所有向量关于$B$正交，和非退化的性质矛盾。

    于是我们可以这样来构建$E_R(\lambda)$中一组关于$B$正交的自由向量，它们构成了$E_R(\lambda)$的基底。首先，可以找到$\mathbf{u}_1 \in E_R(\lambda)$，使得$\phi_B^{\lambda}(\mathbf{u}_1, \mathbf{u}_1)\neq 0$，考虑$E_R(\lambda)$中与$\mathbf{u}_1$关于$B$正交的向量的集合：
    $$ W_1 = \{\mathbf{w}\in E_R(\lambda)\, \mid \, [\mathbf{u}_1]^{\tr}B[\mathbf{w}] = 0\}. $$
    可以证明：$E_R(\lambda) = \langle\mathbf{u}_1\rangle \oplus W_1$：\\
    首先按定义$\langle\mathbf{u}_1\rangle\cap W_1 = \{\mathbf{0}\}$，其次：
    $$ \forall \mathbf{x} \, \in \, E_R(\lambda), \quad \mathbf{x} = \frac{\phi_B^{\lambda}(\mathbf{u}_1,\mathbf{x})}{\phi_B^{\lambda}(\mathbf{u}_1,\mathbf{u}_1)}\mathbf{u}_1 + \mathbf{x} -\frac{\phi_B^{\lambda}(\mathbf{u}_1,\mathbf{x})}{\phi_B^{\lambda}(\mathbf{u}_1,\mathbf{u}_1)}\mathbf{u}_1 = c\mathbf{u}_1 + \mathbf{w}_1.
    $$
    其中$c = \frac{\phi_B^{\lambda}(\mathbf{u}_1,\mathbf{x})}{\phi_B^{\lambda}(\mathbf{u}_1,\mathbf{u}_1)}$为系数，$\mathbf{w}_1 = \mathbf{x} - \frac{\phi_B(\mathbf{u}_1,\mathbf{x})}{\phi_B(\mathbf{u}_1,\mathbf{u}_1)}\mathbf{u}_1$满足：
    \begin{align*}
        \phi_B^{\lambda}(\mathbf{u}_1, \mathbf{w}_1) &= \phi_B^{\lambda}(\mathbf{u}_1, \mathbf{x}) - \frac{\phi_B^{\lambda}(\mathbf{u}_1,\mathbf{x})}{\phi_B^{\lambda}(\mathbf{u}_1,\mathbf{u}_1)} \phi_B^{\lambda}(\mathbf{u}_1,\mathbf{u}_1) = 0.
    \end{align*}
    这说明$\mathbf{w}_1\in \, W_1$。于是$E_R(\lambda) = \langle\mathbf{u}_1\rangle + W_1 = \langle\mathbf{u}_1\rangle \oplus W_1$。
    $W_1$是$E_R(\lambda)$的子空间，如果维度大于$0$，重复以上思路可以再找出$\mathbf{u}_2$，并构造出$W_1 = \langle\mathbf{u}_2\rangle \oplus W_2$。以此类推，可以构造出$E_R(\lambda)$的一组基底，两两关于$B$正交。

    对$R$的每个本征空间都做同样的操作，就得到整个空间的一组基底$\mathcal{B} = (\mathbf{e}_1, \mathbf{e}_2, \cdots, \mathbf{e}_n)$，关于$B$两两正交。因为首先不同本征空间的基向量关于$B$正交，其次按照我们构造的方式，同本征空间的基向量也两两关于$B$正交。于是我们看以这些基向量为列向量在$\mathbb{R}^n$上的矩阵：
    $$ P = \begin{bmatrix}
        \mathbf{e}_1, \mathbf{e}_2, \cdots, \mathbf{e}_n
    \end{bmatrix}.$$
    它就是能同时对角化$A,B$的合同变换矩阵。首先它满足：
    $$ P^{\tr} BP = D_B.$$
    其中$D_B$是对角矩阵。这是因为$D_B$第$i$行第$j$列的元素就是$[\mathbf{e}_i]^{\tr}B[\mathbf{e}_j]$，$i\neq j$时它等于$0$。

    下面再证明$P^{\tr}AP$也是对角矩阵，这样$A,B$就同时合同对角化了。

    计算$RP$。按定义$\mathbf{e}_i$是$T_R$的本征向量（来自$T_R$的本征空间），所以：
    $$ R[\mathbf{e}_i] = \lambda_i [\mathbf{e}_i],$$
    其中$\lambda_i$是$\mathbf{e}_i$对应的本征值。因此，$RP$的第$i$列是列向量$\lambda_i [\mathbf{e}_i]$。考虑主对角线元素是$\lambda_1,\lambda_2, \cdots, \lambda_n$的对角矩阵$D$，则$PD$的第$i$列也是列向量$\lambda_i [\mathbf{e}_i]$。也就是说
    $$ RP = PD. $$
    于是$ P^{\tr}AP = P^{\tr}BRP = P^{\tr}BPD = D_BD$是对角矩阵的乘积，从而是对角矩阵。
\end{proof}

\end{appendix}

%%%%%%%%%%%%%%%%%%%%%%%%%%%%%%%%%%%
%%%       end of document       %%%
%%%%%%%%%%%%%%%%%%%%%%%%%%%%%%%%%%%


\end{document}